\documentclass[11pt]{article}

\usepackage[letterpaper,margin=1in]{geometry}
\usepackage[T1]{fontenc}
\usepackage[utf8]{inputenc}
\usepackage{lmodern}
\usepackage{microtype}
\usepackage{amsmath,amssymb,amsthm,mathtools,bm,mathrsfs}
\usepackage{booktabs,longtable,array,tabularx}
\usepackage{enumitem}
\usepackage{xcolor}
\usepackage[hidelinks]{hyperref}
\usepackage[nameinlink,noabbrev]{cleveref}

\allowdisplaybreaks
\setlength{\parindent}{0pt}
\setlength{\parskip}{0.55em}
\setlist[enumerate]{itemsep=0.25em,topsep=0.35em}
\setlist[itemize]{itemsep=0.2em,topsep=0.35em}
\renewcommand{\arraystretch}{1.12}

\newtheorem{theorem}{Theorem}[section]
\newtheorem{lemma}[theorem]{Lemma}
\newtheorem{proposition}[theorem]{Proposition}
\newtheorem{corollary}[theorem]{Corollary}
\newtheorem{countertheorem}[theorem]{Countertheorem}
\newtheorem{openproblem}[theorem]{Open problem}
\theoremstyle{definition}
\newtheorem{definition}[theorem]{Definition}
\newtheorem{construction}[theorem]{Construction}
\newtheorem{programme}[theorem]{Programme}
\newtheorem{request}[theorem]{Requested packet}
\theoremstyle{remark}
\newtheorem{remark}[theorem]{Remark}
\newtheorem{assessment}[theorem]{Assessment}
\newtheorem{firewall}[theorem]{Firewall}

\newcommand{\R}{\mathbb R}
\newcommand{\C}{\mathbb C}
\newcommand{\N}{\mathbb N}
\newcommand{\Z}{\mathbb Z}
\newcommand{\T}{\mathbb T}
\newcommand{\one}{\mathbf 1}
\newcommand{\SU}{\mathrm{SU}}
\newcommand{\Tr}{\operatorname{Tr}}
\newcommand{\tr}{\operatorname{tr}}
\newcommand{\Ad}{\operatorname{Ad}}
\newcommand{\Ran}{\operatorname{Ran}}
\newcommand{\Ker}{\operatorname{Ker}}
\newcommand{\Span}{\operatorname{span}}
\newcommand{\spec}{\operatorname{spec}}
\newcommand{\Cov}{\operatorname{Cov}}
\newcommand{\Var}{\operatorname{Var}}
\newcommand{\dist}{\operatorname{dist}}
\newcommand{\dd}{\,\mathrm d}
\newcommand{\norm}[1]{\left\lVert#1\right\rVert}
\newcommand{\abs}[1]{\left\lvert#1\right\rvert}
\newcommand{\ip}[2]{\left\langle#1,#2\right\rangle}
\newcommand{\status}[1]{\textsc{#1}}
\newcommand{\gC}{\mathfrak g_{\C}}
\newcommand{\Wtwo}{\mathsf W_2}
\newcommand{\Wfour}{\mathsf W_4}
\newcommand{\omg}{\omega}
\newcommand{\ellone}[1]{\left\lVert#1\right\rVert_1}

\definecolor{deepgreen}{RGB}{28,110,70}
\definecolor{deepred}{RGB}{150,45,45}
\definecolor{deepblue}{RGB}{30,70,145}
\newcommand{\pass}{\textcolor{deepgreen}{\textbf{PASS}}}
\newcommand{\obstruction}{\textcolor{deepred}{\textbf{TYPED OBSTRUCTION}}}
\newcommand{\stopstatus}{\textcolor{deepblue}{\textbf{STOP}}}

\hypersetup{
  pdftitle={Renewal Yang--Mills Reduced-Fibre Wilson Shell Symbol and Local-Vacuum Programme V1},
  pdfsubject={The first central-vacuum axial mode, Ward-relative displacement tails, and an exact Fejer--Riesz shell-margin certificate}
}

\begin{document}

\begin{center}
{\LARGE\bfseries Renewal Yang--Mills Reduced-Fibre Wilson Shell Symbol\\
and Local-Vacuum Programme}\\[0.45em]
{\Large\bfseries Version V1}\\[0.35em]
{\large The first central-vacuum axial card, a Ward-relative kernel tail,\\
and an exact Fej\'er--Riesz margin certificate}\\[0.35em]
{\normalsize 2 September 2026}
\end{center}

\begin{abstract}
This is the first file in a new append-only Yang--Mills programme.  It takes
as an inherited theorem layer the exact dyadic Wilson fibre, the relative
internal tree gauge, the reduced-fibre conditional law, the
Fourier--stabilizer decomposition, the three-query score--Hessian identity,
the connected score structure factor, the local-curvature factorization, and
the pure momentum--colour failure classifier already proved in the attached
Yang--Mills series.  Those results are not reproduced here.

The open object inherited from that series is the actual selected-shell
symbol
\[
 \widehat H_{r,\sigma}^{\rm hi}(k)
 =\mathbb E K_{r,k,\sigma}
  -\operatorname{Cov}(j_{r,k,\sigma}),
\]
compared with the flat Wilson Schur reference on the same physical mode.  The
present version attacks the first local-vacuum card rather than enlarging the
ambient shell machinery.  On the identity central representative and an exactly
axial-little-group-covariant response-matched local packet, the first nonzero
axial Hodge block
is an irreducible module for the axial little group times the adjoint colour
group.  Its deterministic Hessian, connected score covariance, physical
Hessian, and flat reference are therefore each scalar.  One fixed real
transverse colour polarization supplies the complete block; the nonzero
momentum makes the mean-score query vanish.

The second advance is a tail theorem adapted to the Ward zero.  A naive
absolute Fourier-tail bound is not relative to the vanishing curvature
symbol and becomes useless near zero lattice angle.  After the contact term
is repaired so that both head and tail separately annihilate the harmonic
mode, the omitted tail is bounded by
\(
 \omega(q)\tau_R
\), where \(\omega(q)=2(1-\cos q)\) and \(\tau_R\) is a weighted second
moment of the signed displacement kernel.  The finite head is divisible by
\(\omega\); its quotient is an explicit real trigonometric polynomial.
Fej\'er--Riesz factorization turns positivity of that quotient into a finite
positive-semidefinite Gram feasibility problem.  Combining this Gram with
rigorous head-coefficient enclosures and the second-moment tail yields an
exact three-way verdict: \pass, \obstruction, or \stopstatus.  Failure is
returned at an actual axial momentum, while a zero-row-sum failure is
separated from a genuine topological or harmonic survivor.

No numerical Wilson expectation is invented in this file.  The actual scalar
head coefficients and their tail stock remain physical acquisition rows.
No weak-facing-root escape, continuum vacuum uniqueness,
Osterwalder--Schrader mass, or Yang--Mills mass gap is claimed.
\end{abstract}

\tableofcontents
\clearpage

% =============================================================================
\section{Context, inherited theorem layer, and the V1 target}
\label{sec:YMA1-context}
% =============================================================================

\subsection{Append-only convention and attached source layer}

The programme begins with the file
\begin{center}
\path{Renewal_Yang_Mills_Reduced_Fibre_Wilson_Shell_Symbol_}\\[-0.25em]
\path{and_Local_Vacuum_Programme_V1.tex}.
\end{center}
At every later iteration the complete preceding source is to be preserved
verbatim up to its sole final \verb|\end{document}|; new material is appended
immediately before that command and the version number is increased by one.

The source packet used to determine the nonduplicating starting point is
frozen by the following SHA--256 values:
{\footnotesize
\begin{description}[leftmargin=5.0cm,style=nextline]
\item[Programme architecture]
\texttt{44ccab676db008dbd9460f1518d432ef0eb0ea08a4599cf1e3bac5155f72170b}.
\item[Renewal Grand Tensor manuscript]
\texttt{ada139e74cc525a13a5dce3d7ca26f3ce2c9286a922d1df9be739c726bc3f989}.
\item[Dedicated Yang--Mills V39]
\texttt{361efa298325e48b6961b2d5fd7655519d18121923525a692a5e5496c0b88591}.
\item[Wilson boundary/cofinal V11]
\texttt{3feb6eb0fac7652cf1e6c73037c4b57f0c56b01fa51ad7edfd5ae22435b0dfaa}.
\end{description}
}

\begin{firewall}[Imported results are not silently reproved]
\label{fire:YMA1-no-redo}
The attached Dedicated Yang--Mills V39 series already proves the global and
reduced dyadic fibre constructions, the exact score--Hessian identity, the
modewise direct-loss theorem, the connected structure-factor identity, local
curvature factorization, the absolute signed-kernel head/tail bound, the
cofinal pure-mode classifier, and the quadratic-margin plus fused-third-jet
local-convexity handoff.  The boundary/cofinal V11 series has a separate
boundary-sector mandate.  The proofs below address only the additional
central-vacuum scalarization, Ward-relative tail, exact univariate
positivity certificate, and its cutoff transport.  None of the imported
proofs is copied into this file.
\end{firewall}

\subsection{Inherited packet and the first genuinely open row}

The following finite statements are treated as established inputs.
\begin{enumerate}[label=\textup{(I\arabic*)},leftmargin=4em]
\item The exact dyadic block map admits a global product-Haar fibre and an
      exact quotient by internal fine gauge transformations.
\item Gauge-invariant shell integrals and all coarse jets may be computed on
      one reduced compact fibre with no gauge-volume factor.
\item At a commuting flat background, translation and toron-stabilizer
      symmetry decompose the response-matched Hessian into finite
      momentum--colour symbols.
\item On each retained block,
      \[
       \widehat H=\mathbb E K-\Cov(j),
      \]
      and a failed direct floor is attained by one pure mode.
\item The local contractible packet factors through the covariant curvature
      differential and annihilates local harmonic one-forms; wrapping and
      toron-sensitive packets remain separate.
\item A signed displacement kernel with an exponential row stock has a finite
      Fourier head and an absolute tail bound.
\item Adjacent symbols have exact transported physical-Hessian and reference
      defects.
\end{enumerate}

The missing row is not another shell representation.  It is the physical
value of the first reduced symbol and, cofinally, a quantitative estimate
which is uniform in the correct gauge-relative normalization:
\begin{equation}
 \boxed{
 \begin{gathered}
 \text{actual reduced-fibre mode card}
 \longrightarrow
 \text{uniform symbol margin}\\
 \Downarrow\\[-0.2em]
 \text{higher jets and transport}
 \longrightarrow
 \text{local-vacuum contraction}
 \end{gathered}}
 \label{eq:YMA1-programme-arrow}
\end{equation}

\begin{programme}[V1 acquisition target]
\label{prog:YMA1-target}
On the first escaping selected shell for which the central local-vacuum and
hypercubic symmetry audits pass, freeze the first nonzero axial Hodge orbit,
one real transverse direction, and one predeclared unit colour vector.  The
V1 target is to reduce that whole orbit to one scalar physical card and to
supply a terminating rigorous certificate for its margin from a finite
translated-pair head and a Ward-relative tail.
\end{programme}

\begin{firewall}[No trajectory or mass promotion]
\label{fire:YMA1-no-mass}
The central-vacuum card is conditional on the selected weak-facing shell
entering that patch.  A bounded selected-root subsequence remains on the
inherited cylindrical-Haar/contact branch.  A positive ultraviolet shell
symbol is not a transfer gap, a vacuum uniqueness theorem, an
Osterwalder--Schrader lifetime, or a continuum Yang--Mills mass.
\end{firewall}

% =============================================================================
\section{The frozen central-vacuum axial card}
\label{sec:YMA1-card}
% =============================================================================

Fix one selected shell and suppress the shell index temporarily.  Let
\begin{equation}
 \Lambda_M=(\Z/M\Z)^4,
 \qquad M\ge4,
 \label{eq:YMA1-lattice}
\end{equation}
be the coarse periodic lattice.  Let
\(
 G=\SU(3)
\),
\(
 \mathfrak g=\mathfrak{su}(3)
\), and equip \(\mathfrak g\) with a fixed positive multiple of
\(-\Tr(XY)\).  Its complexification is
\(
 \gC\cong\mathfrak{sl}_3(\C)
\).

\begin{definition}[Central flat backgrounds and the local-vacuum representative]
\label{def:YMA1-central-background}
A central flat background is a constant coarse link field
\begin{equation}
 W^z_{x,\mu}=z_\mu,
 \qquad z_\mu\in Z(\SU(3)).
 \label{eq:YMA1-central-background}
\end{equation}
Every contractible plaquette has holonomy one, and the adjoint covariant
translation phases vanish.  The distinguished fully hypercubic representative
is
\begin{equation}
 \boxed{W^0_{x,\mu}=I.}
 \label{eq:YMA1-trivial-central-background}
\end{equation}
Centre-twist and wrapping data at a nontrivial \(W^z\) are not erased; they are
retained in the separate inherited topological packet unless an exact
same-carrier twist transport has been supplied.
\end{definition}

\begin{lemma}[The fully hypercubic central representative is unique]
\label{lem:YMA1-central-fixed}
Among the constant central \(\SU(3)\) backgrounds
\eqref{eq:YMA1-central-background}, the only one fixed by every coordinate
reflection and permutation in \(B_4\) is \(W^0\).
\end{lemma}

\begin{proof}
Reflection of the \(\mu\)-th coordinate reverses the corresponding oriented
link and sends \(z_\mu\) to \(z_\mu^{-1}\).  A fixed background therefore
satisfies \(z_\mu=z_\mu^{-1}\), or \(z_\mu^2=I\), for each \(\mu\).  The centre
of \(\SU(3)\) is cyclic of order three, so it has no nontrivial element of
order two; hence every \(z_\mu=I\).  Conversely, \(W^0\) is plainly fixed by
all signed coordinate permutations.
\end{proof}

Let \(H^{\rm loc}\) be the Hessian of the fully response-matched
\emph{local contractible} effective action at \(W^0\), after the inherited
rooted gauge, Hodge, resonance, and buffered-high projections.  Let
\(S_0^{\rm loc}\) be the flat Wilson Schur reference on the same carrier.
Write
\begin{equation}
 \widehat H^{\rm loc}(k),
 \qquad
 \widehat S_0^{\rm loc}(k)
 \label{eq:YMA1-symbols}
\end{equation}
for their Fourier symbols.

\begin{definition}[Exact central-axial audit]
\label{def:YMA1-symmetry-audit}
The selected local card is called \emph{central-axial} when the following
finite identities hold before any eigenvalue is inspected.
\begin{enumerate}[label=\textup{(C\arabic*)},leftmargin=3.8em]
\item The background is the fully hypercubic representative
      \eqref{eq:YMA1-trivial-central-background}.  A nontrivial central toron
      is not silently identified with it.
\item The exact response-matched local action, its reduced-fibre law, and the
      flat reference are invariant under coarse translations, the axial
      little group
      \[
       B_3^{(1)}=(\Z_2)^3\rtimes S_3,
      \]
      acting by signed permutations of the three directions transverse to
      the first axis, and constant global gauge conjugation.
\item The buffered-high projection is defined by axial-little-group-invariant
      Hodge, resonance, and physical-momentum criteria and commutes with those
      symmetries.
\item The local packet annihilates the harmonic one-form carrier.  Any
      nonzero harmonic action is assigned to the separately retained
      topological packet.
\item Full \(B_4\) covariance is required only to transport the proved first-axis
      card to all coordinate axes.  Failure of that stronger row does not
      invalidate the first-axis scalarization.
\end{enumerate}
A failure of one of the first four identities is a typed symmetry, patch,
projection, or local/topological-incidence obstruction for the first-axis
card.  Failure only of the final \(B_4\) row blocks direction-orbit propagation.
The physical card is not replaced by its group average.
\end{definition}

Let
\begin{equation}
 q_n=\frac{2\pi n}{M},
 \qquad
 k_n=(n,0,0,0),
 \qquad
 1\le n\le\left\lfloor\frac M2\right\rfloor.
 \label{eq:YMA1-axial-momentum}
\end{equation}
The first axial index \(n_*\) is the least predeclared \(n\) for which the
buffered-high projection retains the nonzero transverse Hodge block.  The
choice is made without reading the shell Hessian.

Put
\begin{equation}
 V_\perp=\Span_{\C}\{e_2,e_3,e_4\}\subset\C^4,
 \qquad
 \mathcal E_{n}^{\rm ax}=V_\perp\otimes\gC.
 \label{eq:YMA1-transverse-module}
\end{equation}
The longitudinal line \(\C e_1\otimes\gC\) is the nonzero-momentum exact
gauge direction and is removed by the Hodge projection.

% =============================================================================
\section{Exact scalarization of the first momentum--colour block}
\label{sec:YMA1-scalarization}
% =============================================================================

\begin{lemma}[Irreducibility of the axial transverse little-group module]
\label{lem:YMA1-B3-irreducible}
Let
\(
 B_3=(\Z_2)^3\rtimes S_3
\)
act on \(V_\perp\cong\C^3\) by signed permutations of
\(e_2,e_3,e_4\).  This representation is irreducible.
\end{lemma}

\begin{proof}
Let \(W\subseteq\C^3\) be a nonzero invariant subspace and choose
\(v=(v_2,v_3,v_4)\in W\setminus\{0\}\).  The coordinate projection onto
\(\C e_j\) is a linear combination of the eight diagonal sign changes:
\[
 P_j=\frac18\sum_{\varepsilon\in\{\pm1\}^3}
       \varepsilon_j\,\operatorname{diag}(\varepsilon_2,
       \varepsilon_3,\varepsilon_4).
\]
Hence \(v_je_j=P_jv\in W\).  Some \(v_j\ne0\), so \(e_j\in W\).
Permutation invariance then gives all three coordinate vectors, and
\(W=\C^3\).
\end{proof}

\begin{lemma}[Irreducibility of the complexified adjoint colour module]
\label{lem:YMA1-adjoint-irreducible}
The adjoint representation of \(\SU(3)\) on
\(\gC\cong\mathfrak{sl}_3(\C)\) is irreducible.
\end{lemma}

\begin{proof}
A complex subspace invariant under \(\Ad(\SU(3))\) is invariant under the
differentiated adjoint action of \(\mathfrak{su}(3)\), and hence, after
complexification, under \(\operatorname{ad}(\mathfrak{sl}_3(\C))\).  It is
therefore an ideal of the simple Lie algebra \(\mathfrak{sl}_3(\C)\).  The
only possibilities are zero and the whole algebra.
\end{proof}

\begin{theorem}[Central-vacuum axial scalarization]
\label{thm:YMA1-axial-scalarization}
Assume the central-axial audit of
\Cref{def:YMA1-symmetry-audit}.  For every retained nonzero axial momentum
\(k_n\), there are real numbers
\begin{equation}
 d_n\in\R,
 \qquad c_n\ge0,
 \qquad h_n=d_n-c_n,
 \qquad s_n>0,
 \label{eq:YMA1-four-scalars}
\end{equation}
such that on \(\mathcal E_n^{\rm ax}\),
\begin{equation}
 \boxed{
 \begin{aligned}
 \mathbb E\widehat K(k_n)&=d_n I,\\
 \widehat{\mathcal C}(k_n)&=c_n I,\\
 \widehat H^{\rm loc}(k_n)&=h_n I,\\
 \widehat S_0^{\rm loc}(k_n)&=s_n I.
 \end{aligned}}
 \label{eq:YMA1-scalar-blocks}
\end{equation}
Here \(\widehat{\mathcal C}\) is the connected score structure factor.
The retained projection is either zero or the identity on the whole
\(24\)-dimensional complex block \(V_\perp\otimes\gC\).

Consequently the complete normalized direct margin on that block is the one
scalar
\begin{equation}
 \boxed{
 \rho_n^{\rm ax}=\frac{h_n}{s_n}=\frac{d_n-c_n}{s_n}.}
 \label{eq:YMA1-axial-margin}
\end{equation}
There is no hidden polarization eigenproblem inside the first axial block.
\end{theorem}

\begin{proof}
The subgroup of the hypercubic little group fixing \(k_n\) contains the
signed permutations \(B_3\) of the transverse coordinates.  The trivial central
background is fixed by constant global gauge conjugation, so the colour
factor carries the adjoint \(\SU(3)\) action.  The product representation of
\(B_3\times\SU(3)\) on
\(V_\perp\otimes\gC\) is irreducible.  Indeed, by
\Cref{lem:YMA1-B3-irreducible,lem:YMA1-adjoint-irreducible} and the
finite-dimensional complex Schur lemma, an operator commuting with
\(B_3\otimes I\) lies in \(I\otimes\operatorname{End}(\gC)\), and
commutation also with \(I\otimes\Ad(\SU(3))\) reduces this space to the
scalars.  Complete reducibility then leaves only one irreducible summand.

The physical Hessian and flat reference commute with this product action by
the exact symmetry audit.  The reduced-fibre deterministic Hessian mean and
the connected score covariance commute separately: a symmetry acts on the
reduced coordinates by a measure-preserving reparametrization and acts
unitarily on the coarse tangent, so change of variables in the expectation
intertwines each operator.  Schur's lemma therefore gives the four scalar
forms in \eqref{eq:YMA1-scalar-blocks}.  Positivity of a covariance gives
\(c_n\ge0\), and the retained flat reference gives \(s_n>0\).  The imported
score--Hessian identity gives \(h_n=d_n-c_n\).

The buffered-high projection is an orthogonal projection commuting with the
same irreducible product action.  Schur's lemma makes it a scalar projection,
therefore either zero or the identity.  The ratio
\eqref{eq:YMA1-axial-margin} follows.
\end{proof}

\begin{corollary}[One polarization and no mean-score query]
\label{cor:YMA1-one-polarization}
Fix before acquisition a unit colour vector
\(X_*\in\mathfrak g\) and the transverse direction \(e_2\).  Let
\(v_{n,*}\) be the normalized complex plane wave of momentum \(k_n\) with
polarization \(e_2\otimes X_*\).  Then
\begin{equation}
 \boxed{
 d_n=\mathbb E\,K[v_{n,*},v_{n,*}],
 \qquad
 c_n=\mathbb E\abs{j[v_{n,*}]}^2,
 \qquad
 \mathbb E j[v_{n,*}]=0.}
 \label{eq:YMA1-one-polarization-card}
\end{equation}
Thus the general three-query mode card reduces on this orbit to one
deterministic second-jet query and one diagonal connected score query; the
mean score vanishes exactly.

For a real physical implementation, let \(v_{n,c}\) and \(v_{n,s}\) be the
normalized cosine and sine waves with the same transverse colour
polarization.  When \(k_n\not\equiv-k_n\pmod M\),
\begin{equation}
 \boxed{
 \begin{aligned}
 d_n&=\frac12\mathbb E\bigl(
 K[v_{n,c},v_{n,c}]+K[v_{n,s},v_{n,s}]\bigr),\\
 c_n&=\frac12\mathbb E\bigl(
 j[v_{n,c}]^2+j[v_{n,s}]^2\bigr).
 \end{aligned}}
 \label{eq:YMA1-real-card}
\end{equation}
The sine--cosine cross terms vanish.  At a self-conjugate Nyquist momentum,
the surviving real character is used directly.
\end{corollary}

\begin{proof}
For every coarse translation \(y\), covariance of the score and invariance of
the reduced-fibre law give
\[
 \mathbb E j[v_{n,*}]
 =\mathbb E j[\tau_yv_{n,*}]
 =e^{iq_n y_1}\mathbb E j[v_{n,*}].
\]
Because the character is nontrivial, some \(y\) has
\(e^{iq_ny_1}\ne1\), so the mean vanishes.  This proves the last identity
in \eqref{eq:YMA1-one-polarization-card}.  The first
two identities are the scalar block values from
\Cref{thm:YMA1-axial-scalarization}.  Complex Fourier synthesis is the
complexification of the real sine--cosine pair.  Translation invariance makes
the two real diagonal values equal and their cross covariance zero, giving
\eqref{eq:YMA1-real-card}.
\end{proof}

\begin{corollary}[Exact first-orbit verdict]
\label{cor:YMA1-first-orbit-verdict}
For a target margin \(\delta\ge0\) fixed before acquisition,
\begin{equation}
 \widehat H^{\rm loc}(k_n)
 \succeq\delta\widehat S_0^{\rm loc}(k_n)
 \quad\Longleftrightarrow\quad
 d_n-c_n-\delta s_n\ge0.
 \label{eq:YMA1-scalar-verdict}
\end{equation}
If the inequality fails, every nonzero vector in the retained axial block has
the same normalized Rayleigh quotient \(\rho_n^{\rm ax}\).  The obstruction
is the axial orbit itself, not an arbitrarily chosen colour or transverse
polarization.
\end{corollary}

\begin{proof}
By \eqref{eq:YMA1-scalar-blocks}, the difference
\(
 \widehat H^{\rm loc}(k_n)-\delta\widehat S_0^{\rm loc}(k_n)
\)
is the scalar \(d_n-c_n-\delta s_n\) times the identity.  This operator is
positive semidefinite exactly when that scalar is nonnegative.  The same
scalar form gives the quotient \(\rho_n^{\rm ax}\) on every nonzero vector,
which proves the failure statement.
\end{proof}

% =============================================================================
\section{Translated-pair acquisition of the signed axial kernel}
\label{sec:YMA1-real-space-card}
% =============================================================================

The scalar block can be acquired directly in momentum space by
\Cref{cor:YMA1-one-polarization}.  For a uniform all-angle certificate it is
more efficient to acquire a finite displacement head.  The following formula
keeps the deterministic and connected contributions fused with their correct
signs.

Let \(\xi_{x}\) denote the unit coarse-link tangent at the link
\((x,2)\), with colour \(X_*\), transported in the frozen right-trivialized
connection.  Let \(J_x=j[\xi_x]\) be its reduced-fibre score.  Define
\begin{align}
 k(x)&:=\mathbb E\,K[\xi_0,\xi_x],
 \label{eq:YMA1-kernel-k}\\
 c(x)&:=\Cov(J_0,J_x),
 \label{eq:YMA1-kernel-c}\\
 s(x)&:=S_0^{\rm loc}[\xi_0,\xi_x].
 \label{eq:YMA1-kernel-s}
\end{align}
For a predeclared target \(\delta\ge0\), put
\begin{equation}
 \boxed{a_\delta(x):=k(x)-c(x)-\delta s(x).}
 \label{eq:YMA1-signed-displacement}
\end{equation}

\begin{theorem}[Exact translated-pair kernel card]
\label{thm:YMA1-translated-pair-card}
On a central-axial local card, the scalar symbol of
\(
 H^{\rm loc}-\delta S_0^{\rm loc}
\)
on the axial transverse block is
\begin{equation}
 \boxed{
 \lambda_\delta(q)
 =\sum_{x\in\Lambda_M}a_\delta(x)e^{-iqx_1}.}
 \label{eq:YMA1-kernel-Fourier}
\end{equation}
The kernel is real and reflection even in the axial coordinate.  Each
coefficient in a finite displacement head is populated by exactly three
same-law translated-pair quantities: the deterministic cross jet \(k(x)\),
the connected score cross moment \(c(x)\), and the known reference
coefficient \(s(x)\).

At fixed displacement radius, hypercubic stabilizer orbits reduce the number
of distinct translated-pair cards to a number depending on the radius but
not on the lattice volume.
\end{theorem}

\begin{proof}
Translation invariance makes each of the three bilinear forms a convolution
operator.  Fourier transformation converts convolution to multiplication,
and evaluation on the fixed transverse colour polarization gives
\eqref{eq:YMA1-kernel-Fourier}.  Reflection in the first coordinate sends
\(x_1\) to \(-x_1\) and preserves the real bilinear card, so the scalar kernel
is real and even.  The displayed definition of \(a_\delta\) is the exact
score--Hessian sign combination; no positive component has been separated or
recharged.  A finite hypercubic group has finitely many orbits in every fixed
displacement ball, proving the last statement.
\end{proof}

\begin{firewall}[The covariance is not a second source price]
\label{fire:YMA1-no-double-charge}
The quantity \(c(x)\) is the connected score covariance already subtracted in
the logarithmic Hessian of the one occurring reduced-fibre Wilson law.  A
boundary, chart, pattern, mediator, or graph disintegration may be used later
to estimate the same coefficient, but its conditional pieces are not added
again to \eqref{eq:YMA1-signed-displacement}.
\end{firewall}

% =============================================================================
\section{Ward-compatible contact repair and relative tail control}
\label{sec:YMA1-Ward-tail}
% =============================================================================

The inherited absolute tail estimate controls
\(
 \abs{\lambda(q)-\lambda^{\le R}(q)}
\)
uniformly.  For a curvature symbol this is not the right quantity near
\(q=0\), because both the physical symbol and its flat reference vanish
quadratically.  We now preserve that Ward zero separately in the head and
tail.

The next theorem is stated on \(\Z^4\).  On a finite torus one uses shortest
periodic representatives and a radius below half the side length; the proof
is identical for the periodized kernel.

The next proposition records that the contact repair is forced by the local
Ward identity; it is not an adjustable counterterm.

\begin{proposition}[Ward row sum and the unique head--tail contact split]
\label{prop:YMA1-Ward-contact}
Let \(a_\delta\) be the translated-pair kernel of
\eqref{eq:YMA1-signed-displacement}.  On the local contractible branch,
\begin{equation}
 \boxed{\sum_x a_\delta(x)=0.}
 \label{eq:YMA1-physical-row-sum}
\end{equation}
Consequently, after the off-contact coefficients have been acquired, the
contact coefficient is fixed uniquely by
\begin{equation}
 \boxed{a_\delta(0)=-\sum_{x\ne0}a_\delta(x).}
 \label{eq:YMA1-contact-fixed}
\end{equation}
For a radius \(R\), define two repaired kernels by retaining the original
off-contact coefficients on, respectively,
\(0<\ellone{x}\le R\) and \(\ellone{x}>R\), setting every other nonzero
coefficient to zero, and setting their contact coefficients to
\begin{align}
 a_{\delta,R}^{\rm hd}(0)
 &:=-\sum_{0<\ellone{x}\le R}a_\delta(x),
 \label{eq:YMA1-head-contact}\\
 a_{\delta,R}^{\rm tl}(0)
 &:=-\sum_{\ellone{x}>R}a_\delta(x).
 \label{eq:YMA1-tail-contact}
\end{align}
Then both repaired kernels have zero row sum and
\begin{equation}
 \boxed{a_\delta=a_{\delta,R}^{\rm hd}+a_{\delta,R}^{\rm tl}.}
 \label{eq:YMA1-contact-kernel-split}
\end{equation}
Thus neither the full contact nor either repaired contact is an independent
physical query once the local Ward branch and all corresponding off-contact
rows have been fixed.
\end{proposition}

\begin{proof}
A constant transverse link tangent is the zero-momentum harmonic one-form.
The local contractible physical Hessian and the flat local reference both
annihilate that carrier by the inherited curvature factorization.  Applying
the convolution operator \(H^{\rm loc}-\delta S_0^{\rm loc}\) to the
constant tangent therefore gives
\(
 (\sum_xa_\delta(x))v=0
\)
for every fixed transverse colour vector \(v\), which proves
\eqref{eq:YMA1-physical-row-sum}.  Equation
\eqref{eq:YMA1-contact-fixed} follows immediately.  The two displayed
contact definitions make the two partial row sums vanish.  Their nonzero
coefficients partition those of \(a_\delta\), while the sum of their contact
coefficients equals \(a_\delta(0)\) by
\eqref{eq:YMA1-contact-fixed}; hence
\eqref{eq:YMA1-contact-kernel-split} holds coefficientwise.
\end{proof}

\begin{theorem}[Ward-repaired head and second-moment relative tail]
\label{thm:YMA1-Ward-tail}
Let \(a:\Z^4\to\R\) satisfy
\begin{equation}
 a(x_1,x')=a(-x_1,x'),
 \qquad
 \sum_x\abs{a(x)}(1+x_1^2)<\infty,
 \qquad
 \sum_xa(x)=0.
 \label{eq:YMA1-Ward-kernel-assumptions}
\end{equation}
Define
\begin{equation}
 \lambda(q)=\sum_xa(x)e^{-iqx_1},
 \qquad
 \omg(q)=2(1-\cos q)=4\sin^2(q/2).
 \label{eq:YMA1-omega}
\end{equation}
For an integer \(R\ge1\), define the Fourier symbols of the uniquely
repaired head and tail from \Cref{prop:YMA1-Ward-contact} by
\begin{align}
 \lambda^{\le R}(q)
 &:=\sum_{0<\ellone{x}\le R}
      a(x)\bigl(\cos(qx_1)-1\bigr),
 \label{eq:YMA1-repaired-head}\\
 \lambda^{>R}(q)
 &:=\sum_{\ellone{x}>R}
      a(x)\bigl(\cos(qx_1)-1\bigr).
 \label{eq:YMA1-repaired-tail}
\end{align}
Then
\begin{equation}
 \boxed{
 \lambda=\lambda^{\le R}+\lambda^{>R},
 \qquad
 \lambda^{\le R}(0)=\lambda^{>R}(0)=0.}
 \label{eq:YMA1-exact-Ward-split}
\end{equation}
Moreover, with
\begin{equation}
 \boxed{
 \tau_R^{(2)}
 :=\frac12\sum_{\ellone{x}>R}x_1^2\abs{a(x)},}
 \label{eq:YMA1-second-moment-tail}
\end{equation}
one has the gauge-relative estimate
\begin{equation}
 \boxed{
 \abs{\lambda^{>R}(q)}
 \le\omg(q)\tau_R^{(2)}
 \qquad(-\pi\le q\le\pi).}
 \label{eq:YMA1-relative-tail-bound}
\end{equation}
Thus the omitted tail is controlled relative to the same quadratic Ward
factor as the local curvature symbol.
\end{theorem}

\begin{proof}
Reflection evenness removes the imaginary Fourier part.  The zero-row-sum
identity gives
\[
 \lambda(q)=\sum_xa(x)(\cos(qx_1)-1),
\]
which splits exactly into \eqref{eq:YMA1-repaired-head} and
\eqref{eq:YMA1-repaired-tail}.  Hence each piece vanishes at zero.

For every integer \(m\),
\[
 \abs{\sin(mt)}\le\abs{m}\,\abs{\sin t},
\]
which follows by writing \(e^{2imt}-1\) as a geometric sum times
\(e^{2it}-1\).  Therefore
\[
 1-\cos(mq)
 =2\sin^2(mq/2)
 \le m^2\,2\sin^2(q/2)
 =m^2(1-\cos q).
\]
Using this inequality in the absolute value of
\eqref{eq:YMA1-repaired-tail} gives
\[
 \abs{\lambda^{>R}(q)}
 \le(1-\cos q)
    \sum_{\ellone{x}>R}x_1^2\abs{a(x)}
 =\omg(q)\tau_R^{(2)}.
\]
\end{proof}

\begin{corollary}[Exponential row stock gives an exponential relative tail]
\label{cor:YMA1-exponential-relative-tail}
Assume in addition that, for some \(\mu>0\),
\begin{equation}
 \mathfrak M_a(\mu)
 :=\sum_xe^{\mu\ellone{x}}\abs{a(x)}<\infty.
 \label{eq:YMA1-exponential-stock}
\end{equation}
If \(R\ge2/\mu\), then
\begin{equation}
 \boxed{
 \tau_R^{(2)}
 \le\frac12R^2e^{-\mu R}\mathfrak M_a(\mu).}
 \label{eq:YMA1-exponential-second-tail}
\end{equation}
\end{corollary}

\begin{proof}
The function \(t\mapsto t^2e^{-\mu t}\) is decreasing for
\(t\ge2/\mu\).  Hence, on \(\ellone{x}>R\),
\[
 x_1^2e^{-\mu\ellone{x}}
 \le\ellone{x}^2e^{-\mu\ellone{x}}
 \le R^2e^{-\mu R}.
\]
Multiplication by \(e^{\mu\ellone{x}}\abs{a(x)}\), followed by summation
and division by two, proves the claim.
\end{proof}

\begin{remark}[Why the cutoff is a full displacement ball]
The condition \(\ellone{x}\le R\) makes the translated-pair head finite
independently of the lattice volume.  A longitudinal truncation
\(\abs{x_1}\le R\) obeys the same one-dimensional Ward estimate, but on an
infinite transverse lattice it is not a finite acquisition packet unless a
separate transverse support or tail theorem has already been supplied.
\end{remark}

\subsection{The low-angle stiffness coefficient}

For \(n\ge1\), aggregate the reflected displacement coefficients by
\begin{equation}
 \alpha_n:=\sum_{x':\,(n,x')\in\Z^4}a(n,x').
 \label{eq:YMA1-alpha-n}
\end{equation}
Then
\begin{equation}
 \lambda(q)=2\sum_{n\ge1}\alpha_n(\cos(nq)-1).
 \label{eq:YMA1-scalar-Laurent}
\end{equation}

\begin{theorem}[Second-moment stiffness and quantitative low-angle control]
\label{thm:YMA1-low-angle-stiffness}
Assume
\(
 \sum_{n\ge1}n^4\abs{\alpha_n}<\infty
\).
For \(q\ne0\), put
\begin{equation}
 g(q):=\frac{\lambda(q)}{\omg(q)}.
 \label{eq:YMA1-Ward-quotient}
\end{equation}
Then \(g\) extends continuously to \(q=0\), with
\begin{equation}
 \boxed{
 g(0)=\kappa_a:=-\sum_{n\ge1}n^2\alpha_n
 =-\frac12\sum_xx_1^2a(x).}
 \label{eq:YMA1-stiffness}
\end{equation}
For \(\abs q\le1\),
\begin{equation}
 \boxed{
 \abs{g(q)-\kappa_a}
 \le\frac{q^2}{4}
       \sum_{n\ge1}n^4\abs{\alpha_n}.}
 \label{eq:YMA1-low-angle-remainder}
\end{equation}
Thus the competition between deterministic curvature and connected score
structure at low lattice angle is exactly the signed second displacement
moment of their already fused kernel.
\end{theorem}

\begin{proof}
For \(q\ne0\),
\begin{equation}
 g(q)
 =-\sum_{n\ge1}\alpha_n
   \frac{1-\cos(nq)}{1-\cos q}.
 \label{eq:YMA1-Dirichlet-ratio}
\end{equation}
The ratio tends to \(n^2\) as \(q\to0\), and it is bounded by \(n^2\) by
the sine inequality used in \Cref{thm:YMA1-Ward-tail}.  Dominated convergence
proves the continuous extension and \eqref{eq:YMA1-stiffness}.

For \(\abs q\le1\), Taylor's theorem gives
\[
 \abs{(1-\cos(nq))-\tfrac12n^2q^2}
 \le\frac1{24}n^4q^4,
 \qquad
 \abs{(1-\cos q)-\tfrac12q^2}
 \le\frac1{24}q^4.
\]
Also \(1-\cos q\ge q^2/3\) on this interval.  Therefore
\[
 \left|
 \frac{1-\cos(nq)}{1-\cos q}-n^2
 \right|
 \le\frac14n^4q^2.
\]
Insert this bound in \eqref{eq:YMA1-Dirichlet-ratio} and sum.
\end{proof}

\begin{corollary}[Low-angle pass and obstruction intervals]
\label{cor:YMA1-low-angle-interval}
Let a finite head have stiffness \(\kappa_R\), fourth stock
\(M_{4,R}\), and relative tail \(\tau_R^{(2)}\).  For \(\abs q\le1\), the
full quotient satisfies
\begin{equation}
 \boxed{
 \kappa_R-\tau_R^{(2)}-\frac{q^2}{4}M_{4,R}
 \le g(q)\le
 \kappa_R+\tau_R^{(2)}+\frac{q^2}{4}M_{4,R}.}
 \label{eq:YMA1-low-angle-interval}
\end{equation}
A positive lower endpoint closes a low-angle window.  A negative upper
endpoint at an actual allowed nonzero lattice angle returns a physical
low-angle shell obstruction.  At a non-grid angle it is a continuous
interpolation dip and is promoted only by the grid-landing result below.  An
interval meeting zero is \stopstatus: a larger head or sharper tail is
required, not another abstract carrier.
\end{corollary}

\begin{proof}
Apply \Cref{thm:YMA1-low-angle-stiffness} to the finite head.  Its quotient
differs from \(\kappa_R\) by at most
\(
 q^2M_{4,R}/4
\).  The Ward-relative tail theorem changes the quotient by at most
\(\tau_R^{(2)}\).  The triangle inequality gives both endpoints in
\eqref{eq:YMA1-low-angle-interval}.  Their sign interpretations follow from
the one-sided enclosure and, for a physical obstruction, evaluation at an
allowed nonzero momentum.
\end{proof}

% =============================================================================
\section{Exact Fej\'er--Riesz certificate for the finite axial head}
\label{sec:YMA1-Fejer}
% =============================================================================

Let a contact-repaired finite head have aggregated coefficients
\(\alpha_1,\ldots,\alpha_R\), so
\begin{equation}
 \lambda_R(q)
 =2\sum_{m=1}^{R}\alpha_m(\cos(mq)-1).
 \label{eq:YMA1-finite-head-symbol}
\end{equation}

\begin{lemma}[Explicit division by the Ward factor]
\label{lem:YMA1-Ward-division}
For \(0\le j\le R-1\), define
\begin{equation}
 \boxed{
 b_j:=-\sum_{m=j+1}^{R}(m-j)\alpha_m.}
 \label{eq:YMA1-b-coefficients}
\end{equation}
Then the real trigonometric polynomial
\begin{equation}
 g_R(q)=b_0+2\sum_{j=1}^{R-1}b_j\cos(jq)
 \label{eq:YMA1-finite-quotient}
\end{equation}
satisfies the exact factorization
\begin{equation}
 \boxed{\lambda_R(q)=\omg(q)g_R(q).}
 \label{eq:YMA1-Ward-factorization}
\end{equation}
Moreover,
\begin{equation}
 \boxed{g_R(0)=-\sum_{m=1}^{R}m^2\alpha_m.}
 \label{eq:YMA1-head-stiffness}
\end{equation}
\end{lemma}

\begin{proof}
Set \(b_R=0\) and extend \(b_{-j}=b_j\).  Direct subtraction of the sums in
\eqref{eq:YMA1-b-coefficients} gives, for \(1\le m\le R\),
\[
 2b_m-b_{m-1}-b_{m+1}=\alpha_m.
\]
The constant coefficient is
\[
 2b_0-2b_1=-2\sum_{m=1}^{R}\alpha_m.
\]
These are precisely the Laurent coefficients of
\((2-z-z^{-1})\sum_{j=-(R-1)}^{R-1}b_jz^j\), while
\(2-z-z^{-1}=\omg(q)\) on \(z=e^{iq}\).  This proves
\eqref{eq:YMA1-Ward-factorization}.  Divide by \(q^2\) and let \(q\to0\), or
sum the coefficients, to obtain \eqref{eq:YMA1-head-stiffness}.
\end{proof}

\begin{theorem}[Finite Gram characterization of a uniform axial head margin]
\label{thm:YMA1-Fejer-Gram}
Let
\begin{equation}
 f(q)=c_0+2\sum_{n=1}^{d}c_n\cos(nq)
 \label{eq:YMA1-real-trig-poly}
\end{equation}
be a real trigonometric polynomial.  Then
\begin{equation}
 \boxed{f(q)\ge0\quad\text{for every }q\in\R}
 \label{eq:YMA1-nonnegative-poly}
\end{equation}
if and only if there is a Hermitian matrix
\(Q\in M_{d+1}(\C)\), \(Q\succeq0\), satisfying
\begin{align}
 \sum_{j=0}^{d}Q_{jj}&=c_0,
 \label{eq:YMA1-Gram-constant}\\
 \sum_{j=0}^{d-n}Q_{j,j+n}&=c_n,
 \qquad 1\le n\le d.
 \label{eq:YMA1-Gram-diagonals}
\end{align}
Equivalently, with \(v_d(z)=(1,z,\ldots,z^d)^T\),
\begin{equation}
 f(q)=v_d(e^{iq})^*Qv_d(e^{iq}).
 \label{eq:YMA1-Gram-polynomial}
\end{equation}
For a prescribed scalar \(\eta\), the condition \(f\ge\eta\) is obtained by
replacing \(c_0\) by \(c_0-\eta\).
\end{theorem}

\begin{proof}
If \(Q\succeq0\), expansion of
\eqref{eq:YMA1-Gram-polynomial} gives the diagonal-sum coefficients in
\eqref{eq:YMA1-Gram-constant}--\eqref{eq:YMA1-Gram-diagonals}, so
\(f\ge0\).

Conversely, assume \(f\ge0\).  If \(f\equiv0\), take \(Q=0\).
Otherwise let \(d_0\le d\) be the effective trigonometric degree and write
\(
 F(z)=\sum_{n=-d_0}^{d_0}c_{\abs n}z^n
\)
and \(P(z)=z^{d_0}F(z)\).  The reciprocal-Hermitian identity
\(
 P(z)=z^{2d_0}\overline{P(1/\overline z)}
\)
forces every root off the unit circle to occur with its
reciprocal conjugate.  Every unit-circle root has even multiplicity, since a
real analytic nonnegative function cannot have a zero of odd order.  Select
one root from each reciprocal-conjugate pair and half of every unit-circle
multiplicity.  After choosing the positive scalar factor, this gives a
polynomial \(p(z)=\sum_{j=0}^{d_0}u_jz^j\) such that
\[
 P(z)=z^{d_0}p(z)\overline{p(1/\overline z)}.
\]
On \(|z|=1\), division by \(z^{d_0}\) yields
\[
 f(q)=\abs{p(e^{iq})}^2.
\]
Pad \(u\) by zeros through degree \(d\), put
\(w=(\overline{u_0},\ldots,\overline{u_d})^T\), and take
\(Q=ww^*\).  Then
\(
 v_d(z)^*Qv_d(z)=\abs{p(z)}^2
\)
on the unit circle.  The final assertion applies the same argument to
\(f-\eta\).
\end{proof}

\subsection{Certified lower and upper bounds for the physical axial margin}

Let \(h(q)\) and \(s(q)\) be respectively the scalar axial symbols of the
physical local Hessian and the flat reference, with \(s(q)>0\) for
\(q\ne0\).  Let \(g_{H,R}\) and \(g_{S,R}\) be the Ward quotients of their
contact-repaired heads.  Suppose rigorous relative-tail bounds are available:
\begin{equation}
 \abs{h(q)-\omg(q)g_{H,R}(q)}
 \le\omg(q)\tau_{H,R},
 \qquad
 \abs{s(q)-\omg(q)g_{S,R}(q)}
 \le\omg(q)\tau_{S,R}.
 \label{eq:YMA1-HS-tail-enclosures}
\end{equation}
For a target \(\delta\ge0\), put
\begin{equation}
 g_{\delta,R}:=g_{H,R}-\delta g_{S,R}.
 \label{eq:YMA1-fused-head-quotient}
\end{equation}
The preferred tail is the directly fused signed stock
\begin{equation}
 \boxed{
 \tau_{\delta,R}^{\rm fus}
 :=\frac12\sum_{\ellone{x}>R}x_1^2
       \abs{k(x)-c(x)-\delta s(x)}.}
 \label{eq:YMA1-fused-tail-stock}
\end{equation}
By \Cref{thm:YMA1-Ward-tail}, it satisfies
\begin{equation}
 \abs{h(q)-\delta s(q)-\omg(q)g_{\delta,R}(q)}
 \le\omg(q)\tau_{\delta,R}^{\rm fus}.
 \label{eq:YMA1-fused-tail-enclosure}
\end{equation}
If \(\tau_{H,R}\) and \(\tau_{S,R}\) are the corresponding
second-moment coefficient stocks for the physical and reference kernels, but
the fused tail coefficients have not been retained, the triangle inequality
gives the admissible majorant
\begin{equation}
 \boxed{
 \tau_{\delta,R}^{\rm fus}
 \le \widehat\tau_{\delta,R}
 :=\tau_{H,R}+\delta\tau_{S,R}.}
 \label{eq:YMA1-separate-tail-majorant}
\end{equation}
The direct fused stock is never replaced by this larger bound when its actual
signed coefficients are available.

\begin{theorem}[PASS--obstruction--STOP axial certificate]
\label{thm:YMA1-three-way-certificate}
Fix a target margin \(\delta\ge0\) before reading the card.  Let
\(\tau_{\delta,R}\) be any certified radius for
\eqref{eq:YMA1-fused-tail-enclosure}, preferably
\(\tau_{\delta,R}^{\rm fus}\), and define
\begin{align}
 f^-_{R,\delta}(q)
 &:=g_{\delta,R}(q)-\tau_{\delta,R},
 \label{eq:YMA1-lower-poly}\\
 f^+_{R,\delta}(q)
 &:=g_{\delta,R}(q)+\tau_{\delta,R}.
 \label{eq:YMA1-upper-poly}
\end{align}
Exactly the following certified outcomes are available.
\begin{enumerate}[label=\textup{(A\arabic*)},leftmargin=3.8em]
\item \pass: if a Gram from \Cref{thm:YMA1-Fejer-Gram} certifies
      \(f^-_{R,\delta}\ge0\) on the unit circle, then
      \begin{equation}
       \boxed{h(q)\ge\delta s(q)\quad\text{for every }q.}
       \label{eq:YMA1-pass-floor}
      \end{equation}
\item \obstruction: if an actual allowed lattice momentum \(q_*\ne0\)
      satisfies
      \begin{equation}
       \boxed{f^+_{R,\delta}(q_*)<0,}
       \label{eq:YMA1-obstruction-test}
      \end{equation}
      then \(h(q_*)<\delta s(q_*)\).  The axial orbit at \(q_*\) is an
      attained physical shell witness.
\item \stopstatus: if neither one-sided certificate holds, the present head
      and tail do not decide the target.  The permitted next operation is to
      enlarge the physical translated-pair head or sharpen its tail stock.
      No positive or negative theorem is inferred.
\end{enumerate}
\end{theorem}

\begin{proof}
From the certified fused enclosure,
\[
 h(q)-\delta s(q)
 \ge\omg(q)f^-_{R,\delta}(q).
\]
This proves the pass branch.  The opposite one-sided estimate is
\[
 h(q)-\delta s(q)
 \le\omg(q)f^+_{R,\delta}(q),
\]
which is strictly negative at the displayed nonzero lattice momentum in the
obstruction branch.  If neither inequality has a certified sign, the data
only enclose zero and no theorem verdict follows.
\end{proof}

\begin{remark}[Fused-tail priority and affine optimization]
For one fixed target \(\delta\), the stock
\(\tau_{\delta,R}^{\rm fus}\) is the cancellation-preserving coefficientwise
certificate furnished by the available signed displacement data and may be
strictly smaller than the
separate majorant.  The majorant is retained only because it is affine in
\(\delta\) and therefore turns optimization over the margin into one SDP.
Using it does not assign a second source price to the covariance; it merely
forgets cancellations in a one-sided analytic error bound.
\end{remark}

\begin{definition}[Continuous interpolating and physical discrete margins]
\label{def:YMA1-two-axial-margins}
Let
\[
 \mathcal Q_M
 :=\left\{
  \left(\frac{2\pi n}{M}\right)_{[-\pi,\pi)}:
  0\le n<M
 \right\}
\]
be the centred physical lattice-angle set.  Define
\begin{align}
 \delta_{\rm cont}^{\rm ax}
 &:=\inf_{q\in[-\pi,\pi]\setminus\{0\}}
       \frac{h(q)}{s(q)},
 \label{eq:YMA1-continuous-axial-margin}\\
 \delta_{M,\rm disc}^{\rm ax}
 &:=\min_{q\in\mathcal Q_M\setminus\{0\}}
       \frac{h(q)}{s(q)}.
 \label{eq:YMA1-discrete-axial-margin}
\end{align}
The first is the margin of the continuous Fourier interpolation and the
second is the actual finite-torus axial margin.  Necessarily
\begin{equation}
 \boxed{\delta_{\rm cont}^{\rm ax}
        \le\delta_{M,\rm disc}^{\rm ax}.}
 \label{eq:YMA1-continuous-below-discrete}
\end{equation}
A continuous-circle Gram pass is therefore a sufficient physical finite-card
certificate.  Failure of that Gram alone is not a finite-torus obstruction;
the obstruction branch of \Cref{thm:YMA1-three-way-certificate} still
requires an allowed lattice momentum.
\end{definition}

\begin{corollary}[Finite semidefinite programme and rigorous margin intervals]
\label{cor:YMA1-SDP-margin}
Assume
\begin{equation}
 \inf_q\bigl(g_{S,R}(q)+\tau_{S,R}\bigr)>0.
 \label{eq:YMA1-reference-upper-positive}
\end{equation}
Use the affine separate-tail majorant
\(\widehat\tau_{\delta,R}=\tau_{H,R}+\delta\tau_{S,R}\) from
\eqref{eq:YMA1-separate-tail-majorant}.  For each predeclared
\(\delta\ge0\), the resulting continuous-circle pass test is a finite
semidefinite feasibility problem: all coefficients are affine in
\(\delta\), while \Cref{thm:YMA1-Fejer-Gram} imposes affine diagonal-sum
equations and \(Q\succeq0\).  If the feasible set is nonempty, it is a closed
bounded interval and its largest element is
\begin{equation}
 \boxed{
 \underline\delta_R
 =\inf_q
  \frac{g_{H,R}(q)-\tau_{H,R}}
       {g_{S,R}(q)+\tau_{S,R}}.}
 \label{eq:YMA1-lower-margin}
\end{equation}
Whenever also
\begin{equation}
 \inf_q\bigl(g_{S,R}(q)-\tau_{S,R}\bigr)>0,
 \label{eq:YMA1-reference-lower-positive}
\end{equation}
define
\begin{align}
 \overline\delta_{R,\rm cont}
 &:=\inf_q
 \frac{g_{H,R}(q)+\tau_{H,R}}
      {g_{S,R}(q)-\tau_{S,R}},
 \label{eq:YMA1-upper-margin-cont}\\
 \overline\delta_{R,M,\rm disc}
 &:=\min_{q\in\mathcal Q_M\setminus\{0\}}
 \frac{g_{H,R}(q)+\tau_{H,R}}
      {g_{S,R}(q)-\tau_{S,R}}.
 \label{eq:YMA1-upper-margin-disc}
\end{align}
Then
\begin{equation}
 \boxed{
 \underline\delta_R
 \le\delta_{\rm cont}^{\rm ax}
 \le\overline\delta_{R,\rm cont},
 \qquad
 \underline\delta_R
 \le\delta_{M,\rm disc}^{\rm ax}
 \le\overline\delta_{R,M,\rm disc}.}
 \label{eq:YMA1-margin-interval}
\end{equation}
If the full Ward quotients \(G_H=h/\omg\) and \(G_S=s/\omg\) extend
continuously, \(\inf G_S>0\), the two head quotients converge uniformly to
\(G_H,G_S\), the two tail radii tend to zero, and
\(\delta_{\rm cont}^{\rm ax}>0\), then for all sufficiently large \(R\) the
feasible set is nonempty and
\begin{equation}
 \boxed{
 \underline\delta_R\longrightarrow\delta_{\rm cont}^{\rm ax},
 \qquad
 \overline\delta_{R,\rm cont}
 \longrightarrow\delta_{\rm cont}^{\rm ax}.}
 \label{eq:YMA1-margin-interval-convergence}
\end{equation}
If in addition the physical angle meshes become dense, then the discrete
margins converge to the same continuous margin.  No monotonicity in \(R\) is
asserted without a nested validated acquisition scheme.
\end{corollary}

\begin{proof}
The set inclusion in \eqref{eq:YMA1-continuous-below-discrete} is immediate.
By \Cref{thm:YMA1-Fejer-Gram}, semidefinite feasibility is equivalent to
\(f^-_{R,\delta}(q)\ge0\) for every \(q\).  Under
\eqref{eq:YMA1-reference-upper-positive}, this is equivalent to
\[
 \delta\le
 \frac{g_{H,R}(q)-\tau_{H,R}}
      {g_{S,R}(q)+\tau_{S,R}}
 \quad\text{for every }q,
\]
which proves \eqref{eq:YMA1-lower-margin}, closedness, boundedness, and
attainment whenever the nonnegative feasible set is nonempty.  Every feasible
value is a lower bound for the continuous margin by
\Cref{thm:YMA1-three-way-certificate}, hence also for the discrete margin.

Under \eqref{eq:YMA1-reference-lower-positive}, the tail enclosures give
pointwise
\[
 \frac{h(q)}{s(q)}
 \le
 \frac{g_{H,R}(q)+\tau_{H,R}}
      {g_{S,R}(q)-\tau_{S,R}}.
\]
Taking the infimum over the circle or the minimum over the physical grid gives
the two upper bounds.  Under the convergence hypotheses, both continuous
bounding functions converge uniformly to \(G_H/G_S\), because the
denominators have one common positive floor for all sufficiently large
\(R\).  Infima of uniformly convergent functions on the compact unit circle
converge.  Positivity of \(\delta_{\rm cont}^{\rm ax}\) makes the lower
feasible set nonempty eventually, proving
\eqref{eq:YMA1-margin-interval-convergence}.  Finally, the minimum of a fixed
continuous function over increasingly dense compact grids converges to its
continuous infimum; the same conclusion holds for uniformly convergent
cutoff-dependent quotient ratios.
\end{proof}

\begin{proposition}[Validated coefficient enclosures]
\label{prop:YMA1-validated-coefficients}
Suppose the finite-head coefficients have rigorous enclosures
\begin{equation}
 \abs{\alpha_m-\widetilde\alpha_m}\le\varepsilon_m,
 \qquad 1\le m\le R.
 \label{eq:YMA1-alpha-enclosures}
\end{equation}
Let \(b_j\) and \(\widetilde b_j\) be obtained from
\eqref{eq:YMA1-b-coefficients}.  Then
\begin{equation}
 \abs{b_j-\widetilde b_j}
 \le
 \eta_j:=\sum_{m=j+1}^{R}(m-j)\varepsilon_m,
 \label{eq:YMA1-b-errors}
\end{equation}
and
\begin{equation}
 \boxed{
 \norm{g_R-\widetilde g_R}_{L^\infty(\T)}
 \le\eta_0+2\sum_{j=1}^{R-1}\eta_j.}
 \label{eq:YMA1-polynomial-error}
\end{equation}
This error is simply added to the relative tail in the pass and obstruction
polynomials.  Therefore certified interval integration, exact finite sums, or
validated numerical enclosures can populate the card without promoting
uncontrolled floating-point output to a theorem.
\end{proposition}

\begin{proof}
Apply the triangle inequality term by term to
\eqref{eq:YMA1-b-coefficients}, then use \(\abs{\cos(jq)}\le1\) in
\eqref{eq:YMA1-finite-quotient}.
\end{proof}

% =============================================================================
\section{The correct adjacent-shell transport norm}
\label{sec:YMA1-transport}
% =============================================================================

The absolute operator norm used for a fixed nonzero mode does not control a
family approaching the Ward zero unless the reference has a uniform positive
floor.  The quotient by \(\omg\) identifies the correct weighted displacement
norm.

For a reflected scalar kernel with coefficients \(\alpha_n\), define
\begin{equation}
 \boxed{
 \norm{\alpha}_{\Wtwo}
 :=\sum_{n\ge1}n^2\abs{\alpha_n}.}
 \label{eq:YMA1-W2-norm}
\end{equation}
When fourth-order low-angle control is required, put
\begin{equation}
 \norm{\alpha}_{\Wfour}
 :=\sum_{n\ge1}n^4\abs{\alpha_n}.
 \label{eq:YMA1-W4-norm}
\end{equation}

\begin{theorem}[Uniform quotient control by the second displacement moment]
\label{thm:YMA1-W2-control}
Let \(g_\alpha\) and \(g_\beta\) be the continuously extended Ward quotients
associated with two reflected zero-row-sum kernels.  Then
\begin{equation}
 \boxed{
 \norm{g_\alpha-g_\beta}_{L^\infty(\T)}
 \le\norm{\alpha-\beta}_{\Wtwo}.}
 \label{eq:YMA1-W2-Lipschitz}
\end{equation}
In particular,
\(
 \norm{g_\alpha}_{\infty}\le\norm{\alpha}_{\Wtwo}
\).
\end{theorem}

\begin{proof}
For \(q\ne0\), formula \eqref{eq:YMA1-Dirichlet-ratio} gives
\[
 g_\alpha(q)-g_\beta(q)
 =-\sum_{n\ge1}(\alpha_n-\beta_n)
   \frac{1-\cos(nq)}{1-\cos q}.
\]
The quotient is nonnegative and at most \(n^2\).  The asserted bound follows
by summation.  Continuity passes it to \(q=0\).
\end{proof}

\begin{theorem}[Summable adjacent-shell transport of the axial margin]
\label{thm:YMA1-cofinal-W2}
Let a cofinal sequence of central-axial selected shells be identified by
the inherited physical-momentum, source, and reference transports.  Let
\(\alpha_r^H\) and \(\alpha_r^S\) be the transported axial displacement
coefficients of the local physical Hessian and flat reference.  Assume
\begin{equation}
 \sum_{r=r_0}^{\infty}
 \left(
  \norm{\alpha_{r+1}^H-\alpha_r^H}_{\Wtwo}
  +\norm{\alpha_{r+1}^S-\alpha_r^S}_{\Wtwo}
 \right)<\infty.
 \label{eq:YMA1-summable-W2-transport}
\end{equation}
Then the Ward quotients \(g_r^H\) and \(g_r^S\) converge uniformly on the
whole unit circle, including \(q=0\).  If
\begin{equation}
 \inf_q\bigl(g_{r_0}^H(q)-\delta g_{r_0}^S(q)\bigr)
 >
 \sum_{r=r_0}^{\infty}
 \left(
  \norm{\alpha_{r+1}^H-\alpha_r^H}_{\Wtwo}
  +\delta\norm{\alpha_{r+1}^S-\alpha_r^S}_{\Wtwo}
 \right),
 \label{eq:YMA1-transport-margin}
\end{equation}
then every later axial card satisfies
\begin{equation}
 \boxed{h_r(q)\ge\delta s_r(q)\quad(q\ne0).}
 \label{eq:YMA1-cofinal-axial-floor}
\end{equation}
The estimate remains relative to the curvature zero and requires no
artificial positive reference floor at \(q=0\).
\end{theorem}

\begin{proof}
Apply \Cref{thm:YMA1-W2-control} to each adjacent pair and sum the resulting
uniform bounds.  The two quotient sequences are uniformly Cauchy.  For the
margin, telescope
\(
 g_r^H-\delta g_r^S
\)
and subtract the total adjacent debit in
\eqref{eq:YMA1-transport-margin}.  Multiplication by \(\omg(q)>0\) for
\(q\ne0\) gives \eqref{eq:YMA1-cofinal-axial-floor}.
\end{proof}

\begin{assessment}[Why \(\Wtwo\) is the relevant transport stock]
\label{ass:YMA1-W2-significance}
The Ward quotient contains the factor
\(
 (1-\cos(nq))/(1-\cos q)
\), whose exact uniform envelope is \(n^2\).  Thus the second displacement
moment is not an arbitrary strengthening of the absolute kernel norm; it is
the operator norm of the scalar axial kernel after division by the local
curvature zero.  A zeroth-moment tail can certify fixed angles but cannot by
itself transport the normalized low-angle margin.
\end{assessment}

% =============================================================================
\section{Ward-zero refinement of the inherited first-orbit classifier}
\label{sec:YMA1-failure}
% =============================================================================

\begin{proposition}[Ward preclassification before the inherited mode alternative]
\label{thm:YMA1-failure-classifier}
Consider a cofinal sequence of first-axis cards.  Before applying the
inherited pure-mode classifier, perform the following ordered test.
\begin{enumerate}[label=\textup{(F\arabic*)},leftmargin=3.8em]
\item If the acquired full packet has a nonzero value at \(q=0\), or
      equivalently its proposed local displacement kernel has nonzero row sum,
      it has failed the local contractible Ward incidence.  The first such row
      is exported as a topological/wrapping, patch-transition, or mixed
      local/topological residual; it is not called a local momentum soft mode.
\item On the zero-row-sum branch, the harmonic angle is removed exactly and
      every physical failure returned by
      \Cref{thm:YMA1-three-way-certificate} occurs at an allowed
      \(q_r\ne0\).  With physical momentum
      \begin{equation}
       p_r:=q_r/a_{c,r},
       \label{eq:YMA1-physical-momentum}
      \end{equation}
      that witness is passed without alteration to the already proved
      finite-physical-momentum versus ultraviolet alternative of the Dedicated
      Yang--Mills series.
\item Failure of the identity background, axial little-group action, retained
      projection, or local/topological split is retained before either
      momentum branch.  Failure only of full \(B_4\) covariance blocks export
      to the other coordinate axes, not the first-axis card.
\end{enumerate}
Thus Version~V1 adds the Ward-zero and direction-orbit preclassification; it
does not reprove the inherited subsequence classifier.  None of these outputs
is an infrared mass statement before physical-time and continuum incidence
are supplied.
\end{proposition}

\begin{proof}
The first statement is
\Cref{prop:YMA1-Ward-contact}: a local contractible curvature symbol must
annihilate the constant transverse harmonic mode.  On the passing branch,
\(q=0\) vanishes identically, while the obstruction clause of
\Cref{thm:YMA1-three-way-certificate} explicitly requires an allowed nonzero
lattice angle.  The definition of \(p_r\) is exactly the physical rescaling
used by the inherited cofinal mode theorem, so no new compactness or
subsequence argument is needed.  The last item is the audit order in
\Cref{def:YMA1-symmetry-audit}.
\end{proof}

\begin{corollary}[A continuous finite-head dip lands on discrete fine grids]
\label{cor:YMA1-grid-landing}
Suppose a limiting head/tail certificate gives a continuous quotient
\(g_\infty\) and
\begin{equation}
 g_\infty(q_*)\le-\eta<0
 \label{eq:YMA1-continuous-dip}
\end{equation}
for some \(q_*\in(0,\pi]\).  If the allowed lattice-angle meshes tend to
zero and the transported quotients converge uniformly, then all sufficiently
fine cards contain an allowed angle \(q_r\) with
\begin{equation}
 g_r(q_r)<-\eta/2.
 \label{eq:YMA1-discrete-dip}
\end{equation}
If \(q_*>0\) is fixed while \(a_{c,r}\to0\), this is an ultraviolet witness.
A finite-physical-momentum witness instead requires a dip at angles of order
\(a_{c,r}\).
\end{corollary}

\begin{proof}
Continuity gives a neighbourhood of \(q_*\) on which
\(g_\infty<-3\eta/4\).  A sufficiently fine momentum grid meets that
neighbourhood, and uniform convergence makes \(g_r<-\eta/2\) there.  The
physical-momentum statements follow from division by \(a_{c,r}\).
\end{proof}

% =============================================================================
\section{Nonduplicating handoff to the inherited higher-jet theorem}
\label{sec:YMA1-local-vacuum}
% =============================================================================

The attached Dedicated Yang--Mills series already proves the general
finite-dimensional implication
\begin{equation}
 \boxed{
 \text{positive supported quadratic margin}
 +\text{ fully fused third-jet stock}
 \Longrightarrow
 \text{local strong convexity}.}
 \label{eq:YMA1-inherited-higher-jet-handoff}
\end{equation}
It also proves the corresponding contraction estimate for the
reference-metric gradient flow.  Repeating that argument here would not
advance YM--A.

\begin{assessment}[Exact interface exported by Version V1]
\label{ass:YMA1-higher-jet-interface}
The new object supplied by this file is a rigorous lower certificate for the
first central-vacuum axial quadratic margin in the same flat Wilson reference
metric used by the inherited higher-jet theorem.  After all remaining
momentum little-group orbits have been acquired and the local/topological
split has passed, the minimum of their certified supported margins is the
quadratic input to the already proved fused third-jet handoff.  Version~V1
does not reopen the graph-leg envelope theorem, its Brascamp--Lieb/mediator
fusion, or its local-contraction proof.
\end{assessment}

\begin{firewall}[No higher-jet substitution]
\label{fire:YMA1-higher-jet-no-redo}
A positive first axial card is not an all-orbit quadratic margin.  Conversely,
the inherited higher-jet theorem does not populate a missing axial symbol.
The order remains: acquire the physical mode cards, take their supported
minimum, acquire the already defined fully fused higher stock, and only then
invoke the inherited local-vacuum result.
\end{firewall}

% =============================================================================
\section{Version V1 terminal theorem, ledger, and next attack}
\label{sec:YMA1-terminal}
% =============================================================================

\begin{theorem}[Version V1 reduced axial-card theorem]
\label{thm:YMA1-terminal}
Relative to the inherited Dedicated Yang--Mills V39 and boundary/cofinal V11
layers, Version~V1 proves the following new statements.
\begin{enumerate}[label=\textup{(E1.\arabic*)},leftmargin=4.8em]
\item Among constant central \(\SU(3)\) backgrounds, full hypercubic
      invariance fixes the identity representative \(W^0\).  On an exactly
      central-axial selected shell at \(W^0\), requiring only the first-axis
      little group, every retained nonzero axial Hodge block is an irreducible
      \(B_3\times\SU(3)\) module.
\item The mean deterministic Hessian, connected score structure factor,
      physical Hessian, and flat reference are each scalar on that block.
\item One predeclared transverse colour polarization determines the complete
      block, and nonzero momentum removes the mean-score query.
\item A finite translated-pair bank reconstructs the signed scalar
      displacement head after deterministic curvature and connected score
      covariance have been fused with their correct signs.
\item The Ward row sum uniquely fixes the repaired head and tail contacts, so
      both pieces annihilate the zero mode.  The omitted signed tail is
      controlled relative to \(\omg(q)=2(1-\cos q)\) by its fused weighted
      second displacement moment; separate stocks are only an optional affine
      majorant.
\item The low-angle normalized stiffness is the signed second moment of the
      fused kernel, with an explicit fourth-moment remainder.
\item The finite head has an explicit quotient by \(\omg\), and
      Fej\'er--Riesz factorization is equivalent to a finite positive Gram
      certificate for its complete continuous-angle margin.
\item Rigorous coefficient enclosures and relative tails give an exact
      \pass--\obstruction--\stopstatus decision and a lower/upper interval for
      the continuous and physical discrete axial symbol margins.
\item The \(\Wtwo\) displacement norm is the exact uniform transport norm for
      the Ward quotient; summable adjacent defects preserve the axial margin
      through the gauge zero.
\item Before the inherited finite-physical/ultraviolet classifier is invoked,
      the first orbit now has an exact Ward/topological and direction-orbit
      preclassification; a continuous head dip is landed on physical grids
      only by a separate density argument.
\item The resulting all-orbit quadratic margin feeds directly into the
      already proved fused-third-jet local-vacuum theorem; that theorem is
      imported rather than reproved.
\end{enumerate}
No actual Wilson scalar value is supplied by these structural reductions.
\end{theorem}

\begin{proof}
The central fixed-point statement in item~\textup{(E1.1)} is
\Cref{lem:YMA1-central-fixed}.  The remaining assertions in
items~\textup{(E1.1)--(E1.3)} are
\Cref{lem:YMA1-B3-irreducible,lem:YMA1-adjoint-irreducible},
\Cref{thm:YMA1-axial-scalarization}, and
\Cref{cor:YMA1-one-polarization}.
Item~\textup{(E1.4)} is \Cref{thm:YMA1-translated-pair-card}.
Items~\textup{(E1.5)--(E1.6)} are
\Cref{prop:YMA1-Ward-contact,thm:YMA1-Ward-tail,thm:YMA1-low-angle-stiffness}.
Items~\textup{(E1.7)--(E1.8)} are
\Cref{lem:YMA1-Ward-division,thm:YMA1-Fejer-Gram},
\Cref{thm:YMA1-three-way-certificate,cor:YMA1-SDP-margin}, and
\Cref{prop:YMA1-validated-coefficients}.
Item~\textup{(E1.9)} is
\Cref{thm:YMA1-W2-control,thm:YMA1-cofinal-W2}.
Item~\textup{(E1.10)} is
\Cref{thm:YMA1-failure-classifier,cor:YMA1-grid-landing} and imports the
previously proved finite-physical/ultraviolet subsequence alternative.
Item~\textup{(E1.11)} is the nonduplicating interface recorded in
\Cref{ass:YMA1-higher-jet-interface,fire:YMA1-higher-jet-no-redo}.
\end{proof}

\begingroup\small
\begin{longtable}{>{\raggedright\arraybackslash}p{0.29\textwidth}
                  >{\raggedright\arraybackslash}p{0.19\textwidth}
                  >{\raggedright\arraybackslash}p{0.44\textwidth}}
\caption{Version V1 proof-status ledger.}
\label{tab:YMA1-ledger}\\
\toprule
\textbf{Row}&\textbf{Status}&\textbf{Exact statement}\\
\midrule
\endfirsthead
\toprule
\textbf{Row}&\textbf{Status}&\textbf{Exact statement}\\
\midrule
\endhead
Central-vacuum axial scalarization
&\status{Proved conditionally on the exact symmetry audit}
&The complete first axial transverse colour block is one scalar with complex
multiplicity twenty-four.\\[0.35em]
One-polarization physical card
&\status{Proved exact}
&One fixed real sine/cosine colour polarization supplies the deterministic and
connected-score scalars; the nonzero-mode mean is zero.\\[0.35em]
Translated-pair head acquisition
&\status{Exact finite card}
&Every head coefficient is deterministic cross curvature minus connected
score cross covariance minus the target reference coefficient.\\[0.35em]
Ward-relative tail
&\status{Proved}
&The Ward row sum fixes both repaired contacts; the directly fused signed tail
is at most \(\omg(q)\) times its second longitudinal displacement moment.\\[0.35em]
Fej\'er--Riesz margin compiler
&\status{Proved exact}
&The finite axial quotient is nonnegative exactly when one finite Hermitian
Gram with prescribed diagonal sums is positive semidefinite.\\[0.35em]
Validated margin interval
&\status{Proved exact}
&Head intervals and tail stocks give one-sided pass and obstruction tests,
a rigorous interval for the continuous interpolating margin, and a separate
upper interval on the actual finite momentum grid.\\[0.35em]
Ward-quotient cutoff transport
&\status{Proved sufficient}
&Summable \(\Wtwo\) coefficient defects give uniform cofinal transport,
including the zero-angle extension.\\[0.35em]
Actual first-shell scalar values
&\status{Open physical acquisition}
&The selected reduced-fibre law has not yet supplied the numerical or analytic
values of \(d_{n_*},c_{n_*},s_{n_*}\), or the translated-pair head.\\[0.35em]
Actual central-axial audit
&\status{Open on the first escaping shell}
&The physical selected shell must verify the trivial central representative,
exact axial little-group invariance, invariant high projection, and local/topological
split.\\[0.35em]
Weak-facing-root escape
&\status{Open trajectory row}
&A bounded subsequence remains on the inherited cylindrical-Haar/contact
branch.\\[0.35em]
All-orbit margin and higher jets
&\status{Open physical values}
&Only after the first axial card passes should the next momentum little-group
orbit and the fully fused third-jet transport be acquired.\\[0.35em]
Vacuum uniqueness, OS contraction, and mass gap
&\status{Not proved}
&No continuum reflected vacuum or full source-complete contraction has been
constructed.\\
\bottomrule
\end{longtable}
\endgroup

\subsection{Exact Version V2 attack}

\begin{openproblem}[Version V2: populate the first scalar axial Wilson card]
\label{op:YMA1-V2}
Preserve Version~V1 verbatim.  Keep the physical response, weak-facing
selectors, exact dyadic block map, reduced fibre, rooted gauge and Hodge
conventions, local/topological split, source/follower accounting, and the
predeclared first axial orbit fixed.  Proceed in this order.
\begin{enumerate}[label=\textup{(V2.\arabic*)},leftmargin=4.8em]
\item Test weak-facing escape.  A bounded subsequence is returned to the
      inherited contact branch without opening this card.
\item On the first escaping shell, test the central-axial audit at
      \(W^0\).  Keep any nontrivial central-toron card in the topological
      packet; do not average away a failed symmetry or stabilizer row.  Test
      full \(B_4\) covariance separately only before exporting the card to
      other axes.
\item Freeze \(n_*\), \(e_2\), and \(X_*\) before reading any score or
      Hessian value.
\item Acquire the exact reduced-fibre scalars
      \[
       d_{n_*}=\mathbb EK[v_*,v_*],
       \qquad
       c_{n_*}=\mathbb E\abs{j[v_*]}^2,
       \qquad
       s_{n_*}=S_0[v_*,v_*].
      \]
      Report the rigorous interval for \((d_{n_*}-c_{n_*})/s_{n_*}\).
\item If a uniform axial-angle statement is required, acquire the translated
      pair coefficients through the first predeclared radius \(R\), preserving
      the signed deterministic-minus-connected combination.
      Derive the repaired contact from the zero-row-sum identity rather than
      querying or fitting it independently.
\item Certify the second-moment tail in the Ward-relative normalization and
      include all coefficient-enclosure errors.
\item Run the Fej\'er--Riesz Gram test.  Return exactly \pass,
      \obstruction with its actual lattice angle, or \stopstatus with the
      first unresolved coefficient/tail interval.
\item Open the next momentum orbit only after the first axial card passes.
      Open graph, boundary, mediator, or pattern attribution only on a frozen
      failed mode or to prove the required physical tail.
\end{enumerate}
\end{openproblem}

\begin{assessment}[Exact V1 endpoint]
\label{ass:YMA1-endpoint}
The first local-vacuum momentum--colour block is no longer a matrix search.
On the central-axial branch it is the scalar
\[
 \boxed{
 \rho_{n_*}^{\rm ax}
 =\frac{
   \mathbb E K[v_*,v_*]
   -\mathbb E\abs{j[v_*]}^2
  }{S_0[v_*,v_*]}.}
\]
For the complete axial orbit, its finite translated-pair head is divided
exactly by the Ward factor \(2(1-\cos q)\), its omitted tail is priced by a
second displacement moment, and its margin is a finite positive-Gram
feasibility problem.  The next work is therefore literal population of these
physical scalars and coefficients.  Another unreduced fibre, arbitrary bank,
identity-normalized eigenvalue, or global contraction theorem before that
acquisition would move the same open row into a larger carrier.
\end{assessment}

\begin{firewall}[Claims not made in Version V1]
\label{fire:YMA1-claims}
Version~V1 does not prove that the weak-facing roots escape, that the actual
selected shell lies in the central-axial patch, that its first axial
margin is positive, that the signed kernel has the required cofinal
second-moment tail, that every momentum orbit passes, that the topological
packet is harmless, that higher jets transport, that a unique local or global
vacuum exists, that the full gauge-invariant local algebra contracts in
physical time, or that four-dimensional \(\SU(3)\) Yang--Mills theory has a
mass gap.
\end{firewall}

% =============================================================================
\section*{Version V1 append-only continuation rule}
\addcontentsline{toc}{section}{Version V1 append-only continuation rule}
% =============================================================================

For Version~V2, preserve the complete present source verbatim.  Remove only
the sole final command
\begin{center}
\verb|\end{document}|
\end{center}
Append new mathematical material immediately before that command and save the
result as
\begin{center}
\path{Renewal_Yang_Mills_Reduced_Fibre_Wilson_Shell_Symbol_}\\[-0.25em]
\path{and_Local_Vacuum_Programme_V2.tex}.
\end{center}
Do not alter the inherited reduced-fibre convention, central-axial audit,
predeclared axial orbit, fixed transverse colour polarization, deterministic
minus connected-score sign, contact repair, Ward quotient, second-moment tail,
Fej\'er--Riesz Gram semantics, \(\Wtwo\) transport norm, local/topological
separation, or \pass--\obstruction--\stopstatus meaning after reading a
physical scalar, translated-pair coefficient, tail interval, Gram, momentum
witness, patch transition, higher jet, vacuum, contraction, or mass datum.


% =============================================================================
% START APPENDED VERSION V2 MATERIAL
% =============================================================================
\clearpage
\hypersetup{
 pdftitle={Renewal Yang--Mills Reduced-Fibre Wilson Shell Symbol and Local-Vacuum Programme V2},
 pdfsubject={Literal Wilson fibre coercivity, a uniform plaquette interval, a tree-gauge jet correction, and the evaluated sixteen-alias axial reference}
}

\section{Version V2: an evaluated Wilson row and a necessary jet correction}
\label{sec:YMA2-direction}

Version~V1 is preserved above.  This continuation does not repeat its
Fourier positivity compiler or the inherited general shell theorems.  It
computes quantities on the actual parity-tree fibre of Dedicated
Yang--Mills V39.  The main new conclusions are an explicit non-Abelian
filling inequality with constant twenty, a cutoff-independent interval for
the mean corner plaquette, the complete deterministic Wilson Hessian, and a
closed rational formula for the axial flat reference.

The calculation also detects a qualification missing from the proof of
\Cref{thm:YMA1-axial-scalarization}.  In the fixed parity-tree section,
symmetry of the integrated action does not imply separate symmetry of the
raw deterministic Hessian and score covariance.  Reflections restore that
section by a fibre transformation which depends on the coarse field.
Their difference is still equivariant.  The explicit Wilson countercalculation
and the replacement statement are given in
\Cref{sec:YMA2-jet-correction}; no earlier source text has been changed.

\subsection{Exact scope and normalization}

Write the fine side as $N=2M$, with $M\ge4$, and keep the forest, block map,
and reduced variables already fixed in Dedicated V39:
\begin{equation}
 U_e(z,W)=
 \begin{cases}
 I,&e\in\mathscr T_N,\\
 W_{y,\mu},&e=(2y+\widehat e_\mu,2y+2\widehat e_\mu),\\
 z_e,&e\in\mathscr E_N^{\rm red}.
 \end{cases}
 \label{eq:YMA2-section}
\end{equation}
There are $45M^4$ reduced links and $96M^4$ fine plaquettes.  The Lie-algebra
metric in this continuation is
\begin{equation}
 \ip{X}{Y}=-\Tr(XY),\qquad X,Y\in\mathfrak{su}(3).
 \label{eq:YMA2-Lie-metric}
\end{equation}
The Frobenius matrix norm is denoted $\|\cdot\|_F$.

Define the unit Wilson action and its literal reduced-fibre integral by
\begin{align}
 V_M(z,W)&=\sum_{p}\left(1-\frac13\Re\Tr U_p(z,W)\right),
 \label{eq:YMA2-Wilson-action}\\
 Z_{\beta,\Psi,M}(W)&=\int
 e^{-\beta V_M(z,W)-\Psi(z,W)}\,\dd m_{\rm red}(z),
 \label{eq:YMA2-cylinder}\\
 \mathcal A_{\beta,\Psi,M}(W)&=-\log Z_{\beta,\Psi,M}(W).
 \label{eq:YMA2-effective-action}
\end{align}
Here $\beta>0$ is an independently fixed coefficient.  The function $\Psi$
is the \emph{actual} remaining fine action pulled through
\eqref{eq:YMA2-section}, when such terms are present.  It includes generated
interactions and any separately frozen marginal retuning not absorbed into
$\beta$.  The notation $\Psi=0$ always means the pure Wilson law, not an
approximation silently made to the response-matched law.

All coarse derivatives below are taken at $W=I$ in product exponential
coordinates, with the reduced variables held fixed.  Global conjugation
invariance makes the first derivative of the integrated action vanish at
$I$.  Its Hessian can consequently be computed before the coarse Hodge
projection, without an ambiguity from acceleration of the coarse coordinate
curve.  For statements involving nonzero $\Psi$, gauge invariance, coarse
translation invariance, and the reflection/permutation invariances actually
used are stated explicitly.

Let $\mathsf S_M$ denote the \emph{unit-coefficient} flat Wilson Schur form,
with the inherited characterization
\begin{equation}
 \ip{a}{\mathsf S_Ma}
 =\frac13\min_{\mathcal B b=a}\|d_fb\|_2^2.
 \label{eq:YMA2-unit-reference}
\end{equation}
Thus the reference in unscaled group-angle coordinates for $\beta V_M$ is
$\beta\mathsf S_M$.  A common physical field calibration is applied to both
forms by the same congruence; it does not change their generalized ratio.
No bare coefficient is reinserted after shell integration.

\begin{firewall}[Full blocked action and local contractible action]
\label{fire:YMA2-full-local}
The integral \eqref{eq:YMA2-cylinder} is the complete finite-torus integral.
If the inherited support decomposition gives
$\mathcal A=\mathcal A^{\rm loc}+\mathcal A^{\rm top}$, its raw score and
second-jet formula computes $H^{\rm loc}+H^{\rm top}$, not $H^{\rm loc}$
alone.  The topological Hessian must be retained and subtracted on the same
carrier before the zero-harmonic row of
\Cref{prop:YMA1-Ward-contact} is invoked.  A constant transverse coarse field
can carry a toron response in the full finite-torus integral.  This is not
removed by fitting a contact term.
\end{firewall}

\section{Forty-five literal link fillings and the constant twenty}
\label{sec:YMA2-filling}

Every fine vertex has a unique expression $x=2y+s$ with
$s\in\{0,1\}^4$.  Let $S(s)=\{\nu:s_\nu=1\}$ and write $P_s$ for the
word listing these directions in decreasing order.  This is the path from
the coarse root to $x$ along the inherited parity forest: reversing the
parent rule which removes the least odd coordinate gives precisely this
word.

For the positive fine edge $(x,\mu)$ compare the two positive paths
\begin{equation}
 A_{s,\mu}=P_s\,\mu,
 \qquad
 B_{s,\mu}=
 \begin{cases}
 P_{s+\widehat e_\mu},&s_\mu=0,\\
 \mu\,\mu\,P_{s-\widehat e_\mu},&s_\mu=1.
 \end{cases}
 \label{eq:YMA2-path-pair}
\end{equation}
In the second case the last forest path starts at the neighboring coarse
root.  The two words have the same fine endpoint.  At $W=I$, the holonomy
of $B_{s,\mu}$ is $I$, while that of $A_{s,\mu}$ is the edge variable in
\eqref{eq:YMA2-section}.

\begin{construction}[A fixed adjacent-swap filling]
\label{con:YMA2-swaps}
Transform $A_{s,\mu}$ into $B_{s,\mu}$ as follows.  Scan the desired target
word from the left.  At position $i$, locate the first still available copy
of the desired direction and move it to position $i$ by adjacent swaps.
Each swap of distinct directions $a,b$ at the current prefix vertex is the
elementary $ab$ plaquette there.  Retain repetitions with their literal
multiplicity.  The resulting plaquette list is $\Sigma_{s,\mu}$ and its
length is $m_{s,\mu}$.
\end{construction}

\begin{lemma}[Complete filling count]
\label{lem:YMA2-filling-count}
The filling lengths are
\begin{equation}
 m_{s,\mu}=
 \begin{cases}
 |\{\nu<\mu:s_\nu=1\}|,&s_\mu=0,\\
 2|\{\nu>\mu:s_\nu=1\}|+
 |\{\nu<\mu:s_\nu=1\}|,&s_\mu=1.
 \end{cases}
 \label{eq:YMA2-filling-length}
\end{equation}
Exactly nineteen of the sixty-four edge types have zero length: the fifteen
forest types and four designated second constituents.  The other forty-five
are exactly the reduced-link types.  Moreover,
\begin{equation}
 \max m_{s,\mu}=6,\qquad
 \sum_{s,\mu}m_{s,\mu}=96,\qquad
 \sum_{s,\mu}m_{s,\mu}^2=272.
 \label{eq:YMA2-filling-totals}
\end{equation}
If every translated filling contributes its length $m_{s,\mu}$ to each of
its plaquette occurrences, the largest resulting weighted plaquette
occupancy is exactly twenty.
\end{lemma}

\begin{proof}
For $s_\mu=0$, the appended $\mu$ must pass exactly the smaller occupied
directions to reach decreasing order.  For $s_\mu=1$, move the first copy
of $\mu$ to the first position and the appended copy to the second.
Every larger direction is crossed twice and every smaller direction once.
This proves \eqref{eq:YMA2-filling-length} and identifies the zero cases.
The complete finite tables below evaluate its sums and the weighted
occupancies.

A plaquette type consists of its base parity and its increasing direction
pair.  In the occupancy table, contributions from fillings rooted in all
neighboring coarse cells are included by reducing the plaquette base modulo
two.  A fixed plaquette of that type receives exactly the listed
contributions: its coarse translation determines the root of every
contributing filling uniquely.  This proves the interpretation of the table,
not just a count in one isolated cell.  Its largest entry is twenty.
\end{proof}

\begin{center}
\small
\begin{tabular}{c|rrrr@{\qquad}c|rrrr}
\toprule
$s$&$m_1$&$m_2$&$m_3$&$m_4$&$s$&$m_1$&$m_2$&$m_3$&$m_4$\\
\midrule
0000&0&0&0&0&1000&0&1&1&1\\
0001&0&0&0&0&1001&2&1&1&1\\
0010&0&0&0&1&1010&2&1&1&2\\
0011&0&0&2&1&1011&4&1&3&2\\
0100&0&0&1&1&1100&2&1&2&2\\
0101&0&2&1&1&1101&4&3&2&2\\
0110&0&2&1&2&1110&4&3&2&3\\
0111&0&4&3&2&1111&6&5&4&3\\
\bottomrule
\end{tabular}
\end{center}

\begin{center}
\small
\begin{tabular}{c|rrrrrr}
\toprule
Plaquette base parity&12&13&14&23&24&34\\
\midrule
0000&3&7&17&8&17&20\\
0001&5&11&1&12&3&8\\
0010&5&1&2&3&5&12\\
0011&7&3&2&7&5&0\\
0100&1&2&2&5&14&0\\
0101&3&2&2&9&0&0\\
0110&3&2&3&0&0&0\\
0111&5&4&3&0&0&0\\
1000&2&6&16&0&0&0\\
1001&4&10&0&0&0&0\\
1010&4&0&0&0&0&0\\
1011&6&0&0&0&0&0\\
1100&0&0&0&0&0&0\\
1101&0&0&0&0&0&0\\
1110&0&0&0&0&0&0\\
1111&0&0&0&0&0&0\\
\bottomrule
\end{tabular}
\end{center}

\begin{theorem}[Global non-Abelian coercivity of the pinned Wilson fibre]
\label{thm:YMA2-global-fibre-coercivity}
For every $M\ge4$ and every reduced configuration at $W=I$, let
\begin{equation}
 D(z):=\sum_{e\in\mathscr E_N^{\rm red}}\|z_e-I\|_F^2.
 \label{eq:YMA2-chord-stock}
\end{equation}
Then
\begin{equation}
 \boxed{
 D(z)\le20\sum_p\|U_p(z,I)-I\|_F^2,
 \qquad
 \frac1{120}D(z)\le V_M(z,I)\le4D(z).}
 \label{eq:YMA2-global-coercivity}
\end{equation}
In particular the reduced Wilson action has a unique global minimum,
$z_e=I$ for every reduced link.  In product exponential coordinates
$z_e=e^{x_e}$ its vertical Hessian $A_M=D_x^2V_M(0,I)$ obeys
\begin{equation}
 \boxed{\frac1{60}I\preceq A_M\preceq8I.}
 \label{eq:YMA2-vertical-Hessian-window}
\end{equation}
These constants do not depend on the lattice side.  They concern the pinned
vertical fibre, not the retained coarse field.
\end{theorem}

\begin{proof}
For unitary products, insertion and removal of one factor preserve the
Frobenius norm.  Hence an adjacent path swap changes the path holonomy by
exactly $\|U_p-I\|_F$, with an irrelevant unitary conjugation and possible
orientation reversal.  Telescoping the swaps of
\Cref{con:YMA2-swaps} gives
\[
 \|z_e-I\|_F\le\sum_{p\in\Sigma_e}\|U_p-I\|_F.
\]
Cauchy--Schwarz, retaining plaquette multiplicities, gives
\[
 \|z_e-I\|_F^2\le m_e\sum_{p\in\Sigma_e}\|U_p-I\|_F^2.
\]
Summation and the weighted occupancy twenty prove the first inequality.
The exact identity
\[
 1-\tfrac13\Re\Tr U=\tfrac16\|U-I\|_F^2
 \quad(U\in\SU(3))
\]
gives the lower action bound.  Conversely, each fine plaquette contains at
most four reduced-link factors, since every other factor equals $I$.
The same product inequality and Cauchy--Schwarz bound its squared chord by
four times the sum of those link chords.  A reduced link belongs to six
fine plaquettes.  Thus
$V_M\le(4\cdot6/6)D=4D$.

The lower bound makes the zero set precisely $z=I$.  Finally,
$D(e^{tx})=t^2\|x\|_2^2+O(t^4)$ and
$V_M(e^{tx},I)=\frac12t^2\ip{x}{A_Mx}+O(t^3)$.
Comparison of the quadratic coefficients proves
\eqref{eq:YMA2-vertical-Hessian-window}.
\end{proof}

\begin{remark}[Why this is stronger than a local flat-patch assertion]
The first inequality in \eqref{eq:YMA2-global-coercivity} holds on the
\emph{whole compact reduced fibre}.  It rules out a second flat minimum
inside that fibre and does not assume that a logarithm chart contains the
sampled field.  It does not assert global convexity in logarithm coordinates,
a uniform Gibbs covariance bound, or a coarse-field mass.  Those would be
strictly different assertions.
\end{remark}

\section{An explicit all-volume interval for the Wilson corner plaquette}
\label{sec:YMA2-plaquette-interval}

A fine plaquette is a \emph{corner plaquette} if it touches a coarse vertex.
There are exactly $24M^4$ such plaquettes: six direction pairs and four
quadrants at each coarse vertex, with no double counting.  Every plaquette
incident to a designated second constituent is a corner plaquette.

Assume for the moment that the full conditional law in
\eqref{eq:YMA2-cylinder} is invariant under coarse translations and all
coordinate reflections and permutations.  These symmetries preserve the
constraint $\mathbf B_NU=I$ and act transitively on corner plaquettes.  Put
\begin{equation}
 w_{\beta,\Psi,M}:=
 \mathbb E_{\beta,\Psi,M}\left[\tfrac13\Re\Tr U_{p_*}\right]
 \label{eq:YMA2-corner-w}
\end{equation}
for any one fixed corner plaquette $p_*$.  This is a gauge-invariant scalar
of the actual conditional law.  Write
\begin{equation}
 \operatorname{osc}\Psi_M
 :=\sup_z\Psi(z,I)-\inf_z\Psi(z,I).
 \label{eq:YMA2-osc-Psi}
\end{equation}

\subsection{A completely explicit one-link comparison integral}

\begin{lemma}[Two-sided $\SU(3)$ chord integral]
\label{lem:YMA2-one-link-chord}
For $t>0$ define
\begin{equation}
 I(t):=\int_{\SU(3)}e^{-t\|U-I\|_F^2}\,\dd U.
 \label{eq:YMA2-I-def}
\end{equation}
With normalized Haar probability, the following explicit bounds hold for
$t\ge1$:
\begin{equation}
 \boxed{C_-t^{-4}\le I(t)\le C_+t^{-4},\qquad
 C_-:=\frac{3e^{-43/32}}{1024\pi^8},\quad
 C_+:=\frac{25\pi^7}{512}.}
 \label{eq:YMA2-one-link-bounds}
\end{equation}
The upper bound is valid for every $t>0$.
\end{lemma}

\begin{proof}
In the maximal torus use eigenangles
$(\theta_1,\theta_2,\theta_3)=(\theta_1,\theta_2,-\theta_1-\theta_2)$.
The conjugation change of variables has Jacobian
$\prod_{i<j}|e^{i\theta_i}-e^{i\theta_j}|^2$: on each real two-dimensional
root plane, conjugation contributes the squared modulus of
$e^{i\theta_i}-e^{i\theta_j}$.  Its regular part covers a conjugacy class
six times.  The singular set has Haar measure zero.  Normalizing the torus
measure gives the Weyl formula in the precise convention needed here,
\begin{equation}
 I(t)=\frac1{24\pi^2}\int_{[-\pi,\pi]^2}
 e^{-t[6-2(\cos\theta_1+\cos\theta_2+\cos(\theta_1+\theta_2))]}
 \prod_{i<j}|e^{i\theta_i}-e^{i\theta_j}|^2\,
 d\theta_1d\theta_2.
 \label{eq:YMA2-Weyl-integral}
\end{equation}
For completeness, the constant is also checked by expanding the squared
Vandermonde determinant at $t=0$.  Its six diagonal permutation terms
integrate to six; all cross terms have a nontrivial torus character and
integrate to zero.  Thus the displayed measure integrates the constant to
one.

Let $r^2=\theta_1^2+\theta_2^2$.  The exponent's chord satisfies
$\|U-I\|_F^2\ge4r^2/\pi^2$, using only the first two eigenangles.  The
three linear angle differences have coefficient norms
$\sqrt2,\sqrt5,\sqrt5$.  Since $|e^{ia}-e^{ib}|\le|a-b|$, their squared
product is at most $50r^6$.  Extend the integral to $\R^2$ to obtain
\[
 I(t)\le\frac{50}{24\pi^2}
 \int_{\R^2}r^6e^{-4tr^2/\pi^2}\,d\theta
 =\frac{25\pi^7}{512}\,t^{-4}.
\]

For the lower bound, retain only the rectangle
\[
 \frac1{2\sqrt t}\le\theta_1\le\frac3{4\sqrt t},
 \qquad 0\le\theta_2\le\frac1{8\sqrt t}.
\]
It has area $1/(32t)$ and lies inside the integration square for $t\ge1$.
The three positive angle differences are at least
$3/(8\sqrt t),1/\sqrt t,1/(2\sqrt t)$, respectively, and are at most
$13/8<\pi$.  The inequality
$|e^{ia}-e^{ib}|\ge2|a-b|/\pi$ on this range gives a squared Vandermonde
at least $9/(4\pi^6t^3)$.  Also
\[
 \|U-I\|_F^2\le\theta_1^2+\theta_2^2+
 (\theta_1+\theta_2)^2\le\frac{43}{32t}.
\]
Multiplying these lower bounds and the normalization in
\eqref{eq:YMA2-Weyl-integral} gives $C_-t^{-4}$.
\end{proof}

\begin{theorem}[Cutoff-independent acquisition of the deterministic plaquette row]
\label{thm:YMA2-plaquette-interval}
Define
\begin{equation}
 \boxed{C_{\rm W}:=\frac{15}{4}\log\left(
 \frac{50}{3}e^{43/32}\pi^{15}960^4\right).}
 \label{eq:YMA2-CW}
\end{equation}
For every $M\ge4$, every $\beta\ge240$, and every continuous $\Psi$ with
the corner-plaquette symmetries just stated,
\begin{equation}
 \boxed{
 1-\frac{C_{\rm W}}{\beta}
   -\frac{\operatorname{osc}\Psi_M}{12\beta M^4}
 \le w_{\beta,\Psi,M}\le1.}
 \label{eq:YMA2-w-interval}
\end{equation}
In particular, for the pure Wilson law,
\begin{equation}
 \boxed{1-C_{\rm W}/\beta\le w_{\beta,0,M}\le1}
 \label{eq:YMA2-pure-w-interval}
\end{equation}
uniformly in $M$.  The constant is approximately $182.984408$; its exact
value is \eqref{eq:YMA2-CW}.  The elementary bound $C_{\rm W}<201$ is
sufficient to make the pure lower endpoint strictly positive throughout
$\beta\ge240$.  More generally,
$\operatorname{osc}\Psi_M\le KM^4$ adds at most $K/(12\beta)$ to the
same all-volume error.
\end{theorem}

\begin{proof}
For this proof alone vary an auxiliary coefficient $t$, keeping the
\emph{submitted} $\Psi$ fixed.  Write
\[
 Z_\Psi(t)=\int e^{-tV_M(z,I)-\Psi(z,I)}\,dm_{\rm red}(z),
 \qquad f(t)=\log Z_\Psi(t).
\]
Then $-f'(t)=\mathbb E_{t,\Psi,M}V_M$ and
$f''(t)=\Var_{t,\Psi,M}(V_M)\ge0$.  Consequently
\begin{equation}
 \mathbb E_{\beta,\Psi,M}V_M
 \le\frac2\beta\bigl(f(\beta/2)-f(\beta)\bigr).
 \label{eq:YMA2-pressure-sec}
\end{equation}
This is an inequality for the one law at the submitted value $\beta$; it
neither retunes a physical trajectory nor assumes that $\Psi$ is constant
along that trajectory.

The global comparisons in \eqref{eq:YMA2-global-coercivity}, applied under
product Haar measure on $45M^4$ links, imply
\[
 e^{-\sup\Psi_M}I(4t)^{45M^4}
 \le Z_\Psi(t)\le
 e^{-\inf\Psi_M}I(t/120)^{45M^4}.
\]
For $\beta\ge240$, \Cref{lem:YMA2-one-link-chord} therefore gives
\[
 f(\beta/2)-f(\beta)
 \le\operatorname{osc}\Psi_M+
 45M^4\log\left(\frac{C_+}{C_-}960^4\right).
\]
Every Wilson plaquette defect is nonnegative.  Keeping only the $24M^4$
corner plaquettes in its expectation yields
\[
 24M^4(1-w_{\beta,\Psi,M})\le
 \mathbb E_{\beta,\Psi,M}V_M.
\]
The preceding three displays give \eqref{eq:YMA2-w-interval}; the upper
endpoint follows from $\Re\Tr U\le3$ pointwise.  Substitution of $C_\pm$
gives \eqref{eq:YMA2-CW}.  For a coarse elementary bound, use
$\log(50/3)<3$, $\log\pi<\log4<1.4$, $\log960<\log1024<7$, and
$43/32<1.35$.  These give $C_{\rm W}<201<240$.  The final assertion follows
by substitution of the oscillation density bound.
\end{proof}

\begin{remark}[What has actually been acquired]
The interval \eqref{eq:YMA2-w-interval} is a theorem about an actual Wilson
expectation, rather than a requested future integration panel.  Its constants
are crude but explicit and do not grow with the cutoff.  For a nonzero
$\Psi$, its oscillation density and its symmetry are independent conditions,
not inferred from the name ``generated action.''  Without symmetry the proof
still bounds the \emph{average} corner plaquette, but does not identify that
average with the selected single-corner value.
\end{remark}

\section{The complete deterministic Wilson Hessian is now explicit}
\label{sec:YMA2-direct-Hessian}

Let $K^{\rm W}=D_W^2(\beta V_M)$ be the raw second jet in
\eqref{eq:YMA2-section}.  The expectation in the following result may be
under the full law $\nu_{\beta,\Psi,M}$; it need not be under the pure law.
Assume its gauge and corner-plaquette symmetries.  For a coarse angle
$q\in(2\pi/M)\Z^4$, set
\begin{equation}
 u(q)=(e^{iq_1},e^{iq_2},e^{iq_3},e^{iq_4})^T,
 \qquad
 \mathsf C(q):=\frac13\bigl(7I_4-u(q)u(q)^*\bigr).
 \label{eq:YMA2-C-symbol}
\end{equation}
The Fourier convention is that of Dedicated V39 and V1, with a positive
phase in the inverse plane-wave synthesis.

\begin{theorem}[Exact Wilson second-jet symbol]
\label{thm:YMA2-Wilson-direct-symbol}
In the metric \eqref{eq:YMA2-Lie-metric},
\begin{equation}
 \boxed{
 \mathbb E_{\beta,\Psi,M}\widehat K^{\rm W}(q)
 =\beta w_{\beta,\Psi,M}\,
 \mathsf C(q)\otimes I_{\mathfrak g_\C}.}
 \label{eq:YMA2-exact-direct-symbol}
\end{equation}
The direction eigenvalues of $\mathsf C(q)$ are $1$ once and $7/3$ three
times.  In particular, for an axial angle $q=(q_1,0,0,0)$, its transverse
block is
\begin{equation}
 \boxed{
 \mathbb E\widehat K^{\rm W}_\perp(q_1)
 =\frac{\beta w_{\beta,\Psi,M}}3
  (7I_3-\mathbf 1_3\mathbf 1_3^T)\otimes I_8.}
 \label{eq:YMA2-transverse-direct}
\end{equation}
Its three transverse diagonal entries are all $2\beta w_{\beta,\Psi,M}$,
but its off-diagonal entries are $-\beta w_{\beta,\Psi,M}/3$.
For $\Psi=0$ and $\beta\ge240$, the entire direct matrix has the explicit
uniform enclosure
\begin{equation}
 \boxed{
 (\beta-C_{\rm W})\mathsf C(q)\otimes I_8
 \preceq\mathbb E\widehat K^{\rm W}(q)
 \preceq\beta\mathsf C(q)\otimes I_8.}
 \label{eq:YMA2-direct-interval}
\end{equation}
\end{theorem}

\begin{proof}
Only designated second fine constituents vary when a coarse link varies in
\eqref{eq:YMA2-section}.  A designated link belongs to six fine plaquettes,
all touching its coarse endpoint.  Each fine plaquette contains at most two
designated links.  Two occur together exactly when they end at the same
coarse vertex, in distinct directions.

For one coarse tangent $X$, the corresponding second derivative of a
plaquette defect is $-\frac13\Re\Tr(X^2Q)$, with $Q$ its based holonomy,
up to conjugation or inversion.  Global gauge conjugation at $W=I$ implies
$\mathbb E Q=\alpha I_3$, and
$\Re\alpha=\mathbb E\Re\Tr Q/3=w_{\beta,\Psi,M}$.  Thus every such
plaquette contributes $w\|X\|^2/3$ after averaging.  The six incidences
give the diagonal $2w$.

When two designated links meet, base the plaquette at their common coarse
endpoint.  The two varying factors are adjacent there, with opposite
orientations.  Its mixed second derivative averages to
$-w\ip{X}{Y}/3$.  There is exactly one such plaquette per coarse endpoint
and unordered direction pair.  The two coarse links start at
$y-\widehat e_\mu$ and $y-\widehat e_\nu$.  Fourier synthesis supplies
the phase $e^{i(q_\mu-q_\nu)}$ in row $\mu$, column $\nu$.
Consequently the diagonal is $6/3$ and the off-diagonal is
$-e^{i(q_\mu-q_\nu)}/3$, times $\beta w$ and the colour identity.  This
is \eqref{eq:YMA2-exact-direct-symbol}.

The vector $u$ has norm squared four; hence $7I-uu^*$ has eigenvalues three
and seven.  Restriction to the transverse axial directions sets their phases
to one and proves \eqref{eq:YMA2-transverse-direct}.  Finally use the pure
interval \eqref{eq:YMA2-pure-w-interval} and positivity of $\mathsf C$.
\end{proof}

\begin{corollary}[The same-direction off-site deterministic kernel is zero]
\label{cor:YMA2-no-offsite-direct}
For any fixed direction $\mu$ and distinct coarse sites $x,y$,
\begin{equation}
 \mathbb E K^{\rm W}[(x,\mu),(y,\mu)]=0.
 \label{eq:YMA2-no-offsite-K}
\end{equation}
Thus, in the pure Wilson law, every noncontact displacement coefficient in a
fixed-direction physical Hessian comes from minus the connected score
covariance.  A generated interaction can add its own deterministic
noncontact coefficient, but does not change this Wilson incidence identity.
\end{corollary}

\begin{proof}
No fine plaquette contains two distinct designated links of the same
direction, so the mixed derivative vanishes before integration.
\end{proof}

\section{Correction of separate jet scalarization in the parity-tree gauge}
\label{sec:YMA2-jet-correction}

The integrated action is independent of the auxiliary section used to
represent its fibre.  The two raw terms in its Hessian decomposition are
not.  For a coarse symmetry $g$, the parity-tree restoring map generally has
the form $z\mapsto\chi_g(W,z)$, not $z\mapsto\chi_g(z)$.  Differentiating
an identity involving $F(\chi_g(W,z),gW)$ produces the extra vertical term
$D_zF\,D_W\chi_g$.  Haar preservation at each fixed $W$ does not make
that term vanish pointwise.  This is the missing distinction in the
separate-commutation step of the V1 proof.

\begin{countertheorem}[The literal Wilson raw jets are not separately scalar]
\label{cth:YMA2-raw-jets-not-scalar}
On the pure Wilson reduced fibre at $W=I$, take $\beta\ge240$ and any
nonzero axial angle.  The physical Hessian and the flat reference are scalar
on the transverse--adjoint block, but the raw mean deterministic Hessian is
not scalar there.  Specifically, for two different transverse directions
$\mu,\nu$ and the same unit colour $X$,
\begin{equation}
 \boxed{
 \mathbb E K^{\rm W}[v_{\mu,X},v_{\nu,X}]
 =-\frac{\beta w_{\beta,0,M}}3\ne0.}
 \label{eq:YMA2-mixed-raw-nonzero}
\end{equation}
The raw connected score covariance has exactly the same nonzero mixed
entry.  Consequently invariance of the fibre law at the central point,
even under the full physical little group, is insufficient to justify the
separate scalar identities for these raw jets in
\eqref{eq:YMA1-scalar-blocks}.
\end{countertheorem}

\begin{proof}
The pure Wilson action and block map are equivariant under coarse
translations, coordinate reflections/permutations, and constant global
conjugation.  At $W=I$, the integrated Hessian therefore commutes with the
little-group representation.  The irreducibility argument already proved
in V1 applies to that \emph{integrated} Hessian and gives a scalar on the
transverse--adjoint block.  It applies to the flat reference as well.

In contrast, \Cref{thm:YMA2-Wilson-direct-symbol} gives
\eqref{eq:YMA2-mixed-raw-nonzero}.  Its value is nonzero because
\Cref{thm:YMA2-plaquette-interval} gives $w>0$ for $\beta\ge240$.
A scalar matrix cannot have this mixed entry.  In the exact identity
$H=\mathbb EK-\Cov(j)$ the physical mixed entry is zero, so the covariance
mixed entry equals the deterministic one.  The latter is negative but is
compatible with a positive covariance \emph{matrix}.
\end{proof}

\begin{proposition}[Correct scalar card without separate raw-jet equivariance]
\label{prop:YMA2-correct-scalar-card}
Let the integrated action, reference, and retained projection satisfy the
physical central-axial symmetry requirements of V1.  No separate covariance
symmetry is assumed.  On one retained axial block write
$\widehat H_\perp(q)=h(q)I$ and
$\widehat S_0(q)=s(q)I$.  For any predeclared unit polarization $v_*$,
\begin{equation}
 \boxed{h(q)=\mathbb E K[v_*,v_*]-\Var(j[v_*]).}
 \label{eq:YMA2-correct-card}
\end{equation}
Thus one-polarization acquisition of the \emph{difference} remains exact.
Separate scalarization of $\mathbb EK$ and $\Cov(j)$ may be used only after
an additional differentiated, fixed-fibre equivariance test, or a verified
replacement for it, has passed.

In the pure Wilson law, with $c_{\beta,M}(q)$ the diagonal variance of the
fixed unit transverse colour polarization,
\begin{equation}
 \boxed{h_{\beta,M}(q)=2\beta w_{\beta,0,M}-c_{\beta,M}(q).}
 \label{eq:YMA2-pure-scalar-card}
\end{equation}
The complete transverse covariance is therefore
\begin{equation}
 \boxed{\widehat{\Cov(j)}_\perp(q)
 =\frac{\beta w_{\beta,0,M}}3
 (7I_3-\mathbf1_3\mathbf1_3^T)\otimes I_8-h_{\beta,M}(q)I_{24}.}
 \label{eq:YMA2-complete-axial-covariance}
\end{equation}
\end{proposition}

\begin{proof}
Compress the inherited exact score--Hessian identity to $v_*$.  The physical
Hessian is scalar by the stated symmetry, so its compression equals $h$.
For a nonzero translation character the mean score is zero in the frozen
translation-covariant parity section, as already proved in Dedicated V39.
The real sine/cosine implementation is unchanged.  Substitution of the
Wilson diagonal and then of its full transverse direct block gives the last
two identities.
\end{proof}

\begin{assessment}[What the correction does and does not invalidate]
The physical scalar ratio, the one-polarization difference, and the
Ward-relative Fourier estimates of V1 survive on their stated local and
symmetry branches.  The assertions that the two raw matrices are separately
scalar do not apply to the literal parity-tree Wilson jets without the
stronger equivariance condition.  The nonzero mixed entry is a kinematic
section effect, not a negative physical shell symbol, not a new source, and
not a reason to abandon the scalar physical card.  Only their fully assembled
difference is used henceforth.
\end{assessment}

\subsection{A precise same-law use with generated interactions}

Retain the complete score $j_\Phi=j_{\beta V}+j_\Psi$ in the actual law
$\nu_{\beta,\Psi,M}$.  The inherited response-matched mixed-jet theorem
already specifies all of its cross covariances and is not reproved here.
The new Wilson calculation reduces its axial diagonal to
\begin{equation}
 \boxed{
 h_{\beta,\Psi,M}(q)
 =2\beta w_{\beta,\Psi,M}
  +\mathbb E_{\beta,\Psi,M}K^\Psi[v_*,v_*]
  -\Var_{\beta,\Psi,M}(j_\Phi[v_*]).}
 \label{eq:YMA2-generated-card}
\end{equation}
In particular, the plaquette interval applies in the same full law whenever
its oscillation-density hypothesis passes.  Neither a pure-law covariance
nor a separately sampled generated score is substituted into
\eqref{eq:YMA2-generated-card}.

\section{Evaluation of the complete sixteen-alias axial reference}
\label{sec:YMA2-reference}

The variational flat Wilson reference was already proved in Dedicated V34.
Here its first axial symbol is evaluated, rather than replaced by another
lower-bound compiler.  Work at $W=I$ in the metric
\eqref{eq:YMA2-Lie-metric}.  For the coarse angle $(q,0,0,0)$ the fine
aliases are
\begin{equation}
 p_\sigma=(q/2,0,0,0)+\pi\sigma,
 \qquad\sigma\in\{0,1\}^4.
 \label{eq:YMA2-sixteen-aliases}
\end{equation}
Let $d_\mu(p)=e^{ip_\mu}-1$, $\omega_f(p)=\sum_\mu|d_\mu(p)|^2$, and
$B_\sigma=\operatorname{diag}(1+e^{ip_{\sigma,\mu}})$.

\begin{theorem}[Exact axial flat Wilson coefficient]
\label{thm:YMA2-reference-symbol}
For every nonzero allowed coarse axial angle, put $t=\sin^2(q/2)$.  The
unit reference \eqref{eq:YMA2-unit-reference} has the scalar transverse
symbol
\begin{equation}
 \boxed{s_0(q)=\frac{t(t+8)(t+24)}{9(t^2+18t+16)}.}
 \label{eq:YMA2-reference-rational}
\end{equation}
Its Ward quotient is
\begin{equation}
 \boxed{\frac{s_0(q)}{\omega(q)}
 =\frac{(t+8)(t+24)}{36(t^2+18t+16)},
 \qquad \omega(q)=4t,}
 \label{eq:YMA2-reference-quotient}
\end{equation}
and satisfies the sharp continuous-angle bounds
\begin{equation}
 \boxed{\frac5{28}\le\frac{s_0(q)}{\omega(q)}\le\frac13.}
 \label{eq:YMA2-reference-endpoints}
\end{equation}
The upper endpoint is its $q\to0$ limit; the lower endpoint is attained at
$q=\pi$.  Useful exact checks are
\begin{equation}
 \boxed{s_0(\pi/2)=\frac{833}{1818},\qquad s_0(\pi)=\frac57.}
 \label{eq:YMA2-reference-checks}
\end{equation}
These are four-dimensional, straight-link dyadic reference values; they are
not the older three-dimensional selected-surface trace coefficients.
\end{theorem}

\begin{proof}
The fine unit-coefficient curvature Hessian has, on a nonzero Fourier mode,
\[
 L_f(p)=\tfrac13\bigl(\omega_f(p)I-d(p)d(p)^*\bigr).
\]
Its inverse on the transverse support is
$3\omega_f(p)^{-1}(I-d(p)d(p)^*/\omega_f(p))$.  Fine pure gauges map under
the block differential to coarse pure gauges, since
$(1+e^{ip_\mu})(e^{ip_\mu}-1)=e^{iq_\mu}-1$ for each alias.  Thus the
coarse transverse compliance is obtained by projecting the blocked fine
compliance to the coarse transverse space.  The unitary transforms on
$(2M)^4$ and $M^4$ give a factor $1/4$ in the block map and hence $1/16$
in the compliance.  Equivalently, the constrained quadratic minimum in
\eqref{eq:YMA2-unit-reference}, by Lagrange multipliers modulo the fine
gauge kernel, has transverse compliance
\begin{equation}
 \frac3{16}\sum_\sigma
 B_\sigma\frac{I-d(p_\sigma)d(p_\sigma)^*/\omega_f(p_\sigma)}
 {\omega_f(p_\sigma)}B_\sigma^*.
 \label{eq:YMA2-alias-compliance}
\end{equation}
Only its coarse transverse compression is inverted.

For the fixed transverse direction $2$, the factor
$1+e^{ip_{\sigma,2}}$ is zero for $\sigma_2=1$ and equals two for
$\sigma_2=0$.  In the latter case $d_2(p_\sigma)=0$, so the transverse
projector does not reduce the diagonal.  Thus all sixteen aliases are
accounted for, eight by exact zeros and eight by the following sum.  Put
$a=\sin^2(q/4)$ and $b=\cos^2(q/4)$.  The other two transverse parity
indices contribute $n=0,1,2$ with multiplicities $1,2,1$, and the first
parity index gives $a$ or $b$.  Therefore
\begin{equation}
 s_0(q)^{-1}=\frac3{16}\sum_{n=0}^2\binom2n
 \left(\frac1{a+n}+\frac1{b+n}\right).
 \label{eq:YMA2-axial-compliance-sum}
\end{equation}
Transverse off-diagonal entries vanish: when both corresponding block
factors are nonzero, both corresponding fine differences are zero.  The
longitudinal--transverse compliance entries vanish by the same argument.

Since $a+b=1$ and $ab=t/4$, the sum without its $3/16$ prefactor is
\[
 \frac1{ab}+\frac6{2+ab}+\frac5{6+ab}.
\]
Inversion gives \eqref{eq:YMA2-reference-rational}.  Differentiation of its
Ward quotient as a rational function of $t$ gives
\[
 -\frac{7t^2+176t+1472}{18(t^2+18t+16)^2}<0
 \quad(0\le t\le1).
\]
Its values at zero and one are $1/3$ and $5/28$, proving the bounds.
Substitution at $t=1/2$ and $t=1$ gives
\eqref{eq:YMA2-reference-checks}.
\end{proof}

\begin{corollary}[The leading low-angle coefficient is evaluated]
\label{cor:YMA2-reference-small-angle}
The reference has the expansion
\begin{equation}
 s_0(q)=\frac{q^2}{3}-\frac{31q^4}{288}+O(q^6)
 \quad(q\to0).
 \label{eq:YMA2-reference-small-q}
\end{equation}
In particular, the first axial reference in V1 can be evaluated at its
already frozen $q_{n_*}$ by \eqref{eq:YMA2-reference-rational}, without an
unresolved flat-reference integration query.
\end{corollary}

\begin{proof}
The quotient in \eqref{eq:YMA2-reference-quotient} has Taylor expansion
$1/3-23t/72+O(t^2)$.  Multiply by $4t$ and use
$t=q^2/4-q^4/48+O(q^6)$.
\end{proof}

\section{Fixed-card weak-coupling evaluation of the full compact integral}
\label{sec:YMA2-Laplace}

The all-volume plaquette bound above did not use a Laplace approximation.
The next result evaluates the remaining leading covariance cancellation on
\emph{each fixed finite card}.  Its constants are not declared uniform in
$M$.  This distinction is essential: even rigorous Laplace expansions in
growing dimension require additional control of their coefficients and
remainders; see the primary analysis in \cite{YMA2-Katsevich}.
The proof below establishes the particular finite Wilson statement directly.

In product logarithm coordinates $x$ for the reduced links and coarse
coordinates $v$, let $x_*(v)$ be the local critical point of the unit Wilson
action and put
\begin{equation}
 E_*(v):=V_M(x_*(v),v),\qquad
 A(v):=D_x^2V_M(x_*(v),v),\qquad
 dm_{\rm red}=J(x)\,dx.
 \label{eq:YMA2-Laplace-objects}
\end{equation}
The coordinate density $J$ is fixed, rather than fitted to an asymptotic
answer.

\begin{theorem}[Literal finite Wilson two-jet asymptotics]
\label{thm:YMA2-finite-Laplace}
For every fixed $M\ge4$, there is a coarse neighborhood on which $x_*(v)$
is the unique global minimizing reduced configuration, is analytic, and
$A(v)>0$.  With $d_M=360M^4$, one has on a smaller neighborhood
\begin{equation}
 \mathcal A_{\beta,0,M}(v)
 =\beta E_*(v)+\frac{d_M}{2}\log\frac\beta{2\pi}
 +\Gamma_M(v)+R_{\beta,M}(v),
 \label{eq:YMA2-Laplace-action}
\end{equation}
where
\begin{equation}
 \boxed{\Gamma_M(v)=\frac12\log\det A(v)-\log J(x_*(v)),
 \qquad \|R_{\beta,M}\|_{C^2}\le C_M\beta^{-1}}
 \label{eq:YMA2-Laplace-correction}
\end{equation}
for all sufficiently large $\beta$.  The finite constant $C_M$ and the
threshold are permitted to depend on the whole card.  At $v=0$,
\begin{equation}
 \boxed{H_{\beta,M}=\beta\mathsf S_M+D^2\Gamma_M(0)
   +O_M(\beta^{-1})}
 \label{eq:YMA2-finite-Hessian-expansion}
\end{equation}
in the finite-dimensional operator norm.  Thus for every fixed nonzero
axial angle,
\begin{equation}
 \boxed{h_{\beta,M}(q)=\beta s_0(q)+\gamma_M(q)+O_M(\beta^{-1}),
 \qquad \frac{h_{\beta,M}(q)}{\beta s_0(q)}\longrightarrow1.}
 \label{eq:YMA2-finite-axial-limit}
\end{equation}
Here $\gamma_M(q)$ is the same physical polarization of $D^2\Gamma_M(0)$.
The zero-angle classical coefficient is zero; a finite one-loop harmonic
coefficient is not thereby excluded.
\end{theorem}

\begin{proof}
By \Cref{thm:YMA2-global-fibre-coercivity}, the unique global minimizer at
$v=0$ is $x=0$ and its vertical Hessian is positive.  The implicit function
theorem gives an analytic critical point $x_*(v)$ in a small logarithm
neighborhood.  Shrink that neighborhood so that the vertical Hessian there
is positive.  Compactness of the complement and uniqueness of the minimum
at $v=0$ give a strictly positive action gap outside it.  Continuity preserves
a smaller gap for sufficiently small $v$.  The local critical point is
therefore the unique global minimizer as well.

We detail why two coarse derivatives of the asymptotic remainder are
controlled.  Around the minimizing point set $\xi=x-x_*(v)$ and write
\[
 V_M(x_*(v)+\xi,v)-E_*(v)
 =\tfrac12\xi^T\widetilde A(\xi,v)\xi,
 \qquad
 \widetilde A(\xi,v)=2\int_0^1(1-s)
 D_x^2V_M(x_*(v)+s\xi,v)\,ds.
\]
The positive analytic square root defines the local change of variables
$y=\widetilde A(\xi,v)^{1/2}\xi$.  Its derivative at $\xi=0$ is
$A(v)^{1/2}$, so it is an analytic local diffeomorphism with an analytic
parameter-dependent inverse.  The action becomes exactly
$E_*(v)+|y|^2/2$.  The transformed Haar density $a(y,v)$ is analytic and
\[
 a(0,v)=J(x_*(v))/\sqrt{\det A(v)}>0.
\]
Choose a smooth radial cutoff in a fixed small $y$ ball, equal to one in a
smaller ball, uniformly in the coarse parameters.  Expand $a$ through degree
three in $y$, with a fourth-order remainder bounded together with its first
two $v$ derivatives.  Odd Gaussian moments vanish on this radial domain.
The quadratic term has order $\beta^{-1}$ relative to the leading Gaussian
integral; the fourth-order remainder has order $\beta^{-2}$.  Gaussian tails
outside the fixed ball are exponentially small.  On the remaining part of
the compact original fibre, the action gap gives the same exponential
control even after two $v$ derivatives, which introduce at most a quadratic
polynomial in $\beta$.

Consequently, uniformly with two coarse derivatives,
\[
 Z_{\beta,0,M}(v)
 =e^{-\beta E_*(v)}(2\pi/\beta)^{d_M/2}
 \bigl(a(0,v)+\beta^{-1}a_2(v)+O_{C^2,M}(\beta^{-2})\bigr).
\]
Taking the logarithm on the uniformly positive leading-amplitude branch
proves \eqref{eq:YMA2-Laplace-action}--\eqref{eq:YMA2-Laplace-correction}.

Differentiate the critical-point equation at zero.  With
$B=D_xD_vV_M(0,0)$ and $C=D_v^2V_M(0,0)$,
$x_*'(0)=-A_M^{-1}B$ and
$D^2E_*(0)=C-B^*A_M^{-1}B$.  This is exactly the inherited flat Wilson
Schur form \eqref{eq:YMA2-unit-reference}, independently of the section.
Taking two derivatives in the action expansion proves
\eqref{eq:YMA2-finite-Hessian-expansion}.  The physical axial symmetry and
\Cref{thm:YMA2-reference-symbol} give its scalar form.
\end{proof}

\begin{proposition}[Explicit finite formula for the first correction]
\label{prop:YMA2-one-loop-formula}
Let $L_a=x_*'(0)a=-A_M^{-1}Ba$.  Denote the first and mixed second total
coarse derivatives of $A(v)$ at zero by $A_a$ and $A_{ab}$, respectively;
these derivatives include the motion of $x_*(v)$.  In the Lie metric
\eqref{eq:YMA2-Lie-metric},
\begin{equation}
 \boxed{
 D^2\Gamma_M(0)[a,b]
 =\frac12\Tr\bigl(A_M^{-1}A_{ab}
 -A_M^{-1}A_bA_M^{-1}A_a\bigr)
 +\frac12\ip{L_a}{L_b}.}
 \label{eq:YMA2-one-loop-Hessian}
\end{equation}
This is a finite formula in the literal Wilson derivatives and Haar density.
It is not a claimed numerical evaluation of $\gamma_M(q)$.
\end{proposition}

\begin{proof}
Twice differentiate $\frac12\log\det A(v)$.  In local exponential
coordinates on $\SU(3)$,
\[
 \log J(x)-\log J(0)
 =\frac1{24}\sum_e\Tr_{\rm ad}(\operatorname{ad}x_e)^2
 +O(\|x\|^4)
 =-\frac14\sum_e\|x_e\|^2+O(\|x\|^4).
\]
Indeed the exponential differential has factor
$(1-e^{-\operatorname{ad}X})/\operatorname{ad}X$; its log determinant has
quadratic term $\Tr_{\rm ad}(\operatorname{ad}X)^2/24$ and zero linear
term.  For $\mathfrak{su}(3)$ the adjoint trace form is
$6\Tr(X^2)=-6\|X\|^2$.  Hence $D\log J(0)=0$ and
$-D^2\log J(0)=\frac12I$.  Applying the chain rule to
$-\log J(x_*(v))$ gives precisely the last term in
\eqref{eq:YMA2-one-loop-Hessian}; no $x_*''$ term survives from this
part because the first derivative of $\log J$ is zero.  The matrix
$A_{ab}$, in contrast, still contains all such minimizer-motion terms.
\end{proof}

\begin{corollary}[The leading connected covariance is evaluated, including its cross terms]
\label{cor:YMA2-leading-covariance}
For fixed $M$ and fixed nonzero axial $q$, under the pure Wilson law,
\begin{align}
 \beta^{-1}\mathbb E\widehat K^{\rm W}_\perp(q)
 &\longrightarrow\tfrac13(7I_3-\mathbf1_3\mathbf1_3^T)\otimes I_8,
 \label{eq:YMA2-K-limit}\\
 \beta^{-1}\widehat{\Cov(j)}_\perp(q)
 &\longrightarrow
 \left[\tfrac13(7I_3-\mathbf1_3\mathbf1_3^T)-s_0(q)I_3\right]\otimes I_8.
 \label{eq:YMA2-cov-limit}
\end{align}
In particular,
\begin{equation}
 \boxed{\beta^{-1}\mathbb E K^{\rm W}[v_*,v_*]\to2,
 \qquad \beta^{-1}c_{\beta,M}(q)\to2-s_0(q).}
 \label{eq:YMA2-scalar-limits}
\end{equation}
The mixed transverse covariance divided by $\beta$ tends to $-1/3$, not
zero.  Along a small-angle family the \emph{leading Gaussian} diagonal
return fraction is $1-s_0(q)/2=1-q^2/6+O(q^4)$.  The latter expression is
a reference calculation, not a uniform claim about the interacting
covariance when $M$ varies with $\beta$.
\end{corollary}

\begin{proof}
The uniform plaquette bound gives $w\to1$.  Substitute it in
\eqref{eq:YMA2-transverse-direct} and then subtract the physical Hessian
limit \eqref{eq:YMA2-finite-axial-limit}.  This proves every matrix and
scalar limit, including the off-diagonal entry.  The last expansion follows
from \Cref{cor:YMA2-reference-small-angle}.
\end{proof}

\begin{corollary}[Fixed-card positivity, not a cofinal margin]
\label{cor:YMA2-fixed-card-pass}
For every fixed $M$ and every $0<\delta<1$, the pure Wilson physical Hessian
on its nonzero axial blocks obeys
\begin{equation}
 h_{\beta,M}(q)\ge\delta\beta s_0(q)
 \label{eq:YMA2-fixed-card-floor}
\end{equation}
for all sufficiently large $\beta$.  More explicitly, if a constant $B_M$
bounds $|h_{\beta,M}(q)-\beta s_0(q)|$ over the finite nonzero axial bank
on the large-$\beta$ range under consideration, then
\begin{equation}
 \boxed{
 \left|\frac{h_{\beta,M}(q)}{\beta s_0(q)}-1\right|
 \le\frac{7B_M}{5\beta\sin^2(\pi/M)}.}
 \label{eq:YMA2-fixed-card-ratio-bound}
\end{equation}
A validated numerical threshold requires an actual bound for $B_M$; only its
fixed-card finiteness is proved by the asymptotic argument here.
\end{corollary}

\begin{proof}
There are finitely many nonzero axial angles at fixed $M$.  Their reference
values satisfy
$s_0(q)\ge\frac5{28}\omega(q)\ge\frac57\sin^2(\pi/M)>0$ by
\Cref{thm:YMA2-reference-symbol}.  The uniform-in-this-finite-bank
$O_M(1)$ Hessian remainder follows from
\Cref{thm:YMA2-finite-Laplace}.  Division gives the estimate and eventual
positivity.
\end{proof}

\begin{remark}[A fixed bounded generated action in the finite-card expansion]
If $\Psi$ is a fixed analytic function as $\beta\to\infty$ at fixed $M$,
repeat the displayed proof with transformed amplitude
$e^{-\Psi(x,v)}J(x)$.  The order-one correction becomes
$\Gamma_M(v)+\Psi(x_*(v),v)$ and the leading Schur coefficient is unchanged.
This does not cover an arbitrary response-selected family
$\Psi_{\beta,M}$: that family requires its own uniform derivative and
amplitude estimates.  It does not license suppressing an order-$\beta$
generated term.
\end{remark}

\section{The remaining physical row after the literal acquisitions}
\label{sec:YMA2-frontier}

The new Wilson-specific statements isolate rather than conceal the unresolved
quantity.  On the complete physical law, the deterministic Wilson diagonal
is given by a single mean plaquette with the explicit interval
\eqref{eq:YMA2-w-interval}.  Its raw mixed jets are known.  The flat axial
reference is \eqref{eq:YMA2-reference-rational}.  The remaining competition
is the connected covariance of the \emph{complete} score, plus the actual
generated second jet and the inherited topological subtraction.

To state the next target without changing that law, write
\begin{align}
 k_{\Psi,M}(q)&:=\mathbb E K^\Psi[v_*,v_*],
 &c_{\Phi,M}(q)&:=\Var(j_\Phi[v_*]),
 \label{eq:YMA2-next-two-rows}\\
 t_{\Phi,M}(q)&:=H_{\Phi,M}^{\rm top}[v_*,v_*].
 \label{eq:YMA2-top-row}
\end{align}
The exact local residual relative to the evaluated unit reference is
\begin{equation}
 \boxed{
 \begin{aligned}
 r_{\Phi,M}(q)
 &: =h_{\Phi,M}^{\rm loc}(q)-\beta s_0(q)\\
 &=2\beta(w_{\Phi,M}-1)+k_{\Psi,M}(q)\\
 &\quad-\bigl[c_{\Phi,M}(q)-\beta(2-s_0(q))\bigr]
 -t_{\Phi,M}(q).
 \end{aligned}}
 \label{eq:YMA2-fused-residual}
\end{equation}
This is one signed residual of the same source.  The covariance term is not
priced a second time because it appears in this formula and in the original
Hessian identity.

For a predeclared $0<\delta<1$, the desired local axial floor is now exactly
\begin{equation}
 \boxed{r_{\Phi,M}(q)\ge-(1-\delta)\beta s_0(q)}
 \label{eq:YMA2-final-margin-row}
\end{equation}
on the predeclared retained angles.  A two-sided sufficient route is to bound
the \emph{fused local} residual divided by $\omega(q)$ by less than
$(1-\delta)\beta\cdot5/28$.  V1 already supplies the appropriate weighted
second-displacement mechanism for that task; no new version of that theorem
is introduced here.

\begin{firewall}[Why the new all-volume bound is not yet the shell margin]
\label{fire:YMA2-uniformity-boundary}
The uniform plaquette interval controls the deterministic Wilson row, not the
source-relative dispersion of its covariance.  Even the exact vertical
Hessian floor $1/60$ is not a retained-field floor.  The fixed-$M$ expansion
has constants $C_M,B_M$ whose growth has not been bounded.  In particular,
$O_M(1)$ in an absolute Hessian norm need not be small relative to
$\beta\omega(q)$ when the smallest lattice angle tends to zero.  Nor has the
actual generated interaction or the topological packet been set to zero.
Consequently \eqref{eq:YMA2-final-margin-row} has not been proved on the
selected cofinal response-matched trajectory.
\end{firewall}

\subsection{Two technical qualifications retained for future use of V1}

First, $\Wtwo$ in \Cref{thm:YMA1-W2-control} gives a sharp \emph{universal
operator bound} from the weighted coefficient space to continuous Ward
quotients.  It is not equal to the supremum norm of each individual signed
quotient.  For example, $\alpha_1=4$, $\alpha_2=-1$ gives
$g(q)=-2+2\cos q$, with $\|g\|_\infty=4$ but
$\|\alpha\|_{\Wtwo}=8$.  All V1 inequalities using the upper bound remain
valid; an equality for an individual kernel is not used.

Second, a quotient enclosure with an uncertain positive denominator must
respect the numerator sign.  If $L_H\le G_H\le U_H$ and
$0<L_S\le G_S\le U_S$, then the unconditional pointwise enclosure is
\begin{equation}
 \min\left\{\frac{L_H}{L_S},\frac{L_H}{U_S},
                 \frac{U_H}{L_S},\frac{U_H}{U_S}\right\}
 \le\frac{G_H}{G_S}\le
 \max\left\{\frac{L_H}{L_S},\frac{L_H}{U_S},
                 \frac{U_H}{L_S},\frac{U_H}{U_S}\right\}.
 \label{eq:YMA2-sign-safe-ratio}
\end{equation}
This follows by monotonicity on the positive-denominator rectangle, with a
possible sign change of the numerator.  The simpler upper expression
$U_H/L_S$ used in \Cref{cor:YMA1-SDP-margin} is valid on the nonnegative
numerator branch, in particular when the stated nonnegative feasibility
set is nonempty.  Outside that branch use
\eqref{eq:YMA2-sign-safe-ratio}; the direct signed tests for
$H-\delta S$ require no quotient and are unchanged.

\section{Exact replay, proof-status ledger, and the V3 task}
\label{sec:YMA2-terminal}

The accompanying replay is deterministic and uses exact integer and rational
symbolic arithmetic for its certificate rows.  It verifies the sixty-four
word-pair boundaries, the forty-five nontrivial fillings, every entry of the
weighted occupancy table, the complete $4\times4$ Laurent direct-Hessian
symbol, the rational axial reference, its two endpoint values, and the
strict sign of its derivative.  Its terminal identities are
\begin{equation}
 \boxed{\#\mathscr E^{\rm red}_{\rm types}=45,
 \quad\sum m_e=96,\quad\sum m_e^2=272,
 \quad\max_p\sum_{e:p\in\Sigma_e}m_e=20.}
 \label{eq:YMA2-replay-identities}
\end{equation}
It does not sample Wilson fields, numerically estimate a covariance, or
certify a cofinal physical margin.

\begingroup\small
\begin{longtable}{>{\raggedright\arraybackslash}p{0.28\textwidth}
                  >{\raggedright\arraybackslash}p{0.20\textwidth}
                  >{\raggedright\arraybackslash}p{0.44\textwidth}}
\caption{Version V2 proof-status ledger.}\label{tab:YMA2-ledger}\\
\toprule
\textbf{Row}&\textbf{Status}&\textbf{Exact scope}\\
\midrule
\endfirsthead
\toprule
\textbf{Row}&\textbf{Status}&\textbf{Exact scope}\\
\midrule
\endhead
Parity-tree non-Abelian filling&Proved and exactly replayed
&Forty-five reduced-link types and weighted occupancy twenty.\\[0.3em]
Pinned bare Wilson action&Proved, cutoff independent
&$D/120\le V\le4D$ on the full compact fibre; unique flat minimum.\\[0.3em]
Vertical flat Hessian&Proved, cutoff independent
&$I/60\preceq A_M\preceq8I$; this is not coarse coercivity.\\[0.3em]
Corner plaquette&Explicit physical interval
&For $\beta\ge240$, $w\ge1-C_{\rm W}/\beta$, plus the displayed
oscillation-density term for a symmetric generated action.\\[0.3em]
Deterministic Wilson symbol&Exactly evaluated
&$\mathbb E\widehat K^{\rm W}=\beta w(7I-uu^*)/3$ in the frozen section.\\[0.3em]
Separate raw scalarization&Corrected
&Fails on the literal parity-tree Wilson jets; their physical difference
still scalarizes.\\[0.3em]
Four-dimensional axial reference&Exactly evaluated
&The complete sixteen-alias sum gives
$s_0=t(t+8)(t+24)/(9(t^2+18t+16))$.\\[0.3em]
Actual leading covariance&Proved at each fixed card
&$\Cov(j)_\perp/\beta\to\mathsf C_\perp-s_0I$, including mixed entry
$-1/3$.\\[0.3em]
Finite-card weak Wilson margin&Proved asymptotically
&$h/(\beta s_0)\to1$ at fixed $M$; no validated finite $\beta$ threshold
without $B_M$.\\[0.3em]
One-loop coefficient&Finite exact derivative formula
&Its displayed trace/Haar expression is not yet numerically evaluated.\\[0.3em]
Selected response-matched cofinal margin&Open
&Requires the fused covariance/generated/topological residual in
\eqref{eq:YMA2-fused-residual}.\\[0.3em]
Weak-facing escape, higher jets, and mass&Not supplied
&No trajectory, physical-time contraction, or continuum mass conclusion is
inferred.\\
\bottomrule
\end{longtable}
\endgroup

\begin{openproblem}[V3: the connected remainder on the same first axial card]
\label{op:YMA2-V3}
Preserve V2 and all inherited choices.  On the actual selected shell retain
its submitted $\beta$, full generated action $\Psi$, and local/topological
split.  Use the evaluated $s_0(q_{n_*})$ and the Wilson deterministic row
above.  The next acquisition is the \emph{fused} scalar
\eqref{eq:YMA2-fused-residual}, first at the already frozen axial angle and
then in a Ward-relative displacement norm only if an all-angle or cofinal
statement is needed.

The useful next analytic route is a source-specific bound for the connected
score remainder after the exact pinned-fibre Gaussian return has been
subtracted.  The new vertical gap and explicit plaquette interval are
available inputs, but they are not substituted for that estimate.  On a
failing direct mode, attribute the surviving scalar to its generated,
connected, or topological term without double charging.  Do not open another
abstract positivity, generic shell, or full-gap construction before this
same-law scalar is decided.
\end{openproblem}

\begin{assessment}[V2 termination]
\label{ass:YMA2-termination}
The literal bare Wilson deterministic row and flat reference are
\pass\ at their stated scopes.  The proposed \emph{separate} raw-jet
scalarization has an explicit typed kinematic obstruction and is replaced
by the correct difference scalarization.  For the complete selected cofinal
symbol the current verdict remains \stopstatus: its connected/generated/
topological residual has not yet been bounded in the required relative
norm.  This is a named remaining physical estimate, not an assertion of a
negative Yang--Mills symbol or a mass-gap counterexample.
\end{assessment}

\begin{thebibliography}{9}
\bibitem{YMA2-Katsevich}
A.~Katsevich,
\emph{The Laplace asymptotic expansion in high dimensions},
arXiv:2406.12706v2 (2024).
Used only for external context on dimension-dependent Laplace remainders;
the finite Wilson asymptotic statement and its proof are supplied here.
\end{thebibliography}

\section*{Version V2 append-only continuation record}
\addcontentsline{toc}{section}{Version V2 append-only continuation record}
The entire preceding V1 source is retained verbatim up to its single final
\verb|\end{document}| command.  Its SHA--256 digest is
\begin{center}\scriptsize\ttfamily
4a1403b3a962506dc6a10bef9e87327d284c6f685454a93252d5a53710aaa9bd.
\end{center}
The corrections in this continuation are explicit additions, not retroactive
changes.  For V3, preserve this complete source, append immediately before
the final document-closing command, and save with suffix
\texttt{\_V3.tex}.  Keep the same physical response, block map, parity-tree
section, source normalization, first axial angle, and local/topological
support convention.  The next link should point to that full accumulated
source, not to an isolated addendum.

% =============================================================================
% END APPENDED VERSION V2 MATERIAL
% =============================================================================


% =============================================================================
% START APPENDED VERSION V3 MATERIAL
% =============================================================================
\clearpage
\hypersetup{
 pdftitle={Renewal Yang--Mills Reduced-Fibre Wilson Shell Symbol and Local-Vacuum Programme V3},
 pdfsubject={Haar harmonic compensation, an exact first-axial one-loop interval, and a volume-uniform Ward-relative one-loop bound}
}
\providecommand{\HS}{\mathrm{HS}}
\providecommand{\op}{\mathrm{op}}

\section{Version V3: the connected remainder is evaluated to its first nonzero order}
\label{sec:YMA3-direction}

The entire V2 source is preserved above.  Its exact Wilson diagonal and
sixteen-alias reference are imported, as are its global pinned-fibre bound
and fixed-card two-derivative Laplace theorem.  This continuation does not
reprove the general score--Hessian identity, the axial representation theory,
the V1 Ward-tail theorem, or the higher-jet/local-vacuum handoff.

The new calculation has three outputs.  First, a literal product of Haar
translations, determined by the already fixed flat harmonic lift, removes
the linear Gaussian part of the Wilson score \emph{before} its covariance
is taken.  The physical partition function and the full generated action
are unchanged.  Second, the first nonzero covariance after this compensation
and the accompanying deterministic correction are computed on the
$M=4$, $q=\pi/2$ axial calibration card.  Their difference satisfies the
rigorous interval
\begin{equation}
 \boxed{-1.089801<\gamma_4(\pi/2)<-1.089787.}
 \label{eq:YMA3-main-numerical-interval}
\end{equation}
Third, the entire family of pure-Wilson one-loop coefficients has a
cutoff-independent Ward-relative bound on the nonzero finite momentum grids.
The full finite-torus harmonic coefficient is exponentially small in $M$
at this order, rather than silently set to zero.

The calibration card is the smallest side $M\ge4$ already admitted in V2,
with its first nonzero axial angle, a reference value already evaluated in
V2.  It is \emph{not} identified with an unspecified selected physical shell
or with a retained high mode whose buffer has not been submitted.  All
formulas below retain the original $q_{n_*}$ as their physical input; the
calibration evaluates one member of that same explicit family.

\begin{firewall}[The order-one coefficient is not the full finite-coupling answer]
\label{fire:YMA3-scope}
The number $\gamma_M(q)$ is the coefficient defined by
\Cref{thm:YMA2-finite-Laplace}:
\[
 h_{\beta,M}(q)=\beta s_0(q)+\gamma_M(q)+O_M(\beta^{-1}).
\]
An interval for $\gamma_M(q)$ is not an interval for $h_{\beta,M}(q)$ at
an arbitrary finite $\beta$.  A uniform bound on this coefficient does not
make the remainder uniform.  The response-selected generated action,
local/topological incidence, weak-facing-root escape, and higher nonlinear
remainder remain separate rows.  No transfer gap or continuum mass is inferred.
\end{firewall}

\section{A literal Haar compensation of the flat Gaussian return}
\label{sec:YMA3-compensation}

\subsection{The fixed incidence matrices}

Use one real colour component in the metric
\eqref{eq:YMA2-Lie-metric}.  Positive fine-edge types are $(s,\mu)$ with
$s\in\{0,1\}^4$ and $1\le\mu\le4$.  Write $\mathcal R$ for the
forty-five types with $m_{s,\mu}>0$ in
\eqref{eq:YMA2-filling-length}; the four second-constituent types are
$(\widehat e_\mu,\mu)$ and the remaining fifteen are forest types.
For a fine plaquette based at $2y+s$ in directions $\mu<\nu$, retain the
ordered list
\begin{equation}
 \bigl((2y+s,\mu),+1\bigr),\quad
 \bigl((2y+s+\widehat e_\mu,\nu),+1\bigr),\quad
 \bigl((2y+s+\widehat e_\nu,\mu),-1\bigr),\quad
 \bigl((2y+s,\nu),-1\bigr).
 \label{eq:YMA3-plaquette-list}
\end{equation}
Each positive edge in this list has a type $e_i$ and a coarse-cell offset
$d_i$, with $|d_i|_1\le1$.  Its sign is denoted $\epsilon_i$.
The $96\times64$ incidence symbol is explicitly
\begin{equation}
 D(k)_{p,e}=\sum_{i:e_i=e}\epsilon_i e^{ik\cdot d_i}.
 \label{eq:YMA3-incidence-symbol}
\end{equation}
Let $D_R$ and $D_C$ be its reduced and second-constituent column blocks.
On the corresponding finite lattice, or on finitely supported infinite
fields, put
\begin{equation}
 A=\tfrac13D_R^*D_R,\qquad
 B=\tfrac13D_R^*D_C,\qquad
 C=\tfrac13D_C^*D_C,\qquad G=A^{-1},\qquad L=-GB.
 \label{eq:YMA3-ABCL}
\end{equation}
The scalar-coordinate $A$ is the colour factor of the V2 vertical Hessian.
In particular, $I/60\preceq A\preceq8I$.

For a real coarse scalar field $a$, define a real fine field $b=b[a]$ by
\begin{equation}
 \boxed{b|_{\mathcal R}=La,\qquad b|_C=a,\qquad b|_{\mathscr T}=0.}
 \label{eq:YMA3-fine-lift}
\end{equation}
Thus
\begin{equation}
 D_R^*Db=0,\qquad \tfrac13\|Db\|_2^2
 =\langle a,\mathsf S_Ma\rangle.
 \label{eq:YMA3-harmonic-equations}
\end{equation}
This is the existing flat Schur minimizer written on the literal section,
not a new coarse-graining map.

\begin{lemma}[Size of the compensated fine direction]
\label{lem:YMA3-lift-size}
For arbitrary $a$,
\[
 \|b[a]\|_2^2\le141\|a\|_2^2.
\]
For a field in one transverse axial polarization, the sharper bound is
\begin{equation}
 \boxed{\|b[a]\|_2^2\le121\|a\|_2^2.}
 \label{eq:YMA3-axial-lift-size}
\end{equation}
\end{lemma}
\begin{proof}
The minimization equation gives
$\langle La,A La\rangle=\langle a,(C-\mathsf S_M)a\rangle$.
The lower bound on $A$ and positivity of the Schur form yield
$\|La\|^2\le60\langle a,Ca\rangle$.  V2 evaluates the maximum eigenvalue
of $C$ as $7/3$, and its diagonal on one transverse axial polarization as
$2$.  The reduced and second-constituent supports are disjoint, so
$\|b\|^2=\|La\|^2+\|a\|^2$.  These facts give both constants.
\end{proof}

Fix a real $X\in\mathfrak{su}(3)$ with $-\Tr X^2=1$ and use the same
colour $X$ on every nonzero component of $a$.  For a real parameter $t$,
let
\begin{equation}
 W_{y,\mu}(t)=e^{t a_{y,\mu}X},\qquad
 (T_tz)_e=e^{t b_eX}z_e\quad(e\in\mathcal R).
 \label{eq:YMA3-Haar-shift}
\end{equation}
Every $T_t$ is a product of left Haar translations.  The compensated action
is
\begin{equation}
 \widetilde\Phi_b(z,t)
 =\beta V_M(T_tz,W(t))+\Psi(T_tz,W(t)).
 \label{eq:YMA3-shifted-action}
\end{equation}
Denote its first two $t$ derivatives at zero by
$\widetilde j_b$ and $\widetilde K_b$, and write
$\widetilde j_b^V,\widetilde K_b^V$ for those of the \emph{unit} Wilson
action.  The corresponding generated derivatives are
$\widetilde j_b^\Psi,\widetilde K_b^\Psi$.

\begin{theorem}[Same-law compensation and disappearance of the Gaussian score]
\label{thm:YMA3-Haar-compensation}
At every finite $\beta>0$, with the full submitted $\Psi$ retained,
\begin{equation}
 \int e^{-\widetilde\Phi_b(z,t)}\,dm_{\rm red}(z)
 =Z_{\beta,\Psi,M}(W(t)).
 \label{eq:YMA3-exact-same-partition}
\end{equation}
Consequently the physical directional Hessian is exactly
\begin{equation}
 \boxed{H_\Phi[aX,aX]
 =\mathbb E_\Phi\widetilde K_b-\Var_\Phi(\widetilde j_b).}
 \label{eq:YMA3-compensated-Hessian}
\end{equation}
The expectation is under the original law at $W=I$, since $T_0$ is the
identity.  If that law and action have the stated global-conjugation
symmetry, $\mathbb E_\Phi\widetilde j_b=0$.

In product logarithms $z_e=e^{x_e}$, the pure Wilson terms satisfy
\begin{equation}
 \boxed{
 \widetilde j_b^V(0)=0,\quad D_x\widetilde j_b^V(0)=0,\quad
 \widetilde K_b^V(0)=\langle a,\mathsf S_Ma\rangle,\quad
 D_x\widetilde K_b^V(0)=0.}
 \label{eq:YMA3-score-cancellation}
\end{equation}
The first nonconstant term of the compensated Wilson score is therefore
quadratic in the fine fluctuations, rather than linear.
\end{theorem}
\begin{proof}
Change variables $z'=T_tz$ using Haar invariance.  This proves
\eqref{eq:YMA3-exact-same-partition} with no Jacobian.  Differentiating the
same positive compact integral gives the already established logarithmic
Hessian identity in the new coordinates.  Global conjugation has no
nonzero invariant linear functional on the adjoint colour tangent, so its
integrated first derivative is zero.

At $x=0$, all fine backgrounds $e^{t b_eX}$ commute.  Set
$f_p=\sum_{i=1}^4\epsilon_i b_{e_i}$ for one plaquette.  A first insertion
$Y\in\mathfrak{su}(3)$ on a reduced edge of that plaquette has derivative
\[
 -\frac{\epsilon_i}{3}\Re\Tr(e^{t f_pX}Y).
\]
Conjugations by the commuting suffix have disappeared from this trace.
Its coefficient of $t$ is the scalar plaquette incidence applied to $b$,
and its coefficient of $t^2$ is zero because
$\Re\Tr(X^2Y)=0$: $X^2$ is Hermitian and $Y$ is anti-Hermitian.
After summing the first coefficient over incident plaquettes,
$D_R^*Db=0$ makes it vanish as well.  Thus the vertical gradient along the
shifted background is $O(t^3)$.  Differentiating it once and twice at zero
gives the second and fourth identities in
\eqref{eq:YMA3-score-cancellation}.  The first follows from stationarity
of each flat plaquette.  Finally,
$-\Re\Tr(e^{t f_pX})/3$ has second derivative $f_p^2/3$ at zero;
summing gives the third identity by
\eqref{eq:YMA3-harmonic-equations}.
\end{proof}

\begin{remark}[A change of integration coordinates, not of physical source]
The lift is fixed by $A,B$ before the Wilson expectation is evaluated.  No
response value is used to fit it.  Equation
\eqref{eq:YMA3-exact-same-partition} makes the new score a derivative of
the \emph{same} unnormalized cylinder.  It is not a new source birth and
its covariance is not added to the old raw covariance.  The two complete
Hessian expressions are equal; their two separate terms need not be equal.
This respects, rather than reverses, the V2 correction of separate raw-jet
scalarization.
\end{remark}

\section{A finite-coupling deterministic bound after compensation}
\label{sec:YMA3-deterministic-bound}

Let
\begin{equation}
 B_p(b)=\sum_{i=1}^4 b_{e_i}^2,\qquad
 D_p(z)=\sum_{i:e_i\in\mathcal R}\|z_{e_i}-I\|_F^2,
 \qquad
 \eta_{\Psi,M}=C_{\rm W}+\frac{\operatorname{osc}\Psi_M}{12M^4}.
 \label{eq:YMA3-local-stocks}
\end{equation}
Edges and coefficients here are the translated literal ones in the
plaquette, not independent copies.

\begin{proposition}[Pointwise second-order Wilson error]
\label{prop:YMA3-pointwise-second-order}
For every reduced configuration, including configurations outside any
small-field chart,
\begin{equation}
 \boxed{
 \left|\widetilde K_b^V(z)-\langle a,\mathsf S_Ma\rangle\right|
 \le2\pi^2\sum_p B_p(b)D_p(z).}
 \label{eq:YMA3-pointwise-K-error}
\end{equation}
\end{proposition}
\begin{proof}
For each reduced unitary choose a principal anti-Hermitian logarithm $Y_e$
with eigenangles in $[-\pi,\pi]$.  It need not be traceless.  This causes
no problem: the unitary word and its real trace extend to $U(3)$, and the
identity $\Re\Tr(X^2Y)=0$ also holds for these logarithms.  Spectral
calculus gives
$\|Y_e\|_F\le(\pi/2)\|z_e-I\|_F$.

Interpolate the at most four reduced factors of one plaquette by
$e^{sY_e}$, $0\le s\le1$, keeping the $t$ factors of
\eqref{eq:YMA3-Haar-shift}.  Two $t$ derivatives and two $s$ derivatives
of its normalized real trace have absolute value at most
\[
 \left(\sum_i|b_{e_i}|\right)^2
 \left(\sum_{i:e_i\in\mathcal R}\|Y_{e_i}\|_F\right)^2.
\]
Indeed the product rule distributes the inserted matrices among unitary
factors, the sums displayed count those distributions, and
$|\Tr Z|/3\le\|Z\|_{\op}$ with $\|X\|_{\op}\le1$ bounds each term.
The first $s$ derivative at zero is a sum of traces containing two copies
of $X$ and one anti-Hermitian $Y$ and hence vanishes, by cyclicity and
$\Re\Tr(X^2Y)=0$.
Taylor's integral remainder and Cauchy--Schwarz now bound the error of this
plaquette by
\[
 \tfrac12(4B_p)\left(4\sum_i\|Y_{e_i}\|_F^2\right)
 \le2\pi^2 B_p D_p.
\]
Sum over all plaquettes and use the flat value in
\eqref{eq:YMA3-score-cancellation}.
\end{proof}

\begin{theorem}[An actual volume-uniform compensated Wilson mean]
\label{thm:YMA3-deterministic-uniform}
Assume coarse-translation invariance of the full law and $\beta\ge240$.
Then
\begin{equation}
 \boxed{
 \left|\beta\mathbb E_\Phi\widetilde K_b^V
       -\beta\langle a,\mathsf S_Ma\rangle\right|
 \le138240\pi^2\eta_{\Psi,M}\,\|b[a]\|_2^2.}
 \label{eq:YMA3-deterministic-mean-error}
\end{equation}
In particular the right side is independent of $M$ for a unit axial
polarization when $\operatorname{osc}\Psi_M/M^4$ is uniformly bounded.
It is an absolute, not a Ward-relative, estimate.
\end{theorem}
\begin{proof}
The pressure comparison in the proof of
\Cref{thm:YMA2-plaquette-interval} gives
$\mathbb E_\Phi V_M\le24\eta_{\Psi,M}M^4/\beta$.
The global filling inequality gives
$\mathbb E_\Phi D\le2880\eta_{\Psi,M}M^4/\beta$.
Every one of the forty-five edge-type expectations is repeated $M^4$ times
by translation invariance and is nonnegative.  Hence each such expectation
is at most $2880\eta_{\Psi,M}/\beta$ and
$\mathbb E D_p\le11520\eta_{\Psi,M}/\beta$.
Every fine edge belongs to six plaquettes, so
$\sum_p B_p=6\|b\|^2$.  Apply
\eqref{eq:YMA3-pointwise-K-error} and multiply by $\beta$.
\end{proof}

\begin{firewall}[The surviving covariance is not bounded by the energy comparison]
\label{fire:YMA3-no-Poincare-from-pinning}
Global quadratic pinning and the preceding first-moment bound do not prove
a uniform Poincar\'e inequality or control the spatial covariance of the
quadratic compensated score.  They have removed the order-$\beta$
deterministic cancellation and bounded its mean error.  They have not
supplied the remaining same-law variance in
\eqref{eq:YMA3-compensated-Hessian}.
\end{firewall}

\section{The literal colour-reduced one-loop vertices}
\label{sec:YMA3-vertices}

The next computation evaluates that remaining variance to its first
nonzero order.  It uses the ordered four-edge word, including the different
positions of the fluctuation on a positive and a negative edge.
For a real $b$, on each plaquette define
\begin{equation}
 f=\sum_{i=1}^4\epsilon_i b_{e_i},\qquad
 s_i=\sum_{j>i}\epsilon_j b_{e_j}
          +\mathbf1_{\{\epsilon_i=-1\}}\epsilon_i b_{e_i},\qquad
 d_{ij}=s_i-s_j-\frac f2\quad(i<j).
 \label{eq:YMA3-suffix-data}
\end{equation}
Only rows and columns associated with reduced edges are retained in the
matrices below.  Repeated contributions are added on their actual lattice
indices.

Define a real skew matrix $R(b)$ and a real symmetric matrix $Q(b)$ by
summing the following plaquette entries:
\begin{align}
 R_{e_i e_j}^{(p)}&=\epsilon_i\epsilon_j d_{ij},&
 R_{e_j e_i}^{(p)}&=-\epsilon_i\epsilon_j d_{ij}\quad(i<j),
 \label{eq:YMA3-R-vertex}\\
 Q_{e_i e_i}^{(p)}&=-\frac89 f^2,&
 Q_{e_i e_j}^{(p)}&=Q_{e_j e_i}^{(p)}
 =\epsilon_i\epsilon_j\left(-2d_{ij}^2-\frac7{18}f^2\right)
 \quad(i<j).
 \label{eq:YMA3-Q-vertex}
\end{align}
The signs in $Q$ are not replaced by their absolute values in any physical
calculation.

\begin{theorem}[Exact $SU(3)$ vertex contraction]
\label{thm:YMA3-colour-vertices}
Let $M_b(t)$ be the vertical logarithmic-coordinate Hessian of the unit
Wilson action at the shifted background, with the new fluctuation set to
zero.  Then $M_b(0)=A\otimes I_8$ and
\begin{equation}
 \boxed{M_b'(0)=\tfrac13 R(b)\otimes J_X,\qquad
        \Tr_{\rm colour}M_b''(0)=Q(b),}
 \label{eq:YMA3-M-derivatives}
\end{equation}
where $J_X=\operatorname{ad}X$ in an orthonormal real adjoint basis and
$\Tr J_X^2=-6$.  In particular, the complete colour trace is performed
by the rational coefficients in
\eqref{eq:YMA3-R-vertex}--\eqref{eq:YMA3-Q-vertex}; no choice of numerical
Gell--Mann matrices remains in the finite integration panel.
\end{theorem}
\begin{proof}
Write $T_1,\ldots,T_8$ for an orthonormal anti-Hermitian basis.  Move all
commuting background factors to the left of the plaquette word.  The
fluctuation at position $i$ becomes
$\epsilon_i\operatorname{Ad}_{e^{-t s_iX}}x_{e_i}$.  The extra term in
$s_i$ for a negative edge is essential: its inverse word is
$e^{-x_{e_i}}e^{-t b_{e_i}X}$, not the reverse order.

For $i<j$, the bilinear colour matrix before suffix rotations is
\[
 T_f(t)_{ab}=-\tfrac13\Re\Tr(e^{t fX}T_aT_b).
\]
It commutes with $J_X$.  Since
$\Re\Tr(XT_aT_b)=\frac12\langle T_a,[X,T_b]\rangle$,
\[
 T_f(0)=\tfrac13I,\qquad
 T_f'(0)=-\tfrac f6J_X,\qquad
 T_f''(0)=\tfrac{f^2}{3}P_X,
 \quad (P_X)_{ab}=-\Re\Tr(X^2T_aT_b).
\]
The suffix rotations give the cross block
$\epsilon_i\epsilon_j e^{t(s_i-s_j)J_X}T_f(t)$.
Its first derivative is
$\epsilon_i\epsilon_j d_{ij}J_X/3$.
Its second derivative is
\[
 \frac{\epsilon_i\epsilon_j}{3}
 \left(((s_i-s_j)^2-f(s_i-s_j))J_X^2+f^2P_X\right).
\]
The reverse block is its transpose.  A diagonal block has second derivative
$f^2P_X/3$ and zero first derivative.

The elementary traceless-matrix completeness relation gives
$\sum_aT_a^2=-(8/3)I_3$.  For clarity, take the six normalized off-diagonal
anti-Hermitian matrix units and the two normalized diagonal traceless
matrices; summing their squares gives this identity entry by entry.
It follows that
$\Tr P_X=(8/3)\Tr X^2=-8/3$.
Also the adjoint trace form on traceless $3\times3$ matrices is
$\Tr_{\rm ad}(\operatorname{ad}X)^2=6\Tr X^2=-6$; this follows by
applying $Y\mapsto XY-YX$ on all matrix units, where the scalar subspace
contributes zero.  Taking the colour trace of the displayed second
cross derivative therefore gives
\[
 \epsilon_i\epsilon_j
 \left[-2\bigl((s_i-s_j)^2-f(s_i-s_j)\bigr)-\tfrac89 f^2\right]
 =\epsilon_i\epsilon_j\left(-2d_{ij}^2-\tfrac7{18}f^2\right).
\]
The diagonal trace is $-8f^2/9$.  This proves every entry asserted.
\end{proof}

\begin{theorem}[The first nonzero compensated covariance and deterministic correction]
\label{thm:YMA3-one-loop-limits}
For fixed finite $M$ and fixed real one-colour $a$, define
\begin{equation}
 \mathcal T_M[a]=\tfrac12\Tr(GQ(b[a])),\qquad
 \mathcal B_M[a]= -\tfrac13\Tr(GR(b[a])GR(b[a]))\ge0.
 \label{eq:YMA3-tadpole-bubble}
\end{equation}
Under the actual pure Wilson law,
\begin{align}
 \lim_{\beta\to\infty}
 \left(\beta\mathbb E\widetilde K_b^V
              -\beta\langle a,\mathsf S_Ma\rangle\right)
 &=\mathcal T_M[a],
 \label{eq:YMA3-tadpole-limit}\\
 \lim_{\beta\to\infty}\Var(\beta\widetilde j_b^V)
 &=\mathcal B_M[a].
 \label{eq:YMA3-bubble-limit}
\end{align}
The one-loop Hessian coefficient already defined in V2 is exactly
\begin{equation}
 \boxed{D^2\Gamma_M(0)[aX,aX]
 =\mathcal T_M[a]-\mathcal B_M[a]
 =\tfrac12\Tr(GQ(b[a]))+\tfrac13\Tr(GR(b[a])GR(b[a])).}
 \label{eq:YMA3-exact-one-loop}
\end{equation}
\end{theorem}
\begin{proof}
The cancellation theorem shows that the vertical gradient of the shifted
background is $O(t^3)$.  The implicit-function theorem and the positive
vertical Hessian imply that the minimizing new fluctuation is $O(t^3)$.
Therefore its first two derivatives at zero vanish.  In these coordinates
the product Haar density has no $t$ dependence.  Two derivatives of the
Laplace leading amplitude thus give
\[
 D^2\Gamma_M(0)[aX,aX]
 =\tfrac12\Tr\left((G\otimes I_8)M_b''(0)
 -(G\otimes I_8)M_b'(0)(G\otimes I_8)M_b'(0)\right).
\]
Substitute \eqref{eq:YMA3-M-derivatives} and $\Tr J_X^2=-6$ to obtain
\eqref{eq:YMA3-exact-one-loop}.  The Haar/minimizer terms present in V2's
original logarithm coordinates have not been omitted: their coordinate
change is precisely the Haar translation used here, and the new minimizing
fluctuation starts only at order $t^3$.

For the two separate limits, the fixed-card Laplace proof in V2 gives
convergence, with every fixed polynomial moment, of $\sqrt\beta\,x$ to
the centered Gaussian with covariance $G\otimes I_8$.  Its exponential
local majorant and the positive action gap outside the chart also control
the finitely many polynomially growing score insertions used here.
By \eqref{eq:YMA3-score-cancellation}, the rescaled score converges in
second moment to $\frac12\xi^TM_b'(0)\xi$, and the compensated second jet
minus its flat value converges in first moment to
$\frac12\xi^TM_b''(0)\xi$.  For a centered Gaussian of covariance $C$,
differentiating $\exp(u^TCu/2)$ four times gives
\[
 \Var\left(\tfrac12\xi^TD\xi\right)
 =\tfrac12\Tr(CDCD).
\]
These identities give \eqref{eq:YMA3-tadpole-limit} and
\eqref{eq:YMA3-bubble-limit}.  Finally,
$G^{1/2}RG^{1/2}$ is real skew, so
\[
 \Tr(GRGR)=-\|G^{1/2}RG^{1/2}\|_{\HS}^2\le0.
\]
\end{proof}

\begin{remark}[A fixed generated action versus a selected generated family]
For a fixed analytic $\Psi$ at fixed $M$, the leading fluctuation law is
unchanged, and the order-one correction receives the additional
$D_t^2\Psi(T_tI,W(t))|_{t=0}$, equivalently the V2 second derivative of
$\Psi$ along the Wilson minimizer.  For a response-selected
$\Psi_{\beta,M}$ none of these uniformity properties is automatic.  The
exact compensated formula \eqref{eq:YMA3-compensated-Hessian}, with all
mixed scores retained, remains the applicable statement for that family.
\end{remark}

\section{An exact first-axial acquisition on the \texorpdfstring{$M=4$}{M=4} Wilson card}
\label{sec:YMA3-calibration}

\subsection{The harmonic vector is rational over the fourth roots of unity}

Take $M=4$, $q=(\pi/2,0,0,0)$, transverse direction $e_2$, and the real
coarse field
\begin{equation}
 a_{y,2}=\sqrt{\frac{2}{4^4}}\cos(\pi y_1/2),\qquad
 a_{y,\mu}=0\quad(\mu\ne2).
 \label{eq:YMA3-real-calibration}
\end{equation}
Its norm is one.  One may take
$X=i\operatorname{diag}(1,-1,0)/\sqrt2$; the colour trace theorem makes
the answer independent of this particular unit adjoint polarization.
Let $\ell$ be the fine harmonic amplitude with coarse complex amplitude
$e_2$ before the real normalization is applied.  The following table gives
\emph{all} its entries, with common denominator $1212$.

\begin{center}
\small
\begin{tabular}{c|rrrr}
\toprule
$s$&$1212\ell_{s,1}$&$1212\ell_{s,2}$&$1212\ell_{s,3}$&$1212\ell_{s,4}$\\
\midrule
0000&0&0&0&0\\
0001&0&0&0&0\\
0010&0&0&0&0\\
0011&0&0&0&0\\
0100&0&1212&232&232\\
0101&0&748&62&-232\\
0110&0&748&-232&62\\
0111&0&624&-62&-62\\
1000&0&$-369+237i$&0&0\\
1001&0&$-170+204i$&0&0\\
1010&0&$-170+204i$&0&0\\
1011&0&$-124+188i$&0&0\\
1100&$606-606i$&$843+237i$&232&232\\
1101&$374-374i$&$578+204i$&62&-232\\
1110&$374-374i$&$578+204i$&-232&62\\
1111&$312-312i$&$500+188i$&-62&-62\\
\bottomrule
\end{tabular}
\end{center}

\begin{lemma}[Exact calibration vector and metric]
\label{lem:YMA3-exact-calibration-lift}
The displayed vector has the required forest and second-constituent values
and satisfies
\begin{equation}
 \boxed{D_R(q)^*D(q)\ell=0,\qquad
  \tfrac13\|D(q)\ell\|^2=\frac{833}{1818},\qquad
  \|\ell\|^2=\frac{592915}{122412}.}
 \label{eq:YMA3-lift-exact-checks}
\end{equation}
It is the unique harmonic lift of the specified coarse amplitude.  The last
number is also $\|b[a]\|^2$ for the normalized real field
\eqref{eq:YMA3-real-calibration}.
\end{lemma}
\begin{proof}
Insert each of the sixteen parity rows in
\eqref{eq:YMA3-incidence-symbol}.  All phases are powers of $i$, so the
first equation is an equality of forty-five Gaussian-integer numerators
with denominator $1212$.  They are all zero.  The sums of squared
plaquette and edge numerators give respectively the second and third
fractions.  These finite integer identities are also checked directly by
the accompanying verifier, which reconstructs the incidence rather than
accepting a submitted harmonic residual.  Positivity of $A(q)$ makes the
solution unique.  Since $q\not\equiv-q$, the sine and cosine components
have equal norms and their cross terms vanish on the four-point grid;
the real normalization in \eqref{eq:YMA3-real-calibration} therefore
preserves the amplitude norm.
\end{proof}

\subsection{The complete finite loop sum}

For every $k\in(\pi/2)\{0,1,2,3\}^4$ put
$G(k)=[D_R(k)^*D_R(k)/3]^{-1}$.  In one plaquette let
\[
 z_i=\ell_{e_i}e^{iq\cdot d_i},\quad
 f=\sum_i\epsilon_i z_i,\quad
 s_i=\sum_{j>i}\epsilon_j z_j+
             \mathbf1_{\{\epsilon_i=-1\}}\epsilon_i z_i,\quad
 d_{ij}=s_i-s_j-f/2.
\]
Define $\mathcal R_q(k)$ by adding, for every reduced pair $i<j$,
\begin{align}
 [\mathcal R_q(k)]_{e_i e_j}
 &\mathrel{+}=\epsilon_i\epsilon_j d_{ij}
       e^{ik\cdot d_j-i(k+q)\cdot d_i},\nonumber\\
 [\mathcal R_q(k)]_{e_j e_i}
 &\mathrel{+}=-\epsilon_i\epsilon_j d_{ij}
       e^{ik\cdot d_i-i(k+q)\cdot d_j}.
 \label{eq:YMA3-Fourier-R}
\end{align}
Define the zero-external-frequency part $\mathcal Q_0(k)$ by the diagonal
and paired cross additions
\begin{align}
 [\mathcal Q_0(k)]_{e_i e_i}&\mathrel{+}=-\tfrac49|f|^2,\nonumber\\
 [\mathcal Q_0(k)]_{e_i e_j}&\mathrel{+}=
 \epsilon_i\epsilon_j\left(-|d_{ij}|^2-\tfrac7{36}|f|^2\right)
 e^{ik\cdot(d_j-d_i)},
 \label{eq:YMA3-Fourier-Q}
\end{align}
with the Hermitian conjugate in the reverse cross entry.  These are
$45\times45$ matrices, with all ninety-six plaquette types included.

\begin{proposition}[Exactly normalized $256$-momentum formula]
\label{prop:YMA3-finite-loop-sum}
For the unit real calibration field,
\begin{equation}
 \boxed{
 \begin{aligned}
 \mathcal T_4&=\frac1{256}\sum_k\Tr(G(k)\mathcal Q_0(k)),\\
 \mathcal B_4&=\frac1{3\cdot256}\sum_k
 \Tr\bigl(G(k+q)\mathcal R_q(k)G(k)\mathcal R_q(k)^*\bigr),\\
 \gamma_4(\pi/2)&=\mathcal T_4-\mathcal B_4.
 \end{aligned}}
 \label{eq:YMA3-finite-trace-sum}
\end{equation}
Every entry is rational over $\mathbb Q(i)$ before inversion, and the final
answer is a real rational number.  The formula involves no sampling of the
Wilson measure and no continuum loop integral.
\end{proposition}
\begin{proof}
For the unnormalized cosine field, the linear matrix $R$ is the sum of the
$q$ and $-q$ insertion matrices divided by two.  A closed trace of two such
insertions has precisely the two opposite-sign combinations; the two
same-sign combinations have nonzero total external momentum and zero
trace.  Since $R$ is real skew, each surviving trace equals minus the
positive trace in the second line of
\eqref{eq:YMA3-finite-trace-sum}.  Thus
$\Tr(GRGR)=-\frac12\sum_k\Tr(G(k+q)\mathcal R_q(k)G(k)\mathcal R_q(k)^*)$.

The constant part of a product of two real cosine amplitudes is one half
the real part of the corresponding complex amplitude product.  Applying
this to \eqref{eq:YMA3-Q-vertex} gives exactly
\eqref{eq:YMA3-Fourier-Q}; the other external frequencies again have zero
trace.  The unnormalized coarse cosine has squared norm $256/2$.
Dividing \eqref{eq:YMA3-tadpole-bubble} by that norm gives both prefactors
in \eqref{eq:YMA3-finite-trace-sum}.  Gaussian-rational inversion and
finite summation preserve $\mathbb Q(i)$; Hermiticity makes the result
real.  This proof is for $q\not\equiv-q$; a Nyquist mode must use its
single real character rather than these two-character factors.
\end{proof}

\subsection{A rigorous enclosure without trusting floating-point inverses}

The inverse witnesses are Hermitian matrices
$T_k=2^{-40}(T_k^{\rm re}+iT_k^{\rm im})$ with integer numerators.
Write $\mathsf D=2^{40}$, $\mathsf B=1212$,
$\mathsf R=2424$, and $\mathsf Q=52881984=36\mathsf B^2$.
Then $\mathcal R_q(k)=R_k/\mathsf R$ and
$\mathcal Q_0(k)=Q_k/\mathsf Q$ with Gaussian-integer matrices $R_k,Q_k$.
Let $\mathsf A_k=D_R(k)^*D_R(k)$, so $A(k)=\mathsf A_k/3$.
From the exact Gaussian-integer residual
\[
 E_k=3\mathsf D I-\mathsf A_k(T_k^{\rm re}+iT_k^{\rm im})
\]
form $r_k$, the larger of the maximum row sum and the maximum column sum
of $|\Re E_k|+|\Im E_k|$.  The stored witnesses have
\begin{equation}
 \boxed{\max_k r_k=205,\qquad
 \|G(k)-T_k\|_{\op}\le\varepsilon_k:=\frac{20r_k}{\mathsf D}.}
 \label{eq:YMA3-inverse-certification}
\end{equation}

\begin{theorem}[Certified first-axial one-loop interval]
\label{thm:YMA3-certified-value}
Reconstructing the incidence, vertices, and residuals specified above gives
\begin{align}
 31.3522957&<\mathcal T_4<31.3522962,\nonumber\\
 32.4420835&<\mathcal B_4<32.4420965,\nonumber\\
 \boxed{-1.089801<\gamma_4(\pi/2)<-1.089787.}
 &\label{eq:YMA3-certified-components}
\end{align}
In particular the first correction is strictly negative, although both the
classical reference and the limiting compensated covariance are positive.
Its reference-normalized coefficient obeys
\begin{equation}
 \boxed{-2.378461<\frac{\gamma_4(\pi/2)}{s_0(\pi/2)}<-2.378431.}
 \label{eq:YMA3-certified-relative}
\end{equation}
These intervals concern the first correction, not a finite-$\beta$ sign
of the full Hessian.
\end{theorem}
\begin{proof}
The inequality $|x+iy|\le|x|+|y|$ and the row/column Schur estimate give
$\|I-A(k)T_k\|_{\op}\le r_k/(3\mathsf D)$.
Since $\|G(k)\|_{\op}\le60$, multiplication by $G(k)$ proves
\eqref{eq:YMA3-inverse-certification}.
For any Hermitian $Q$,
\[
 |\Tr((G-T)Q)|\le\varepsilon\sum_{i,j}|Q_{ij}|.
\]
For the covariance trace, adding and subtracting the mixed product gives
\[
 \begin{split}
 &\left|\Tr(G_+RGR^*)-\Tr(T_+RTR^*)\right|\\
 &\qquad\le
 \bigl(60(\varepsilon_++\varepsilon)+\varepsilon_+\varepsilon\bigr)
 \|R\|_{\HS}^2.
 \end{split}
\]
Indeed the two error terms use respectively $\|G\|\le60$ and
$\|T_+\|\le60+\varepsilon_+$.  Summing these rational error bounds in
\eqref{eq:YMA3-finite-trace-sum} encloses the two traces and their
\emph{signed} difference.

For reproducibility, the exact center of the resulting signed interval is
\[
 c_*=-\frac{185789244645400601928014154112285}
 {170481055600127216487320543821824},
\]
and its exact radius is
\[
 e_*=
 \frac{8802359571059978815958215}
 {1331883246875993878807191748608}.
\]
They satisfy
$c_*\approx-1.0897940770684784$ and
$e_*<0.000006609$.
The two separate rational trace centers and their error numerators are
retained in the replay certificate.  Integer summation gives
$c_*-e_*>-1089801/10^6$ and
$c_*+e_*<-1089787/10^6$, as well as the displayed component intervals.
Division of the exact endpoints by the positive rational $833/1818$
gives \eqref{eq:YMA3-certified-relative}.

The replay accepts the inverse numerators only as witnesses: it recomputes
every residual with integer arithmetic, verifies the full harmonic table
against the incidence, and recomputes the trace bounds.  Floating-point
linear algebra is used only to propose the witnesses and is absent from
the replay inference.  This is an arithmetic certificate relative to the
specified finite matrices and the proved inverse bound, not a formal
verification of arbitrary software or of the nonlinear Gibbs remainder.
\end{proof}

\section{Uniform locality of the actual Wilson one-loop coefficient}
\label{sec:YMA3-one-loop-locality}

The evaluated number is only one finite card.  The next estimates control
this \emph{same coefficient family} as $M$ grows.  No assertion about a
non-Gaussian covariance is inserted into them.  Gaussian propagators have
long been used in lattice-gauge renormalization; see
\cite{YMA3-Balaban}.  For this pinned fibre the existing uniform gap gives
a direct proof, so no background-field propagator theorem is imported.

For a matrix convolution kernel $F$ on $\Z^4$, write
\begin{equation}
 |F|_\mu=\sum_{x\in\Z^4}\sum_{i,j}
          e^{\mu|x|_1}|F_{ij}(x)|.
 \label{eq:YMA3-entrywise-stock}
\end{equation}
This entrywise norm is submultiplicative under compatible convolution and
is unchanged by taking the adjoint kernel.  The indices in each use below
are the fixed finite edge types, not colours.
Put
\begin{equation}
 \rho=\frac{479}{481},\qquad
 \mu=-\frac14\log\rho>0,\qquad
 \begin{aligned}
 g_\mu&=60\cdot45^2
       \left(\frac{1+\rho^{1/4}}{1-\rho^{1/4}}\right)^4,\\
 h_\mu&=4+128e^{2\mu}g_\mu,\\
 R_0&=2304,\qquad Q_0=\frac{65536}{3}.
 \end{aligned}
 \label{eq:YMA3-fixed-locality-constants}
\end{equation}
All these constants are numerical and independent of the coupling and
volume.  They are deliberately conservative.

\begin{lemma}[The $45$-type inverse and harmonic router have explicit tails]
\label{lem:YMA3-Green-locality}
The infinite-lattice $A$ is invertible, its coarse-cell range is at most
two, and its inverse obeys
\begin{equation}
 \boxed{\|G(x,y)\|_{\op}\le60\rho^{\lceil|x-y|_1/2\rceil},
 \qquad |G|_\mu\le g_\mu.}
 \label{eq:YMA3-Green-tail}
\end{equation}
If $\mathcal H$ denotes the complete map $a\mapsto b[a]$, then
$|\mathcal H|_\mu\le h_\mu$.  On every finite torus,
\begin{equation}
 G_M=\operatorname{Per}_M G,\qquad
 \mathcal H_M=\operatorname{Per}_M\mathcal H,
 \label{eq:YMA3-exact-periodization}
\end{equation}
where periodization sums all lattice translates by $M\Z^4$.
\end{lemma}
\begin{proof}
The V2 spectral window extends to every finitely supported infinite field:
embed its support in a sufficiently large torus before using that finite
window.  It consequently extends to $\ell^2$ by density.  The plaquette
list shows that $D_R$ and $D_C$ use offsets of length at most one, so
$A$ and $B$ have range at most two.  With $c=481/120$, the spectral window
gives $\|I-A/c\|\le479/481=\rho$.  The norm-convergent series
\[
 G=c^{-1}\sum_{n\ge0}(I-A/c)^n
\]
has no block at distance $d$ in the terms $2n<d$.  Summing the remaining
geometric series, with $c^{-1}/(1-\rho)=60$, proves the block bound.
Each of its $45^2$ entries is bounded by that block norm.  Moreover,
$e^{\mu d}\rho^{d/2}=\rho^{d/4}$ and
$\sum_{x\in\Z^4}t^{|x|_1}=((1+t)/(1-t))^4$ for $0<t<1$.
This proves the entrywise bound.

For the mixed matrix $B$, one plaquette with $r$ reduced and $c'$ coarse
constituents contributes at most $rc'/3$ to the sum of absolute entries.
Since $r+c'\le4$, $rc'\le4$.  There are ninety-six plaquette types;
therefore $|B|_\mu\le128e^{2\mu}$.  The direct coarse embedding has four
unit entries, and $L=-GB$, proving the bound for $\mathcal H$.
Finally all sums are absolutely convergent.  Periodizing $AG=I$ gives
$A_M\operatorname{Per}_MG=I_M$.  The finite inverse is unique, proving
the first identity in \eqref{eq:YMA3-exact-periodization}; the second
follows from $L=-GB$ and the local embeddings.
\end{proof}

Define the explicit constants
\begin{align}
 K_\mu&=h_\mu^2\left(30Q_0e^{2\mu}
              +\frac{R_0^2}{3}e^{4\mu}g_\mu^2\right),
 \label{eq:YMA3-K-mu}\\
 E_\mu&=h_\mu^2e^{4\mu}
          \left(\frac{Q_0}{2}g_\mu+\frac{R_0^2}{3}g_\mu^2\right).
 \label{eq:YMA3-E-mu}
\end{align}

\begin{theorem}[Wilson one-loop kernel and exponentially small winding difference]
\label{thm:YMA3-loop-kernel}
The quadratic form \eqref{eq:YMA3-exact-one-loop} on finitely supported
infinite coarse fields defines an absolutely summable real symmetric
convolution kernel $\mathcal J_\infty$ with
\begin{equation}
 \boxed{|\mathcal J_\infty|_\mu\le K_\mu.}
 \label{eq:YMA3-one-loop-stock}
\end{equation}
For the full finite-torus one-loop Hessian $\mathcal J_M=D^2\Gamma_M(0)$,
restricted to one unit colour component,
\begin{equation}
 \boxed{
 \|\mathcal J_M-\operatorname{Per}_M\mathcal J_\infty\|_{\op}
 \le E_\mu e^{-\mu M}.}
 \label{eq:YMA3-winding-bound}
\end{equation}
This is an estimate on the full finite one-loop coefficient; no harmonic
contact is removed in its construction.
\end{theorem}
\begin{proof}
Expand $R(b)$ linearly and $Q(b)$ quadratically in the real fine
coefficients $b_e$.  For a fixed ordered reduced pair $i<j$,
$s_i-s_j$ is a sum of a subset of the four oriented $b$'s.  Thus each
coefficient of $d_{ij}=s_i-s_j-f/2$ has magnitude $1/2$.
There are at most twelve ordered off-diagonal reduced pairs per plaquette
and four source coefficients.  Their total absolute linear coefficient
stock is at most $12\cdot4/2=24$ per plaquette, or $R_0=2304$ per coarse
cell.

In $Q$, every coefficient of an ordered source pair is at most
$2(1/2)^2+7/18=8/9$ in magnitude.  The diagonal formula has the same bound.
There are at most sixteen ordered vertical pairs and sixteen ordered
source pairs.  The resulting per-cell stock is at most
$96\cdot16\cdot16\cdot8/9=65536/3=Q_0$.
All endpoints of a local vertex are within one coarse step of its root.

Before substituting $b=\mathcal H a$, the tadpole contraction has a
finite-range source kernel.  Its two source endpoints are at distance at
most two and its internal Green entry has magnitude at most sixty.
Consequently its weighted entrywise stock is at most
$30Q_0e^{2\mu}$.
For the two-vertex covariance contraction, let $y,y'$ be its two roots.
The distance between its source endpoints is at most $|y-y'|_1+2$,
whereas either Green leg has length at least $|y-y'|_1-2$.
The source weight is therefore bounded by $e^{4\mu}$ times the weight
on that Green leg.  Summing absolute values over all compatible Green
endpoints is bounded by the product of their unrestricted weighted sums,
$g_\mu^2$.  The two local coefficient sums contribute $R_0^2$ and the
prefactor is $1/3$.  This proves the fine-source stock bound
\[
 30Q_0e^{2\mu}+(R_0^2/3)e^{4\mu}g_\mu^2.
\]
Convolution on both sides with $\mathcal H$ gives
\eqref{eq:YMA3-one-loop-stock}.  The same absolute sums justify every
interchange of sums and define the infinite quadratic form.

It remains to compare finite and periodized diagrams.  For a tadpole,
its internal Green displacement has length at most two.  The nonzero
periodization translates therefore have length at least $M-2$; their sum
is bounded by $e^{-\mu(M-2)}g_\mu$.
For a covariance loop, write the two unshifted Green displacements as
$u,v$.  They have $|u+v|_1\le4$.  The finite product contains all translates
$G(u+Mn_1)G(v+Mn_2)$.  The periodization of the infinite loop contains
exactly the terms with $n_1+n_2=0$, after lifting the second vertex to the
covering lattice.  On every other term,
\[
 |u+Mn_1|_1+|v+Mn_2|_1\ge M-4.
\]
For fixed local vertex indices the map from the second root modulo $M$
and the two translation indices to the two lifted Green displacements is
injective: either displacement modulo $M$ fixes that root, and then both
indices are fixed.  Thus the sum of all such omitted products is at most
$e^{-\mu(M-4)}g_\mu^2$, without a volume factor.

The finite fine-source difference is consequently bounded in its unweighted
entrywise kernel norm by
\[
 e^{4\mu}e^{-\mu M}
 \left(\tfrac12 Q_0g_\mu+\tfrac13R_0^2g_\mu^2\right).
\]
By \eqref{eq:YMA3-exact-periodization}, the harmonic routers themselves
periodize exactly, and periodization commutes with their convolution.
Multiplying by their two kernel norms gives
\eqref{eq:YMA3-winding-bound}.  The unweighted entrywise kernel norm
dominates the operator norm by the row/column Schur estimate.
\end{proof}

\section{The Ward-relative one-loop bound is uniform on the physical grids}
\label{sec:YMA3-Ward-loop}

\begin{lemma}[The infinite one-loop contact is fixed by gauge invariance]
\label{lem:YMA3-infinite-Ward}
The kernel $\mathcal J_\infty$ inherits coarse gauge invariance, translations,
reflections, and the physical axial little-group covariance.  Its full
four-direction zero-momentum symbol vanishes:
\begin{equation}
 \boxed{\widehat{\mathcal J}_\infty(0)=0.}
 \label{eq:YMA3-infinite-harmonic-zero}
\end{equation}
\end{lemma}
\begin{proof}
The complete finite partition function is gauge invariant.  The unique
fixed-$M$ asymptotic expansion in V2 implies the same invariance for its
classical and one-loop coefficients separately.  The first derivative of
the latter at the identity vanishes by global conjugation, so its Hessian
annihilates every coarse exact gauge direction.  Thus at every finite
allowed momentum $k$,
$\widehat{\mathcal J}_M(k)d_c(k)=0$, where
$d_c(k)_\mu=e^{ik_\mu}-1$.

The exponential comparison in \eqref{eq:YMA3-winding-bound} and the
continuity supplied by \eqref{eq:YMA3-one-loop-stock} pass this identity
to every real momentum, by approximating it with finite allowed momenta as
$M\to\infty$.  They pass the other physical symmetries in the same way.
Along $k=q\widehat e_\mu$, $q\ne0$, the exact gauge direction is a
nonzero multiple of $e_\mu$.  Hence
$\widehat{\mathcal J}_\infty(q\widehat e_\mu)e_\mu=0$.
Letting $q\to0$ proves that every coordinate vector is in the kernel of
$\widehat{\mathcal J}_\infty(0)$.  This proves the full zero symbol and
not just a fitted scalar row sum.
\end{proof}

\begin{theorem}[Uniform one-loop margin and its finite harmonic correction]
\label{thm:YMA3-uniform-one-loop-Ward}
Let $\gamma_\infty(q)$ be the transverse axial scalar symbol of
$\mathcal J_\infty$ and let $\gamma_M(q)$ be the full finite coefficient.
With
\begin{equation}
 T_\mu=\frac{2K_\mu}{\mu^2},\qquad
 U_\mu=\frac{2K_\mu+E_\mu/4}{\mu^2},
 \label{eq:YMA3-U-mu}
\end{equation}
one has
\begin{align}
 |\gamma_\infty(q)|&\le T_\mu\,\omega(q),
 &&q\in[-\pi,\pi],
 \label{eq:YMA3-infinite-Ward-bound}\\
 |\gamma_M(q)-\gamma_\infty(q)|&\le E_\mu e^{-\mu M},
 &&q\text{ on the finite grid},
 \label{eq:YMA3-loop-grid-difference}\\
 \boxed{|\gamma_M(q)|\le U_\mu\,\omega(q)}
 &\quad &&q\ne0\text{ on that grid},
 \label{eq:YMA3-finite-Ward-bound}\\
 |\gamma_M(0)|&\le E_\mu e^{-\mu M}.
 \label{eq:YMA3-loop-harmonic-bound}
\end{align}
All constants are independent of $M\ge4$.
Consequently the classical-plus-one-loop axial form satisfies
\begin{equation}
 \boxed{
 \frac{\beta s_0(q)+\gamma_M(q)}{\beta s_0(q)}
 \ge1-\frac{28U_\mu}{5\beta}}
 \label{eq:YMA3-truncated-uniform-ratio}
\end{equation}
on every nonzero axial grid mode.  This last assertion concerns the
specified truncation only.
\end{theorem}
\begin{proof}
The real even axial coefficient kernel of $\mathcal J_\infty$ has zero
row sum by \Cref{lem:YMA3-infinite-Ward}.  Its weighted entrywise stock is
at most $K_\mu$.  Therefore
\[
 \frac12\sum_x x_1^2|\mathcal J_{\infty,22}(x)|
 \le\frac{2K_\mu}{\mu^2},
\]
using $r^2e^{-\mu r}\le4/\mu^2$.  Applying the already proved V1
Ward-relative estimate to this \emph{evaluated} kernel gives
\eqref{eq:YMA3-infinite-Ward-bound}.  Fourier diagonalization of the
periodized convolution and \eqref{eq:YMA3-winding-bound} give
\eqref{eq:YMA3-loop-grid-difference} and, at zero,
\eqref{eq:YMA3-loop-harmonic-bound}.

For a nonzero axial grid angle at side $M\ge4$,
\[
 \omega(q)\ge4\sin^2(\pi/M)\ge16/M^2.
\]
Thus the winding error divided by $\omega(q)$ is at most
$(E_\mu/16)M^2e^{-\mu M}\le E_\mu/(4\mu^2)$.
Adding the infinite-kernel bound proves
\eqref{eq:YMA3-finite-Ward-bound}.  Finally V2 gives
$s_0(q)\ge(5/28)\omega(q)$; division yields
\eqref{eq:YMA3-truncated-uniform-ratio}.
\end{proof}

\begin{firewall}[Winding control at one loop does not delete the nonlinear topological packet]
\label{fire:YMA3-winding-scope}
The difference in \eqref{eq:YMA3-winding-bound} is defined by lifting the
explicit one-loop contractions to the infinite covering lattice.  Its
exponential estimate is proved without assuming that it equals an arbitrary
submitted local/topological split of the full effective action.  In
particular \eqref{eq:YMA3-loop-harmonic-bound} rules out an order-one
harmonic survivor \emph{in this coefficient family}; it does not rule out
a nonlinear, generated, or nonperturbative toron survivor.  The same physical
support convention must still be applied to the complete action before the
V1 local contact identity is used.
\end{firewall}

\section{The exact remaining nonlinear row and the V4 attack}
\label{sec:YMA3-frontier}

For the pure Wilson law define, without an asymptotic convention,
\begin{equation}
 \mathcal N_{\beta,M}(q)
 :=h_{\beta,M}(q)-\beta s_0(q)-\gamma_M(q).
 \label{eq:YMA3-exact-nonlinear-residual}
\end{equation}
At each fixed card V2 proves $\mathcal N_{\beta,M}=O_M(\beta^{-1})$.
The new result controls the entire preceding coefficient uniformly:
\begin{equation}
 \boxed{
 \frac{h_{\beta,M}(q)}{\beta s_0(q)}
 =1+\frac{\gamma_M(q)}{\beta s_0(q)}
    +\frac{\mathcal N_{\beta,M}(q)}{\beta s_0(q)},\qquad
 \left|\frac{\gamma_M(q)}{\beta s_0(q)}\right|
 \le\frac{28U_\mu}{5\beta}.}
 \label{eq:YMA3-separated-remainder}
\end{equation}
There is no claim here that the last term is uniform.  In particular, this
identity does not interchange the large-$\beta$ and large-$M$ limits.

For the full response-matched law, the compensation gives the exact
alternative expression for the same V2 local residual:
\begin{equation}
 \boxed{
 \begin{aligned}
 r_{\Phi,M}(q)
 ={}&\beta\mathbb E_\Phi\bigl(\widetilde K_b^V-s_0(q)\bigr)
     +\mathbb E_\Phi\widetilde K_b^\Psi\\
 &-\Var_\Phi\bigl(\beta\widetilde j_b^V+
                         \widetilde j_b^\Psi\bigr)-t_{\Phi,M}(q).
 \end{aligned}}
 \label{eq:YMA3-full-compensated-residual}
\end{equation}
Here $aX$ is a unit real implementation of the already retained physical
axial mode, and $b=b[a]$ is its frozen flat lift.  The entire generated
score and its cross covariance are retained, and the topological term is
exactly the inherited one.  The first term has the actual finite-coupling
bound \eqref{eq:YMA3-deterministic-mean-error}; the second and third have
not been replaced by their pure-Wilson limits.

\begin{assessment}[The quantitative frontier has moved]
\label{ass:YMA3-progress}
The formerly unevaluated pure-Wilson one-loop trace is now an explicit
$45$-type, $96$-plaquette contraction, with a rigorous negative interval on
the first $M=4$ axial card and a cutoff-uniform Ward-relative bound on the
whole nonzero axial family.  The compensated source has no linear Gaussian
part.  What remains is the covariance of its genuinely nonlinear quadratic
and higher fluctuation terms, the selected generated interaction, and their
physical local/topological incidence.  An uncontrolled $O_M(1)$ is no longer
needed for the first correction; an uncontrolled higher remainder remains.
\end{assessment}

\begin{openproblem}[V4: control the post-one-loop same-law remainder]
\label{op:YMA3-V4}
Keep the actual selected $\beta$, $\Psi$, physical mode, and support split.
Use \eqref{eq:YMA3-full-compensated-residual}, not a fresh raw covariance
with its leading Gaussian return still present.  On the pure-Wilson
calibration branch the first unclosed quantity is
$\mathcal N_{\beta,M}(q)$; on the actual branch it is the corresponding
full compensated generated/covariance/topological remainder.

The concrete analytic target is a source-specific connected bound after
subtracting the two explicit contractions
\eqref{eq:YMA3-tadpole-bubble}, with constants uniform in the physical
retained family and in the Ward normalization when $q\to0$.  The full
compact measure must be retained: a small-field expansion without its
large-field error does not answer this question.  A successful bound must
be inserted in the already fixed scalar inequality
$r_{\Phi,M}(q)\ge-(1-\delta)\beta s_0(q)$.  A failure must return the same
mode with its signed generated, nonlinear covariance, or topological
remainder, not the separately large tadpole or covariance limit.
\end{openproblem}

\section{Version V3 proof ledger and exact replay boundary}
\label{sec:YMA3-ledger}

\begingroup\small
\begin{longtable}{>{\raggedright\arraybackslash}p{0.27\textwidth}
                  >{\raggedright\arraybackslash}p{0.20\textwidth}
                  >{\raggedright\arraybackslash}p{0.45\textwidth}}
\caption{Version V3 proof-status ledger.}\label{tab:YMA3-ledger}\\
\toprule
\textbf{Row}&\textbf{Status}&\textbf{Exact scope}\\
\midrule
\endfirsthead
\toprule
\textbf{Row}&\textbf{Status}&\textbf{Exact scope}\\
\midrule
\endhead
Haar harmonic compensation&Proved, exact same law
&No change of $Z$ or of the generated action; linear Gaussian score removed.\\[0.3em]
Compensated Wilson mean&Proved, all volumes
&An explicit $O(1)$ absolute deterministic error under bounded oscillation density.\\[0.3em]
Complete colour vertices&Proved and explicit
&The $-8/9$ diagonal and $-2d^2-7f^2/18$ cross coefficients include every colour.\\[0.3em]
First compensated covariance&Evaluated at fixed-card leading order
&Its limit is the positive trace $\mathcal B_M$; it is not a new source price.\\[0.3em]
First $M=4$ axial correction&Certified strict interval
&$-1.089801<\gamma_4(\pi/2)<-1.089787$, with exact integer replay.\\[0.3em]
Infinite one-loop kernel&Proved, explicit weighted stock
&The literal pinned inverse and local Wilson vertices give $|\mathcal J_\infty|_\mu\le K_\mu$.\\[0.3em]
Finite one-loop winding&Proved, exponential in $M$
&The full coefficient differs from its covering-lattice periodization by at most $E_\mu e^{-\mu M}$.\\[0.3em]
One-loop axial Ward margin&Proved, cutoff independent
&$|\gamma_M(q)|\le U_\mu\omega(q)$ on every nonzero finite axial grid.\\[0.3em]
Complete selected shell margin&Open
&The nonlinear and generated same-law remainder has no proved cofinal relative bound.\\[0.3em]
Higher jets, local-vacuum contraction, mass&Not inferred
&No downstream conclusion is promoted from the coefficient calculation.\\
\bottomrule
\end{longtable}
\endgroup

The replay package contains the complete harmonic table, all $256$ Hermitian
dyadic inverse witnesses, the literal incidence and vertex generators, and
an exact-integer verifier.  The verifier checks harmonic incidence, its
reference metric and norm, every inverse residual, the two trace sums, the
outward error radii, and the strict interval.  A separate direct finite-matrix
determinant difference was used only as a numerical diagnostic of the
orientation and colour conventions; it is not an input to the interval
proof.  No Monte Carlo confidence interval, inferred finite-$\beta$ error,
or claimed formal verification is included.

\begin{assessment}[V3 termination]
The explicit first-correction acquisitions and their stated uniform
one-loop bounds are \pass.  The negative coefficient in
\eqref{eq:YMA3-certified-components} is not a negative full Wilson symbol:
it is an order-one debit against the positive order-$\beta$ reference.
For the complete selected response-matched physical shell, the verdict
remains \stopstatus\ pending the exact post-one-loop remainder in
\Cref{op:YMA3-V4}.  No alternative physical law has been substituted to
manufacture a pass.
\end{assessment}

\begin{thebibliography}{9}
\bibitem{YMA3-Balaban}
T.~Ba{\l}aban,
\emph{Propagators and renormalization transformations for lattice gauge
theories. I}, Communications in Mathematical Physics \textbf{95}
(1984), 17--40. DOI: 10.1007/BF01215753.
Cited for historical context only; the fixed-fibre inverse, Wilson vertices,
and quantitative bounds used in this continuation are proved above.
\end{thebibliography}

\section*{Version V3 append-only continuation record}
\addcontentsline{toc}{section}{Version V3 append-only continuation record}
The complete V2 source is preserved byte for byte up to its sole final
\verb|\end{document}| command.  Its SHA--256 digest is
\begin{center}\scriptsize\ttfamily
3dc6337f56f432f01e569607ef14087a4cbe462516bf111e175e0c444e31fd3b.
\end{center}
For V4, preserve this accumulated source, append new material immediately
before the final document-closing command, and increment the filename to
\texttt{\_V4.tex}.  The submitted physical response, exact block map,
parity-tree convention, source normalization, selected axial mode, and
local/topological packet remain unchanged.

% =============================================================================
% END APPENDED VERSION V3 MATERIAL
% =============================================================================


% =============================================================================
% START APPENDED VERSION V4 MATERIAL
% =============================================================================
\clearpage
\hypersetup{
 pdftitle={Renewal Yang--Mills Reduced-Fibre Wilson Shell Symbol and Local-Vacuum Programme V4},
 pdfsubject={Renormalized beta-flow cumulants, a finite two-loop coefficient card, and an exact generated-action interpolation}
}
\pdfinfo{
 /Title (Renewal Yang--Mills Reduced-Fibre Wilson Shell Symbol and Local-Vacuum Programme V4)
 /Subject (Renormalized beta-flow cumulants, a finite two-loop coefficient card, and an exact generated-action interpolation)
}
\providecommand{\cum}{\operatorname{cum}}
\providecommand{\osc}{\operatorname{osc}}

\section{Version V4: the post-one-loop remainder as a renormalized flow card}
\label{sec:YMA4-direction}

The complete Version~V3 source is preserved above.  Its Haar compensation,
first nonzero Wilson coefficient, certified $M=4$ axial interval, and
cutoff-independent one-loop Ward bound are imported without repetition.
The present continuation attacks only the remainder left in
\eqref{eq:YMA3-exact-nonlinear-residual} and the generated/topological terms
in \eqref{eq:YMA3-full-compensated-residual}.

The first new point is that the pure-Wilson nonlinear remainder is not an
unnamed difference of two large finite-coupling quantities.  It is the
improper integral of one signed \emph{renormalized coupling-flow cumulant}.
For the already compensated source this flow is populated by a single
positive two-replica Wilson cylinder carrying only the unit action, its first
and second source jets, and their common action insertion.  Its four signed
pieces have an exact Gaussian homogeneity cancellation.  Consequently the
flow begins at order $\beta^{-2}$ on every fixed finite card, while the
remainder itself begins at order $\beta^{-1}$.

The second new point is a finite formula for that first post-one-loop
coefficient.  After rescaling the literal reduced-link logarithms, it is one
Gaussian cumulant expression involving only the cubic and quartic fluctuation
potentials and the first three compensated source vertices.  This is the
actual two-loop acquisition card; it is not a new shell representation.

Finally, the submitted generated action is joined to the pure Wilson law by
one positive interpolation.  Differentiation of the same compensated
Hessian gives a second exact two-replica flow.  The complete local symbol is
therefore decomposed as
\begin{equation}
 \boxed{
 \begin{aligned}
 h_{\beta,\Psi,M}^{\rm loc}(q)-\beta s_0(q)
 ={}&\gamma_M(q)+\mathcal N_{\beta,M}(q)\\
 &+\int_0^1\mathfrak G_{\theta,M}(q)\,d\theta
 -t_{\beta,\Psi,M}(q).
 \end{aligned}}
 \label{eq:YMA4-master-decomposition-preview}
\end{equation}
Every term belongs to the same frozen coarse source.  The one-loop term is
already controlled in V3.  This version gives exact finite acquisition and
relative-norm criteria for the other three terms.  It does not assert that
the required uniform connected-flow stock has already been populated on the
selected renormalized trajectory.

\begin{firewall}[A coupling probe is not a retuned physical trajectory]
\label{fire:YMA4-beta-probe}
The auxiliary variables $u\ge\beta$ and $0\le\theta\le1$ below parametrize
positive unnormalized cylinders obtained by changing, respectively, the
coefficient of the already occurring Wilson action and the strength of the
already submitted generated action.  They are response probes of the frozen
finite action grammar.  They do not assert that the physical renormalization
trajectory passes through those auxiliary laws, and no beta function is
inferred from them.  Their role is to evaluate an exact identity for the
single selected finite-coupling symbol.
\end{firewall}

% =============================================================================
\section{The exact pure-Wilson beta flow}
\label{sec:YMA4-beta-flow}
% =============================================================================

Fix a finite side $M$, one real coarse tangent $aX$, and its flat harmonic
lift $b=b[a]$ from \eqref{eq:YMA3-fine-lift}.  In this section the law is the
pure Wilson reduced-fibre law at $W=I$,
\begin{equation}
 d\nu_{\beta,M}(z)=Z_{\beta,0,M}(I)^{-1}
 e^{-\beta V_M(z,I)}\,dm_{\rm red}(z).
 \label{eq:YMA4-pure-law}
\end{equation}
Suppress $M$ and write
\begin{equation}
 V=V_M(z,I),\qquad
 J_a=\widetilde j^V_{b[a]}(z),\qquad
 K_a=\widetilde K^V_{b[a]}(z),\qquad
 s_a=\langle a,\mathsf S_Ma\rangle.
 \label{eq:YMA4-VJKs}
\end{equation}
All three observables are derivatives of the same compensated unit Wilson
cylinder.  In particular, $J_a$ is not the uncompensated score used before
V3.

For bounded random variables under a law $\nu$, use the joint-cumulant
notation
\begin{equation}
 \langle A_1;\ldots;A_n\rangle_{\nu,c}
 =\left.\partial_{t_1}\cdots\partial_{t_n}
 \log\mathbb E_\nu
 e^{t_1A_1+\cdots+t_nA_n}\right|_{t=0}.
 \label{eq:YMA4-cumulant-definition}
\end{equation}
Thus the one- and two-point cumulants are expectation and covariance.  The
compactness of the finite fibre makes every cumulant below finite and makes
the dependence on $\beta>0$ real analytic.

\begin{theorem}[Exact renormalized beta-flow identity]
\label{thm:YMA4-beta-flow}
Let
\begin{equation}
 F_{\beta,M}[a]
 :=H_{\beta,0,M}[aX,aX]-\beta s_a,
 \label{eq:YMA4-F-beta}
\end{equation}
where $H_{\beta,0,M}$ is the full pure-Wilson physical Hessian.  Then
\begin{equation}
 \boxed{\frac{d}{d\beta}F_{\beta,M}[a]
 =\mathfrak D_{\beta,M}[a],}
 \label{eq:YMA4-flow-derivative}
\end{equation}
with
\begin{equation}
 \boxed{
 \begin{aligned}
 \mathfrak D_{\beta,M}[a]
 ={}&\langle K_a-s_a\rangle_{\beta}
 -\beta\langle K_a-s_a;V\rangle_{\beta,c}\\
 &-2\beta\langle J_a;J_a\rangle_{\beta,c}
 +\beta^2\langle J_a;J_a;V\rangle_{\beta,c}.
 \end{aligned}}
 \label{eq:YMA4-beta-flow-card}
\end{equation}
The first line is the scale-homogeneity defect of the compensated
second jet.  The second line is the scale-homogeneity defect of its complete
connected score variance.  Neither line is assigned a separate physical
price.
\end{theorem}

\begin{proof}
The exact compensated Hessian identity gives
\begin{equation}
 H_{\beta,0,M}[aX,aX]
 =\beta\langle K_a\rangle_\beta
  -\beta^2\langle J_a;J_a\rangle_{\beta,c}.
 \label{eq:YMA4-H-beta-cumulant}
\end{equation}
For every $\beta$-independent family of observables,
\begin{equation}
 \frac{d}{d\beta}
 \langle A_1;\ldots;A_n\rangle_{\beta,c}
 =-\langle A_1;\ldots;A_n;V\rangle_{\beta,c}.
 \label{eq:YMA4-cumulant-beta-derivative}
\end{equation}
Indeed both sides are the corresponding mixed derivatives of
$\log\int e^{-\beta V+\sum_jt_jA_j}\,dm_{\rm red}$, and mixed derivatives
commute on the finite compact fibre.  Differentiate
\eqref{eq:YMA4-H-beta-cumulant}, subtract $s_a$, and use that a constant has
zero covariance.  This gives exactly \eqref{eq:YMA4-beta-flow-card}.
\end{proof}

\begin{corollary}[One two-replica acquisition of the complete beta flow]
\label{cor:YMA4-beta-two-replica}
Let $z_1,z_2$ be independent samples from \eqref{eq:YMA4-pure-law}, and
write $\Delta A=A(z_1)-A(z_2)$.  Global colour symmetry gives
$\langle J_a\rangle_\beta=0$.  Consequently
\begin{equation}
 \boxed{\mathfrak D_{\beta,M}[a]
 =\mathbb E_{\beta}^{\otimes2}\mathscr D_{\beta,a}(z_1,z_2),}
 \label{eq:YMA4-beta-replica-expectation}
\end{equation}
where the single signed replica writer is
\begin{equation}
 \boxed{
 \begin{aligned}
 \mathscr D_{\beta,a}
 ={}&\frac12\bigl(K_a(z_1)+K_a(z_2)\bigr)-s_a\\
 &-\frac\beta2\,\Delta K_a\,\Delta V
 -\beta(\Delta J_a)^2
 +\frac{\beta^2}{2}\,\Delta(J_a^2)\,\Delta V.
 \end{aligned}}
 \label{eq:YMA4-beta-replica-writer}
\end{equation}
Thus the four cancellations in \eqref{eq:YMA4-beta-flow-card} can be
acquired on one common two-copy cylinder.  Separate confidence intervals
for the four large terms are not required and need not be summed.
\end{corollary}

\begin{proof}
For independent identical copies,
\[
 \frac12\mathbb E(\Delta A\,\Delta B)=\Cov(A,B),
 \qquad
 \frac12\mathbb E(\Delta A)^2=\Var(A).
\]
Since $\mathbb E J_a=0$, the third cumulant in
\eqref{eq:YMA4-beta-flow-card} is $\Cov(J_a^2,V)$, and
\[
 \frac12\mathbb E\bigl(\Delta(J_a^2)\,\Delta V\bigr)
 =\Cov(J_a^2,V).
\]
Substitution proves the formula.  The mean-score identity follows either
from constant global conjugation, as in V3, or from the nontrivial
translation character on the frozen axial mode.
\end{proof}

\begin{theorem}[Exact tail integral and dyadic telescope for the nonlinear remainder]
\label{thm:YMA4-beta-tail}
Let $\gamma_M[a]$ be the one-loop coefficient of V3 and put
\begin{equation}
 \mathcal N_{\beta,M}[a]
 =F_{\beta,M}[a]-\gamma_M[a].
 \label{eq:YMA4-N-functional}
\end{equation}
Then, as an improper finite-card integral,
\begin{equation}
 \boxed{
 \mathcal N_{\beta,M}[a]
 =-\int_\beta^\infty\mathfrak D_{u,M}[a] \,du.}
 \label{eq:YMA4-N-beta-integral}
\end{equation}
Equivalently, with $\beta_j=2^j\beta$,
\begin{equation}
 \boxed{
 \mathcal N_{\beta,M}[a]
 =\sum_{j=0}^{\infty}
 \bigl(F_{\beta_j,M}[a]-F_{\beta_{j+1},M}[a]\bigr).}
 \label{eq:YMA4-dyadic-telescope}
\end{equation}
Both identities use the same pure Wilson source at different auxiliary
couplings.  They introduce no source decomposition and no attribution
between the two terms of the Hessian.
\end{theorem}

\begin{proof}
V2 and V3 prove, at every fixed finite card,
$F_{u,M}[a]\to\gamma_M[a]$ as $u\to\infty$.  The fundamental theorem of
calculus on $[\beta,B]$ gives
$F_{B,M}-F_{\beta,M}=\int_\beta^B\mathfrak D_{u,M}\,du$.
Letting $B\to\infty$ proves \eqref{eq:YMA4-N-beta-integral}; no prior
absolute-integrability claim is needed.  Telescoping the finite sum from
$j=0$ to $J-1$ gives $F_{\beta,M}-F_{\beta_J,M}$, which converges to
$F_{\beta,M}-\gamma_M$.
\end{proof}

% =============================================================================
\section{Gaussian homogeneity cancellation and the finite two-loop card}
\label{sec:YMA4-two-loop}
% =============================================================================

The beta-flow formula is exact at finite coupling.  We now identify why it
is better conditioned than the original raw Hessian.  This section is
finite-card only; no constant is declared uniform in $M$.

Use product exponential coordinates $z_e=e^{x_e}$ near the unique flat
minimum and put $x=\varepsilon\xi$, $\varepsilon=\beta^{-1/2}$.  Let
$J_{\rm Haar}(x)$ be the product Haar density in these coordinates.  On a
fixed symmetric chart define the complete rescaled fluctuation potential
\begin{equation}
 \mathscr U_\varepsilon(\xi)
 =\varepsilon^{-2}V_M(e^{\varepsilon\xi},I)
  -\log\frac{J_{\rm Haar}(\varepsilon\xi)}{J_{\rm Haar}(0)}.
 \label{eq:YMA4-rescaled-potential}
\end{equation}
The factor $\varepsilon^{d_M}$ from $dx$ is independent of the coarse
source and cancels from normalized expectations.  Define also
\begin{align}
 \mathscr X_\varepsilon(\xi)
 &=\varepsilon^{-2}
   \bigl(\widetilde K_b^V(e^{\varepsilon\xi})-s_a\bigr),
 \label{eq:YMA4-rescaled-X}\\
 \mathscr Y_\varepsilon(\xi)
 &=\varepsilon^{-2}\widetilde j_b^V(e^{\varepsilon\xi}).
 \label{eq:YMA4-rescaled-Y}
\end{align}
Their local Taylor coefficients are fixed by
\begin{align}
 \mathscr U_\varepsilon
 &=E_0+\varepsilon U_1+\varepsilon^2U_2+O(\varepsilon^3),
 \label{eq:YMA4-U-expansion}\\
 \mathscr X_\varepsilon
 &=X_0+\varepsilon X_1+\varepsilon^2X_2+O(\varepsilon^3),
 \label{eq:YMA4-X-expansion}\\
 \mathscr Y_\varepsilon
 &=Y_0+\varepsilon Y_1+\varepsilon^2Y_2+O(\varepsilon^3).
 \label{eq:YMA4-Y-expansion}
\end{align}
Here
\begin{equation}
 E_0=\frac12\langle\xi,(A\otimes I_8)\xi\rangle,
 \quad
 X_0=\frac12\langle\xi,M_b''(0)\xi\rangle,
 \quad
 Y_0=\frac12\langle\xi,M_b'(0)\xi\rangle,
 \label{eq:YMA4-leading-polynomials}
\end{equation}
with the literal V3 matrices.  The coefficient $U_2$ includes the quadratic
Haar-density term; it is not only the quartic Wilson vertex.
Let $\langle\cdot\rangle_0$ and $\langle\cdot;\ldots;\cdot\rangle_{0,c}$
denote moments and cumulants of the centered Gaussian with energy $E_0$.

\begin{lemma}[Parity and integer weak-coupling powers on each fixed card]
\label{lem:YMA4-parity}
The coefficients in \eqref{eq:YMA4-U-expansion}--
\eqref{eq:YMA4-Y-expansion} obey
\begin{equation}
 U_j(-\xi)=(-1)^jU_j(\xi),\qquad
 X_j(-\xi)=(-1)^jX_j(\xi),\qquad
 Y_j(-\xi)=(-1)^jY_j(\xi).
 \label{eq:YMA4-coefficient-parity}
\end{equation}
For every fixed $M$ and fixed coarse source, the complete finite-card
correction has an expansion
\begin{equation}
 \boxed{
 F_{\beta,M}[a]
 =\gamma_M[a]+\frac{\eta_M[a]}{\beta}+O_M(\beta^{-2}).}
 \label{eq:YMA4-fixed-two-loop-expansion}
\end{equation}
In particular
\begin{equation}
 \boxed{
 \mathfrak D_{\beta,M}[a]
 =-\frac{\eta_M[a]}{\beta^2}+O_M(\beta^{-3}).}
 \label{eq:YMA4-fixed-flow-asymptotic}
\end{equation}
\end{lemma}

\begin{proof}
Replacing $\varepsilon$ by $-\varepsilon$ in each defining expression is
the same as replacing $\xi$ by $-\xi$.  Comparison of Taylor coefficients
gives \eqref{eq:YMA4-coefficient-parity}.  The Gaussian base law is even,
so every odd total coefficient has zero integral.

For completeness, the fixed-card Laplace argument of V2 may be continued
two further orders.  Choose the same symmetric logarithm neighborhood of
the unique minimum and a symmetric cutoff equal to one on a smaller
neighborhood.  Analytic Taylor expansion of the action, source jets, and
Haar density through the required order has a remainder dominated by a
fixed Gaussian polynomial after the rescaling $x=\varepsilon\xi$.  The
complement of the chart has a positive action gap and is
$O_M(e^{-c_M/\varepsilon^2})$ together with all source derivatives used
here.  Termwise Gaussian integration is therefore valid through order
$\varepsilon^4$.  Odd coefficients vanish by
\eqref{eq:YMA4-coefficient-parity}, giving an expansion in
$\varepsilon^2=\beta^{-1}$.  Its constant coefficient is the one-loop
quantity already identified in V3.  Denote the next coefficient by
$\eta_M[a]$; this proves \eqref{eq:YMA4-fixed-two-loop-expansion}.
Differentiation of the same finite asymptotic expansion, or direct
termwise differentiation of its Gaussian integrals, gives
\eqref{eq:YMA4-fixed-flow-asymptotic}.
\end{proof}

\begin{lemma}[Exact Gaussian radial identities]
\label{lem:YMA4-Gaussian-homogeneity}
For the quadratic variables in \eqref{eq:YMA4-leading-polynomials},
\begin{equation}
 \boxed{
 \langle X_0;E_0\rangle_{0,c}=\langle X_0\rangle_0,
 \qquad
 \langle Y_0;Y_0;E_0\rangle_{0,c}
 =2\langle Y_0;Y_0\rangle_{0,c}.}
 \label{eq:YMA4-Gaussian-radial-identities}
\end{equation}
Consequently the nominal order-$\beta^{-1}$ part of the exact beta flow
\eqref{eq:YMA4-beta-flow-card} vanishes term for term after the two natural
pairs are assembled.
\end{lemma}

\begin{proof}
For $\alpha>0$, let the Gaussian energy be $\alpha E_0$.  A homogeneous
quadratic variable has expectation proportional to $\alpha^{-1}$ and the
covariance of two homogeneous quadratic variables is proportional to
$\alpha^{-2}$.  Differentiating these two scaling identities at
$\alpha=1$ and using
\[
 \frac d{d\alpha}\langle Z\rangle_\alpha
 =-\langle Z;E_0\rangle_{\alpha,c},
 \qquad
 \frac d{d\alpha}\langle Z;W\rangle_{\alpha,c}
 =-\langle Z;W;E_0\rangle_{\alpha,c}
\]
gives \eqref{eq:YMA4-Gaussian-radial-identities}.  Under the rescaling,
$K_a-s_a=\beta^{-1}X_0+o(\beta^{-1})$,
$J_a=\beta^{-1}Y_0+o(\beta^{-1})$, and
$V=\beta^{-1}E_0+o(\beta^{-1})$.  Substitution into
\eqref{eq:YMA4-beta-flow-card} leaves, at order $\beta^{-1}$,
\[
 \langle X_0\rangle_0-\langle X_0;E_0\rangle_{0,c}
 -2\langle Y_0;Y_0\rangle_{0,c}
 +\langle Y_0;Y_0;E_0\rangle_{0,c}=0.
\]
\end{proof}

\begin{theorem}[Finite exact formula for the first post-one-loop coefficient]
\label{thm:YMA4-two-loop-formula}
With the Taylor convention in
\eqref{eq:YMA4-U-expansion}--\eqref{eq:YMA4-Y-expansion}, the coefficient
$\eta_M[a]$ in \eqref{eq:YMA4-fixed-two-loop-expansion} is
\begin{equation}
 \boxed{
 \begin{aligned}
 \eta_M[a]={}&
 \langle X_2\rangle_0
 -\langle X_1;U_1\rangle_{0,c}
 -\langle X_0;U_2\rangle_{0,c}
 +\frac12\langle X_0;U_1;U_1\rangle_{0,c}\\
 &-2\langle Y_0;Y_2\rangle_{0,c}
 -\langle Y_1;Y_1\rangle_{0,c}
 +2\langle Y_0;Y_1;U_1\rangle_{0,c}\\
 &+\langle Y_0;Y_0;U_2\rangle_{0,c}
 -\frac12\langle Y_0;Y_0;U_1;U_1\rangle_{0,c}.
 \end{aligned}}
 \label{eq:YMA4-two-loop-formula}
\end{equation}
Every object in this formula is a finite Gaussian contraction of the
literal parity-tree Wilson incidence, the ordered plaquette words, and the
product-Haar Jacobian.  Only cubic and quartic fluctuation vertices and
source vertices through the displayed orders are needed.  No finite-coupling
Gibbs sampling and no new coarse carrier enter this coefficient.
\end{theorem}

\begin{proof}
For any observable expansion
$Z_\varepsilon=Z_0+\varepsilon Z_1+\varepsilon^2Z_2+O(\varepsilon^3)$,
normalized perturbation of the Gaussian law by
$\exp(-\varepsilon U_1-\varepsilon^2U_2+O(\varepsilon^3))$ gives
\begin{equation}
 \begin{aligned}
 \langle Z_\varepsilon\rangle_\varepsilon
 ={}&\langle Z_0\rangle_0\\
 &+\varepsilon\bigl(
   \langle Z_1\rangle_0-\langle Z_0;U_1\rangle_{0,c}\bigr)\\
 &+\varepsilon^2\bigl(
   \langle Z_2\rangle_0-\langle Z_1;U_1\rangle_{0,c}
   -\langle Z_0;U_2\rangle_{0,c}
   +\tfrac12\langle Z_0;U_1;U_1\rangle_{0,c}\bigr)
 +O(\varepsilon^3).
 \end{aligned}
 \label{eq:YMA4-expectation-perturbation}
\end{equation}
This follows by expanding the numerator and denominator and collecting
connected terms, or by differentiating a joint cumulant-generating
function.

The same cumulant differentiation gives
\begin{equation}
 \begin{aligned}
 \langle\mathscr Y_\varepsilon;
             \mathscr Y_\varepsilon\rangle_{\varepsilon,c}
 ={}&\langle Y_0;Y_0\rangle_{0,c}\\
 &+\varepsilon\bigl(
  2\langle Y_0;Y_1\rangle_{0,c}
  -\langle Y_0;Y_0;U_1\rangle_{0,c}\bigr)\\
 &+\varepsilon^2\bigl(
  2\langle Y_0;Y_2\rangle_{0,c}
  +\langle Y_1;Y_1\rangle_{0,c}
  -2\langle Y_0;Y_1;U_1\rangle_{0,c}\\
 &\hspace{6.2em}
  -\langle Y_0;Y_0;U_2\rangle_{0,c}
  +\tfrac12\langle Y_0;Y_0;U_1;U_1\rangle_{0,c}\bigr)
 +O(\varepsilon^3).
 \end{aligned}
 \label{eq:YMA4-covariance-perturbation}
\end{equation}
By the exact compensated Hessian identity,
$F_{\beta,M}=\langle\mathscr X_\varepsilon\rangle_\varepsilon-
\langle\mathscr Y_\varepsilon;\mathscr Y_\varepsilon\rangle_{
\varepsilon,c}$.  The constant term is
$\gamma_M=\langle X_0\rangle_0-
\langle Y_0;Y_0\rangle_{0,c}$, exactly the V3 tadpole minus bubble.
The coefficient of $\varepsilon$ vanishes by parity.  Subtracting the two
$\varepsilon^2$ coefficients gives \eqref{eq:YMA4-two-loop-formula}.
\end{proof}

\begin{assessment}[What the two-loop formula accomplishes]
\label{ass:YMA4-two-loop-scope}
The first unknown coefficient after V3 is now a finite, explicit contraction
rather than an unspecified $O_M(\beta^{-1})$.  On the $M=4$, $q=\pi/2$
card, the same $256$ momentum fibres and exact inverse witnesses already
used in V3 suffice in principle; the new data are the ordered cubic and
quartic Wilson/Haar vertices and the source vertices $X_1,X_2,Y_1,Y_2$.
The value and sign of \eqref{eq:YMA4-two-loop-formula} have not been
invented in this version.  More importantly, fixed-card finiteness of this
contraction does not give a cutoff-uniform bound for the complete beta-flow
tail.
\end{assessment}

% =============================================================================
\section{The generated action as one exact positive interpolation}
\label{sec:YMA4-generated-flow}
% =============================================================================

Return to the submitted finite generated action $\Psi$ and retain the same
compensated coarse path.  Assume the inherited coarse-translation,
reflection, gauge, and global-colour symmetries whenever a scalar momentum
symbol is invoked.  The differential identities themselves remain valid as
quadratic-form identities without that scalarization.  For
$0\le\theta\le1$ put
\begin{equation}
 \Phi_\theta=\beta V_M+\theta\Psi,
 \qquad
 d\nu_\theta=Z_\theta^{-1}e^{-\Phi_\theta}\,dm_{\rm red}.
 \label{eq:YMA4-theta-law}
\end{equation}
Let
\begin{equation}
 J_\theta=\beta J_a+\theta J_a^\Psi,
 \qquad
 K_\theta=\beta K_a+\theta K_a^\Psi,
 \label{eq:YMA4-theta-jets}
\end{equation}
where $J_a^\Psi=\widetilde j_b^\Psi$ and
$K_a^\Psi=\widetilde K_b^\Psi$.  These are the complete first and second
source derivatives of one action $\Phi_\theta$; no pure-law covariance is
inserted into a generated-law identity.

\begin{theorem}[Exact generated-action flow]
\label{thm:YMA4-generated-flow}
Let
\begin{equation}
 H_{\theta,M}[a]
 =\langle K_\theta\rangle_\theta
  -\langle J_\theta;J_\theta\rangle_{\theta,c}.
 \label{eq:YMA4-H-theta}
\end{equation}
Then
\begin{equation}
 \boxed{\frac d{d\theta}H_{\theta,M}[a]
 =\mathfrak G_{\theta,M}[a],}
 \label{eq:YMA4-theta-derivative}
\end{equation}
where
\begin{equation}
 \boxed{
 \begin{aligned}
 \mathfrak G_{\theta,M}[a]
 ={}&\langle K_a^\Psi\rangle_\theta
 -\langle K_\theta;\Psi\rangle_{\theta,c}\\
 &-2\langle J_\theta;J_a^\Psi\rangle_{\theta,c}
 +\langle J_\theta;J_\theta;\Psi\rangle_{\theta,c}.
 \end{aligned}}
 \label{eq:YMA4-generated-flow-card}
\end{equation}
Consequently the exact full-Hessian change caused by the submitted
generated action is
\begin{equation}
 \boxed{
 H_{\beta,\Psi,M}[a]-H_{\beta,0,M}[a]
 =\int_0^1\mathfrak G_{\theta,M}[a] \,d\theta.}
 \label{eq:YMA4-generated-integral}
\end{equation}
This integral includes the Wilson--generated cross covariance and the
change of the Wilson expectation under the generated law.
\end{theorem}

\begin{proof}
For $\theta$-independent observables,
$d\langle A_1;\ldots;A_n\rangle_{\theta,c}/d\theta
=-\langle A_1;\ldots;A_n;\Psi\rangle_{\theta,c}$.
The explicit derivatives are $dK_\theta/d\theta=K_a^\Psi$ and
$dJ_\theta/d\theta=J_a^\Psi$.  Apply the product rule to
\eqref{eq:YMA4-H-theta}.  This gives
\eqref{eq:YMA4-generated-flow-card}.  Integration from zero to one proves
\eqref{eq:YMA4-generated-integral}.
\end{proof}

\begin{corollary}[Two-replica generated-flow writer]
\label{cor:YMA4-generated-two-replica}
On the frozen nonzero axial colour mode, translation/global-colour symmetry
gives $\langle J_\theta\rangle_\theta=0$.  For two independent
$\nu_\theta$ replicas,
\begin{equation}
 \boxed{\mathfrak G_{\theta,M}[a]
 =\mathbb E_\theta^{\otimes2}\mathscr G_{\theta,a}^\Psi,}
 \label{eq:YMA4-generated-replica-expectation}
\end{equation}
with
\begin{equation}
 \boxed{
 \begin{aligned}
 \mathscr G_{\theta,a}^\Psi
 ={}&\frac12\bigl(K_a^\Psi(z_1)+K_a^\Psi(z_2)\bigr)
 -\frac12\Delta K_\theta\,\Delta\Psi\\
 &-\Delta J_\theta\,\Delta J_a^\Psi
 +\frac12\Delta(J_\theta^2)\,\Delta\Psi.
 \end{aligned}}
 \label{eq:YMA4-generated-replica-writer}
\end{equation}
Thus one common two-replica panel acquires the generated second jet, the
law-change covariance, every Wilson--generated score cross term, and the
third connected action insertion.
\end{corollary}

\begin{proof}
Use the same two-copy covariance identities as in
\Cref{cor:YMA4-beta-two-replica}.  On the zero-mean branch,
$\langle J_\theta;J_\theta;\Psi\rangle_c=
\Cov(J_\theta^2,\Psi)$, which is the final replica difference in
\eqref{eq:YMA4-generated-replica-writer}.
\end{proof}

\begin{corollary}[Exact assembled residual after pure and generated flows]
\label{cor:YMA4-master-decomposition}
Let $t_{\beta,\Psi,M}(q)$ be the inherited topological Hessian contribution
on the same axial polarization, and let $\mathfrak G_{\theta,M}(q)$ be the
symbol of the polarized form in
\eqref{eq:YMA4-generated-flow-card}.  Then for every nonzero allowed axial
angle,
\begin{equation}
 \boxed{
 \begin{aligned}
 h_{\beta,\Psi,M}^{\rm loc}(q)-\beta s_0(q)
 ={}&\gamma_M(q)
 -\int_\beta^\infty\mathfrak D_{u,M}(q)\,du\\
 &+\int_0^1\mathfrak G_{\theta,M}(q)\,d\theta
 -t_{\beta,\Psi,M}(q).
 \end{aligned}}
 \label{eq:YMA4-master-decomposition}
\end{equation}
This is \eqref{eq:YMA4-master-decomposition-preview} with
$\mathcal N_{\beta,M}$ evaluated by its exact beta-flow integral.
\end{corollary}

\begin{proof}
The pure full Hessian is
$\beta s_0+\gamma_M+\mathcal N_{\beta,M}$ by definition.  Apply
\Cref{thm:YMA4-beta-tail,thm:YMA4-generated-flow} and then subtract the
final local/topological contribution on the same carrier.  No intermediate
term is separately shorted or repriced.
\end{proof}

\begin{firewall}[Interpolation does not erase support incidence]
\label{fire:YMA4-generated-support}
The integral in \eqref{eq:YMA4-generated-integral} compares full finite
Hessians.  The local/topological support split is applied to the complete
final form, as in \eqref{eq:YMA4-master-decomposition}.  A generated
interpolation whose endpoint is local does not prove that each intermediate
integrand is separately local, and a vanishing generated scalar does not
remove an independently occurring topological or harmonic source.
\end{firewall}

% =============================================================================
\section{A harmonic-inclusive Ward router for the flow kernels}
\label{sec:YMA4-Ward-flow}
% =============================================================================

The beta and generated flow cards are exact scalars at each momentum.  For
a cofinal statement, one must distinguish a genuine zero-mode contribution
from the curvature-relative displacement stock.  The following elementary
router extends the V1 Ward estimate to a kernel whose row sum has not yet
been proved to vanish.

For a periodic displacement $x$, use the shortest representative of $x_1$
in the interval $[-M/2,M/2]$.  Let
$\omega(q)=2(1-\cos q)$ as before.

\begin{lemma}[Harmonic plus Ward-repaired periodic kernel]
\label{lem:YMA4-harmonic-Ward-router}
Let $r_M(x)$ be a real reflection-even scalar convolution kernel on the
coarse $M$-torus, with symbol $\widehat r_M(q)$ on the axial grid.  Put
\begin{equation}
 z_M=\widehat r_M(0)=\sum_xr_M(x),
 \qquad
 r_M^\circ(x)=r_M(x)-z_M\mathbf1_{\{x=0\}},
 \label{eq:YMA4-repaired-kernel}
\end{equation}
and
\begin{equation}
 \mathsf W_{2,M}(r_M^\circ)
 =\frac12\sum_xx_1^2|r_M^\circ(x)|.
 \label{eq:YMA4-periodic-W2}
\end{equation}
Then $\sum_xr_M^\circ(x)=0$ and, for every nonzero axial grid angle,
\begin{equation}
 \boxed{
 \frac{|\widehat r_M(q)|}{\omega(q)}
 \le\mathsf W_{2,M}(r_M^\circ)+\frac{M^2}{16}|z_M|.}
 \label{eq:YMA4-harmonic-Ward-bound}
\end{equation}
The first term is the curvature-relative displacement stock.  The second
is the exact finite-grid price of a harmonic or topological scalar.
\end{lemma}

\begin{proof}
The repaired kernel has zero row sum and
$\widehat r_M(q)=z_M+\widehat r_M^\circ(q)$.  Reflection evenness and the
integer sine inequality used in V1 give
$|\widehat r_M^\circ(q)|\le
\omega(q)\mathsf W_{2,M}(r_M^\circ)$.
For a nonzero $M$-grid angle,
$\omega(q)\ge4\sin^2(\pi/M)\ge16/M^2$.  Divide the preceding decomposition
by $\omega(q)$.
\end{proof}

For a predeclared retained nonzero axial set $\mathcal Q_M$, define the
exact finite-card relative stocks
\begin{align}
 d_{u,M}&=\max_{q\in\mathcal Q_M}
   \frac{|\mathfrak D_{u,M}(q)|}{\omega(q)},
 \label{eq:YMA4-d-stock}\\
 g_{\theta,M}&=\max_{q\in\mathcal Q_M}
   \frac{|\mathfrak G_{\theta,M}(q)|}{\omega(q)},
 \label{eq:YMA4-g-stock}\\
 \tau_{\beta,\Psi,M}^{\rm top}&=\max_{q\in\mathcal Q_M}
   \frac{|t_{\beta,\Psi,M}(q)|}{\omega(q)}.
 \label{eq:YMA4-top-stock}
\end{align}
They are values of the actual selected cards, not universal norms imposed on
an unrelated carrier.  Put
\begin{equation}
 A_{\beta,M}=\int_\beta^\infty d_{u,M}\,du,
 \qquad
 B_{\beta,\Psi,M}=\int_0^1g_{\theta,M}\,d\theta,
 \qquad
 T_{\beta,\Psi,M}=\tau_{\beta,\Psi,M}^{\rm top}.
 \label{eq:YMA4-ABT-stocks}
\end{equation}
At fixed $M$, the first integral is finite by
\eqref{eq:YMA4-fixed-flow-asymptotic}.  Its uniformity in $M$ is a separate
statement.

\begin{proposition}[Displacement and harmonic sufficient stocks]
\label{prop:YMA4-kernel-sufficient}
Suppose the beta-flow and generated-flow forms have real even periodic
kernels $\dot f_{u,M}$ and $g_{\theta,M}^{\rm ker}$.  Let their zero symbols
be $z_{u,M}^{\beta}$ and $z_{\theta,M}^{\Psi}$ and repair their contacts as
in \Cref{lem:YMA4-harmonic-Ward-router}.  Then
\begin{align}
 d_{u,M}
 &\le \mathsf W_{2,M}(\dot f_{u,M}^{\circ})
      +\frac{M^2}{16}|z_{u,M}^{\beta}|,
 \label{eq:YMA4-d-kernel-bound}\\
 g_{\theta,M}
 &\le \mathsf W_{2,M}((g_{\theta,M}^{\rm ker})^{\circ})
      +\frac{M^2}{16}|z_{\theta,M}^{\Psi}|.
 \label{eq:YMA4-g-kernel-bound}
\end{align}
On an independently verified local-curvature branch the corresponding zero
symbol vanishes and only the weighted displacement moment remains.  If that
moment is bounded while the $M^2z$ term persists, the failure is harmonic or
topological rather than an ordinary local-curvature mode.
\end{proposition}

\begin{proof}
Apply \Cref{lem:YMA4-harmonic-Ward-router} to each parameter value and then
take the maximum over the retained grid.
\end{proof}

\begin{theorem}[Exact sufficient margin from the three remaining physical stocks]
\label{thm:YMA4-margin}
Let $0\le\delta<1$.  On every retained nonzero axial mode,
\begin{equation}
 \boxed{
 \begin{aligned}
 h_{\beta,\Psi,M}^{\rm loc}(q)-\delta\beta s_0(q)
 \ge{}&\Bigl[
 (1-\delta)\frac{5\beta}{28}
 -U_\mu-A_{\beta,M}\\
 &\hspace{5em}-B_{\beta,\Psi,M}
 -T_{\beta,\Psi,M}\Bigr]\omega(q).
 \end{aligned}}
 \label{eq:YMA4-margin-lower}
\end{equation}
Consequently
\begin{equation}
 \boxed{
 U_\mu+A_{\beta,M}+B_{\beta,\Psi,M}+T_{\beta,\Psi,M}
 <(1-\delta)\frac{5\beta}{28}}
 \label{eq:YMA4-margin-pass-condition}
\end{equation}
is a rigorous \pass\ certificate for the selected local axial bank.

For a cofinal shell family $(\beta_r,M_r,\Psi_r)$, the same conclusion is
uniform whenever the left side of
\eqref{eq:YMA4-margin-pass-condition} has an $r$-independent upper bound
strictly below the displayed right side.  In particular, a bound
$A_{\beta_r,M_r}\le C/\beta_r$ has the correct post-one-loop scale, but
neither that estimate nor bounds for the generated and topological stocks
are asserted merely from fixed-card asymptotics.
\end{theorem}

\begin{proof}
Use the exact decomposition
\eqref{eq:YMA4-master-decomposition}.  V3 gives
$|\gamma_M(q)|\le U_\mu\omega(q)$.  The beta-flow integral is bounded by
$A_{\beta,M}\omega(q)$, the generated interpolation by
$B_{\beta,\Psi,M}\omega(q)$, and the topological term by
$T_{\beta,\Psi,M}\omega(q)$.  Finally V2 gives
$s_0(q)\ge(5/28)\omega(q)$.  Subtracting the absolute debits from
$(1-\delta)\beta s_0(q)$ proves
\eqref{eq:YMA4-margin-lower} and its consequences.
\end{proof}

\begin{corollary}[A terminating harmonic-versus-curvature alternative]
\label{cor:YMA4-flow-alternative}
Assume finite outward intervals are supplied for the two replica writers
\eqref{eq:YMA4-beta-replica-writer} and
\eqref{eq:YMA4-generated-replica-writer}, their parameter integrals, and
the final topological card.  Then the selected finite bank terminates in
one of the following three mathematical outcomes.
\begin{enumerate}[label=\textup{(V4.\arabic*)},leftmargin=4.7em]
\item \pass: the lower endpoint of
      $h^{\rm loc}-\delta\beta s_0$ is positive on every retained mode.
\item \obstruction: at one retained mode the upper endpoint of the
      \emph{assembled} signed difference is negative.  Only after this direct
      failure may its beta-flow, generated-flow, or topological component be
      named as the surviving debit.
\item \stopstatus: an integral tail, replica interval, support split, or
      zero-mode versus repaired-kernel stock is unresolved.
\end{enumerate}
On the second branch, bounded repaired $\mathsf W_2$ stock together with a
nonvanishing $M^2z$ term is a typed harmonic/topological survivor.  A failed
repaired $\mathsf W_2$ stock with controlled zero mode is a local nonlinear
or generated curvature survivor.  Failure of the sufficient absolute bound
\eqref{eq:YMA4-margin-pass-condition} alone is not an obstruction.
\end{corollary}

\begin{proof}
The direct interval alternatives are exhaustive.  The attribution in the
last paragraph is exactly the decomposition of
\Cref{lem:YMA4-harmonic-Ward-router,prop:YMA4-kernel-sufficient}.  Since
\eqref{eq:YMA4-margin-pass-condition} used triangle inequalities, its
failure only removes that sufficient certificate and cannot reverse the
sign of the assembled physical card.
\end{proof}

% =============================================================================
\section{Finite acquisition, tail closure, and the next literal calculation}
\label{sec:YMA4-acquisition}
% =============================================================================

The beta integral may be certified without sampling an unbounded continuum
of couplings.  Let $\beta_j=2^j\beta$.  Suppose an interval computation or
analytic estimate gives
\begin{equation}
 d_{u,M}\le D_{j,M}
 \quad(\beta_j\le u\le\beta_{j+1}),
 \qquad
 \int_{\beta_J}^{\infty}d_{u,M}\,du\le R_{J,M}.
 \label{eq:YMA4-dyadic-input}
\end{equation}
Then
\begin{equation}
 \boxed{
 A_{\beta,M}
 \le\sum_{j=0}^{J-1}\beta_jD_{j,M}+R_{J,M}.}
 \label{eq:YMA4-dyadic-bound}
\end{equation}
The analogous finite partition of $[0,1]$ applies to the generated flow.
Each $D_{j,M}$ is obtained from the signed two-replica writer before its
four terms are separated.  A tail estimate $d_{u,M}\le C_Mu^{-2}$ gives
$R_{J,M}\le C_M/\beta_J$ at fixed card; a cofinal result requires a common
$C_M$ bound or a direct replacement.

\begin{proposition}[Finite head plus analytic parameter-tail certificate]
\label{prop:YMA4-finite-parameter-certificate}
Assume \eqref{eq:YMA4-dyadic-input}, finite interval enclosures for
$\int_0^1\mathfrak G_{\theta,M}(q)d\theta$, and a finite interval for
$t_{\beta,\Psi,M}(q)$ on every retained mode.  Then
\eqref{eq:YMA4-master-decomposition} gives a rigorous finite interval for
the exact local symbol.  If the lower endpoint clears
$\delta\beta s_0(q)$ on all modes, the result is a finite \pass\ theorem.
If an upper endpoint lies below it, the corresponding mode is a direct
\obstruction.  Otherwise the correct verdict is \stopstatus, with the first
unclosed parameter interval or physical support row recorded.
\end{proposition}

\begin{proof}
Integrate the interval enclosures using
\eqref{eq:YMA4-dyadic-bound}, add them with their actual signs in
\eqref{eq:YMA4-master-decomposition}, and use outward-rounded interval
arithmetic.  The sign conclusions are then ordinary interval inclusion.
\end{proof}

\begin{assessment}[The actual V4 advance]
\label{ass:YMA4-progress}
Version~V3 left one uncontrolled object called the post-one-loop remainder.
Version~V4 resolves its internal structure.  It is the tail integral of the
single scale-homogeneity card \eqref{eq:YMA4-beta-flow-card}; that card has
one exact two-replica implementation, its apparent order-$\beta^{-1}$
terms cancel by Gaussian radial homogeneity, and its first coefficient is
the explicit finite contraction \eqref{eq:YMA4-two-loop-formula}.  The
submitted generated interaction has the independent exact flow
\eqref{eq:YMA4-generated-flow-card}, also on one two-replica cylinder.  The
zero-mode/repaired-kernel router distinguishes a true harmonic survivor from
a curvature-tail failure before either is divided by the Ward factor.

What remains is quantitative rather than structural: evaluate the two-loop
card on the first literal axial calibration, and then establish or acquire
a cutoff-uniform integrable bound for the complete beta-flow, together with
the generated-flow and topological stocks of the actual selected shell.
No general shell, graph, or mediator theorem substitutes for those values.
\end{assessment}

\begin{openproblem}[V5: evaluate the first two-loop card and test the scale-flow tail]
\label{op:YMA4-V5}
Preserve V4 and the frozen $M=4$, $q=\pi/2$ calibration used in V3.  Expand
the same ninety-six ordered plaquette words and product-Haar density to
obtain $U_1,U_2,X_1,X_2,Y_1,Y_2$.  Evaluate
\eqref{eq:YMA4-two-loop-formula} using the existing $256$ momentum fibres
and exact inverse witnesses, with a rigorous outward interval for
$\eta_4(\pi/2)$.  This is the first literal V5 number.

In parallel, populate the signed beta-flow writer
\eqref{eq:YMA4-beta-replica-writer} on a predeclared finite coupling panel,
or prove a Wilson-specific connected estimate which gives
$d_{u,M}\le C u^{-2}$ uniformly on the retained physical family.  Apply the
harmonic-inclusive router before claiming a Ward-relative tail.  Only after
the pure flow closes should the generated interpolation
\eqref{eq:YMA4-generated-replica-writer} and the final topological card be
inserted.  Do not return to a raw order-$\beta$ covariance or introduce a
larger auxiliary shell carrier.
\end{openproblem}

% =============================================================================
\section{Version V4 proof ledger and replay boundary}
\label{sec:YMA4-ledger}
% =============================================================================

\begingroup\small
\begin{longtable}{>{\raggedright\arraybackslash}p{0.27\textwidth}
                  >{\raggedright\arraybackslash}p{0.20\textwidth}
                  >{\raggedright\arraybackslash}p{0.45\textwidth}}
\caption{Version V4 proof-status ledger.}\label{tab:YMA4-ledger}\\
\toprule
\textbf{Row}&\textbf{Status}&\textbf{Exact scope}\\
\midrule
\endfirsthead
\toprule
\textbf{Row}&\textbf{Status}&\textbf{Exact scope}\\
\midrule
\endhead
Pure Wilson beta flow&Proved exactly
&Derivative of the compensated physical correction is the four-term
connected card \eqref{eq:YMA4-beta-flow-card}.\\[0.3em]
Beta-flow replica acquisition&Proved exactly
&One two-copy signed writer acquires all cancellations; no componentwise
source price is introduced.\\[0.3em]
Nonlinear remainder integral&Proved exactly
&$\mathcal N_{\beta,M}=-\int_\beta^\infty\mathfrak D_u\,du$ on every fixed
finite card.\\[0.3em]
Gaussian leading-flow cancellation&Proved exactly
&Radial homogeneity cancels both apparent order-$\beta^{-1}$ pairs.\\[0.3em]
First post-one-loop coefficient&Finite exact formula
&Equation \eqref{eq:YMA4-two-loop-formula} is the complete two-loop
Gaussian cumulant card; its numerical value is unexecuted.\\[0.3em]
Generated-action interpolation&Proved exactly
&One positive $\theta$ family gives the complete Wilson/generated Hessian
change and all cross covariances.\\[0.3em]
Harmonic-inclusive Ward router&Proved exactly
&Relative stock splits into repaired $\mathsf W_2$ and the explicit
$M^2|z|/16$ harmonic price.\\[0.3em]
Selected local margin&Sufficient finite theorem
&Equation \eqref{eq:YMA4-margin-pass-condition} closes the bank once the
three remaining physical stocks are populated.\\[0.3em]
Uniform beta-flow tail&Open physical estimate
&Fixed-card $O_M(u^{-2})$ is proved; a common connected bound is not.\\[0.3em]
Generated/topological cofinal stocks&Open physical values
&No submitted $\Psi$ or topological term is silently set to zero.\\[0.3em]
Higher jets, local-vacuum contraction, mass&Not inferred
&The inherited handoff opens only after the actual quadratic margin passes.\\
\bottomrule
\end{longtable}
\endgroup

The accompanying exact replay checks three algebraic layers only.  First,
it verifies the beta and generated two-replica identities on rational finite
sample spaces.  Second, it verifies the Gaussian homogeneity identities by
exact differentiation of a quadratic Gaussian determinant.  Third, it
expands a rational finite-space perturbation through second order and checks
every sign and coefficient in \eqref{eq:YMA4-two-loop-formula}.  These are
replays of the finite algebra, not a formal proof of the analytic Laplace
remainder and not a numerical acquisition of the Wilson two-loop value.

\begin{assessment}[V4 termination]
The exact flow reductions, Gaussian cancellation, two-loop formula,
generated interpolation, and harmonic router are \pass\ at their stated
finite mathematical scopes.  The selected response-matched cofinal symbol
remains \stopstatus: neither the uniform beta-flow tail nor the submitted
generated/topological stocks have been populated.  A negative direct mode
has not been produced.  The next calculation is the finite
$\eta_4(\pi/2)$ card, followed by the scale-flow tail on the same source.
\end{assessment}

\section*{Version V4 append-only continuation record}
\addcontentsline{toc}{section}{Version V4 append-only continuation record}
The complete V3 source is preserved byte for byte up to its sole final
\verb|\end{document}| command.  Its SHA--256 digest is
\begin{center}\scriptsize\ttfamily
205c477d5a6366c96f2eca72dfd1985916c3d739839ed469a3378cd357a4531d.
\end{center}
For V5, preserve this accumulated source, append new material immediately
before the final document-closing command, and increment the filename to
\texttt{\_V5.tex}.  Keep the physical response, exact block map,
parity-tree section, Haar compensation, selected axial mode, source
normalization, and local/topological support convention unchanged.

% =============================================================================
% END APPENDED VERSION V4 MATERIAL
% =============================================================================

% =============================================================================
% START APPENDED VERSION V5 MATERIAL
% =============================================================================
\clearpage
\hypersetup{
 pdftitle={Renewal Yang--Mills Reduced-Fibre Wilson Shell Symbol and Local-Vacuum Programme V5},
 pdfsubject={Exact fibre-gap audit, sharp dyadic reference floor, gauge-shorted Witten residual, and small-field inverse locality}
}
\pdfinfo{
 /Title (Renewal Yang--Mills Reduced-Fibre Wilson Shell Symbol and Local-Vacuum Programme V5)
 /Subject (Exact fibre-gap audit, sharp dyadic reference floor, gauge-shorted Witten residual, and small-field inverse locality)
}
\providecommand{\Hhi}{\widehat H^{\mathrm{hi}}}
\providecommand{\Shi}{\widehat S^{\mathrm{hi}}_0}
\providecommand{\Ric}{\operatorname{Ric}}
\providecommand{\Hess}{\operatorname{Hess}}
\providecommand{\Sch}{\operatorname{Sch}}
\providecommand{\rank}{\operatorname{rank}}
\providecommand{\supp}{\operatorname{supp}}
\providecommand{\Per}{\operatorname{Per}}
\providecommand{\cell}{\operatorname{cell}}

\section{Version V5: exact fibre-gap audit, sharp reference floor, and the gauge-shorted Witten residual}
\label{sec:YMA5-direction}

The complete Version~V4 source is preserved above.  The attached
\emph{Yang--Mills margin report} is used as a new comparison source.  Its
central structural contribution is the exact Bochner--Schur identity for the
same reduced-fibre score--Hessian card, together with a claimed exact gap for
the flat parity-tree fibre.  Its SHA--256 digest is
\begin{center}\scriptsize\ttfamily
b070312fad2643646de1e2881b9ec9fa08dbc04ae146942091859c62152c54ee.
\end{center}
The report is not copied into this continuation.  Instead, Version~V5 does
four things which are needed before that material can be used in the present
append-only programme.

First, the exact fibre-gap calculation is reconstructed from the literal
$45$ reduced-link types and $96$ plaquette rows, and is supplied with an
independent exact-arithmetic replay.  Second, the Wilson normalization is
aligned with Versions~V2--V4: the number $\beta/6$ is the coefficient of the
quadratic Taylor term, whereas the Hessian multiplier is $\beta/3$.  Third,
the formerly numerical all-momentum reference lower constant and gauge-lift
norm are proved exactly:
\begin{equation}
 \boxed{
 \mathsf S^{\rm tree}(k)|_{\perp}
 \succeq\frac14|p(k)|^2I,
 \qquad
 \norm{a_0}^2=7\norm{d_0^{\rm c}\xi}^2.}
 \label{eq:YMA5-two-exact-constants}
\end{equation}
The constant $1/4$ is sharp.  Fourth, the positive Witten residual in the
Bochner--Schur formula is shorted against the gauge orbit exactly.  This
shows that the gauge-covariant floor differs from the physical Hessian by a
specific positive Schur residual.  It also corrects the logical status of a
residual inequality in the report: dropping this positive term gives a
sufficient certificate, not an equivalent characterization.

The exact fibre gap then yields a Chebyshev--Demko inverse bound which is much
sharper than the Neumann-series bound used in V3.  The same argument gives a
cutoff-independent spatial decay theorem for the curvature-augmented inverse
on the Wilson small-field region.  This closes the inverse-locality part of
the report's small-field row.  It does not determine the sign of the
expected gauge-shorted head and it does not control the large-field Witten
current.  Those remain the two physical acquisitions.

\begin{firewall}[The report changes the order, not the physical law]
\label{fire:YMA5-report-scope}
The new comparison source does not replace the selected response-matched
law, the parity-tree section, the local/topological split, or the V4
coupling-flow identity.  In particular, the finite two-loop number
$\eta_4(\pi/2)$ is no longer the first compulsory calculation: it becomes a
useful calibration only after the gauge-shorted small-field head has been
formed.  No coefficient marked merely as computed in the report is promoted
to an exact theorem below without a proof or an exact replay.
\end{firewall}

% =============================================================================
\section{Normalization audit for the flat Wilson Hessian}
\label{sec:YMA5-normalization}
% =============================================================================

Retain the unit Wilson action of V2,
\begin{equation}
 V_M(U)=\sum_p\left(1-\frac13\Re\Tr U_p\right),
 \label{eq:YMA5-unit-Wilson}
\end{equation}
and let $d_1$ be the abelianized fine curl in one normalized colour
direction.  Write $\Sigma=d_1^*d_1$ for the corresponding link form.  Its
reduced/restrained block decomposition is
\begin{equation}
 \Sigma(k)=
 \begin{pmatrix}
  \Sigma_{rr}(k)&\Sigma_{rc}(k)\\
  \Sigma_{cr}(k)&\Sigma_{cc}(k)
 \end{pmatrix},
 \label{eq:YMA5-Sigma-block}
\end{equation}
where the four constrained coordinates are the designated second
constituents carrying the coarse-link variation.

\begin{proposition}[Hessian multiplier versus quadratic Taylor coefficient]
\label{prop:YMA5-Wilson-normalization}
At a flat configuration, in the metric
$\ip{X}{Y}=-\Tr(XY)$,
\begin{equation}
 \boxed{
 D^2(\beta V_M)(0)[a,b]
 =\frac{\beta}{3}\ip{d_1a}{d_1b}.}
 \label{eq:YMA5-Hessian-normalization}
\end{equation}
Equivalently, if
\begin{equation}
 \alpha_\beta:=\frac{\beta}{3},
 \label{eq:YMA5-alpha-beta}
\end{equation}
then the vertical Hessian and the flat coarse Schur reference are
\begin{equation}
 \boxed{
 A_{\beta,M}^{\rm vert}=\alpha_\beta\Sigma_{rr},
 \qquad
 S_{\beta,M}^{\rm flat}=\alpha_\beta\mathsf S^{\rm tree},
 \qquad
 \mathsf S^{\rm tree}
 =\Sigma_{cc}-\Sigma_{cr}\Sigma_{rr}^{-1}\Sigma_{rc}.}
 \label{eq:YMA5-normalized-blocks}
\end{equation}
The number $\beta/6$ is the coefficient of
$\norm{d_1a}^2$ in the Taylor expansion of the action, not the multiplier of
its bilinear Hessian.  With the product-group Ricci constant
$\kappa=3/2$, the dimensionless curvature ratio is therefore
\begin{equation}
 \boxed{\frac{\kappa}{\alpha_\beta}=\frac{9}{2\beta}.}
 \label{eq:YMA5-Ricci-ratio}
\end{equation}
\end{proposition}

\begin{proof}
For a skew-Hermitian traceless matrix $X$,
\[
 1-\frac13\Re\Tr(e^{tX})
 =-\frac{t^2}{6}\Tr(X^2)+O(t^3)
 =\frac{t^2}{6}\norm{X}^2+O(t^3).
\]
At a flat plaquette the first logarithmic variation is $(d_1a)_p$; hence
\[
 \beta V_M(e^{ta})
 =\frac{\beta t^2}{6}\norm{d_1a}^2+O(t^3).
\]
Twice differentiating gives
$D^2(\beta V_M)(0)[a,a]=(\beta/3)\norm{d_1a}^2$.
Polarization gives \eqref{eq:YMA5-Hessian-normalization}.  Restriction to the
vertical block and minimization over it give
\eqref{eq:YMA5-normalized-blocks}.  Finally
$(3/2)/(\beta/3)=9/(2\beta)$.
\end{proof}

\begin{corollary}[Alignment with the V2 convention]
\label{cor:YMA5-V2-alignment}
The unit forms used in V2 satisfy
\begin{equation}
 A_M=\frac13\Sigma_{rr},
 \qquad
 \mathsf S_M=\frac13\mathsf S^{\rm tree}.
 \label{eq:YMA5-V2-unit-forms}
\end{equation}
Consequently every generalized ratio in V2--V4 is unchanged.  Only formulas
which identify $\beta/6$ itself with a Hessian multiplier require the factor
two correction above.
\end{corollary}

\begin{proof}
Set $\beta=1$ in \eqref{eq:YMA5-normalized-blocks} and compare with the
variational definition \eqref{eq:YMA2-unit-reference}.  Multiplying both
numerator and reference by the same positive factor leaves every generalized
Rayleigh quotient invariant.
\end{proof}

% =============================================================================
\section{Independent exact replay of the selected-shell fibre gap}
\label{sec:YMA5-gap}
% =============================================================================

Let $\lambda_*$ denote the bottom of $\Sigma_{rr}(k)$ over the continuous
coarse Brillouin torus.  The cohomological kernel argument in the attached
report shows that every $\Sigma_{rr}(k)$ is strictly positive.  We now make
the claimed sharp value an executable theorem in this series.

For one parity cell $B=\{0,1\}^4$, let $D_B$ be the integer incidence matrix
whose $96$ rows are the fine plaquettes based in $B$.  Delete every tree and
designated-second-constituent column.  The remaining columns split into the
$45$ reduced links based in $B$ and the $56$ reduced links based in adjacent
cells:
\begin{equation}
 D_B=(D_c\;D_n).
 \label{eq:YMA5-cell-incidence}
\end{equation}
The neighboring block has rank $38$.  If $B_n$ is any full-rank choice of
$38$ columns spanning $\Ran D_n$, define the exact free-boundary cell short
\begin{equation}
 \mathsf C_B
 =D_c^*D_c-D_c^*B_n(B_n^*B_n)^{-1}B_n^*D_c.
 \label{eq:YMA5-cell-short}
\end{equation}
The right side is independent of the chosen basis of $\Ran D_n$.

\begin{theorem}[Exact parity-tree fibre gap]
\label{thm:YMA5-exact-gap}
Put
\begin{equation}
 \begin{aligned}
 q_8(\lambda)={}&\lambda^8-16\lambda^7+98\lambda^6-296\lambda^5
 +472\lambda^4\\
 &-396\lambda^3+168\lambda^2-32\lambda+2.
 \end{aligned}
 \label{eq:YMA5-q8}
\end{equation}
Then
\begin{equation}
 \boxed{
 \lambda_*
 =\min_{k\in\T^4}\lambda_{\min}(\Sigma_{rr}(k))}
 \label{eq:YMA5-gap-definition}
\end{equation}
is the unique root of $q_8$ in $(1/9,3/20)$ and obeys
\begin{equation}
 \boxed{
 0.11760418512646926265
 <\lambda_*<
 0.11760418512646926269.}
 \label{eq:YMA5-gap-interval}
\end{equation}
For every finite even torus and every commuting toron colour character,
\begin{equation}
 \boxed{
 \Sigma_{rr}\succeq\lambda_* I,
 \qquad
 A_M=\frac13\Sigma_{rr}\succeq\frac{\lambda_*}{3}I.}
 \label{eq:YMA5-exact-vertical-floor}
\end{equation}
In particular,
\begin{equation}
 0.039201395042156420883
 <\frac{\lambda_*}{3}<
 0.039201395042156420897,
 \label{eq:YMA5-gap-third}
\end{equation}
which strictly strengthens the V2 lower bound $1/60$ at the flat point.
\end{theorem}

\begin{proof}
The cell localization argument gives
$\Sigma_{rr}(k)\succeq\mathsf C_B$ for every Bloch momentum: in physical
space, the plaquette rows based in all parity cells partition the full curl
energy, while the $45$ core-link norms partition the reduced-link norm.
Minimizing the $56$ neighboring variables in one cell gives exactly
\eqref{eq:YMA5-cell-short}.

The accompanying exact replay reconstructs $D_B$ from the parity rule,
checks the counts $(96,45,56)$ and $\rank D_n=38$, and performs fraction-free
elimination.  It obtains an integer $45\times45$ matrix $\mathsf C_B$ with
characteristic polynomial
\begin{equation}
 \begin{aligned}
 \det(\lambda I-\mathsf C_B)
 ={}&(\lambda-2)(\lambda-8)^2(\lambda-1)^2
 (\lambda^2-4\lambda+2)\\
 &\times(\lambda^4-8\lambda^3+18\lambda^2-12\lambda+2)
 q_8(\lambda)r_{26}(\lambda),
 \end{aligned}
 \label{eq:YMA5-cell-charpoly}
\end{equation}
where
\begin{equation}
\begin{aligned}
 r_{26}(\lambda)={}&\lambda^{26}-70\lambda^{25}+2294\lambda^{24}
 -46790\lambda^{23}+666127\lambda^{22}-7036382\lambda^{21}\\
 &+57233266\lambda^{20}-367229404\lambda^{19}
 +1889127345\lambda^{18}-7877197032\lambda^{17}\\
 &+26814942584\lambda^{16}-74839585336\lambda^{15}
 +171585326638\lambda^{14}-323156367496\lambda^{13}\\
 &+499066626657\lambda^{12}-629850714950\lambda^{11}
 +646351813884\lambda^{10}-535706275134\lambda^9\\
 &+355485437358\lambda^8-186754350404\lambda^7
 +76540646288\lambda^6-23995177542\lambda^5\\
 &+5598038393\lambda^4-933456610\lambda^3
 +104318104\lambda^2-6945892\lambda+206763.
\end{aligned}
\label{eq:YMA5-r26}
\end{equation}
Exact Sturm counts give no root of any factor other than $q_8$ in
$(0,3/20]$.  For $q_8$, the Sturm variation is eight at $0$ and $1/9$, and
seven at $3/20$.  It is eight at the lower rational endpoint in
\eqref{eq:YMA5-gap-interval} and seven at the upper endpoint.  Thus
$\lambda_{\min}(\mathsf C_B)$ is the unique displayed root and
$\Sigma_{rr}(k)\succeq\lambda_*I$.

For the reverse inequality take
\begin{equation}
 k_*=(\pi,\pi,\pi,0).
 \label{eq:YMA5-kstar}
\end{equation}
All Bloch phases are integral signs and exact elimination gives
\begin{equation}
 \begin{aligned}
 \det(\lambda I-\Sigma_{rr}(k_*))
 ={}&(\lambda-10)^{14}(\lambda-8)^2(\lambda-6)^{14}
 (\lambda-2)\\
 &\times(\lambda^2-4\lambda+2)
 (\lambda^4-8\lambda^3+18\lambda^2-12\lambda+2)q_8(\lambda).
 \end{aligned}
 \label{eq:YMA5-Bloch-charpoly}
\end{equation}
The same Sturm count and positivity of $\Sigma_{rr}(k_*)$ show that its
least eigenvalue is $\lambda_*$.  Hence the cell lower bound is sharp and
\eqref{eq:YMA5-gap-definition} follows.  A commuting toron character merely
translates $k$ on $\T^4$, so it has the same minimum.  Finally use
\eqref{eq:YMA5-V2-unit-forms}.
\end{proof}

\begin{firewall}[Meaning of the exact computation]
\label{fire:YMA5-gap-replay}
The theorem is a finite exact-arithmetic statement about the declared
parity-tree incidence.  The replay reconstructs its matrices rather than
reading stored characteristic polynomials.  It is not a formal verification
of the computer algebra implementation, and the flat vertical gap is not a
Witten one-form gap, a coarse-field margin, or a Yang--Mills mass.
\end{firewall}

% =============================================================================
\section{A sharp all-momentum dyadic reference floor}
\label{sec:YMA5-reference}
% =============================================================================

The report supplied a grid value $0.2500$ for the lower ratio of
$\mathsf S^{\rm tree}$ to coarse curl.  The value is in fact exact and
requires no momentum scan.

Let $d_1^{\rm c}$ be the coarse abelian curl.  If $b$ is a fine one-form
whose straight two-link block is the coarse one-form $h$, then on every
coarse plaquette $P$ the linear Stokes identity is
\begin{equation}
 (d_1^{\rm c}h)_P
 =\sum_{p\subset P}(d_1b)_p,
 \label{eq:YMA5-four-plaquette-Stokes}
\end{equation}
where the sum contains the four fine plaquettes tiling $P$ with their
induced orientations.

\begin{theorem}[Sharp dyadic Stokes floor]
\label{thm:YMA5-sharp-reference-floor}
For every coarse momentum and every coarse one-form,
\begin{equation}
 \boxed{
 \mathsf S^{\rm tree}(k)
 \succeq\frac14\,d_1^{\rm c}(k)^*d_1^{\rm c}(k).}
 \label{eq:YMA5-reference-operator-floor}
\end{equation}
Consequently, on the nonzero transverse carrier,
\begin{equation}
 \boxed{
 \mathsf S^{\rm tree}(k)|_\perp
 \succeq\frac14|p(k)|^2I,
 \qquad
 S_{\beta,M}^{\rm flat}(k)|_\perp
 \succeq\frac{\beta}{12}|p(k)|^2I.}
 \label{eq:YMA5-reference-transverse-floor}
\end{equation}
The constant $1/4$ is sharp.  At $k_*$ of
\eqref{eq:YMA5-kstar},
\begin{equation}
 \mathsf S^{\rm tree}(k_*)=
 \begin{pmatrix}
  5/2&-5/4&-5/4&0\\
  -5/4&5/2&-5/4&0\\
  -5/4&-5/4&5/2&0\\
  0&0&0&3
 \end{pmatrix},
 \label{eq:YMA5-reference-kstar}
\end{equation}
and the transverse vector $e_4$ has quotient exactly $1/4$ relative to
$d_1^{\rm c*}d_1^{\rm c}$.
\end{theorem}

\begin{proof}
By Cauchy--Schwarz in \eqref{eq:YMA5-four-plaquette-Stokes},
\[
 \abs{(d_1^{\rm c}h)_P}^2
 \le4\sum_{p\subset P}\abs{(d_1b)_p}^2.
\]
The four-plaquette tiles of distinct coarse plaquettes are disjoint.  They
form only a subset of all fine plaquettes, so summation gives
\begin{equation}
 \norm{d_1^{\rm c}h}^2\le4\norm{d_1b}^2.
 \label{eq:YMA5-Stokes-energy}
\end{equation}
Minimize the right side over every fine completion of $h$.  The parity-tree
reduced parametrization is exact, and the minimizer exists uniquely in its
vertical variables by \Cref{thm:YMA5-exact-gap}; this proves
\eqref{eq:YMA5-reference-operator-floor}.  On a nonzero Fourier mode,
$d_1^{\rm c*}d_1^{\rm c}=|p|^2I-d_0^{\rm c}d_0^{\rm c*}$, hence it equals
$|p|^2I$ on the transverse carrier.  The physical normalization follows
from \Cref{prop:YMA5-Wilson-normalization}.

It remains to prove sharpness without a floating-point minimization.  On the
Bloch character $k_*$ take coarse input $h=e_4$ and define a reduced
completion by the following cell-type rule, with
$s=(s_1,s_2,s_3,s_4)\in\{0,1\}^4$:
\begin{equation}
 a(s,\mu)=
 \begin{cases}
  -\frac12,&\mu=4,\ s_4=0,\ (s_1,s_2,s_3)\ne0,\\
  \frac12,&\mu=4,\ s_4=1,\ (s_1,s_2,s_3)\ne0,\\
  1,&1\le\mu\le3,\ s_\mu=s_4=1,\\
  0,&\text{otherwise}.
 \end{cases}
 \label{eq:YMA5-sharp-completion}
\end{equation}
Extend it by the Bloch sign $(-1)^{m_1+m_2+m_3}$ from cell to cell.  Direct
substitution in the $96$ plaquette rows gives curl $-1/2$ precisely on the
twelve plaquettes $(\mu,4)$ based at
$s=t e_\mu+u e_4$, where $1\le\mu\le3$ and $t,u\in\{0,1\}$, and gives zero
on all other plaquettes.  Therefore
\begin{equation}
 \norm{d_1(a,h)}^2=12\left(\frac12\right)^2=3.
 \label{eq:YMA5-sharp-fine-energy}
\end{equation}
The coarse curl has three entries equal to $-2$, so
$\norm{d_1^{\rm c}h}^2=12$.  Equality holds in
\eqref{eq:YMA5-Stokes-energy}; hence the constant is sharp.  Solving the
normal equations, or evaluating the bilinear form on the remaining three
coordinate inputs, gives the exact matrix
\eqref{eq:YMA5-reference-kstar}.
\end{proof}

\begin{corollary}[Exact replacement for the computed lower reference constant]
\label{cor:YMA5-cminus}
In every later use of the report's lower reference constant one may set
\begin{equation}
 \boxed{c_-=\frac14}
 \label{eq:YMA5-cminus}
\end{equation}
for the unscaled form $\mathsf S^{\rm tree}$.  For the V2 unit action the
corresponding all-momentum lower constant is $1/12$.  The axial formula
\eqref{eq:YMA2-reference-rational} remains the sharper exact statement on
that one-dimensional orbit.
\end{corollary}

% =============================================================================
\section{The exact gauge completion and the corrected Ricci scale}
\label{sec:YMA5-gauge-lift}
% =============================================================================

For a nonzero coarse Bloch character let
$h=d_0^{\rm c}\xi$.  Let $a_0$ be the unique reduced component for which
$(a_0,h)$ is a closed fine form vanishing on the parity-tree links.  By the
normal equations,
\begin{equation}
 a_0=-\Sigma_{rr}^{-1}\Sigma_{rc}h.
 \label{eq:YMA5-gauge-completion}
\end{equation}

\begin{theorem}[Seven-copy gauge completion identity]
\label{thm:YMA5-gauge-seven}
For every nonzero Bloch momentum, every colour component, and every
commuting toron character,
\begin{equation}
 \boxed{\norm{a_0}^2=7\norm{h}^2.}
 \label{eq:YMA5-gauge-seven}
\end{equation}
Thus the value seven observed in the margin report is exact at every
momentum, rather than merely a grid maximum.
\end{theorem}

\begin{proof}
Let the fine gauge potential be constant on each parity tree:
\[
 \varphi(2m+s)=e^{ik\cdot m}\xi,
 \qquad s\in\{0,1\}^4.
\]
A positive fine link based at $2m+s$ in direction $\mu$ has zero gradient
when $s_\mu=0$, because both endpoints have the same coarse root.  When
$s_\mu=1$, its endpoints have roots $m$ and $m+e_\mu$, so its gradient is
$e^{ik\cdot m}h_\mu$.  Among the eight parity types with $s_\mu=1$, the type
$s=e_\mu$ is the designated second constituent.  The other seven types are
exactly reduced links: each has at least one additional occupied parity
coordinate and hence a positive filling length.  Therefore each coarse
component $h_\mu$ appears on exactly seven reduced links per cell, with only
unit-modulus Bloch or toron phases.  Summation over directions, cells, and
colours proves \eqref{eq:YMA5-gauge-seven}.  The constructed completion has
zero curl; uniqueness follows from the positive vertical gap in
\Cref{thm:YMA5-exact-gap}, so it equals
\eqref{eq:YMA5-gauge-completion}.
\end{proof}

At the flat fibre point, augment the vertical block by the product-group
Ricci curvature.  With $\alpha_\beta=\beta/3$ and
$\varepsilon_\beta=\kappa/\alpha_\beta$, the resulting pointwise coarse
short is
\begin{equation}
 \begin{aligned}
 \mathsf P_{\beta,\kappa}(k)
 &:={}
 \alpha_\beta\left(
 \Sigma_{cc}-\Sigma_{cr}
 (\Sigma_{rr}+\varepsilon_\beta I)^{-1}\Sigma_{rc}
 \right)\\
 &=\alpha_\beta\mathsf S^{\rm tree}(k)
   +\Delta_{\Ric,\beta}(k),
 \end{aligned}
 \label{eq:YMA5-Ricci-short}
\end{equation}
where
\begin{equation}
 \boxed{
 \Delta_{\Ric,\beta}
 =\kappa\Sigma_{cr}\Sigma_{rr}^{-1}
  (\Sigma_{rr}+\varepsilon_\beta I)^{-1}\Sigma_{rc}
 \succeq0.}
 \label{eq:YMA5-Ricci-bonus}
\end{equation}

\begin{proposition}[Exact gauge artefact bound and the method-level infrared window]
\label{prop:YMA5-corrected-IR}
On a gauge direction $h=d_0^{\rm c}\xi\ne0$,
\begin{equation}
 \boxed{
 0<\Delta_{\Ric,\beta}[h]
 \le7\kappa\norm{h}^2
 =\frac{21}{2}\norm{h}^2.}
 \label{eq:YMA5-Ricci-seven-bound}
\end{equation}
Suppose a pointwise-curvature argument, before the exact gauge-current short
of the next section, yields only an absolute transverse estimate
\begin{equation}
 \Hhi(k)\succeq
 \alpha_\beta\mathsf S^{\rm tree}(k)-\Xi I.
 \label{eq:YMA5-absolute-floor}
\end{equation}
Then it proves a target margin $\delta\in(0,1)$ on every mode satisfying
\begin{equation}
 \boxed{
 |p(k)|^2\ge
 \frac{12\Xi}{(1-\delta)\beta}.}
 \label{eq:YMA5-corrected-window}
\end{equation}
If the proof separately pays the worst possible Ricci gauge term
$\Xi=7\kappa$, its unclosed window is
\begin{equation}
 \boxed{|p(k)|^2<\frac{126}{(1-\delta)\beta}.}
 \label{eq:YMA5-126-window}
\end{equation}
This is a limitation of that termwise proof, not a lower bound on the exact
physical loss and not an infrared Yang--Mills obstruction.
\end{proposition}

\begin{proof}
For a gauge direction,
$\Sigma_{rc}h=-\Sigma_{rr}a_0$.  Substitution into
\eqref{eq:YMA5-Ricci-bonus} gives
\[
 \Delta_{\Ric,\beta}[h]
 =\kappa\ip{a_0}{
  \Sigma_{rr}(\Sigma_{rr}+\varepsilon_\beta I)^{-1}a_0}.
\]
The middle operator has spectrum in $(0,1)$, proving strict positivity and
the upper bound $\kappa\norm{a_0}^2$.  Apply
\Cref{thm:YMA5-gauge-seven}.  For the second assertion,
\Cref{thm:YMA5-sharp-reference-floor} gives
$\alpha_\beta\mathsf S^{\rm tree}\succeq
(\beta/12)|p|^2I$ on the transverse carrier.  The inequality
$\Xi\le(1-\delta)(\beta/12)|p|^2$ is exactly
\eqref{eq:YMA5-corrected-window}; inserting
$7\kappa=21/2$ gives \eqref{eq:YMA5-126-window}.
\end{proof}

\begin{firewall}[The Ward cancellation may remove the absolute debit]
\label{fire:YMA5-IR-not-obstruction}
The gauge Ward identity forces the expected pointwise Schur row, the
Witten-current residual, and the large-field term to cancel exactly on the
gauge subspace.  It does not force the assembled transverse loss to have
size $7\kappa$.  Accordingly, \eqref{eq:YMA5-126-window} applies only to a
proof which bounds the flat Ricci bonus and the fluctuation correction
separately in absolute norm.  A gauge-shorted or Ward-relative estimate may
close below that scale.
\end{firewall}

% =============================================================================
\section{The gauge-covariant floor is an exact Witten-current short}
\label{sec:YMA5-current-short}
% =============================================================================

We now use the exact Bochner--Schur identity established in the attached
margin report.  On one fixed nonzero momentum--colour block it has the form
\begin{equation}
 \boxed{H=Q+\mathscr G,\qquad
 Q=\mathbb E\Sch_{\mathcal R}-\mathscr T_\Theta,}
 \label{eq:YMA5-Bochner-import}
\end{equation}
where $H=\Hhi$ annihilates the coarse gauge space and
$\mathscr G\succeq0$ is the quadratic form
\begin{equation}
 \mathscr G[h]
 =\norm{\nabla\omega_h}_{L^2(\nu)}^2
 +\norm{\mathcal R^{1/2}
  (\omega_h-\mathcal R^{-1}\mathscr B h)}_{L^2(\nu)}^2.
 \label{eq:YMA5-G-form}
\end{equation}
Here $\omega_h=L_1^{-1}\mathscr B h$ is the exact Witten current.  Let
$E=E_\perp\oplus E_g$ be the orthogonal transverse/gauge decomposition and
write the corresponding blocks of $\mathscr G$ as
$G_{\perp\perp},G_{\perp g},G_{g\perp},G_{gg}$.

\begin{theorem}[Exact gauge-current Schur identity]
\label{thm:YMA5-current-short}
Define the report's gauge-covariant floor by
\begin{equation}
 \mathfrak F_{\mathcal R}[h]
 :=\sup_{g\in E_g}Q[h+g],
 \qquad h\in E_\perp.
 \label{eq:YMA5-covariant-floor}
\end{equation}
Then the range condition
$\Ran G_{g\perp}\subseteq\Ran G_{gg}$ holds and
\begin{equation}
 \boxed{
 \begin{aligned}
 \mathscr G_{\perp\mid g}
 &:={}
 G_{\perp\perp}
 -G_{\perp g}G_{gg}^{\dagger}G_{g\perp}
 \succeq0,\\
 \mathfrak F_{\mathcal R}|_{E_\perp}
 &=H|_{E_\perp}-\mathscr G_{\perp\mid g},\\
 H|_{E_\perp}
 &=\mathfrak F_{\mathcal R}|_{E_\perp}
   +\mathscr G_{\perp\mid g}.
 \end{aligned}}
 \label{eq:YMA5-exact-current-short}
\end{equation}
Equivalently,
\begin{equation}
 \boxed{
 \mathscr G_{\perp\mid g}[h]
 =\inf_{g\in E_g}\mathscr G[h+g],}
 \label{eq:YMA5-current-infimum}
\end{equation}
and a minimizing gauge correction is
$g_*=-G_{gg}^{\dagger}G_{g\perp}h$ modulo $\Ker G_{gg}$.
\end{theorem}

\begin{proof}
Because $H$ is self-adjoint and $E_g\subseteq\Ker H$, its gauge row and
column vanish.  From $Q=H-\mathscr G$ one therefore has
\begin{equation}
 Q_{\perp\perp}=H_{\perp\perp}-G_{\perp\perp},
 \qquad
 Q_{\perp g}=-G_{\perp g},
 \qquad
 Q_{gg}=-G_{gg}.
 \label{eq:YMA5-Q-blocks}
\end{equation}
Positivity of the block form $\mathscr G$ implies
$\Ran G_{g\perp}\subseteq\Ran G_{gg}$: if
$z\in\Ker G_{gg}$, positivity of
$\mathscr G[(th,z)]$ for all real $t$ forces
$\ip{z}{G_{g\perp}h}=0$.

For fixed $h\in E_\perp$, gauge invariance of $H$ gives
\[
 Q[h+g]=H[h]-\mathscr G[h+g].
\]
Taking the supremum in $g$ proves
$\mathfrak F_{\mathcal R}[h]=H[h]-\inf_g\mathscr G[h+g]$.
The Moore--Penrose completion of the positive block square gives
\[
 \inf_g\mathscr G[h+g]
 =\ip{h}{
  (G_{\perp\perp}-G_{\perp g}G_{gg}^{\dagger}G_{g\perp})h},
\]
with the stated minimizer.  This proves all identities.
\end{proof}

\begin{corollary}[Correct logical status of the covariant floor]
\label{cor:YMA5-floor-status}
For every $\delta\ge0$,
\begin{equation}
 \mathfrak F_{\mathcal R}\succeq\delta\Shi
 \quad\Longrightarrow\quad
 H\succeq\delta\Shi.
 \label{eq:YMA5-floor-sufficient}
\end{equation}
The converse need not hold: a negative direction of
$\mathfrak F_{\mathcal R}-\delta\Shi$ is an obstruction only if the positive
current short $\mathscr G_{\perp\mid g}$ is too small to repair it.  Thus a
failed covariant-floor test is \stopstatus, not a physical negative-mode
certificate, unless the residual in \eqref{eq:YMA5-exact-current-short} is
also acquired.
\end{corollary}

\begin{proof}
The implication and qualification are immediate from the last identity in
\eqref{eq:YMA5-exact-current-short} and positivity of
$\mathscr G_{\perp\mid g}$.
\end{proof}

\begin{theorem}[Exact and sufficient residual certificates]
\label{thm:YMA5-corrected-certificate}
Write the imported decomposition relative to the physical flat reference as
\begin{equation}
 H=\Shi+\mathfrak E+\mathscr G-\mathscr T_\Theta.
 \label{eq:YMA5-reference-decomposition}
\end{equation}
Then the target margin is exactly equivalent to
\begin{equation}
 \boxed{
 \mathscr T_\Theta-\mathfrak E-\mathscr G
 \preceq(1-\delta)\Shi.}
 \label{eq:YMA5-exact-residual-certificate}
\end{equation}
Either of the stronger conditions
\begin{equation}
 \boxed{
 \mathscr T_\Theta-\mathfrak E
 \preceq(1-\delta)\Shi}
 \label{eq:YMA5-sufficient-raw-certificate}
\end{equation}
or
\begin{equation}
 \boxed{\mathfrak F_{\mathcal R}\succeq\delta\Shi}
 \label{eq:YMA5-sufficient-covariant-certificate}
\end{equation}
is sufficient.  Neither is equivalent to the target without an additional
identity which removes the positive Witten-current residual.
\end{theorem}

\begin{proof}
Subtract $\delta\Shi$ from
\eqref{eq:YMA5-reference-decomposition}; this gives
\eqref{eq:YMA5-exact-residual-certificate} with no inequality.  Dropping
$\mathscr G\succeq0$ proves
\eqref{eq:YMA5-sufficient-raw-certificate}.  The final sufficient condition
is \Cref{cor:YMA5-floor-status}.
\end{proof}

\begin{firewall}[The current short is not a second source price]
\label{fire:YMA5-current-no-double-charge}
The form $\mathscr G_{\perp\mid g}$ is the positive residual left by the same
Witten current already appearing in the exact covariance representation.
It is not added to the physical Hessian as a new source.  Equation
\eqref{eq:YMA5-exact-current-short} is a decomposition of one existing card,
and any later replica acquisition must assemble the common current before
shorting its gauge component.
\end{firewall}

% =============================================================================
\section{Sharp flat Green decay and small-field inverse locality}
\label{sec:YMA5-locality}
% =============================================================================

The exact gap has a second use: it supplies a substantially better spatial
inverse bound and turns the requested small-field Combes--Thomas row into a
finite theorem.  Distances below are coarse-cell graph distances.  Two
reduced-link types coupled by one entry of $\Sigma_{rr}$ have cell distance
at most two.

\begin{lemma}[Exact upper edge of the reduced flat curl]
\label{lem:YMA5-flat-upper}
For every Bloch momentum,
\begin{equation}
 \boxed{\lambda_*I\preceq\Sigma_{rr}(k)\preceq16I.}
 \label{eq:YMA5-flat-spectral-window}
\end{equation}
\end{lemma}

\begin{proof}
The full fine-lattice curl symbol is
$|d(p)|^2I-d(p)d(p)^*$ and has nonzero eigenvalue
$|d(p)|^2=4\sum_{\mu=1}^4\sin^2(p_\mu/2)\le16$.
The parity-cell Bloch form is unitarily the direct sum of its sixteen alias
symbols.  Compression to the reduced-link coordinates cannot increase the
operator norm.  The lower bound is
\Cref{thm:YMA5-exact-gap}.
\end{proof}

\begin{lemma}[Chebyshev inverse localization]
\label{lem:YMA5-Chebyshev-localization}
Let $A$ be a positive self-adjoint block operator on a finite torus or on
$\ell^2(\Z^4;\C^m)$, with propagation at most $R_0$ and
\begin{equation}
 0<mI\preceq A\preceq LI.
 \label{eq:YMA5-general-window}
\end{equation}
Put
\begin{equation}
 \rho(m,L)=\frac{\sqrt L-\sqrt m}{\sqrt L+\sqrt m}\in[0,1).
 \label{eq:YMA5-Chebyshev-rho}
\end{equation}
Then for blocks at graph distance $d$,
\begin{equation}
 \boxed{
 \norm{(A^{-1})_{xy}}_{\op}
 \le\frac{2}{m}
 \rho(m,L)^{\lceil d/R_0\rceil}.}
 \label{eq:YMA5-Chebyshev-block}
\end{equation}
\end{lemma}

\begin{proof}
For $n\ge0$, set
\[
 a=\frac{L+m}{L-m}>1,
 \qquad
 r_{n+1}(t)=
 \frac{T_{n+1}((L+m-2t)/(L-m))}{T_{n+1}(a)},
\]
where $T_j$ is the Chebyshev polynomial.  Since $r_{n+1}(0)=1$, the
polynomial $p_n(t)=(1-r_{n+1}(t))/t$ has degree $n$.  On $[m,L]$,
\[
 \abs{t^{-1}-p_n(t)}
 \le\frac1m\abs{T_{n+1}(a)}^{-1}
 \le\frac2m\rho(m,L)^{n+1};
\]
the last identity follows from
$a+\sqrt{a^2-1}=\rho(m,L)^{-1}$.
Functional calculus gives the same operator-norm bound for
$A^{-1}-p_n(A)$.  The block $p_n(A)_{xy}$ vanishes when $d>nR_0$.
Take $n=\lceil d/R_0\rceil-1$; for $d=0$ use
$\norm{A^{-1}}\le1/m$.
\end{proof}

\begin{theorem}[Improved exact flat reduced-fibre propagator]
\label{thm:YMA5-flat-Green}
Define
\begin{equation}
 \boxed{
 \rho_*:=\frac{4-\sqrt{\lambda_*}}{4+\sqrt{\lambda_*}}
 \in(0,1).}
 \label{eq:YMA5-rhostar}
\end{equation}
Then, on the infinite parity-cell lattice and on every finite torus with
shortest periodic distance,
\begin{align}
 \norm{\Sigma_{rr}^{-1}(x,y)}_{\op}
 &\le\frac{2}{\lambda_*}
 \rho_*^{\lceil\dist_{\cell}(x,y)/2\rceil},
 \label{eq:YMA5-Sigma-inverse-tail}\\
 \boxed{
 \norm{G_M(x,y)}_{\op}
 \le\frac{6}{\lambda_*}
 \rho_*^{\lceil\dist_{\cell}(x,y)/2\rceil},}
 \qquad G_M=A_M^{-1}.
 \label{eq:YMA5-unit-Green-tail}
\end{align}
Numerically,
\begin{equation}
 \rho_*<0.842073,
 \qquad
 \frac6{\lambda_*}<51.019,
 \qquad
 -\frac12\log\rho_*>0.08594.
 \label{eq:YMA5-Green-numerics}
\end{equation}
Thus the exponential rate in the V3 one-loop locality proof may be replaced
by any $\mu< -\frac12\log\rho_*$.
\end{theorem}

\begin{proof}
The reduced curl has propagation two.  Apply
\Cref{lem:YMA5-Chebyshev-localization} to
\eqref{eq:YMA5-flat-spectral-window}, with $m=\lambda_*$, $L=16$, and
$R_0=2$.  This gives \eqref{eq:YMA5-Sigma-inverse-tail}.  Since
$A_M=\Sigma_{rr}/3$, its inverse is $3\Sigma_{rr}^{-1}$, proving
\eqref{eq:YMA5-unit-Green-tail}.  The numerical enclosures follow by
outward rational evaluation of \eqref{eq:YMA5-gap-interval}; they are also
checked by the exact replay.
\end{proof}

\begin{corollary}[Sharper V3 weighted inverse stock]
\label{cor:YMA5-V3-stock}
Set
\begin{equation}
 \mu_*=-\frac14\log\rho_*,
 \qquad
 g_*=\frac{6}{\lambda_*}\,45^2
 \left(\frac{1+\rho_*^{1/4}}{1-\rho_*^{1/4}}\right)^4.
 \label{eq:YMA5-improved-stock}
\end{equation}
Then the V3 entrywise norm satisfies $|G|_{\mu_*}\le g_*$.  Every later V3
one-loop kernel and winding estimate remains valid after replacing its
Neumann-series inverse stock by \eqref{eq:YMA5-improved-stock}; no new
one-loop diagram is introduced.
\end{corollary}

\begin{proof}
Each of the $45^2$ scalar entries of a block is bounded by its operator norm.
Moreover
$e^{\mu_*d}\rho_*^{\lceil d/2\rceil}
\le(\rho_*^{1/4})^d$, and
\[
 \sum_{x\in\Z^4}t^{\ellone{x}}
 =\left(\frac{1+t}{1-t}\right)^4.
\]
Insert \eqref{eq:YMA5-unit-Green-tail}.  The remaining V3 estimates use only
submultiplicativity of this weighted stock.
\end{proof}

We next state the spatial inverse result needed by the report's small-field
row.  It is deliberately restricted to the pure Wilson vertical Hessian;
a submitted generated interaction must contribute its own local Hessian
stock.

\begin{proposition}[Uniform small-field curvature window]
\label{prop:YMA5-small-field-window}
There is a finite geometry constant $C_{\rm H}>0$, independent of the torus
side and of a commuting toron background, with the following property.
Suppose every plaquette used by the V2 filling of each reduced link is within
Frobenius chord $\varepsilon$ of its flat value.  For the pure Wilson action,
put
\begin{equation}
 \mathcal R_{\beta}(z)
 =\kappa I+\Hess_z(\beta V_M)(z,I).
 \label{eq:YMA5-small-R}
\end{equation}
If $C_{\rm H}\varepsilon<\lambda_*$, then
\begin{equation}
 \boxed{
 \rho_{\beta,\varepsilon} I
 \preceq\mathcal R_{\beta}(z)
 \preceq L_\beta I,}
 \label{eq:YMA5-small-window}
\end{equation}
where
\begin{equation}
 \rho_{\beta,\varepsilon}
 =\kappa+\alpha_\beta(\lambda_*-C_{\rm H}\varepsilon),
 \qquad
 L_\beta=\kappa+24\alpha_\beta.
 \label{eq:YMA5-small-endpoints}
\end{equation}
The operator has coarse-cell propagation at most two.
\end{proposition}

\begin{proof}
The V2 filling has length at most six.  Telescoping its plaquette holonomies
shows that each reduced link differs from its flat or toron value by at most
$6\varepsilon$ in chord norm.  Chord and geodesic distance are uniformly
equivalent on the compact group.  The Hessian of one plaquette function is a
smooth function of its four links, so on this finite cell grammar its
left-trivialized block coefficients differ from their flat coefficients by
at most $C_{\rm H}\varepsilon$ after increasing one universal constant.
Each link occurs in six plaquettes, and the same finite constant may be
chosen after summing the incident blocks.  The exact flat lower edge is
$\alpha_\beta\lambda_*$ by \Cref{thm:YMA5-exact-gap}; this gives the lower
endpoint.

For the upper endpoint, a second derivative block of
$\beta(1-\Re\Tr U_p/3)$ has operator norm at most
$\alpha_\beta$: after left trivialization it is the real trace of two unit
Lie-algebra insertions separated by unitary transports, and
$|\Tr(XUYV)|\le\norm{X}\norm{Y}$.  A fixed link is coupled to at most four
link slots in each of six plaquettes.  The block Schur row sum is therefore
at most $24\alpha_\beta$, so the self-adjoint Hessian has norm at most that
number.  Add $\kappa I$.  Two links in one plaquette have coarse-cell bases
at distance at most two, proving the propagation statement.
\end{proof}

\begin{theorem}[Cutoff-independent small-field inverse decay]
\label{thm:YMA5-small-field-inverse}
Assume the hypotheses of \Cref{prop:YMA5-small-field-window} and strengthen
them to
\begin{equation}
 C_{\rm H}\varepsilon\le\frac12\lambda_*.
 \label{eq:YMA5-half-gap-small}
\end{equation}
Define
\begin{equation}
 \rho_{\rm sf,*}
 :=\frac{\sqrt{48/\lambda_*}-1}
         {\sqrt{48/\lambda_*}+1}<1.
 \label{eq:YMA5-rho-small-field}
\end{equation}
Then, uniformly in $M$, $\beta>0$, the flat/toron background, and every
small-field configuration in the stated set,
\begin{equation}
 \boxed{
 \norm{\mathcal R_{\beta}(z)^{-1}(x,y)}_{\op}
 \le\frac{2}{\rho_{\beta,\varepsilon}}
 \rho_{\rm sf,*}^{\lceil\dist_{\cell}(x,y)/2\rceil}.}
 \label{eq:YMA5-small-inverse-tail}
\end{equation}
\end{theorem}

\begin{proof}
By \eqref{eq:YMA5-half-gap-small},
$\rho_{\beta,\varepsilon}\ge\kappa+
\alpha_\beta\lambda_*/2$.  Therefore
\[
 \frac{L_\beta}{\rho_{\beta,\varepsilon}}
 \le\max\left\{1,\frac{48}{\lambda_*}\right\}
 =\frac{48}{\lambda_*}.
\]
The function $(\sqrt r-1)/(\sqrt r+1)$ increases in the condition number
$r$.  Apply \Cref{lem:YMA5-Chebyshev-localization} with propagation two and
the spectral window \eqref{eq:YMA5-small-window}.
\end{proof}

\begin{corollary}[The small-field pointwise Schur kernel has a uniform finite head]
\label{cor:YMA5-small-Schur-tail}
Let $\mathscr B_\beta(z)$ be the mixed reduced/coarse Wilson Hessian and
$K_\beta(z)$ the coarse/coarse Wilson Hessian at a small-field configuration
of \Cref{thm:YMA5-small-field-inverse}.  Their propagation is at most two,
and each block column has absolute block sum at most $24\alpha_\beta$.
For cell displacement $x$ with $\ellone{x}>4$,
\begin{equation}
 \boxed{
 \norm{\bigl(\mathscr B_\beta^*
 \mathcal R_\beta^{-1}\mathscr B_\beta\bigr)(0,x)}_{\op}
 \le
 \frac{1152\alpha_\beta^2}{\rho_{\beta,\varepsilon}}
 \rho_{\rm sf,*}^{\lceil(\ellone{x}-4)/2\rceil}.}
 \label{eq:YMA5-small-Schur-decay}
\end{equation}
Consequently, if
$\Sch_{\mathcal R}=K_\beta-
\mathscr B_\beta^*\mathcal R_\beta^{-1}\mathscr B_\beta$, then for every
$R>4$ its omitted second-displacement stock is bounded by
\begin{equation}
 \boxed{
 \frac12\sum_{\ellone{x}>R}x_1^2
 \norm{\Sch_{\mathcal R}(0,x)}_{\op}
 \le
 \frac{9216\alpha_\beta^2}{\rho_{\beta,\varepsilon}}
 \sum_{n>R}n^2(n+1)^3
 \rho_{\rm sf,*}^{\lceil(n-4)/2\rceil}.}
 \label{eq:YMA5-small-Schur-second-tail}
\end{equation}
The series is finite and decays exponentially in $R$, uniformly in the
cutoff.  Thus long-distance propagation through
$\mathcal R^{-1}$ is no longer an open part of the small-field row.
\end{corollary}

\begin{proof}
For fixed coarse cells $0$ and $x$, expand the kernel as a sum over reduced
links $e$ within distance two of $0$ and $f$ within distance two of $x$.
Then
$\dist_{\cell}(e,f)\ge\ellone{x}-4$.  Apply
\eqref{eq:YMA5-small-inverse-tail}, multiply by the two absolute block-column
sums, and use
$2(24\alpha_\beta)^2=1152\alpha_\beta^2$.  The direct block $K_\beta$ has
range two and contributes nothing when $R>4$.  Finally the number of points
of $\Z^4$ at $\ell^1$ distance $n$ is at most $16(n+1)^3$, while
$x_1^2\le n^2$.  Multiplication by the factor $1/2$ gives the displayed
constant.
\end{proof}

\begin{firewall}[Inverse locality is not the signed physical margin]
\label{fire:YMA5-locality-not-sign}
The last theorem controls the spatial tail of the pointwise small-field
Schur operator.  It does not determine the sign of its expected finite head,
does not by itself impose the Ward zero on a separately truncated term, and
does not control the current on the large-field region.  The head must be
assembled with the gauge-current short of
\Cref{thm:YMA5-current-short}; only the complete physical form may then be
passed to the V1 Ward-relative tail theorem.
\end{firewall}

% =============================================================================
\section{Revised high/low-mode handoff and the next physical acquisition}
\label{sec:YMA5-frontier}
% =============================================================================

The report and the preceding corrections change the economical attack.  The
flat vertical algebra is now complete:
\begin{equation}
 \boxed{
 \lambda_*\text{ is exact},\qquad
 \mathsf S^{\rm tree}|_\perp\succeq\tfrac14|p|^2I,
 \qquad
 \norm{a_0}^2=7\norm{d_0^{\rm c}\xi}^2.}
 \label{eq:YMA5-flat-packet-complete}
\end{equation}
The small-field inverse also has a cutoff-independent exponential tail.
What is not known is the signed finite head after expectation and gauge
shorting, together with the large-field current return.

\begin{theorem}[Exact two-regime assembly]
\label{thm:YMA5-two-regime}
Fix $0<\delta<1$ and a momentum threshold $p_*^2>0$.  Suppose the selected
physical bank is partitioned into
\begin{equation}
 \mathcal Q_{\rm hi}=\{k:|p(k)|^2\ge p_*^2\},
 \qquad
 \mathcal Q_{\rm lo}=\{k:0<|p(k)|^2<p_*^2\}.
 \label{eq:YMA5-mode-split}
\end{equation}
Assume:
\begin{enumerate}[label=\textup{(R\arabic*)},leftmargin=4.2em]
\item on $\mathcal Q_{\rm hi}$, an assembled Bochner--Schur estimate gives
      $H(k)\succeq\alpha_\beta\mathsf S^{\rm tree}(k)-\Xi_{\rm hi}I$;
\item on $\mathcal Q_{\rm lo}$, the V4 physical beta/generated/topological
      flow is bounded in the Ward-relative form
      $H(k)\succeq\alpha_\beta\mathsf S^{\rm tree}(k)
      -\Xi_{\rm lo}|p(k)|^2I$.
\end{enumerate}
If
\begin{equation}
 \boxed{
 \Xi_{\rm hi}\le(1-\delta)\frac{\beta}{12}p_*^2,
 \qquad
 \Xi_{\rm lo}\le(1-\delta)\frac{\beta}{12},}
 \label{eq:YMA5-two-regime-condition}
\end{equation}
then $H(k)\succeq\delta S_{\beta,M}^{\rm flat}(k)$ on the entire nonzero
bank.
\end{theorem}

\begin{proof}
On the high bank,
\Cref{thm:YMA5-sharp-reference-floor} and the first inequality in
\eqref{eq:YMA5-two-regime-condition} give
$\Xi_{\rm hi}I\preceq(1-\delta)S_{\beta,M}^{\rm flat}(k)$.
The same reference floor and the second inequality give the corresponding
statement on the low bank.  Insert the two assumed estimates.
\end{proof}

\begin{assessment}[Why the two-loop card is no longer first]
\label{ass:YMA5-new-order}
The coefficient $\eta_4(\pi/2)$ remains a legitimate exact finite
calibration of the V4 beta flow.  The margin report shows, however, that a
pointwise curvature expansion carries a gauge artefact and that the
physically useful lower object is the gauge-shorted expected form.  The next
calculation should therefore acquire
\begin{equation}
 \boxed{
 \mathfrak F_{\mathcal R}(q_{n_*})
 \quad\text{and}\quad
 \mathscr G_{\perp\mid g}(q_{n_*})}
 \label{eq:YMA5-next-two-scalars}
\end{equation}
on the same first selected axial card, using one common Witten-current
cylinder.  Their sum is the already existing physical Hessian by
\eqref{eq:YMA5-exact-current-short}.  Only if the resulting interval still
meets the target threshold is the finite $\eta_4$ card the next useful
low-coupling discriminator.
\end{assessment}

\begin{openproblem}[V6: gauge-shorted finite head and large-field return]
\label{op:YMA5-V6}
Preserve all preceding choices.  On the first retained axial mode, proceed
in the following order.
\begin{enumerate}[label=\textup{(V6.\arabic*)},leftmargin=5.2em]
\item Use the exact Bochner current $\omega_h=L_1^{-1}\mathscr B h$ and its
      gauge companions to acquire the complete positive Gram $\mathscr G$
      on transverse plus gauge directions.  Form
      $\mathscr G_{\perp\mid g}$ only after that common Gram is assembled.
\item Populate the finite small-field head of
      $\mathbb E\Sch_{\mathcal R}-\mathscr T_\Theta$ using
      \Cref{cor:YMA5-small-Schur-tail}; retain an outward interval for the
      omitted tail rather than an absolute all-volume norm.
\item Isolate the large-field current by one stopped or IMS-localized Witten
      cylinder.  Its return is charged once to the same score source.  A
      successful estimate must be relative to the physical reference or be
      used only on the high-mode side of
      \Cref{thm:YMA5-two-regime}.
\item Compare the resulting physical interval with the exact V4 coupling
      flow.  Evaluate $\eta_4(\pi/2)$ only when it is the first unresolved
      coefficient of that same interval.
\end{enumerate}
A negative eigenvector of the assembled sum in
\eqref{eq:YMA5-exact-current-short} is a physical pure-mode obstruction.  A
negative gauge-covariant floor without the current short, failure of an
absolute pointwise estimate below \eqref{eq:YMA5-126-window}, or an
unpopulated large-field return is \stopstatus, not an obstruction.
\end{openproblem}

% =============================================================================
\section{Version V5 proof ledger and exact replay boundary}
\label{sec:YMA5-ledger}
% =============================================================================

\begingroup\small
\begin{longtable}{>{\raggedright\arraybackslash}p{0.27\textwidth}
                  >{\raggedright\arraybackslash}p{0.20\textwidth}
                  >{\raggedright\arraybackslash}p{0.45\textwidth}}
\caption{Version V5 proof-status ledger.}\label{tab:YMA5-ledger}\\
\toprule
\textbf{Row}&\textbf{Status}&\textbf{Exact scope}\\
\midrule
\endfirsthead
\toprule
\textbf{Row}&\textbf{Status}&\textbf{Exact scope}\\
\midrule
\endhead
Wilson Hessian normalization&Corrected exactly
&The actual multiplier is $\beta/3$; $\beta/6$ is the Taylor quadratic
coefficient.\\[0.3em]
Parity-tree fibre gap&Proved and exactly replayed
&$\lambda_*$ is the displayed octic root; the cell lower bound is attained
at $k=(\pi,\pi,\pi,0)$.\\[0.3em]
Flat vertical Hessian&Strengthened exactly
&$A_M\succeq(\lambda_*/3)I$, improving the former $I/60$ flat bound.\\[0.3em]
All-momentum tree reference&Proved sharply
&$\mathsf S^{\rm tree}\succeq\frac14d_c^*d_c$; an exact Bloch completion
attains equality.\\[0.3em]
Gauge completion stock&Proved exactly
&Every coarse gauge component is copied to exactly seven reduced crossing
links.\\[0.3em]
Ricci pointwise window&Corrected
&A termwise absolute proof pays at most $7\kappa$ and has the method-level
window $126/((1-\delta)\beta)$, not $252/((1-\delta)\beta)$.\\[0.3em]
Gauge-covariant floor&Sharpened to an identity
&Physical Hessian equals the covariant floor plus the positive gauge-shorted
Witten-current residual.\\[0.3em]
Residual certificate&Corrected logically
&The condition retaining $-\mathscr G$ is equivalent; dropping the positive
term is sufficient only.\\[0.3em]
Flat propagator locality&Strengthened exactly
&Chebyshev decay uses $\rho_*<0.842073$ instead of the V3 Neumann ratio.\\[0.3em]
Small-field inverse locality&Proved uniformly
&The curvature-augmented vertical inverse and its pointwise Schur tail decay
exponentially, uniformly in cutoff.\\[0.3em]
Gauge-shorted finite head&Open physical value
&The common Witten Gram and expected finite Schur head have not been
populated on the selected law.\\[0.3em]
Large-field current return&Open physical estimate
&No Balaban/IMS large-field bound is inferred from the flat fibre gap.\\[0.3em]
Selected shell margin&\stopstatus
&No negative physical mode is obtained; the two rows in
\eqref{eq:YMA5-next-two-scalars} remain unacquired.\\[0.3em]
Higher jets, local-vacuum contraction, mass&Not inferred
&The inherited downstream handoff remains closed until the quadratic card
passes.\\
\bottomrule
\end{longtable}
\endgroup

The accompanying replay reconstructs both exact finite matrices from the
parity rule, verifies the characteristic-polynomial factorizations,
performs the rational Sturm counts and endpoint sign test, checks the sharp
reference completion on all $96$ plaquette rows, and checks the seven-copy
gauge count.  The Stokes lower bound, Chebyshev theorem, Witten-current
short, and small-field analytic estimates are proved in the text and are not
replaced by numerical tests.

\begin{assessment}[V5 termination]
\label{ass:YMA5-termination}
The exact flat fibre gap, sharp all-momentum reference floor, exact gauge
lift, normalization repair, current-short identity, and small-field inverse
locality are \pass\ at their stated scopes.  The report's useful structural
input has therefore been absorbed without importing its two normalization
and certificate overstatements.  The complete selected response-matched
symbol remains \stopstatus: the expected gauge-shorted finite head and the
large-field Witten return are still physical values.  Version~V6 should
acquire those two quantities on the first axial card before returning to the
finite two-loop number.
\end{assessment}

\section*{Version V5 append-only continuation record}
\addcontentsline{toc}{section}{Version V5 append-only continuation record}
The complete Version~V4 source above is preserved byte for byte up to its
sole final \verb|\end{document}| command.  Its SHA--256 digest is
\begin{center}\scriptsize\ttfamily
75caa09e37dacb584b0f027d2260ce138110657531b4189d73b504441abaf477.
\end{center}
For Version~V6, preserve this accumulated source, append new material
immediately before the final document-closing command, and increment the
filename to \texttt{\_V6.tex}.  Keep the exact block map, parity-tree
section, physical response, source normalization, first retained axial
mode, Haar compensation, and local/topological support convention fixed.

% =============================================================================
% END APPENDED VERSION V5 MATERIAL
% =============================================================================


% =============================================================================
% START APPENDED VERSION V6 MATERIAL
% =============================================================================
\clearpage
\hypersetup{
 pdftitle={Renewal Yang--Mills Reduced-Fibre Wilson Shell Symbol and Local-Vacuum Programme V6},
 pdfsubject={Return-renormalized Witten card, local curvature-defect channel, common gauge short, and comparison-current locality}
}
\pdfinfo{
 /Title (Renewal Yang--Mills Reduced-Fibre Wilson Shell Symbol and Local-Vacuum Programme V6)
 /Subject (Return-renormalized Witten card, local curvature-defect channel, common gauge short, and comparison-current locality)
}
\providecommand{\Short}{\operatorname{Short}}
\providecommand{\Pos}{\operatorname{Pos}}
\providecommand{\esssup}{\operatorname*{ess\,sup}}

\section{Version V6: return-renormalized Witten card and the local curvature-defect channel}
\label{sec:YMA6-direction}

Version~V5 is preserved above.  Its exact conclusion was that the flat
parity-tree algebra and the small-field pointwise inverse tail are no longer
the missing rows.  The first unresolved object is the physical axial card
after a common gauge short, with the large-field Witten return retained on
the same source.  The attached margin report organized this as an expected
pointwise Schur row, a positive Witten-current residual, and a negative
large-field current energy.  That organization is exact, but it contains one
further cancellable follower: part of the positive current residual is the
return generated by the same negative-curvature channel which appears in the
large-field term.

The present continuation removes that follower algebraically before any
estimate.  On the actual finite reduced fibre it proves the exact chain
\begin{equation}
 \boxed{
 \begin{gathered}
 \text{local Wilson curvature majorant}
 \longrightarrow
 \text{positive comparison Witten operator }\mathcal A_\eta
 \\
 \longrightarrow
 \text{comparison head }\mathbb H_\eta^{\rm cmp}
 \;-
 \text{source-visible curvature-defect return }\mathbb L_\eta
 \\
 \longrightarrow
 \text{one common transverse--gauge card}
 \longrightarrow
 \text{physical reduced-fibre symbol}.
 \end{gathered}}
 \label{eq:YMA6-programme-chain}
\end{equation}
The defect return has an exact Woodbury representation through a positive
compact channel.  Its source-cyclic moments give a monotone finite acquisition
and a rigorous tail once the channel edge is bounded.  The positive
comparison Witten inverse inherits cutoff-independent spatial decay from the
exact fibre gap of V5.  Thus both sides of the remaining card now have a
finite physical-radius head/tail architecture.  No value of the first
response-matched head or of the defect-channel edge is invented below.

\begin{firewall}[This is the actual shell card, not a third Yang--Mills programme]
\label{fire:YMA6-no-new-programme}
The operators introduced below are exact rearrangements of the same
reduced-fibre Hessian
\(
 \mathbb E K-\Cov(j)
\)
on the already frozen shell.  The comparison operator is not a new physical
dynamics, the curvature-defect channel is not a second source birth, and its
Neumann words are not an independent boundary programme.  Their sole purpose
is to acquire and bound the Wilson-specific finite head and return which V5
left open.
\end{firewall}

% =============================================================================
\section{The frozen Witten block and a Wilson-local curvature majorant}
\label{sec:YMA6-majorant}
% =============================================================================

Fix one finite regulator, the identity central coarse background, one
nonzero momentum--colour block, and first the exact local reduced law chosen
for the pure/local comparison card
\begin{equation}
 \nu(\dd z)=Z^{-1}e^{-F(z)}\,\dd m_{\rm red}(z)
 \label{eq:YMA6-law}
\end{equation}
on the parity-tree fibre
\(
 Y=\SU(3)^{\mathscr E^{\rm red}}
\).
Let
\(
 \mathscr H_1=L^2(\nu;T^*Y)
\)
be the one-form carrier, with the product Levi--Civita connection.  Let
\begin{equation}
 \mathscr B:E\longrightarrow\mathscr H_1
 \label{eq:YMA6-mixed-source}
\end{equation}
be the mixed vertical/coarse Hessian source on the frozen finite coarse bank
\(E\), and let
\(
 \mathbb D=\mathbb E_\nu K
\)
be its deterministic coarse Hessian.  The exact Witten card is
\begin{equation}
 \boxed{
 \mathbb H
 =\mathbb D-\mathscr B^*L_1^{-1}\mathscr B.}
 \label{eq:YMA6-exact-card}
\end{equation}
The inverse is taken on the exact one-form source carrier.  On the present
finite fibre it is an ordinary inverse: the weighted Witten one-form kernel
has dimension \(b_1(Y)=0\), since a finite product of \(\SU(3)\)'s is simply
connected.

For this exact local card the complete action is written as
\begin{equation}
 F^{\rm loc}=\sum_{p}F_p^{\rm W}+\sum_{a\in\mathcal A_{\Psi}}F_a^{\Psi}.
 \label{eq:YMA6-local-action-sum}
\end{equation}
The first sum is over the literal fine Wilson plaquettes.  The second is a
submitted finite-range generated head only when that head is explicitly part
of the present local cylinder.  A quasilocal generated tail and the inherited
topological packet are not silently included in this law; they retain the
independent V4 insertion cards.  The algebraic resolvent identities below
apply to a complete action after all of its Hessian terms are put into the
curvature split, but the explicit locality constants are asserted only for
this local branch.

For a local term \(\gamma\), let
\(
 \mathcal H_\gamma(z)
\)
be its vertical Levi--Civita Hessian, embedded in \(T_z^*Y\).  For a Wilson
plaquette let
\(
 \mathcal H_p^0=\mathcal H_p(e)
\)
be its flat value, with the physical coefficient \(\alpha_\beta=\beta/3\)
already fixed in V5.  For a self-adjoint operator \(A\), write
\(
 A_+=(|A|+A)/2
\)
and
\(
 A_-=(|A|-A)/2
\).

\begin{definition}[Local flat-majorant curvature split]
\label{def:YMA6-local-majorant}
Fix \(0<\eta<1\).  For every Wilson plaquette and generated local term put
\begin{align}
 \Theta_{p,\eta}(z)
 &:=\bigl((1-\eta)\mathcal H_p^0-\mathcal H_p(z)\bigr)_+,
 \label{eq:YMA6-Wilson-defect}\\
 \Theta_a^{\Psi}(z)
 &:=\bigl(-\mathcal H_a^{\Psi}(z)\bigr)_+.
 \label{eq:YMA6-generated-defect}
\end{align}
Define
\begin{align}
 \Theta_\eta
 &:=\sum_p\Theta_{p,\eta}+\sum_a\Theta_a^{\Psi},
 \label{eq:YMA6-total-defect}\\
 \mathcal R_\eta
 &:=\kappa I+\sum_p(\mathcal H_p+\Theta_{p,\eta})
       +\sum_a(\mathcal H_a^{\Psi}+\Theta_a^{\Psi}),
 \qquad \kappa=\frac32.
 \label{eq:YMA6-majorant-R}
\end{align}
All sums retain the literal local supports and multiplicities of the
submitted action.
\end{definition}

\begin{theorem}[Exact Wilson-local majorant and regulator-independent floor]
\label{thm:YMA6-local-majorant}
The split of \Cref{def:YMA6-local-majorant} is a measurable local curvature
split in the sense of the attached Bochner--Schur formula:
\begin{equation}
 \boxed{
 \kappa I+\Hess_zF^{\rm loc}
 =\mathcal R_\eta-\Theta_\eta,
 \qquad
 \mathcal R_\eta\succ0,
 \qquad
 \Theta_\eta\succeq0.}
 \label{eq:YMA6-exact-curvature-split}
\end{equation}
On the reduced vertical carrier,
\begin{equation}
 \boxed{
 \mathcal R_\eta(z)
 \succeq
 \kappa I+(1-\eta)\alpha_\beta\Sigma_{rr}
 \succeq
 \rho_{\beta,\eta}I,
 \qquad
 \rho_{\beta,\eta}
 :=\kappa+(1-\eta)\alpha_\beta\lambda_* .}
 \label{eq:YMA6-majorant-floor}
\end{equation}
The constants are independent of the torus side.  The correction is local:
each Wilson channel acts only on the four link slots of one plaquette, and a
generated channel has exactly the support of its declared local term.

For the pure Wilson head, the coarse-cell propagation of
\(\mathcal R_\eta\) is at most two and its off-diagonal absolute block-row
sum is at most
\begin{equation}
 \boxed{J_{\beta,\eta}^{\rm W}\le216\alpha_\beta.}
 \label{eq:YMA6-majorant-row}
\end{equation}
A submitted generated head adds its independently certified finite-range row
constant; it is not absorbed into the number \(216\).
\end{theorem}

\begin{proof}
For any self-adjoint \(A,B\),
\(
 (B-A)_+\succeq B-A
\), because
\(
 (B-A)_+-(B-A)=(B-A)_-\succeq0
\).
Consequently
\[
 \mathcal H_p+\bigl((1-\eta)\mathcal H_p^0-\mathcal H_p\bigr)_+
 \succeq(1-\eta)\mathcal H_p^0,
\]
and
\(
 \mathcal H_a^\Psi+(-\mathcal H_a^\Psi)_+
 = (\mathcal H_a^\Psi)_+\succeq0
\).
Summing gives both the exact difference
\eqref{eq:YMA6-exact-curvature-split} and the first inequality in
\eqref{eq:YMA6-majorant-floor}.  The flat Wilson vertical Hessian is
\(
 \alpha_\beta\Sigma_{rr}
\), and V5 proved
\(
 \Sigma_{rr}\succeq\lambda_*I
\), giving the second inequality.

For the row estimate, one plaquette second-derivative block has norm at most
\(\alpha_\beta\), as in V5.  Its four-by-four link-slot operator therefore
has norm at most \(4\alpha_\beta\).  Hence
\[
 \norm{(1-\eta)\mathcal H_p^0-\mathcal H_p}
 \le8\alpha_\beta,
 \qquad
 \norm{\Theta_{p,\eta}}\le8\alpha_\beta.
\]
Every block of the positive part has norm at most its operator norm.  Thus a
row inside one plaquette contributes at most
\(
 4(\alpha_\beta+8\alpha_\beta)=36\alpha_\beta
\).
A fine link belongs to six plaquettes, giving
\(
 6\cdot36\alpha_\beta=216\alpha_\beta
\).
The Ricci term is diagonal and does not enter the off-diagonal row.  The
support statement gives propagation two.
\end{proof}

\begin{remark}[What the defect channel means]
The local operator \(\Theta_{p,\eta}\) measures failure of the actual
plaquette Hessian to dominate the fixed fraction
\((1-\eta)\mathcal H_p^0\).  It is zero at the flat configuration and is
small under a small Hessian perturbation, but it is not defined by an
arbitrary pointwise field-amplitude threshold.  It therefore contains the
precise local curvature deficit which a large-field/IMS proof must control.
A later amplitude partition may split these channels further, but it may not
change their sum or count either part as a new source.
\end{remark}

% =============================================================================
\section{Return-renormalized Bochner--Schur identity}
\label{sec:YMA6-resolvent-split}
% =============================================================================

Let
\begin{equation}
 \boxed{
 \mathcal A_\eta
 :=\nabla_\nu^*\nabla+\mathcal R_\eta,
 \qquad
 L_1=\mathcal A_\eta-\Theta_\eta.}
 \label{eq:YMA6-A-and-L}
\end{equation}
The first operator is the positive comparison Witten operator.  It is not a
second law: both operators act in \(L^2(\nu;T^*Y)\), and the expectation law
\(\nu\) is unchanged.  Define three source Grams on the same coarse bank:
\begin{equation}
 \boxed{
 \begin{aligned}
 \mathbb C_{\mathcal R}
 &:=\mathscr B^*\mathcal R_\eta^{-1}\mathscr B,\\
 \mathbb C_{\mathcal A}
 &:=\mathscr B^*\mathcal A_\eta^{-1}\mathscr B,\\
 \mathbb C_{1}
 &:=\mathscr B^*L_1^{-1}\mathscr B.
 \end{aligned}}
 \label{eq:YMA6-three-source-Grams}
\end{equation}
The first inverse is pointwise in the field variable; the other two are
Witten resolvents.

\begin{theorem}[Three-Gram identity and the net curvature-defect return]
\label{thm:YMA6-three-Gram}
Put
\begin{align}
 \mathbb G_\eta^{\rm cmp}
 &:=\mathbb C_{\mathcal R}-\mathbb C_{\mathcal A},
 \label{eq:YMA6-Gcmp}\\
 \mathbb L_\eta^{\rm ret}
 &:=\mathbb C_{1}-\mathbb C_{\mathcal A},
 \label{eq:YMA6-Lret}\\
 \mathbb H_\eta^{\rm cmp}
 &:=\mathbb D-\mathbb C_{\mathcal A}.
 \label{eq:YMA6-Hcmp}
\end{align}
Then
\begin{equation}
 \boxed{
 \mathbb G_\eta^{\rm cmp}\succeq0,
 \qquad
 \mathbb L_\eta^{\rm ret}\succeq0,}
 \label{eq:YMA6-two-positive-forms}
\end{equation}
and the physical card has both exact decompositions
\begin{equation}
 \boxed{
 \begin{aligned}
 \mathbb H
 &=\mathbb E\Sch_{\mathcal R_\eta}
   +\mathbb G_\eta^{\rm cmp}
   -\mathbb L_\eta^{\rm ret},\\
 &=\mathbb H_\eta^{\rm cmp}-\mathbb L_\eta^{\rm ret}.
 \end{aligned}}
 \label{eq:YMA6-master-resolvent-decomposition}
\end{equation}
Here
\(
 \mathbb E\Sch_{\mathcal R_\eta}
 =\mathbb D-\mathbb C_{\mathcal R}
\)
is the expected pointwise curvature-augmented Schur form.

For \(h\in E\), let
\begin{equation}
 b_h=\mathscr Bh,
 \qquad
 \eta_h=\mathcal A_\eta^{-1}b_h.
 \label{eq:YMA6-comparison-current}
\end{equation}
Then the positive comparison residual is the literal Gram
\begin{equation}
 \boxed{
 \begin{aligned}
 \mathbb G_\eta^{\rm cmp}[h,k]
 ={}&\ip{\nabla\eta_h}{\nabla\eta_k}_{\nu}\\
 &+\ip{\mathcal R_\eta^{1/2}
  (\eta_h-\mathcal R_\eta^{-1}b_h)}
 {\mathcal R_\eta^{1/2}
  (\eta_k-\mathcal R_\eta^{-1}b_k)}_{\nu}.
 \end{aligned}}
 \label{eq:YMA6-Gcmp-Gram}
\end{equation}
Thus one common bank of comparison currents acquires its complete
transverse--gauge matrix before any short is taken.
\end{theorem}

\begin{proof}
The form order
\(
 \mathcal A_\eta=\mathcal R_\eta+\nabla_\nu^*\nabla
 \succeq\mathcal R_\eta
\)
gives
\(
 \mathcal R_\eta^{-1}\succeq\mathcal A_\eta^{-1}
\).
Likewise
\(
 \mathcal A_\eta=L_1+\Theta_\eta\succeq L_1
\)
gives
\(
 L_1^{-1}\succeq\mathcal A_\eta^{-1}
\).
Compression by \(\mathscr B\) proves
\eqref{eq:YMA6-two-positive-forms}.  Algebra then gives
\[
 \mathbb D-\mathbb C_1
 = (\mathbb D-\mathbb C_{\mathcal R})
   +(\mathbb C_{\mathcal R}-\mathbb C_{\mathcal A})
   -(\mathbb C_1-\mathbb C_{\mathcal A}),
\]
which is the first decomposition; the second follows by grouping its first
two terms.

For the Gram formula, consider the positive quadratic functional
\[
 \mathcal J_h(\zeta)
 =\norm{\nabla\zeta}_{\nu}^2
  +\norm{\mathcal R_\eta^{1/2}
    (\zeta-\mathcal R_\eta^{-1}b_h)}_{\nu}^2.
\]
After expansion,
\[
 \mathcal J_h(\zeta)
 =\ip{\zeta}{\mathcal A_\eta\zeta}_\nu
  -2\operatorname{Re}\ip{b_h}{\zeta}_\nu
  +\ip{b_h}{\mathcal R_\eta^{-1}b_h}_\nu.
\]
Its unique minimizer is \(\eta_h=\mathcal A_\eta^{-1}b_h\), and its
minimum is
\(
 \ip{b_h}{(\mathcal R_\eta^{-1}-\mathcal A_\eta^{-1})b_h}
\).
Polarization gives \eqref{eq:YMA6-Gcmp-Gram}.
\end{proof}

\begin{theorem}[Exact cancellation of the duplicated Witten follower]
\label{thm:YMA6-follower-cancellation}
Let
\(
 \omega_h=L_1^{-1}b_h
\)
be the exact Witten current and put
\(
 d_h=\omega_h-\eta_h
\).
Let \(\mathscr G\) and \(\mathscr T_\Theta\) be the two forms in the
Bochner--Schur decomposition imported in V5.  Define the positive follower
Gram
\begin{equation}
 \mathbb D_\eta^{\rm fol}[h,k]
 :=\ip{d_h}{\mathcal A_\eta d_k}_\nu
 =\ip{\Theta_\eta\omega_h}
        {\mathcal A_\eta^{-1}\Theta_\eta\omega_k}_\nu.
 \label{eq:YMA6-follower-Gram}
\end{equation}
Then
\begin{equation}
 \boxed{
 \mathscr G
 =\mathbb G_\eta^{\rm cmp}+\mathbb D_\eta^{\rm fol},
 \qquad
 \mathscr T_\Theta
 =\mathbb L_\eta^{\rm ret}+\mathbb D_\eta^{\rm fol}.}
 \label{eq:YMA6-follower-two-appearances}
\end{equation}
Consequently
\begin{equation}
 \boxed{
 \mathscr G-\mathscr T_\Theta
 =\mathbb G_\eta^{\rm cmp}-\mathbb L_\eta^{\rm ret}.}
 \label{eq:YMA6-follower-cancels}
\end{equation}
The raw large-field energy \(\mathscr T_\Theta\) is therefore not the net
physical debit.  Its comparison-return follower
\(\mathbb D_\eta^{\rm fol}\) must first be removed.
\end{theorem}

\begin{proof}
The Euler equations are
\[
 \mathcal A_\eta\eta_h=b_h,
 \qquad
 (\mathcal A_\eta-\Theta_\eta)\omega_h=b_h,
\]
so
\begin{equation}
 \mathcal A_\eta d_h=\Theta_\eta\omega_h.
 \label{eq:YMA6-d-equation}
\end{equation}
The quadratic functional \(\mathcal J_h\) used in the preceding proof has
minimizer \(\eta_h\).  Completion of the square gives
\[
 \mathcal J_h(\omega_h)
 =\mathcal J_h(\eta_h)
  +\norm{d_h}_{\mathcal A_\eta}^2.
\]
The left side is exactly the imported form \(\mathscr G[h]\), the first
term is \(\mathbb G_\eta^{\rm cmp}[h]\), and polarization gives the first
identity in \eqref{eq:YMA6-follower-two-appearances}.

For the second identity, use
\(
 \mathbb L_\eta^{\rm ret}[h,k]
 =\ip{b_h}{d_k}_\nu
\).
By \eqref{eq:YMA6-d-equation},
\[
 \ip{b_h}{d_k}
 =\ip{\eta_h}{\mathcal A_\eta d_k}
 =\ip{\eta_h}{\Theta_\eta\omega_k}.
\]
On the other hand,
\[
 \mathscr T_\Theta[h,k]-\mathbb D_\eta^{\rm fol}[h,k]
 =\ip{\omega_h-d_h}{\Theta_\eta\omega_k}
 =\ip{\eta_h}{\Theta_\eta\omega_k}.
\]
This proves the second identity and hence the cancellation.
\end{proof}

\begin{firewall}[Source-first accounting]
\label{fire:YMA6-follower-not-new-source}
The same coarse score source \(b_h\) generates \(\eta_h\), \(\omega_h\),
\(d_h\), and every form in
\eqref{eq:YMA6-follower-two-appearances}.  The follower
\(\mathbb D_\eta^{\rm fol}\) appears once on each side of the old
Bochner split and cancels.  It is not a second trajectory price and is not
retained as a new current source after
\eqref{eq:YMA6-follower-cancels}.
\end{firewall}

% =============================================================================
\section{The return-renormalized gauge-current short}
\label{sec:YMA6-gauge-short}
% =============================================================================

Define the signed return-renormalized floor form
\begin{equation}
 \overline Q_\eta
 :=\mathbb E\Sch_{\mathcal R_\eta}-\mathbb L_\eta^{\rm ret}.
 \label{eq:YMA6-Qbar}
\end{equation}
Then
\begin{equation}
 \mathbb H=\overline Q_\eta+\mathbb G_\eta^{\rm cmp}.
 \label{eq:YMA6-Qbar-plus-G}
\end{equation}
Let \(E=E_\perp\oplus E_g\) be the physical transverse/gauge decomposition
on the nonzero momentum block.

\begin{theorem}[Return-renormalized exact current short]
\label{thm:YMA6-return-current-short}
Define
\begin{equation}
 \overline{\mathfrak F}_{\eta}[h]
 :=\sup_{g\in E_g}\overline Q_\eta[h+g],
 \qquad h\in E_\perp.
 \label{eq:YMA6-return-floor}
\end{equation}
Let
\(
 G^{\rm cmp}_{\perp\perp},G^{\rm cmp}_{\perp g},
 G^{\rm cmp}_{g\perp},G^{\rm cmp}_{gg}
\)
be the blocks of \(\mathbb G_\eta^{\rm cmp}\).  Then
\begin{equation}
 \boxed{
 \begin{aligned}
 \mathbb G_{\eta,\perp\mid g}^{\rm cmp}
 &:={}
 G^{\rm cmp}_{\perp\perp}
 -G^{\rm cmp}_{\perp g}
  (G^{\rm cmp}_{gg})^\dagger
  G^{\rm cmp}_{g\perp}\succeq0,\\
 \mathbb H|_{E_\perp}
 &=\overline{\mathfrak F}_{\eta}
   +\mathbb G_{\eta,\perp\mid g}^{\rm cmp}.
 \end{aligned}}
 \label{eq:YMA6-return-short-identity}
\end{equation}
Equivalently,
\begin{equation}
 \mathbb G_{\eta,\perp\mid g}^{\rm cmp}[h]
 =\inf_{g\in E_g}\mathbb G_\eta^{\rm cmp}[h+g].
 \label{eq:YMA6-Gcmp-short-inf}
\end{equation}
Thus the two V5 acquisition rows may be replaced, without changing their
sum, by the easier pair
\begin{equation}
 \boxed{
 \overline{\mathfrak F}_{\eta}(q_{n_*})
 \quad\text{and}\quad
 \mathbb G_{\eta,\perp\mid g}^{\rm cmp}(q_{n_*}),}
 \label{eq:YMA6-new-two-rows}
\end{equation}
where the positive comparison-current Gram involves only
\(\mathcal A_\eta\), and the entire negative-curvature effect is already
contained once in \(\mathbb L_\eta^{\rm ret}\).
\end{theorem}

\begin{proof}
The proof is the V5 block completion applied to
\eqref{eq:YMA6-Qbar-plus-G}.  Since the physical Hessian annihilates
\(E_g\), its gauge row and column vanish.  Hence the gauge blocks of
\(\overline Q_\eta\) are the negatives of the corresponding blocks of the
positive Gram \(\mathbb G_\eta^{\rm cmp}\).  Positivity gives the
Moore--Penrose range condition.  For fixed \(h\in E_\perp\),
\[
 \overline Q_\eta[h+g]
 =\mathbb H[h]-\mathbb G_\eta^{\rm cmp}[h+g].
\]
Taking the supremum in \(g\), and evaluating the minimum of the positive
quadratic by the supported Schur complement, gives
\eqref{eq:YMA6-return-short-identity}--\eqref{eq:YMA6-Gcmp-short-inf}.
\end{proof}

\begin{remark}[Why this is sharper than bounding \(\mathscr T_\Theta\)]
The old covariant floor shorted
\(
 \mathbb E\Sch_{\mathcal R_\eta}-\mathscr T_\Theta
\)
and repaired it with the short of the full current Gram \(\mathscr G\).
The new identity moves the common follower
\(\mathbb D_\eta^{\rm fol}\) through both terms before shorting.  It asks
for the net return \(\mathbb L_\eta^{\rm ret}\), not the larger raw energy
\(\mathscr T_\Theta\), and for the comparison-current residual, not the
full current residual.  No inequality is used in this replacement.
\end{remark}

% =============================================================================
\section{The local curvature-defect channel and its source-cyclic return}
\label{sec:YMA6-defect-channel}
% =============================================================================

The local decomposition of \(\Theta_\eta\) gives a canonical channel
factorization.  Let
\begin{equation}
 \mathscr C_\eta
 =\bigoplus_p L^2(\nu;\Ran\Theta_{p,\eta})
  \oplus
  \bigoplus_a L^2(\nu;\Ran\Theta_a^\Psi)
 \label{eq:YMA6-channel-space}
\end{equation}
with the zero spectral fibres omitted, and define
\begin{equation}
 \mathcal V_\eta:\mathscr C_\eta\longrightarrow\mathscr H_1,
 \qquad
 \mathcal V_\eta(y_p,y_a)
 =\sum_p\Theta_{p,\eta}^{1/2}y_p
  +\sum_a(\Theta_a^\Psi)^{1/2}y_a.
 \label{eq:YMA6-channel-synthesis}
\end{equation}
The injections into the global one-form bundle are understood.  Then
\begin{equation}
 \boxed{\mathcal V_\eta\mathcal V_\eta^*=\Theta_\eta.}
 \label{eq:YMA6-VVstar}
\end{equation}
Define the defect-channel return and source entrance
\begin{equation}
 \boxed{
 \mathcal K_\eta
 :=\mathcal V_\eta^*\mathcal A_\eta^{-1}\mathcal V_\eta,
 \qquad
 \mathcal X_\eta
 :=\mathcal V_\eta^*\mathcal A_\eta^{-1}\mathscr B.}
 \label{eq:YMA6-K-and-X}
\end{equation}

\begin{theorem}[Exact Woodbury landing of the large-field return]
\label{thm:YMA6-Woodbury}
On every finite reduced-fibre card,
\begin{equation}
 \boxed{0\preceq\mathcal K_\eta\prec I}
 \label{eq:YMA6-K-contraction}
\end{equation}
on its supported channel carrier, and
\begin{equation}
 \boxed{
 L_1^{-1}
 =\mathcal A_\eta^{-1}
  +\mathcal A_\eta^{-1}\mathcal V_\eta
   (I-\mathcal K_\eta)^{-1}
   \mathcal V_\eta^*\mathcal A_\eta^{-1}.}
 \label{eq:YMA6-Woodbury-resolvent}
\end{equation}
Consequently the complete net curvature-defect debit is the positive
source-relative return
\begin{equation}
 \boxed{
 \mathbb L_\eta^{\rm ret}
 =\mathcal X_\eta^*(I-\mathcal K_\eta)^{-1}\mathcal X_\eta.}
 \label{eq:YMA6-return-channel}
\end{equation}
It depends only on the source-cyclic channel
\begin{equation}
 \mathscr C_{\eta,\mathcal X}
 :=\overline{\Span}\{
   \mathcal K_\eta^n\Ran\mathcal X_\eta:n\ge0\}.
 \label{eq:YMA6-source-cyclic-channel}
\end{equation}
A defect-channel direction orthogonal to this subspace carries no price in
the selected Wilson card.
\end{theorem}

\begin{proof}
Consider the positive block form
\begin{equation}
 \mathcal M_\eta
 =\begin{pmatrix}
   \mathcal A_\eta&\mathcal V_\eta\\
   \mathcal V_\eta^*&I
  \end{pmatrix}.
 \label{eq:YMA6-block-M}
\end{equation}
Its Schur complement over the channel block is
\(
 \mathcal A_\eta-\mathcal V_\eta\mathcal V_\eta^*=L_1
\), which is strictly positive on the one-form carrier.  Hence the other
Schur complement
\(
 I-\mathcal V_\eta^*\mathcal A_\eta^{-1}\mathcal V_\eta
 =I-\mathcal K_\eta
\)
is strictly positive.  The comparison resolvent is compact on the compact
finite-dimensional fibre manifold, and \(\mathcal V_\eta\) is bounded, so
\(\mathcal K_\eta\) is compact.  Strict positivity therefore gives an
operator-norm edge strictly below one on each finite card.

Block inversion of \(\mathcal M_\eta\), or direct multiplication, gives
\eqref{eq:YMA6-Woodbury-resolvent}.  Compressing its difference from
\(\mathcal A_\eta^{-1}\) by \(\mathscr B\) gives
\eqref{eq:YMA6-return-channel}.  Invariance of the cyclic subspace and the
Neumann series show that only the displayed source-cyclic closure is used.
\end{proof}

\begin{corollary}[Monotone return words and a finite depth certificate]
\label{cor:YMA6-return-words}
Let
\begin{equation}
 \mathbb M_{\eta,n}
 :=\mathcal X_\eta^*\mathcal K_\eta^n\mathcal X_\eta
 \succeq0,
 \qquad n\ge0.
 \label{eq:YMA6-return-moments}
\end{equation}
Then
\begin{equation}
 \boxed{
 \mathbb L_\eta^{\rm ret}
 =\sum_{n=0}^{\infty}\mathbb M_{\eta,n}}
 \label{eq:YMA6-Neumann-return}
\end{equation}
in operator norm.  If
\begin{equation}
 \rho_{\eta,\mathcal X}
 :=\norm{\mathcal K_\eta|_{\mathscr C_{\eta,\mathcal X}}}<1,
 \label{eq:YMA6-source-edge}
\end{equation}
then for every \(N\ge0\),
\begin{equation}
 \boxed{
 0\preceq
 \mathbb L_\eta^{\rm ret}-\sum_{n=0}^{N}\mathbb M_{\eta,n}
 \preceq
 \frac{\rho_{\eta,\mathcal X}^{N+1}}
      {1-\rho_{\eta,\mathcal X}}
 \mathbb M_{\eta,0}.}
 \label{eq:YMA6-return-depth-tail}
\end{equation}
Thus a finite bank of positive same-source return moments and one upper
interval for the source-cyclic edge gives an outward interval for the entire
large-field debit.
\end{corollary}

\begin{proof}
Use the norm-convergent Neumann series on the source-cyclic carrier.  Since
\(
 0\preceq\mathcal K_\eta\preceq
 \rho_{\eta,\mathcal X}I
\)
there,
\[
 \sum_{n>N}\mathcal K_\eta^n
 \preceq
 \frac{\rho_{\eta,\mathcal X}^{N+1}}
      {1-\rho_{\eta,\mathcal X}}I.
\]
Compression by \(\mathcal X_\eta\) proves the matrix inequality.
\end{proof}

\begin{corollary}[A source-visible near-unit survivor]
\label{cor:YMA6-near-unit-survivor}
Consider a cofinal family and normalized coarse vectors \(h_r\) for which
\(
 \mathbb M_{\eta,0,r}[h_r]>0
\).
If
\begin{equation}
 \frac{\mathbb L_{\eta,r}^{\rm ret}[h_r]}
      {\mathbb M_{\eta,0,r}[h_r]}
 \longrightarrow\infty,
 \label{eq:YMA6-return-ratio-diverges}
\end{equation}
then
\(
 \rho_{\eta,\mathcal X,r}\to1
\).
Consequently there are unit vectors
\(
 y_r\in\mathscr C_{\eta,\mathcal X,r}
\)
with
\begin{equation}
 \ip{y_r}{\mathcal K_{\eta,r}y_r}\longrightarrow1.
 \label{eq:YMA6-channel-survivor}
\end{equation}
This is a typed source-visible curvature-defect survivor.  It becomes a
physical negative shell mode only if the assembled difference
\eqref{eq:YMA6-master-resolvent-decomposition} fails on the same coarse
vector.
\end{corollary}

\begin{proof}
The tail estimate with \(N=-1\), interpreted as the full geometric-series
bound, gives
\[
 \mathbb L_{\eta,r}^{\rm ret}[h_r]
 \le\frac{1}{1-\rho_{\eta,\mathcal X,r}}
      \mathbb M_{\eta,0,r}[h_r].
\]
The assumed ratio forces \(\rho_{\eta,\mathcal X,r}\to1\).  The variational
characterization of the norm of the positive self-adjoint contraction gives
the unit vectors.  The final qualification follows because the return is
only one signed component of the physical Hessian.
\end{proof}

\begin{theorem}[Exact head--tail Feshbach split of the defect channel]
\label{thm:YMA6-channel-Feshbach}
Let
\(
 \mathscr C_{\eta}=\mathscr C_H\oplus\mathscr C_T
\)
be any orthogonal partition of the local curvature-defect channels, for
example by a physical cell radius.  Write
\begin{equation}
 I-\mathcal K_\eta
 =\begin{pmatrix}A_H&B_{HT}\\B_{TH}&D_T\end{pmatrix},
 \qquad
 \mathcal X_\eta=\binom{X_H}{X_T}.
 \label{eq:YMA6-channel-blocks}
\end{equation}
Then \(D_T\succ0\), the supported head short
\begin{equation}
 S_H=A_H-B_{HT}D_T^{-1}B_{TH}
 \label{eq:YMA6-channel-head-short}
\end{equation}
is positive, and
\begin{equation}
 \boxed{
 \begin{aligned}
 \mathbb L_\eta^{\rm ret}
 ={}&X_T^*D_T^{-1}X_T\\
 &+\bigl(X_H-B_{HT}D_T^{-1}X_T\bigr)^*
   S_H^{-1}
   \bigl(X_H-B_{HT}D_T^{-1}X_T\bigr).
 \end{aligned}}
 \label{eq:YMA6-channel-Feshbach-return}
\end{equation}
The tail therefore contributes both a direct positive return and a
renormalization of the head entrance and head metric.  Shorting the head
before the mixed block is assembled loses exactly these terms.
\end{theorem}

\begin{proof}
The operator \(I-\mathcal K_\eta\) is strictly positive by
\Cref{thm:YMA6-Woodbury}; hence its principal block \(D_T\) and its Schur
complement \(S_H\) are positive.  Complete the square in the inverse block
quadratic form, or use the standard block-inverse formula.  With the signs in
\eqref{eq:YMA6-channel-blocks}, the effective head source is
\(
 X_H-B_{HT}D_T^{-1}X_T
\), giving the displayed identity.
\end{proof}

% =============================================================================
\section{Cutoff-independent locality of the positive comparison Witten inverse}
\label{sec:YMA6-comparison-locality}
% =============================================================================

V5 proved spatial decay of the pointwise inverse
\(\mathcal R_\eta(z)^{-1}\) on a small-field window.  The comparison card in
V6 requires the full elliptic inverse
\(\mathcal A_\eta^{-1}\).  The rough configuration-space Laplacian does not
spoil cell locality because the product connection acts diagonally in the
link-cell label.

Let
\begin{equation}
 \mathscr H_1=\bigoplus_x\mathscr H_{1,x}
 \label{eq:YMA6-cell-oneform-split}
\end{equation}
be the decomposition by coarse-cell base of a reduced link, and let \(P_x\)
be the corresponding orthogonal projection.  Distances are periodic
coarse-cell graph distances.

\begin{theorem}[Weighted-conjugation bound for a comparison Witten operator]
\label{thm:YMA6-Witten-Combes-Thomas}
Let
\begin{equation}
 \mathcal A=\nabla_\nu^*\nabla+\mathcal R
 \succeq\rho I,
 \qquad \rho>0,
 \label{eq:YMA6-general-A-floor}
\end{equation}
where \(\mathcal R(z)\) has cell propagation at most \(r_0\).  Suppose its
off-diagonal block kernel obeys the uniform Schur row and column bound
\begin{equation}
 \max\left\{
 \sup_x\sum_{y\ne x}\norm{\mathcal R_{xy}(z)},
 \sup_y\sum_{x\ne y}\norm{\mathcal R_{xy}(z)}
 \right\}\le J
 \label{eq:YMA6-R-row-J}
\end{equation}
for \(\nu\)-almost every \(z\).  If
\begin{equation}
 \varepsilon_\mu:=J(e^{\mu r_0}-1)<\rho,
 \label{eq:YMA6-CT-condition}
\end{equation}
then
\begin{equation}
 \boxed{
 \norm{P_x\mathcal A^{-1}P_y}_{\op}
 \le\frac{e^{-\mu\dist(x,y)}}{\rho-\varepsilon_\mu}.}
 \label{eq:YMA6-A-inverse-decay}
\end{equation}
The norm is the operator norm between the two \(L^2\) one-form cell
subspaces; no pointwise field kernel is assumed.
\end{theorem}

\begin{proof}
Fix \(y\) and let \(W_y\) multiply the cell component at \(x\) by
\(
 e^{\mu\dist(x,y)}
\).  On a finite torus this is a bounded invertible operator.  Because the
weight is independent of the field variable and scalar on each product-link
tangent, it commutes with the product covariant derivative and its weighted
adjoint.  Therefore
\[
 W_y\mathcal A W_y^{-1}-\mathcal A
 =W_y\mathcal R W_y^{-1}-\mathcal R=:\mathcal E_y.
\]
For a nonzero block \(\mathcal R_{xz}\), the distance difference has
absolute value at most \(r_0\).  The row/column Schur test gives
\(
 \norm{\mathcal E_y}\le J(e^{\mu r_0}-1)=\varepsilon_\mu
\).
Since \(\norm{\mathcal A^{-1}}\le\rho^{-1}\), the resolvent identity or a
Neumann series gives
\[
 \norm{(W_y\mathcal A W_y^{-1})^{-1}}
 \le(\rho-\varepsilon_\mu)^{-1}.
\]
Finally
\[
 P_x\mathcal A^{-1}P_y
 =e^{-\mu\dist(x,y)}
  P_x(W_y\mathcal A W_y^{-1})^{-1}P_y,
\]
which proves the estimate.
\end{proof}

\begin{corollary}[An explicit all-volume comparison-current decay rate]
\label{cor:YMA6-explicit-comparison-decay}
For the pure Wilson local-majorant split of
\Cref{thm:YMA6-local-majorant}, put
\begin{equation}
 \boxed{
 \mu_\eta
 :=\frac12\log\left(
  1+\frac{(1-\eta)\lambda_*}{432}\right)>0.}
 \label{eq:YMA6-mu-eta}
\end{equation}
Then, for every finite torus and every \(\beta>0\),
\begin{equation}
 \boxed{
 \norm{P_x\mathcal A_\eta^{-1}P_y}_{\op}
 \le\frac{2}{\rho_{\beta,\eta}}
 e^{-\mu_\eta\dist(x,y)}.}
 \label{eq:YMA6-explicit-A-decay}
\end{equation}
For \(\eta=1/2\), one may use
\(
 \mu_{1/2}>6.8\times10^{-5}
\).

Let \(\mathscr B_\beta\) be the pure Wilson mixed source.  It has propagation
at most two and absolute block-column sum at most
\(24\alpha_\beta\).  Hence, for coarse source cells \(0,x\) with
\(\ellone{x}>4\),
\begin{equation}
 \boxed{
 \norm{\mathbb C_{\mathcal A}(0,x)}_{\op}
 \le
 \frac{1152\alpha_\beta^2}{\rho_{\beta,\eta}}
 e^{-\mu_\eta(\ellone{x}-4)}.}
 \label{eq:YMA6-CA-tail}
\end{equation}
The comparison Hessian
\(
 \mathbb H_\eta^{\rm cmp}=\mathbb D-\mathbb C_{\mathcal A}
\)
therefore has a finite displacement head and the explicit second-moment tail
\begin{equation}
 \boxed{
 \frac12\sum_{\ellone{x}>R}x_1^2
 \norm{\mathbb H_\eta^{\rm cmp}(0,x)}_{\op}
 \le
 \frac{9216\alpha_\beta^2}{\rho_{\beta,\eta}}
 \sum_{n>R}n^2(n+1)^3e^{-\mu_\eta(n-4)}}
 \label{eq:YMA6-Hcmp-second-tail}
\end{equation}
for \(R>4\).  This is the full positive-comparison Witten head, not merely
the pointwise Schur inverse of V5.
\end{corollary}

\begin{proof}
Use \(r_0=2\), \(J=216\alpha_\beta\), and
\(\rho=\rho_{\beta,\eta}\) in
\Cref{thm:YMA6-Witten-Combes-Thomas}.  The choice
\eqref{eq:YMA6-mu-eta} gives
\[
 J(e^{2\mu_\eta}-1)
 \le216\alpha_\beta
   \frac{(1-\eta)\lambda_*}{432}
 \le\frac12\rho_{\beta,\eta}.
\]
This proves \eqref{eq:YMA6-explicit-A-decay}.  The numerical lower bound at
\(\eta=1/2\) follows from the certified lower endpoint for \(\lambda_*\).

For \eqref{eq:YMA6-CA-tail}, expand
\(
 \mathscr B_\beta^*\mathcal A_\eta^{-1}\mathscr B_\beta
\)
between source cells.  The two mixed-source supports move the endpoints by
at most four cells in total.  Multiply the inverse estimate by
\(
 2(24\alpha_\beta)^2=1152\alpha_\beta^2
\).
The deterministic block has finite range, so it does not contribute beyond
\(R>4\).  The number of points at \(\ell^1\)-distance \(n\) is at most
\(16(n+1)^3\); multiplying by \(x_1^2\le n^2\) and by \(1/2\) gives
\eqref{eq:YMA6-Hcmp-second-tail}.
\end{proof}

\begin{firewall}[The explicit decay rate is not a margin]
\label{fire:YMA6-locality-not-margin}
The comparison inverse and comparison head now have regulator-independent
spatial tails.  Their finite head still contains the actual Wilson
expectation, the comparison-current response, and any submitted generated
head.  The return channel
\(\mathbb L_\eta^{\rm ret}\) must also be inserted.  Exponential locality
of \(\mathcal A_\eta^{-1}\) alone does not determine the sign of the
assembled shell symbol.
\end{firewall}

% =============================================================================
\section{Weighted physical-radius control of the curvature-defect return}
\label{sec:YMA6-weighted-return}
% =============================================================================

Index the local defect channels by their coarse-cell anchors and write
\(
 \mathscr C_\eta=\bigoplus_x\mathscr C_{\eta,x}
\).
For a block kernel \(T=(T_{xy})\), define the symmetric weighted Schur norm
\begin{equation}
 \norm{T}_{\mu,\mathrm{Sch}}
 :=\max\left\{
 \sup_x\sum_y e^{\mu\dist(x,y)}\norm{T_{xy}},
 \sup_y\sum_x e^{\mu\dist(x,y)}\norm{T_{xy}}
 \right\}.
 \label{eq:YMA6-weighted-Schur}
\end{equation}
The same definition is used for rectangular channel/source kernels.  The
triangle inequality for the distance makes this a submultiplicative kernel
norm.

\begin{theorem}[Weighted defect-channel contraction gives a finite physical head]
\label{thm:YMA6-weighted-channel}
Fix \(0<\mu<\mu_\eta\).  Suppose the actual local defect channel satisfies
\begin{equation}
 \boxed{
 \kappa_{\eta,\mu}
 :=\norm{\mathcal K_\eta}_{\mu,\mathrm{Sch}}<1,
 \qquad
 X_{\eta,\mu}
 :=\norm{\mathcal X_\eta}_{\mu,\mathrm{Sch}}<\infty.}
 \label{eq:YMA6-weighted-channel-condition}
\end{equation}
Then
\begin{equation}
 \boxed{
 \norm{(I-\mathcal K_\eta)^{-1}}_{\mu,\mathrm{Sch}}
 \le\frac{1}{1-\kappa_{\eta,\mu}},
 \qquad
 \norm{\mathbb L_\eta^{\rm ret}}_{\mu,\mathrm{Sch}}
 \le\frac{X_{\eta,\mu}^2}{1-\kappa_{\eta,\mu}}.}
 \label{eq:YMA6-weighted-return-bound}
\end{equation}
Moreover, the return-depth truncation obeys
\begin{equation}
 \boxed{
 \left\|
 \mathbb L_\eta^{\rm ret}
 -\sum_{n=0}^{N}\mathbb M_{\eta,n}
 \right\|_{\mu,\mathrm{Sch}}
 \le
 \frac{X_{\eta,\mu}^2\kappa_{\eta,\mu}^{N+1}}
      {1-\kappa_{\eta,\mu}}.}
 \label{eq:YMA6-weighted-depth-tail}
\end{equation}
In particular, its omitted displacement second moment satisfies
\begin{equation}
 \boxed{
 \frac12\sum_{\ellone{x}>R}x_1^2
 \norm{\mathbb L_\eta^{\rm ret}(0,x)}_{\op}
 \le
 \frac{8X_{\eta,\mu}^2}{1-\kappa_{\eta,\mu}}
 \sum_{n>R}n^2(n+1)^3e^{-\mu n}.}
 \label{eq:YMA6-return-spatial-tail}
\end{equation}
Thus the single weighted contraction
\eqref{eq:YMA6-weighted-channel-condition} is a sufficient Row~B
certificate for the physical-radius tail of the net return.
\end{theorem}

\begin{proof}
For block kernels, submultiplicativity follows from
\(
 e^{\mu\dist(x,z)}
 \le e^{\mu\dist(x,y)}e^{\mu\dist(y,z)}
\)
and the row/column Schur estimates.  Therefore the Neumann series gives
\[
 \norm{(I-\mathcal K_\eta)^{-1}}_{\mu,\mathrm{Sch}}
 \le\sum_{n\ge0}\kappa_{\eta,\mu}^n
 =\frac1{1-\kappa_{\eta,\mu}}.
\]
Multiplication on both sides by the rectangular kernel
\(\mathcal X_\eta\) proves the return bound.  Starting the series at
\(N+1\) proves \eqref{eq:YMA6-weighted-depth-tail}.  The weighted norm gives
\(
 \norm{\mathbb L_\eta^{\rm ret}(0,x)}
 \le X_{\eta,\mu}^2(1-\kappa_{\eta,\mu})^{-1}
 e^{-\mu\ellone{x}}
\).
Count lattice points as in V5 and the preceding corollary to obtain
\eqref{eq:YMA6-return-spatial-tail}.
\end{proof}

\begin{assessment}[What remains physical in the weighted condition]
\label{ass:YMA6-weighted-physical}
The comparison inverse makes every block of \(\mathcal K_\eta\) and
\(\mathcal X_\eta\) exponentially local.  That fact alone does not force
the weighted norm of \(\mathcal K_\eta\) below one.  The strict unweighted
edge is automatic on each finite card, but a uniform weighted edge is the
actual no-near-unit curvature-defect return estimate.  Its failure may be
caused by a genuine source-visible long return, or merely by the crudeness of
the weighted Schur norm.  Only a near-unit source-cyclic spectral witness or
a failed assembled Hessian is a typed mathematical obstruction.
\end{assessment}

% =============================================================================
\section{The literal axial four-direction card}
\label{sec:YMA6-axial-card}
% =============================================================================

Now freeze the first retained axial momentum
\(
 q=(q_1,0,0,0)
e0
\)
and one predeclared unit colour vector.  Before the Hodge short, its coarse
direction bank is
\begin{equation}
 E_q\cong\C^4,
 \qquad
 E_{g,q}=\Span\{d(q)\},
 \qquad
 d(q)=(e^{iq_1}-1,0,0,0),
 \label{eq:YMA6-axial-direction-bank}
\end{equation}
so
\(
 E_{\perp,q}=\Span\{e_2,e_3,e_4\}
\).
The global colour conjugation symmetry commutes with the local spectral
functional calculus in \Cref{def:YMA6-local-majorant}.  Hence every block
below is a four-by-four direction matrix tensored with the adjoint colour
identity.  No separate direction scalarization is assumed.

Let
\begin{equation}
 P_{\eta,q}:=\widehat{\mathbb H_\eta^{\rm cmp}}(q),
 \qquad
 R_{\eta,q}:=\widehat{\mathbb L_\eta^{\rm ret}}(q),
 \qquad
 H_q=P_{\eta,q}-R_{\eta,q}.
 \label{eq:YMA6-P-R-H}
\end{equation}
The return matrix is positive semidefinite.  Write their blocks with respect
to \(E_{\perp,q}\oplus E_{g,q}\).

\begin{theorem}[Common-gauge supported reduction of head minus return]
\label{thm:YMA6-common-gauge-card}
The exact Ward identity \(H_qE_{g,q}=0\) implies
\begin{equation}
 \boxed{
 P_{\eta,gg}=R_{\eta,gg},
 \qquad
 P_{\eta,\perp g}=R_{\eta,\perp g},
 \qquad
 P_{\eta,g\perp}=R_{\eta,g\perp}.}
 \label{eq:YMA6-common-gauge-blocks}
\end{equation}
Positivity of \(R_{\eta,q}\) gives the support incidence
\(
 \Ran R_{\eta,g\perp}\subseteq\Ran R_{\eta,gg}
\).
Define the two common-gauge reductions
\begin{align}
 P_{\eta,\perp\mid g}
 &:={}
 P_{\eta,\perp\perp}
 -P_{\eta,\perp g}R_{\eta,gg}^{\dagger}
  P_{\eta,g\perp},
 \label{eq:YMA6-P-short}\\
 R_{\eta,\perp\mid g}
 &:={}
 R_{\eta,\perp\perp}
 -R_{\eta,\perp g}R_{\eta,gg}^{\dagger}
  R_{\eta,g\perp}\succeq0.
 \label{eq:YMA6-R-short}
\end{align}
Then
\begin{equation}
 \boxed{
 H_q|_{E_{\perp,q}}
 =P_{\eta,\perp\mid g}-R_{\eta,\perp\mid g}.}
 \label{eq:YMA6-physical-short-difference}
\end{equation}
If \(R_{\eta,gg}=0\), positivity forces the mixed return blocks to vanish,
and the formula is read with zero pseudoinverse.

The axial little group makes only the assembled left side a scalar multiple
of \(I_3\).  Neither \(P_{\eta,\perp\mid g}\) nor
\(R_{\eta,\perp\mid g}\) is separately replaced by its group average.
\end{theorem}

\begin{proof}
Since \(H_q=P_{\eta,q}-R_{\eta,q}\) and the physical Hessian has zero gauge
row and column, the three block equalities are immediate.  If
\(z\in\Ker R_{\eta,gg}\), positivity of
\(
 R_{\eta,q}[(th,z)]
\)
for all real \(t\) forces
\(
 z\perp\Ran R_{\eta,g\perp}
\), which is the support incidence.  Subtract
\eqref{eq:YMA6-R-short} from \eqref{eq:YMA6-P-short}; the supported correction
terms are identical by \eqref{eq:YMA6-common-gauge-blocks} and cancel,
leaving
\(
 P_{\eta,\perp\perp}-R_{\eta,\perp\perp}=H_{q,\perp\perp}
\).
The zero-block statement is the same positivity argument.  The final
symmetry qualification is the raw-jet correction already established in
V2: section-dependent components need not inherit the separate physical
little-group scalarization.
\end{proof}

\begin{corollary}[Finite margin and direct obstruction certificate]
\label{cor:YMA6-finite-margin-card}
Let
\(
 S_q=\widehat S_{\beta,M}^{\rm flat}(q)|_{E_{\perp,q}}
\succ0
\)
and fix \(\delta\ge0\).  A rigorous interval enclosure for the complete
matrices in \eqref{eq:YMA6-P-R-H} gives the following terminating outcomes.
\begin{enumerate}[label=\textup{(C\arabic*)},leftmargin=4em]
\item If the lower endpoint of
      \begin{equation}
       S_q^{-1/2}
       (P_{\eta,\perp\mid g}-R_{\eta,\perp\mid g})
       S_q^{-1/2}-\delta I
       \label{eq:YMA6-direct-normalized-card}
      \end{equation}
      is positive, the selected card is \pass.
\item If an explicit transverse vector \(v\) has a negative upper endpoint
      for
      \begin{equation}
       v^*(H_q-\delta S_q)v,
       \label{eq:YMA6-direct-negative-witness}
      \end{equation}
      it is a direct physical \obstruction.
\item If the comparison head, source-cyclic return tail, Ward matching, or
      Moore--Penrose support is unresolved, the verdict is \stopstatus.
\end{enumerate}
A convenient sufficient, but not necessary, pass condition is
\begin{equation}
 \boxed{
 \lambda_{\min}
  (S_q^{-1/2}P_{\eta,\perp\mid g}S_q^{-1/2})
 -\lambda_{\max}
  (S_q^{-1/2}R_{\eta,\perp\mid g}S_q^{-1/2})
 \ge\delta.}
 \label{eq:YMA6-separated-sufficient-margin}
\end{equation}
Failure of this separated sufficient bound is not an obstruction; the two
matrices may have favorable eigenvector mismatch in their assembled
difference.
\end{corollary}

\begin{proof}
The first two outcomes are the ordinary spectral alternatives for the exact
identity \eqref{eq:YMA6-physical-short-difference}, with outward interval
arithmetic.  The support condition is part of the Moore--Penrose Schur
identity and cannot be guessed.  The sufficient condition follows from
\(
 \lambda_{\min}(A-B)
 \ge\lambda_{\min}(A)-\lambda_{\max}(B)
\).
The inequality need not be sharp, proving the final qualification.
\end{proof}

\begin{request}[Literal first-card acquisition packet]
\label{req:YMA6-first-card}
On the frozen first axial block, the smallest complete packet is now:
\begin{enumerate}[label=\textup{(A\arabic*)},leftmargin=4em]
\item the four-by-four deterministic mean \(\widehat{\mathbb D}(q)\);
\item the four-by-four comparison source Gram
      \(\widehat{\mathbb C_{\mathcal A}}(q)\), obtained from the positive
      comparison Witten solve and its finite displacement head;
\item the positive return moments
      \(\mathbb M_{\eta,0},\ldots,\mathbb M_{\eta,N}\) on the same four
      sources, together with either a source-cyclic edge interval or the
      weighted channel certificate of
      \Cref{thm:YMA6-weighted-channel};
\item the exact Ward residual
      \((P_{\eta,q}-R_{\eta,q})d(q)\), which must vanish before any supported
      gauge inverse is used;
\item the common-gauge reductions
      \eqref{eq:YMA6-P-short}--\eqref{eq:YMA6-R-short} and the direct
      comparison with the already evaluated reference \(s_0(q)I_3\).
\end{enumerate}
All entries are acquired on one law and one source bank.  A separately
sampled pure-law covariance, a pointwise curvature majorant from another
split, or a boundary channel from YM-B is not substituted into this card.
\end{request}

% =============================================================================
\section{V6 theorem, proof ledger, and exact next attack}
\label{sec:YMA6-terminal}
% =============================================================================

\begin{theorem}[Version V6 return-renormalized shell theorem]
\label{thm:YMA6-terminal}
Versions~V1--V5 are retained.  Version~V6 additionally proves:
\begin{enumerate}[label=\textup{(V6.\arabic*)},leftmargin=5.2em]
\item a local flat-majorant curvature split exists on the literal Wilson
      parity-tree fibre and has the regulator-independent floor
      \(\kappa+(1-\eta)\alpha_\beta\lambda_*\);
\item the physical Witten card is exactly a positive-comparison head minus
      one positive curvature-defect return;
\item the follower generated by comparing the exact and positive Witten
      currents occurs once in the old positive residual and once in the raw
      large-field energy, and cancels exactly;
\item after that cancellation, the physical transverse Hessian is the
      return-renormalized covariant floor plus the gauge-shorted positive
      comparison-current Gram;
\item the net curvature-defect return factors through the compact positive
      channel \(\mathcal K_\eta\) and one source entrance
      \(\mathcal X_\eta\);
\item positive source-cyclic return moments increase monotonically to the
      exact return and have a finite geometric tail once the source-cyclic
      edge is bounded;
\item the positive comparison Witten inverse has a cutoff-independent
      cell-decay estimate, and its complete comparison Hessian has an
      explicit finite second-displacement head/tail;
\item a weighted defect-channel contraction gives the missing physical-radius
      tail of the net return;
\item the literal nonzero axial acquisition is a four-by-four direction card
      per colour; head and return have the same gauge row and must be reduced
      with one common supported pseudoinverse before their difference is
      scalarized;
\item the selected card has the exact \pass--\obstruction--\stopstatus\ alternative of \Cref{cor:YMA6-finite-margin-card}.
\end{enumerate}
No item populates the first response-matched comparison head, the
source-cyclic channel edge, the submitted generated/topological corrections,
weak-facing-root escape, higher physical jets, local-vacuum contraction, or
a continuum Yang--Mills mass.
\end{theorem}

\begin{proof}
Items~\textup{(V6.1)}--\textup{(V6.4)} are
\Cref{thm:YMA6-local-majorant,thm:YMA6-three-Gram,thm:YMA6-follower-cancellation,thm:YMA6-return-current-short}.
Items~\textup{(V6.5)}--\textup{(V6.6)} are
\Cref{thm:YMA6-Woodbury,cor:YMA6-return-words}.  Items~\textup{(V6.7)}--
\textup{(V6.8)} are
\Cref{thm:YMA6-Witten-Combes-Thomas,cor:YMA6-explicit-comparison-decay,thm:YMA6-weighted-channel}.
Items~\textup{(V6.9)}--\textup{(V6.10)} are
\Cref{thm:YMA6-common-gauge-card,cor:YMA6-finite-margin-card}.
The limitations are the unpopulated entries in
\Cref{req:YMA6-first-card} and the inherited programme firewalls.
\end{proof}

\begingroup\small
\begin{longtable}{>{\raggedright\arraybackslash}p{0.27\textwidth}
                  >{\raggedright\arraybackslash}p{0.20\textwidth}
                  >{\raggedright\arraybackslash}p{0.45\textwidth}}
\caption{Version V6 proof-status ledger.}\label{tab:YMA6-ledger}\\
\toprule
\textbf{Row}&\textbf{Status}&\textbf{Exact scope}\\
\midrule
\endfirsthead
\toprule
\textbf{Row}&\textbf{Status}&\textbf{Exact scope}\\
\midrule
\endhead
Local Wilson curvature majorant&Proved exactly
&The local spectral deficit gives an exact curvature split with floor
$\kappa+(1-\eta)\alpha_\beta\lambda_*$ and finite propagation.\\[0.3em]
Comparison/return resolvent split&Proved exactly
&The physical card is $\mathbb H_\eta^{\rm cmp}-\mathbb L_\eta^{\rm ret}$,
with both covariance differences positive.\\[0.3em]
Duplicated Witten follower&Removed exactly
&The same positive follower occurs in $\mathscr G$ and
$\mathscr T_\Theta$ and cancels before estimation.\\[0.3em]
Return-renormalized gauge short&Proved exactly
&The physical Hessian is the new covariant floor plus the shorted comparison-current Gram.\\[0.3em]
Curvature-defect channel&Proved exactly
&$\mathbb L^{\rm ret}=\mathcal X^*(I-\mathcal K)^{-1}\mathcal X$ on the source-cyclic channel.\\[0.3em]
Return-depth tail&Finite exact certificate
&Positive moments plus a source-cyclic edge interval give the complete return interval.\\[0.3em]
Comparison Witten locality&Proved cutoff independent
&The local majorant and exact fibre gap give exponential decay of
$\mathcal A_\eta^{-1}$ and the comparison-head tail.\\[0.3em]
Physical-radius return tail&Conditional quantitative theorem
&A weighted channel norm below one gives the explicit spatial and return-depth tails.\\[0.3em]
Literal axial carrier&Reduced exactly
&One colour requires the complete four-direction head/return card; only their assembled transverse difference scalarizes.\\[0.3em]
First axial comparison head&Open physical value
&No expectation or Witten solve on the selected response-matched law is invented.\\[0.3em]
Source-cyclic defect edge&Open physical value
&Finite-card strictness is automatic; a useful outward or cofinal edge has not been acquired.\\[0.3em]
Generated/topological rows&Open physical values
&The submitted local generated head may enter the split, but its quasilocal tail and topological packet remain separate.\\[0.3em]
Selected shell margin&\stopstatus
&No negative assembled mode is obtained and the literal packet of
\Cref{req:YMA6-first-card} is not yet populated.\\[0.3em]
Higher jets, local-vacuum contraction, mass&Not inferred
&The downstream handoff remains closed until the physical quadratic card passes.\\
\bottomrule
\end{longtable}
\endgroup

The accompanying exact-rational replay checks the finite algebra of the
three-Gram decomposition, the duplicated-follower cancellation, the
Woodbury return channel, the positive return-word tail, the channel
head--tail Feshbach identity, and the common-gauge short.  It uses a separate
rational positive matrix model satisfying the theorem hypotheses.  It does
not evaluate a Wilson expectation, solve the selected Witten equation,
certify a channel edge, or populate a shell margin.

\begin{assessment}[The actual V6 advance]
\label{ass:YMA6-progress}
The open large-field row is no longer the raw quantity
\(\mathscr T_\Theta\).  The exact source-visible debit is
\[
 \mathbb L_\eta^{\rm ret}
 =\mathscr B^*(L_1^{-1}-\mathcal A_\eta^{-1})\mathscr B
 =\mathcal X_\eta^*(I-\mathcal K_\eta)^{-1}\mathcal X_\eta,
\]
a positive return of the original Wilson score through local curvature-defect
channels.  The positive follower shared with the Witten residual has been
removed once.  The comparison side is now the full Witten head
\(
 \mathbb D-\mathscr B^*\mathcal A_\eta^{-1}\mathscr B
\), not merely a pointwise Schur approximation, and it has a uniform finite
spatial tail.  What remains is literal population of one four-by-four head,
one source-cyclic return panel, and their common gauge short.
\end{assessment}

\begin{openproblem}[V7: execute the first return-renormalized axial card]
\label{op:YMA6-V7}
Preserve V6 and the frozen first axial mode.  Use the local-majorant split
with one predeclared \(\eta\), initially \(\eta=1/2\), and proceed in the
following order.
\begin{enumerate}[label=\textup{(V7.\arabic*)},leftmargin=5.2em]
\item Acquire the complete four-direction comparison Witten Gram
      \(\widehat{\mathbb C_{\mathcal A}}(q_{n_*})\) on the actual selected local cylinder, without replacing it by a
      Gaussian law.  Use the finite displacement head and the certified tail
      \eqref{eq:YMA6-Hcmp-second-tail}; do not substitute the flat Gaussian
      covariance.
\item Acquire the first source-cyclic defect moments
      \(\mathbb M_{\eta,n}\) on the same four sources.  Certify either a
      direct source-cyclic edge interval or the weighted contraction
      \eqref{eq:YMA6-weighted-channel-condition}; otherwise record the first
      unresolved return depth or channel radius.
\item Form \(P_{\eta,q}\) and \(R_{\eta,q}\), verify the Ward row
      \eqref{eq:YMA6-common-gauge-blocks}, and only then apply the common
      Moore--Penrose gauge reduction.
\item Compare the assembled transverse matrix with the exact reference
      \(s_0(q_{n_*})I_3\).  A negative upper interval on one transverse
      vector is the first physical pure-mode obstruction.  A near-unit
      source-cyclic defect channel without a negative assembled interval is
      a typed return survivor but not a negative Yang--Mills symbol.
\item Insert the submitted generated-flow and topological cards only after
      the pure local return is closed.  Evaluate the V4 coefficient
      \(\eta_4(\pi/2)\) only if it remains the first unresolved coefficient
      of this same assembled interval.
\end{enumerate}
No larger shell, graph, boundary, or generic positivity compiler is opened
before these values are decided.
\end{openproblem}

\begin{assessment}[V6 termination]
\label{ass:YMA6-termination}
The local majorant, return-renormalized resolvent identity, follower
cancellation, source-channel factorization, comparison-Witten locality,
weighted return theorem, and common-gauge axial reduction are \pass\ at
their stated mathematical scopes.  The complete selected response-matched
symbol remains \stopstatus: its first comparison head and source-cyclic
curvature-defect return have not yet been numerically or analytically
populated.  No direct physical negative mode has been produced.
\end{assessment}

\section*{Version V6 append-only continuation record}
\addcontentsline{toc}{section}{Version V6 append-only continuation record}
The complete Version~V5 source above is preserved byte for byte up to its
sole final \verb|\end{document}| command.  Its SHA--256 digest is
\begin{center}\scriptsize\ttfamily
a626dbf37b9511c67af8ed77d5e486ed7a97cb644c5bf8c4a9e44046efd29178.
\end{center}
For Version~V7, preserve this accumulated source, append new material
immediately before the final document-closing command, and increment the
filename to \texttt{\_V7.tex}.  Keep the exact block map, parity-tree
section, physical response, source normalization, first retained axial
mode, Haar compensation, local/topological support convention, and the
return-renormalized curvature split fixed before reading the first-card
values.

% =============================================================================
% END APPENDED VERSION V6 MATERIAL
% =============================================================================



% =============================================================================
% START APPENDED VERSION V7 MATERIAL
% =============================================================================

\clearpage
\hypersetup{
 pdftitle={Renewal Yang--Mills Reduced-Fibre Wilson Shell Symbol and Local-Vacuum Programme V7},
 pdfsubject={Boundary-source/Witten incidence, common-source Ward repair, and low-island return-tail control}
}
\pdfinfo{/Title (Renewal Yang--Mills Reduced-Fibre Wilson Shell Symbol and Local-Vacuum Programme V7)
 /Subject (Boundary-source/Witten incidence, common-source Ward repair, and low-island return-tail control)}

\section{Version V7: boundary-source/Witten incidence and the selected return island}
\label{sec:YMA7-direction}
% =============================================================================

Version~V6 is preserved above.  The newly attached boundary/cofinal and
shared survivor programmes change the order of the first-card attack, but not
its target.  The boundary programme proves a sharp weak-facing alternative:
a bounded selected-root subsequence is already a rough-Haar/contact trajectory
obstruction, while an escaping visible subsequence leaves an outer old/fine
source line and an independent radial four-leaf card.  In particular, the
cutoff-alignment coefficient and the radial attenuation coefficient are not
interchangeable.  The shared survivor programme proves, independently, that a
soft source-normalized Yang--Mills head carries a noncollapsed Witten current
only after the deterministic window and the same-history source/current
incidence have been established.

The first missing bridge is therefore not another estimate for the scalar
boundary contraction.  It is the physical incidence
\begin{equation}
 \boxed{
 \text{occurring outer boundary source}
 \quad\longleftrightarrow\quad
 \text{comparison/return split of the same mode score}.}
 \label{eq:YMA7-new-bridge}
\end{equation}
Once this bridge is formed, the return channel of Version~V6 can be shorted by
the source which actually occurs, its Ward row can be checked before gauge
elimination, and any persistent near-unit return can be localized on the
actual low physical mode island rather than on the whole source bank.

\begin{assessment}[Imported interfaces and nonduplication]
\label{ass:YMA7-imports}
The following statements are imported and are not reproved.
\begin{enumerate}[label=\textup{(I\arabic*)},leftmargin=4em]
\item The boundary programme's bounded/escaping weak-facing alternative,
      extinction of the centre-fibre source on the escaping branch, the
      positive outer replica card, and the independence of the outer cutoff
      card from the radial four-leaf card.
\item The shared survivor programme's total source/current Gram, its ancestry
      Schur residual, the direct-window fork, the four low-island relative
      moment rows, and the cofinal survivor theorem.
\item The Version~V6 identities
      \(
       L_1=\mathcal A_\eta-\mathcal V_\eta\mathcal V_\eta^*
      \),
      \(
       \mathcal K_\eta=\mathcal V_\eta^*\mathcal A_\eta^{-1}
                        \mathcal V_\eta
      \), and
      \(
       \mathbb L_\eta^{\rm ret}
       =\mathcal X_\eta^*(I-\mathcal K_\eta)^{-1}\mathcal X_\eta
      \).
\end{enumerate}
The new results below are the exact-form realization of the missing incidence,
its comparison/return arm decomposition, the resulting common-source Ward
identity, its quantitative cutoff handoff, and a low-island replacement for
the global weighted-channel condition of Version~V6.  The attached session
report's finite-volume one-loop table is retained only in its stated
\emph{measured} status: it contains no certified selected-root
boundary-source/mode-score incidence and is not used as an outward physical
interval here.
\end{assessment}

\begin{firewall}[No boundary coefficient is inserted as a shell eigenvalue]
\label{fire:YMA7-no-substitution}
The old/fine fan residual, the selected radial coefficient, the exact flat
fibre gap, and the reduced-fibre Wilson Hessian are different objects.  No
scalar from the boundary programme is substituted for
\(
 \mathbb EK-\Cov(j)
\), for the source-cyclic return edge, or for the physical reference.  The
boundary packet enters Version~V7 only through an occurring source line and,
when available, its adjacent-cutoff overlap.
\end{firewall}

% =============================================================================
\section{A canonical exact-form lift of the occurring boundary source}
\label{sec:YMA7-exact-lift}
% =============================================================================

Fix one finite response-matched reduced-fibre law \(\nu\) and one nonzero
momentum--colour card.  Let
\begin{equation}
 \mathcal H_0=L_0^2(\nu)
 :=\left\{f\in L^2(\nu):\mathbb Ef=0\right\},
 \qquad
 L_0=d_\nu^*d
 \label{eq:YMA7-scalar-Witten}
\end{equation}
be the mean-zero scalar Witten space and scalar Laplacian.  On one-forms retain
\begin{equation}
 L_1=dd_\nu^*+d_\nu^*d
 =\mathcal A_\eta-\mathcal V_\eta\mathcal V_\eta^*
 \label{eq:YMA7-one-form-split}
\end{equation}
from Version~V6.  All inverses below are taken on their supported ranges.  On
a finite compact reduced fibre, \(L_0\) is strictly positive on
\(\mathcal H_0\), and \(L_1\) is strictly positive on the exact one-form
range generated by the card.

\begin{lemma}[Exact-form Witten isometry]
\label{lem:YMA7-exact-isometry}
For every \(f,g\in\mathcal H_0\),
\begin{equation}
 \boxed{
 \ip{df}{L_1^{-1}dg}_{\nu}
 =\ip{f}{g}_{\nu}.}
 \label{eq:YMA7-exact-isometry}
\end{equation}
Equivalently, \(L_1^{-1/2}d\) is an isometry from the mean-zero scalar
space onto its exact one-form range.
\end{lemma}

\begin{proof}
The weighted de~Rham identities give
\(
 L_1d=dL_0
\)
on the smooth core.  Put \(u=L_0^{-1}g\).  Then
\(
 L_1du=dL_0u=dg
\), so the supported exact-form inverse satisfies
\(
 L_1^{-1}dg=du
\).  Weighted integration by parts therefore gives
\[
 \ip{df}{L_1^{-1}dg}_{\nu}
 =\ip{df}{du}_{\nu}
 =\ip f{d_\nu^*du}_{\nu}
 =\ip f{L_0u}_{\nu}
 =\ip{f}{g}_{\nu}.
\]
Density closes the identity on the form domain.
\end{proof}

Define the comparison and return lifts
\begin{align}
 \mathscr U_{\eta}^{\rm cmp}f
 &:=\mathcal A_\eta^{-1/2}df,
 \label{eq:YMA7-Ucmp}\\
 \mathscr U_{\eta}^{\rm ret}f
 &:=(I-\mathcal K_\eta)^{-1/2}
    \mathcal V_\eta^*\mathcal A_\eta^{-1}df.
 \label{eq:YMA7-Uret}
\end{align}
The second vector lies in the supported curvature-defect channel.

\begin{theorem}[Comparison/return Pythagoras for every occurring scalar source]
\label{thm:YMA7-source-Pythagoras}
For all \(f,g\in\mathcal H_0\),
\begin{equation}
 \boxed{
 \ip{f}{g}_{\nu}
 =\ip{\mathscr U_{\eta}^{\rm cmp}f}
          {\mathscr U_{\eta}^{\rm cmp}g}
  +\ip{\mathscr U_{\eta}^{\rm ret}f}
          {\mathscr U_{\eta}^{\rm ret}g}.}
 \label{eq:YMA7-source-Pythagoras}
\end{equation}
Thus
\begin{equation}
 \boxed{
 \mathscr U_\eta f
 :=\mathscr U_{\eta}^{\rm cmp}f
    \oplus\mathscr U_{\eta}^{\rm ret}f}
 \label{eq:YMA7-total-source-lift}
\end{equation}
is the canonical isometry associated with the fixed, predeclared Version~V6
comparison/return split.  Its direct-sum target is the comparison one-form
space plus the curvature-defect return channel.

Let
\(
 \mathscr F_q:E_q\to\mathcal H_0
\)
be the centered mode-score synthesis and
\(
 \mathscr B_q=d\mathscr F_q
\).
Define
\begin{align}
 Q_{\eta,q}^{\rm cmp}
 &:=\mathcal A_\eta^{-1/2}\mathscr B_q,
 \label{eq:YMA7-Qcmp}\\
 Q_{\eta,q}^{\rm ret}
 &:=(I-\mathcal K_\eta)^{-1/2}
    \mathcal V_\eta^*\mathcal A_\eta^{-1}\mathscr B_q.
 \label{eq:YMA7-Qret}
\end{align}
Then the complete connected mode covariance has the exact orthogonal split
\begin{equation}
 \boxed{
 \mathscr F_q^*\mathscr F_q
 =(Q_{\eta,q}^{\rm cmp})^*Q_{\eta,q}^{\rm cmp}
 +(Q_{\eta,q}^{\rm ret})^*Q_{\eta,q}^{\rm ret}.}
 \label{eq:YMA7-covariance-split}
\end{equation}
The second summand is exactly the Version~V6 return
\(
 \mathbb L_\eta^{\rm ret}(q)
\).
\end{theorem}

\begin{proof}
Version~V6 gives the Woodbury identity
\[
 L_1^{-1}
 =\mathcal A_\eta^{-1}
  +\mathcal A_\eta^{-1}\mathcal V_\eta
   (I-\mathcal K_\eta)^{-1}
   \mathcal V_\eta^*\mathcal A_\eta^{-1}.
\]
Insert this identity between \(df\) and \(dg\), and use
\Cref{lem:YMA7-exact-isometry}.  The two resulting terms are precisely the
two inner products in \eqref{eq:YMA7-source-Pythagoras}.  Applying the same
identity to the columns \(d\mathscr F_qh\) and polarizing over
\(h\in E_q\) proves \eqref{eq:YMA7-covariance-split}.  Its return term is
\[
 \mathscr B_q^*\mathcal A_\eta^{-1}\mathcal V_\eta
 (I-\mathcal K_\eta)^{-1}
 \mathcal V_\eta^*\mathcal A_\eta^{-1}\mathscr B_q
 =\mathcal X_\eta^*(I-\mathcal K_\eta)^{-1}\mathcal X_\eta.
\]
\end{proof}

\begin{assessment}[Why this removes an arbitrary ambient router]
\label{ass:YMA7-canonical-router}
The boundary and reduced-fibre notes previously placed the occurring source
and the Witten current in separately presented Hilbert spaces and therefore
required a physical current-to-boundary identification.  Once the boundary
writer itself occurs as a scalar function on the same reduced-fibre law, its
exact-form Witten lift \eqref{eq:YMA7-total-source-lift} is canonical.  The
mixed source--score row is then the ordinary covariance on that law.  This
does not identify an arbitrary half-slab amplitude, boundary instrument, or
spatial route; it gives exactly the minimum incidence needed by the selected
scalar source and mode score.
\end{assessment}

% =============================================================================
\section{The arm-resolved source/score Gram and its ancestry residual}
\label{sec:YMA7-arm-Gram}
% =============================================================================

Let \(m\in\mathcal H_0\) be the centered occurring outer boundary source on
the selected card, after projection to the same retained translation--colour
character as the mode-score bank, and normalize it by
\begin{equation}
 \norm m_{L^2(\nu)}=1.
 \label{eq:YMA7-m-normalized}
\end{equation}
If its variance vanishes, the source-visible card stops before this
normalization.  Put
\begin{equation}
 s_A:=\mathscr U_{\eta}^{\rm cmp}m,
 \qquad
 s_R:=\mathscr U_{\eta}^{\rm ret}m,
 \label{eq:YMA7-source-arms}
\end{equation}
\begin{equation}
 a:=\norm{s_A}^2,
 \qquad
 b:=\norm{s_R}^2,
 \qquad
 a+b=1,
 \label{eq:YMA7-ab}
\end{equation}
and define the two mixed rows
\begin{equation}
 x_Ah:=\ip{s_A}{Q_{\eta,q}^{\rm cmp}h},
 \qquad
 x_Rh:=\ip{s_R}{Q_{\eta,q}^{\rm ret}h}.
 \label{eq:YMA7-xA-xR}
\end{equation}
Also put
\begin{equation}
 C_A:=(Q_{\eta,q}^{\rm cmp})^*Q_{\eta,q}^{\rm cmp},
 \qquad
 C_R:=(Q_{\eta,q}^{\rm ret})^*Q_{\eta,q}^{\rm ret},
 \label{eq:YMA7-CA-CR}
\end{equation}
and let
\begin{equation}
 y h:=\ip m{\mathscr F_qh}_{\nu}
 \label{eq:YMA7-total-mixed-row}
\end{equation}
be the direct same-history boundary-source/mode-score row.

\begin{theorem}[Arm-resolved incidence, assembly dispersion, and minimum missing source]
\label{thm:YMA7-arm-incidence}
The quantities above satisfy
\begin{equation}
 \boxed{y=x_A+x_R,}
 \label{eq:YMA7-y-split}
\end{equation}
and the complete split Gram is positive:
\begin{equation}
 \boxed{
 \mathbb G_{\eta,q}^{\rm split}
 =\begin{pmatrix}
    a&0&x_A\\
    0&b&x_R\\
    x_A^*&x_R^*&C_A+C_R
   \end{pmatrix}\succeq0.}
 \label{eq:YMA7-split-Gram}
\end{equation}
Let \(a^\dagger\) and \(b^\dagger\) be the scalar Moore--Penrose inverses and
set
\begin{equation}
 \mathbb I_A:=C_A-x_A^*a^\dagger x_A,
 \qquad
 \mathbb I_R:=C_R-x_R^*b^\dagger x_R.
 \label{eq:YMA7-arm-residuals}
\end{equation}
Both are positive semidefinite.  The residual after shorting by the
\emph{assembled} occurring source is
\begin{equation}
 \boxed{
 \mathbb I_{\partial,q}^{\rm sc\mid src}
 :=C_A+C_R-y^*y\succeq0.}
 \label{eq:YMA7-total-ancestry-residual}
\end{equation}
If \(a,b>0\), then
\begin{equation}
 \boxed{
 \mathbb I_{\partial,q}^{\rm sc\mid src}
 =\mathbb I_A+\mathbb I_R+\Delta_{\eta,q}^{\rm route},}
 \label{eq:YMA7-dispersion-identity}
\end{equation}
where the router-dispersion Gram is
\begin{equation}
 \boxed{
 \Delta_{\eta,q}^{\rm route}
 =\left(
   \sqrt{\frac ba}\,x_A-\sqrt{\frac ab}\,x_R
  \right)^*
  \left(
   \sqrt{\frac ba}\,x_A-\sqrt{\frac ab}\,x_R
  \right)\succeq0.}
 \label{eq:YMA7-route-dispersion}
\end{equation}
If one of \(a,b\) vanishes, positivity forces the corresponding mixed row to
vanish and \eqref{eq:YMA7-dispersion-identity} holds with
\(\Delta_{\eta,q}^{\rm route}=0\).

Writing \(Q=Q_{\eta,q}^{\rm cmp}\oplus Q_{\eta,q}^{\rm ret}\) and
\(s=s_A\oplus s_R\), one has
\begin{equation}
 \mathbb I_{\partial,q}^{\rm sc\mid src}
 =Q^*(I-ss^*)Q.
 \label{eq:YMA7-residual-projection}
\end{equation}
Its rank is the minimum number of additional orthogonal response-source
directions required to synthesize the complete selected mode score after the
outer boundary source has been retained.
\end{theorem}

\begin{proof}
By \Cref{thm:YMA7-source-Pythagoras},
\[
 \ip m{\mathscr F_qh}
 =\ip{s_A}{Q_{\eta,q}^{\rm cmp}h}
  +\ip{s_R}{Q_{\eta,q}^{\rm ret}h},
\]
which proves \eqref{eq:YMA7-y-split}.  The matrix
\eqref{eq:YMA7-split-Gram} is the Gram of the two orthogonal source vectors
\(s_A\oplus0\), \(0\oplus s_R\), and the score synthesis
\(Q_{\eta,q}^{\rm cmp}\oplus Q_{\eta,q}^{\rm ret}\).  Its positivity is
therefore exact.  Supported Schur complementation by the scalar blocks gives
\(\mathbb I_A\succeq0\) and \(\mathbb I_R\succeq0\).  If \(a=0\), then
\(s_A=0\) and hence \(x_A=0\); the same argument applies to \(b\).

The assembled source vector \(s=s_A\oplus s_R\) has unit norm by
\eqref{eq:YMA7-ab}.  Thus the Gram of the score component orthogonal to
\(s\) is
\[
 Q^*(I-ss^*)Q=C_A+C_R-(x_A+x_R)^*(x_A+x_R),
\]
which is \eqref{eq:YMA7-total-ancestry-residual}.  For \(a,b>0\), expand
\eqref{eq:YMA7-route-dispersion}.  The coefficient of \(x_A^*x_A\) in
\(\mathbb I_A+\Delta^{\rm route}\) is
\(
 -a^{-1}+b/a=-1
\), because \(a+b=1\).  The coefficient of \(x_R^*x_R\) is likewise
\(-1\), and the two mixed terms are
\(-x_A^*x_R-x_R^*x_A\).  This proves
\eqref{eq:YMA7-dispersion-identity}.  The zero-arm cases were already
identified.

Finally, \(Q^*(I-ss^*)Q\) is the Gram of the orthogonal residual synthesis
\((I-ss^*)Q\).  Its rank is therefore the dimension of its range, which is
the minimum number of orthogonal source coordinates needed to represent that
residual.
\end{proof}

\begin{corollary}[Exact source-ancestry decomposition of the physical mode Hessian]
\label{cor:YMA7-H-ancestry}
Let \(D_q\) be the deterministic mode Hessian on the same coefficient bank.
Then
\begin{equation}
 \boxed{
 H_q
 =D_q-y^*y-\mathbb I_{\partial,q}^{\rm sc\mid src}.}
 \label{eq:YMA7-H-ancestry}
\end{equation}
The term \(y^*y\) is the score covariance carried by the occurring outer
boundary source.  Only the positive residual
\(\mathbb I_{\partial,q}^{\rm sc\mid src}\) lies outside that source line.
In particular,
\begin{equation}
 C_R=x_R^*b^\dagger x_R+\mathbb I_R
 \label{eq:YMA7-return-ancestry}
\end{equation}
separates the source-following curvature-defect return from its orthogonal
return-source residual.
\end{corollary}

\begin{proof}
The score--Hessian identity is
\(
 H_q=D_q-(C_A+C_R)
\).  Substitute
\eqref{eq:YMA7-total-ancestry-residual}.  The return identity is the
definition of \(\mathbb I_R\).
\end{proof}

\begin{firewall}[An ancestry residual is not automatically a new primitive]
\label{fire:YMA7-ancestry-not-primitive}
A positive residual in \eqref{eq:YMA7-total-ancestry-residual} is the minimum
additional response-source dimension relative to the selected outer source
line.  It is not automatically a new microscopic source, a positive-range
boundary survivor, or an Osterwalder--Schrader excitation.  Those promotions
require, respectively, an occurrence interpretation, physical-radius
transport, and an independently calibrated time-semigroup incidence.
\end{firewall}

% =============================================================================
\section{Ward cancellation requires assembly before source shorting}
\label{sec:YMA7-Ward}
% =============================================================================

\begin{theorem}[Common-source Ward identity and the exact separate-short debit]
\label{thm:YMA7-Ward-assembly}
Let \(g=d_0^{\rm c}\xi\in E_{g,q}\) be a coarse gauge direction.  Assume the
outer boundary writer \(m\) is gauge invariant on the same physical
cylinder.  Then
\begin{equation}
 \boxed{y(g)=0,
 \qquad x_A(g)+x_R(g)=0.}
 \label{eq:YMA7-Ward-cross}
\end{equation}
If \(a,b>0\), the route-dispersion term on the gauge line is
\begin{equation}
 \boxed{
 \Delta_{\eta,q}^{\rm route}[g]
 =\frac{|x_A(g)|^2}{a}+\frac{|x_R(g)|^2}{b}.}
 \label{eq:YMA7-Ward-dispersion}
\end{equation}
Consequently
\begin{equation}
 \boxed{
 \mathbb I_{\partial,q}^{\rm sc\mid src}[g]
 =C_A[g]+C_R[g],}
 \label{eq:YMA7-Ward-total-residual}
\end{equation}
whereas shorting the comparison and return arms separately and then adding
would give the smaller quantity
\begin{equation}
 \mathbb I_A[g]+\mathbb I_R[g]
 =C_A[g]+C_R[g]-\Delta_{\eta,q}^{\rm route}[g].
 \label{eq:YMA7-Ward-separate-short}
\end{equation}
The positive difference is exactly the artificial follower created by
forgetting that the two arm covariances cancel on one gauge-invariant source.
\end{theorem}

\begin{proof}
Let \(f_g=\mathscr F_qg\) be the centered score in the gauge direction.
Differentiate the expectation of the gauge-invariant writer \(m\) along the
coarse gauge orbit.  Invariance of both the law and the writer gives
\[
 0=\left.\frac{\dd}{\dd t}\right|_{t=0}\mathbb E_t m
  =\ip m{f_g}_{\nu}=y(g).
\]
Equation \eqref{eq:YMA7-y-split} then gives
\(x_A(g)+x_R(g)=0\).  Substitute \(x_R(g)=-x_A(g)\) into
\eqref{eq:YMA7-route-dispersion}.  Since \(a+b=1\),
\[
 \left(\sqrt{b/a}+\sqrt{a/b}\right)^2|x_A(g)|^2
 =\left(a^{-1}+b^{-1}\right)|x_A(g)|^2,
\]
which is \eqref{eq:YMA7-Ward-dispersion}.  The total residual formula follows
from \(y(g)=0\); the separate-short formula is
\eqref{eq:YMA7-dispersion-identity} rearranged on \(g\).
\end{proof}

\begin{countertheorem}[Total marginals do not determine return ancestry]
\label{cth:YMA7-total-not-return}
There are two finite source/score cards with the same comparison and return
marginal variances, the same total source--score row, and the same total score
covariance, but different return ancestry.

Take \(a=b=1/2\), one scalar mode, and
\begin{equation}
 C_A=C_R=\frac12,
 \qquad y=0.
 \label{eq:YMA7-two-card-marginals}
\end{equation}
In the first card set \(x_A=x_R=0\).  In the second set
\begin{equation}
 x_A=\frac12,
 \qquad x_R=-\frac12.
 \label{eq:YMA7-two-card-crosses}
\end{equation}
Both split Grams \eqref{eq:YMA7-split-Gram} are positive, but
\begin{equation}
 \boxed{
 \mathbb I_R^{(0)}=\frac12,
 \qquad
 \mathbb I_R^{(1)}=0.}
 \label{eq:YMA7-return-nonidentifiability}
\end{equation}
Thus the total source--score covariance and the two score-arm marginals do not
identify whether the return is a follower of the occurring boundary source.
The split mixed row \(x_R\), or an equivalent same-history analyser, is an
independent physical acquisition.
\end{countertheorem}

\begin{proof}
The first split Gram is diagonal with entries \(1/2,1/2,1\).  The second is
the Gram of
\[
 s_A=2^{-1/2}e_A,
 \qquad
 s_R=2^{-1/2}e_R,
 \qquad
 Q=2^{-1/2}e_A-2^{-1/2}e_R
\]
in \(\R e_A\oplus\R e_R\), so it is positive.  Both have
\(y=x_A+x_R=0\) and \(C_A+C_R=1\).  The scalar supported Schur formula gives
\[
 \mathbb I_R^{(0)}=\frac12,
 \qquad
 \mathbb I_R^{(1)}
 =\frac12-\frac{(1/2)^2}{1/2}=0.
\]
\end{proof}

% =============================================================================
\section{The target-minimal six-by-six first axial incidence card}
\label{sec:YMA7-six-card}
% =============================================================================

Return to the first axial direction bank \(E_q\cong\C^4\) of Version~V6 and
one fixed colour polarization.  The two scalar source arms together with the
four mode-score columns form a \(6\times6\) Hermitian Gram.  It is not
necessary to populate two unrelated source experiments.

Let
\begin{equation}
 \eta_m:=\mathcal A_\eta^{-1}dm.
 \label{eq:YMA7-eta-m}
\end{equation}
Then
\begin{equation}
 a=\ip{dm}{\eta_m}_{\nu},
 \qquad
 x_Ah=\ip{\eta_m}{\mathscr B_qh}_{\nu}.
 \label{eq:YMA7-comparison-source-rows}
\end{equation}
The total mixed row is acquired directly on two independent copies of the
same cylinder:
\begin{equation}
 \boxed{
 y(h)=\frac12\mathbb E\left[
  \overline{m_1-m_2}
  \bigl(f_h^{(1)}-f_h^{(2)}\bigr)
 \right],
 \qquad f_h=\mathscr F_qh.}
 \label{eq:YMA7-two-replica-y}
\end{equation}
The real sine--cosine implementation uses ordinary real polarization of this
identity.

\begin{theorem}[One comparison solve and one direct mixed row determine the split card]
\label{thm:YMA7-minimal-card}
Assume the Version~V6 comparison and return mode Grams \(C_A,C_R\) have been
acquired on the actual selected law, and the normalized boundary source
\(m\) occurs on that same law.  Then the following data determine the entire
split Gram \eqref{eq:YMA7-split-Gram}:
\begin{equation}
 \boxed{
 C_A,\quad C_R,\quad
 a=\ip{dm}{\mathcal A_\eta^{-1}dm},\quad
 x_A=\eta_m^*\mathscr B_q,\quad
 y=m^*\mathscr F_q.}
 \label{eq:YMA7-minimal-data}
\end{equation}
Indeed,
\begin{equation}
 \boxed{b=1-a,
 \qquad x_R=y-x_A.}
 \label{eq:YMA7-derived-rows}
\end{equation}
Thus one comparison Witten solve for the boundary source and one direct
same-history boundary-source/mode-score row are sufficient.  A separately
sampled return-source line is optional as a held-out check, not a new source
birth.

The card must pass all of the exact identities
\begin{equation}
 \boxed{
 a+b=1,
 \qquad x_A+x_R=y,
 \qquad
 \mathbb G_{\eta,q}^{\rm split}\succeq0,
 \qquad y\,d(q)=0.}
 \label{eq:YMA7-card-checks}
\end{equation}
Failure of an algebraic identity or of positivity rejects the submitted
mixed card as ill typed; it is not a negative Yang--Mills mode.
\end{theorem}

\begin{proof}
Equation \eqref{eq:YMA7-comparison-source-rows} is the definition of the
comparison lift written without a square root.  The first identity in
\eqref{eq:YMA7-derived-rows} is the source Pythagoras
\eqref{eq:YMA7-ab}; the second is \eqref{eq:YMA7-y-split}.  The two-replica
formula follows by expanding the product and using independence together
with \(\mathbb Em=\mathbb Ef_h=0\).  Positivity and the Ward row are
\Cref{thm:YMA7-arm-incidence,thm:YMA7-Ward-assembly}.
\end{proof}

\begin{proposition}[Validated interval for the total ancestry residual]
\label{prop:YMA7-residual-interval}
Suppose Hermitian approximants satisfy
\begin{equation}
 \norm{C_A-\widehat C_A}\le\varepsilon_A,
 \qquad
 \norm{C_R-\widehat C_R}\le\varepsilon_R,
 \qquad
 \norm{y-\widehat y}\le\varepsilon_y.
 \label{eq:YMA7-approx-errors}
\end{equation}
Put
\begin{equation}
 \widehat I
 :=\widehat C_A+\widehat C_R-\widehat y^*\widehat y,
 \qquad
 \varepsilon_I
 :=\varepsilon_A+\varepsilon_R
   +(2\norm{\widehat y}+\varepsilon_y)\varepsilon_y.
 \label{eq:YMA7-Ihat}
\end{equation}
Then
\begin{equation}
 \boxed{
 \norm{\mathbb I_{\partial,q}^{\rm sc\mid src}-\widehat I}
 \le\varepsilon_I.}
 \label{eq:YMA7-I-error}
\end{equation}
Every ordered eigenvalue of the physical ancestry residual therefore lies
within \(\varepsilon_I\) of the corresponding eigenvalue of
\(\widehat I\).  If \(\lambda_{\min}(\widehat I)<-\varepsilon_I\), the card
is inconsistent with a common positive source/score cylinder.  If
\(\widehat I-\varepsilon_I I\succeq0\), it supplies a certified outward
ancestry interval.
\end{proposition}

\begin{proof}
The Gram estimate
\[
 \norm{y^*y-\widehat y^*\widehat y}
 \le(\norm y+\norm{\widehat y})\norm{y-\widehat y}
 \le(2\norm{\widehat y}+\varepsilon_y)\varepsilon_y
\]
combined with the two covariance errors proves
\eqref{eq:YMA7-I-error}.  Weyl's Hermitian perturbation inequality gives the
eigenvalue intervals.  Exact positivity of
\eqref{eq:YMA7-total-ancestry-residual} gives the inconsistency test.
\end{proof}

\begin{request}[Literal first-card acquisition after the new attachments]
\label{req:YMA7-first-card}
On the frozen first axial mode and one colour, acquire in this order:
\begin{enumerate}[label=\textup{(A\arabic*)},leftmargin=4em]
\item the outer-source variance and the centre-fibre fraction supplied by the
      boundary card; stop if the variance has no positive lower endpoint;
\item the comparison solve \(\eta_m=\mathcal A_\eta^{-1}dm\), giving
      \(a\) and \(x_A\);
\item the direct two-replica row \(y\) in
      \eqref{eq:YMA7-two-replica-y};
\item the already requested Version~V6 matrices \(C_A,C_R\), after which
      \(b,x_R,\mathbb I_A,\mathbb I_R,
       \Delta^{\rm route},\mathbb I_{\partial}^{\rm sc\mid src}\)
      are formed without another source query;
\item the exact checks \eqref{eq:YMA7-card-checks}, followed by the common
      gauge reduction of Version~V6 and only then the physical margin test.
\end{enumerate}
The finite card contains two source columns and four score columns.  It is
therefore a six-by-six Gram per colour, not an unbounded boundary-sector or
current basis.
\end{request}

% =============================================================================
\section{Quantitative native-to-outer and adjacent-cutoff handoff}
\label{sec:YMA7-handoff}
% =============================================================================

The boundary programme distinguishes the native fine writer from its
centre-orbit conditional line.  The following estimate states exactly when
that scalar convergence can be inserted into the mode incidence.

\begin{theorem}[Native-to-outer transfer of the mixed mode row]
\label{thm:YMA7-native-outer}
Let \(\widehat F,\widehat M\in L_0^2(\nu)\) be the normalized native and
outer conditional source lines on one escaping visible card, and suppose
\begin{equation}
 \norm{\widehat F-\widehat M}^2\le2t_{\rm cf}.
 \label{eq:YMA7-FM-input}
\end{equation}
For the mode-score synthesis \(\mathscr F_q:E_q\to L_0^2(\nu)\), put
\begin{equation}
 y_F=\widehat F^*\mathscr F_q,
 \qquad
 y_M=\widehat M^*\mathscr F_q.
 \label{eq:YMA7-yF-yM}
\end{equation}
Then
\begin{equation}
 \boxed{
 \norm{y_F-y_M}
 \le\norm{\mathscr F_q}\sqrt{2t_{\rm cf}},}
 \label{eq:YMA7-y-transfer}
\end{equation}
\begin{equation}
 \boxed{
 \norm{y_F^*y_F-y_M^*y_M}
 \le2\norm{\mathscr F_q}^2\sqrt{2t_{\rm cf}}.}
 \label{eq:YMA7-follower-transfer}
\end{equation}
Consequently the native and outer source followers coalesce on an escaping
branch whenever \(t_{{\rm cf},r}\to0\) and the selected score-synthesis norms
are uniformly bounded.  This conclusion does not populate \(y_M\), its
comparison/return split, or the independent radial coefficient.
\end{theorem}

\begin{proof}
For every unit \(h\in E_q\), Cauchy--Schwarz gives
\[
 |(y_F-y_M)h|
 =|\ip{\widehat F-\widehat M}{\mathscr F_qh}|
 \le\norm{\widehat F-\widehat M}\norm{\mathscr F_q},
\]
which proves \eqref{eq:YMA7-y-transfer}.  Since both source vectors are unit,
\(\norm{y_F},\norm{y_M}\le\norm{\mathscr F_q}\).  The standard Gram
estimate
\(
 \norm{X^*X-Y^*Y}\le(\norm X+\norm Y)\norm{X-Y}
\)
then gives \eqref{eq:YMA7-follower-transfer}.
\end{proof}

Now consider adjacent cutoffs.  Let \(m_r\in\mathcal H_r\) and
\(m_{r+1}\in\mathcal H_{r+1}\) be unit boundary source lines, let
\(U_r:\mathcal H_r\to\mathcal H_{r+1}\) be the physical occurrence
isometry, and let \(T_r:E_r\to E_{r+1}\) be the coefficient identification.
Set
\begin{equation}
 \alpha_r:=\ip{m_{r+1}}{U_rm_r},
 \qquad
 R_r^{\partial}:=1-|\alpha_r|^2.
 \label{eq:YMA7-boundary-overlap}
\end{equation}
The first scalar is exactly the type of direct old/fine mixed fan coefficient
isolated by the boundary programme, but it may be used here only when the fan
line has been physically identified with the source line \(m_r\) of the
mode card.

\begin{theorem}[Boundary-overlap plus score transport controls incidence transport]
\label{thm:YMA7-incidence-transport}
There is a unit phase \(\zeta_r\) such that
\begin{equation}
 \norm{m_{r+1}-\zeta_rU_rm_r}
 =\sqrt{2(1-|\alpha_r|)}
 \le\sqrt{2R_r^{\partial}}.
 \label{eq:YMA7-source-line-distance}
\end{equation}
Let
\begin{equation}
 e_r^{\rm sc}
 :=\norm{\mathscr F_{r+1}T_r-U_r\mathscr F_r},
 \qquad
 y_r=m_r^*\mathscr F_r.
 \label{eq:YMA7-score-defect}
\end{equation}
Then
\begin{equation}
 \boxed{
 \norm{y_{r+1}T_r-\overline{\zeta_r}y_r}
 \le e_r^{\rm sc}
     +\sqrt{2R_r^{\partial}}\,\norm{\mathscr F_r}.}
 \label{eq:YMA7-y-cutoff-bound}
\end{equation}
Hence a uniform score ceiling together with
\begin{equation}
 \sum_r e_r^{\rm sc}<\infty,
 \qquad
 \sum_r\sqrt{R_r^{\partial}}<\infty
 \label{eq:YMA7-y-summable}
\end{equation}
transports the mixed source--score row after the predeclared orientation
phases are applied.

The scalar boundary overlap alone does not transport \(x_A\), \(x_R\), or
\(\mathbb I_R\).  Their transport additionally requires the comparison and
return Witten arms to belong to the same adjacent-cutoff rectangle.
\end{theorem}

\begin{proof}
Use the convention that the Hilbert inner product is conjugate-linear in its
first entry and choose \(\zeta_r\) so that
\(
 \zeta_r\alpha_r=|\alpha_r|
\), with arbitrary phase if \(\alpha_r=0\).  Expanding the squared distance
gives the equality in \eqref{eq:YMA7-source-line-distance}; the inequality
uses \(1-t\le1-t^2\) for \(0\le t\le1\).

Insert \(m_{r+1}^*U_r\mathscr F_r\):
\[
 \begin{aligned}
 y_{r+1}T_r-\overline{\zeta_r}y_r
 ={}&m_{r+1}^*(\mathscr F_{r+1}T_r-U_r\mathscr F_r)\\
 &+(m_{r+1}^*U_r-\overline{\zeta_r}m_r^*)\mathscr F_r.
 \end{aligned}
\]
The first term has norm at most \(e_r^{\rm sc}\).  The second is bounded by
\[
 \norm{U_r^*m_{r+1}-\zeta_rm_r}\norm{\mathscr F_r}
 \le\norm{m_{r+1}-\zeta_rU_rm_r}\norm{\mathscr F_r},
\]
which proves \eqref{eq:YMA7-y-cutoff-bound}.  Summability gives a Cauchy
incidence row in the oriented direct limit.  The final qualification follows
because a total source line and total score synthesis do not determine their
comparison/return mixed rows, as shown by
\Cref{cth:YMA7-total-not-return}.
\end{proof}

\begin{firewall}[The fan overlap is used only on the identified source line]
\label{fire:YMA7-fan-identification}
The degree-twelve fan overlap from YM--B is a boundary transport datum.  It
may enter \eqref{eq:YMA7-y-cutoff-bound} only after the selected mode source
has been shown to be that occurring fan line, with common normalization and
orientation.  Equality of loop labels, dimensions, or separate variances is
not such an identification.  If the escaping outer line is a different
boundary source, its own adjacent mixed coefficient must be acquired.
\end{firewall}

% =============================================================================
\section{Low-island return control without a global weighted-channel norm}
\label{sec:YMA7-low-island}
% =============================================================================

Version~V6 gave the sufficient global condition
\(
 \norm{\mathcal K_\eta}_{\mu,{\rm Sch}}<1
\).
The shared survivor programme shows that this is stronger than necessary on
a selected soft sequence.  Only the return generated by the collapsing
physical island must be tight.

After the Version~V6 common-gauge reduction and reference whitening, write
\begin{equation}
 A_r=D_r-Q_r^*Q_r\succeq0,
 \label{eq:YMA7-normalized-return-card}
\end{equation}
where \(D_r\) is the normalized comparison head and
\(Q_r:E_{\perp,r}\to\mathscr C_{\eta,r}\) is the physical, gauge-reduced return
synthesis, so that \(Q_r^*Q_r\) is the normalized return matrix.  This
section applies only on the nonnegative branch.  A certified negative mode
has already terminated as a direct obstruction.

Let
\begin{equation}
 P_r=\one_{[0,\rho_r]}(A_r)
 \label{eq:YMA7-low-island-P}
\end{equation}
be a separated nonzero low-island projection.  Put
\begin{equation}
 H_{\eta,r}:=I-\mathcal K_{\eta,r}>0
 \label{eq:YMA7-H-channel}
\end{equation}
and let \(R_{\partial,r}\) be the physical boundary-range operator applied
after the return has been formed.  Define the two positive source forms
\begin{equation}
 K_{{\rm edge},r}:=Q_r^*H_{\eta,r}^{-1}Q_r,
 \qquad
 K_{\partial,r}:=Q_r^*R_{\partial,r}^2Q_r.
 \label{eq:YMA7-two-return-moments}
\end{equation}

\begin{theorem}[Two relative edges control the selected near-unit and radius tails]
\label{thm:YMA7-low-island-return}
Assume a direct window on the low island,
\begin{equation}
 \boxed{
 d_-P_r\preceq P_rD_rP_r\preceq d_+P_r,
 \qquad
 0\le\rho_r<d_-.}
 \label{eq:YMA7-direct-window}
\end{equation}
Then
\begin{equation}
 \boxed{
 (d_- -\rho_r)P_r
 \preceq P_rQ_r^*Q_rP_r
 \preceq d_+P_r.}
 \label{eq:YMA7-return-window}
\end{equation}
Suppose the two finite generalized relative edges obey
\begin{equation}
 \boxed{
 P_rK_{{\rm edge},r}P_r
 \preceq\kappa_{{\rm edge},r}P_rQ_r^*Q_rP_r,
 \qquad
 P_rK_{\partial,r}P_r
 \preceq\kappa_{\partial,r}P_rQ_r^*Q_rP_r.}
 \label{eq:YMA7-relative-return-edges}
\end{equation}
For
\begin{equation}
 E_{r,\varepsilon}
 :=\one_{(0,\varepsilon]}(H_{\eta,r}),
 \qquad
 \Pi_{r,R}:=\one_{[0,R]}(R_{\partial,r}),
 \label{eq:YMA7-return-screens}
\end{equation}
one has
\begin{equation}
 \boxed{
 \norm{E_{r,\varepsilon}Q_rP_r}
 \le\sqrt{\varepsilon\kappa_{{\rm edge},r}d_+},}
 \label{eq:YMA7-near-unit-tail}
\end{equation}
\begin{equation}
 \boxed{
 \norm{(I-\Pi_{r,R})Q_rP_r}
 \le\frac{\sqrt{\kappa_{\partial,r}d_+}}{R}.}
 \label{eq:YMA7-range-tail}
\end{equation}
Thus uniform bounds for the two relative edges on the actual low island give
tightness against a near-unit curvature-defect channel and against positive
boundary-range escape, without requiring a weighted contraction bound on
source directions outside that island.

On each finite card the sharp constants are the attained generalized edges
\begin{equation}
 \boxed{
 \begin{aligned}
 \kappa_{{\rm edge},r}^{\sharp}
 &=\lambda_{\max}\!
 \left[(P_rQ_r^*Q_rP_r)^{-1/2}
       P_rK_{{\rm edge},r}P_r
       (P_rQ_r^*Q_rP_r)^{-1/2}\right],\\
 \kappa_{\partial,r}^{\sharp}
 &=\lambda_{\max}\!
 \left[(P_rQ_r^*Q_rP_r)^{-1/2}
       P_rK_{\partial,r}P_r
       (P_rQ_r^*Q_rP_r)^{-1/2}\right].
 \end{aligned}}
 \label{eq:YMA7-sharp-return-edges}
\end{equation}
All inverses are supported on \(\Ran P_r\).
\end{theorem}

\begin{proof}
Because \(P_r\) reduces \(A_r\) and
\(0\preceq P_rA_rP_r\preceq\rho_rP_r\),
\[
 P_rQ_r^*Q_rP_r=P_rD_rP_r-P_rA_rP_r,
\]
which proves \eqref{eq:YMA7-return-window}.  On the spectral support of
\(E_{r,\varepsilon}\),
\(
 I\preceq\varepsilon H_{\eta,r}^{-1}
\).  Hence
\[
 \begin{aligned}
 \norm{E_{r,\varepsilon}Q_rP_r}^2
 &\le\varepsilon\norm{P_rQ_r^*H_{\eta,r}^{-1}Q_rP_r}\\
 &\le\varepsilon\kappa_{{\rm edge},r}
      \norm{P_rQ_r^*Q_rP_r}
 \le\varepsilon\kappa_{{\rm edge},r}d_+.
 \end{aligned}
\]
Likewise
\(
 I-\Pi_{r,R}\preceq R^{-2}R_{\partial,r}^2
\), so
\[
 \norm{(I-\Pi_{r,R})Q_rP_r}^2
 \le R^{-2}\kappa_{\partial,r}d_+.
\]
This proves the two tails.  The lower endpoint in
\eqref{eq:YMA7-return-window} is positive, so the generalized eigenvalue
problems in \eqref{eq:YMA7-sharp-return-edges} are ordinary finite Hermitian
problems on \(\Ran P_r\), and their maxima are attained.
\end{proof}

\begin{corollary}[Return-survivor alternative on a persistent soft island]
\label{cor:YMA7-return-survivor}
Consider a cofinal family satisfying the direct window
\eqref{eq:YMA7-direct-window}, with separated low islands and the same-head
transport hypotheses of the shared survivor programme.
\begin{enumerate}[label=\textup{(R\arabic*)},leftmargin=4em]
\item If the sharp relative edges in
      \eqref{eq:YMA7-sharp-return-edges} are uniformly bounded and the
      corresponding joint physical screens are finite rank with summable
      transport tails, the return carriers \(Q_rP_r\) have a nonzero
      strongly convergent subsequence in the direct limit.
\item If the edge relative bound fails, a finite-card generalized
      eigenvector in the low island is the first near-unit-channel
      extremizer.  If the radius relative bound fails, an attained low-island
      polarization is the first boundary-range extremizer.
\item Cofinal divergence of either moment must still be separated into
      order-one escaping stock and high-cost dust by the shared survivor
      theorem.  Moment divergence alone is not a positive-range physical
      survivor.
\end{enumerate}
Under Mosco convergence, the compact branch produces a nonzero return carrier
attached to a limiting null polarization.  This is a classified failure
object, not a Yang--Mills mass-gap statement.
\end{corollary}

\begin{proof}
The tail bounds in \Cref{thm:YMA7-low-island-return} give uniform properness
of the restricted return carrier.  The imported low-island survivor theorem
then supplies the diagonal finite-head exhaustion and cofinal promotion under
the stated transport hypotheses.  Failure of either finite generalized edge
returns an attained eigenvector by finite-dimensional spectral calculus.  The
stock-versus-dust and Mosco conclusions are exactly the imported survivor
alternatives.
\end{proof}

\begin{assessment}[Economy relative to the Version V6 weighted theorem]
\label{ass:YMA7-low-island-economy}
The weighted Schur norm of \(\mathcal K_\eta\) remains a useful sufficient
all-source estimate.  It is no longer the first required acquisition after a
specific physical low island has been isolated.  The two matrices in
\eqref{eq:YMA7-two-return-moments}, compressed to that island and normalized
by its actual return Gram, are the target-minimal return tests.  Source
directions orthogonal to the failing island are not charged merely because a
global weighted norm is large.
\end{assessment}

% =============================================================================
\section{Integrated weak-facing and mode-card alternative}
\label{sec:YMA7-integrated}
% =============================================================================

\begin{theorem}[Boundary-to-mode completion and typed failure alternatives]
\label{thm:YMA7-integrated}
Fix the predeclared weak-facing selected-root sequence and the frozen first
axial mode.  After passage, when necessary, to a cofinal subsequence on
which the native source variance has a limit, the following ordered decision
tree is exhaustive.
\begin{enumerate}[label=\textup{(V7.\arabic*)},leftmargin=5.4em]
\item \textbf{Bounded root.}  A bounded cofinal subsequence carries the
      imported rough-Haar/contact trajectory obstruction.  The raw smooth
      local-vacuum branch stops there; no boundary source coefficient is
      promoted to a reduced-fibre mode margin.
\item \textbf{Escaping root with visibility collapse.}  The centre-fibre
      source vanishes and the total native source variance also tends to
      zero.  Source whitening is unavailable; this is a typed source-
      visibility failure.
\item \textbf{Escaping visible outer source, physical card unpopulated.}
      The normalized native line coalesces with its outer conditional line,
      but at least one of the direct mixed row \(y\), comparison source solve,
      comparison/return mode Grams, comparison head, canonical support, or
      Ward row is absent, or the outward assembled interval is still
      inconclusive before nonnegativity is certified.  The correct verdict is
      \(\stopstatus\) at the first missing or undecided row of the six-by-six
      incidence and four-direction mode card.
\item \textbf{Incidence and mode card populated with a separated verdict.}
      The identities \eqref{eq:YMA7-card-checks} give the exact
      boundary-source follower, the source-orthogonal score residual, and the
      comparison/return ancestry split.  The complete Version~V6 head and
      return are then reduced by their common gauge row and compared with the
      exact reference.  A cofinally uniform positive lower interval is
      \(\pass\); a negative upper interval on an explicit transverse vector is
      a physical pure-mode \(\obstruction\).
\item \textbf{Nonnegative soft island.}  If every assembled finite card is
      certified nonnegative but no cofinally uniform target margin closes and
      its normalized floor collapses, the return is tested by the two
      low-island relative edges
      \eqref{eq:YMA7-sharp-return-edges}.  Their bounded branch gives compact
      return transport; their failing branches return a near-unit-channel or
      boundary-range extremizer, followed by the imported stock-versus-dust
      classification.
\end{enumerate}
The selected radial four-leaf card of YM--B remains independent throughout.
It may classify spatial contraction, contact survival, or positive-range
survival only after its own source stock and physical-radius transport are
populated.  None of the five branches supplies an Osterwalder--Schrader time
gap or a continuum Yang--Mills mass.
\end{theorem}

\begin{proof}
The first two branches are the imported weak-facing alternative and source-
visibility dichotomy after the stated subsequence selection.  On the visible
branch, \Cref{thm:YMA7-native-outer} controls replacement of the normalized
native line by the outer line once the score ceiling is supplied, but it does
not construct the mixed row, the arm Grams, the comparison head, or their
support and Ward checks.  This gives branch~\textup{(V7.3)}.  The finite
incidence and common-source Ward identities are
\Cref{thm:YMA7-arm-incidence,thm:YMA7-Ward-assembly,thm:YMA7-minimal-card};
combining them with the Version~V6 common-gauge card gives
branch~\textup{(V7.4)}.  The remaining nonnegative collapsing case is
\Cref{thm:YMA7-low-island-return,cor:YMA7-return-survivor}.  Independence of
the radial card and all final firewalls are imported from YM--B and retained
in \Cref{fire:YMA7-no-substitution}.
\end{proof}

\begin{corollary}[Exact stopping semantics for the first axial card]
\label{cor:YMA7-termination}
After the newly attached interfaces are imposed, the first axial card has
only the following mathematical outcomes:
\begin{equation}
 \boxed{
 \pass
 \quad\big|\quad
 \obstruction
 \quad\big|\quad
 \stopstatus.}
 \label{eq:YMA7-three-outcomes}
\end{equation}
An ill-typed common-cylinder, failed Gram positivity, false support
projection, or violated Ward identity is rejected before this trichotomy.  A
failed sufficient global weighted norm is \(\stopstatus\) unless the
low-island generalized edge or the assembled physical matrix supplies an
actual witness.
\end{corollary}

\begin{proof}
The pass and obstruction branches are the sign alternatives for the assembled
physical Hermitian card.  Every missing physical value or inconclusive
one-sided estimate is a stop.  Domain and incidence failures do not define a
Hermitian card and therefore are input rejections rather than theorem
verdicts.
\end{proof}

% =============================================================================
\section{Version V7 theorem, ledger, and the next literal acquisition}
\label{sec:YMA7-terminal}
% =============================================================================

\begin{theorem}[Version V7 boundary-incidence and low-island return theorem]
\label{thm:YMA7-terminal}
Versions~V1--V6 are retained.  Version~V7 additionally proves:
\begin{enumerate}[label=\textup{(E7.\arabic*)},leftmargin=5.2em]
\item every centered scalar source on the selected reduced-fibre law has a
      canonical exact-form Witten lift;
\item the lift splits orthogonally into the positive comparison arm and the
      curvature-defect return arm of Version~V6;
\item the complete score covariance splits into the corresponding two mode
      Grams, with the return arm equal to
      \(\mathbb L_\eta^{\rm ret}\);
\item one occurring outer boundary source yields a positive six-by-six
      source/score Gram per colour on the first axial direction bank;
\item the assembled source-ancestry residual is the sum of the two arm
      residuals plus a positive router-dispersion Gram;
\item this dispersion is exactly the correction required for Ward
      cancellation when the comparison and return source overlaps cancel on
      a gauge direction;
\item one comparison solve for the boundary source and one direct two-replica
      source--score row determine the complete split incidence card;
\item the boundary programme's native-to-outer error transfers to the mixed
      mode row under a finite score ceiling;
\item an adjacent boundary-line overlap and the score-synthesis defect give
      an explicit transport estimate for the source--score incidence row;
\item two generalized relative edges on the collapsing physical island
      control the near-unit return-channel and boundary-range tails, replacing
      a global weighted-channel condition for that selected failure;
\item persistent failure is classified as a physical negative mode, source
      visibility collapse, near-unit return extremizer, boundary-range
      extremizer, compact Witten-current survivor, or unresolved physical
      acquisition, with the radial YM--B card retained independently.
\end{enumerate}
No item evaluates the actual six-by-six Gram, the first comparison head, a
selected-root radial coefficient, the generated/topological correction, or
the continuum mass gap.
\end{theorem}

\begin{proof}
Items~\textup{(E7.1)--(E7.3)} are
\Cref{lem:YMA7-exact-isometry,thm:YMA7-source-Pythagoras}.  Items
\textup{(E7.4)--(E7.6)} are
\Cref{thm:YMA7-arm-incidence,thm:YMA7-Ward-assembly}.  Item~\textup{(E7.7)}
is \Cref{thm:YMA7-minimal-card}.  Items~\textup{(E7.8)--(E7.9)} are
\Cref{thm:YMA7-native-outer,thm:YMA7-incidence-transport}.  Item
\textup{(E7.10)} is \Cref{thm:YMA7-low-island-return}.  The final
classification is
\Cref{cor:YMA7-return-survivor,thm:YMA7-integrated,cor:YMA7-termination}.
The limitations are the unpopulated entries in
\Cref{req:YMA7-first-card} and the retained boundary, trajectory, and mass
firewalls.
\end{proof}

\begingroup\small
\begin{longtable}{>{\raggedright\arraybackslash}p{0.27\textwidth}
                  >{\raggedright\arraybackslash}p{0.20\textwidth}
                  >{\raggedright\arraybackslash}p{0.45\textwidth}}
\caption{Version V7 proof-status ledger.}\label{tab:YMA7-ledger}\\
\toprule
\textbf{Row}&\textbf{Status}&\textbf{Exact scope}\\
\midrule
\endfirsthead
\toprule
\textbf{Row}&\textbf{Status}&\textbf{Exact scope}\\
\midrule
\endhead
Weak-facing boundary alternative&Imported exactly
&Bounded roots give the rough-Haar/contact trajectory obstruction; escaping
roots leave visibility collapse or an outer source line plus an independent
radial card.\\[0.3em]
Exact-form source lift&Proved exactly
&Every centered same-law source is isometric to a comparison arm plus a
curvature-defect return arm.\\[0.3em]
Score covariance arm split&Proved exactly
&The return arm is the V6 net return and the comparison arm is the positive
comparison Witten covariance.\\[0.3em]
Boundary-source/mode-score incidence&Reduced to a finite physical card
&One normalized outer source and four axial score columns give a six-by-six
Gram per colour.\\[0.3em]
Minimum card population&Proved as a reduction
&One comparison solve and one direct two-replica mixed row determine both
source arms and their cross rows.\\[0.3em]
Source ancestry&Proved exactly
&The complete source-orthogonal score Gram is
$C_A+C_R-y^*y$; its rank is the minimum missing response-source dimension.\\[0.3em]
Assembly dispersion&Proved exactly
&Separate comparison/return shorts miss one positive router-dispersion Gram.\\[0.3em]
Common-source Ward row&Proved exactly
&The two arm overlaps cancel on gauge directions, and the dispersion restores
the correct assembled residual.\\[0.3em]
Native-to-outer mode handoff&Quantitative conditional theorem
&The mixed row error is at most the score norm times the square root of the
centre-fibre fraction.\\[0.3em]
Adjacent-cutoff incidence transport&Quantitative conditional theorem
&Boundary-line overlap plus score transport controls the mixed source--score
row; split-arm transport remains separate.\\[0.3em]
Low-island return tails&Proved conditionally from two finite edges
&Near-unit channel and physical-radius tails are controlled only on the
collapsing physical island.\\[0.3em]
Actual outer source/score row&Open physical value
&No selected common-cylinder expectation or deterministic enclosure is
invented.\\[0.3em]
Actual comparison and return mode Grams&Open physical values
&The Version V6 first-head and return panels remain to be populated.\\[0.3em]
YM--B radial four-leaf card&Independent open physical value
&It is not used as a mode covariance, return edge, or reference coefficient.\\[0.3em]
Selected shell margin&\stopstatus
&No assembled negative mode is obtained; the six-by-six incidence and
four-direction head/return matrices remain unpopulated.\\[0.3em]
Higher jets, local-vacuum contraction, mass&Not inferred
&The downstream handoff remains closed until the actual quadratic card passes.\\
\bottomrule
\end{longtable}
\endgroup

The accompanying exact replay files are
\begin{center}
\path{Renewal_YMA_V7_boundary_incidence_replay.py},\\
\path{Renewal_YMA_V7_boundary_incidence_replay.json}.
\end{center}
The script SHA--256 digest is
\begin{center}\scriptsize\ttfamily
5d903a3aa61b63afe413ed3598223dfbd76ede1138a184d8fa3bfb956abac265.
\end{center}
It verifies, in exact symbolic arithmetic, one rational Woodbury/source
Pythagoras model, the arm-resolved dispersion identity, the two-card return-
ancestry counterexample, an adjacent boundary-overlap transport instance,
and both low-island tail inequalities.  It does not evaluate a Wilson
integral, solve the selected physical Witten equation, or populate a mass
margin.

\begin{assessment}[The actual V7 advance]
\label{ass:YMA7-progress}
Version~V6 isolated the net curvature-defect return but still treated its
source entrance only through the four coarse score columns.  The new boundary
and survivor notes show that this is insufficient for the required source
ancestry and cutoff handoff.  Version~V7 supplies the missing exact interface:
once the outer boundary writer occurs on the same reduced-fibre law, its
canonical exact-form lift resolves into comparison and return arms, one
six-by-six Gram decides which part of the score is its follower, and the Ward
correction generated by assembling those arms is explicit.  The global
weighted return condition is also reduced to two generalized moments on the
actual soft island.  The remaining work is now literal population of this
finite card, not another source, shell, or survivor theorem.
\end{assessment}

\begin{openproblem}[V8: populate the boundary-incidence and return-renormalized mode card]
\label{op:YMA7-V8}
Preserve V7 and the frozen first axial mode.  Proceed only through the
following physical acquisitions.
\begin{enumerate}[label=\textup{(V8.\arabic*)},leftmargin=5.2em]
\item Apply the boundary weak-facing alternative.  On a bounded branch export
      the rough-Haar/contact trajectory obstruction and stop.  On an escaping
      branch certify a positive native visibility floor and the centre-fibre
      error before normalizing the outer source.
\item Place that outer source \(m\) and the four mode scores on one actual
      response-matched cylinder.  Acquire the direct two-replica row \(y\) and
      verify \(y\,d(q)=0\).
\item Solve once for \(\eta_m=\mathcal A_\eta^{-1}dm\).  Combine its
      \(a,x_A\) rows with the Version~V6 comparison and return matrices to
      construct the full split Gram, \(x_R\), the route dispersion, and the
      total and return ancestry residuals.
\item Reject any card which fails positivity, support, source normalization,
      or the exact identities \eqref{eq:YMA7-card-checks}.  On a passing card,
      perform the Version~V6 common-gauge reduction and compare the assembled
      physical matrix with the exact flat reference.
\item If the assembled matrix is nonnegative but has a collapsing normalized
      low island, acquire the two sharp relative edges in
      \eqref{eq:YMA7-sharp-return-edges}.  Use their attained vectors to
      return either a compact return tail, a near-unit-channel extremizer, or
      a boundary-range extremizer.
\item For adjacent cutoffs, acquire the boundary-line mixed coefficient and
      the score-synthesis defect in the same rectangle.  Transport the split
      arms only after their comparison and return Witten operators have also
      been identified.
\item Run the selected radial four-leaf card independently.  It is inserted
      only into the boundary contraction/contact/positive-range trichotomy,
      never into the shell Hessian.
\item Insert generated and topological corrections only after the pure local
      incidence card closes.  Evaluate the finite two-loop coefficient only
      if it remains the first unresolved coefficient of this same assembled
      interval.
\end{enumerate}
If no physical observation, deterministic integral enclosure, or validated
finite solve is supplied, the exact verdict remains \(\stopstatus\) at the
first named row.  No general shell or boundary theorem is reopened.
\end{openproblem}

\begin{assessment}[V7 termination]
\label{ass:YMA7-termination}
The exact-form lift, comparison/return source Pythagoras, arm-resolved
incidence, common-source Ward correction, target-minimal six-by-six card,
cutoff handoff, and low-island return-tail theorem are \(\pass\) at their
stated mathematical scopes.  The complete selected response-matched symbol
remains \(\stopstatus\): the outer source/mode mixed row, comparison head,
return matrix, and selected radial physical values have not been populated.
No physical negative mode or continuum Yang--Mills mass has been produced.
\end{assessment}

\section*{Version V7 append-only continuation record}
\addcontentsline{toc}{section}{Version V7 append-only continuation record}
The complete Version~V6 source above is preserved byte for byte up to its
sole final \verb|\end{document}| command.  Its SHA--256 digest is
\begin{center}\scriptsize\ttfamily
6f5ab78e172514f3b2653cd0f20ceea3695ddc1cf0a54f54cc9c57095c0395af.
\end{center}
For Version~V8, preserve this accumulated source, append new material
immediately before the final document-closing command, and increment the
filename to \texttt{\_V8.tex}.  Keep the exact block map, parity-tree
section, physical response, first axial mode, source normalization,
local/topological support convention, curvature-defect split, boundary source
line, and comparison/return incidence convention fixed before reading any
mixed, return, radial, or symbol value.

% =============================================================================
% END APPENDED VERSION V7 MATERIAL
% =============================================================================



% =============================================================================
% START APPENDED VERSION V8 MATERIAL
% =============================================================================

\clearpage
\hypersetup{
 pdftitle={Renewal Yang--Mills Reduced-Fibre Wilson Shell Symbol and Local-Vacuum Programme V8},
 pdfsubject={The unnormalized partition-curvature germ, an exact symmetric secant decomposition, and the target-minimal first axial card}
}
\pdfinfo{
 /Title (Renewal Yang--Mills Reduced-Fibre Wilson Shell Symbol and Local-Vacuum Programme V8)
 /Subject (The unnormalized partition-curvature germ, an exact symmetric secant decomposition, and the target-minimal first axial card)
}

\section{Version V8: the unnormalized partition-curvature germ is upstream}
\label{sec:YMA8-direction}

Version~V7 is preserved above.  It constructed the exact common-history
boundary-source/Witten-current incidence needed to determine source ancestry,
route dispersion, and low-island return tails.  Those results remain valid.
They do not, however, populate the signed reduced-fibre Hessian.  The present
continuation moves one arrow upstream and identifies the minimum physical datum
which determines the first axial symbol itself.

Fix one finite selected shell, its exact dyadic block map, the relative
parity-tree section, the rooted Hodge carrier, the response-matched marginal
coefficient, and one nonzero first axial momentum before any value is read.
Let
\[
 v_*
\]
be the predeclared real unit cosine polarization in direction~$2$ and one unit
colour direction; at a self-conjugate momentum use the corresponding real
character.  Let $W_t$ be the product-exponential coarse-field curve with
$W_0=I$ and tangent $v_*$.  The marginal response coefficient is frozen before
this curve is formed.  If the intended physical response instead requires a
$t$-dependent retuning, that retuned coefficient and all of its derivatives
must be included in the action family below; it may not be re-solved silently
at the displaced points.

Let $Y_M^{\rm red}$ be the already fixed compact reduced fibre with normalized
product Haar measure $m_M$.  Write
\begin{equation}
 Z_M(t)=\int_{Y_M^{\rm red}}e^{-F_M(z,t)}\,dm_M(z),
 \qquad
 \mathcal S_M^{\rm full}(t)=-\log Z_M(t),
 \label{eq:YMA8-partition-action}
\end{equation}
where $F_M$ is the \emph{complete submitted shell action} on the fixed fibre:
pure Wilson, generated, marginally retuned, and every other term which belongs
to the declared full finite-torus cylinder.  Its normalized law is
\begin{equation}
 d\nu_{M,t}(z)=Z_M(t)^{-1}e^{-F_M(z,t)}\,dm_M(z).
 \label{eq:YMA8-normalized-law}
\end{equation}
At the central background the integrated full action is critical by the frozen
reflection and global-colour symmetries.  Thus
\begin{equation}
 H_M^{\rm full}[v_*,v_*]
 =\left.\frac{d^2}{dt^2}\right|_{t=0}\mathcal S_M^{\rm full}(t)
 \label{eq:YMA8-full-Hessian-direction}
\end{equation}
is the geometric Hessian in this direction, not a coordinate-acceleration
correction.

If the inherited support theorem supplies an exact effective-action split
\begin{equation}
 \mathcal S_M^{\rm full}
 =\mathcal S_M^{\rm loc}+\mathcal S_M^{\rm top},
 \label{eq:YMA8-effective-support-split}
\end{equation}
then the target local coefficient is
\begin{equation}
 \boxed{
 h_M^{\rm loc}(q)
 =H_M^{\rm full}[v_*,v_*]-t_M(q),
 \qquad
 t_M(q):=H_M^{\rm top}[v_*,v_*].}
 \label{eq:YMA8-local-from-full}
\end{equation}
The topological term is therefore retained and subtracted once.  The Ward row
sum of the local contractible packet is invoked only after this subtraction.

\begin{firewall}[No repetition of the existing score and ancestry compilers]
\label{fire:YMA8-no-redo}
The exact three-query identity
$H=\mathbb EK-\operatorname{Cov}(j)$, the Witten representation of the
covariance, the comparison/return decomposition, and the Version~V7
boundary-source incidence are inherited.  The new issue is logically prior:
all of those normalized-law objects can be identical while the direct
partition curvature, and therefore the physical symbol, is different.  The
proof below does not construct another boundary, Witten, graph, or shell
carrier.  It constructs the one scalar unnormalized germ required by the
already frozen first axial block.
\end{firewall}

\begin{theorem}[Upstream order of the first axial physical verdict]
\label{thm:YMA8-upstream-order}
Assume the central-axial symmetry audit and the exact local/topological support
convention.  The first physical quadratic verdict on the retained first axial
block is the sign of
\begin{equation}
 \boxed{h_M^{\rm loc}(q)-\delta\,s_M(q),}
 \label{eq:YMA8-first-physical-verdict}
\end{equation}
where $s_M(q)>0$ is the already evaluated physical reference coefficient
($s_M(q)=\beta s_0(q)$ on the pure-Wilson normalization of Version~V2).
This number is determined by the unnormalized partition germ
\eqref{eq:YMA8-partition-action} together with the one topological subtraction
in \eqref{eq:YMA8-local-from-full}.  A boundary-source Gram, current-ancestry
short, or radial message is not a prerequisite for defining it.  Those packets
become relevant only to attribute or transport an already acquired direct
card, or when their independent boundary target is being run.
\end{theorem}

\begin{proof}
The exact integrated Hessian is the second derivative of
$-\log Z_M$.  Equation~\eqref{eq:YMA8-local-from-full} then gives the local
coefficient.  On the passing central-axial audit the integrated local Hessian
and its reference are scalar on the whole transverse-adjoint orbit, so the
predeclared polarization evaluates the complete block.  The countertheorem in
the next section proves that normalized laws, their centered scores, Witten
currents, and every Gram formed from them do not determine this second
derivative.  Hence such downstream packets cannot precede the partition germ
in the logical order.
\end{proof}

% =============================================================================
\section{The coarse action-curvature gauge and the exact nonidentifiability theorem}
\label{sec:YMA8-curvature-gauge}
% =============================================================================

The distinction between a normalized score experiment and an unnormalized
action experiment is not a matter of numerical conditioning.  It is an exact
finite-dimensional nonidentifiability.

Let $E$ be any finite real coarse coefficient space and let
$F(z,v)$ be a smooth action on one fixed compact fibre, with
\[
 Z(v)=\int e^{-F(z,v)}\,dm(z),
 \qquad
 \mathcal S(v)=-\log Z(v),
 \qquad
 d\nu_v=Z(v)^{-1}e^{-F(z,v)}dm(z).
\]
For $h\in E$, denote the centered action score by
\begin{equation}
 f_{v,h}(z)
 =D_vF(z,v)[h]-\mathbb E_{\nu_v}D_vF(\cdot,v)[h].
 \label{eq:YMA8-centered-action-score}
\end{equation}
\begin{definition}[Coarse action-curvature gauge]
\label{def:YMA8-curvature-gauge}
For a smooth scalar $c:E\to\mathbb R$, independent of the fibre variable,
put
\begin{equation}
 F^c(z,v)=F(z,v)+c(v).
 \label{eq:YMA8-curvature-gauge-action}
\end{equation}
The transformation is called a \emph{coarse action-curvature gauge}.  It is
not declared physically redundant: it changes the unnormalized cylinder.
The name records only that normalized conditional laws cannot see it.
\end{definition}

\begin{countertheorem}[Normalized source geometry does not determine the physical Hessian]
\label{cth:YMA8-normalized-no-Hessian}
Under \eqref{eq:YMA8-curvature-gauge-action} the following objects are
unchanged for every $v$:
\begin{enumerate}[label=\textup{(N\arabic*)},leftmargin=4.2em]
\item the normalized law $\nu_v$ and all of its conditional expectations;
\item every centered score $f_{v,h}$ and every score covariance;
\item the fibre gradients $d_zf_{v,h}$, the scalar and one-form Witten
      operators, their minimum currents, and every comparison/return Gram
      constructed from them;
\item every Gram involving a fixed normalized boundary writer and the
      preceding centered score or current columns.
\end{enumerate}
In contrast,
\begin{equation}
 \boxed{
 \mathcal S^c(v)=\mathcal S(v)+c(v),
 \qquad
 \operatorname{Hess}\mathcal S^c(v)
 =\operatorname{Hess}\mathcal S(v)+\operatorname{Hess}c(v).}
 \label{eq:YMA8-Hessian-gauge-shift}
\end{equation}
Consequently, for every self-adjoint $B:E\to E$, the choice
$c(v)=\frac12\langle v,Bv\rangle$ preserves the complete normalized source,
Witten, return, and boundary-incidence packet but shifts the physical shell
Hessian by $B$.  No packet formed solely from those normalized objects can
determine the sign of the selected-shell margin.
\end{countertheorem}

\begin{proof}
The scalar factor $e^{-c(v)}$ cancels between the numerator and denominator of
$\nu_v$.  Moreover,
\[
 D_vF^c(z,v)[h]
 =D_vF(z,v)[h]+Dc(v)[h],
\]
and the second term is constant in $z$, so conditional centering removes it.
Its fibre differential is also zero.  Hence the weighted fibre operators,
centered score columns, exact Witten solutions, and every Gram built from them
are unchanged.  The partition function transforms as
$Z^c(v)=e^{-c(v)}Z(v)$, which gives
\eqref{eq:YMA8-Hessian-gauge-shift}.  The quadratic choice of $c$ realizes an
arbitrary $B$.
\end{proof}

\begin{corollary}[Zero source-ancestry residual is not a symbol margin]
\label{cor:YMA8-zero-ancestry-not-margin}
Even if the complete Version~V7 source/current ancestry residual vanishes,
the comparison and return routes satisfy every Ward and positivity identity,
and the entire connected covariance is a deterministic follower of the
occurring boundary source, the sign of
$H_M^{\rm loc}-\delta S_M$ remains undetermined until the coarse
partition-curvature gauge is fixed by an unnormalized action germ.
\end{corollary}

\begin{proof}
Apply \Cref{cth:YMA8-normalized-no-Hessian} with $B$ positive or negative on
the retained transverse block.  All stated normalized packets remain fixed,
while the physical difference crosses the target threshold.
\end{proof}

\begin{assessment}[What the upstream failure means physically]
\label{ass:YMA8-curvature-meaning}
The freedom in \Cref{cth:YMA8-normalized-no-Hessian} is not a licence to add a
counterterm after seeing the answer.  The actual Wilson/generated cylinder
fixes one member of the family.  The theorem says that this member must be
read through its unnormalized relative action or an equivalent deterministic
second-action writer.  Normalized sampling, score whitening, source ancestry,
and boundary overlap cannot reconstruct that missing datum.
\end{assessment}

% =============================================================================
\section{Exact symmetric partition secant and its direct--Hellinger decomposition}
\label{sec:YMA8-secant}
% =============================================================================

We now construct the target-minimal unnormalized Read.  Suppress $M$ for this
section and write
\[
 F_0=F(\cdot,0),
 \qquad
 F_+=F(\cdot,\varepsilon),
 \qquad
 F_-=F(\cdot,-\varepsilon),
 \qquad \varepsilon>0.
\]
All three actions are evaluated on the same compact fibre, with the same Haar
measure, response coefficient, source normalization, and support convention.

\begin{definition}[Symmetric partition curvature]
\label{def:YMA8-partition-curvature}
Define
\begin{equation}
 \boxed{
 Q_\varepsilon
 =\frac{\mathcal S(\varepsilon)+\mathcal S(-\varepsilon)
        -2\mathcal S(0)}{\varepsilon^2}
 =-\frac1{\varepsilon^2}
   \log\frac{Z(\varepsilon)Z(-\varepsilon)}{Z(0)^2}.}
 \label{eq:YMA8-Qepsilon}
\end{equation}
The quotient uses only relative partition values.  An unknown common positive
gain multiplying all three unnormalized cylinders cancels.
\end{definition}

Let
\begin{align}
 A_\varepsilon
 &:=\mathbb E_{\nu_0}
 \exp\left[-\frac12\bigl(F_++F_--2F_0\bigr)\right],
 \label{eq:YMA8-direct-midpoint-factor}\\
 B_\varepsilon
 &:=\int
 \sqrt{\frac{d\nu_+}{dm}\frac{d\nu_-}{dm}}\,dm,
 \label{eq:YMA8-Hellinger-affinity}
\end{align}
and define
\begin{equation}
 D_\varepsilon=-\frac2{\varepsilon^2}\log A_\varepsilon,
 \qquad
 I_\varepsilon=-\frac2{\varepsilon^2}\log B_\varepsilon.
 \label{eq:YMA8-DI-secant}
\end{equation}
The number $B_\varepsilon$ is the Hellinger affinity of the two normalized
neighbouring laws.  It lies in $(0,1]$, so $I_\varepsilon\ge0$.

\begin{theorem}[Exact finite-secant direct--Hellinger identity]
\label{thm:YMA8-secant-Pythagoras}
For every finite positive card,
\begin{equation}
 \boxed{
 Q_\varepsilon=D_\varepsilon-I_\varepsilon.}
 \label{eq:YMA8-secant-D-minus-I}
\end{equation}
Equivalently,
\begin{equation}
 \boxed{
 \frac{\sqrt{Z(\varepsilon)Z(-\varepsilon)}}{Z(0)}
 =\frac{A_\varepsilon}{B_\varepsilon}.}
 \label{eq:YMA8-affinity-partition-identity}
\end{equation}
If $F(z,t)$ is $C^2$ jointly on the compact fibre, then as
$\varepsilon\downarrow0$,
\begin{equation}
 \boxed{
 \begin{aligned}
 D_\varepsilon&\longrightarrow
   \mathbb E_{\nu_0}\,\partial_t^2F(\cdot,0),\\
 I_\varepsilon&\longrightarrow
   \operatorname{Var}_{\nu_0}(\partial_tF(\cdot,0)),\\
 Q_\varepsilon&\longrightarrow
   \mathcal S''(0)
   =\mathbb E_{\nu_0}\partial_t^2F
    -\operatorname{Var}_{\nu_0}(\partial_tF).
 \end{aligned}}
 \label{eq:YMA8-secant-limits}
\end{equation}
Thus \eqref{eq:YMA8-secant-D-minus-I} is the exact finite-displacement
ancestor of the inherited deterministic-Hessian minus connected-covariance
card.
\end{theorem}

\begin{proof}
Direct substitution gives
\[
 A_\varepsilon
 =\frac1{Z(0)}\int e^{-(F_++F_-)/2}\,dm,
\]
whereas
\[
 B_\varepsilon
 =\frac1{\sqrt{Z(\varepsilon)Z(-\varepsilon)}}
   \int e^{-(F_++F_-)/2}\,dm.
\]
Their quotient is \eqref{eq:YMA8-affinity-partition-identity}; taking
$-2\varepsilon^{-2}$ times its logarithm proves
\eqref{eq:YMA8-secant-D-minus-I}.  Positivity of the densities gives
$B_\varepsilon>0$, and Cauchy--Schwarz gives $B_\varepsilon\le1$.

For the limits, compactness and $C^2$ regularity give the uniform expansion
\[
 F_++F_--2F_0
 =\varepsilon^2\partial_t^2F(\cdot,0)+o(\varepsilon^2).
\]
Expanding the logarithm of the bounded positive expectation proves the first
limit.  Differentiation of $-\log Z(t)$ under the finite integral gives
\[
 \mathcal S''(0)
 =\mathbb E\partial_t^2F
  -\operatorname{Var}(\partial_tF).
\]
The third limit is ordinary symmetric differentiation, and the second then
follows from \eqref{eq:YMA8-secant-D-minus-I}.
\end{proof}

\begin{corollary}[The normalized law determines only the Hellinger debit]
\label{cor:YMA8-normalized-only-I}
The family of normalized laws $(\nu_t)$ determines $B_\varepsilon$ and hence
$I_\varepsilon$.  It does not determine $A_\varepsilon$, $D_\varepsilon$, or
$Q_\varepsilon$.  Under the coarse action-curvature gauge $F^c=F+c(t)$,
\begin{equation}
 \boxed{
 \begin{aligned}
 B_\varepsilon^c&=B_\varepsilon,\\
 D_\varepsilon^c
 &=D_\varepsilon+
   \frac{c(\varepsilon)+c(-\varepsilon)-2c(0)}{\varepsilon^2},\\
 Q_\varepsilon^c
 &=Q_\varepsilon+
   \frac{c(\varepsilon)+c(-\varepsilon)-2c(0)}{\varepsilon^2}.
 \end{aligned}}
 \label{eq:YMA8-secant-gauge-shift}
\end{equation}
\end{corollary}

\begin{proof}
The normalized laws are unchanged.  In
\eqref{eq:YMA8-direct-midpoint-factor}, the integrand acquires the constant
factor
$\exp[-(c(\varepsilon)+c(-\varepsilon)-2c(0))/2]$.
Taking logarithms gives the remaining identities.
\end{proof}

\begin{remark}[Why the full action is assembled before the positive split]
The physical sign may be read directly from $Q_\varepsilon$ without measuring
$D_\varepsilon$ and $I_\varepsilon$ separately.  The decomposition is useful
for diagnosis, but separately enclosing its two potentially large terms can
lose their cancellation.  The target-minimal writer below acquires their
assembled ratio first.
\end{remark}

% =============================================================================
\section{One positive two-replica writer for the assembled partition ratio}
\label{sec:YMA8-two-replica}
% =============================================================================

Define positive same-fibre likelihood writers
\begin{equation}
 L_\pm(z)
 :=\exp[-(F_\pm(z)-F_0(z))].
 \label{eq:YMA8-Lpm}
\end{equation}
They are not normalized likelihood ratios: their means carry the relative
partition values.

\begin{theorem}[One two-replica acquisition of the complete symmetric secant]
\label{thm:YMA8-two-replica-partition}
Let $z_1,z_2$ be independent samples from the single base law $\nu_0$.  Then
\begin{equation}
 \boxed{
 \mathscr R_\varepsilon
 :=\mathbb E_{\nu_0}^{\otimes2}
   [L_+(z_1)L_-(z_2)]
 =\frac{Z(\varepsilon)Z(-\varepsilon)}{Z(0)^2}>0,}
 \label{eq:YMA8-R-two-replica}
\end{equation}
and therefore
\begin{equation}
 \boxed{Q_\varepsilon=-\varepsilon^{-2}\log\mathscr R_\varepsilon.}
 \label{eq:YMA8-Q-from-R}
\end{equation}
This one positive expectation contains the direct second-action term, the
complete covariance, every Wilson--generated score cross term, and their
finite-displacement nonlinear corrections with their actual signs.  No
boundary-source column or Witten-current decomposition is required to define
it.
\end{theorem}

\begin{proof}
The base-law change-of-measure identity gives
\[
 \mathbb E_{\nu_0}L_\pm
 =\frac1{Z(0)}\int e^{-F_\pm}\,dm
 =\frac{Z(\pm\varepsilon)}{Z(0)}.
\]
Independence multiplies the two means, proving
\eqref{eq:YMA8-R-two-replica}; \eqref{eq:YMA8-Q-from-R} is
\eqref{eq:YMA8-Qepsilon}.
\end{proof}

\begin{corollary}[A common gain is harmless; a source-dependent gain is physical]
\label{cor:YMA8-gain-scope}
Multiplication of every unnormalized cylinder by one common unknown constant
leaves $L_\pm$, $\mathscr R_\varepsilon$, and $Q_\varepsilon$ unchanged.
Multiplication by a $t$-dependent common factor changes
$Q_\varepsilon$ by its symmetric logarithmic curvature.  Thus the required
absolute datum is a \emph{relative partition germ}, not an independently
calibrated value of $Z(0)$; nevertheless it cannot be reconstructed from
normalized neighbouring laws.
\end{corollary}

\begin{proof}
A $t$-independent factor cancels in every ratio.  A $t$-dependent factor is
exactly the transformation in
\Cref{cor:YMA8-normalized-only-I}, with the sign determined by whether it is
written in the weight or in the action.
\end{proof}

\begin{firewall}[Common base law and common action are mandatory]
\label{fire:YMA8-common-cylinder}
The two means in \eqref{eq:YMA8-R-two-replica} must use the same base law,
coarse source normalization, fibre, Haar measure, response coefficient, and
action convention.  Multiplying estimates from separately retuned or
separately normalized experiments does not form the physical partition ratio.
A positive-extension or same-history occurrence check is an input-domain
audit; failure rejects the card and is not a negative Yang--Mills mode.
\end{firewall}

% =============================================================================
\section{Rigorous curvature enclosure from a finite displacement}
\label{sec:YMA8-curvature-enclosure}
% =============================================================================

The symmetric secant is an exact finite quantity.  A fourth-derivative stock
turns it into a rigorous interval for the infinitesimal Hessian.

\begin{theorem}[Sharp central-second-difference remainder]
\label{thm:YMA8-central-difference}
Let $\mathcal S\in C^4([ -\varepsilon,\varepsilon])$ and set
\[
 M_4(\varepsilon)
 :=\sup_{|t|\le\varepsilon}|\mathcal S^{(4)}(t)|.
\]
Then
\begin{equation}
 \boxed{
 \left|Q_\varepsilon-\mathcal S''(0)\right|
 \le\frac{\varepsilon^2}{12}M_4(\varepsilon).}
 \label{eq:YMA8-central-error}
\end{equation}
The constant $1/12$ is sharp, as shown by a pure quartic perturbation.
\end{theorem}

\begin{proof}
Taylor's formula with integral remainder, applied at $\varepsilon$ and
$-\varepsilon$, cancels the constant, linear, and cubic terms and gives
\[
 \mathcal S(\varepsilon)+\mathcal S(-\varepsilon)-2\mathcal S(0)
 =\varepsilon^2\mathcal S''(0)+R,
 \qquad |R|\le\frac{\varepsilon^4}{12}M_4(\varepsilon).
\]
Divide by $\varepsilon^2$.  For
$\mathcal S(t)=at^2/2+bt^4/24$, equality holds with
$|b|=M_4$.
\end{proof}

The derivative stock can itself be bounded directly from the finite action
writers.  For $j=1,2,3,4$ put
\begin{equation}
 X_j(t,z)=\partial_t^jF(z,t),
 \qquad
 B_j(\varepsilon)=
 \sup_{|t|\le\varepsilon}\sup_z|X_j(t,z)|.
 \label{eq:YMA8-action-derivative-bounds}
\end{equation}
All suprema are finite on a finite compact card.  They are physical action
bounds, not consequences of the normalized score law.

\begin{proposition}[Exact fourth derivative and an executable action bound]
\label{prop:YMA8-fourth-derivative}
Let $\kappa_{\nu_t}$ denote joint cumulants under $\nu_t$.  Then
\begin{equation}
 \boxed{
 \begin{aligned}
 \mathcal S^{(4)}(t)
 ={}&\mathbb E X_4
 -4\kappa(X_1,X_3)-3\kappa(X_2,X_2)\\
 &+6\kappa(X_1,X_1,X_2)-\kappa(X_1,X_1,X_1,X_1).
 \end{aligned}}
 \label{eq:YMA8-S-fourth-cumulant}
\end{equation}
Consequently
\begin{equation}
 \boxed{
 \begin{aligned}
 M_4(\varepsilon)
 \le{}&B_4+8B_1B_3+6B_2^2
       +36B_1^2B_2+26B_1^4,
 \end{aligned}}
 \label{eq:YMA8-M4-bound}
\end{equation}
where the $B_j$ are evaluated at $\varepsilon$.
\end{proposition}

\begin{proof}
For a $t$-dependent observable $G_t$,
\[
 \frac d{dt}\mathbb E_{\nu_t}G_t
 =\mathbb E G_t'-\kappa(G_t,X_1).
\]
More generally, differentiation of a joint cumulant differentiates each entry
and subtracts the cumulant with one additional $X_1$ entry.  Starting from
$\mathcal S'=\mathbb EX_1$ gives successively
\[
 \mathcal S''=\mathbb EX_2-\kappa(X_1,X_1),
\]
\[
 \mathcal S'''=\mathbb EX_3-3\kappa(X_1,X_2)
 +\kappa(X_1,X_1,X_1),
\]
and then \eqref{eq:YMA8-S-fourth-cumulant}.

For bounded random variables $|Y_i|\le b_i$, the moment--cumulant formula
implies
\[
 |\kappa_2(Y_1,Y_2)|\le2b_1b_2,
 \quad
 |\kappa_3(Y_1,Y_2,Y_3)|\le6b_1b_2b_3,
\]
\[
 |\kappa_4(Y_1,Y_2,Y_3,Y_4)|\le26b_1b_2b_3b_4.
\]
Indeed, the sums of the absolute partition coefficients are respectively
$2$, $6$, and
$1+7+2\cdot6+6=26$.  Insert these bounds into
\eqref{eq:YMA8-S-fourth-cumulant}.
\end{proof}

\begin{lemma}[Safe Wilson insertion bound]
\label{lem:YMA8-Wilson-insertion-bound}
Let $U_0,\ldots,U_m\in\mathrm U(3)$ and let $X_1,\ldots,X_m$ be arbitrary
$3\times3$ matrices.  Then
\begin{equation}
 \boxed{
 \frac13\left|\Tr(U_0X_1U_1\cdots X_mU_m)\right|
 \le\prod_{j=1}^m\|X_j\|_{\rm op}
 \le\prod_{j=1}^m\|X_j\|_F.}
 \label{eq:YMA8-trace-word-bound}
\end{equation}
Therefore every $B_j$ for a finite Wilson/generated action is bounded by a
finite sum of explicitly enumerable insertion words and the corresponding
generated-action derivative intervals.  Equality in
\eqref{eq:YMA8-trace-word-bound} is not asserted merely from unit Frobenius
normalization of the colour generators.
\end{lemma}

\begin{proof}
The product inside the trace has operator norm at most the product of the
operator norms of the inserted matrices.  The inequality
$|\Tr A|\le3\|A\|_{\rm op}$ proves the first bound; the second is standard.
Differentiating a finite ordered plaquette word produces finitely many terms
of this form.  Summing their bounds and the submitted generated-action
bounds gives the last statement.
\end{proof}

\begin{theorem}[Validated interval for the full directional Hessian]
\label{thm:YMA8-full-interval}
Suppose a rigorous positive interval
\begin{equation}
 0<\underline R_\varepsilon
 \le\mathscr R_\varepsilon
 \le\overline R_\varepsilon<\infty
 \label{eq:YMA8-R-interval}
\end{equation}
is supplied for the two-replica writer, and let
$M_4(\varepsilon)$ be bounded by
\Cref{prop:YMA8-fourth-derivative} or by any sharper valid estimate.  Put
\begin{align}
 \underline H_M^{\rm full}
 &:=-\varepsilon^{-2}\log\overline R_\varepsilon
   -\frac{\varepsilon^2}{12}M_4(\varepsilon),
 \label{eq:YMA8-Hfull-lower}\\
 \overline H_M^{\rm full}
 &:=-\varepsilon^{-2}\log\underline R_\varepsilon
   +\frac{\varepsilon^2}{12}M_4(\varepsilon).
 \label{eq:YMA8-Hfull-upper}
\end{align}
Then
\begin{equation}
 \boxed{
 \underline H_M^{\rm full}
 \le H_M^{\rm full}[v_*,v_*]
 \le\overline H_M^{\rm full}.}
 \label{eq:YMA8-Hfull-interval}
\end{equation}
An enclosure whose lower endpoint for $\mathscr R_\varepsilon$ meets zero is
\stopstatus: the logarithmic card is not yet certified.  It is not a negative
physical mode.
\end{theorem}

\begin{proof}
Monotonicity of $-\log r$ on $(0,\infty)$ and
\eqref{eq:YMA8-Q-from-R} give
\[
 -\varepsilon^{-2}\log\overline R_\varepsilon
 \le Q_\varepsilon
 \le-\varepsilon^{-2}\log\underline R_\varepsilon.
\]
Add the two-sided truncation radius from
\Cref{thm:YMA8-central-difference}.
\end{proof}

\begin{remark}[No cutoff-uniform claim from a fixed-card supremum]
The bound \eqref{eq:YMA8-M4-bound} is deliberately finite-card and may grow
with the volume.  It is nevertheless a rigorous terminating acquisition for
the first selected card.  A cutoff-uniform theorem requires a separate local
connected estimate or summable scalar transport.  It is not inferred from
compactness of each individual fibre.
\end{remark}

% =============================================================================
\section{The local/topological subtraction and the first axial margin interval}
\label{sec:YMA8-local-card}
% =============================================================================

The partition writer in \Cref{thm:YMA8-two-replica-partition} evaluates the
full finite-torus action.  The local contractible card is obtained only after
the inherited topological sector is removed on the same polarization.

\begin{theorem}[Topological-safe local interval]
\label{thm:YMA8-local-interval}
Assume the full interval \eqref{eq:YMA8-Hfull-interval} and a certified
same-carrier topological Hessian interval
\begin{equation}
 \underline t_M(q)\le t_M(q)\le\overline t_M(q).
 \label{eq:YMA8-topological-interval}
\end{equation}
Then
\begin{equation}
 \boxed{
 \underline h_M^{\rm loc}(q)
 :=\underline H_M^{\rm full}-\overline t_M(q)
 \le h_M^{\rm loc}(q)
 \le
 \overline H_M^{\rm full}-\underline t_M(q)
 =:\overline h_M^{\rm loc}(q).}
 \label{eq:YMA8-local-interval}
\end{equation}
For an exact reference interval
$0<\underline s_M(q)\le s_M(q)\le\overline s_M(q)$ and
$\delta\ge0$, the physical margin has the sign-safe enclosure
\begin{equation}
 \boxed{
 \begin{aligned}
 \underline m_M(q)
 &=\underline h_M^{\rm loc}(q)-\delta\overline s_M(q),\\
 \overline m_M(q)
 &=\overline h_M^{\rm loc}(q)-\delta\underline s_M(q).
 \end{aligned}}
 \label{eq:YMA8-margin-interval}
\end{equation}
Thus $\underline m_M(q)>0$ is a finite \pass, while
$\overline m_M(q)<0$ returns the explicit polarization $v_*$ as a direct
\obstruction.  An interval meeting zero is \stopstatus.
\end{theorem}

\begin{proof}
Subtract the interval \eqref{eq:YMA8-topological-interval} from
\eqref{eq:YMA8-Hfull-interval}; subtraction reverses the topological endpoint
used on each side.  Since $\delta\ge0$ and the reference interval is positive,
the smallest margin uses its upper endpoint and the largest margin its lower
endpoint.  The sign conclusions are ordinary interval inclusion.
\end{proof}

\begin{corollary}[One real polarization closes the whole first axial block]
\label{cor:YMA8-one-polarization-block}
If the integrated local action, the topological split, and the retained
projection pass the exact central-axial symmetry audit, then
\begin{equation}
 H_M^{\rm loc}(q)=h_M^{\rm loc}(q)I_{24},
 \qquad
 S_M(q)=s_M(q)I_{24}.
 \label{eq:YMA8-physical-scalar-block}
\end{equation}
Consequently the single interval in
\eqref{eq:YMA8-margin-interval} is the complete first-orbit verdict.  If the
symmetry or support audit fails, the same computation is only a directional
card and is not promoted to the full block.
\end{corollary}

\begin{proof}
Apply the already proved integrated-action little-group and global-colour
Schur argument, retaining the Version~V2 qualification that the two raw terms
need not scalarize separately.  The direct partition secant approximates their
assembled physical difference, so no separate raw-jet symmetry is used.
\end{proof}

\begin{corollary}[Exact support-split secant, when available]
\label{cor:YMA8-support-split-secant}
If the effective split \eqref{eq:YMA8-effective-support-split} is evaluable at
$t=0,\pm\varepsilon$, define
\[
 Q_\varepsilon^{\rm top}
 =\frac{\mathcal S^{\rm top}(\varepsilon)
       +\mathcal S^{\rm top}(-\varepsilon)
       -2\mathcal S^{\rm top}(0)}{\varepsilon^2}.
\]
Then
\begin{equation}
 \boxed{
 Q_\varepsilon^{\rm loc}
 =Q_\varepsilon^{\rm full}-Q_\varepsilon^{\rm top}}
 \label{eq:YMA8-local-secant-exact}
\end{equation}
holds exactly before either secant is sent to zero.  A fourth-derivative bound
for $\mathcal S^{\rm loc}$ may then replace the separate full and topological
Hessian intervals.  This is a tightening of
\Cref{thm:YMA8-local-interval}, not permission to omit the topological row.
\end{corollary}

\begin{proof}
Symmetric second differences are linear in the effective action.  Apply them
to \eqref{eq:YMA8-effective-support-split}.
\end{proof}

\begin{firewall}[The full torus is not declared locally Ward zero]
\label{fire:YMA8-full-not-local-Ward}
A constant coarse one-form can carry wrapping or nonlinear commutator
response in the full finite-torus cylinder.  Neither the contact coefficient
nor the zero symbol of the full partition secant is set to zero by hand.  The
local Ward identity applies only after the support split has been physically
landed and the topological contribution has been removed.
\end{firewall}

% =============================================================================
\section{Target-native transport of the partition-curvature scalar}
\label{sec:YMA8-transport}
% =============================================================================

The intended output at this stage is one scalar or one finite mode matrix, not
transport of a complete sampling instrument.  The relative partition germ has
its own exact adjacent-cutoff identity.

Let $r$ and $r+1$ be adjacent selected cards.  Transport the predeclared
polarization by the already frozen coefficient identification and use one
common finite displacement $\varepsilon_r$ on the two cards.  Let
$\mathscr R_{r,\varepsilon_r}$ and
$\mathscr R_{r+1,\varepsilon_r}$ be their positive partition ratios.

\begin{theorem}[Exact adjacent-cutoff partition-curvature defect]
\label{thm:YMA8-partition-transport}
The full symmetric secants satisfy
\begin{equation}
 \boxed{
 Q_{r+1,\varepsilon_r}^{\rm full}
 -Q_{r,\varepsilon_r}^{\rm full}
 =-\varepsilon_r^{-2}
  \log\frac{\mathscr R_{r+1,\varepsilon_r}}
            {\mathscr R_{r,\varepsilon_r}}.}
 \label{eq:YMA8-partition-transport}
\end{equation}
If $M_{4,r}$ and $M_{4,r+1}$ bound the two fourth derivatives, then
\begin{equation}
 \boxed{
 \begin{aligned}
 |H_{r+1}^{\rm full}[v_{r+1},v_{r+1}]
   -H_r^{\rm full}[v_r,v_r]|
 \le{}&\varepsilon_r^{-2}
  \left|\log\frac{\mathscr R_{r+1,\varepsilon_r}}
                 {\mathscr R_{r,\varepsilon_r}}\right|\\
 &+\frac{\varepsilon_r^2}{12}
   (M_{4,r+1}+M_{4,r}).
 \end{aligned}}
 \label{eq:YMA8-Hessian-transport-bound}
\end{equation}
After adding a valid interval for the adjacent topological Hessian defect, the
same statement holds for the local coefficient.
\end{theorem}

\begin{proof}
Subtract the two identities \eqref{eq:YMA8-Q-from-R}.  Each exact Hessian
differs from its secant by at most the corresponding radius in
\Cref{thm:YMA8-central-difference}; the triangle inequality gives
\eqref{eq:YMA8-Hessian-transport-bound}.  Equation
\eqref{eq:YMA8-local-from-full} supplies the final subtraction.
\end{proof}

\begin{corollary}[Cofinal scalar landing without instrument convergence]
\label{cor:YMA8-scalar-cofinal}
Suppose along a selected cofinal sequence there are displacements
$\varepsilon_r>0$ and certified numbers $e_r$ which dominate the right side of
\eqref{eq:YMA8-Hessian-transport-bound} together with the local/topological
defect, and
\begin{equation}
 \sum_r e_r<\infty.
 \label{eq:YMA8-summable-scalar-defects}
\end{equation}
Then the transported local directional Hessians are Cauchy and have a unique
limit.  This scalar conclusion does not require convergence of the complete
boundary-source instrument, Witten-current dictionary, or path law.  Those
stronger transports remain necessary when their carriers, rather than the
selected Hessian value, are themselves the target.
\end{corollary}

\begin{proof}
Equation \eqref{eq:YMA8-Hessian-transport-bound} and the topological defect
bound give
$|h_{r+1}^{\rm loc}-h_r^{\rm loc}|\le e_r$.  Summability gives the Cauchy
property.
\end{proof}

\begin{countertheorem}[Score transport does not imply partition-curvature transport]
\label{cth:YMA8-score-transport-no-partition}
There are adjacent-card families with zero normalized score-loss Gram, exact
transport of every Witten and boundary-source incidence row, and arbitrary
prescribed partition-curvature defect.  Hence the scalar in
\eqref{eq:YMA8-partition-transport} is an independent adjacent-cutoff row.
\end{countertheorem}

\begin{proof}
Take one normalized family and multiply the fine unnormalized weights by
$e^{-c(v)}$, leaving the old card unchanged.  Choose $c(0)=Dc(0)=0$ and an
arbitrary Hessian on the transported coefficient space.  By
\Cref{cth:YMA8-normalized-no-Hessian}, every normalized score and incidence
row transports exactly, while the fine physical Hessian acquires that
arbitrary defect.
\end{proof}

% =============================================================================
\section{Finite executable card and the revised programme order}
\label{sec:YMA8-executable}
% =============================================================================

\begin{request}[Literal first selected-shell partition card]
\label{req:YMA8-first-card}
On the first selected nonzero axial mode, acquire the following data in this
order.
\begin{enumerate}[label=\textup{(A\arabic*)},leftmargin=4.2em]
\item Freeze the complete unnormalized action family $F_M(z,t)$ on the exact
      reduced fibre, including the response-matched generated terms.  Verify
      that no coefficient is refitted between $t=0$ and
      $t=\pm\varepsilon$.
\item Acquire one outward positive interval for the two-replica writer
      $\mathscr R_\varepsilon$ in
      \eqref{eq:YMA8-R-two-replica}.
\item Enumerate or otherwise certify the four action derivative stocks
      $B_1,\ldots,B_4$ on the occupied displacement interval, and form the
      truncation radius in \eqref{eq:YMA8-M4-bound}.  A sharper analytic or
      interval bound may replace it.
\item Acquire the topological Hessian or exact topological secant on the same
      polarization.  Do not infer it by dividing by a vanishing curl symbol.
\item Form \eqref{eq:YMA8-margin-interval}.  Use a second predeclared
      displacement, for example $\varepsilon/2$, as a held-out consistency
      card; it may tighten but does not fit the first interval.
\end{enumerate}
No boundary-source Gram is required unless the assembled interval is then to
be attributed or transported through the boundary programme.
\end{request}

\begin{theorem}[Terminating finite-card semantics]
\label{thm:YMA8-termination-alternative}
After the input-domain checks in
\Cref{req:YMA8-first-card}, the first axial card has exactly one of the
following mathematical verdicts.
\begin{enumerate}[label=\textup{(V8.\arabic*)},leftmargin=4.8em]
\item \pass: the lower endpoint in
      \eqref{eq:YMA8-margin-interval} is strictly positive.
\item \obstruction: the upper endpoint is strictly negative; the frozen real
      momentum--colour polarization is an explicit direct physical witness.
\item \stopstatus: the ratio interval touches zero, the fourth-derivative or
      topological row is absent, or the final interval meets zero.
\end{enumerate}
A failed same-history, support, response-freezing, or positivity check rejects
an ill-typed input before these three verdicts.  It is not a fourth
mathematical outcome and is not a Yang--Mills obstruction.
\end{theorem}

\begin{proof}
The interval in \eqref{eq:YMA8-margin-interval} contains the exact physical
margin.  Strict positivity or strict negativity therefore gives the first two
claims.  Every remaining valid interval meets the threshold or lacks a
required enclosing row, which is precisely the unresolved branch.
\end{proof}

\begin{theorem}[Revised YM--A order after moving upstream]
\label{thm:YMA8-revised-order}
The economical order of the programme is now:
\begin{equation}
 \boxed{
 \begin{gathered}
 \text{weak-facing bounded/escaping trajectory alternative}\\
 \Downarrow\\
 \text{complete unnormalized relative partition germ on the selected mode}\\
 \Downarrow\\
 \text{topological subtraction and direct physical margin interval}\\
 \Downarrow\\
 \begin{matrix}
 \text{if the direct card passes: higher jets and adjacent-cutoff transport},\\
 \text{if it fails or remains narrow-source unresolved: source/current attribution}
 \end{matrix}\\
 \Downarrow\\
 \text{local-vacuum contraction only after the quadratic packet passes.}
 \end{gathered}}
 \label{eq:YMA8-revised-order}
\end{equation}
The Version~V7 incidence card and the independent YM--B radial card are
retained at their proper downstream interfaces.  Neither is discarded, and
neither substitutes for the missing partition-curvature value.
\end{theorem}

\begin{proof}
The bounded/escaping branch is the inherited trajectory alternative.
\Cref{cth:YMA8-normalized-no-Hessian} places the unnormalized germ before all
normalized source and current cards.  The local/topological order is
\Cref{thm:YMA8-local-interval}.  The final two arrows are the standing
quadratic-before-higher-jet and spatial-channel-not-temporal-gap firewalls.
\end{proof}

\begin{assessment}[Why this is not another general shell theorem]
\label{ass:YMA8-not-general}
The output is one positive expectation, one finite derivative radius, one
already required support subtraction, and one scalar interval on the frozen
physical polarization.  It does not enlarge the shell Hilbert space or ask for
a full boundary, path, or state reconstruction.  Its countertheorem also
explains why the preceding normalized incidence programme could continue
indefinitely without deciding the physical sign: it was invariant under the
first unclosed action-curvature row.
\end{assessment}

% =============================================================================
\section{Version V8 terminal theorem, proof ledger, and next acquisition}
\label{sec:YMA8-terminal}
% =============================================================================

\begin{theorem}[Version V8 partition-curvature theorem]
\label{thm:YMA8-terminal}
Versions~V1--V7 are retained.  Version~V8 additionally proves:
\begin{enumerate}[label=\textup{(E8.\arabic*)},leftmargin=5.2em]
\item normalized reduced-fibre laws, centered scores, Witten currents,
      comparison/return Grams, and normalized boundary-source incidence are
      invariant under an arbitrary coarse action-curvature gauge;
\item the physical shell Hessian shifts by the arbitrary Hessian of that gauge,
      so the normalized packet cannot determine the symbol margin;
\item the symmetric partition curvature has the exact finite-secant
      decomposition $Q_\varepsilon=D_\varepsilon-I_\varepsilon$, where the
      second term is the nonnegative Hellinger debit and the first is the
      independent direct-action row;
\item one positive two-replica writer on the base physical cylinder acquires
      the complete assembled partition ratio and hence $Q_\varepsilon$;
\item a sharp fourth-derivative estimate turns that finite secant into a
      rigorous Hessian interval;
\item the fourth derivative has the explicit finite cumulant formula
      \eqref{eq:YMA8-S-fourth-cumulant} and the executable bound
      \eqref{eq:YMA8-M4-bound};
\item ordered Wilson insertions obey the safe trace inequality
      \eqref{eq:YMA8-trace-word-bound}; no unproved equality is used;
\item the full finite-torus interval becomes the local contractible interval
      only after the independently supported topological Hessian is subtracted;
\item under the physical axial audit, one real polarization gives the complete
      first transverse-adjoint block;
\item adjacent-cutoff transport of this scalar is the exact logarithmic ratio
      \eqref{eq:YMA8-partition-transport} plus the displayed finite-difference
      and topological defects.
\end{enumerate}
No item supplies the numerical partition ratio on the first selected shell,
proves weak-facing-root escape, closes the nonaxial orbit bank, controls higher
jets, or establishes a local-vacuum or continuum mass theorem.
\end{theorem}

\begin{proof}
Items~\textup{(E8.1)--(E8.2)} are
\Cref{cth:YMA8-normalized-no-Hessian}.  Item~\textup{(E8.3)} is
\Cref{thm:YMA8-secant-Pythagoras}.  Item~\textup{(E8.4)} is
\Cref{thm:YMA8-two-replica-partition}.  Items~\textup{(E8.5)--(E8.7)} are
\Cref{thm:YMA8-central-difference,prop:YMA8-fourth-derivative,lem:YMA8-Wilson-insertion-bound}.
Items~\textup{(E8.8)--(E8.9)} are
\Cref{thm:YMA8-local-interval,cor:YMA8-one-polarization-block}.  Item
\textup{(E8.10)} is \Cref{thm:YMA8-partition-transport}.
\end{proof}

\begingroup\small
\begin{longtable}{>{\raggedright\arraybackslash}p{0.27\textwidth}
                  >{\raggedright\arraybackslash}p{0.20\textwidth}
                  >{\raggedright\arraybackslash}p{0.45\textwidth}}
\caption{Version V8 proof-status ledger.}\label{tab:YMA8-ledger}\\
\toprule
\textbf{Row}&\textbf{Status}&\textbf{Exact scope}\\
\midrule
\endfirsthead
\toprule
\textbf{Row}&\textbf{Status}&\textbf{Exact scope}\\
\midrule
\endhead
Upstream action-curvature row&Proved necessary
&An arbitrary fibre-independent coarse quadratic action changes the Hessian
while preserving every normalized score, Witten, return, and boundary-source
card.\\[0.3em]
Symmetric partition secant&Proved exactly
&$Q_\varepsilon=-\varepsilon^{-2}\log[Z_+Z_-/Z_0^2]$ and
$Q_\varepsilon=D_\varepsilon-I_\varepsilon$.\\[0.3em]
Normalized Hellinger debit&Proved positive
&$I_\varepsilon=-2\varepsilon^{-2}\log B_\varepsilon\ge0$ and converges to
the complete score variance.\\[0.3em]
Direct action secant&Proved independent
&$D_\varepsilon$ converges to the deterministic second-action expectation and
carries the invisible action-curvature gauge.\\[0.3em]
Two-replica partition writer&Target-minimal exact card
&One positive base-law expectation gives the assembled ratio without
separating large direct and covariance terms.\\[0.3em]
Finite-displacement error&Proved sharply
&The Hessian error is at most $\varepsilon^2M_4/12$.\\[0.3em]
Fourth derivative stock&Exact formula and finite bound
&Four action derivatives and finite cumulant coefficients give
\eqref{eq:YMA8-M4-bound}; no cutoff-uniformity is inferred.\\[0.3em]
Wilson insertion control&Proved
&Trace words are bounded by operator/Frobenius norm products; equality is not
assumed from colour normalization.\\[0.3em]
Local/topological support&Kept exact
&The full partition curvature is not declared locally Ward zero; the supported
topological block is subtracted once.\\[0.3em]
First axial orbit&Finite terminating compiler
&After the physical symmetry audit, one real polarization gives PASS,
TYPED OBSTRUCTION, or STOP.\\[0.3em]
Scalar cutoff transport&Proved exactly
&Adjacent secants differ by one logarithmic partition-ratio rectangle; fourth-
derivative and topological errors are explicit.\\[0.3em]
Selected physical value&\stopstatus
&No outward interval for $\mathscr R_\varepsilon$ or the same-card topological
coefficient has yet been supplied.\\[0.3em]
Boundary ancestry and radial card&Retained downstream
&They attribute or transport an acquired direct mode and remain independent
physical targets.\\[0.3em]
Higher jets, contraction, mass&Not inferred
&The downstream handoff remains closed until the complete quadratic card
passes.\\
\bottomrule
\end{longtable}
\endgroup

\begin{openproblem}[V9: populate the first absolute partition-curvature card]
\label{op:YMA8-V9}
Preserve Version~V8 and the frozen physical choices.  Proceed in this order.
\begin{enumerate}[label=\textup{(V9.\arabic*)},leftmargin=5.2em]
\item Resolve the weak-facing trajectory alternative.  A bounded cofinal
      subsequence exports the inherited rough-Haar/contact obstruction.  On an
      escaping branch retain the first predeclared physical shell.
\item On that shell freeze the complete response-matched unnormalized action
      $F_M(z,t)$ and select two displacements
      $0<\varepsilon_2<\varepsilon_1$ before integration.
\item Acquire rigorous outward intervals for
      $\mathscr R_{\varepsilon_1}$ and
      $\mathscr R_{\varepsilon_2}$ by deterministic integration, validated
      quadrature, or an actual calibrated partition-ratio Read on the same
      base cylinder.  A statistical interval must carry its declared
      confidence semantics and is not an exact theorem interval.
\item Populate $B_1,\ldots,B_4$ or a sharper fourth-derivative enclosure and
      acquire the same-carrier topological Hessian/secant.  Form the two
      independent local margin intervals.  Their overlap is a held-out
      consistency check, not a fit.
\item If the lower endpoint clears the target, proceed to the next
      predeclared momentum orbit and then the fused higher-jet row.  If the
      upper endpoint is negative, freeze the explicit physical polarization.
      If the interval remains unresolved, tighten this same scalar card rather
      than opening a larger auxiliary source space.
\item Open the Version~V7 boundary-source/Witten incidence only to attribute
      an actual failing or unresolved direct covariance, or when the
      independent YM--B cutoff/radial transport target requires it.
\end{enumerate}
The first missing physical value is therefore the symmetric unnormalized
partition ratio of the already occurring Wilson/generated source, not another
normalized Gram.
\end{openproblem}

\begin{assessment}[V8 termination]
\label{ass:YMA8-termination}
The upstream order, action-curvature nonidentifiability, exact finite-secant
identity, positive two-replica partition writer, fourth-derivative certificate,
topological-safe local interval, and scalar cutoff transport are \pass\ at
their stated finite mathematical scopes.  The selected response-matched
symbol remains \stopstatus: its partition-ratio and topological intervals have
not been populated.  No negative physical mode has been produced.  The next
iteration should read that one direct scalar before adding any further
boundary, current, or survivor machinery.
\end{assessment}

\section*{Version V8 append-only continuation record}
\addcontentsline{toc}{section}{Version V8 append-only continuation record}
The complete Version~V7 source above is preserved byte for byte up to its sole
final \verb|\end{document}| command.  Its SHA--256 digest is
\begin{center}\scriptsize\ttfamily
84b206129d62c4ebced82841d3806001a19347a313cb7351a79e3bb1bff2a7ec.
\end{center}
For Version~V9, preserve this accumulated source, append new material
immediately before the final document-closing command, and increment the
filename to \texttt{\_V9.tex}.  Keep the physical response, weak-facing
selector, exact block map, parity-tree section, rooted Hodge carrier, first
axial mode, source normalization, response-coefficient freezing rule, and
local/topological support convention fixed before reading any partition,
score, ancestry, radial, or transport value.

% =============================================================================
% END APPENDED VERSION V8 MATERIAL
% =============================================================================


% =============================================================================
% START APPENDED VERSION V9 MATERIAL
% =============================================================================

\clearpage
\hypersetup{
  pdftitle={Renewal Yang--Mills Reduced-Fibre Wilson Shell Symbol and Local-Vacuum Programme V9},
  pdfsubject={Action-germ incidence, elimination of the native scalar proxy, and two-radius partition-curvature acquisition}
}

\section{Version V9: action-germ incidence, proxy elimination, and two-radius partition curvature}
\label{sec:YMA9-direction}
% =============================================================================

Version~V8 moved the selected-shell problem upstream from normalized score,
Witten-current, and boundary-source geometry to the absolute unnormalized
partition curvature of the already frozen coarse mode.  The newly supplied
programmes sharpen that move in one important respect.  A finite response may
be perfectly well defined on the wrong deformation family.  A later Gram,
current, lag, or tail theorem cannot repair that earlier source-incidence
error.  The Navier--Stokes continuations express the same point as a missing
material--caloric route; the Einstein--Standard--Model continuation exhibits
it by differentiating the actual kinetic law and finding a nonzero mismatch
with an inserted-potential surrogate.  Those domain-specific results are not
imported into Yang--Mills.  Their common audit order is.

The attached boundary Version~V8 gives an even more direct Yang--Mills
correction.  Its unrescaled native scalar fan converges in absolute
$L^2$ to a constant on every escaping periodic Wilson branch.  Moreover, on a
central colour-invariant law, every equivariant cross block between that
scalar singlet and an adjoint-valued mode-score or Witten-current bank
vanishes.  Hence that fan and all of its unrescaled conditional descendants
are removed as candidates for the first YM--A action tangent.  The original
physical-radius Creutz source remains a separate YM--B target, and the actual
adjoint-valued reflected half-score remains a possible downstream ancestry
source.  Neither substitutes for the absolute coarse-link deformation used by
the shell Hessian.

The resulting order is
\begin{equation}
 \boxed{
 \begin{gathered}
 \text{actual response-matched coarse-link action germ}\\
 \Downarrow\\[-0.2em]
 \text{one two-copy action-germ incidence Read}\\
 \Downarrow\\[-0.2em]
 \text{two positive partition-ratio Reads at predeclared radii}\\
 \Downarrow\\[-0.2em]
 \text{fourth-order secant cancellation and topological subtraction}\\
 \Downarrow\\[-0.2em]
 \text{physical mode margin or an attained pure-mode failure}.
 \end{gathered}}
 \label{eq:YMA9-order}
\end{equation}
The boundary/current packets of Versions~V6--V7 are opened only after this
direct card has a physical value to attribute or transport.

\begin{firewall}[No proxy promotion]
\label{fire:YMA9-no-proxy}
A scalar Wilson-loop fan, a normalized Hellinger family, a Witten-current
Gram, a measured one-loop coefficient, and an inserted potential are all
legitimate objects on their stated carriers.  None is called the selected
axial shell deformation merely because it has the same symmetry-averaged
scalar dimension or gives a numerically similar Hessian.  Equality must be
proved at the action germ or corrected by the exact mismatch card below.
\end{firewall}

% =============================================================================
\section{The actual action deformation is an independent finite row}
\label{sec:YMA9-germ}
% =============================================================================

Let $Y$ be one finite compact reduced fibre with probability Haar measure
$m$.  Let
\begin{equation}
 F,G:(-\tau,\tau)\times Y\longrightarrow\R
 \label{eq:YMA9-two-actions}
\end{equation}
be two $C^2$ action families.  The intended physical family is $F$; $G$ is a
candidate family supplied by an inserted source, an auxiliary instrument, a
proxy action, or a second implementation.  Write
\begin{equation}
 Z_F(t)=\int_Ye^{-F(t,z)}\,dm(z),
 \qquad
 \Phi_F(t)=-\log Z_F(t),
 \label{eq:YMA9-partitions}
\end{equation}
and similarly for $G$.  Assume that
\begin{equation}
 G(0,z)-F(0,z)=c_0
 \quad\text{for all }z\in Y
 \label{eq:YMA9-common-base}
\end{equation}
with one constant $c_0$.  Thus the two normalized base laws agree.  Subtracting
$c_0$ from $G$ changes no derivative below, so put
\begin{equation}
 R(t,z):=G(t,z)-F(t,z)-c_0,
 \qquad R(0,z)=0.
 \label{eq:YMA9-remainder-action}
\end{equation}
All expectations and covariances in this section are under the common base law
\begin{equation}
 \nu_0(dz)=Z_F(0)^{-1}e^{-F(0,z)}\,dm(z).
 \label{eq:YMA9-base-law}
\end{equation}
For $j=1,2$ write $F_j=\partial_t^jF(0,\cdot)$ and
$R_j=\partial_t^jR(0,\cdot)$.

\begin{definition}[Target-curvature incidence defect]
\label{def:YMA9-germ-defect}
The signed scalar
\begin{equation}
 \boxed{
 \Delta_{G/F}^{(2)}
 :=\Phi_G''(0)-\Phi_F''(0)}
 \label{eq:YMA9-germ-defect}
\end{equation}
is the target-curvature incidence defect of the candidate deformation relative
to the physical deformation.  Vanishing of this number identifies the one
selected Hessian target.  It does not assert equality of the two full action
families or license reuse for a different mode, conditioning, or finite
secant.
\end{definition}

\begin{theorem}[Exact common-base action-germ mismatch]
\label{thm:YMA9-germ-mismatch}
Under \eqref{eq:YMA9-common-base},
\begin{equation}
 \boxed{
 \Delta_{G/F}^{(2)}
 =\mathbb E R_2
  -2\Cov(F_1,R_1)-\Var(R_1).}
 \label{eq:YMA9-germ-mismatch}
\end{equation}
Equivalently,
\begin{equation}
 \boxed{
 \Delta_{G/F}^{(2)}
 =\mathbb E R_2-\Cov(F_1+G_1,R_1),
 \qquad G_1=F_1+R_1.}
 \label{eq:YMA9-germ-mismatch-symmetric}
\end{equation}
In particular:
\begin{enumerate}[label=\textup{(G\arabic*)},leftmargin=4.0em]
\item if $G=F+c(t)$ with $c$ independent of $z$, then
      $\Delta_{G/F}^{(2)}=c''(0)$;
\item equality of the normalized base law and equality of the centered first
      scores do not determine the physical curvature, because the direct row
      $\mathbb ER_2$ remains;
\item equality of the second action means alone does not determine the
      curvature, because the two covariance terms remain.
\end{enumerate}
\end{theorem}

\begin{proof}
Differentiation under the compact integral gives the standard exact identity
\[
 \Phi_F''(0)=\mathbb EF_2-\Var(F_1).
\]
The same formula for $G$ is evaluated under the same base law because of
\eqref{eq:YMA9-common-base}.  Since $G_j=F_j+R_j$,
\[
 \begin{aligned}
 \Phi_G''(0)-\Phi_F''(0)
 &=\mathbb ER_2-
   \bigl[\Var(F_1+R_1)-\Var(F_1)\bigr]\\
 &=\mathbb ER_2-2\Cov(F_1,R_1)-\Var(R_1),
 \end{aligned}
\]
which is \eqref{eq:YMA9-germ-mismatch}.  The polarization identity
\[
 \Var(F_1+R_1)-\Var(F_1)
 =\Cov(F_1+G_1,R_1)
\]
gives \eqref{eq:YMA9-germ-mismatch-symmetric}.  The three consequences follow
by setting the indicated terms to zero.  The first is exactly the infinitesimal
coarse action-curvature gauge isolated in Version~V8.
\end{proof}

\begin{corollary}[The strongest action-germ quotient relevant to one curvature]
\label{cor:YMA9-germ-quotient}
For the selected symmetric curvature target, an affine fibre-independent
addition $a+bt$ is invisible, whereas a quadratic fibre-independent addition
is physical.  More generally, two common-base germs determine the same target
curvature exactly when the right side of
\eqref{eq:YMA9-germ-mismatch} vanishes.  This scalar equality is weaker than
physical equality of the deformation writers and is used only for the one
predeclared Hessian target.
\end{corollary}

\begin{proof}
An affine addition has $R_1=b$ constant and $R_2=0$, so every term in
\eqref{eq:YMA9-germ-mismatch} vanishes.  A quadratic addition has nonzero
$R_2$.  The final assertion is the theorem.
\end{proof}

\begin{firewall}[A different base law requires an explicit bridge]
\label{fire:YMA9-base-bridge}
If \eqref{eq:YMA9-common-base} fails, the two Hessians are computed in
different normalized laws and \eqref{eq:YMA9-germ-mismatch} is not the correct
comparison.  One must first supply the common unnormalized cylinder or an
explicit positive Radon--Nikodym bridge with its normalization.  Matching
marginal moments or separately whitened score Grams does not create that
bridge.
\end{firewall}

% =============================================================================
\section{One two-copy writer acquires the complete germ mismatch}
\label{sec:YMA9-germ-writer}
% =============================================================================

For two independent base samples $z_1,z_2\sim\nu_0$, write
\begin{equation}
 \Delta A:=A(z_1)-A(z_2)
 \label{eq:YMA9-pair-difference}
\end{equation}
for every scalar writer $A$.

\begin{theorem}[Target-minimal two-replica action-germ Read]
\label{thm:YMA9-germ-replica}
Define the signed two-copy writer
\begin{equation}
 \boxed{
 \begin{aligned}
 \mathscr M_{G/F}(z_1,z_2)
 :={}&\frac12\bigl(R_2(z_1)+R_2(z_2)\bigr)\\
 &-\Delta F_1\,\Delta R_1
 -\frac12(\Delta R_1)^2.
 \end{aligned}}
 \label{eq:YMA9-germ-replica-writer}
\end{equation}
Then
\begin{equation}
 \boxed{
 \mathbb E_{\nu_0}^{\otimes2}\mathscr M_{G/F}
 =\Delta_{G/F}^{(2)}.}
 \label{eq:YMA9-germ-replica-identity}
\end{equation}
Equivalently,
\begin{equation}
 \mathscr M_{G/F}
 =\frac12(R_2^{(1)}+R_2^{(2)})
  -\frac12\Delta(F_1+G_1)\,\Delta R_1.
 \label{eq:YMA9-germ-replica-symmetric}
\end{equation}
Thus one same-base two-copy panel acquires the direct second-action mismatch,
the physical--proxy mixed covariance, and the proxy innovation with their
actual signs before an interval is taken.
\end{theorem}

\begin{proof}
For independent identically distributed copies,
\[
 \mathbb E[\Delta A\,\Delta B]=2\Cov(A,B),
 \qquad
 \frac12\mathbb E(\Delta B)^2=\Var(B).
\]
The expectation of the first line in
\eqref{eq:YMA9-germ-replica-writer} is $\mathbb ER_2$.  Substitution gives
\eqref{eq:YMA9-germ-mismatch}, proving
\eqref{eq:YMA9-germ-replica-identity}.  Since
$\Delta(F_1+G_1)=2\Delta F_1+\Delta R_1$, the two displayed writers agree
pointwise.
\end{proof}

\begin{corollary}[Outward correction of a candidate partition curvature]
\label{cor:YMA9-germ-interval}
Suppose an outward interval
\begin{equation}
 \underline\Delta\le\Delta_{G/F}^{(2)}\le\overline\Delta
 \label{eq:YMA9-germ-interval}
\end{equation}
is supplied by the writer in
\eqref{eq:YMA9-germ-replica-writer}, and suppose
\(
 \underline H_G\le\Phi_G''(0)\le\overline H_G
\).
Then the physical curvature satisfies
\begin{equation}
 \boxed{
 \underline H_G-\overline\Delta
 \le\Phi_F''(0)\le
 \overline H_G-\underline\Delta.}
 \label{eq:YMA9-corrected-curvature-interval}
\end{equation}
If a candidate is submitted under the stronger claim that it is the physical
deformation itself, an interval for
$\Delta_{G/F}^{(2)}$ excluding zero rejects that claim.  It is an
action-incidence error, not a negative Yang--Mills mode.
\end{corollary}

\begin{proof}
By definition,
$\Phi_F''(0)=\Phi_G''(0)-\Delta_{G/F}^{(2)}$.  Interval subtraction reverses
the endpoints of the defect interval.  The final statement concerns the
domain of the claimed identification rather than the sign of the corrected
physical Hessian.
\end{proof}

\begin{corollary}[Gauge-orbit germ held-out]
\label{cor:YMA9-gauge-germ}
Let $F$ be the actual deformation along a residual coarse gauge orbit.  Then
$\Phi_F$ is constant and $\Phi_F''(0)=0$.  Any candidate $G$ for that gauge
orbit must therefore satisfy
\begin{equation}
 \boxed{
 \Phi_G''(0)=\Delta_{G/F}^{(2)}.}
 \label{eq:YMA9-gauge-germ-check}
\end{equation}
The partition-ratio value one from Version~V8 and the germ-mismatch writer
must agree with this identity.  Failure is a common-action or gauge-incidence
error.
\end{corollary}

\begin{proof}
Gauge invariance makes the physical partition constant along the orbit.  Use
\Cref{thm:YMA9-germ-mismatch}.
\end{proof}

% =============================================================================
\section{Elimination of the native scalar fan as an axial action proxy}
\label{sec:YMA9-native-pruning}
% =============================================================================

The supplied boundary Version~V8 proves two results which are imported here
without repeating their combinatorial or colour-twirl proofs.
\begin{enumerate}[label=\textup{(B8.\arabic*)},leftmargin=4.5em]
\item On every escaping periodic Wilson root sequence, the unrescaled native
      fan $f_N$ converges in $L^2$ to its constant maximum, and every
      conditional descendant has vanishing absolute stock.
\item On a central-vacuum invariant law, an equivariant linear cross block
      between a colour-singlet scalar and an adjoint-valued mode-score or
      Witten-current synthesis is exactly zero.
\end{enumerate}
The same source also proves that the first nonconstant local fan response is a
quadratic colour singlet rather than an adjoint linear source.

\begin{theorem}[Representation-typed action-germ gate]
\label{thm:YMA9-colour-germ-gate}
On the central-vacuum axial branch, let
\begin{equation}
 X\longmapsto F_{1,X}^{\rm phys}\in L^2(\nu_0)
 \label{eq:YMA9-physical-adjoint-score}
\end{equation}
be the actual linear coarse-link action tangent in one transverse direction.
It is an equivariant adjoint-indexed family.  Let a proposed proxy tangent be
constructed from colour-singlet scalar writers only, with no additional
adjoint occurrence index.  Then its equivariant linear map in $X$ is zero.
Consequently it cannot establish equality with
\eqref{eq:YMA9-physical-adjoint-score} unless the physical tangent itself
vanishes.  A scalar quadratic response may contribute to a second-action
number, but it does not identify the missing linear score or its covariance.
\end{theorem}

\begin{proof}
The vanishing of the singlet--adjoint cross block and of every equivariant
linear singlet-to-adjoint map is the imported colour-selection theorem of the
boundary Version~V8.  The physical coarse-link derivative transforms in the
adjoint index by simultaneous global conjugation.  Therefore a proxy assembled
only from singlets has zero equivariant linear tangent.  A quadratic singlet
belongs to the symmetric square of the adjoint coefficient and can supply a
scalar second-order response, but it contains no identification of the
adjoint linear score.  The latter distinction is exactly what enters
\eqref{eq:YMA9-germ-mismatch} through $R_1$ and its covariance terms.
\end{proof}

\begin{corollary}[The escaping native branch is pruned before the partition Read]
\label{cor:YMA9-native-pruned}
On an escaping weak-facing Wilson branch, no unrescaled native-fan Read or
conditional native-fan descendant can populate the Version~V8 absolute axial
partition curvature.  A deterministic rescaling large enough to prevent its
variance collapse is a new calibration row, and it remains colour singlet.
The admissible first deformation is therefore the actual response-matched
coarse-link family
\begin{equation}
 \boxed{
 F_{M,q,X}^{\rm phys}(t,z)
 :=F_M^{\rm RG}\bigl(z,W_{q,X}(t)\bigr),}
 \label{eq:YMA9-actual-action-family}
\end{equation}
with every generated, marginal, and support term retained and with
$W_{q,X}(t)$ frozen before the partition values are inspected.
\end{corollary}

\begin{proof}
Absolute collapse is imported from boundary Version~V8.  A rescaling changes
the physical writer and needs a new scale anchor; it does not change the
trivial colour representation.  The representation gate then excludes it as
the adjoint linear tangent.  The family
\eqref{eq:YMA9-actual-action-family} is the unnormalized deformation already
required by Version~V8.
\end{proof}

\begin{firewall}[The physical-radius Creutz source remains independent]
\label{fire:YMA9-radial-independent}
The boundary Version~V8 correctly returns to its original growing physical
rectangle source after pruning the native fan.  Its annular attenuation ratio
is a YM--B spatial target.  It is not inserted into
\eqref{eq:YMA9-actual-action-family}, and failure to evaluate it does not
block the finite YM--A partition curvature at one already loaded shell.
Conversely, a positive YM--A mode margin does not evaluate that radial ratio.
\end{firewall}

% =============================================================================
\section{Two-radius cancellation of the finite-secant bias}
\label{sec:YMA9-Richardson}
% =============================================================================

Retain the symmetric partition curvature $Q_\varepsilon$ of
\Cref{def:YMA8-partition-curvature}.  Version~V8 requested two independently
predeclared displacements.  Their first use is to cancel the entire quartic
central-difference bias rather than merely check overlap of two
$O(\varepsilon^2)$ intervals.

\begin{theorem}[General two-radius Richardson curvature]
\label{thm:YMA9-Richardson}
Let $\Phi\in C^6([-\varepsilon_1,\varepsilon_1])$ and let
$0<\varepsilon_2<\varepsilon_1$.  Put
\begin{equation}
 Q_{\varepsilon}
 =\frac{\Phi(\varepsilon)+\Phi(-\varepsilon)-2\Phi(0)}{\varepsilon^2}
 \label{eq:YMA9-Q-def}
\end{equation}
and define
\begin{equation}
 \boxed{
 \mathcal Q_{\varepsilon_1,\varepsilon_2}^{[4]}
 :=\frac{\varepsilon_1^2Q_{\varepsilon_2}
          -\varepsilon_2^2Q_{\varepsilon_1}}
         {\varepsilon_1^2-\varepsilon_2^2}.}
 \label{eq:YMA9-Richardson}
\end{equation}
If
\begin{equation}
 M_6(\varepsilon_1)
 :=\sup_{|t|\le\varepsilon_1}|\Phi^{(6)}(t)|,
 \label{eq:YMA9-M6}
\end{equation}
then
\begin{equation}
 \boxed{
 \left|
 \mathcal Q_{\varepsilon_1,\varepsilon_2}^{[4]}-\Phi''(0)
 \right|
 \le
 \frac{\varepsilon_1^2\varepsilon_2^2}{360}
 M_6(\varepsilon_1).}
 \label{eq:YMA9-Richardson-error}
\end{equation}
The constant is sharp.  For the dyadic choice
$\varepsilon_2=\varepsilon_1/2$ this becomes
\begin{equation}
 \boxed{
 \left|
 \frac{4Q_{\varepsilon_1/2}-Q_{\varepsilon_1}}3-\Phi''(0)
 \right|
 \le\frac{\varepsilon_1^4}{1440}M_6(\varepsilon_1).}
 \label{eq:YMA9-dyadic-Richardson}
\end{equation}
\end{theorem}

\begin{proof}
Let $a=\varepsilon_1$, $b=\varepsilon_2$, and let $\mathcal L$ be the
left side of \eqref{eq:YMA9-Richardson} minus $\Phi''(0)$.  This linear
functional annihilates every polynomial of degree at most five.  Its sixth-
order Peano kernel is
\[
 K(s)=\mathcal L\left(\frac{(t-s)_+^5}{5!}\right),
 \qquad
 \mathcal L(\Phi)=\int_{-a}^{a}K(s)\Phi^{(6)}(s)\,ds.
\]
The functional is reflection invariant, so $K$ is even.  For $0\le s\le b$,
\begin{equation}
 K(s)=\frac1{120(a^2-b^2)}
 \left[
  \frac{a^2}{b^2}(b-s)^5-
  \frac{b^2}{a^2}(a-s)^5
 \right],
 \label{eq:YMA9-Peano-inner}
\end{equation}
whereas for $b\le s\le a$,
\begin{equation}
 K(s)=-\frac{b^2(a-s)^5}{120a^2(a^2-b^2)}.
 \label{eq:YMA9-Peano-outer}
\end{equation}
Now $(b-s)/(a-s)\le b/a$ for $0\le s\le b$.  Hence the first bracket is
nonpositive; the outer formula is also nonpositive.  Therefore $K\le0$.
Finally,
\[
 \int_{-a}^{a}K(s)\,ds
 =\mathcal L(t^6/6!)=-\frac{a^2b^2}{360}.
\]
Thus
\[
 |\mathcal L(\Phi)|
 \le M_6(a)\int|K(s)|\,ds
 =\frac{a^2b^2}{360}M_6(a),
\]
which proves \eqref{eq:YMA9-Richardson-error}.  Equality is attained by a
pure sixth-degree perturbation with constant sixth derivative.  The dyadic
constant is $1/1440$.
\end{proof}

\begin{corollary}[Outward interval from two positive partition ratios]
\label{cor:YMA9-Richardson-interval}
Suppose Version~V8's positive two-copy Reads give
\begin{equation}
 0<\underline R_i\le\mathscr R_{\varepsilon_i}
 \le\overline R_i<\infty,
 \qquad i=1,2.
 \label{eq:YMA9-two-ratio-intervals}
\end{equation}
Put
\begin{equation}
 \underline Q_i=-\varepsilon_i^{-2}\log\overline R_i,
 \qquad
 \overline Q_i=-\varepsilon_i^{-2}\log\underline R_i,
 \label{eq:YMA9-Q-intervals}
\end{equation}
and
\begin{align}
 \underline Q^{[4]}
 &:=\frac{\varepsilon_1^2\underline Q_2
          -\varepsilon_2^2\overline Q_1}
         {\varepsilon_1^2-\varepsilon_2^2},
 \label{eq:YMA9-Richardson-lower}\\
 \overline Q^{[4]}
 &:=\frac{\varepsilon_1^2\overline Q_2
          -\varepsilon_2^2\underline Q_1}
         {\varepsilon_1^2-\varepsilon_2^2}.
 \label{eq:YMA9-Richardson-upper}
\end{align}
With
\begin{equation}
 E_6:=\frac{\varepsilon_1^2\varepsilon_2^2}{360}
       M_6(\varepsilon_1),
 \label{eq:YMA9-E6}
\end{equation}
one has
\begin{equation}
 \boxed{
 \underline Q^{[4]}-E_6
 \le\Phi''(0)\le
 \overline Q^{[4]}+E_6.}
 \label{eq:YMA9-Richardson-H-interval}
\end{equation}
The two ratios remain independent physical Reads.  Their combination is fixed
before their values are inspected; it is not a post hoc extrapolation fit.
\end{corollary}

\begin{proof}
The map $r\mapsto-\log r$ is decreasing on $(0,\infty)$, giving
\eqref{eq:YMA9-Q-intervals}.  In
\eqref{eq:YMA9-Richardson}, the coefficient of $Q_{\varepsilon_2}$ is
positive and that of $Q_{\varepsilon_1}$ is negative, which gives the two
outward endpoints.  Apply \Cref{thm:YMA9-Richardson}.
\end{proof}

\begin{remark}[Why two radii are now acquisition rather than redundancy]
A single radius and a fourth-derivative bound remain valid by Version~V8.
The two-radius formula pays the sixth derivative and removes the complete
quartic bias.  It is useful precisely when the ratio intervals are accurate
enough that the $O(\varepsilon^4)$ analytic radius is smaller than the
statistical or deterministic interval width.  Otherwise the one-radius card
may still be sharper.
\end{remark}

% =============================================================================
\section{The sixth derivative is one finite cumulant card}
\label{sec:YMA9-sixth}
% =============================================================================

Let $F(t,z)$ be the actual finite action, let
$X_j(t,z)=\partial_t^jF(t,z)$, and let $\kappa_t$ denote joint cumulants under
the normalized law with density proportional to $e^{-F(t,z)}$.  Denote by
$\Pi_n$ the set of set partitions of $\{1,\ldots,n\}$.

\begin{lemma}[Partition formula for every partition derivative]
\label{lem:YMA9-cumulant-partitions}
For every $n\ge1$ for which the derivatives exist,
\begin{equation}
 \boxed{
 \Phi^{(n)}(t)
 =\sum_{\pi\in\Pi_n}(-1)^{|\pi|-1}
   \kappa_t\bigl(X_{|B|}(t):B\in\pi\bigr),
 \qquad \Phi=-\log Z.}
 \label{eq:YMA9-cumulant-partition-formula}
\end{equation}
The entries of a cumulant are indexed by the block sizes of $\pi$.
\end{lemma}

\begin{proof}
For a $t$-dependent joint cumulant,
\begin{equation}
 \frac d{dt}\kappa_t(A_1,\ldots,A_r)
 =\sum_{j=1}^r\kappa_t(A_1,\ldots,A_j',\ldots,A_r)
  -\kappa_t(A_1,\ldots,A_r,X_1).
 \label{eq:YMA9-cumulant-derivative}
\end{equation}
This follows by differentiating the logarithm of the joint moment-generating
function under the exponentially tilted law.  Formula
\eqref{eq:YMA9-cumulant-partition-formula} is true for $n=1$ because
$\Phi'=\mathbb EX_1$.  Assume it at order $n$.  Differentiating one existing
entry enlarges exactly one block of a partition by the new label $n+1$ and
keeps its sign.  The last term of
\eqref{eq:YMA9-cumulant-derivative} adds the singleton block
$\{n+1\}$ and changes the sign.  Every partition of $n+1$ occurs uniquely in
one of these two ways.  This proves the induction.
\end{proof}

\begin{proposition}[Exact sixth partition derivative]
\label{prop:YMA9-sixth-cumulant}
The sixth derivative is
\begin{equation}
 \boxed{
 \begin{aligned}
 \Phi^{(6)}={}&\mathbb EX_6
 -6\kappa(X_5,X_1)
 -15\kappa(X_4,X_2)
 -10\kappa(X_3,X_3)\\
 &+15\kappa(X_4,X_1,X_1)
 +60\kappa(X_3,X_2,X_1)
 +15\kappa(X_2,X_2,X_2)\\
 &-20\kappa(X_3,X_1,X_1,X_1)
 -45\kappa(X_2,X_2,X_1,X_1)\\
 &+15\kappa(X_2,X_1,X_1,X_1,X_1)
 -\kappa(X_1,X_1,X_1,X_1,X_1,X_1).
 \end{aligned}}
 \label{eq:YMA9-sixth-cumulant}
\end{equation}
All expectations and cumulants are at the same occupied parameter $t$.
\end{proposition}

\begin{proof}
Group the set partitions in
\eqref{eq:YMA9-cumulant-partition-formula} by block-size type.  The numbers of
partitions of types
\[
 6,\ 5+1,\ 4+2,\ 3+3,\ 4+1+1,\ 3+2+1,\ 2+2+2,
\]
\[
 3+1+1+1,\ 2+2+1+1,\ 2+1+1+1+1,\ 1+1+1+1+1+1
\]
are respectively
\[
 1,6,15,10,15,60,15,20,45,15,1.
\]
The sign is $(-1)^{r-1}$ for $r$ blocks, which gives the displayed formula.
\end{proof}

For a bounded real writer $A$, write
\begin{equation}
 \operatorname{rad}(A)
 :=\frac12\bigl(\sup_zA(z)-\inf_zA(z)\bigr).
 \label{eq:YMA9-osc-radius}
\end{equation}
Cumulants of order at least two are unchanged when a constant is added to any
entry.  Hence oscillation radii, rather than absolute action offsets, control
the connected terms.

\begin{theorem}[Executable sixth-derivative bound with correct direct row]
\label{thm:YMA9-M6-bound}
For $1\le j\le6$ define
\begin{equation}
 B_j=\sup_{|t|\le\varepsilon_1}\sup_z|X_j(t,z)|,
 \qquad
 b_j=\sup_{|t|\le\varepsilon_1}\operatorname{rad}(X_j(t,\cdot)).
 \label{eq:YMA9-Bb}
\end{equation}
Then
\begin{equation}
 \boxed{
 \begin{aligned}
 M_6(\varepsilon_1)\le{}&B_6
 +12b_5b_1+30b_4b_2+20b_3^2\\
 &+90b_4b_1^2+360b_3b_2b_1+90b_2^3\\
 &+520b_3b_1^3+1170b_2^2b_1^2\\
 &+2250b_2b_1^4+1082b_1^6.
 \end{aligned}}
 \label{eq:YMA9-M6-bound}
\end{equation}
The absolute direct sixth-action row $B_6$ is retained.  Every higher-order
cumulant uses only fibre oscillation and is therefore insensitive to an
irrelevant constant offset, while remaining sensitive to a physical
$t$-dependent partition gauge through its direct derivative.
\end{theorem}

\begin{proof}
For bounded variables $|Y_j|\le c_j$, the moment--cumulant formula gives
\begin{equation}
 |\kappa(Y_1,\ldots,Y_r)|
 \le C_r\prod_{j=1}^rc_j,
 \qquad
 C_r:=\sum_{\pi\in\Pi_r}(|\pi|-1)!.
 \label{eq:YMA9-cumulant-bound}
\end{equation}
Indeed each partition term has magnitude at most the common product.  The
needed constants are
\begin{equation}
 C_2=2,\qquad C_3=6,\qquad C_4=26,
 \qquad C_5=150,\qquad C_6=1082.
 \label{eq:YMA9-cumulant-constants}
\end{equation}
For every cumulant in \eqref{eq:YMA9-sixth-cumulant}, subtract the midpoint of
the range from each entry; this does not change a cumulant of order at least
two and leaves the centered entry bounded by $b_j$.  Apply
\eqref{eq:YMA9-cumulant-bound} term by term.  Multiplying the partition-type
coefficients by $C_r$ gives respectively
\[
 12,30,20,90,360,90,520,1170,2250,1082,
\]
which proves \eqref{eq:YMA9-M6-bound}.  The one-block term is
$\mathbb EX_6$ and is bounded by $B_6$; it cannot be replaced by an oscillation
without losing the absolute partition curvature.
\end{proof}

\begin{corollary}[Literal Wilson insertion stocks through order six]
\label{cor:YMA9-Wilson-six-stocks}
Use the parity-tree section and vary one coarse link direction only, with a
unit Frobenius colour generator $X$ and a real normalized axial wave $h_y$.
Every fine plaquette contains at most one designated second constituent in
that fixed direction, and each designated link belongs to six plaquettes.
Consequently, for the pure Wilson part,
\begin{equation}
 \boxed{
 B_j^{\rm W}\le6\beta\sum_{y\in\Lambda_M}|h_y|^j,
 \qquad 1\le j\le6.}
 \label{eq:YMA9-Wilson-j-bound}
\end{equation}
For a non-self-conjugate normalized sine or cosine mode,
$\|h\|_2=1$ and $\|h\|_\infty\le\sqrt2/M^2$, so in particular
\begin{equation}
 \boxed{
 B_5^{\rm W}\le\frac{12\sqrt2\,\beta}{M^6},
 \qquad
 B_6^{\rm W}\le\frac{24\beta}{M^8}.}
 \label{eq:YMA9-Wilson-56}
\end{equation}
The earlier bounds
\begin{equation}
 B_1^{\rm W}\le6\beta M^2,
 \quad B_2^{\rm W}\le6\beta,
 \quad B_3^{\rm W}\le\frac{6\sqrt2\beta}{M^2},
 \quad B_4^{\rm W}\le\frac{12\beta}{M^4}
 \label{eq:YMA9-Wilson-14}
\end{equation}
are recovered.  If the submitted generated action has derivative bounds
$B_j^\Psi$, the full physical stocks obey
\begin{equation}
 B_j\le B_j^{\rm W}+B_j^\Psi.
 \label{eq:YMA9-generated-six-stocks}
\end{equation}
No generated contribution is set to zero.
\end{corollary}

\begin{proof}
Along the chosen direction a varying fine plaquette word contains at most one
factor $\exp(th_yX)$.  Its $j$-th derivative inserts $h_y^jX^j$.
\Cref{lem:YMA8-Wilson-insertion-bound} bounds the normalized trace by
$|h_y|^j\|X\|_{\rm op}^j\le|h_y|^j$.  Summing the six plaquettes incident to
each designated link proves \eqref{eq:YMA9-Wilson-j-bound}.  For $j\ge2$,
\[
 \sum_y|h_y|^j
 \le\|h\|_\infty^{j-2}\sum_y|h_y|^2.
\]
Substitution gives \eqref{eq:YMA9-Wilson-56} and the stated earlier stocks;
Cauchy--Schwarz gives $\sum_y|h_y|\le M^2$ for $j=1$.  Derivatives add for the
full action, proving the last assertion.
\end{proof}

\begin{remark}[Finite-card rigor versus cofinal usefulness]
The bound \eqref{eq:YMA9-M6-bound} is fully explicit once the finite action
writers are enumerated, but powers of the absolute first-derivative stock can
be conservative.  The oscillation-centred form already removes irrelevant
volume offsets.  A cutoff-uniform result may still require connected
clustering or a local cumulant estimate.  Failure of the crude bound to be
small is \stopstatus, not a physical soft mode.
\end{remark}

% =============================================================================
\section{A rigorous positive-writer sampling interval}
\label{sec:YMA9-sampling}
% =============================================================================

The ratio writer of Version~V8 is positive and bounded on every finite compact
card.  The next theorem records exact confidence semantics for a direct sample
of that one scalar.  It is optional when deterministic quadrature or exact
integration is available.

\begin{theorem}[Finite-sample interval for one partition ratio]
\label{thm:YMA9-Hoeffding-ratio}
Let $W_1,\ldots,W_n$ be independent copies of the positive writer
$L_+(z_1)L_-(z_2)$ in
\Cref{thm:YMA8-two-replica-partition}.  Suppose deterministic bounds
\begin{equation}
 0<\ell\le W_i\le u<\infty
 \label{eq:YMA9-writer-range}
\end{equation}
hold.  Put $\overline W_n=n^{-1}\sum_iW_i$.  For
$0<\alpha<1$, define
\begin{equation}
 r_{n,\alpha}=(u-\ell)
 \sqrt{\frac{\log(2/\alpha)}{2n}}.
 \label{eq:YMA9-Hoeffding-radius}
\end{equation}
Then, with probability at least $1-\alpha$,
\begin{equation}
 \boxed{
 \max\{\ell,\overline W_n-r_{n,\alpha}\}
 \le\mathscr R_\varepsilon\le
 \min\{u,\overline W_n+r_{n,\alpha}\}.}
 \label{eq:YMA9-ratio-confidence}
\end{equation}
If
\begin{equation}
 D_\pm=\sup_z|F(\pm\varepsilon,z)-F(0,z)|,
 \label{eq:YMA9-Dpm}
\end{equation}
then one may take
\begin{equation}
 \boxed{
 \ell=e^{-(D_++D_-)},
 \qquad u=e^{D_++D_-}.}
 \label{eq:YMA9-writer-range-action}
\end{equation}
In particular $D_\pm\le\varepsilon B_1(\varepsilon)$.
\end{theorem}

\begin{proof}
Hoeffding's exponential-moment argument for bounded independent variables
gives
\[
 \Pr\bigl(|\overline W_n-\mathbb EW_1|\ge r\bigr)
 \le2\exp\left(-\frac{2nr^2}{(u-\ell)^2}\right).
\]
Set the right side equal to $\alpha$ and intersect the resulting interval with
the deterministic range.  Version~V8 identifies
$\mathbb EW_1=\mathscr R_\varepsilon$.  The exponent in $W_i$ lies between
$-(D_++D_-)$ and $D_++D_-$, proving
\eqref{eq:YMA9-writer-range-action}.  The final bound is the fundamental
theorem of calculus.
\end{proof}

\begin{corollary}[Joint two-radius confidence card]
\label{cor:YMA9-two-radius-confidence}
Use any common base samples at the two radii and assign failure probabilities
$\alpha_1,\alpha_2>0$.  With probability at least
$1-\alpha_1-\alpha_2$, both ratio intervals
\eqref{eq:YMA9-ratio-confidence} hold simultaneously.  On that event,
\Cref{cor:YMA9-Richardson-interval} is a valid confidence interval for the
curvature.  Dependence between the two radius estimators is allowed; only the
within-radius independent sample assumption and the union bound are used.
\end{corollary}

\begin{proof}
Apply \Cref{thm:YMA9-Hoeffding-ratio} to each radius and use the union bound.
The Richardson transformation is deterministic on the event that both input
intervals cover their means.
\end{proof}

\begin{firewall}[Confidence is not exact enclosure]
\label{fire:YMA9-confidence}
A $1-\alpha$ coverage theorem is a rigorous statistical statement, but it is
not an exact interval valid on every sample outcome.  The final card must
retain its confidence level, sampling law, independence convention, and any
multiple-radius allocation.  A deterministic PASS or obstruction is reported
only from deterministic integration or from an exact proof-carrying Read.
\end{firewall}

% =============================================================================
\section{Complete physical interval with germ and topological correction}
\label{sec:YMA9-complete-card}
% =============================================================================

The preceding rows now assemble without opening a larger source carrier.
Let $G$ be the action family used by the available partition-ratio instrument
and $F$ the frozen physical family
\eqref{eq:YMA9-actual-action-family}.  The preferred branch is $G=F$, in which
case the germ correction below is identically zero.  A proxy is usable for
this one target only when its common-base relation to $F$ and its mismatch
interval are both supplied.

\begin{theorem}[Two-radius, germ-corrected, topological-safe axial interval]
\label{thm:YMA9-complete-interval}
Assume the following finite data on one predeclared axial polarization:
\begin{enumerate}[label=\textup{(A\arabic*)},leftmargin=4.0em]
\item positive partition-ratio intervals
      \eqref{eq:YMA9-two-ratio-intervals} for the candidate family $G$;
\item a sixth-derivative bound $M_6^G(\varepsilon_1)$;
\item a common-base action-germ mismatch interval
      $[\underline\Delta,\overline\Delta]$ for
      $\Delta_{G/F}^{(2)}$;
\item a same-carrier physical topological interval
      $[\underline t,\overline t]$ for the full-to-local subtraction;
\item a positive reference interval
      $0<\underline s\le s_M(q)\le\overline s$.
\end{enumerate}
Let $\underline Q^{[4]},\overline Q^{[4]}$ be
\eqref{eq:YMA9-Richardson-lower}--\eqref{eq:YMA9-Richardson-upper} and let
$E_6$ be \eqref{eq:YMA9-E6}.  Then the physical local Hessian obeys
\begin{equation}
 \boxed{
 \begin{aligned}
 \underline h_F^{\rm loc}
 &:=\underline Q^{[4]}-E_6-\overline\Delta-\overline t,\\
 \overline h_F^{\rm loc}
 &:=\overline Q^{[4]}+E_6-\underline\Delta-\underline t,
 \end{aligned}}
 \label{eq:YMA9-complete-local-interval}
\end{equation}
that is,
\begin{equation}
 \underline h_F^{\rm loc}
 \le h_F^{\rm loc}(q)\le
 \overline h_F^{\rm loc}.
 \label{eq:YMA9-complete-local-inclusion}
\end{equation}
For a target $\delta\ge0$, put
\begin{equation}
 \boxed{
 \underline m_F=\underline h_F^{\rm loc}-\delta\overline s,
 \qquad
 \overline m_F=\overline h_F^{\rm loc}-\delta\underline s.}
 \label{eq:YMA9-complete-margin}
\end{equation}
Then
\begin{equation}
 \boxed{
 \begin{array}{rcl}
 \underline m_F>0&\Longrightarrow&\pass,\\[0.2em]
 \overline m_F<0&\Longrightarrow&\obstruction,\\[0.2em]
 0\in[\underline m_F,\overline m_F]&\Longrightarrow&\stopstatus.
 \end{array}}
 \label{eq:YMA9-verdict}
\end{equation}
On the obstruction branch, the frozen real axial polarization is the attained
physical witness.  A rejected common-base or action-germ claim is an input
error before these three mathematical outcomes.
\end{theorem}

\begin{proof}
\Cref{cor:YMA9-Richardson-interval} bounds the candidate full curvature.
\Cref{cor:YMA9-germ-interval} subtracts the candidate-minus-physical germ
interval, and the physical local card is the resulting full curvature minus
the topological Hessian.  Two successive interval subtractions give
\eqref{eq:YMA9-complete-local-interval}.  Since $\delta\ge0$ and the reference
is positive, the lower margin uses $\overline s$ and the upper margin uses
$\underline s$.  The sign alternatives are ordinary interval inclusion and
the axial scalarization inherited from Versions~V2 and V8.
\end{proof}

\begin{corollary}[Exact physical-action specialization]
\label{cor:YMA9-physical-specialization}
If the ratio writer is constructed directly from the physical family
$G=F=F_{M,q,X}^{\rm phys}$, take
\begin{equation}
 \underline\Delta=\overline\Delta=0.
 \label{eq:YMA9-zero-germ-correction}
\end{equation}
No normalized boundary-source or current-incidence row is then required for
the direct quadratic verdict.  Such rows remain available to attribute a
small covariance or to transport a passing card.
\end{corollary}

\begin{proof}
The action difference $R$ is identically zero, so
\Cref{thm:YMA9-germ-mismatch} gives zero.  Apply
\Cref{thm:YMA9-complete-interval}.
\end{proof}

\begin{theorem}[Root-first, proxy-free selected-shell alternative]
\label{thm:YMA9-root-alternative}
For the frozen weak-facing selector, the next YM--A step has the following
ordered alternatives.
\begin{enumerate}[label=\textup{(V9.\arabic*)},leftmargin=5.0em]
\item A bounded cofinal root subsequence exports the inherited
      rough-Haar/contact trajectory obstruction and terminates the intended raw
      smooth local-vacuum branch.
\item On an escaping branch, the boundary native scalar fan and all of its
      unrescaled descendants are pruned by
      \Cref{cor:YMA9-native-pruned}; no scalar-to-adjoint overlap panel is
      opened.
\item If the physical unnormalized action family
      \eqref{eq:YMA9-actual-action-family} is absent, the verdict is
      \stopstatus\ at the loading row.
\item If a candidate ratio instrument is supplied on a common base, first
      acquire \eqref{eq:YMA9-germ-replica-identity}.  A failed claim of zero
      mismatch rejects the proxy; a certified mismatch may instead be retained
      as the correction in \eqref{eq:YMA9-complete-local-interval}.
\item Once the two physical or corrected ratios, sixth-derivative stock, and
      topological interval are present, the finite card terminates by
      \eqref{eq:YMA9-verdict}.
\item Boundary/current ancestry, physical-radius attenuation, higher jets, and
      adjacent-cutoff transport open only after the direct verdict or when
      independently required by their own target.
\end{enumerate}
No fourth branch called ``develop more shell machinery'' is available.
\end{theorem}

\begin{proof}
The bounded/escaping dichotomy and the bounded-branch consequence are inherited.
The native-fan pruning is the supplied boundary Version~V8 together with
\Cref{thm:YMA9-colour-germ-gate}.  The loading and candidate-incidence branches
are \Cref{fire:YMA9-base-bridge,thm:YMA9-germ-replica}.  The finite verdict is
\Cref{thm:YMA9-complete-interval}.  The final ordering is exactly the upstream
physical-loading discipline of Versions~V7--V8.
\end{proof}

% =============================================================================
\section{Version V9 terminal theorem, proof ledger, and next acquisition}
\label{sec:YMA9-terminal}
% =============================================================================

\begin{theorem}[Version V9 action-germ and two-radius theorem]
\label{thm:YMA9-terminal}
Versions~V1--V8 are retained.  Version~V9 additionally proves:
\begin{enumerate}[label=\textup{(E9.\arabic*)},leftmargin=5.2em]
\item two common-base unnormalized deformation families differ at the selected
      Hessian by the exact signed action-germ defect
      \eqref{eq:YMA9-germ-mismatch};
\item one two-copy writer acquires that complete defect without separately
      pricing its direct, mixed, and innovation terms;
\item a nonzero defect rejects a claimed physical-deformation identification
      but can be retained as an exact correction for the one scalar target;
\item the newly supplied boundary result removes the unrescaled native scalar
      fan as an escaping-branch source and forbids its equivariant linear
      identification with the adjoint mode-score bank;
\item two predeclared partition-ratio radii yield the Richardson combination
      \eqref{eq:YMA9-Richardson}, which cancels the quartic secant bias and has
      the explicit sixth-derivative error
      \eqref{eq:YMA9-Richardson-error};
\item the sixth partition derivative is the finite cumulant card
      \eqref{eq:YMA9-sixth-cumulant}, with the executable oscillation-centred
      bound \eqref{eq:YMA9-M6-bound};
\item the literal one-direction Wilson insertion stocks extend through orders
      five and six with the explicit bounds
      \eqref{eq:YMA9-Wilson-56};
\item bounded positive two-copy samples admit the exact coverage interval
      \eqref{eq:YMA9-ratio-confidence}, and two radii combine by a union-bound
      coverage theorem;
\item candidate curvature, germ correction, topological subtraction, and the
      positive reference assemble into the terminating physical interval
      \eqref{eq:YMA9-complete-margin};
\item the weak-facing continuation is root first, proxy free, and has only
      PASS, TYPED OBSTRUCTION, STOP, or a prior rejected-input outcome.
\end{enumerate}
No item supplies an observed selected-root partition ratio, proves root escape,
evaluates the physical-radius Creutz attenuation, controls the nonaxial orbit
bank, or establishes higher-jet contraction or a continuum mass gap.
\end{theorem}

\begin{proof}
Items~\textup{(E9.1)--(E9.3)} are
\Cref{thm:YMA9-germ-mismatch,thm:YMA9-germ-replica,cor:YMA9-germ-interval}.
Item~\textup{(E9.4)} is
\Cref{thm:YMA9-colour-germ-gate,cor:YMA9-native-pruned} using the supplied
boundary Version~V8.  Item~\textup{(E9.5)} is
\Cref{thm:YMA9-Richardson,cor:YMA9-Richardson-interval}.
Items~\textup{(E9.6)--(E9.7)} are
\Cref{prop:YMA9-sixth-cumulant,thm:YMA9-M6-bound,cor:YMA9-Wilson-six-stocks}.
Item~\textup{(E9.8)} is
\Cref{thm:YMA9-Hoeffding-ratio,cor:YMA9-two-radius-confidence}.
Items~\textup{(E9.9)--(E9.10)} are
\Cref{thm:YMA9-complete-interval,thm:YMA9-root-alternative}.
\end{proof}

\begingroup\small
\begin{longtable}{>{\raggedright\arraybackslash}p{0.27\textwidth}
                  >{\raggedright\arraybackslash}p{0.20\textwidth}
                  >{\raggedright\arraybackslash}p{0.45\textwidth}}
\caption{Version V9 proof-status ledger.}\label{tab:YMA9-ledger}\\
\toprule
\textbf{Row}&\textbf{Status}&\textbf{Exact scope}\\
\midrule
\endfirsthead
\toprule
\textbf{Row}&\textbf{Status}&\textbf{Exact scope}\\
\midrule
\endhead
Physical action-germ incidence&Proved exactly
&Candidate-minus-physical curvature is the signed direct/mixed/innovation
formula \eqref{eq:YMA9-germ-mismatch}.\\[0.3em]
Two-copy germ Read&Target-minimal exact card
&One same-base pair writer acquires the complete scalar mismatch before termwise
bounds.\\[0.3em]
Native scalar fan&Pruned on escaping branch
&Its absolute stock vanishes and its singlet-to-adjoint linear bridge is zero by
the supplied YM--B V8 theorem.\\[0.3em]
Physical rectangle source&Retained independently
&The YM--B growing-radius Creutz source is not erased and is not substituted
into the shell Hessian.\\[0.3em]
Two-radius secant&Proved fourth-order accurate
&The quartic bias cancels and the remaining error is the explicit sixth-
derivative radius.\\[0.3em]
Sixth cumulant card&Proved exactly
&Eleven partition types give \eqref{eq:YMA9-sixth-cumulant}; connected terms use
oscillation rather than action offsets.\\[0.3em]
Wilson fifth/sixth stocks&Proved explicitly
&For one normalized axial direction, $B_5^{\rm W}\le12\sqrt2\beta/M^6$ and
$B_6^{\rm W}\le24\beta/M^8$.\\[0.3em]
Positive-writer sampling&Finite coverage theorem
&Hoeffding intervals retain their confidence semantics; they are not exact
outcome-wise enclosures.\\[0.3em]
Complete first axial interval&Finite terminating theorem
&Ratio, germ, topological, and reference intervals give PASS, TYPED
OBSTRUCTION, or STOP.\\[0.3em]
Actual selected-root ratio&\stopstatus
&No physical numerical or exact ratio Read was supplied in the attachments.\\[0.3em]
Weak-facing escape&Open trajectory value
&The bounded branch is typed; the escaping branch remains conditional on the
selector.\\[0.3em]
Boundary radial attenuation&Open in YM--B
&The original physical-radius source and its denominator remain independent
selected-root values.\\[0.3em]
Higher jets, contraction, mass&Not inferred
&The inherited downstream handoff opens only after the direct quadratic family
passes.\\
\bottomrule
\end{longtable}
\endgroup

\begin{assessment}[What changed after the supplied notes]
\label{ass:YMA9-change}
The new boundary note removes one false escape route: the scalar native fan
cannot be normalized into the adjoint YM--A mode source without a new
calibration and a new colour-carrying occurrence.  The Navier--Stokes and
Einstein--Standard--Model notes independently show why a later positive Gram
cannot repair a missing physical route or a wrong deformation tangent.  V9
therefore places one exact action-germ incidence number before the partition
Read.  It then uses the two radii already requested by V8 to improve the same
direct scalar rather than opening another boundary or current space.
\end{assessment}

\begin{openproblem}[V10: execute the actual adjoint action germ and dyadic partition pair]
\label{op:YMA9-V10}
Preserve Version~V9 and every frozen physical choice.  Proceed in this order.
\begin{enumerate}[label=\textup{(V10.\arabic*)},leftmargin=5.5em]
\item Decide the weak-facing branch.  On a bounded cofinal subsequence export
      the rough-Haar/contact trajectory obstruction.  On an escaping branch
      retain the first predeclared shell and do not use the native scalar fan.
\item Write the complete physical family
      $F_{M,q,X}^{\rm phys}(t,z)$ of
      \eqref{eq:YMA9-actual-action-family}, including every generated,
      marginal, support, and topological term.  Verify the gauge-orbit held-out.
\item If a ratio instrument uses any candidate family $G\ne F$, acquire the
      single two-copy mismatch writer
      \eqref{eq:YMA9-germ-replica-writer}.  Otherwise freeze the exact zero
      correction \eqref{eq:YMA9-zero-germ-correction}.
\item Choose a dyadic pair $\varepsilon,\varepsilon/2$ before sampling or
      integration.  Acquire outward positive intervals for both Version~V8
      ratio writers on one common base cylinder.
\item Enumerate $B_1,\ldots,B_6$, or prove a sharper connected sixth-derivative
      bound, and acquire the same-carrier topological interval.  Apply
      \eqref{eq:YMA9-complete-margin}.
\item If the direct margin passes, continue to the next predeclared orbit and
      then the fused higher jet.  If it fails, freeze the attained physical
      polarization and only then open the V7 ancestry/current attribution.  If
      the interval remains undecided, tighten these same scalar Reads.
\item Run the YM--B physical-radius card independently.  Its result may be
      transported alongside the shell card but is not used as its deformation
      tangent or reference denominator.
\end{enumerate}
If the actual unnormalized action family or its ratio writer remains absent,
the exact output is \stopstatus\ at that loading row, not another abstract
shell theorem.
\end{openproblem}

\begin{assessment}[V9 termination]
\label{ass:YMA9-termination}
The action-germ mismatch identity, its two-copy acquisition, scalar-proxy
elimination, two-radius fourth-order curvature formula, sixth-cumulant bound,
Wilson fifth/sixth insertion stocks, statistical coverage theorem, and complete
physical interval are \pass\ at their stated finite mathematical scopes.  The
selected response-matched symbol remains \stopstatus: no actual selected-root
partition-ratio or topological interval was supplied.  No negative physical
mode has been produced.  The next iteration must execute the actual adjoint
action germ and its two positive ratio Reads before returning to current or
boundary attribution.
\end{assessment}

\section*{Version V9 append-only continuation record}
\addcontentsline{toc}{section}{Version V9 append-only continuation record}
The complete Version~V8 source above is preserved byte for byte up to its sole
final \verb|\end{document}| command.  Its SHA--256 digest is
\begin{center}\scriptsize\ttfamily
fcd91052363c1ccc722ae30da911680aa6b95838c08cb097e971b92998cdf07a.
\end{center}
For Version~V10, preserve this accumulated source, append new material
immediately before the final document-closing command, and increment the
filename to \texttt{\_V10.tex}.  Keep the physical response, weak-facing
selector, exact block map, parity-tree section, rooted Hodge carrier, first
axial mode, colour normalization, response coefficient, action-germ family,
and local/topological support convention fixed before reading any ratio,
score, ancestry, radial, or transport value.

% =============================================================================
% END APPENDED VERSION V9 MATERIAL
% =============================================================================


% =============================================================================
% START APPENDED VERSION V10 MATERIAL
% =============================================================================

\clearpage
\hypersetup{
 pdftitle={Renewal Yang--Mills Reduced-Fibre Wilson Shell Symbol and Local-Vacuum Programme V10},
 pdfsubject={Exact finite Wilson-Haar evaluation, selected-root isolation, and the separation of finite computability from cofinal control}
}
\pdfinfo{
 /Title (Renewal Yang--Mills Reduced-Fibre Wilson Shell Symbol and Local-Vacuum Programme V10)
 /Subject (Exact finite Wilson-Haar evaluation, selected-root isolation, and the separation of finite computability from cofinal control)
}
\providecommand{\PoisTail}{\mathsf T}
\providecommand{\Inv}{\operatorname{Inv}}

\section{Version V10: the finite Wilson card is specified, and its exact Haar evaluation}
\label{sec:YMA10-direction}

Version~V9 is preserved above.  It correctly proved that normalized fibre
laws and their centred score descendants do not determine the absolute
partition curvature, and it supplied a positive two-radius acquisition when
one has access to an unnormalized ratio writer.  Its terminal assessment,
however, grouped two different failures under the phrase ``the ratio is
absent.''  They must now be separated:
\begin{equation}
 \boxed{
 \begin{array}{c}
 \text{the finite unnormalized action is not specified}\\[0.2em]
 \neq\quad
 \text{the specified finite Haar integral has not been evaluated}\\[0.2em]
 \neq\quad
 \text{no cutoff-uniform estimate is known.}
 \end{array}}
 \label{eq:YMA10-three-levels}
\end{equation}
The first is a physical loading defect.  The second is a finite calculation.
The third is the cofinal Yang--Mills estimate.

For the literal pure Wilson shell and for every finite tower obtained from it
by the inherited exact Haar pushforwards, the unnormalized cylinder is already
specified mathematically.  Exact shell associativity permits all generated
interactions and support births to be evaluated by returning to one finest
Wilson cylinder.  No independently observed partition ratio is needed to
define that finite number.  What remained absent in V9 was an executed
certificate, not the finite object itself.

The present continuation gives such a certificate architecture without
introducing a new shell carrier.  It proves:
\begin{enumerate}[label=\textup{(A10.\arabic*)},leftmargin=5.3em]
\item a direct finest-lattice representation of the selected shell, including
      every exactly generated interaction;
\item an absolutely convergent Wilson--Haar moment expansion with an explicit
      Poisson-tail error for the partition function, Wilson-loop numerators,
      and the exact source Hessian numerator;
\item an exact finite invariant-tensor algorithm for every \(\SU(3)\) Haar
      coefficient;
\item a logarithm-free selected-root equation whose weak-facing root has a
      finite complex-zero certificate;
\item a sign-safe interval for the physical Hessian directly at the selected
      root, with no finite-difference bias and no separately estimated
      order-\(\beta\) cancellation;
\item a marked connected-coefficient recursion which removes every vacuum
      component before source attribution;
\item a finite strict-margin decision theorem and an explicit explanation of
      why the resulting brute-force expansion is not a cofinal estimate.
\end{enumerate}

Thus V10 goes upstream of the unexecuted ratio Read, but it does not pretend
that finite computability supplies a practical or cutoff-uniform proof.  The
next genuinely analytic problem is to compress the same exact coefficient
system by a source-connected low-temperature expansion, or to produce an
equivalent Wilson-specific uniform estimate.

\begin{firewall}[No replacement of the selected action]
\label{fire:YMA10-no-action-replacement}
Every coefficient below is computed from the already frozen Wilson action,
block map, response root, coarse-link deformation, and Haar measure.  A
coarse-only marginal term is retained as an explicit additive curvature.  An
additional phenomenological or topological term not specified as a finite
writer remains an unloaded input and is not manufactured by the Haar
algorithm.  The construction is therefore an evaluator of the declared card,
not a new model and not a fitted surrogate.
\end{firewall}

% =============================================================================
\section{One finest-lattice cylinder contains every exactly generated shell term}
\label{sec:YMA10-flattening}
% =============================================================================

Let \(X_L\) be a finite finest-link configuration space with product Haar
probability \(m_L\), and let
\begin{equation}
 \mathbf B_{L\to r}:X_L\longrightarrow X_r
 \label{eq:YMA10-composed-block}
\end{equation}
be the composition of the already frozen exact dyadic block maps.  Let
\(\Phi_L\) be the complete action on that finest card.  It may contain the
pure Wilson term and any finite analytic terms which have actually been
inserted before the trace.  For a coarse configuration \(W\in X_r\), use the
inherited global or relative-tree Haar trivialization to write
\begin{equation}
 \mathcal Z_{L\to r}(W)
 =\int_{Y_{L\to r}}
   e^{-F_{L\to r}(z,W)}\,\dd m_{L\to r}(z),
 \qquad
 \mathcal S_{L\to r}(W)=-\log\mathcal Z_{L\to r}(W).
 \label{eq:YMA10-direct-cylinder}
\end{equation}

\begin{proposition}[Exact flattening of a finite generated-action tower]
\label{prop:YMA10-flattening}
Suppose the level-\(r\) action is obtained from \(\Phi_L\) by any finite
sequence of the inherited exact Haar pushforwards.  Then its unnormalized
coarse density is exactly \eqref{eq:YMA10-direct-cylinder}.  Direct and staged
integrations agree, and every local, quasilocal, or support-birth term created
by an intermediate logarithm is already contained in this one finest-lattice
integral.

If explicit coarse-only terms \(C_j\) were inserted at intermediate levels,
then after pullback to \(X_L\),
\begin{equation}
 \Phi_L^{\rm tot}(U)
 =\Phi_L(U)+\sum_j C_j(\mathbf B_{L\to j}U),
 \label{eq:YMA10-pulled-coarse-terms}
\end{equation}
and the same direct formula applies.  If a term depends only on the final
coarse variable, \(C_r(W)\), it factors as
\begin{equation}
 \boxed{
 \mathcal Z_{L\to r}^{\rm tot}(W)
 =e^{-C_r(W)}\mathcal Z_{L\to r}(W),
 \qquad
 \operatorname{Hess}\mathcal S^{\rm tot}
 =\operatorname{Hess}\mathcal S+
   \operatorname{Hess}C_r.}
 \label{eq:YMA10-coarse-factor}
\end{equation}
\end{proposition}

\begin{proof}
The inherited product-Haar fibre theorem disintegrates \(m_L\) into the
coarse Haar variable and a normalized fibre Haar variable at every stage.
Applying Fubini to two consecutive stages shows that their iterated fibre
integral is the integral over the fibre of the composed block map.  Induction
proves the first assertion.  Intermediate generated interactions are the
logarithm of partial integrals of the same positive density, so reinserting
them and integrating the remaining variables returns the direct integral.
Pulling an inserted \(C_j\) back through the deterministic block map gives
\eqref{eq:YMA10-pulled-coarse-terms}.  A final coarse-only term is independent
of the fibre coordinate and factors out.  Taking minus the logarithm and two
coarse derivatives proves \eqref{eq:YMA10-coarse-factor}.
\end{proof}

\begin{assessment}[What is already specified]
\label{ass:YMA10-loaded-scope}
For a fixed cutoff, a fixed selected root \(\beta_r\), and the literal Wilson
fine action, \eqref{eq:YMA10-direct-cylinder} is an unnormalized physical
shell cylinder.  An explicit list of all coefficients in its generated
coarse action is not an additional entrance requirement.  Conversely, if a
response-matched term is only named but neither given as a writer nor derived
by an exact pushforward, then that term is still absent.  V10 closes only the
former finite-evaluation issue.
\end{assessment}

% =============================================================================
\section{An absolute Wilson--Haar expansion with an explicit tail}
\label{sec:YMA10-Haar-expansion}
% =============================================================================

We first treat one literal pure Wilson fibre.  The flattened version is
identical after replacing the plaquette set by the finest one and pulling the
coarse path through the composed block map.

Let \(\mathscr P\) be the finite plaquette set and put
\begin{equation}
 P:=|\mathscr P|,
 \qquad
 A_W(z):=\frac13\sum_{p\in\mathscr P}
                 \operatorname{Re}\Tr U_p(z,W),
 \qquad
 V_W(z)=P-A_W(z).
 \label{eq:YMA10-AW}
\end{equation}
For \(\SU(3)\),
\begin{equation}
 -\frac12\le\frac13\operatorname{Re}\Tr U\le1,
 \qquad U\in\SU(3),
 \label{eq:YMA10-trace-range}
\end{equation}
so \(\abs{A_W(z)}\le P\).  The reduced-fibre partition function is
\begin{equation}
 Z_\beta(W)
 =e^{-\beta P}\int e^{\beta A_W(z)}\,\dd m_{\rm red}(z).
 \label{eq:YMA10-Z}
\end{equation}

For \(x\ge0\), define the upper Poisson tail
\begin{equation}
 \PoisTail_K(x):=e^{-x}\sum_{n=K+1}^{\infty}\frac{x^n}{n!}.
 \label{eq:YMA10-Poisson-tail}
\end{equation}

\begin{theorem}[Uniform finite-card Wilson moment expansion]
\label{thm:YMA10-Wilson-moment-expansion}
Define the exact moments
\begin{equation}
 \mu_n(W)=\int A_W(z)^n\,\dd m_{\rm red}(z)
 \label{eq:YMA10-moments}
\end{equation}
and the degree-\(K\) approximation
\begin{equation}
 Z_{\beta,K}(W)
 :=e^{-\beta P}\sum_{n=0}^{K}
     \frac{\beta^n}{n!}\mu_n(W).
 \label{eq:YMA10-ZK}
\end{equation}
Then, for every \(\beta\ge0\) and every coarse field \(W\),
\begin{equation}
 \boxed{
 \abs{Z_\beta(W)-Z_{\beta,K}(W)}
 \le \PoisTail_K(\beta P)
 \le\frac{(\beta P)^{K+1}}{(K+1)!}.}
 \label{eq:YMA10-Z-tail}
\end{equation}
The bound is uniform in the fibre coordinate, background, toron sector, and
coarse deformation.

More generally, if \(O\) is any bounded insertion with
\(\norm{O}_\infty\le C_O\), then
\begin{equation}
 N_{O,\beta}(W)
 :=\int O(z,W)e^{-\beta V_W(z)}\,\dd m_{\rm red}(z)
 \label{eq:YMA10-inserted-numerator}
\end{equation}
has the corresponding truncated moment expansion with error at most
\begin{equation}
 \boxed{C_O\PoisTail_K(\beta P).}
 \label{eq:YMA10-insertion-tail}
\end{equation}
\end{theorem}

\begin{proof}
For real \(x\), Taylor's theorem and the exponential series give
\begin{equation*}
 \left|e^x-\sum_{n=0}^{K}\frac{x^n}{n!}\right|
 \le\sum_{n=K+1}^{\infty}\frac{|x|^n}{n!}.
\end{equation*}
Set \(x=\beta A_W(z)\), multiply by \(e^{-\beta P}\), and use
\(\abs{A_W}\le P\).  The result is bounded by
\[
 e^{-\beta P}\sum_{n>K}\frac{(\beta P)^n}{n!}
 =\PoisTail_K(\beta P).
\]
Integration proves the first inequality.  Since
\(\sum_{n>K}x^n/n!\le e^x x^{K+1}/(K+1)!\), the second follows.  Multiplying
the pointwise remainder by \(\abs O\le C_O\) proves
\eqref{eq:YMA10-insertion-tail}.
\end{proof}

\begin{corollary}[A priori positive denominator window]
\label{cor:YMA10-Z-positive-window}
For real \(\beta\ge0\),
\begin{equation}
 \boxed{
 e^{-3\beta P/2}\le Z_\beta(W)\le1.}
 \label{eq:YMA10-Z-positive-window}
\end{equation}
Thus the denominator of the exact Hessian ratio has a certified positive
lower endpoint before any moment coefficient is evaluated.
\end{corollary}

\begin{proof}
Each plaquette defect belongs to \([0,3/2]\) by
\eqref{eq:YMA10-trace-range}.  Hence
\(0\le V_W\le3P/2\) pointwise.  Integrate
\(e^{-3\beta P/2}\le e^{-\beta V_W}\le1\) against normalized Haar
measure.
\end{proof}

\begin{corollary}[Complex-disk polynomial enclosure]
\label{cor:YMA10-complex-disk}
Put \(V_{\max}=3P/2\).  For a bounded insertion \(O\), define
\begin{equation}
 \theta_n(O):=\int O(z,W)V_W(z)^n\,\dd m_{\rm red}(z).
 \label{eq:YMA10-V-moments}
\end{equation}
Then, for every complex \(\beta\) with \(\abs\beta\le R\),
\begin{equation}
 \boxed{
 \left|N_{O,\beta}(W)-
       \sum_{n=0}^{K}\frac{(-\beta)^n}{n!}\theta_n(O)\right|
 \le C_Oe^{RV_{\max}}
       \frac{(RV_{\max})^{K+1}}{(K+1)!}.}
 \label{eq:YMA10-complex-tail}
\end{equation}
Every coefficient on the right is an exact Haar contraction.  This is the
complex-uniform enclosure used by the root-counting argument below; the
sharper Poisson tail \eqref{eq:YMA10-insertion-tail} is retained on the
positive real axis.
\end{corollary}

\begin{proof}
Expand \(e^{-\beta V_W}\) directly.  Since
\(\abs{\beta V_W}\le RV_{\max}\), the omitted exponential series is
bounded by
\(e^{RV_{\max}}(RV_{\max})^{K+1}/(K+1)!\).  Multiply by
\(\abs O\le C_O\), integrate, and use
\Cref{thm:YMA10-exact-Haar-algorithm} after expanding
\(V_W^n=(P-A_W)^n\).
\end{proof}

\begin{corollary}[A simple explicit truncation depth]
\label{cor:YMA10-simple-depth}
If \(K+1\ge2e\beta P\), then
\begin{equation}
 \boxed{
 \PoisTail_K(\beta P)\le2^{-(K+1)}.}
 \label{eq:YMA10-simple-tail}
\end{equation}
Thus every fixed finite Wilson integral is determined to absolute error
\(\epsilon\) by a finite degree satisfying, for example,
\begin{equation}
 K+1\ge
 \max\left\{2e\beta P,\frac{\log(1/\epsilon)}{\log2}\right\}.
 \label{eq:YMA10-simple-K}
\end{equation}
\end{corollary}

\begin{proof}
The factorial bound \((K+1)!\ge((K+1)/e)^{K+1}\), inserted in
\eqref{eq:YMA10-Z-tail}, gives
\[
 \PoisTail_K(\beta P)
 \le\left(\frac{e\beta P}{K+1}\right)^{K+1}.
\]
The stated depth makes the parenthesis at most \(1/2\).
\end{proof}

\begin{firewall}[Finite convergence is not a useful cofinal norm]
\label{fire:YMA10-global-series-cost}
For one dyadic shell, \(P=96M^4\).  The sufficient global degree in
\eqref{eq:YMA10-simple-K} is therefore of order \(\beta M^4\), before the
number of trace monomials is counted.  This proves finite determination but
not computational feasibility, a polynomial-time evaluator, a uniform
cluster norm, or a cofinal symbol margin.  Vacuum-volume cancellation must be
used before this expansion becomes an analytic continuum tool.
\end{firewall}

% =============================================================================
\section{Every coefficient is an exact finite SU(3) invariant-tensor contraction}
\label{sec:YMA10-invariant-tensors}
% =============================================================================

Let \(V=\C^3\) be the defining representation.  For one reduced link and
nonnegative integers \(p,q\), put
\begin{equation}
 \mathcal P_{p,q}
 :=\int_{\SU(3)}
   U^{\otimes p}\otimes\overline U^{\otimes q}\,\dd U
 \quad\text{on}\quad
 V^{\otimes p}\otimes\overline V^{\otimes q}.
 \label{eq:YMA10-Haar-projector}
\end{equation}

\begin{lemma}[Haar integration is the invariant orthogonal projection]
\label{lem:YMA10-Haar-projector}
The operator \(\mathcal P_{p,q}\) is the orthogonal projection onto
\begin{equation}
 \Inv_{p,q}
 :=\left(V^{\otimes p}\otimes\overline V^{\otimes q}\right)^{\SU(3)}.
 \label{eq:YMA10-invariant-space}
\end{equation}
In particular,
\begin{equation}
 p-q\not\equiv0\pmod3
 \quad\Longrightarrow\quad
 \mathcal P_{p,q}=0.
 \label{eq:YMA10-centre-selection}
\end{equation}
\end{lemma}

\begin{proof}
Left or right multiplication of the integral by the tensor representation
leaves it unchanged by Haar invariance.  Hence its range is invariant and it
acts as the identity on invariant vectors.  Self-adjointness follows from
inversion invariance, and idempotence follows by applying Haar invariance to
two independent group variables.  Thus it is the orthogonal invariant
projection.  The central element \(\zeta I\), \(\zeta^3=1\), acts by
\(\zeta^{p-q}\); a nonzero invariant vector therefore requires
\(p-q\equiv0\pmod3\).
\end{proof}

\begin{theorem}[Finite exact \(\delta\)--\(\varepsilon\) evaluation algorithm]
\label{thm:YMA10-exact-Haar-algorithm}
For every \(p,q\), enumerate the contraction tensors obtained from copies of
\begin{equation}
 \delta^i_j,\qquad
 \varepsilon_{ijk},\qquad
 \varepsilon^{ijk}
 \label{eq:YMA10-delta-epsilon}
\end{equation}
with the prescribed \(p\) fundamental and \(q\) antifundamental free slots.
Let \(C_{p,q}\) be the matrix whose columns are the resulting tensors after
duplicates are allowed, and let \(G_{p,q}=C_{p,q}^*C_{p,q}\).  Then
\begin{equation}
 \boxed{
 \mathcal P_{p,q}
 =C_{p,q}G_{p,q}^{\dagger}C_{p,q}^*.}
 \label{eq:YMA10-projection-formula}
\end{equation}
All entries of this projector are rational numbers.  Consequently every
coefficient \(\mu_n(W)\), and every coefficient of a bounded Wilson-loop or
source-derivative numerator, is obtained by a finite sequence of exact
rational operations and contractions with the fixed coarse matrices appearing
in \(W\).

If those coarse matrices have algebraic entries, the coefficient is
algebraic and exactly representable.  For arbitrary computable entries, an
outward interval of any prescribed width is obtained by interval evaluation
of only those final matrix entries; the Haar part remains exact.
\end{theorem}

\begin{proof}
The first fundamental invariant theorem for the defining and dual
\(\mathrm{SL}_3\)-modules says that the tensors in
\eqref{eq:YMA10-delta-epsilon} span \(\Inv_{p,q}\).  The compact real form
\(\SU(3)\) has the same invariant tensors after complexification.  For any
matrix \(C\), the operator \(C(C^*C)^\dagger C^*\) is the orthogonal
projection onto \(\Ran C\), giving \eqref{eq:YMA10-projection-formula}.
The contraction tensors have entries in \(\{0,\pm1\}\), so their Gram matrix
has integer entries.  A Moore--Penrose inverse of a rational matrix is
rational: choose a rational rank factorization and use the ordinary inverses
of its nonsingular rational Gram factors.  Hence the projection is rational.

Expanding \(A_W^n\) produces finitely many products of traces of ordered
plaquette words and their conjugates.  Expanding each trace into matrix
entries gives a finite tensor monomial in every reduced-link variable.
Integrate one link at a time using \eqref{eq:YMA10-projection-formula}.  The
remaining coefficients are polynomials with rational coefficients in the
fixed coarse matrix entries.  The same argument applies after inserting a
Wilson loop or differentiating a fixed coarse-link path finitely many times.
\end{proof}

\begin{corollary}[Exact low-degree checks]
\label{cor:YMA10-low-Haar-checks}
The first two nontrivial one-link projectors are
\begin{align}
 \int U_{ij}\overline{U_{k\ell}}\,\dd U
 &=\frac13\delta_{ik}\delta_{j\ell},
 \label{eq:YMA10-two-Haar}\\
 \int U_{i_1j_1}U_{i_2j_2}U_{i_3j_3}\,\dd U
 &=\frac16\varepsilon_{i_1i_2i_3}
             \varepsilon_{j_1j_2j_3}.
 \label{eq:YMA10-three-Haar}
\end{align}
Writing \(T=\Tr U\), one obtains
\begin{equation}
 \boxed{
 \int\left(\frac13\operatorname{Re}T\right)^2\dd U=\frac1{18},
 \qquad
 \int\left(\frac13\operatorname{Re}T\right)^3\dd U=\frac1{108}.}
 \label{eq:YMA10-low-trace-moments}
\end{equation}
\end{corollary}

\begin{proof}
The invariant lines in \(V\otimes\overline V\) and \(V^{\otimes3}\) are
spanned by \(\delta\) and \(\varepsilon\), whose squared norms are three and
six.  This gives \eqref{eq:YMA10-two-Haar}--\eqref{eq:YMA10-three-Haar}.
Centre selection gives
\(\int T^2=\int T^2\overline T=0\), while
\(\int T\overline T=1\) and \(\int T^3=\int\overline T^3=1\).
Expanding \((T+\overline T)^2/4\) and
\((T+\overline T)^3/8\), followed by the factor \(3^{-n}\), proves
\eqref{eq:YMA10-low-trace-moments}.
\end{proof}

% =============================================================================
\section{Direct exact source curvature without a secant bias}
\label{sec:YMA10-direct-Hessian}
% =============================================================================

Let \(W(t)\) be the frozen real axial coarse-link path and put
\begin{equation}
 A_t(z):=A_{W(t)}(z),
 \qquad
 A=A_0,
 \qquad
 \dot A=\partial_tA_t|_{t=0},
 \qquad
 \ddot A=\partial_t^2A_t|_{t=0}.
 \label{eq:YMA10-A-derivatives}
\end{equation}
The central translation and colour symmetries give
\begin{equation}
 Z_\beta'(0)=0
 \label{eq:YMA10-first-zero}
\end{equation}
for the nontrivial axial character.  Define
\begin{equation}
 L_1:=\sup_z|\dot A(z)|,
 \qquad
 L_2:=\sup_z|\ddot A(z)|.
 \label{eq:YMA10-L12}
\end{equation}

\begin{theorem}[Exact source-Hessian moment series and Poisson tail]
\label{thm:YMA10-Hessian-series}
For \(n\ge0\), define the exact marked Haar moments
\begin{equation}
 \nu_n:=\int A^n\ddot A\,\dd m_{\rm red},
 \qquad
 \xi_n:=\int A^n(\dot A)^2\,\dd m_{\rm red}.
 \label{eq:YMA10-marked-moments}
\end{equation}
Then
\begin{equation}
 \boxed{
 \begin{aligned}
 Z_\beta''(0)
 =e^{-\beta P}\biggl[
  &\beta\sum_{n=0}^{\infty}\frac{\beta^n}{n!}\nu_n\\
  &+\beta^2\sum_{n=0}^{\infty}\frac{\beta^n}{n!}\xi_n
 \biggr].
 \end{aligned}}
 \label{eq:YMA10-Zpp-series}
\end{equation}
If both series are truncated after \(n=K\), the total error obeys
\begin{equation}
 \boxed{
 |R_{K}^{(2)}|
 \le
 \bigl(\beta L_2+\beta^2L_1^2\bigr)
 \PoisTail_K(\beta P).}
 \label{eq:YMA10-Zpp-tail}
\end{equation}
At the symmetric background the exact physical directional Hessian is
\begin{equation}
 \boxed{
 H_{\beta,M}[aX,aX]
 =-\frac{Z_\beta''(0)}{Z_\beta(0)}.}
 \label{eq:YMA10-H-from-Z}
\end{equation}
Every coefficient in \eqref{eq:YMA10-Zpp-series} is evaluated by
\Cref{thm:YMA10-exact-Haar-algorithm}.
\end{theorem}

\begin{proof}
Differentiate \eqref{eq:YMA10-Z} twice under the finite compact integral:
\[
 Z_\beta''(0)
 =e^{-\beta P}\int e^{\beta A}
   \bigl(\beta\ddot A+\beta^2(\dot A)^2\bigr)\,\dd m_{\rm red}.
\]
Expanding \(e^{\beta A}\) proves \eqref{eq:YMA10-Zpp-series}.  The first
omitted series is bounded by
\[
 \beta L_2e^{-\beta P}
 \sum_{n>K}\frac{(\beta P)^n}{n!}
 =\beta L_2\PoisTail_K(\beta P),
\]
and the second in the same way, proving \eqref{eq:YMA10-Zpp-tail}.
For \(\mathcal S=-\log Z\),
\[
 \mathcal S''=-Z''/Z+(Z'/Z)^2.
\]
Use \eqref{eq:YMA10-first-zero}.  The coefficient claim follows from the
finite tensor expansion of \(A^n\ddot A\) and \(A^n(\dot A)^2\).
\end{proof}

\begin{corollary}[Literal axial derivative stocks]
\label{cor:YMA10-axial-L12}
For the one-direction normalized real axial mode used in Versions~V3--V9,
\begin{equation}
 \boxed{L_1\le6M^2,
 \qquad L_2\le6.}
 \label{eq:YMA10-axial-L12}
\end{equation}
Hence
\begin{equation}
 \boxed{
 |R_K^{(2)}|
 \le\bigl(6\beta+36\beta^2M^4\bigr)
       \PoisTail_K(96\beta M^4)}
 \label{eq:YMA10-axial-Zpp-tail}
\end{equation}
for one dyadic shell.
\end{corollary}

\begin{proof}
For \(F=\beta(P-A_t)\), the V9 Wilson insertion bounds give
\(\sup|\partial_tF|\le6\beta M^2\) and
\(\sup|\partial_t^2F|\le6\beta\).  Divide by \(\beta\) and use
\(P=96M^4\) in \eqref{eq:YMA10-Zpp-tail}.
\end{proof}

\begin{corollary}[Sign-safe Hessian interval]
\label{cor:YMA10-H-interval}
Suppose exact arithmetic and the tails above give
\begin{equation}
 0<z_-\le Z_\beta(0)\le z_+,
 \qquad
 d_-\le Z_\beta''(0)\le d_+.
 \label{eq:YMA10-Z-Zpp-interval}
\end{equation}
Then
\begin{equation}
 \boxed{
 \min_{d\in\{d_-,d_+\},\ z\in\{z_-,z_+\}}
       \left(-\frac dz\right)
 \le H_{\beta,M}[aX,aX]\le
 \max_{d\in\{d_-,d_+\},\ z\in\{z_-,z_+\}}
       \left(-\frac dz\right).}
 \label{eq:YMA10-H-corners}
\end{equation}
The corner formula remains valid if the interval for \(Z''\) crosses zero.
\end{corollary}

\begin{proof}
The function \((d,z)\mapsto-d/z\) is continuous on the compact rectangle
with \(z>0\), and is affine in \(d\) and monotone in \(z\) on each fixed-sign
half.  Its extrema are therefore attained at corners.  Apply
\eqref{eq:YMA10-H-from-Z}.
\end{proof}

\begin{assessment}[Relation to the V8--V9 ratio card]
\label{ass:YMA10-ratio-relation}
The positive two-radius partition ratio remains the preferable physical or
statistical instrument when direct source derivatives are unavailable.  The
Wilson--Haar evaluator has the complete action formula and therefore computes
\(Z\) and \(Z''\) directly.  It has no finite-difference bias, needs no sixth
derivative majorant, and keeps the direct second insertion and paired first
insertions inside one exact numerator.  These are two evaluations of the same
curvature, not two source charges.
\end{assessment}

% =============================================================================
\section{The selected response root without logarithms or normalized denominators}
\label{sec:YMA10-root}
% =============================================================================

Let \(O_{\epsilon\eta}\) be the four frozen normalized Wilson-loop writers
used in the Creutz response, and put
\begin{equation}
 N_{\epsilon\eta,N}(\beta)
 :=\int O_{\epsilon\eta}(U)e^{-\beta V_N(U)}\,\dd m_N(U).
 \label{eq:YMA10-loop-numerators}
\end{equation}
Their common partition denominator cancels from the response ratio.  For a
fixed response value \(u>0\), define the entire real function
\begin{equation}
 \boxed{
 D_{N,u}(\beta)
 :=N_{11,N}(\beta)N_{00,N}(\beta)
   -e^{-u}N_{10,N}(\beta)N_{01,N}(\beta).}
 \label{eq:YMA10-root-function}
\end{equation}

\begin{proposition}[Logarithm-free root and orientation]
\label{prop:YMA10-root-equivalence}
On every interval on which the four loop expectations are positive,
\begin{equation}
 \boxed{
 \chi_N(\beta)=u
 \iff D_{N,u}(\beta)=0.}
 \label{eq:YMA10-root-equivalence}
\end{equation}
Moreover, a stable down-crossing of \(\chi_N-u\) is exactly a
negative-to-positive sign crossing of \(D_{N,u}\).
Every value and derivative of \(D_{N,u}\) on a compact \(\beta\)-interval has
an arbitrarily sharp outward interval obtained from the moment expansion
\eqref{eq:YMA10-insertion-tail}.
\end{proposition}

\begin{proof}
Writing \(W_{\epsilon\eta}=N_{\epsilon\eta}/Z\), the two factors of \(Z\)
cancel:
\[
 e^{-\chi_N}
 =\frac{W_{11}W_{00}}{W_{10}W_{01}}
 =\frac{N_{11}N_{00}}{N_{10}N_{01}}.
\]
Positivity of the denominator proves the zero equivalence.  Immediately to
the left of a stable down-crossing one has \(\chi_N>u\), hence
\(e^{-\chi_N}<e^{-u}\) and \(D_{N,u}<0\); the sign reverses on the right.
The final assertion is \Cref{thm:YMA10-Wilson-moment-expansion} applied to the
bounded loop insertions and to their \(\beta\)-derivatives, which merely add
bounded powers of \(V_N\).
\end{proof}

\begin{lemma}[A certified all-coupling weak-endpoint corridor]
\label{lem:YMA10-tail-corridor}
Let \(O_1,\ldots,O_4\) be the four normalized real Wilson-loop writers.
For \(0<\eta<1\), put
\begin{equation}
 \mathscr K_\eta
 :=\{U:\min_iO_i(U)\le1-\eta\}.
 \label{eq:YMA10-bad-loop-set}
\end{equation}
Every flat configuration lies outside \(\mathscr K_\eta\).  If this set is
nonempty, compactness gives
\begin{equation}
 \delta_\eta:=\min_{\mathscr K_\eta}V_N>0.
 \label{eq:YMA10-bad-action-gap}
\end{equation}
If it is empty, the loop floor below is immediate and the exponential term
may be omitted.  On the nonempty branch, choose a positive-Haar-measure set \(\mathscr V_\eta\) on which
\(V_N\le\delta_\eta/2\), and write
\(m_\eta=m_N(\mathscr V_\eta)>0\).  Then, for all \(\beta\ge0\),
\begin{equation}
 \boxed{
 \mu_{N,\beta}(\mathscr K_\eta)
 \le m_\eta^{-1}e^{-\beta\delta_\eta/2},}
 \label{eq:YMA10-tail-probability}
\end{equation}
and each selected loop expectation satisfies
\begin{equation}
 \boxed{
 W_i(\beta)\ge
 1-\eta-\frac32m_\eta^{-1}e^{-\beta\delta_\eta/2}.}
 \label{eq:YMA10-loop-tail-floor}
\end{equation}
If \(B\) is chosen so that
\begin{equation}
 s_B:=\eta+\frac32m_\eta^{-1}e^{-B\delta_\eta/2}
 <1-e^{-u/4},
 \label{eq:YMA10-tail-B-condition}
\end{equation}
then
\begin{equation}
 \boxed{\abs{\chi_N(\beta)}<u
 \quad\text{for every }\beta\ge B.}
 \label{eq:YMA10-global-tail-corridor}
\end{equation}
Rational lower bounds for \(\delta_\eta,m_\eta\) and a rational
\(B\) satisfying \eqref{eq:YMA10-tail-B-condition} are finite endpoint
witnesses.
\end{lemma}

\begin{proof}
On \(\mathscr K_\eta\), the Wilson action is at least
\(\delta_\eta\).  The denominator of the Gibbs law is bounded below by
\(m_\eta e^{-\beta\delta_\eta/2}\), while its numerator on
\(\mathscr K_\eta\) is at most \(e^{-\beta\delta_\eta}\).  This
gives \eqref{eq:YMA10-tail-probability}.  Every normalized real
\(\SU(3)\) trace is at least \(-1/2\).  Splitting the loop expectation
over \(\mathscr K_\eta\) and its complement gives
\eqref{eq:YMA10-loop-tail-floor}.  The upper bound \(W_i\le1\) and
\eqref{eq:YMA10-tail-B-condition} imply
\(1-s_B\le W_i(\beta)\le1\) for all \(\beta\ge B\).  The Creutz
response is a signed sum of four logarithms, so
\[
 \abs{\chi_N(\beta)}
 \le4[-\log(1-s_B)]<u.
\]
On the nonempty branch, compactness gives \(\delta_\eta>0\): the zero
set of \(V_N\) consists of flat configurations, on which every selected contractible loop equals
one.  A sufficiently small product neighborhood of one flat configuration
supplies \(\mathscr V_\eta\) and a positive Haar-volume witness.
\end{proof}

\begin{definition}[Finite complex-zero certificate for the weak-facing root]
\label{def:YMA10-root-certificate}
A weak-facing root certificate on \([b_0,B]\subset(0,\infty)\) consists of:
\begin{enumerate}[label=\textup{(R\arabic*)},leftmargin=3.8em]
\item outward analytic intervals proving positivity of the four loop
      numerators on the interval;
\item disjoint rational rectangles \(Q_1,\ldots,Q_s\subset\C\), invariant as
      a family under conjugation, whose interiors contain all zeros of
      \(D_{N,u}\) meeting \([b_0,B]\);
\item for each \(Q_j\), a certified nonzero boundary image and the integer
      winding number of \(D_{N,u}(\partial Q_j)\) about zero;
\item certified real signs between the real zero boxes, including a negative
      sign near \(b_0\) and a positive sign after the final stable crossing;
\item endpoint witnesses \(\eta,\underline\delta_\eta,
      \underline m_\eta,B\) satisfying the hypotheses of
      \Cref{lem:YMA10-tail-corridor}, and hence the certified corridor
      \(\chi_N(\beta)<u\) for every \(\beta\ge B\).
\end{enumerate}
The largest real box at which the sign changes from negative to positive is
the certified weak-facing root interval.
\end{definition}

\begin{theorem}[Every fixed weak-facing root admits finite certified isolation]
\label{thm:YMA10-root-isolation}
Assume \(u>0\) is a computable real number, the four-loop positive branch is
fixed, and \(D_{N,u}\) is not identically zero.  Every finite card has a root
certificate in the sense of \Cref{def:YMA10-root-certificate}; its
weak-facing interval can be refined to arbitrary width.
\end{theorem}

\begin{proof}
The loop numerators are entire in \(\beta\): the finite action exponential
has an everywhere absolutely convergent power series, and termwise
integration is justified by compactness.  The complex-disk bound in
\Cref{cor:YMA10-complex-disk} gives uniform polynomial enclosures on every
compact rectangle.  The inherited strong endpoint gives a small positive
point with \(\chi_N>u\).  Choose \(\eta\) and then the finite endpoint
witnesses of \Cref{lem:YMA10-tail-corridor}; they certify
\(\chi_N(\beta)<u\) on the complete half-line \(\beta\ge B\), not merely
at one endpoint.

An analytic function which is not identically zero has finitely many zeros in
the compact rectangle between the strong point and \(B\).  Choose pairwise
disjoint rational rectangles around those zeros with zero-free boundaries.
Uniform polynomial enclosure on each boundary eventually becomes sharper
than the positive minimum of \(\abs{D_{N,u}}\) there.  Rouch\'e's theorem
then transfers a certified polynomial winding number to \(D_{N,u}\).  The
number can be obtained by exact rational argument-principle arithmetic when
the coefficients are algebraic, and otherwise by outward complex intervals
whose boundary image excludes zero.  Subdivide the real boxes and refine
until the signs on their complementary real intervals are certified.  The
global tail corridor excludes every later root, so the last negative-to-
positive crossing is exactly the weak-facing root.  Repeating the construction
in a smaller rectangle around that selected zero gives arbitrary refinement.
\end{proof}

\begin{firewall}[The response value is still a frozen scheme input]
\label{fire:YMA10-u-input}
The theorem computes the root of the declared response equation.  It does not
select \(u\), prove that two different response schemes have the same finite
symbol, or make a noncomputable symbolic parameter into a numerical card.  A
numerical execution requires a certified value or interval for the already
predeclared \(u_*\).  This is calibration of the chosen scheme, not a new
Yang--Mills observable.
\end{firewall}

% =============================================================================
\section{Root-box propagation into the exact shell margin}
\label{sec:YMA10-root-to-margin}
% =============================================================================

Let \(I_\beta=[\beta_-,\beta_+]\) be a certified weak-facing root interval.
The moment coefficients in \Cref{thm:YMA10-Hessian-series} do not depend on
\(\beta\).  Evaluating the finite polynomials by outward interval arithmetic
on \(I_\beta\), and bounding every tail at \(\beta_+\), gives intervals.
The latter step is valid because
\(\PoisTail_K'(x)=e^{-x}x^K/K!\ge0\).  Thus
\begin{equation}
 0<z_-\le Z_\beta(0)\le z_+,
 \qquad
 d_-\le Z_\beta''(0)\le d_+
 \quad(\beta\in I_\beta).
 \label{eq:YMA10-root-box-Z}
\end{equation}
Use \Cref{cor:YMA10-H-interval} to obtain
\begin{equation}
 H_-(I_\beta,q)\le H_{\beta,M}(q)\le H_+(I_\beta,q).
 \label{eq:YMA10-root-box-H}
\end{equation}

\begin{theorem}[Complete finite selected-root margin certificate]
\label{thm:YMA10-selected-margin}
Fix the exact axial reference coefficient \(s_0(q)>0\), a target
\(\delta\ge0\), and a same-carrier topological interval
\begin{equation}
 t_{\rm top}(q)\in[t_-,t_+].
 \label{eq:YMA10-top-interval}
\end{equation}
Include every explicit final coarse-only curvature in
\(H_\pm(I_\beta,q)\).  Then the local physical margin obeys
\begin{equation}
 \boxed{
 \begin{aligned}
 m_-&:=H_-(I_\beta,q)-t_+-\delta\beta_+s_0(q),\\
 m_+&:=H_+(I_\beta,q)-t_- -\delta\beta_-s_0(q),
 \end{aligned}}
 \label{eq:YMA10-margin-interval}
\end{equation}
and
\begin{equation}
 \boxed{m_-\le
 h_{\beta,M}^{\rm loc}(q)-\delta\beta s_0(q)
 \le m_+.}
 \label{eq:YMA10-margin-inclusion}
\end{equation}
Therefore
\begin{equation}
 \boxed{
 \begin{array}{rcl}
 m_->0&\Longrightarrow&\pass,\\[0.2em]
 m_+<0&\Longrightarrow&\obstruction,\\[0.2em]
 0\in[m_-,m_+]&\Longrightarrow&\stopstatus.
 \end{array}}
 \label{eq:YMA10-verdict}
\end{equation}
On the obstruction branch, the frozen real axial polarization is an attained
finite physical witness.

If the topological writer is a finite linear combination of trace monomials,
or has a separately supplied uniformly convergent trace expansion with a
certified tail, its numerator and derivatives are evaluated by the same
Haar-moment algorithm.  An arbitrary analytic writer requires its own
computable coefficient or enclosure scheme.  If only its name or support
class is supplied, the full finite-torus Hessian is computable but the declared
local subtraction remains \stopstatus.
\end{theorem}

\begin{proof}
For every \(\beta\in I_\beta\), subtract the upper topological endpoint and
the upper positive reference endpoint to obtain the lower bound, and subtract
the corresponding lower endpoints to obtain the upper bound.  This gives
\eqref{eq:YMA10-margin-inclusion}.  The three sign alternatives are ordinary
interval logic.  All operations are on the one selected action and therefore
introduce no candidate-action mismatch term.  The final paragraph follows because a finite trace-polynomial insertion is
covered by
\Cref{thm:YMA10-Wilson-moment-expansion,thm:YMA10-exact-Haar-algorithm}; a
certified convergent trace expansion follows by adding its own tail.  An
unspecified writer has no coefficients to integrate.
\end{proof}

\begin{corollary}[Finite strict-margin decidability]
\label{cor:YMA10-strict-decision}
On every fixed finite card with computable \(u_*\), a finite trace-polynomial action grammar (or an explicitly certified
computable expansion), explicit local/topological split, and exact nonzero
target margin, repeated
root-box, moment-degree, and interval refinement terminates in either \pass\
or \obstruction.  If the exact margin equals zero, the nested intervals may
converge to zero and the correct mathematical status is \stopstatus.  Failure
of the action, support, root, or positive-denominator domain checks is a
rejected input, not an obstruction verdict.
\end{corollary}

\begin{proof}
\Cref{thm:YMA10-root-isolation} shrinks \(I_\beta\) arbitrarily.  The Poisson
tails tend to zero as \(K\to\infty\), and interval evaluation of the finitely
many exact coefficients may be refined arbitrarily.  Hence the margin interval
shrinks to its exact value.  A nonzero real number has a neighborhood of fixed
sign, so a strict branch is eventually certified.  At equality no finite
strict sign certificate is forced.  Domain rejection is logically prior to
\eqref{eq:YMA10-verdict}.
\end{proof}

% =============================================================================
\section{Marked connected coefficients and the nonduplicating implementation}
\label{sec:YMA10-connected}
% =============================================================================

The exact moment series is intentionally redundant: powers of the plaquette
sum contain vacuum components which cancel after division by \(Z\).  The
following formal recursion identifies the source-marked coefficients that an
implementation should retain.

Write ordinary power-series coefficients at the central background as
\begin{equation}
 Z_\beta(0)=e^{-\beta P}\sum_{n\ge0}z_n\beta^n,
 \qquad
 -Z_\beta''(0)=e^{-\beta P}\sum_{n\ge0}r_n\beta^n,
 \qquad z_0=1.
 \label{eq:YMA10-zr-series}
\end{equation}
Define \(c_n\) recursively by
\begin{equation}
 \boxed{
 c_0=r_0,
 \qquad
 c_n=r_n-\sum_{j=1}^{n}z_jc_{n-j}
 \quad(n\ge1).}
 \label{eq:YMA10-connected-recursion}
\end{equation}

\begin{theorem}[Source-marked linked-coefficient cancellation]
\label{thm:YMA10-linked-coefficients}
As a formal series,
\begin{equation}
 \boxed{
 -\frac{Z_\beta''(0)}{Z_\beta(0)}
 =\sum_{n\ge0}c_n\beta^n.}
 \label{eq:YMA10-formal-H-series}
\end{equation}
Expand each \(z_n\) and \(r_n\) into local plaquette trace monomials.  Under
product Haar measure, every contribution whose reduced-link incidence graph
contains a connected component disjoint from all first- or second-source
insertions cancels from \(c_n\).  Equivalently, \(c_n\) is the sum of the
Haar cumulants of connected plaquette families carrying either one marked
second insertion or two marked first insertions.

The union of plaquette supports of each surviving term is explicit.  At every
finite order it may therefore be tagged as contractible or wrapping before it
is added to a local or topological coefficient bank.  This support tag is an
exact finite combinatorial statement; no convergence of the resulting
connected series at the selected weak coupling is asserted.
\end{theorem}

\begin{proof}
Multiplying \(\sum z_n\beta^n\) by \(\sum c_n\beta^n\) and comparing the
coefficient of \(\beta^n\) gives precisely
\eqref{eq:YMA10-connected-recursion}; hence the quotient identity holds
formally.  The moment--cumulant formula expresses a logarithmic derivative of
a partition function as a sum of joint cumulants.  Product Haar measure makes
two families of trace monomials independent whenever their reduced-link
supports are disjoint.  A joint cumulant vanishes when its arguments split
into two independent nonempty subfamilies.  The second source derivative
marks either one argument twice or two arguments once.  Therefore only a
single incidence-connected component containing the marks survives.  The
support and its winding class are read directly from that finite plaquette
family.
\end{proof}

\begin{firewall}[Formal connected coefficients do not solve weak coupling]
\label{fire:YMA10-formal-not-weak}
The quotient recursion is exact coefficient by coefficient and is the correct
nonduplicating implementation of the finite Haar calculation.  The Taylor
series of the logarithm is limited by complex zeros of the partition
function and need not converge at a selected large \(\beta\).  The rigorous
all-\(\beta\) evaluator remains the separate entire numerator and denominator
with Poisson tails.  A useful cofinal theorem needs a low-temperature or
renormalized connected expansion with its own uniform convergence proof.
\end{firewall}

\section{Exact replay boundary}
\label{sec:YMA10-replay}
The accompanying verifier
\begin{center}
\path{YMA_V10_exact_Haar_replay.py}
\end{center}
uses exact rational and symbolic arithmetic to check the rank-one
\(V\otimes\overline V\) and \(V^{\otimes3}\) Haar projectors, the trace
moments \(1/18\) and \(1/108\), the redundant-column Moore--Penrose
projection, the marked Hessian series on a rational finite model, the
negative-to-positive root orientation, both interval corner rules, the
connected cancellation of an independent vacuum factor, and the factorial
inequality underlying the Poisson tail.

The replay does not evaluate a nontrivial lattice Wilson integral, locate the
physical selected root, certify the local/topological split, or prove a
cutoff-uniform connected expansion.  Those are precisely the execution and
cofinal boundaries retained below.

% =============================================================================
\section{Upstream conclusion, proof ledger, and the V11 attack}
\label{sec:YMA10-terminal}
% =============================================================================

\begin{theorem}[Version V10 finite Wilson-card theorem]
\label{thm:YMA10-terminal}
Versions~V1--V9 are retained.  Version~V10 additionally proves:
\begin{enumerate}[label=\textup{(E10.\arabic*)},leftmargin=5.4em]
\item a finite exact pushforward tower is evaluated on one finest-lattice Haar
      cylinder, so exactly generated interactions do not require a separately
      supplied coefficient list;
\item the Wilson partition function and every bounded insertion have an exact
      Haar-moment expansion with the Poisson-tail bounds
      \eqref{eq:YMA10-Z-tail}--\eqref{eq:YMA10-insertion-tail};
\item every finite coefficient is an exact rational \(\SU(3)\) invariant-tensor
      contraction, with coarse matrices retained algebraically or by outward
      intervals;
\item the exact source Hessian is the ratio \(-Z''/Z\) at the symmetric axial
      background and has the marked expansion and tail
      \eqref{eq:YMA10-Zpp-series}--\eqref{eq:YMA10-Zpp-tail};
\item the Creutz root is the zero of the unnormalized entire function
      \eqref{eq:YMA10-root-function}; every fixed weak-facing root admits a
      finite endpoint corridor, complex-zero certificate, and arbitrary
      interval refinement;
\item root uncertainty, full shell curvature, an explicit topological row, and
      the exact flat reference assemble into the terminating interval
      \eqref{eq:YMA10-margin-interval};
\item every strict finite-card margin is decidable by refinement, while exact
      threshold contact remains unresolved and malformed loading is rejected;
\item the formal quotient recursion \eqref{eq:YMA10-connected-recursion}
      cancels every deformation-blind vacuum component and leaves only
      source-marked connected supports;
\item brute-force finite determination has a sufficient degree of order
      \(\beta M^4\) and therefore supplies neither a practical cofinal
      algorithm nor a uniform Wilson symbol estimate.
\end{enumerate}
No item proves weak-facing-root escape, supplies an efficient selected-root
numerical evaluation, bounds the explicit response-matched topological packet
uniformly, controls the complete nonaxial orbit bank, or establishes
higher-jet contraction or a continuum Yang--Mills mass gap.
\end{theorem}

\begin{proof}
Item~\textup{(E10.1)} is \Cref{prop:YMA10-flattening}.
Item~\textup{(E10.2)} is
\Cref{thm:YMA10-Wilson-moment-expansion}.
Item~\textup{(E10.3)} is \Cref{thm:YMA10-exact-Haar-algorithm}.
Item~\textup{(E10.4)} follows from
\Cref{thm:YMA10-Hessian-series} and
\Cref{cor:YMA10-H-interval}.
Item~\textup{(E10.5)} follows from
\Cref{prop:YMA10-root-equivalence} and
\Cref{thm:YMA10-root-isolation}.
Item~\textup{(E10.6)} is \Cref{thm:YMA10-selected-margin}, and
item~\textup{(E10.7)} is \Cref{cor:YMA10-strict-decision}.
Item~\textup{(E10.8)} is \Cref{thm:YMA10-linked-coefficients}.
Item~\textup{(E10.9)} follows from
\Cref{cor:YMA10-simple-depth},
\Cref{fire:YMA10-global-series-cost}, and
\Cref{fire:YMA10-formal-not-weak}.
\end{proof}

\begingroup\small
\begin{longtable}{>{\raggedright\arraybackslash}p{0.27\textwidth}
                  >{\raggedright\arraybackslash}p{0.20\textwidth}
                  >{\raggedright\arraybackslash}p{0.45\textwidth}}
\caption{Version V10 proof-status ledger.}\label{tab:YMA10-ledger}\\
\toprule
\textbf{Row}&\textbf{Status}&\textbf{Exact scope}\\
\midrule
\endfirsthead
\toprule
\textbf{Row}&\textbf{Status}&\textbf{Exact scope}\\
\midrule
\endhead
Finite exact Wilson cylinder&Specified conditional on the frozen root
&The original fine action, block constraint, Haar measure, and selected root
determine the pure Wilson-generated shell without an observed ratio.\\[0.3em]
Generated-action flattening&Proved exactly
&Every finite exact pushforward is one finest-lattice integral; explicit
coarse-only terms are retained additively.\\[0.3em]
Wilson--Haar expansion&Proved with explicit tail
&Partition, loop, and marked source numerators have uniform Poisson-tail
errors at every finite coupling.\\[0.3em]
\(\SU(3)\) coefficient evaluation&Exact finite algorithm
&Haar averaging is a rational invariant projector generated by
\(\delta\) and \(\varepsilon\) contractions.\\[0.3em]
Direct source curvature&Exact finite interval
&At the symmetric axial point the card is \(-Z''/Z\); no secant or sixth-order
bias is required.\\[0.3em]
Weak-facing root&Finite certified isolation
&The logarithm-free root function has a certified all-coupling endpoint
corridor and arbitrarily refinable complex-zero certificates.\\[0.3em]
Strict finite margin&Decidable by refinement
&An explicit action and topological split terminate in PASS or TYPED
OBSTRUCTION whenever the exact margin is nonzero.\\[0.3em]
Vacuum-component cancellation&Proved coefficientwise
&The quotient recursion leaves only source-marked connected Haar cumulants and
retains exact winding tags.\\[0.3em]
Practical selected-card value&Unexecuted
&The complete tensor contraction has not been run on a nontrivial selected
root; no numerical margin is invented.\\[0.3em]
Cofinal Wilson estimate&\stopstatus
&Global truncation is volumetric; no source-connected weak-coupling convergence
bound is proved.\\[0.3em]
Weak-facing escape&Open trajectory row
&Finite root computability does not imply that the weak-facing roots diverge.\\[0.3em]
Higher jets, contraction, mass&Not inferred
&The inherited downstream handoff remains closed until a uniform physical
quadratic family passes.\\
\bottomrule
\end{longtable}
\endgroup

\begin{assessment}[The upstream correction]
\label{ass:YMA10-upstream-correction}
The first finite Wilson card is not waiting for a new physical observable.
Its exact value is already fixed by the unnormalized Wilson cylinder and can
be certified by finite Haar algebra.  The programme was becoming circular
when it treated the absence of a convenient ratio instrument as absence of
the mathematical card and then added further incidence machinery around that
instrument.  V10 removes that circle.

The remaining obstruction is now sharper.  Either execute the finite tensor
contraction on a declared selected root, or prove a source-connected
low-temperature expansion whose constants are uniform on the selected family.
A new normalized-current, boundary-overlap, or finite-difference compiler does
not address that obstruction.
\end{assessment}

\begin{openproblem}[V11: source-connected low-temperature compression or one executed root]
\label{op:YMA10-V11}
Preserve Version~V10 and all frozen physical choices.  Proceed through exactly
one of the following two branches.
\begin{enumerate}[label=\textup{(V11.\arabic*)},leftmargin=5.6em]
\item \textbf{Finite execution branch.}  Supply a computable value of the
      predeclared \(u_*\), construct the weak-facing root certificate for the
      first nontrivial cutoff, and evaluate the exact coefficient banks
      \(\mu_n,\nu_n,\xi_n\) until \eqref{eq:YMA10-margin-interval} terminates.
      Use sparse invariant contractions and the connected recursion, but do
      not replace their certified tails by floating-point cancellation.
\item \textbf{Uniform analytic branch.}  Starting from the exact flat gap and
      the V5--V6 small-field inverse locality, construct a source-marked
      low-temperature polymer or flow expansion for the same numerator and
      denominator.  Prove a convergence norm which replaces the global
      \(96M^4\) volume in \eqref{eq:YMA10-Zpp-tail} by a summable connected
      source radius and controls the large-field remainder on the actual
      selected law.
\item In either branch retain every explicit coarse-only marginal curvature,
      the same local/topological split, and the exact Ward row.  If the direct
      axial interval fails, freeze that mode before opening V7 ancestry or
      YM--B radial attribution.
\item Do not append another general shell, current, ratio, or root theorem
      unless it supplies one of the two quantities in items~\textup{(V11.1)}
      or~\textup{(V11.2)}.
\end{enumerate}
The finite branch returns \pass, \obstruction, or \stopstatus\ on one literal
card.  The uniform branch returns a cofinal margin or a typed connected
large-field/infrared survivor.  Failure merely to improve the volumetric
Taylor cost is \stopstatus, not a physical obstruction.
\end{openproblem}

\begin{assessment}[V10 termination]
\label{ass:YMA10-termination}
The flattening theorem, Wilson--Haar moment and marked-Hessian expansions,
exact \(\SU(3)\) invariant projector, logarithm-free root certificate,
selected-root interval, strict finite decision, and connected cancellation
are \pass\ at their stated finite mathematical scopes.  The earlier loading
STOP is refined: the pure Wilson-generated part is an unexecuted finite
computation once the frozen response value is numerically calibrated, whereas
any separately named action or topological writer remains an actual loading
defect.  The selected cofinal response-matched margin remains \stopstatus:
neither a full selected-root contraction nor a uniform source-connected
low-temperature estimate has been supplied.  No negative physical mode has
been produced.
\end{assessment}

\section*{Version V10 append-only continuation record}
\addcontentsline{toc}{section}{Version V10 append-only continuation record}
The complete Version~V9 source above is preserved byte for byte up to its sole
final \verb|\end{document}| command.  Its SHA--256 digest is
\begin{center}\scriptsize\ttfamily
9ef28ffed46383a9ecbe0d54cb6c04f89d042e08317209fd6e96e23a56887c80.
\end{center}
For Version~V11, preserve this accumulated source, append new material
immediately before the final document-closing command, and increment the
filename to \texttt{\_V11.tex}.  Keep the physical response, weak-facing
selector, exact block map, parity-tree section, rooted Hodge carrier, first
axial mode, colour normalization, response coefficient, full unnormalized
action, and local/topological support convention fixed before reading any
root, coefficient, ratio, current, ancestry, radial, or transport value.

% =============================================================================
% END APPENDED VERSION V10 MATERIAL
% =============================================================================

% =============================================================================
% START APPENDED VERSION V11 MATERIAL
% =============================================================================

\clearpage
\hypersetup{
 pdftitle={Renewal Yang--Mills Reduced-Fibre Wilson Shell Symbol and Local-Vacuum Programme V11},
 pdfsubject={Weak-facing root multiplicity, response-chart regularity, and the complete marginally matched finite card}
}

\section{Version V11: weak-facing root multiplicity, response-chart regularity,
and the complete marginally matched finite card}
\label{sec:YMA11-direction}

Version~V10 is preserved above.  It proved that a fixed finite Wilson cylinder,
its weak-facing root, and its marked physical Hessian are in principle
computable by exact Haar contraction.  The next obstruction is upstream of the
volumetric Taylor cost.  The local action entering the selected shell is not
merely the pure Wilson action at a root; it is the \emph{response-matched}
action obtained after the generated cutoff insertion is corrected in the
plaquette marginal direction.  The inherited interaction-chart construction
selects the coefficient
\begin{equation}
 b^{\rm marg}
 =-\frac{\mathscr R_\Phi(\Psi)}{\mathscr R_\Phi(P)}
 \label{eq:YMA11-inherited-b}
\end{equation}
only when the response denominator is nonzero.  A root certificate which does
not decide this denominator does not yet define the first-order matched action.

The corrected order is
\begin{equation}
 \boxed{
 \begin{gathered}
 \text{unnormalized weak-facing root function and its multiplicity}\\
 \Downarrow\\
 \text{response slope in the plaquette direction}\\
 \Downarrow\\
 \text{generated-insertion response numerator and marginal coefficient}\\
 \Downarrow\\
 \text{same-law mixed deterministic/covariance jet card}\\
 \Downarrow\\
 \text{physical axial margin or attained response-chart/mode obstruction}.
 \end{gathered}}
 \label{eq:YMA11-corrected-order}
\end{equation}
This continuation does not reprove the general response differential, the
interaction-versus-field carrier distinction, or the mixed-jet Hessian theorem
already established in the dedicated Yang--Mills series.  It specializes those
results to the logarithm-free root and exact Haar evaluator of V10, proves the
missing multiplicity and conditioning estimates, and gives a terminating
finite response-matched card.

\begin{firewall}[A selected root is not yet a regular response coordinate]
\label{fire:YMA11-root-not-chart}
The scalar equation defining a weak-facing root determines a point on a
response level.  It does not imply that the plaquette coupling is transverse
to that level.  At a ramified root the ordinary inverse response slope is
unavailable; at a nearly ramified root the marginal coefficient can be large
or numerically ill conditioned.  Neither phenomenon is repaired by a more
accurate pure-Wilson Hessian evaluated before the marginal retuning.
\end{firewall}

% =============================================================================
\section{The response differential in unnormalized numerator coordinates}
\label{sec:YMA11-response-numerators}
% =============================================================================

Fix one finite card and four real Wilson-loop writers
\(\mathcal O_\alpha\),
\(\alpha\in\{11,00,10,01\}\), on the same reduced Haar fibre.  Let
\begin{equation}
 \varepsilon_{11}=\varepsilon_{00}=-1,
 \qquad
 \varepsilon_{10}=\varepsilon_{01}=1.
 \label{eq:YMA11-Creutz-signs}
\end{equation}
For an analytic action family \(\Phi_{\beta,s}\), define the unnormalized
loop numerators
\begin{equation}
 N_\alpha(\beta,s)
 :=\int \mathcal O_\alpha(z)
      e^{-\Phi_{\beta,s}(z)}\,\dd m_{\rm red}(z).
 \label{eq:YMA11-loop-numerators}
\end{equation}
On the positive loop branch put
\begin{equation}
 \chi(\beta,s)
 :=\sum_\alpha\varepsilon_\alpha\log N_\alpha(\beta,s)
 =-\log\frac{N_{11}N_{00}}{N_{10}N_{01}}.
 \label{eq:YMA11-Creutz-response}
\end{equation}
The partition denominator cancels because
\(\sum_\alpha\varepsilon_\alpha=0\).

For a bounded action insertion \(\Xi\), write
\begin{equation}
 N_{\alpha;\Xi}(\beta,s)
 :=\int \mathcal O_\alpha(z)\Xi(z)
      e^{-\Phi_{\beta,s}(z)}\,\dd m_{\rm red}(z).
 \label{eq:YMA11-inserted-numerator}
\end{equation}

\begin{proposition}[Unnormalized four-numerator response card]
\label{prop:YMA11-response-numerator-card}
Suppose
\(\partial_t\Phi_{\beta,s,t}|_{t=0}=\Xi\), with the loop writers held fixed.
Then
\begin{equation}
 \boxed{
 \left.\partial_t\chi[\Phi_{\beta,s,t}]\right|_{t=0}
 =-\sum_\alpha\varepsilon_\alpha
   \frac{N_{\alpha;\Xi}(\beta,s)}{N_\alpha(\beta,s)}.}
 \label{eq:YMA11-response-from-numerators}
\end{equation}
In particular, if
\begin{equation}
 \Phi_{\beta,s}=\Phi_*+(\beta-\beta_*)P+s\Psi,
 \label{eq:YMA11-two-parameter-action}
\end{equation}
then at \((\beta_*,0)\)
\begin{align}
 d&:=\partial_\beta\chi
   =-\sum_\alpha\varepsilon_\alpha
     \frac{N_{\alpha;P}}{N_\alpha},
 \label{eq:YMA11-response-denominator}\\
 n&:=\partial_s\chi
   =-\sum_\alpha\varepsilon_\alpha
     \frac{N_{\alpha;\Psi}}{N_\alpha}.
 \label{eq:YMA11-response-numerator}
\end{align}
Thus \(d=\mathscr R_{\Phi_*}(P)\),
\(n=\mathscr R_{\Phi_*}(\Psi)\), and the regular marginal coefficient is
\begin{equation}
 \boxed{b^{\rm marg}=-n/d.}
 \label{eq:YMA11-b-ratio}
\end{equation}
Every entry in
\eqref{eq:YMA11-response-denominator}--\eqref{eq:YMA11-response-numerator}
is an unnormalized insertion numerator of the kind evaluated by the V10 Haar
algorithm whenever \(P\), \(\Psi\), and the loop writers belong to its
explicit trace-polynomial or certified expansion class.
\end{proposition}

\begin{proof}
Differentiation under the finite compact integral gives
\[
 \partial_tN_\alpha=-N_{\alpha;\Xi}.
\]
Therefore
\[
 \partial_t\log N_\alpha
 =-N_{\alpha;\Xi}/N_\alpha.
\]
Multiply by the signs in \eqref{eq:YMA11-Creutz-signs} and sum.  No partition
function appears because it had already cancelled in
\eqref{eq:YMA11-Creutz-response}.  Substitution of the two coordinate
directions in \eqref{eq:YMA11-two-parameter-action} gives the last formulas.
\end{proof}

\begin{remark}[What V10 must acquire in addition to the root]
The loop numerators used to locate the root are already present in V10.  The
response chart adds only two marked banks: the four \(P\)-inserted numerators
and the four \(\Psi\)-inserted numerators.  It does not require a normalized
boundary source, a Witten-current ancestry panel, or a new response instrument.
Those objects remain optional attribution layers after the coefficient
\eqref{eq:YMA11-b-ratio} and the physical mixed jets have been assembled.
\end{remark}

% =============================================================================
\section{The logarithm-free root determines slope and multiplicity}
\label{sec:YMA11-root-multiplicity}
% =============================================================================

Fix the target response \(u\) and abbreviate
\begin{equation}
 A(\beta)=N_{11}(\beta)N_{00}(\beta),
 \qquad
 B(\beta)=N_{10}(\beta)N_{01}(\beta)>0.
 \label{eq:YMA11-AB}
\end{equation}
The V10 root function is
\begin{equation}
 D_u(\beta)=A(\beta)-e^{-u}B(\beta).
 \label{eq:YMA11-D}
\end{equation}

\begin{proposition}[Exact root factorization and slope identity]
\label{prop:YMA11-root-slope}
On the positive loop branch,
\begin{equation}
 \boxed{
 D_u(\beta)
 =B(\beta)\bigl(e^{-\chi(\beta)}-e^{-u}\bigr).}
 \label{eq:YMA11-D-factor}
\end{equation}
At every root \(\beta_*\) with \(\chi(\beta_*)=u\),
\begin{equation}
 \boxed{
 D_u'(\beta_*)=-e^{-u}B(\beta_*)\chi'(\beta_*),
 \qquad
 d=\chi'(\beta_*)
   =-\frac{e^uD_u'(\beta_*)}{B(\beta_*)}.}
 \label{eq:YMA11-slope-from-D}
\end{equation}
The multiplicity of \(\beta_*\) as a zero of \(D_u\) equals its multiplicity
as a zero of \(\chi-u\).  In particular, if
\begin{equation}
 \chi(\beta)-u
 =-c(\beta-\beta_*)^m v(\beta),
 \qquad c>0,
 \qquad v(\beta_*)=1,
 \label{eq:YMA11-stable-normal-form}
\end{equation}
then
\begin{equation}
 \boxed{
 \frac{D_u^{(m)}(\beta_*)}{m!}
 =e^{-u}B(\beta_*)c>0.}
 \label{eq:YMA11-D-leading-coefficient}
\end{equation}
A stable down-crossing has odd \(m\).  It is simple exactly when
\(D_u'(\beta_*)>0\), equivalently \(d<0\).
\end{proposition}

\begin{proof}
The response definition gives
\(A/B=e^{-\chi}\), proving \eqref{eq:YMA11-D-factor}.  Differentiate it.
At a root the factor in parentheses vanishes, so the derivative of \(B\)
drops out and \eqref{eq:YMA11-slope-from-D} follows.

The analytic function
\[
 z\longmapsto e^{-z}-1
\]
has a simple zero at \(z=0\).  Hence composition with
\(z=\chi-u\) preserves zero multiplicity, and multiplication by the nonzero
analytic factor \(e^{-u}B\) does not change it.  Substituting
\eqref{eq:YMA11-stable-normal-form} into
\[
 e^{-\chi}-e^{-u}
 =e^{-u}\bigl(e^{c(\beta-\beta_*)^mv(\beta)}-1\bigr)
\]
gives the leading coefficient in
\eqref{eq:YMA11-D-leading-coefficient}.  A real analytic function changes
sign across an isolated zero exactly when the zero order is odd.  The final
sign equivalences follow from \eqref{eq:YMA11-slope-from-D}.
\end{proof}

\begin{corollary}[The V10 argument-principle box is also a ramification card]
\label{cor:YMA11-root-count-multiplicity}
Suppose a complex root rectangle or disk has zero-free boundary and contains
one geometric real zero \(\beta_*\) of \(D_u\).  Its argument-principle winding
number equals the multiplicity \(m\) in
\eqref{eq:YMA11-stable-normal-form}.  A winding number one certifies the
regular response branch.  A larger odd winding number certifies a stable
ramified weak-facing root and blocks the ordinary coefficient
\eqref{eq:YMA11-b-ratio} whenever the generated response numerator does not
vanish.
\end{corollary}

\begin{proof}
The argument principle counts zeros with multiplicity.  By hypothesis there
is only one geometric zero in the region, and
\Cref{prop:YMA11-root-slope} identifies its root-function multiplicity with
the response-level multiplicity.  The final statement is the linear equation
\(n+bd=0\) with \(d=0\).
\end{proof}

\subsection{A quantitative slope floor from the existing root disk}

\begin{theorem}[Boundary-modulus certificate for the response slope]
\label{thm:YMA11-root-disk-slope}
Let \(D_u\) be analytic on the closed disk
\(\overline{\mathbb D(c,R)}\).  Assume:
\begin{enumerate}[label=\textup{(S\arabic*)},leftmargin=3.5em]
\item the disk contains exactly one zero counting multiplicity, so that the
      zero \(\beta_*\) is simple;
\item \(\abs{\beta_*-c}\le r<R\);
\item on the boundary,
      \begin{equation}
       0<m_\partial\le\abs{D_u(z)}\le M_\partial;
       \label{eq:YMA11-boundary-modulus}
      \end{equation}
\item at the root,
      \begin{equation}
       0<B_-\le B(\beta_*)\le B_+.
       \label{eq:YMA11-B-window}
      \end{equation}
\end{enumerate}
Then
\begin{equation}
 \boxed{
 \frac{m_\partial}{R+r}
 \le\abs{D_u'(\beta_*)}
 \le\frac{M_\partial}{R-r}.}
 \label{eq:YMA11-Dprime-window}
\end{equation}
If the real zero is the stable down-crossing, its response denominator obeys
\begin{equation}
 \boxed{
 -\frac{e^uM_\partial}{(R-r)B_-}
 \le d
 \le-\frac{e^um_\partial}{(R+r)B_+}<0.}
 \label{eq:YMA11-d-window}
\end{equation}
Thus the same complex boundary computation used to isolate the root can
supply a nonzero response-slope margin without differentiating a noisy
normalized response curve.
\end{theorem}

\begin{proof}
Factor
\[
 D_u(z)=(z-\beta_*)g(z).
\]
Because the disk contains exactly one simple zero, \(g\) is analytic and
nonvanishing there.  On \(\abs{z-c}=R\),
\[
 R-r\le\abs{z-\beta_*}\le R+r.
\]
Hence
\[
 \frac{m_\partial}{R+r}
 \le\abs{g(z)}
 \le\frac{M_\partial}{R-r}.
\]
Apply the maximum-modulus principle to \(g\) for the upper bound at
\(\beta_*\), and to \(1/g\) for the lower bound.  Since
\(D_u'(\beta_*)=g(\beta_*)\), this proves
\eqref{eq:YMA11-Dprime-window}.  The stable orientation makes \(D_u'>0\), so
\eqref{eq:YMA11-slope-from-D} and
\eqref{eq:YMA11-B-window} give
\eqref{eq:YMA11-d-window} with the indicated endpoint directions.
\end{proof}

\begin{remark}[No derivative-of-log instability]
The quantities in \eqref{eq:YMA11-boundary-modulus} are values of the entire
unnormalized function \(D_u\).  The positive product \(B\) is evaluated
separately.  No subtraction of nearly equal logarithms is used.  A loose
boundary modulus can make the resulting interval conservative, but it cannot
reverse its certified sign.
\end{remark}

% =============================================================================
\section{Regular and ramified response-matching branches}
\label{sec:YMA11-response-branches}
% =============================================================================

Let
\begin{equation}
 F(\beta,s):=\chi(\beta,s)-u,
 \qquad F(\beta_*,0)=0,
 \label{eq:YMA11-F}
\end{equation}
and retain \(d=F_\beta(\beta_*,0)\),
\(n=F_s(\beta_*,0)\).

\begin{theorem}[Regular selected-root chart and its first two jets]
\label{thm:YMA11-regular-root-chart}
If \(d\ne0\), there is a unique real-analytic root graph
\(s\mapsto\beta(s)\) near zero satisfying
\begin{equation}
 \beta(0)=\beta_*,
 \qquad F(\beta(s),s)=0.
 \label{eq:YMA11-root-graph}
\end{equation}
Its first derivative is exactly the response-matching coefficient,
\begin{equation}
 \boxed{
 \beta'(0)=-\frac nd=b^{\rm marg}.}
 \label{eq:YMA11-root-first-jet}
\end{equation}
Its second derivative is
\begin{equation}
 \boxed{
 \beta''(0)
 =-\frac{
   F_{ss}+2b^{\rm marg}F_{\beta s}
   +(b^{\rm marg})^2F_{\beta\beta}}
  {d}\bigg|_{(\beta_*,0)}.}
 \label{eq:YMA11-root-second-jet}
\end{equation}
Consequently, the first-order interaction
\(\Psi+b^{\rm marg}P\) is the tangent to the actual response level, while
\eqref{eq:YMA11-root-second-jet} is the first curvature correction required
if the finite deformation rather than only its Hessian tangent is used.
\end{theorem}

\begin{proof}
The analytic implicit-function theorem gives the unique graph.  Differentiate
\(F(\beta(s),s)=0\) once to obtain
\[
 d\,\beta'(0)+n=0.
\]
Differentiate a second time:
\[
 d\,\beta''(0)
 +F_{\beta\beta}(\beta'(0))^2
 +2F_{\beta s}\beta'(0)+F_{ss}=0.
\]
Substitute \eqref{eq:YMA11-root-first-jet} and solve for
\(\beta''(0)\).
\end{proof}

\begin{theorem}[Ramified perturbation alternative]
\label{thm:YMA11-ramified-perturbation}
Assume that the selected stable root has odd multiplicity \(m\ge3\) and local
normal form
\begin{equation}
 F(\beta_*+x,0)=-cx^m+O(x^{m+1}),
 \qquad c>0.
 \label{eq:YMA11-ramified-unperturbed}
\end{equation}
Then exactly the following first-order alternatives occur.
\begin{enumerate}[label=\textup{(R\arabic*)},leftmargin=3.6em]
\item If \(n=F_s(\beta_*,0)\ne0\), no finite scalar \(b\) solves
      \(n+bd=0\).  For every real root branch
      \(x=x(s)\to0\) of \(F(\beta_*+x,s)=0\),
      \begin{equation}
       \boxed{
       x(s)^m=\frac nc\,s+o(s),
       \qquad
       \abs{x(s)}\asymp\abs{s}^{1/m}.}
       \label{eq:YMA11-Puiseux-scale}
      \end{equation}
      At least one such real branch exists for each sufficiently small real
      \(s\), but no branch has a finite derivative at zero.  This is a typed
      response-chart obstruction, not a field-space soft mode.
\item If \(n=0\), then every coefficient \(b\) makes the first-order response
      of \(\Psi+bP\) vanish.  The first-order card does not select a marginal
      coefficient.  The least nonzero mixed response jet in
      \((x,s)\) is required; until it is supplied, the matched action is
      \stopstatus.
\end{enumerate}
\end{theorem}

\begin{proof}
Here \(d=F_\beta(\beta_*,0)=0\).  The linear response equation is therefore
\(n+bd=n\), which proves the two algebraic alternatives.

For the first branch, Taylor expansion gives
\begin{equation}
 F(\beta_*+x,s)
 =-cx^m+ns+o(\abs{x}^m+\abs{s}).
 \label{eq:YMA11-ramified-Taylor}
\end{equation}
Choose a fixed \(K>(2\abs n/c)^{1/m}+1\).  At
\(x=\pm K\abs{s}^{1/m}\), the leading odd powers have opposite signs and
dominate the remainder for small \(s\).  The intermediate value theorem
therefore gives a real root between them.  Every root tending to zero first
satisfies \(\abs{x}^m=O(\abs{s})\): otherwise divide
\eqref{eq:YMA11-ramified-Taylor} by \(\abs{x}^m\) and obtain the contradiction
\(-c\operatorname{sgn}(x)^m+o(1)=0\).  Substitution of that estimate back into
\eqref{eq:YMA11-ramified-Taylor} yields
\(cx^m=ns+o(s)\), proving \eqref{eq:YMA11-Puiseux-scale}.  Since
\(1/m-1<0\), \(\abs{x(s)/s}\to\infty\), so no finite derivative exists.
\end{proof}

\begin{firewall}[Ramification is prior to the mixed field Hessian]
\label{fire:YMA11-ramification-prior}
The response coefficient is selected on the interaction carrier before the
coarse field jets are formed.  If the regular coefficient does not exist, one
may still evaluate the pure pre-retuning Hessian, but it is not the Hessian of
a first-order response-matched action.  A ramification card is a trajectory or
coordinate obstruction.  It becomes a physical field-mode obstruction only
if an independently defined finite matched branch is supplied and its
assembled field Hessian has a negative witness.
\end{firewall}

\begin{countertheorem}[A root and the pre-retuning Hessian do not determine the matched card]
\label{cth:YMA11-root-pure-not-matched}
For every \(\epsilon>0\), let the response germ be
\begin{equation}
 F_\epsilon(\beta,s)
 =-\epsilon(\beta-\beta_*)+s.
 \label{eq:YMA11-counter-response}
\end{equation}
All cards have the same stable selected root \(\beta_*\) at \(s=0\), and let
all have the same one-dimensional pre-retuning physical Hessian
\(H_{\rm pre}=1\).  Take
\begin{equation}
 H_P=0,
 \qquad
 \mathbb V=
 \begin{pmatrix}0&0\\0&1\end{pmatrix}\succeq0.
 \label{eq:YMA11-counter-jets}
\end{equation}
Then
\begin{equation}
 d=-\epsilon,
 \qquad n=1,
 \qquad b^{\rm marg}=\epsilon^{-1},
 \qquad
 \boxed{H_{\rm RG}=1-\epsilon^{-2}.}
 \label{eq:YMA11-counter-H}
\end{equation}
Thus identical selected roots and identical pre-retuning Hessians are
compatible with arbitrarily negative response-matched Hessians.  The response
slope, generated response numerator, and plaquette-score covariance are
independent load-bearing entries.
\end{countertheorem}

\begin{proof}
The first three values follow by differentiation of
\eqref{eq:YMA11-counter-response}.  Compressing
\eqref{eq:YMA11-counter-jets} by \((1,b)\) gives \(b^2\), so the exact
mixed-jet formula gives \eqref{eq:YMA11-counter-H}.  Positivity of the
covariance panel is immediate.
\end{proof}

% =============================================================================
\section{Validated marginal coefficient and complete mixed-jet card}
\label{sec:YMA11-matched-card}
% =============================================================================

Work henceforth on the already rooted, Hodge-reduced, support-stable physical
field bank \(E\).  Response matching remains on the interaction carrier.  Once
\(b\) has been selected, all expectations in the following card are taken in
the \emph{same final law} with action
\begin{equation}
 \Phi_b=\Phi_0+bP.
 \label{eq:YMA11-final-law}
\end{equation}
Let
\begin{equation}
 H_0=\mathbb E_bD_V^2\Phi_0,
 \qquad
 H_P=\mathbb E_bD_V^2P,
 \label{eq:YMA11-deterministic-blocks}
\end{equation}
and let the centred score covariance be
\begin{equation}
 \mathbb V_b=
 \begin{pmatrix}
 C_{00}&C_{0P}\\
 C_{P0}&C_{PP}
 \end{pmatrix}\succeq0.
 \label{eq:YMA11-covariance-block}
\end{equation}
For \(Q_bv=(v,bv)\), the inherited exact formula is
\begin{equation}
 \boxed{
 H_b^{\rm RG}=H_0+bH_P-Q_b^*\mathbb V_bQ_b.}
 \label{eq:YMA11-exact-matched-H}
\end{equation}
The point of this section is not to reprove
\eqref{eq:YMA11-exact-matched-H}, but to make its response coefficient and
outward error propagation executable.

\begin{proposition}[Sign-safe interval for the marginal coefficient]
\label{prop:YMA11-b-interval}
Suppose outward intervals satisfy
\begin{equation}
 n\in[n_-,n_+],
 \qquad
 d\in[d_-,d_+],
 \qquad
 0\notin[d_-,d_+].
 \label{eq:YMA11-nd-interval}
\end{equation}
Set
\begin{equation}
 \mathcal C_{n,d}
 :=\left\{-\frac{n_i}{d_j}:
 n_i\in\{n_-,n_+\},\ d_j\in\{d_-,d_+\}\right\}.
 \label{eq:YMA11-corner-set}
\end{equation}
Then
\begin{equation}
 \boxed{
 b^{\rm marg}\in
 [b_-,b_+]
 :=[\min\mathcal C_{n,d},\max\mathcal C_{n,d}].}
 \label{eq:YMA11-b-interval}
\end{equation}
Moreover, if \(b=-n/d\), \(\widehat b=-\widehat n/\widehat d\), and
\begin{equation}
 \abs d,\abs{\widehat d}\ge d_*>0,
 \label{eq:YMA11-slope-floor}
\end{equation}
then
\begin{equation}
 \boxed{
 \abs{b-\widehat b}
 \le\frac{\abs{n-\widehat n}}{d_*}
 +\frac{\abs{\widehat n}\abs{d-\widehat d}}{d_*^2}.}
 \label{eq:YMA11-b-conditioning}
\end{equation}
Thus a uniform slope floor is the quantitative condition which turns the
marginal retuning into a stable finite input.
\end{proposition}

\begin{proof}
The function \((n,d)\mapsto-n/d\) is continuous on the rectangle in
\eqref{eq:YMA11-nd-interval}; because the denominator has fixed sign, its
extrema occur at the four corners.  For the second claim,
\[
 \left|\frac nd-\frac{\widehat n}{\widehat d}\right|
 \le\frac{|n-\widehat n|}{|d|}
 +|\widehat n|
  \frac{|d-\widehat d|}{|d\widehat d|},
\]
and \eqref{eq:YMA11-slope-floor} gives the result.
\end{proof}

\begin{theorem}[Validated response-matched Hessian enclosure]
\label{thm:YMA11-matched-enclosure}
Let \(\bar b\in\R\), \(r_b\ge0\), and suppose
\begin{equation}
 \abs{b-\bar b}\le r_b,
 \qquad
 B:=\abs{\bar b}+r_b,
 \qquad
 q_B:=\sqrt{1+B^2},
 \qquad
 q_0:=\sqrt{1+\bar b^2}.
 \label{eq:YMA11-b-radius}
\end{equation}
Assume Hermitian approximants satisfy
\begin{equation}
 \norm{H_0-\widehat H_0}\le\epsilon_0,
 \qquad
 \norm{H_P-\widehat H_P}\le\epsilon_P,
 \qquad
 \norm{\mathbb V_b-\widehat{\mathbb V}}\le\epsilon_V.
 \label{eq:YMA11-jet-errors}
\end{equation}
Define
\begin{equation}
 \widehat H_{\rm RG}
 :=\widehat H_0+\bar b\widehat H_P
   -Q_{\bar b}^*\widehat{\mathbb V}Q_{\bar b}
 \label{eq:YMA11-Hhat}
\end{equation}
and
\begin{equation}
 \boxed{
 \epsilon_{\rm RG}
 :=\epsilon_0+B\epsilon_P+r_b\norm{\widehat H_P}
   +q_B^2\epsilon_V
   +r_b(q_B+q_0)\norm{\widehat{\mathbb V}}.}
 \label{eq:YMA11-RG-error}
\end{equation}
Then
\begin{equation}
 \boxed{
 \norm{H_b^{\rm RG}-\widehat H_{\rm RG}}
 \le\epsilon_{\rm RG}.}
 \label{eq:YMA11-RG-enclosure}
\end{equation}

Let \(T\) be the topological Hessian and \(S\succ0\) the selected flat
reference on the supported physical bank.  If
\begin{equation}
 \norm{T-\widehat T}\le\epsilon_T,
 \qquad
 \norm{S-\widehat S}\le\epsilon_S,
 \label{eq:YMA11-TS-errors}
\end{equation}
and
\begin{equation}
 \widehat M
 :=\widehat H_{\rm RG}-\widehat T-\delta\widehat S,
 \qquad
 \epsilon_M:=\epsilon_{\rm RG}+\epsilon_T+\delta\epsilon_S,
 \label{eq:YMA11-margin-approx}
\end{equation}
then
\begin{equation}
 \boxed{
 \begin{array}{rcl}
 \lambda_{\min}(\widehat M)-\epsilon_M>0
 &\Longrightarrow&\pass,\\[0.25em]
 \lambda_{\min}(\widehat M)+\epsilon_M<0
 &\Longrightarrow&\obstruction,\\[0.25em]
 \text{otherwise}&\Longrightarrow&\stopstatus.
 \end{array}}
 \label{eq:YMA11-matched-verdict}
\end{equation}
On the obstruction branch, a unit minimum eigenvector of
\(\widehat M\) is an explicit physical witness after its already verified
support and gauge reduction.
\end{theorem}

\begin{proof}
The deterministic part obeys
\[
 \norm{bH_P-\bar b\widehat H_P}
 \le |b|\epsilon_P+|b-\bar b|\norm{\widehat H_P}
 \le B\epsilon_P+r_b\norm{\widehat H_P}.
\]
Also \(\norm{Q_b}\le q_B\) and
\(\norm{Q_b-Q_{\bar b}}=|b-\bar b|\le r_b\).  Therefore
\begin{align*}
 &\norm{Q_b^*\mathbb V_bQ_b
       -Q_{\bar b}^*\widehat{\mathbb V}Q_{\bar b}}\\
 &\quad\le
 q_B^2\epsilon_V
 +\norm{(Q_b-Q_{\bar b})^*\widehat{\mathbb V}Q_b}
 +\norm{Q_{\bar b}^*\widehat{\mathbb V}(Q_b-Q_{\bar b})}\\
 &\quad\le q_B^2\epsilon_V
  +r_b(q_B+q_0)\norm{\widehat{\mathbb V}}.
\end{align*}
Adding the \(H_0\) error proves
\eqref{eq:YMA11-RG-enclosure}.  The margin error is then at most
\(\epsilon_M\).  Weyl's inequality gives the positive branch.  If
\(v\) is a unit minimum eigenvector of \(\widehat M\), then
\[
 \ip{v}{(H_b^{\rm RG}-T-\delta S)v}
 \le\lambda_{\min}(\widehat M)+\epsilon_M,
\]
which proves the obstruction branch and its witness.  The remaining interval
contains zero and is unresolved.
\end{proof}

\begin{definition}[Response-regular executable manifest]
\label{def:YMA11-response-manifest}
A finite selected-shell response manifest consists of
\begin{equation}
 \boxed{
 \mathfrak M_{\rm resp}
 =\bigl(
 u,\ \mathcal D_*,\ m_*,\ [B_-,B_+],\ [d_-,d_+],\ [n_-,n_+],
 [b_-,b_+],\ \widehat H_0,\widehat H_P,\widehat{\mathbb V},
 \widehat T,\widehat S;\ \boldsymbol\epsilon
 \bigr),}
 \label{eq:YMA11-response-manifest}
\end{equation}
where \(\mathcal D_*\) is the certified root region, \(m_*\) its
argument-principle multiplicity, every interval is outward rounded, the jet
blocks are taken in the same final law, and \(\boldsymbol\epsilon\) contains
the stated operator error radii.  The manifest is \emph{regular} exactly when
\(m_*=1\) and \(0\notin[d_-,d_+]\).  It is \emph{ramified} when
\(m_*>1\), and is \emph{ill typed} when its loop positivity, root count,
same-law, support, Hermiticity, covariance positivity, or Ward checks fail.
\end{definition}

\begin{corollary}[Finite response-matched decision]
\label{cor:YMA11-finite-decision}
Every regular manifest with refinable root, insertion, and mixed-jet intervals
terminates in \pass\ or \obstruction\ whenever the exact selected margin is
nonzero.  Exact threshold contact remains \stopstatus.  A ramified manifest
terminates first in the response-chart alternative of
\Cref{thm:YMA11-ramified-perturbation}; it is not passed to
\eqref{eq:YMA11-matched-verdict} unless an independently specified finite
matched branch replaces the absent ordinary tangent.
\end{corollary}

\begin{proof}
On the regular branch, the slope interval is separated from zero, so
\Cref{prop:YMA11-b-interval} and
\Cref{thm:YMA11-matched-enclosure} shrink with the underlying certified data.
A nonzero Hermitian spectral margin has a neighborhood of fixed sign.  The
ramified statement is exactly the domain condition of the coefficient
\eqref{eq:YMA11-b-ratio}.
\end{proof}

% =============================================================================
\section{The marginal coefficient as a physical instability amplifier}
\label{sec:YMA11-b-amplifier}
% =============================================================================

For a fixed final law, expand \eqref{eq:YMA11-exact-matched-H} as
\begin{equation}
 \boxed{
 H_b^{\rm RG}=H_{\rm pre}+bL-b^2C,}
 \label{eq:YMA11-quadratic-b}
\end{equation}
where
\begin{equation}
 H_{\rm pre}:=H_0-C_{00},
 \qquad
 L:=H_P-C_{0P}-C_{P0},
 \qquad
 C:=C_{PP}\succeq0.
 \label{eq:YMA11-HLC}
\end{equation}
The negative quadratic term is not an estimate; it is the exact plaquette-score
variance in the final response-matched law.

\begin{lemma}[Covariance support removes cross terms on invisible directions]
\label{lem:YMA11-covariance-support}
For the positive block covariance
\(\mathbb V_b\) in \eqref{eq:YMA11-covariance-block},
\begin{equation}
 \boxed{
 Cv=0\quad\Longrightarrow\quad C_{0P}v=0.}
 \label{eq:YMA11-cross-support}
\end{equation}
Equivalently,
\(\Ran C_{P0}\subseteq\Ran C\).  Thus on a field polarization invisible to
the plaquette score, the mixed covariance term vanishes automatically and the
marginal dependence is purely deterministic.
\end{lemma}

\begin{proof}
A positive semidefinite block matrix defines a semidefinite inner product.
For arbitrary \(u\in E\), its Cauchy--Schwarz inequality gives
\[
 |\ip{u}{C_{0P}v}|^2
 \le\ip{u}{C_{00}u}\ip{v}{Cv}=0.
\]
Hence \(C_{0P}v=0\).  Taking adjoints gives the range formulation in finite
dimension.
\end{proof}

\begin{theorem}[Quantitative marginal-slope blow-up alternative]
\label{thm:YMA11-b-blowup}
Let \(v\) be a unit physical polarization and suppose, on a family of regular
selected cards,
\begin{equation}
 \begin{aligned}
 C[v]&\ge c_*>0,\\
 \abs{L[v]}&\le\ell_*,\\
 (H_{\rm pre}-T-\delta S)[v]&\le h_*.
 \end{aligned}
 \label{eq:YMA11-blowup-hypotheses}
\end{equation}
Then
\begin{equation}
 \boxed{
 (H_b^{\rm RG}-T-\delta S)[v]
 \le h_*+\ell_*\abs b-c_*\abs b^2.}
 \label{eq:YMA11-blowup-upper}
\end{equation}
In particular the mode is a strict physical obstruction whenever
\begin{equation}
 \boxed{
 \abs b>
 b_{\rm crit}:=
 \frac{\ell_*+
  \sqrt{\ell_*^2+4c_*\max\{h_*,0\}}}{2c_*}.}
 \label{eq:YMA11-bcrit}
\end{equation}
Consequently, if along selected regular cards
\begin{equation}
 \abs{d_r}\longrightarrow0,
 \qquad
 \abs{n_r}\ge n_*>0,
 \label{eq:YMA11-slope-collapse-nonzero-numerator}
\end{equation}
while \eqref{eq:YMA11-blowup-hypotheses} holds with common constants on an
attained or transported polarization, then \(\abs{b_r}\to\infty\) and the
assembled response-matched margin is eventually negative on that
polarization.
\end{theorem}

\begin{proof}
Evaluate \eqref{eq:YMA11-quadratic-b} on \(v\), subtract \(T+\delta S\), and
apply \eqref{eq:YMA11-blowup-hypotheses}.  For \(x=|b|\), the right side is
strictly negative whenever
\(c_*x^2-\ell_*x-\max\{h_*,0\}>0\).  Solving the corresponding quadratic
equality gives \eqref{eq:YMA11-bcrit}.  Finally
\(\abs{b_r}=\abs{n_r/d_r}\to\infty\), so the threshold is crossed.
\end{proof}

\begin{corollary}[Three typed near-ramification outcomes]
\label{cor:YMA11-near-ramification-types}
On a sequence with \(d_r\to0\), the following cases are mathematically
distinct.
\begin{enumerate}[label=\textup{(N\arabic*)},leftmargin=3.6em]
\item \textbf{Visible marginal-score amplification.}  If
      \(\abs{n_r}\) stays positive and \(C_r\) has a common positive direction
      satisfying \Cref{thm:YMA11-b-blowup}, an attained negative physical mode
      eventually occurs.
\item \textbf{Deterministic marginal amplification.}  If \(C_rv_r=0\), then
      \Cref{lem:YMA11-covariance-support} removes the mixed covariance on
      \(v_r\), and the sign is decided by the affine deterministic term
      \(b_rH_{P,r}[v_r]\) together with the bounded pre-card.
\item \textbf{Field-invisible response singularity.}  If both the plaquette
      score and deterministic plaquette field jet vanish on the relevant
      physical carrier, the large or undefined interaction coefficient does
      not produce a field Hessian on that carrier.  The output is a response
      chart or transport singularity, not a local-vacuum mode obstruction.
\end{enumerate}
If \(n_r\to0\) together with \(d_r\to0\), boundedness of the numerical ratio
alone does not give a uniform response chart; the higher-order common zero
must still be resolved.
\end{corollary}

\begin{proof}
The first item is \Cref{thm:YMA11-b-blowup}.  The second uses
\Cref{lem:YMA11-covariance-support} in
\eqref{eq:YMA11-quadratic-b}; on that direction \(L=H_P\).  In the third,
all \(b\)-dependent field jets vanish on the carrier.  The final statement
follows because \(n/d\) may have a finite limit along many inequivalent
higher-order response germs while the derivative transverse to the response
level remains zero.
\end{proof}

\begin{proposition}[Uniform response regularity confines the analytic problem]
\label{prop:YMA11-uniform-regularity}
Suppose a cofinal selected family satisfies
\begin{equation}
 \boxed{
 \abs{d_r}\ge d_*>0,
 \qquad
 \abs{n_r}\le n^*<\infty.}
 \label{eq:YMA11-uniform-dn}
\end{equation}
Then
\begin{equation}
 \boxed{\abs{b_r^{\rm marg}}\le n^*/d_*}
 \label{eq:YMA11-uniform-b}
\end{equation}
for every card.  Hence a Wilson-specific low-temperature, flow, or polymer
estimate need only be uniform on one compact marginal-coefficient interval.
Conversely, a proof uniform only for \(b\) in a compact set cannot be applied
before either \eqref{eq:YMA11-uniform-dn} or an equivalent direct bound on the
selected coefficient has been established.
\end{proposition}

\begin{proof}
This is immediate from \(b_r=-n_r/d_r\).  The converse is a domain statement:
without a coefficient bound, the selected final actions need not lie in the
compact family on which the analytic estimate was proved.
\end{proof}

% =============================================================================
\section{Finite acquisition order and the revised V12 target}
\label{sec:YMA11-acquisition}
% =============================================================================

The first selected card can now be executed without a normalized response
instrument.  The necessary finite quantities, in order, are:
\begin{enumerate}[label=\textup{(A11.\arabic*)},leftmargin=5.4em]
\item the four positive loop numerators and a complex zero-counting region for
      \(D_u\);
\item the boundary modulus and positive \(B\)-interval needed for the root
      multiplicity and slope card;
\item the four \(P\)-inserted and four \(\Psi\)-inserted loop numerators,
      yielding \(d,n,b\);
\item the complete same-final-law deterministic pair \((H_0,H_P)\) and joint
      covariance block \(\mathbb V_b\);
\item the local/topological split and exact Ward held-out on the assembled
      matrix;
\item only after the direct verdict, any Witten-current, boundary-source,
      radial, mediator, or ancestry attribution.
\end{enumerate}
The first three rows are scalar exact-Haar insertion banks.  The fourth is the
existing response-matched physical card.  No row may be populated from a pure
law different from \(\Phi_b\) merely because the marginal coefficient was
computed there.

\begin{theorem}[Response-chart-first finite alternative]
\label{thm:YMA11-finite-alternative}
For one predeclared finite selected shell, exactly one of the following ordered
outcomes applies.
\begin{enumerate}[label=\textup{(V11.\arabic*)},leftmargin=5.1em]
\item \textbf{Rejected input.}  The loop-positive branch, root boundary,
      common action, support, covariance positivity, Hermiticity, or Ward row
      fails its defining identity.  This is not a theorem verdict.
\item \textbf{Ramified response-chart obstruction.}  The root has multiplicity
      greater than one and the generated response numerator is nonzero.  No
      ordinary marginal tangent exists.
\item \textbf{Higher response ramification.}  Both first response derivatives
      vanish.  The least nonzero mixed response jet is unpopulated, so the
      matched action is \stopstatus.
\item \textbf{Regular incomplete card.}  The root is simple, but one of the
      numerator, coefficient, same-law mixed-jet, topology, or error-radius
      rows is absent or its interval is inconclusive.  The verdict is
      \stopstatus.
\item \textbf{Regular physical pass.}  The first line of
      \eqref{eq:YMA11-matched-verdict} holds.  The selected physical bank has
      the declared positive margin.
\item \textbf{Regular physical obstruction.}  The second line of
      \eqref{eq:YMA11-matched-verdict} holds.  The returned supported
      transverse eigenvector is the attained pure-mode witness.
\end{enumerate}
A failure of the response-slope floor is never silently relabeled as a
large-field, topological, radial, or massless field mode.  Those attributions
open only after the matched action or a ramified replacement branch has been
physically defined.
\end{theorem}

\begin{proof}
The input checks are the domain of the subsequent formulas.  The two
ramification branches are exhaustive when the response denominator vanishes,
by \Cref{thm:YMA11-ramified-perturbation}.  When it does not vanish,
\Cref{prop:YMA11-b-interval} defines the coefficient and
\Cref{thm:YMA11-matched-enclosure} gives the three interval outcomes.  The
ordered list is therefore exhaustive and disjoint.
\end{proof}

\begin{assessment}[What going upstream changed]
\label{ass:YMA11-upstream-change}
Version~V10 correctly removed the false idea that an extra ratio instrument
was required to define a finite pure Wilson card.  Its V11 fork nevertheless
placed the computational Taylor cost before one earlier domain question: is
the weak-facing root a regular coordinate for the plaquette marginal
retuning?  The answer is encoded in the multiplicity and derivative of the
same unnormalized root function already used by V10.

The new result is not another root theorem detached from the target.  The
slope is the denominator of the actual marginal coefficient, and that
coefficient enters the physical Hessian quadratically through a positive
same-law covariance.  Root ramification therefore either blocks the ordinary
matched action, selects a higher response jet, or, under source visibility,
amplifies an attained physical mode.  Only after this branch is decided does
a source-connected low-temperature estimate have the correct compact family
of final actions on which to operate.
\end{assessment}

\begin{openproblem}[V12: execute the response-regular first card or return its ramification]
\label{op:YMA11-V12}
Preserve Version~V11 and all frozen physical choices.  On the first nontrivial
selected cutoff, proceed in the following order.
\begin{enumerate}[label=\textup{(V12.\arabic*)},leftmargin=5.4em]
\item Load the numerical response target \(u_*\) and use the V10 unnormalized
      root function to produce one zero-counting region, multiplicity, and
      stable-crossing orientation.
\item If the root is simple, acquire the boundary modulus and \(B\)-window of
      \Cref{thm:YMA11-root-disk-slope}; evaluate the marked numerator banks in
      \eqref{eq:YMA11-response-denominator}--\eqref{eq:YMA11-response-numerator}
      and return a certified interval for \(b^{\rm marg}\).
\item If the root is ramified, evaluate the generated response numerator.  A
      nonzero value returns the response-chart obstruction.  A zero value
      opens only the least mixed response jet needed to determine the local
      branch; do not continue with a fictitious linear coefficient.
\item On the regular branch, form the complete final-law mixed jet
      \eqref{eq:YMA11-exact-matched-H}, including the cross covariance and
      explicit topological subtraction, and apply
      \eqref{eq:YMA11-matched-verdict}.
\item If finite exact Haar contraction remains computationally prohibitive,
      the uniform analytic replacement must prove both the response regularity
      bounds \eqref{eq:YMA11-uniform-dn} and a source-connected estimate for
      the mixed-jet card on the resulting compact \(b\)-interval.  A pure
      \(b=0\) expansion alone is insufficient.
\item Open V7 ancestry, YM--B radial transport, higher jets, or local-vacuum
      contraction only after the direct matched card passes or returns an
      attained physical failure needing attribution.
\end{enumerate}
If the numerical target or final generated insertion is still absent, the
correct output is the named loading \stopstatus.  It is not another request
for a general shell, current, or boundary compiler.
\end{openproblem}

% =============================================================================
\section{Version V11 proof ledger and append-only boundary}
\label{sec:YMA11-ledger}
% =============================================================================

The accompanying verifier
\begin{center}
\path{YMA_V11_response_chart_replay.py}
\end{center}
uses exact symbolic and rational arithmetic to check the root-factorization
and slope identity, the leading coefficient at odd multiplicities, one
root-disk slope enclosure, the implicit root jets, ramified model scaling,
the sign-safe coefficient interval, the quadratic mixed-jet identity, the
covariance-support statement on a literal positive Gram, the
near-ramification examples, and the three interval verdicts.  It also runs a
finite scalar check of the matched-card error radius.  The replay does not
evaluate a Wilson integral, select the physical response target, locate a
weak-facing root, or prove a cofinal estimate.

\begingroup\small
\begin{longtable}{>{\raggedright\arraybackslash}p{0.27\textwidth}
                  >{\raggedright\arraybackslash}p{0.20\textwidth}
                  >{\raggedright\arraybackslash}p{0.45\textwidth}}
\caption{Version V11 proof-status ledger.}\label{tab:YMA11-ledger}\\
\toprule
\textbf{Row}&\textbf{Status}&\textbf{Exact scope}\\
\midrule
\endfirsthead
\toprule
\textbf{Row}&\textbf{Status}&\textbf{Exact scope}\\
\midrule
\endhead
Unnormalized response differential&Proved exactly
&Four loop numerators and one inserted bank evaluate every interaction
response derivative; the partition denominator cancels.\\[0.3em]
Root slope and multiplicity&Proved exactly
&The V10 root function has the same zero order as \(\chi-u\), and its first
nonzero derivative fixes the stable crossing orientation.\\[0.3em]
Root-disk response floor&Proved quantitatively
&A one-zero complex disk, its boundary modulus, and the positive loop-product
window give explicit lower and upper bounds for \(\abs{\chi'}\).\\[0.3em]
Regular marginal coefficient&Proved and intervalized
&At a simple root \(b=-n/d\) is the root-path derivative; sign-safe corner
intervals and a slope-conditioned error bound are explicit.\\[0.3em]
Ramified perturbation branch&Proved
&A nonzero generated response gives a Hölder/Puiseux root displacement and no
finite tangent; a doubly zero first jet requires higher ramification data.\\[0.3em]
Complete matched Hessian&Validated finite enclosure
&Same-law deterministic and full mixed covariance blocks, coefficient error,
topology, and reference assemble into PASS/OBSTRUCTION/STOP.\\[0.3em]
Near-ramification amplification&Conditional physical theorem
&Persistent plaquette-score variance and a collapsing response slope force an
attained negative mode once the exact coefficient crosses the displayed
threshold.\\[0.3em]
Plaquette-score invisible branch&Classified exactly
&Positive covariance support removes mixed covariance; deterministic marginal
curvature or a field-invisible chart defect remains.\\[0.3em]
First selected response manifest&Unexecuted physical value
&No numerical \(u_*\), root multiplicity, generated response numerator, or
same-final-law mixed jet has been supplied.\\[0.3em]
Uniform response regularity&Open cofinal estimate
&A common slope floor and generated-response bound are required before a
compact-coefficient low-temperature theorem can be applied.\\[0.3em]
Selected cofinal symbol&\stopstatus
&Neither a complete executed response-matched card nor a uniform regular
mixed-jet estimate is present.\\[0.3em]
Higher jets, local-vacuum contraction, mass&Not inferred
&The downstream handoff remains closed until the regular matched quadratic
family passes.\\
\bottomrule
\end{longtable}
\endgroup

\begin{theorem}[Version V11 response-chart theorem]
\label{thm:YMA11-terminal}
Versions~V1--V10 are retained.  Version~V11 additionally proves:
\begin{enumerate}[label=\textup{(E11.\arabic*)},leftmargin=5.4em]
\item the response denominator and generated numerator are exact ratios of
      unnormalized marked loop numerators evaluated by the V10 Haar machinery;
\item the logarithm-free root function has the same multiplicity as the
      response level, and its derivative gives the response slope with exact
      sign;
\item a one-zero complex disk and its boundary modulus give a quantitative
      nonzero response-slope certificate;
\item a simple selected root has an analytic response-matching graph whose
      first derivative is the marginal coefficient and whose second derivative
      is the displayed mixed response curvature;
\item an odd ramified root with nonzero generated response has
      \(\abs{\beta(s)-\beta_*}\asymp\abs{s}^{1/m}\) and no finite ordinary
      marginal tangent;
\item outward response and same-law mixed-jet intervals give the complete
      finite matched Hessian and a terminating physical verdict;
\item the exact dependence on the marginal coefficient is
      \(H_{\rm pre}+bL-b^2C\), with \(C\succeq0\) and covariance support
      annihilating the mixed block on \(\Ker C\);
\item under persistent marginal-score visibility, collapse of the response
      slope with nonvanishing numerator forces a finite attained negative mode;
\item a uniform slope floor and numerator bound confine the later analytic
      problem to a compact coefficient family.
\end{enumerate}
No item supplies the missing numerical target, proves weak-facing-root escape,
evaluates the first root multiplicity or final-law mixed jet, proves a cofinal
connected expansion, or establishes local-vacuum contraction or a continuum
Yang--Mills mass gap.
\end{theorem}

\begin{proof}
The cited proofs are, in order:
\begin{itemize}
\item item~\textup{(E11.1)}:
      \Cref{prop:YMA11-response-numerator-card};
\item items~\textup{(E11.2)--(E11.3)}:
      \Cref{prop:YMA11-root-slope} and
      \Cref{thm:YMA11-root-disk-slope};
\item items~\textup{(E11.4)--(E11.5)}:
      \Cref{thm:YMA11-regular-root-chart} and
      \Cref{thm:YMA11-ramified-perturbation};
\item item~\textup{(E11.6)}:
      \Cref{prop:YMA11-b-interval} and
      \Cref{thm:YMA11-matched-enclosure};
\item item~\textup{(E11.7)}:
      \Cref{lem:YMA11-covariance-support} and
      \eqref{eq:YMA11-quadratic-b};
\item item~\textup{(E11.8)}:
      \Cref{thm:YMA11-b-blowup};
\item item~\textup{(E11.9)}:
      \Cref{prop:YMA11-uniform-regularity}.
\end{itemize}
\end{proof}

\begin{assessment}[V11 termination]
\label{ass:YMA11-termination}
The root-multiplicity, response-slope, marginal-coefficient, ramified-branch,
matched-card error, covariance-support, and coefficient-amplification results
are \pass\ at their stated finite mathematical scopes.  The programme has gone
one arrow upstream of the V10 computational fork: before compressing a
low-temperature series, it must establish that the selected root actually
supports the marginal response coordinate used to define the final action.
The first selected physical response manifest remains unpopulated, so the
current symbol verdict is \stopstatus.  No negative physical mode has been
produced from the supplied files.
\end{assessment}

\section*{Version V11 append-only continuation record}
\addcontentsline{toc}{section}{Version V11 append-only continuation record}
The complete Version~V10 source above is preserved byte for byte up to its sole
terminal document-closing command.  Its SHA--256 digest is
\begin{center}\scriptsize\ttfamily
9f40377d400527f2683dcaba8b75064a45dec8f04f951397652a2943f8c90618.
\end{center}
For Version~V12, preserve this accumulated source, append new mathematical
material immediately before the final document-closing command, and increment
the filename to \texttt{\_V12.tex}.  Keep the physical response target,
weak-facing selector, exact block map, parity-tree section, rooted Hodge
carrier, first axial mode, colour normalization, full unnormalized action,
generated insertion, marginal direction, and local/topological support
convention fixed before reading any root multiplicity, response slope,
coefficient, mixed jet, current, ancestry, radial, or transport value.

% =============================================================================
% END APPENDED VERSION V11 MATERIAL
% =============================================================================



% =============================================================================
% START APPENDED VERSION V12 MATERIAL
% =============================================================================

\clearpage
\hypersetup{
 pdftitle={Renewal Yang--Mills Reduced-Fibre Wilson Shell Symbol and Local-Vacuum Programme V12},
 pdfsubject={Selector-tangent versus marginal-direction separation, exact RG tangent transport, and the noncircular response chart}
}

\pdfinfo{
 /Title (Renewal Yang--Mills Reduced-Fibre Wilson Shell Symbol and Local-Vacuum Programme V12)
 /Subject (Selector-tangent versus marginal-direction separation, exact RG tangent transport, and the noncircular response chart)
}

% Three-digit section numbers begin in this continuation.  Widen only the
% subsequent section-number box when the table of contents is next read.
\addtocontents{toc}{\protect\makeatletter}
\makeatletter
\addtocontents{toc}{\protect\renewcommand*\protect\l@section{\protect\@dottedtocline{1}{1.5em}{2.8em}}}
\makeatother
\addtocontents{toc}{\protect\makeatother}

\section{Version V12: selector-tangent versus marginal-direction separation,
exact RG tangent transport, and the noncircular response chart}
\label{sec:YMA12-direction}

Version~V11 is preserved above.  Its root-multiplicity calculations are
correct for the one-parameter family which is actually differentiated.  The
specialization to the Yang--Mills shell, however, used one additional
identification which has not been supplied by the preceding programme: it
identified the derivative of the weak-facing selector with the independently
declared plaquette marginal interaction.

Those are generally different writers.  The weak-facing roots are selected
from the bare Wilson response family.  After exact shell integration, the
derivative of that family is the conditional expectation of the complete
fine action tangent.  It contains every generated interaction produced by the
shell.  By contrast, the denominator of the marginal coefficient is the
response to one fixed coarse plaquette interaction on the interaction
carrier.  A root slope measures the first writer; the coefficient
\(b^{\rm marg}\) divides by the response to the second.

The corrected order is
\begin{equation}
 \boxed{
 \begin{gathered}
 \text{frozen weak-facing root and the actual selector tangent}\\
 \Downarrow\\
 d_{\rm sel}=\mathscr R(T_{\rm sel})
 \quad\text{and independently}\quad
 d_P=\mathscr R(P)\\
 \Downarrow\\
 \text{exact tangent incidence }
 T_{\rm sel}=aP+\Psi+\Delta\\
 \Downarrow\\
 \text{selector chart or plaquette-marginal chart, according to the target}\\
 \Downarrow\\
 \text{complete final-law mixed field jet and the physical mode verdict}.
 \end{gathered}}
 \label{eq:YMA12-corrected-order}
\end{equation}
The new step is upstream rather than another root or current compiler.  It
identifies which derivative the existing root certificate has actually
measured, proves exact transport of that derivative through the Wilson shell,
and gives the minimum noncircular response manifest.

\begin{firewall}[A selector coordinate is not the marginal interaction]
\label{fire:YMA12-selector-not-marginal}
A scalar parameter may label a physical trajectory without multiplying the
coarse plaquette interaction alone.  Reparametrizing that scalar changes the
numerical selector slope, whereas the response to the fixed physical writer
\(P\) is unchanged.  Therefore a root derivative cannot be used as
\(\mathscr R(P)\) until an action-tangent incidence identity has been proved.
This is a loading statement, not a matter of choosing a more convenient
normalization after the root is observed.
\end{firewall}

% =============================================================================
\section{Exact transport of the actual selector tangent through a Wilson shell}
\label{sec:YMA12-selector-tangent}
% =============================================================================

Use the global product-Haar fibre already fixed in the inherited shell.  For
this section write the fine configuration space as
\(Y\times X_c\), with product probability
\(m_Y\otimes m_c\), and let
\begin{equation}
 F_\lambda:Y\times X_c\longrightarrow\R
 \label{eq:YMA12-fine-family}
\end{equation}
be a twice continuously differentiable finite action family.  Define
\begin{align}
 Z_\lambda(W)
 &:=\int_Y e^{-F_\lambda(z,W)}\,\dd m_Y(z),
 \label{eq:YMA12-fibre-Z}\\
 S_\lambda(W)
 &:=-\log Z_\lambda(W),
 \label{eq:YMA12-effective-S}\\
 \nu_{\lambda,W}(\dd z)
 &:=Z_\lambda(W)^{-1}e^{-F_\lambda(z,W)}\,\dd m_Y(z).
 \label{eq:YMA12-conditional-law}
\end{align}
Put
\begin{equation}
 T_\lambda:=\partial_\lambda F_\lambda,
 \qquad
 A_\lambda:=\partial_\lambda^2F_\lambda.
 \label{eq:YMA12-fine-tangent-acceleration}
\end{equation}

\begin{theorem}[Exact shell tangent and acceleration]
\label{thm:YMA12-shell-tangent}
For every coarse field \(W\),
\begin{equation}
 \boxed{
 \partial_\lambda S_\lambda(W)
 =\mathbb E_{\nu_{\lambda,W}}T_\lambda.}
 \label{eq:YMA12-effective-tangent}
\end{equation}
Moreover,
\begin{equation}
 \boxed{
 \partial_\lambda^2S_\lambda(W)
 =\mathbb E_{\nu_{\lambda,W}}A_\lambda
  -\operatorname{Var}_{\nu_{\lambda,W}}(T_\lambda).}
 \label{eq:YMA12-effective-acceleration}
\end{equation}
In particular, for the affine fine Wilson family
\(F_\lambda=F_0+\lambda V_f\),
\begin{equation}
 \boxed{
 T_{\rm sel}^{\rm eff}(W)
 :=\partial_\lambda S_\lambda(W)
 =\mathbb E_{\nu_{\lambda,W}}V_f,}
 \label{eq:YMA12-Wilson-selector-tangent}
\end{equation}
and
\begin{equation}
 \partial_\lambda^2S_\lambda(W)
 =-\operatorname{Var}_{\nu_{\lambda,W}}(V_f)\le0.
 \label{eq:YMA12-Wilson-selector-concavity}
\end{equation}
The selected coarse tangent is therefore the conditional fine Wilson energy,
not automatically the coarse plaquette interaction.
\end{theorem}

\begin{proof}
Differentiation under the finite compact integral gives
\[
 \partial_\lambda Z_\lambda(W)
 =-Z_\lambda(W)\mathbb E_{\nu_{\lambda,W}}T_\lambda.
\]
Division by \(-Z_\lambda\) proves
\eqref{eq:YMA12-effective-tangent}.  For any differentiable writer
\(G_\lambda\), direct differentiation of its conditional expectation gives
\[
 \partial_\lambda\mathbb E_{\nu_{\lambda,W}}G_\lambda
 =\mathbb E_{\nu_{\lambda,W}}\partial_\lambda G_\lambda
  -\operatorname{Cov}_{\nu_{\lambda,W}}(G_\lambda,T_\lambda).
\]
Apply this with \(G_\lambda=T_\lambda\) to obtain
\eqref{eq:YMA12-effective-acceleration}.  The affine Wilson specialization has
\(A_\lambda=0\), which gives the remaining formulas.
\end{proof}

\begin{corollary}[Multishell tangent is a conditional-expectation martingale]
\label{cor:YMA12-multishell-tangent}
Let
\(\mathbf B_{L\to j}:X_L\to X_j\) be a finite nested sequence of exact block
maps and let \(S_{\lambda,j}\) be the exact action obtained at level \(j\)
from one terminal action \(F_{\lambda,L}\).  Then
\begin{equation}
 \boxed{
 \partial_\lambda S_{\lambda,j}(W_j)
 =\mathbb E_{\lambda}
   \left[
    \partial_\lambda F_{\lambda,L}
    \mid \mathbf B_{L\to j}U=W_j
   \right].}
 \label{eq:YMA12-multishell-tangent}
\end{equation}
For \(i<j<L\), these tangents obey the tower identity
\begin{equation}
 \boxed{
 \partial_\lambda S_{\lambda,i}
 =\mathbb E_\lambda
   [\partial_\lambda S_{\lambda,j}\mid X_i].}
 \label{eq:YMA12-tangent-tower}
\end{equation}
Thus intermediate generated interactions are already part of the actual
selector tangent; they are not discarded by flattening the tower.
\end{corollary}

\begin{proof}
Apply \Cref{thm:YMA12-shell-tangent} to the direct terminal-to-level-\(j\)
fibre.  Applying it first at level \(j\) and then at level \(i\) gives the
iterated conditional expectation of the same terminal tangent.  The ordinary
tower property proves \eqref{eq:YMA12-tangent-tower}.
\end{proof}

Let \(O_\alpha\), \(\alpha\in\{11,00,10,01\}\), be the four positive loop
writers used by the response, now regarded as functions on \(X_c\).  Define
\begin{equation}
 N_\alpha^{c}(\lambda)
 :=\int_{X_c}O_\alpha(W)e^{-S_\lambda(W)}\,\dd m_c(W).
 \label{eq:YMA12-coarse-loop-numerator}
\end{equation}
Their fine pullbacks give the identical numerators
\begin{equation}
 N_\alpha^{f}(\lambda)
 :=\int_{Y\times X_c}
 O_\alpha(W)e^{-F_\lambda(z,W)}\,
 \dd m_Y(z)\dd m_c(W)
 =N_\alpha^{c}(\lambda).
 \label{eq:YMA12-fine-loop-numerator}
\end{equation}

\begin{proposition}[Direct/staged response-slope identity]
\label{prop:YMA12-response-slope-transport}
Let
\begin{equation}
 T_\lambda^{c}(W):=\partial_\lambda S_\lambda(W).
 \label{eq:YMA12-coarse-selector-tangent}
\end{equation}
Then every marked numerator satisfies
\begin{equation}
 \boxed{
 \int O_\alpha T_\lambda^{c}e^{-S_\lambda}\,\dd m_c
 =\int O_\alpha T_\lambda e^{-F_\lambda}\,
   \dd m_Y\dd m_c.}
 \label{eq:YMA12-marked-tangent-transport}
\end{equation}
Consequently the derivative of the four-loop response may be evaluated
before or after shell integration:
\begin{equation}
 \boxed{
 \partial_\lambda\chi(\lambda)
 =\mathscr R_{S_\lambda}(T_\lambda^{c})
 =\mathscr R_{F_\lambda}(T_\lambda).}
 \label{eq:YMA12-response-slope-transport}
\end{equation}
For the bare Wilson selector, the V10--V11 root derivative is therefore the
response to the conditional fine Wilson energy
\eqref{eq:YMA12-Wilson-selector-tangent}.
\end{proposition}

\begin{proof}
Multiply \eqref{eq:YMA12-effective-tangent} by
\(O_\alpha e^{-S_\lambda}\) and integrate in \(W\).  Substitution of
\eqref{eq:YMA12-conditional-law} gives
\eqref{eq:YMA12-marked-tangent-transport}.  Differentiate the signed sum of
logarithms and use the unnormalized numerator formula of
\Cref{prop:YMA11-response-numerator-card} on the two sides.
\end{proof}

% =============================================================================
\section{Three distinct response rows}
\label{sec:YMA12-three-response-rows}
% =============================================================================

Fix a selected point \(\lambda_*\) and let the coarse action be perturbed by
an independently declared plaquette interaction \(P\) and generated
interaction \(\Psi\):
\begin{equation}
 S_{\lambda,c,s}
 :=S_\lambda+cP+s\Psi.
 \label{eq:YMA12-three-parameter-action}
\end{equation}
All derivatives in this section are evaluated at
\((\lambda,c,s)=(\lambda_*,0,0)\).  Put
\begin{equation}
 T_{\rm sel}:=\partial_\lambda S_\lambda\big|_{\lambda_*}.
 \label{eq:YMA12-selector-writer}
\end{equation}

\begin{definition}[Selector, marginal, and generated response rows]
\label{def:YMA12-response-triple}
The three first response coefficients are
\begin{equation}
 \boxed{
 d_{\rm sel}:=\mathscr R(T_{\rm sel}),
 \qquad
 d_P:=\mathscr R(P),
 \qquad
 n_\Psi:=\mathscr R(\Psi).}
 \label{eq:YMA12-response-triple}
\end{equation}
When the corresponding denominators are nonzero, define
\begin{equation}
 b^{\rm sel}:=-\frac{n_\Psi}{d_{\rm sel}},
 \qquad
 b^{\rm marg}:=-\frac{n_\Psi}{d_P}.
 \label{eq:YMA12-two-coefficients}
\end{equation}
The first is a change in the selector coordinate per unit generated
perturbation.  The second is the coefficient of the physical plaquette
interaction in the response-matched action.
\end{definition}

\begin{theorem}[The root certificate measures only the selector row]
\label{thm:YMA12-root-measures-selector}
For the family \eqref{eq:YMA12-three-parameter-action},
\begin{equation}
 \boxed{
 \partial_\lambda\chi=d_{\rm sel},
 \qquad
 \partial_c\chi=d_P,
 \qquad
 \partial_s\chi=n_\Psi.}
 \label{eq:YMA12-three-partials}
\end{equation}
The zero multiplicity and boundary-modulus estimates of the V10--V11
weak-facing root function determine \(d_{\rm sel}\).  They determine
\(d_P\) only on an additional branch on which
\begin{equation}
 \boxed{T_{\rm sel}=P+q}
 \label{eq:YMA12-T-equals-P-mod-q}
\end{equation}
with a writer \(q\) satisfying \(\mathscr R(q)=0\); equality of the full
action tangents requires the stronger statement that every field jet of
\(q\) used by the target also vanishes or is retained.
\end{theorem}

\begin{proof}
The three equalities are the response derivative formula applied separately
to the three action coordinates.  The weak-facing root varies \(\lambda\)
with \(c=s=0\), so its derivative and multiplicity concern the first partial
only.  Under \eqref{eq:YMA12-T-equals-P-mod-q}, linearity gives
\(d_{\rm sel}=d_P+\mathscr R(q)=d_P\).  Response-nullity of \(q\) is only one
scalar identity and does not annihilate its field jets; this is precisely the
two-carrier distinction proved in the inherited Version~V33 theorem layer.
\end{proof}

\begin{proposition}[Coordinate covariance separates the two coefficients]
\label{prop:YMA12-coordinate-covariance}
Let \(\lambda=g(\tau)\) be a real-analytic local reparametrization with
\(g'(\tau_*)\ne0\).  Then
\begin{equation}
 \boxed{
 d_{\rm sel}^{(\tau)}=g'(\tau_*)d_{\rm sel}^{(\lambda)},
 \qquad
 d_P^{(\tau)}=d_P,
 \qquad
 n_\Psi^{(\tau)}=n_\Psi.}
 \label{eq:YMA12-coordinate-transform}
\end{equation}
Accordingly,
\begin{equation}
 b_{\tau}^{\rm sel}
 =\frac{b_{\lambda}^{\rm sel}}{g'(\tau_*)},
 \qquad
 b^{\rm marg}_{\tau}=b^{\rm marg}_{\lambda}.
 \label{eq:YMA12-coefficient-transform}
\end{equation}
The selector correction transforms as a parameter derivative; the plaquette
coefficient is invariant because its writer has been physically fixed.
\end{proposition}

\begin{proof}
The chain rule gives
\(\partial_\tau S_{g(\tau)}=g'(\tau)T_{\rm sel}\).  Apply the linear response
functional.  The writers \(P\) and \(\Psi\) are unchanged by a relabeling of
the selector coordinate.  Substitution in
\eqref{eq:YMA12-two-coefficients} proves the last identities.
\end{proof}

\begin{countertheorem}[Selector regularity and marginal regularity are independent]
\label{cth:YMA12-selector-marginal-independent}
Each of the following is realized by positive analytic loop-numerator germs.
At the base point \((\lambda,c,s)=(0,0,0)\), take
\(N_{11}=e^{-\chi}\) and
\(N_{00}=N_{10}=N_{01}=1\), so that the four-loop response is exactly
\(\chi\).
\begin{enumerate}[label=\textup{(I\arabic*)},leftmargin=3.8em]
\item For
      \begin{equation}
       \chi_A=u-\lambda+s+c^2,
       \label{eq:YMA12-counter-A}
      \end{equation}
      the selector root is a simple stable down-crossing,
      \(d_{\rm sel}=-1\), but \(d_P=0\) and the plaquette marginal chart is
      singular.
\item For
      \begin{equation}
       \chi_B=u-\lambda^3+c-s,
       \label{eq:YMA12-counter-B}
      \end{equation}
      the selector root has multiplicity three,
      \(d_{\rm sel}=0\), while \(d_P=1\), \(n_\Psi=-1\), and
      \(b^{\rm marg}=1\) is regular.
\item For every \(\epsilon>0\),
      \begin{equation}
       \chi_{C,\epsilon}=u-\lambda+\epsilon c+s
       \label{eq:YMA12-counter-C}
      \end{equation}
      has the uniform selector slope \(-1\), but
      \(b^{\rm marg}=-\epsilon^{-1}\) is unbounded as
      \(\epsilon\downarrow0\).
\item For every \(\epsilon>0\),
      \begin{equation}
       \chi_{D,\epsilon}=u-\epsilon\lambda+c-s
       \label{eq:YMA12-counter-D}
      \end{equation}
      has \(d_{\rm sel}=-\epsilon\to0\), while
      \(d_P=1\), \(n_\Psi=-1\), and \(b^{\rm marg}=1\) remain fixed.
\end{enumerate}
Thus a simple weak-facing root does not certify the marginal denominator, a
ramified selector root does not preclude a regular marginal chart, a uniform
selector-slope floor does not bound the plaquette coefficient, and collapse
of the selector slope does not by itself amplify that coefficient.
\end{countertheorem}

\begin{proof}
All four numerator germs are strictly positive.  Their signed logarithmic
combination is the displayed \(\chi\).  Direct differentiation gives every
coefficient.  The odd negative leading powers in \(\lambda\) give stable
down-crossings in the first two examples; the last two are simple stable
down-crossings for every \(\epsilon>0\).
\end{proof}

% =============================================================================
\section{Two legitimate response charts and their noninterchangeability}
\label{sec:YMA12-two-charts}
% =============================================================================

\begin{theorem}[Selector chart versus plaquette-marginal chart]
\label{thm:YMA12-two-charts}
Assume \(\chi(\lambda_*,0,0)=u\).
\begin{enumerate}[label=\textup{(C\arabic*)},leftmargin=3.8em]
\item If \(d_{\rm sel}\ne0\), there is a unique analytic selector graph
      \(\lambda(s)\) with \(c=0\), and
      \begin{equation}
       \boxed{\lambda'(0)=b^{\rm sel}=-n_\Psi/d_{\rm sel}.}
       \label{eq:YMA12-selector-chart-jet}
      \end{equation}
\item If \(d_P\ne0\), there is a unique analytic marginal graph \(c(s)\)
      with \(\lambda=\lambda_*\), and
      \begin{equation}
       \boxed{c'(0)=b^{\rm marg}=-n_\Psi/d_P.}
       \label{eq:YMA12-marginal-chart-jet}
      \end{equation}
\item If both denominators are nonzero, the two response-null tangent writers
      \begin{equation}
       I_{\rm sel}:=\Psi+b^{\rm sel}T_{\rm sel},
       \qquad
       I_{\rm marg}:=\Psi+b^{\rm marg}P
       \label{eq:YMA12-two-null-writers}
      \end{equation}
      obey
      \begin{equation}
       \boxed{
       \mathscr R(I_{\rm sel})=\mathscr R(I_{\rm marg})=0,
       \qquad
       \mathscr R(I_{\rm marg}-I_{\rm sel})=0.}
       \label{eq:YMA12-null-difference}
      \end{equation}
      Their coefficients satisfy
      \begin{equation}
       \boxed{
       b^{\rm marg}
       =\frac{d_{\rm sel}}{d_P}b^{\rm sel}.}
       \label{eq:YMA12-chart-conversion}
      \end{equation}
\end{enumerate}
The two tangent writers need not have the same field jets.  The programme
therefore uses the marginal chart only because \(P\) is the predeclared
physical interaction direction; the selector chart cannot replace it merely
because both tangents preserve the scalar response.
\end{theorem}

\begin{proof}
The first two statements are the analytic implicit-function theorem applied
to \(\chi-u\) in the indicated coordinate.  Differentiation gives the two
first jets.  Linearity of \(\mathscr R\) and the definitions of the
coefficients prove \eqref{eq:YMA12-null-difference}.  The conversion formula
follows by eliminating \(n_\Psi\).  The final statement follows from the
inherited two-carrier counterexample: a response-null interaction can have a
nonzero positive or negative field Hessian.
\end{proof}

\begin{assessment}[What selector ramification actually obstructs]
\label{ass:YMA12-selector-ramification-scope}
A ramified weak-facing root blocks an ordinary analytic inverse in the bare
selector coordinate and hence affects adjacent bare-coupling transport.  It
does not, by itself, obstruct a marginally response-matched shell action at
the already selected point.  That action is regular whenever
\(d_P\ne0\).  Conversely, a simple root does not rescue a vanishing
\(d_P\).  The trajectory and interaction-chart obstructions are therefore
separate rows.
\end{assessment}

% =============================================================================
\section{The actual tangent-incidence residual and a two-row reduction}
\label{sec:YMA12-tangent-incidence}
% =============================================================================

Choose before acquisition a scalar \(a\), representing the coefficient of
\(P\) in the declared interaction basis, and define the action-tangent
residual
\begin{equation}
 \boxed{
 \Delta_{\rm tan}
 :=T_{\rm sel}-aP-\Psi.}
 \label{eq:YMA12-tangent-residual}
\end{equation}
Its response residual is
\begin{equation}
 r_{\rm tan}:=\mathscr R(\Delta_{\rm tan}).
 \label{eq:YMA12-tangent-response-residual}
\end{equation}

\begin{theorem}[Noncircular tangent-incidence identity]
\label{thm:YMA12-tangent-incidence}
The three response rows satisfy the exact identity
\begin{equation}
 \boxed{
 d_{\rm sel}=a d_P+n_\Psi+r_{\rm tan}.}
 \label{eq:YMA12-tangent-response-identity}
\end{equation}
Assume \(d_P\ne0\).  The coefficient which makes the complete nonmarginal
selector tangent
\(T_{\rm sel}-aP=\Psi+\Delta_{\rm tan}\) response null is
\begin{equation}
 \boxed{
 b_{\rm full}^{\rm marg}
 =-\frac{n_\Psi+r_{\rm tan}}{d_P}
 =a-\frac{d_{\rm sel}}{d_P}.}
 \label{eq:YMA12-full-b}
\end{equation}
The coefficient obtained by matching only the submitted \(\Psi\) is
\begin{equation}
 b_{\Psi}^{\rm marg}=-\frac{n_\Psi}{d_P},
 \label{eq:YMA12-submitted-b}
\end{equation}
and the exact coefficient defect is
\begin{equation}
 \boxed{
 b_{\rm full}^{\rm marg}-b_{\Psi}^{\rm marg}
 =-\frac{r_{\rm tan}}{d_P}.}
 \label{eq:YMA12-b-defect}
\end{equation}
Hence a nonzero scalar incidence residual already rejects the claimed
response matching.  If \(r_{\rm tan}=0\), scalar response matching agrees,
but the physical field card still requires either
\(\Delta_{\rm tan}=0\) modulo a field-independent constant, or explicit
retention of its first and second field jets.
\end{theorem}

\begin{proof}
Apply the linear functional \(\mathscr R\) to
\eqref{eq:YMA12-tangent-residual}.  This gives
\eqref{eq:YMA12-tangent-response-identity}.  The response of
\(\Psi+\Delta_{\rm tan}+bP\) is
\(n_\Psi+r_{\rm tan}+bd_P\); setting it to zero gives the first expression
in \eqref{eq:YMA12-full-b}.  Substitution of
\eqref{eq:YMA12-tangent-response-identity} gives the second.  Subtracting
\eqref{eq:YMA12-submitted-b} proves \eqref{eq:YMA12-b-defect}.  Finally,
scalar response-nullity does not annihilate routed field jets, by the
interaction/field separation already proved in Version~V33 and retained in
Versions~V8--V11.
\end{proof}

\begin{corollary}[Exact two-response-row acquisition]
\label{cor:YMA12-two-row-reduction}
Suppose the flattened exact shell proves the action identity
\begin{equation}
 \boxed{
 T_{\rm sel}=aP+\Psi+c\one}
 \label{eq:YMA12-exact-tangent-decomposition}
\end{equation}
for a scalar \(c\) independent of the fibre and coarse field.  Then
\(r_{\rm tan}=0\), all physical field jets of the residual vanish, and
\begin{equation}
 \boxed{
 n_\Psi=d_{\rm sel}-a d_P,
 \qquad
 b^{\rm marg}=a-\frac{d_{\rm sel}}{d_P}.}
 \label{eq:YMA12-two-row-formula}
\end{equation}
Thus only the selector response row and the independent plaquette response row
are required; a separate generated-insertion numerator bank is redundant on
this exact incidence branch.
\end{corollary}

\begin{proof}
The response of a constant insertion vanishes because the four Creutz signs
sum to zero.  A constant also has zero coarse field derivatives.  Apply
\Cref{thm:YMA12-tangent-incidence}.
\end{proof}

\begin{countertheorem}[Exact shell integration need not preserve a one-plaquette tangent]
\label{cth:YMA12-effective-tangent-not-P}
Let the coarse space be \(X_c=\{0,1,2\}\), the fibre be
\(Y=\{0,1\}\), and both carry uniform measure.  Set
\(P(w)=w\) and define the affine fine action
\(F_\lambda(y,w)=\lambda V(y,w)\) by
\begin{equation}
 \begin{array}{c|cc}
 w&V(0,w)&V(1,w)\\ \hline
 0&0&0\\
 1&1&1\\
 2&0&2
 \end{array}.
 \label{eq:YMA12-finite-V-table}
\end{equation}
Then the exact coarse selector tangent is
\begin{equation}
 \boxed{
 T_{\rm sel}^{\rm eff}(0)=0,
 \qquad
 T_{\rm sel}^{\rm eff}(1)=1,
 \qquad
 T_{\rm sel}^{\rm eff}(2)
 =\frac{2e^{-2\lambda}}{1+e^{-2\lambda}}.}
 \label{eq:YMA12-finite-effective-tangent}
\end{equation}
At \(\lambda=0\) this is the vector \((0,1,1)\), which does not belong to
\(\operatorname{span}\{\one,P\}\).  Thus even an affine microscopic
coupling can generate a coarse selector tangent outside the declared
one-plaquette line.
\end{countertheorem}

\begin{proof}
The first two fibres have constant energies zero and one.  On the third fibre,
the conditional weights are proportional to \(1\) and \(e^{-2\lambda}\), so
\Cref{thm:YMA12-shell-tangent} gives the displayed conditional mean.  At
\(\lambda=0\), an affine function \(c+aP\) agreeing with the first two
entries must have \(c=0\) and \(a=1\), hence value two at \(w=2\), not one.
\end{proof}

\begin{firewall}[A scalar response check is weaker than action incidence]
\label{fire:YMA12-scalar-incidence-scope}
The equality \(r_{\rm tan}=0\) says only that the omitted tangent residual is
in the kernel of one scalar response.  It is sufficient for the scalar
calibration target, but not for the field Hessian.  The full physical card
must use the complete action or retain the residual jets.  A response-null
writer is not deleted from the field carrier.
\end{firewall}

% =============================================================================
\section{Explicit qualification of the Version V11 specialization}
\label{sec:YMA12-V11-qualification}
% =============================================================================

\begin{theorem}[Correct reading of the V11 response-chart results]
\label{thm:YMA12-V11-qualification}
The Version~V11 results remain valid with the following exact typing.
\begin{enumerate}[label=\textup{(Q\arabic*)},leftmargin=3.8em]
\item \Cref{prop:YMA11-root-slope,thm:YMA11-root-disk-slope} compute and
      bound \(d_{\rm sel}=\mathscr R(T_{\rm sel})\), the derivative of the
      action family used by the weak-facing selector.
\item \Cref{thm:YMA11-regular-root-chart} computes
      \(b^{\rm sel}=-n_\Psi/d_{\rm sel}\), the selector-coordinate correction.
      It computes \(b^{\rm marg}\) only under an independently proved
      incidence \(T_{\rm sel}=P\) modulo a response- and field-invisible
      constant.
\item \Cref{thm:YMA11-ramified-perturbation} classifies ramification of the
      selector chart.  It does not by itself classify the plaquette-marginal
      chart.
\item \Cref{prop:YMA11-b-interval,thm:YMA11-matched-enclosure} remain the
      correct finite interval and physical Hessian theorems when their input
      coefficient is the independently acquired
      \(b^{\rm marg}=-n_\Psi/d_P\).
\item The amplification theorem
      \Cref{thm:YMA11-b-blowup} applies when the \emph{marginal denominator}
      \(d_P\) tends to zero with visible numerator and field variance.  A
      collapse of \(d_{\rm sel}\) alone does not trigger it.
\item The compact-coefficient condition of
      \Cref{prop:YMA11-uniform-regularity} must read
      \begin{equation}
       \boxed{
       \abs{d_{P,r}}\ge d_*>0,
       \qquad
       \abs{n_{\Psi,r}}\le n^*<\infty.}
       \label{eq:YMA12-correct-uniform-regularity}
      \end{equation}
\end{enumerate}
No previous formula for an already supplied \(b^{\rm marg}\) or final-law
mixed Hessian is altered.  What is withdrawn is the unsupported use of the
bare selector root slope as the plaquette response denominator.
\end{theorem}

\begin{proof}
Items~\textup{(Q1)--(Q3)} are
\Cref{thm:YMA12-root-measures-selector,thm:YMA12-two-charts}.  Item~\textup{(Q4)}
uses only the algebraic coefficient and same-law jet intervals and is
therefore unchanged.  Items~\textup{(Q5)--(Q6)} follow from the exact
coefficient definition \(b^{\rm marg}=-n_\Psi/d_P\) and the independence
examples in \Cref{cth:YMA12-selector-marginal-independent}.
\end{proof}

% =============================================================================
\section{Outward acquisition and the noncircular finite manifest}
\label{sec:YMA12-manifest}
% =============================================================================

For a positive denominator interval \(0<L\le N\le U\) and a signed inserted
numerator interval \(\ell\le M\le u\), define the sign-safe quotient hull
\begin{equation}
 \mathcal Q([\ell,u],[L,U])
 :=\left[
  \min\left\{\frac\ell L,\frac\ell U,\frac uL,\frac uU\right\},
  \max\left\{\frac\ell L,\frac\ell U,\frac uL,\frac uU\right\}
 \right].
 \label{eq:YMA12-quotient-hull}
\end{equation}
For intervals \(I,J\), use ordinary outward Minkowski addition and scalar
multiplication.

\begin{proposition}[Marked-numerator interval for every response row]
\label{prop:YMA12-marked-response-interval}
Suppose, at one certified root box, each loop numerator and one inserted bank
obey
\begin{equation}
 0<L_\alpha\le N_\alpha\le U_\alpha,
 \qquad
 \ell_{\alpha;X}\le N_{\alpha;X}\le u_{\alpha;X}.
 \label{eq:YMA12-marked-inputs}
\end{equation}
Then the response to the writer \(X\) has the rigorous interval
\begin{equation}
 \boxed{
 \mathscr R(X)
 \in
 -\sum_\alpha\varepsilon_\alpha
 \mathcal Q(
  [\ell_{\alpha;X},u_{\alpha;X}],
  [L_\alpha,U_\alpha]).}
 \label{eq:YMA12-response-interval}
\end{equation}
The writer \(X=T_{\rm sel}\) may be evaluated either by the coarse effective
tangent or by the pulled-back fine tangent, with identical exact value by
\Cref{prop:YMA12-response-slope-transport}.  The writer \(X=P\) is an
independent marked bank unless the tangent-incidence theorem has already
identified it.
\end{proposition}

\begin{proof}
Each ratio lies in the corner hull
\eqref{eq:YMA12-quotient-hull}.  The response is the signed sum in
\Cref{prop:YMA11-response-numerator-card}; outward interval addition proves
the claim.  The direct/staged statement is
\eqref{eq:YMA12-marked-tangent-transport}.
\end{proof}

\begin{theorem}[Noncircular finite response manifest]
\label{thm:YMA12-noncircular-manifest}
A finite selected-root response manifest consists of:
\begin{enumerate}[label=\textup{(M\arabic*)},leftmargin=3.8em]
\item the V10 root box, zero count, stable-crossing orientation, and hence the
      selector multiplicity;
\item the actual selector tangent, supplied either as the flattened coarse
      writer \(T_{\rm sel}\) or as the terminal fine tangent together with
      \Cref{prop:YMA12-response-slope-transport};
\item the independently declared plaquette writer \(P\), generated writer
      \(\Psi\), and any predeclared coefficient \(a\) used in their tangent
      decomposition;
\item outward response intervals for \(d_{\rm sel}\), \(d_P\), and either
      \(n_\Psi\) or the tangent residual \(r_{\rm tan}\);
\item the complete action-incidence status of \(\Delta_{\rm tan}\) for the
      field jets required by the physical Hessian.
\end{enumerate}
After input-domain checks, the first applicable mathematical outcome is:
\begin{enumerate}[label=\textup{(R\arabic*)},leftmargin=3.8em]
\item If an outward interval proves \(d_P\ne0\), compute
      \(b^{\rm marg}\) by the V11 sign-safe ratio or by
      \eqref{eq:YMA12-full-b}, retain every nonzero tangent residual jet, and
      continue to the final-law mixed Hessian.
\item If \(d_P=0\) is proved and \(n_\Psi\ne0\), return the typed
      \emph{plaquette response-chart obstruction}: no finite first-order
      marginal coefficient exists.
\item If \(d_P=n_\Psi=0\) is proved, the first-order marginal coefficient is
      nonunique; acquire the least nonzero mixed response jet and return
      \stopstatus\ until it is populated.
\item If the available intervals meet zero without proving an equality, return
      \stopstatus\ at the first unresolved response row.
\end{enumerate}
Selector ramification is recorded independently.  It neither upgrades nor
downgrades the marginal-chart verdict without the tangent-incidence row.
\end{theorem}

\begin{proof}
The root packet determines only the selector coordinate by
\Cref{thm:YMA12-root-measures-selector}.  The direct/staged tangent identity
makes item~\textup{(M2)} exact.  The response intervals are
\Cref{prop:YMA12-marked-response-interval}.  On a nonzero \(d_P\) branch,
the ordinary implicit-function theorem in the \(P\)-coordinate and the V11
interval arithmetic give item~\textup{(R1)}.  When \(d_P=0\), the first-order
response of \(\Psi+bP\) equals \(n_\Psi\) for every finite \(b\), proving
item~\textup{(R2)}.  If both vanish, every \(b\) is first-order null and no
unique coefficient is selected, proving item~\textup{(R3)}.  The final branch
is the logic of outward intervals.  The independence from selector
ramification is \Cref{cth:YMA12-selector-marginal-independent}.
\end{proof}

\begin{proposition}[Correct cofinal compact-coefficient hypotheses]
\label{prop:YMA12-cofinal-coefficient}
Let a selected cofinal family have marginal response rows
\(d_{P,r}\), generated rows \(n_{\Psi,r}\), and coefficients
\(b_r^{\rm marg}\).  If
\begin{equation}
 \boxed{
 \inf_r\abs{d_{P,r}}\ge d_*>0,
 \qquad
 \sup_r\abs{n_{\Psi,r}}\le n^*<\infty,}
 \label{eq:YMA12-cofinal-dP-n}
\end{equation}
then
\begin{equation}
 \boxed{
 \sup_r\abs{b_r^{\rm marg}}\le n^*/d_*.}
 \label{eq:YMA12-cofinal-b-bound}
\end{equation}
Alternatively, on the exact tangent-incidence branch
\(T_{{\rm sel},r}=a_rP_r+\Psi_r+c_r\one\), the hypotheses
\begin{equation}
 \sup_r\abs{a_r}\le a^*,
 \qquad
 \sup_r\abs{d_{{\rm sel},r}}\le s^*,
 \qquad
 \inf_r\abs{d_{P,r}}\ge d_*>0
 \label{eq:YMA12-cofinal-two-row-input}
\end{equation}
imply
\begin{equation}
 \boxed{
 \sup_r\abs{b_r^{\rm marg}}
 \le a^*+s^*/d_*.}
 \label{eq:YMA12-cofinal-two-row-bound}
\end{equation}
No lower bound on \(\abs{d_{{\rm sel},r}}\) is required for either
coefficient bound.  Such a lower bound belongs to analytic transport of the
bare selector, not to regularity of the plaquette marginal chart.
\end{proposition}

\begin{proof}
The first estimate is the definition
\(b_r^{\rm marg}=-n_{\Psi,r}/d_{P,r}\).  On the exact tangent-incidence
branch, use \eqref{eq:YMA12-two-row-formula} and the triangle inequality.
The independence statement is witnessed explicitly by items~\textup{(I2)}
and \textup{(I4)} of
\Cref{cth:YMA12-selector-marginal-independent}.
\end{proof}

% =============================================================================
\section{The revised first physical acquisition}
\label{sec:YMA12-frontier}
% =============================================================================

\begin{assessment}[What the upstream correction removes]
\label{ass:YMA12-progress}
The weak-facing root machinery need not be enlarged.  Its root box,
multiplicity, and boundary derivative already measure the actual selector
path.  They do not measure the coarse plaquette marginal direction.  V12
therefore removes a circular proposed use of the same derivative for two
different writers and reduces the next finite acquisition to one genuinely
new marked bank: the response to \(P\), together with an exact tangent
incidence or its retained residual.

If the exact flattened tangent decomposes as
\(aP+\Psi+c\one\), the generated response numerator is recovered from two
rows by \eqref{eq:YMA12-two-row-formula}.  If it does not, the residual is an
actual interaction writer and must remain in the field card.  No current,
boundary, radial, or higher-jet theorem can repair that missing action
incidence.
\end{assessment}

\begin{openproblem}[V13: execute the independent plaquette response row]
\label{op:YMA12-V13}
Preserve V12 and all frozen physical choices.  Proceed in this order.
\begin{enumerate}[label=\textup{(V13.\arabic*)},leftmargin=5.7em]
\item On the first selected weak-facing root box, retain the V10 complex root
      certificate as the selector packet.  Do not reinterpret its derivative
      as \(\mathscr R(P)\).
\item Differentiate the exact flattened Wilson shell with respect to the
      actual bare selector and acquire the marked tangent either before or
      after shell integration, checking
      \eqref{eq:YMA12-marked-tangent-transport}.
\item Acquire the four \(P\)-inserted loop numerators and certify an outward
      interval for \(d_P\) by
      \eqref{eq:YMA12-response-interval}.
\item Compare the actual selector tangent with the declared
      \(aP+\Psi\).  If the identity holds modulo a constant, use the two-row
      coefficient \eqref{eq:YMA12-two-row-formula}.  Otherwise acquire
      \(r_{\rm tan}\) and retain every nonzero residual field jet.
\item Apply the marginal-chart alternative of
      \Cref{thm:YMA12-noncircular-manifest}.  A ramified selector root is not
      substituted for a vanishing plaquette denominator, and a simple root is
      not used as its proof.
\item On the regular marginal branch, populate the complete same-final-law
      mixed jet of Version~V11 and return the physical
      \pass/\obstruction/\stopstatus\ verdict.
\end{enumerate}
If neither the coarse plaquette writer nor the exact selector tangent is
available as an executable insertion, return the named action-incidence
\stopstatus.  Do not append another root, ratio, current, or boundary
compiler.
\end{openproblem}

% =============================================================================
\section{Version V12 proof ledger and append-only boundary}
\label{sec:YMA12-ledger}
% =============================================================================

The accompanying exact replay
\begin{center}
\path{YMA_V12_selector_marginal_replay.py}
\end{center}
checks a finite direct/staged conditional-expectation model, the effective
tangent and acceleration formulas, coordinate covariance, all four
selector/marginal independence germs, the two-chart conversion, the tangent
incidence and coefficient-defect identities, the finite counterexample to a
one-plaquette effective tangent, and the interval branch logic.  It does not
evaluate a Wilson path integral, select \(u_*\), populate a physical
plaquette response row, or prove a cofinal margin.

\begingroup\small
\begin{longtable}{>{\raggedright\arraybackslash}p{0.27\textwidth}
                  >{\raggedright\arraybackslash}p{0.20\textwidth}
                  >{\raggedright\arraybackslash}p{0.45\textwidth}}
\caption{Version V12 proof-status ledger.}\label{tab:YMA12-ledger}\\
\toprule
\textbf{Row}&\textbf{Status}&\textbf{Exact scope}\\
\midrule
\endfirsthead
\toprule
\textbf{Row}&\textbf{Status}&\textbf{Exact scope}\\
\midrule
\endhead
Exact shell selector tangent&Proved
&The derivative of the effective action is the conditional terminal action
tangent; its acceleration is mean acceleration minus tangent variance.\\[0.3em]
Direct/staged selector response&Proved exactly
&Every loop-marked selector derivative agrees before and after exact shell
integration.\\[0.3em]
Selector versus marginal rows&Separated exactly
&The root slope is \(\mathscr R(T_{\rm sel})\); the marginal denominator is
\(\mathscr R(P)\).\\[0.3em]
Coordinate covariance&Proved
&Selector slope rescales with its parameter; the physical plaquette
coefficient is invariant.\\[0.3em]
Regularity implications&Disproved in both directions
&Simple selector roots may have singular marginal charts, and ramified
selector roots may have regular marginal charts.\\[0.3em]
Two response charts&Proved
&Selector and plaquette corrections are distinct response-null tangents and
need not have the same field jets.\\[0.3em]
Tangent-incidence identity&Proved exactly
&\(d_{\rm sel}=ad_P+n_\Psi+r_{\rm tan}\), with the exact coefficient defect
\(-r_{\rm tan}/d_P\).\\[0.3em]
Two-row coefficient reduction&Proved conditionally
&On an exact action-tangent decomposition modulo constants,
\(b^{\rm marg}=a-d_{\rm sel}/d_P\).\\[0.3em]
V11 root specialization&Qualified explicitly
&Root-disk estimates concern the selector row; V11 mixed-Hessian formulas
remain valid once the independent marginal coefficient is supplied.\\[0.3em]
Independent plaquette response&Open physical value
&No selected-root \(P\)-inserted numerator bank is present in the supplied
files.\\[0.3em]
Action-tangent incidence&Open physical loading
&The exact flattened selector tangent has not been decomposed into the declared
marginal and generated interaction writers on the first selected card.\\[0.3em]
Selected response-matched symbol&\stopstatus
&The final-law mixed jet cannot be interpreted until the marginal chart and
action incidence are populated.\\[0.3em]
Higher jets, local-vacuum contraction, mass&Not inferred
&No downstream promotion is opened from selector regularity alone.\\
\bottomrule
\end{longtable}
\endgroup

\begin{theorem}[Version V12 selector/marginal separation theorem]
\label{thm:YMA12-terminal}
Versions~V1--V11 are retained.  Version~V12 additionally proves:
\begin{enumerate}[label=\textup{(E12.\arabic*)},leftmargin=5.5em]
\item exact shell integration transports the physical selector tangent by
      conditional expectation and gives its covariance-corrected acceleration;
\item every marked selector response is identical in the direct terminal and
      staged coarse descriptions;
\item the weak-facing root derivative is the response to the actual selector
      tangent, not automatically the response to the coarse plaquette writer;
\item selector and marginal regularity are independent, and their two implicit
      charts have different, correctly transforming coefficients;
\item an exact tangent-incidence residual controls the difference between the
      full marginal correction and the correction inferred from a submitted
      generated writer;
\item on an exact decomposition modulo constants, two response rows determine
      the generated response and the marginal coefficient;
\item the V11 root and amplification statements admit the precise typing in
      \Cref{thm:YMA12-V11-qualification};
\item marked numerator intervals give a finite noncircular
      regular/obstructed/unresolved marginal-chart alternative;
\item the correct cofinal compact-coefficient hypothesis uses a uniform floor
      for \(\abs{\mathscr R(P)}\), not for the bare selector slope.
\end{enumerate}
No item evaluates the independent plaquette response on the selected Wilson
law, proves the tangent decomposition on that card, supplies the final-law
mixed field jet, proves weak-facing-root escape, or establishes local-vacuum
contraction or a continuum Yang--Mills mass gap.
\end{theorem}

\begin{proof}
Items~\textup{(E12.1)--(E12.2)} are
\Cref{thm:YMA12-shell-tangent,cor:YMA12-multishell-tangent,prop:YMA12-response-slope-transport}.
Items~\textup{(E12.3)--(E12.4)} are
\Cref{thm:YMA12-root-measures-selector,prop:YMA12-coordinate-covariance,cth:YMA12-selector-marginal-independent,thm:YMA12-two-charts}.
Items~\textup{(E12.5)--(E12.6)} are
\Cref{thm:YMA12-tangent-incidence,cor:YMA12-two-row-reduction}.
Item~\textup{(E12.7)} is
\Cref{thm:YMA12-V11-qualification}.
Items~\textup{(E12.8)--(E12.9)} are
\Cref{prop:YMA12-marked-response-interval,thm:YMA12-noncircular-manifest,prop:YMA12-cofinal-coefficient}.
\end{proof}

\begin{assessment}[V12 termination]
\label{ass:YMA12-termination}
The selector-tangent transport, three-row separation, coordinate covariance,
two-chart theorem, tangent-incidence identity, V11 qualification, and finite
manifest are \pass\ at their stated mathematical scopes.  This continuation
goes one genuine arrow upstream: it prevents the root derivative from being
used twice as both a trajectory slope and a plaquette response denominator.
The first unpopulated physical quantity is now the independent selected-root
response \(d_P=\mathscr R(P)\), followed by the action-tangent incidence.
The complete selected symbol therefore remains \stopstatus.  No negative
physical mode has been produced.
\end{assessment}

\section*{Version V12 append-only continuation record}
\addcontentsline{toc}{section}{Version V12 append-only continuation record}
The complete Version~V11 source above is preserved byte for byte up to its sole
terminal document-closing command.  Its SHA--256 digest is
\begin{center}\scriptsize\ttfamily
de921f3ab0972b03b62f7a729fc13f38b42358dea41aae4c47f6c4298462489e.
\end{center}
For Version~V13, preserve this accumulated source, append new mathematical
material immediately before the final document-closing command, and increment
the filename to \texttt{\_V13.tex}.  Keep the physical response target,
weak-facing selector, exact block map, parity-tree section, rooted Hodge
carrier, first axial mode, colour normalization, full unnormalized action,
actual selector tangent, plaquette marginal writer, generated insertion, and
local/topological convention fixed before reading any response, coefficient,
field-jet, current, ancestry, radial, or transport value.

% =============================================================================
% END APPENDED VERSION V12 MATERIAL
% =============================================================================


% =============================================================================
% START APPENDED VERSION V13 MATERIAL
% =============================================================================
\clearpage
\hypersetup{
 pdftitle={Renewal Yang--Mills Reduced-Fibre Wilson Shell Symbol and Local-Vacuum Programme V13},
 pdfsubject={The logarithm-free marginal cofactor, the Creutz response source, and a connected plaquette acquisition}
}
\pdfinfo{/Title (Renewal Yang--Mills Reduced-Fibre Wilson Shell Symbol and Local-Vacuum Programme V13)
 /Subject (The logarithm-free marginal cofactor, the Creutz response source, and a connected plaquette acquisition)}

\section{Version V13: the logarithm-free marginal cofactor and the connected plaquette row}
\label{sec:YMA13-entry}

Version~V12 is preserved above.  Its selector--marginal separation is
necessary, but its terminal acquisition still writes the marginal response as
four separately divided inserted loop numerators.  That representation is
exact, yet it is not the most upstream or best-conditioned finite card.  The
same logarithm-free function which isolates the weak-facing root can be
differentiated directly in the physically frozen plaquette direction.  At the
root, its derivative is a common positive factor times the marginal response.
The common factor cancels from the response-matching coefficient.

This continuation makes that reduction explicit.  It proves that the
remaining plaquette row is equivalently
\begin{equation}
 \boxed{
 d_P=\mathscr R(P)
 =-\operatorname{Cov}_{\mu_*}(\mathfrak s_{\rm Cr},P),}
 \label{eq:YMA13-upstream-target}
\end{equation}
where \(\mathfrak s_{\rm Cr}\) is one centered Creutz response source built
from the four already frozen loop writers.  Hence the selected marginal chart
is controlled by one same-history cross covariance, or by one unnormalized
root cofactor, rather than by four unrelated quotient intervals.  The
corresponding local decomposition is a signed sum of connected loop--plaquette
covariances.  This is the literal Wilson-specific quantity which a finite Haar
calculation or a source-connected low-temperature estimate must populate.

\begin{firewall}[No new selector, loop family, or source normalization]
\label{fire:YMA13-frozen-card}
Throughout V13 the target value \(u\), weak-facing root, four loop writers,
complete unnormalized action, coarse plaquette writer \(P\), generated writer
\(\Psi\), parity-tree section, rooted Hodge carrier, and selected axial
momentum--colour polarization are those already frozen in V10--V12.  The
cofactor below is a different evaluation of the same response derivative.  It
is not a new response selector and it does not identify the scalar native fan
of YM--B with the adjoint field score.
\end{firewall}

% =============================================================================
\section{The root cofactor is the marginal response without four divisions}
\label{sec:YMA13-cofactor}
% =============================================================================

Fix the selected finite card and write its base action as \(\Phi_*\).  For a
bounded real action insertion \(X\), set
\begin{equation}
 \Phi_t=\Phi_*+tX,
 \qquad
 N_\alpha(t)=\int \mathcal O_\alpha e^{-\Phi_t}\,\dd m_{\rm red},
 \label{eq:YMA13-X-family}
\end{equation}
with \(\alpha\in\{11,00,10,01\}\).  The loop writers are held fixed under this
derivative.  Put
\begin{equation}
 A(t)=N_{11}(t)N_{00}(t),
 \qquad
 B(t)=N_{10}(t)N_{01}(t),
 \label{eq:YMA13-ABt}
\end{equation}
and retain the V10 logarithm-free root function
\begin{equation}
 D_u(t)=A(t)-e^{-u}B(t).
 \label{eq:YMA13-Dut}
\end{equation}
At the selected root, \(D_u(0)=0\), all four \(N_\alpha(0)\) are positive,
and \(B_*:=B(0)>0\).

\begin{definition}[Root cofactor in an action direction]
\label{def:YMA13-root-cofactor}
The root cofactor of \(X\) is
\begin{equation}
 \boxed{
 \mathfrak c_u[X]
 :=\left.\frac{\dd}{\dd t}D_u(t)\right|_{t=0}.}
 \label{eq:YMA13-root-cofactor}
\end{equation}
It is evaluated at the already selected root and in the already frozen action
chart.
\end{definition}

For the inserted numerators of V11, abbreviate
\begin{align}
 \mathfrak J_{11,00}[X]
 &: =N_{11;X}N_{00}+N_{11}N_{00;X},
 \label{eq:YMA13-Jplus}\\
 \mathfrak J_{10,01}[X]
 &: =N_{10;X}N_{01}+N_{10}N_{01;X}.
 \label{eq:YMA13-Jminus}
\end{align}
All unlabelled quantities in this section are evaluated at \(t=0\).

\begin{theorem}[Logarithm-free cofactor identity]
\label{thm:YMA13-cofactor-identity}
For every bounded action direction \(X\),
\begin{equation}
 \boxed{
 \mathfrak c_u[X]
 =-\mathfrak J_{11,00}[X]
   +e^{-u}\mathfrak J_{10,01}[X].}
 \label{eq:YMA13-cofactor-pair-form}
\end{equation}
At the selected root,
\begin{equation}
 \boxed{
 \mathfrak c_u[X]
 =-e^{-u}B_*\,\mathscr R_{\Phi_*}(X).}
 \label{eq:YMA13-cofactor-response}
\end{equation}
Consequently
\begin{align}
 d_P&=-\frac{e^u}{B_*}\mathfrak c_u[P],
 \label{eq:YMA13-dP-cofactor}\\
 n_\Psi&=-\frac{e^u}{B_*}\mathfrak c_u[\Psi],
 \label{eq:YMA13-nPsi-cofactor}
\end{align}
and, whenever \(\mathfrak c_u[P]\ne0\),
\begin{equation}
 \boxed{
 b^{\rm marg}
 =-\frac{n_\Psi}{d_P}
 =-\frac{\mathfrak c_u[\Psi]}{\mathfrak c_u[P]}.}
 \label{eq:YMA13-b-cofactor}
\end{equation}
Thus the marginal coefficient requires no division by the four loop
numerators and no separate partition-function normalization.
\end{theorem}

\begin{proof}
Differentiation under the finite compact integral gives
\(N_\alpha'(0)=-N_{\alpha;X}\).  The product rule therefore gives
\[
 A'(0)=-\mathfrak J_{11,00}[X],
 \qquad
 B'(0)=-\mathfrak J_{10,01}[X],
\]
which proves \eqref{eq:YMA13-cofactor-pair-form}.

On the positive loop branch, V11 gives
\[
 D_u(t)=B(t)\bigl(e^{-\chi(t)}-e^{-u}\bigr).
\]
At the selected root the parenthesis vanishes.  Differentiation therefore
removes the derivative of \(B\) and gives
\[
 \mathfrak c_u[X]
 =-e^{-u}B_*\,\partial_t\chi(0)
 =-e^{-u}B_*\,\mathscr R_{\Phi_*}(X).
\]
Substitution of \(X=P\) and \(X=\Psi\) gives the two response rows.  Their
common nonzero factor cancels in \(-n_\Psi/d_P\), proving
\eqref{eq:YMA13-b-cofactor}.
\end{proof}

\begin{corollary}[Selector cofactor and the exact tangent-incidence ratio]
\label{cor:YMA13-selector-cofactor}
Let \(\lambda\) be the actual weak-facing selector parameter and
\(T_{\rm sel}=\partial_\lambda\Phi_\lambda|_{\lambda_*}\).  Then
\begin{equation}
 \boxed{
 \mathfrak c_u[T_{\rm sel}]
 =\partial_\lambda D_u(\lambda_*) .}
 \label{eq:YMA13-selector-Dprime}
\end{equation}
At a simple stable down-crossing this number is strictly positive.  If
\begin{equation}
 T_{\rm sel}=aP+\Psi+\Delta_{\rm tan},
 \label{eq:YMA13-tangent-decomposition}
\end{equation}
then
\begin{equation}
 \boxed{
 \mathfrak c_u[T_{\rm sel}]
 =a\mathfrak c_u[P]
  +\mathfrak c_u[\Psi]
  +\mathfrak c_u[\Delta_{\rm tan}].}
 \label{eq:YMA13-cofactor-incidence}
\end{equation}
On the exact response- and field-invisible residual branch,
\begin{equation}
 \boxed{
 b_{\rm full}^{\rm marg}
 =a-\frac{\mathfrak c_u[T_{\rm sel}]}
          {\mathfrak c_u[P]}.}
 \label{eq:YMA13-two-cofactor-b}
\end{equation}
\end{corollary}

\begin{proof}
The first identity is the definition of the action-direction derivative of
\(D_u\).  V11 proves that a simple stable down-crossing has
\(D_u'(\lambda_*)>0\).  Linearity of differentiation in the action direction
gives \eqref{eq:YMA13-cofactor-incidence}.  If the residual may be deleted
from both the response and selected field jets, then
\(\mathfrak c_u[\Delta_{\rm tan}]=0\), and solving the preceding identity for
\(-\mathfrak c_u[\Psi]/\mathfrak c_u[P]\) gives
\eqref{eq:YMA13-two-cofactor-b}.
\end{proof}

\begin{proposition}[Common vacuum and partition factors cancel at the root]
\label{prop:YMA13-common-factor}
Let \(q(t)>0\) be analytic and replace every loop numerator by
\(\widetilde N_\alpha(t)=q(t)N_\alpha(t)\).  Then
\begin{equation}
 \widetilde D_u(t)=q(t)^2D_u(t).
 \label{eq:YMA13-common-factor-D}
\end{equation}
At the selected root,
\begin{equation}
 \boxed{
 \widetilde{\mathfrak c}_u[X]
 =q(0)^2\mathfrak c_u[X].}
 \label{eq:YMA13-common-factor-cofactor}
\end{equation}
Hence the response \eqref{eq:YMA13-cofactor-response}, the regularity verdict
\(\mathfrak c_u[P]\ne0\), and the coefficient ratios
\eqref{eq:YMA13-b-cofactor}--\eqref{eq:YMA13-two-cofactor-b} are invariant
under every common unnormalized factor.  In particular, a vacuum-volume
factor is removed before source attribution rather than estimated as part of
the marginal response.
\end{proposition}

\begin{proof}
Both products \(A\) and \(B\) acquire the factor \(q^2\), proving
\eqref{eq:YMA13-common-factor-D}.  Differentiating gives
\[
 \widetilde D_u'
 =2qq'D_u+q^2D_u'.
\]
The first term vanishes at the root.  The remaining conclusions follow by
cancellation of the common positive factor.
\end{proof}

\begin{assessment}[The first simplification]
\label{ass:YMA13-first-simplification}
Version~V12 requested four \(P\)-inserted numerator ratios.  V13 replaces
that bank by the single signed cofactor
\(\mathfrak c_u[P]\), assembled before division.  The selected marginal
coefficient is the ratio of two cofactor derivatives of the same entire root
function.  This avoids four independent interval divisions and automatically
removes every common partition or vacuum factor at the selected root.
\end{assessment}

% =============================================================================
\section{The Creutz response source and its target-minimal Gram}
\label{sec:YMA13-response-source}
% =============================================================================

Let
\begin{equation}
 \dd\mu_*(z)=Z_*^{-1}e^{-\Phi_*(z)}\,\dd m_{\rm red}(z),
 \qquad
 W_\alpha=\mathbb E_{\mu_*}\mathcal O_\alpha>0.
 \label{eq:YMA13-base-law}
\end{equation}
Define the root-anchored Creutz response source
\begin{equation}
 \boxed{
 \mathfrak s_{\rm Cr}(z)
 :=\sum_\alpha
   \varepsilon_\alpha\frac{\mathcal O_\alpha(z)}{W_\alpha}.}
 \label{eq:YMA13-Creutz-source}
\end{equation}
The normalizations are evaluated on the same selected law and are not
retuned after inspecting the plaquette response.

\begin{theorem}[Every action response is one covariance with the Creutz source]
\label{thm:YMA13-Creutz-covariance}
The source \(\mathfrak s_{\rm Cr}\) is centered,
\begin{equation}
 \mathbb E_{\mu_*}\mathfrak s_{\rm Cr}=0,
 \label{eq:YMA13-Creutz-centered}
\end{equation}
and every bounded action insertion \(X\) satisfies
\begin{equation}
 \boxed{
 \mathscr R_{\Phi_*}(X)
 =-\operatorname{Cov}_{\mu_*}(\mathfrak s_{\rm Cr},X)
 =-\mathbb E_{\mu_*}[\mathfrak s_{\rm Cr}X].}
 \label{eq:YMA13-response-covariance}
\end{equation}
Equivalently,
\begin{equation}
 \boxed{
 \mathfrak c_u[X]
 =e^{-u}B_*\,
  \operatorname{Cov}_{\mu_*}(\mathfrak s_{\rm Cr},X).}
 \label{eq:YMA13-cofactor-covariance}
\end{equation}
In particular,
\begin{equation}
 \boxed{
 d_P=-\operatorname{Cov}(\mathfrak s_{\rm Cr},P),
 \qquad
 n_\Psi=-\operatorname{Cov}(\mathfrak s_{\rm Cr},\Psi).}
 \label{eq:YMA13-dP-n-covariance}
\end{equation}
\end{theorem}

\begin{proof}
Since \(\mathbb E\mathcal O_\alpha=W_\alpha\),
\[
 \mathbb E\mathfrak s_{\rm Cr}
 =\sum_\alpha\varepsilon_\alpha=0.
\]
Under the action deformation \(\Phi_*+tX\),
\[
 \left.\frac{\dd}{\dd t}W_\alpha(t)\right|_{0}
 =-\operatorname{Cov}_{\mu_*}(\mathcal O_\alpha,X).
\]
The common partition derivative is already included in this covariance.
Differentiating
\(\chi=\sum_\alpha\varepsilon_\alpha\log W_\alpha\) gives
\[
 \mathscr R_{\Phi_*}(X)
 =-\sum_\alpha\varepsilon_\alpha
   \frac{\operatorname{Cov}(\mathcal O_\alpha,X)}{W_\alpha}
 =-\operatorname{Cov}(\mathfrak s_{\rm Cr},X).
\]
Centeredness permits the last expectation form.  Combine this with
\eqref{eq:YMA13-cofactor-response} to obtain
\eqref{eq:YMA13-cofactor-covariance}; then substitute \(P\) and \(\Psi\).
\end{proof}

\begin{definition}[Response-source stocks and marginal angle]
\label{def:YMA13-response-stocks}
Put
\begin{equation}
 v_{\rm Cr}=\operatorname{Var}(\mathfrak s_{\rm Cr}),
 \qquad
 v_P=\operatorname{Var}(P),
 \qquad
 v_\Psi=\operatorname{Var}(\Psi),
 \label{eq:YMA13-stocks}
\end{equation}
and, when \(v_{\rm Cr}v_P>0\), define
\begin{equation}
 \boxed{
 \rho_P
 :=\frac{\operatorname{Cov}(\mathfrak s_{\rm Cr},P)}
          {\sqrt{v_{\rm Cr}v_P}}
 =-\frac{d_P}{\sqrt{v_{\rm Cr}v_P}}.}
 \label{eq:YMA13-marginal-angle}
\end{equation}
The number \(\rho_P\in[-1,1]\) is the normalized incidence angle between the
actual Creutz response source and the frozen plaquette writer.
\end{definition}

\begin{theorem}[Positive response Gram and typed marginal failure]
\label{thm:YMA13-response-Gram}
The target-minimal same-law response Gram is
\begin{equation}
 \boxed{
 \mathbb G_{\rm resp}
 =\begin{pmatrix}
 v_{\rm Cr}&-d_P&-n_\Psi\\
 -d_P&v_P&c_{P\Psi}\\
 -n_\Psi&c_{P\Psi}&v_\Psi
 \end{pmatrix}\succeq0,}
 \qquad
 c_{P\Psi}=\operatorname{Cov}(P,\Psi).
 \label{eq:YMA13-response-Gram}
\end{equation}
If \(v_{\rm Cr}>0\), the plaquette stock has the exact orthogonal split
\begin{equation}
 \boxed{
 v_P
 =\frac{d_P^2}{v_{\rm Cr}}
  +\mathbb I_{P\mid{\rm Cr}},
 \qquad
 \mathbb I_{P\mid{\rm Cr}}
 :=v_P-\frac{d_P^2}{v_{\rm Cr}}\ge0.}
 \label{eq:YMA13-P-short}
\end{equation}
Accordingly, the exact marginal alternatives are:
\begin{enumerate}[label=\textup{(A\arabic*)},leftmargin=3.8em]
\item \(v_{\rm Cr}=0\): the selected Creutz response source has collapsed and
      every first action response vanishes;
\item \(v_{\rm Cr}>0\), \(v_P=0\): the declared plaquette writer is invisible
      in the selected law and \(d_P=0\);
\item \(v_{\rm Cr}v_P>0\), \(\rho_P=0\): both stocks occur, but the plaquette
      direction is response orthogonal;
\item \(d_P\ne0\): the one-dimensional plaquette marginal chart is regular.
\end{enumerate}
The first three outcomes are mathematically distinct and are not replaced by
one undifferentiated statement that the denominator is small.
\end{theorem}

\begin{proof}
The displayed matrix is the covariance Gram of the three centered real
writers \(\mathfrak s_{\rm Cr},P,\Psi\), using
\eqref{eq:YMA13-dP-n-covariance}; hence it is positive semidefinite.  Shorting
the first scalar block gives
\(v_P-d_P^2/v_{\rm Cr}\ge0\), which is
\eqref{eq:YMA13-P-short}.  If \(v_{\rm Cr}=0\), the centered source vanishes
in \(L^2\), so every covariance with it vanishes.  If \(v_P=0\), the
centered plaquette writer vanishes and hence so does its covariance with the
response source.  On the positive-stock branch,
\eqref{eq:YMA13-marginal-angle} shows that \(d_P=0\) is precisely
orthogonality.  The last branch is the ordinary implicit-function condition.
\end{proof}

\begin{corollary}[What a simple weak-facing root actually guarantees]
\label{cor:YMA13-simple-root-source-floor}
Let
\(v_{\rm sel}=\operatorname{Var}(T_{\rm sel})\).  At a simple selected root,
\begin{equation}
 \boxed{
 d_{\rm sel}=-\operatorname{Cov}(\mathfrak s_{\rm Cr},T_{\rm sel})\ne0,
 \qquad
 v_{\rm Cr}v_{\rm sel}\ge d_{\rm sel}^2.}
 \label{eq:YMA13-simple-root-stock}
\end{equation}
Thus a simple root proves nonzero Creutz-source stock and nonzero selector
stock.  It proves neither \(v_P>0\) nor \(\rho_P\ne0\), and therefore still
does not prove marginal regularity.
\end{corollary}

\begin{proof}
The first identity is \Cref{thm:YMA13-Creutz-covariance}.  A simple root has
\(d_{\rm sel}\ne0\).  Cauchy--Schwarz gives the product lower bound and
forces both variances to be positive.  No term involving \(P\) appears in
this argument.
\end{proof}

\begin{corollary}[One two-replica acquisition of the entire response Gram]
\label{cor:YMA13-two-replica-Gram}
For independent \(z_1,z_2\sim\mu_*\), write
\(\Delta X=X(z_1)-X(z_2)\).  Then
\begin{equation}
 \boxed{
 \begin{aligned}
 v_{\rm Cr}&=\frac12\mathbb E(\Delta\mathfrak s_{\rm Cr})^2,\\
 v_P&=\frac12\mathbb E(\Delta P)^2,\\
 v_\Psi&=\frac12\mathbb E(\Delta\Psi)^2,\\
 d_P&=-\frac12\mathbb E[
       \Delta\mathfrak s_{\rm Cr}\,\Delta P],\\
 n_\Psi&=-\frac12\mathbb E[
       \Delta\mathfrak s_{\rm Cr}\,\Delta\Psi],\\
 c_{P\Psi}&=\frac12\mathbb E[\Delta P\,\Delta\Psi].
 \end{aligned}}
 \label{eq:YMA13-two-replica-Gram}
\end{equation}
Hence one common two-copy cylinder acquires the stocks, the marginal angle,
and the response numerator with their relative orientation intact.
\end{corollary}

\begin{proof}
For any two square-integrable real writers \(X,Y\), independence gives
\[
 \frac12\mathbb E[(X_1-X_2)(Y_1-Y_2)]
 =\mathbb E[XY]-\mathbb EX\,\mathbb EY
 =\operatorname{Cov}(X,Y).
\]
Apply this identity to the three writers and use
\eqref{eq:YMA13-dP-n-covariance}.
\end{proof}

\begin{firewall}[A response short is not a field short]
\label{fire:YMA13-response-not-field}
The nonnegative innovation \(\mathbb I_{P\mid{\rm Cr}}\) is the part of the
plaquette writer invisible to the single Creutz response source.  It is not
deleted from the physical action.  It may have nonzero coarse-field first and
second jets and may contribute to the final Wilson Hessian.  Response
matching removes one scalar response component; it does not short the action
or the field carrier by \(\mathfrak s_{\rm Cr}\).
\end{firewall}

% =============================================================================
\section{One unnormalized two-pair writer and exact source-connected cancellation}
\label{sec:YMA13-two-pair}
% =============================================================================

The normalized response Gram is well conditioned when the loop expectations
are already certified.  The exact Haar route can avoid those divisions
altogether.  For \(\alpha,\beta\in\{11,00,10,01\}\), define the two-copy
marked cylinder
\begin{equation}
 \boxed{
 \begin{aligned}
 \mathfrak J_{\alpha,\beta}[X]
 :={}&\int_{Y^2}
 \mathcal O_\alpha(z_1)\mathcal O_\beta(z_2)\\
 &\qquad\times\bigl(X(z_1)+X(z_2)\bigr)
 e^{-\Phi_*(z_1)-\Phi_*(z_2)}
 \,\dd m(z_1)\dd m(z_2).
 \end{aligned}}
 \label{eq:YMA13-two-pair-writer}
\end{equation}
This notation agrees with
\eqref{eq:YMA13-Jplus}--\eqref{eq:YMA13-Jminus}.

\begin{proposition}[Exact two-pair acquisition]
\label{prop:YMA13-two-pair-acquisition}
For every bounded insertion \(X\),
\begin{equation}
 \boxed{
 \mathfrak c_u[X]
 =-\mathfrak J_{11,00}[X]
  +e^{-u}\mathfrak J_{10,01}[X].}
 \label{eq:YMA13-two-pair-cofactor}
\end{equation}
Thus \(d_P\) requires two pair-marked unnormalized integrals, assembled with
their actual signs before division.  If the loop writers are not pointwise
nonnegative, the cylinder in \eqref{eq:YMA13-two-pair-writer} is a signed
writer on a positive base measure; no positive-replica interpretation is
asserted.  Exact Haar contraction and outward signed intervals remain valid.
\end{proposition}

\begin{proof}
Fubini gives
\[
 \mathfrak J_{\alpha,\beta}[X]
 =N_{\alpha;X}N_\beta+N_\alpha N_{\beta;X}.
\]
Substitution into \eqref{eq:YMA13-cofactor-pair-form} proves the identity.
The qualification about signs is immediate from the displayed integrand.
\end{proof}

\begin{theorem}[Normalized connected form of the cofactor]
\label{thm:YMA13-connected-cofactor}
Let \(W_\alpha=N_\alpha/Z_*\).  Define
\begin{equation}
 \mathscr D_u^{\rm norm}
 :=W_{11}W_{00}-e^{-u}W_{10}W_{01}.
 \label{eq:YMA13-normalized-D}
\end{equation}
Then
\begin{equation}
 D_u=Z_*^2\mathscr D_u^{\rm norm}.
 \label{eq:YMA13-D-Z2}
\end{equation}
At the selected root,
\begin{equation}
 \boxed{
 \mathfrak c_u[X]
 =Z_*^2\left.\frac{\dd}{\dd t}
     \mathscr D_u^{\rm norm}[\Phi_*+tX]
   \right|_{0}.}
 \label{eq:YMA13-root-connected-D}
\end{equation}
Moreover,
\begin{equation}
 \boxed{
 \begin{aligned}
 Z_*^{-2}\mathfrak c_u[X]
 ={}&-W_{00}\operatorname{Cov}(\mathcal O_{11},X)
     -W_{11}\operatorname{Cov}(\mathcal O_{00},X)\\
 &+e^{-u}W_{01}\operatorname{Cov}(\mathcal O_{10},X)
  +e^{-u}W_{10}\operatorname{Cov}(\mathcal O_{01},X).
 \end{aligned}}
 \label{eq:YMA13-connected-cofactor-expanded}
\end{equation}
Every unmarked partition derivative cancels.  The root cofactor is therefore
a source-connected response quantity even though its exact evaluator uses
unnormalized numerators.
\end{theorem}

\begin{proof}
The first identity follows from \(N_\alpha=Z_*W_\alpha\).  Differentiate it
along \(\Phi_*+tX\):
\[
 \mathfrak c_u[X]
 =2Z_*Z_*'\mathscr D_u^{\rm norm}
  +Z_*^2(\mathscr D_u^{\rm norm})'.
\]
The first term vanishes at the root.  Finally
\(W_\alpha'=-\operatorname{Cov}(\mathcal O_\alpha,X)\); applying the product
rule to \eqref{eq:YMA13-normalized-D} gives
\eqref{eq:YMA13-connected-cofactor-expanded}.
\end{proof}

\begin{assessment}[No double payment of vacuum bubbles]
\label{ass:YMA13-vacuum-cancellation}
The unnormalized cofactor and the connected covariance are two exact
representations of one quantity.  The Haar evaluator may use the former; the
analytic head--tail estimate may use the latter.  Their equality means that a
vacuum factor, partition derivative, or source-disjoint bubble is not charged
once in the unnormalized numerator and again in a connected tail.
\end{assessment}

% =============================================================================
\section{The exact connected loop--plaquette response kernel}
\label{sec:YMA13-plaquette-kernel}
% =============================================================================

Write the frozen marginal writer as a sum of local real plaquette terms,
\begin{equation}
 P=\sum_{p\in\mathscr P_c}p_p,
 \label{eq:YMA13-P-local-sum}
\end{equation}
with the coefficient convention already fixed by V12.  Define
\begin{equation}
 \boxed{
 \kappa_P(p)
 :=-\operatorname{Cov}_{\mu_*}(\mathfrak s_{\rm Cr},p_p).}
 \label{eq:YMA13-kappa-p}
\end{equation}

\begin{theorem}[Exact connected plaquette response kernel]
\label{thm:YMA13-connected-P-kernel}
On every finite selected card,
\begin{equation}
 \boxed{
 d_P=\sum_{p\in\mathscr P_c}\kappa_P(p).}
 \label{eq:YMA13-dP-kernel-sum}
\end{equation}
More explicitly,
\begin{equation}
 \boxed{
 \kappa_P(p)
 =-\sum_\alpha
  \frac{\varepsilon_\alpha}{W_\alpha}
  \operatorname{Cov}_{\mu_*}(\mathcal O_\alpha,p_p).}
 \label{eq:YMA13-kernel-four-loops}
\end{equation}
Thus the independent marginal denominator is not an arbitrary extensive
expectation.  It is the signed total mass of one connected
loop--plaquette correlation kernel on the selected law.
\end{theorem}

\begin{proof}
Linearity of covariance and the finite sum
\eqref{eq:YMA13-P-local-sum} give
\[
 d_P=-\operatorname{Cov}(\mathfrak s_{\rm Cr},P)
     =-\sum_p\operatorname{Cov}(\mathfrak s_{\rm Cr},p_p).
\]
This is \eqref{eq:YMA13-dP-kernel-sum}.  Expanding
\(\mathfrak s_{\rm Cr}\) by \eqref{eq:YMA13-Creutz-source} gives the second
formula.
\end{proof}

Let \(\mathscr S_\alpha\) be the set of coarse cells meeting the support of
\(\mathcal O_\alpha\), and let \(s_\alpha=|\mathscr S_\alpha|\).  Distance
between a coarse plaquette and a support cell is the periodic \(\ell^1\)
cell distance.

\begin{proposition}[Exact Markov-blanket factorization of a distant kernel entry]
\label{prop:YMA13-Markov-blanket}
Fix \(\alpha\) and a plaquette \(p\).  Suppose a physical boundary record
\(\mathscr B\) separates the loop support from \(p\) for the finite-range
selected law, so that \(\mathcal O_\alpha\) and \(p_p\) are conditionally
independent given \(\mathscr B\).  Then
\begin{equation}
 \boxed{
 \operatorname{Cov}(\mathcal O_\alpha,p_p)
 =\operatorname{Cov}\bigl(
    \mathbb E[\mathcal O_\alpha\mid\mathscr B],
    \mathbb E[p_p\mid\mathscr B]
   \bigr).}
 \label{eq:YMA13-Markov-factor}
\end{equation}
Hence every distant term in \eqref{eq:YMA13-kernel-four-loops} factors
through the actual Creutz boundary message and the actual plaquette score
message.  A scalar native fan or an unrelated radial coefficient cannot
replace either message.
\end{proposition}

\begin{proof}
Conditional independence gives
\[
 \mathbb E[\mathcal O_\alpha p_p\mid\mathscr B]
 =\mathbb E[\mathcal O_\alpha\mid\mathscr B]
  \mathbb E[p_p\mid\mathscr B].
\]
Take expectations and subtract
\(\mathbb E\mathcal O_\alpha\,\mathbb Ep_p\).  The result is precisely the
covariance on the right of \eqref{eq:YMA13-Markov-factor}.
\end{proof}

\begin{corollary}[Finite connected head plus an explicit exponential tail]
\label{cor:YMA13-connected-tail}
Assume that, on the selected family, there are \(\mu>0\) and constants
\(C_\alpha<\infty\) such that
\begin{equation}
 \abs{\operatorname{Cov}(\mathcal O_\alpha,p_p)}
 \le C_\alpha
 e^{-\mu\operatorname{dist}(p,\mathscr S_\alpha)}
 \label{eq:YMA13-connected-decay-assumption}
\end{equation}
for every coarse plaquette \(p\).  Define the radius-\(R\) head
\begin{equation}
 d_P^{\le R}
 :=\sum_{p:\,\min_\alpha\operatorname{dist}(p,\mathscr S_\alpha)\le R}
   \kappa_P(p).
 \label{eq:YMA13-dP-head}
\end{equation}
Then
\begin{equation}
 \boxed{
 \abs{d_P-d_P^{\le R}}
 \le96\sum_\alpha\frac{C_\alpha s_\alpha}{W_\alpha}
       \sum_{n>R}(n+1)^3e^{-\mu n}.}
 \label{eq:YMA13-dP-tail}
\end{equation}
The bound is uniform in the torus size whenever the displayed constants and
the loop-visibility floors are uniform.  It does not assert the decay
hypothesis; it identifies exactly the connected estimate which closes the
marginal row.
\end{corollary}

\begin{proof}
Equation \eqref{eq:YMA13-kernel-four-loops} and the assumed bound give
\[
 \abs{\kappa_P(p)}
 \le\sum_\alpha\frac{C_\alpha}{W_\alpha}
 e^{-\mu\operatorname{dist}(p,\mathscr S_\alpha)}.
\]
For one support cell, the number of coarse lattice sites at periodic
\(\ell^1\) distance \(n\) is at most \(16(n+1)^3\).  There are six plaquette
orientations and at most \(s_\alpha\) support cells.  A union bound therefore
gives at most \(96s_\alpha(n+1)^3\) plaquettes at distance \(n\) from
\(\mathscr S_\alpha\).  Sum the preceding pointwise estimate over
\(n>R\).
\end{proof}

\begin{firewall}[What YM--B may export to this row]
\label{fire:YMA13-YMB-interface}
The Markov-blanket identity permits YM--B to export an actual boundary-message
estimate for the same four growing physical loop writers.  It does not permit
the extinct fixed native scalar fan, an old/fine overlap coefficient, or an
independent radial four-leaf attenuation to be inserted as \(d_P\).  The
head \(d_P^{\le R}\), the boundary tail, and the loop visibility denominators
must belong to the same selected root and the same response-matched law.
\end{firewall}

% =============================================================================
\section{Exact Haar execution of the two-pair cofactor}
\label{sec:YMA13-Haar}
% =============================================================================

The cofactor does not require a new group-integration calculus.  On the pure
Wilson branch write the nonnegative unit action as \(V\), so that the base
weight is \(e^{-\beta V}\).  For one loop writer define
\begin{align}
 m_{\alpha,n}
 &: =\int\mathcal O_\alpha V^n\,\dd m,
 \label{eq:YMA13-m-alpha-n}\\
 m_{\alpha;X,n}
 &: =\int\mathcal O_\alpha X V^n\,\dd m.
 \label{eq:YMA13-m-alpha-X-n}
\end{align}
Every coefficient is a finite Wilson word contraction and is evaluated by
the invariant-tensor projectors already proved in V10.

\begin{theorem}[Exact convolution formula for the pair-marked moments]
\label{thm:YMA13-pair-convolution}
For indices
\[
 \alpha,\beta\in\{11,00,10,01\},
\]
set
\begin{equation}
 \boxed{
 j_{\alpha,\beta;X,n}
 :=\sum_{r=0}^{n}\binom nr
 \left(
  m_{\alpha;X,r}m_{\beta,n-r}
  +m_{\alpha,r}m_{\beta;X,n-r}
 \right).}
 \label{eq:YMA13-pair-coefficients}
\end{equation}
Then
\begin{equation}
 \boxed{
 \mathfrak J_{\alpha,\beta}[X](\beta)
 =\sum_{n=0}^{\infty}
  \frac{(-\beta)^n}{n!}
  j_{\alpha,\beta;X,n}.}
 \label{eq:YMA13-pair-series}
\end{equation}
Consequently the root cofactor has coefficients
\begin{equation}
 \boxed{
 g_{X,n}
 =-j_{11,00;X,n}
  +e^{-u}j_{10,01;X,n},
 \qquad
 \mathfrak c_u[X](\beta)
 =\sum_{n=0}^{\infty}
   \frac{(-\beta)^n}{n!}g_{X,n}.}
 \label{eq:YMA13-cofactor-series}
\end{equation}
Two replicas therefore introduce no new \(\SU(3)\) Haar projector: their
coefficients are binomial convolutions of the one-copy coefficients already
required by V10--V12.
\end{theorem}

\begin{proof}
Use \eqref{eq:YMA13-two-pair-writer} with the two-copy action
\(V(z_1)+V(z_2)\).  Expanding its exponential gives
\[
 \mathfrak J_{\alpha,\beta}[X]
 =\sum_{n\ge0}\frac{(-\beta)^n}{n!}
 \int\!\!\int
 \mathcal O_\alpha(z_1)\mathcal O_\beta(z_2)
 (X_1+X_2)(V_1+V_2)^n\,\dd m_1\dd m_2.
\]
The binomial theorem and Fubini split the coefficient into the two terms in
\eqref{eq:YMA13-pair-coefficients}.  The cofactor formula follows by the
linear combination in \eqref{eq:YMA13-two-pair-cofactor}.
\end{proof}

\begin{proposition}[Explicit finite-degree remainder for the cofactor]
\label{prop:YMA13-cofactor-tail}
Assume
\begin{equation}
 0\le V\le V_{\max},
 \qquad
 \abs{\mathcal O_\alpha}\le1,
 \qquad
 \abs X\le L_X.
 \label{eq:YMA13-Haar-bounds}
\end{equation}
For \(0\le\beta\le\beta_+\), let
\(\mathfrak J_{\alpha,\beta}^{(K)}[X]\) and
\(\mathfrak c_u^{(K)}[X]\) denote the degree-\(K\) truncations of
\eqref{eq:YMA13-pair-series} and
\eqref{eq:YMA13-cofactor-series}.  Then
\begin{align}
 \boxed{
 \abs{
 \mathfrak J_{\alpha,\beta}[X]
 -\mathfrak J_{\alpha,\beta}^{(K)}[X]}
 \le
 2L_X\frac{(2\beta_+V_{\max})^{K+1}}{(K+1)!},}
 \label{eq:YMA13-pair-tail}\\
 \boxed{
 \abs{
 \mathfrak c_u[X]-\mathfrak c_u^{(K)}[X]}
 \le
 2L_X(1+e^{-u})
 \frac{(2\beta_+V_{\max})^{K+1}}{(K+1)!}.}
 \label{eq:YMA13-cofactor-tail-bound}
\end{align}
For a pure Wilson action with \(N_{\rm plaq}\) plaquettes,
\(V_{\max}=3N_{\rm plaq}/2\).  The estimate is finite-card rigorous but may
be computationally poor at large volume; the connected kernel of
\Cref{thm:YMA13-connected-P-kernel} is the correct route for a uniform
source-local estimate.
\end{proposition}

\begin{proof}
For \(s\ge0\), Taylor's theorem gives
\[
 \left|
 e^{-s}-\sum_{n=0}^{K}\frac{(-s)^n}{n!}
 \right|
 \le\frac{s^{K+1}}{(K+1)!}.
\]
In the two-copy integral, \(s=\beta(V_1+V_2)\le2\beta_+V_{\max}\), while the
absolute value of the marked writer is at most \(2L_X\).  This proves the
first bound.  Add the two pair errors with coefficients \(1\) and \(e^{-u}\)
to obtain the second.  The Wilson value of \(V_{\max}\) is the pointwise
plaquette bound already used in V10.
\end{proof}

\begin{assessment}[Finite versus cofinal execution]
\label{ass:YMA13-finite-cofinal}
At fixed card the V10 invariant-tensor engine, the convolution
\eqref{eq:YMA13-pair-coefficients}, and the factorial tail
\eqref{eq:YMA13-cofactor-tail-bound} determine
\(\mathfrak c_u[P]\) to arbitrary accuracy.  A cofinal theorem should not
use the resulting volume-scale degree.  It should instead prove or acquire
the source-connected head and tail of
\eqref{eq:YMA13-dP-kernel-sum}.  These are two evaluations of one cofactor,
not two different marginal rows.
\end{assessment}

% =============================================================================
\section{Outward cofactor intervals, incidence control, and marginal ramification}
\label{sec:YMA13-intervals}
% =============================================================================

Let \(I_P=[c_P^-,c_P^+]\),
\(I_\Psi=[c_\Psi^-,c_\Psi^+]\), and
\(I_{\rm sel}=[c_{\rm sel}^-,c_{\rm sel}^+]\) be outward intervals for the
three root cofactors on one certified root box.  Let
\(0<B_-\le B_*\le B_+\).  Reuse the sign-safe quotient hull
\(\mathcal Q\) from \eqref{eq:YMA12-quotient-hull}.

\begin{theorem}[Sign-safe cofactor certificate]
\label{thm:YMA13-cofactor-certificate}
The independent marginal response satisfies
\begin{equation}
 \boxed{
 d_P\in-e^u\mathcal Q(I_P,[B_-,B_+]).}
 \label{eq:YMA13-dP-interval}
\end{equation}
If \(0\notin I_P\), then
\begin{align}
 \boxed{
 b^{\rm marg}
 \in-\mathcal Q(I_\Psi,I_P),}
 \label{eq:YMA13-b-interval}\\
 \boxed{
 b_{\rm full}^{\rm marg}
 \in a-\mathcal Q(I_{\rm sel},I_P)}
 \label{eq:YMA13-bfull-interval}
\end{align}
for the submitted and exact-tangent coefficients, respectively, on their
stated incidence branches.  These coefficient intervals require neither a
loop-numerator denominator nor a partition-function interval.

After input checks, the marginal response card terminates as follows.
\begin{enumerate}[label=\textup{(C\arabic*)},leftmargin=3.8em]
\item If \(0\notin I_P\), the chart is regular and the coefficient interval
      is passed to the V11 final-law mixed Hessian.
\item If exact arithmetic proves both
      \(\mathfrak c_u[P]=0\) and \(\mathfrak c_u[\Psi]\ne0\), then no finite
      first-order plaquette coefficient preserves the response.  This is the
      plaquette response-chart obstruction.
\item If both first cofactors vanish, the chart is first-order nonunique and
      the first nonzero higher cofactor in the \(P\)-direction is required.
\item If the outward interval meets zero without proving equality, the verdict
      is \(\stopstatus\), not obstruction.
\end{enumerate}
\end{theorem}

\begin{proof}
Equation \eqref{eq:YMA13-dP-cofactor} and positivity of \(B_*\) give the first
interval by monotonicity on a positive-denominator rectangle.  Equations
\eqref{eq:YMA13-b-cofactor} and
\eqref{eq:YMA13-two-cofactor-b} give the next two intervals, with the signs of
the quotient coefficients retained by the corner hull.  The four outcomes
are the ordinary implicit-function alternatives written in cofactor
coordinates; failure of an interval to separate zero proves no equality.
\end{proof}

\begin{proposition}[Same-law tangent-incidence bound]
\label{prop:YMA13-incidence-bound}
Let
\begin{equation}
 \Delta_{\rm tan}=T_{\rm sel}-aP-\Psi,
 \qquad
 v_\Delta=\operatorname{Var}(\Delta_{\rm tan}).
 \label{eq:YMA13-Delta-var}
\end{equation}
Then
\begin{equation}
 \boxed{
 \abs{d_{\rm sel}-ad_P-n_\Psi}^2
 \le v_{\rm Cr}v_\Delta.}
 \label{eq:YMA13-incidence-CS}
\end{equation}
Equivalently, if \(a\ne0\),
\begin{equation}
 \boxed{
 \left|
 d_P-\frac{d_{\rm sel}-n_\Psi}{a}
 \right|
 \le\frac{\sqrt{v_{\rm Cr}v_\Delta}}{\abs a}.}
 \label{eq:YMA13-dP-transfer}
\end{equation}
Thus a same-law positive variance interval for the literal tangent residual
can certify a sign and floor for \(d_P\) without treating scalar response
agreement as exact action equality.  The bound concerns only the response
row; nonzero residual field jets remain in the final Hessian.
\end{proposition}

\begin{proof}
By \Cref{thm:YMA13-Creutz-covariance},
\[
 d_{\rm sel}-ad_P-n_\Psi
 =-\operatorname{Cov}(\mathfrak s_{\rm Cr},\Delta_{\rm tan}).
\]
Cauchy--Schwarz gives \eqref{eq:YMA13-incidence-CS}; division by
\(\abs a\) gives \eqref{eq:YMA13-dP-transfer}.
\end{proof}

\begin{proposition}[Higher root cofactors decide marginal ramification]
\label{prop:YMA13-higher-cofactor}
Let \(s\mapsto\Phi_*+sP\), and write \(D_u(s)\) and \(\chi(s)\) for the
corresponding root and response functions.  Suppose
\begin{equation}
 \chi^{(1)}(0)=\cdots=\chi^{(m-1)}(0)=0
 \label{eq:YMA13-marginal-vanishing}
\end{equation}
for some \(m\ge2\).  Then
\begin{equation}
 \boxed{
 D_u^{(m)}(0)
 =-e^{-u}B_*\chi^{(m)}(0).}
 \label{eq:YMA13-higher-cofactor-relation}
\end{equation}
In particular, the first nonzero derivative of \(D_u\) in the \(P\)-direction
has the same order as the first nonzero marginal response derivative and the
opposite sign.  If a generated parameter \(r\) has
\(\partial_rD_u(0,0)\ne0\), the local response-level equation is a genuine
H\"older/Puiseux branch of order \(1/m\), rather than an ordinary linear
marginal coefficient.
\end{proposition}

\begin{proof}
Use
\(D_u(s)=B(s)(e^{-\chi(s)}-e^{-u})\).  At the root the second factor
vanishes.  Under \eqref{eq:YMA13-marginal-vanishing}, its first
\(m-1\) derivatives vanish as well, while its \(m\)-th derivative is
\(-e^{-u}\chi^{(m)}(0)\).  Leibniz' rule therefore leaves only the zeroth
derivative of \(B\), proving the identity.  If a second parameter has a
nonzero first derivative, the leading local equation has the form
\(a s^m+b r+o(|s|^m+|r|)=0\); solving its leading balance gives
\(|s|\asymp|r|^{1/m}\) on every real branch allowed by the sign and parity.
\end{proof}

\begin{firewall}[Cofactor zero versus physical symbol zero]
\label{fire:YMA13-cofactor-not-symbol}
A zero marginal cofactor means only that the frozen plaquette writer is
first-order orthogonal to the selected Creutz response source.  It is not a
zero mode of the reduced-fibre Wilson Hessian.  Conversely, a nonzero
cofactor regularizes the response chart but does not prove a positive field
margin.  The complete same-final-law deterministic and score jets of V11
remain mandatory after the coefficient is acquired.
\end{firewall}

% =============================================================================
\section{The revised literal acquisition and termination}
\label{sec:YMA13-frontier}
% =============================================================================

\begin{theorem}[Target-minimal V13 acquisition]
\label{thm:YMA13-target-minimal}
On the first selected root, the independent plaquette response and the
marginal coefficient are completely determined by either of the following
equivalent finite packets.
\begin{enumerate}[label=\textup{(P\arabic*)},leftmargin=3.8em]
\item \emph{Unnormalized cofactor packet:}
      \(B_*\), \(\mathfrak c_u[P]\), and
      \(\mathfrak c_u[\Psi]\), or
      \(\mathfrak c_u[T_{\rm sel}]\) on the exact tangent-incidence branch.
\item \emph{Normalized response-Gram packet:}
      the loop means \(W_\alpha\) and the covariance Gram
      \eqref{eq:YMA13-response-Gram}, acquired on one selected law.
\item \emph{Connected spatial packet:}
      a finite head of \(\kappa_P\) together with an outward bound for its
      omitted connected tail and the corresponding packet for \(\Psi\) or
      the exact selector tangent.
\end{enumerate}
The three packets give the same \(d_P\) and coefficient.  They may be used as
finite exact, same-law replica, and cofinal analytic implementations,
respectively; they are not three source prices.
\end{theorem}

\begin{proof}
Packet~\textup{(P1)} determines the quantities by
\Cref{thm:YMA13-cofactor-identity,cor:YMA13-selector-cofactor}.
Packet~\textup{(P2)} determines them by
\Cref{thm:YMA13-Creutz-covariance,thm:YMA13-response-Gram}.
Packet~\textup{(P3)} determines \(d_P\) by
\Cref{thm:YMA13-connected-P-kernel}; the identical construction applies to
any local decomposition of \(\Psi\) or \(T_{\rm sel}\).  Equality of the
three values is \eqref{eq:YMA13-cofactor-covariance} and
\eqref{eq:YMA13-dP-kernel-sum}.
\end{proof}

\begin{assessment}[What V13 resolves]
\label{ass:YMA13-progress}
The first V12 row is no longer an unspecified collection of four quotient
queries.  It is one root cofactor, one response-source cross Gram, or one
connected loop--plaquette kernel.  The exact Haar route is reduced to
one-copy invariant-tensor coefficients and their binomial convolution.  The
cofinal route is reduced to a finite connected head and a boundary-factored
tail.  A simple selector root is proved to supply Creutz-source visibility but
not the marginal angle.  If the marginal angle vanishes, V13 distinguishes
source collapse, plaquette invisibility, and genuine response orthogonality.
\end{assessment}

\begin{openproblem}[V14: execute the cofactor before opening another compiler]
\label{op:YMA13-V14}
Preserve V13 and every frozen physical choice.  On the first actual selected
root, proceed in this order.
\begin{enumerate}[label=\textup{(V14.\arabic*)},leftmargin=5.7em]
\item Evaluate \(\mathfrak c_u[P]\) directly from the two pair-marked
      unnormalized cylinders in \eqref{eq:YMA13-two-pair-cofactor}, using the
      V10 Haar projectors and the convolution
      \eqref{eq:YMA13-pair-coefficients}; do not first form four response
      ratios.
\item Independently assemble the normalized two-replica Gram
      \eqref{eq:YMA13-response-Gram} when a same-law numerical cylinder is
      available.  Verify
      \eqref{eq:YMA13-cofactor-covariance} as the held-out normalization
      check.
\item If a cofinal statement is intended, split
      \(d_P=\sum_p\kappa_P(p)\) at a predeclared physical collar and acquire
      the finite head plus the actual Creutz boundary-message tail.  Do not
      substitute the extinct native fan or the radial four-leaf coefficient.
\item Evaluate \(\mathfrak c_u[\Psi]\), or use
      \(\mathfrak c_u[T_{\rm sel}]\) only after the tangent-incidence residual
      is bounded by \eqref{eq:YMA13-incidence-CS}.  Return the sign-safe
      coefficient interval of \Cref{thm:YMA13-cofactor-certificate}.
\item On a regular branch, insert that coefficient into the complete
      same-final-law mixed Hessian of V11 and return
      \(\pass/\obstruction/\stopstatus\) on the predeclared axial block.
\item If the first cofactor vanishes exactly, acquire the first nonzero higher
      \(P\)-cofactor before introducing any alternative marginal chart.
\end{enumerate}
If no executable unnormalized pair cylinder, common-law response Gram, or
connected head--tail packet is supplied, the correct verdict is the named
plaquette-response \(\stopstatus\).  No further root, ratio, current,
boundary, or ancestry theorem is an advance before this cofactor is read.
\end{openproblem}

% =============================================================================
\section{Version V13 proof ledger and append-only boundary}
\label{sec:YMA13-ledger}
% =============================================================================

The accompanying replay
\begin{center}
\path{YMA_V13_root_cofactor_replay.py}
\end{center}
checks the root-factor derivative, the common-factor cancellation, the
Creutz-source covariance identity, the two-replica Gram, the cofactor ratio,
the pair-moment binomial convolution, a finite Markov-blanket factorization,
and the higher-cofactor sign.  Its script and exact-output digests are
\begin{center}\scriptsize\ttfamily
 60a004c867e1c52a0195a5d389c9b944aef99c2996dfda5e56cde74c70bbf9e8,\\[-0.15em]
 92a3b0bc5f1b675dd87b6ceb4e80349e4a9b3eb658a29c35ca4cb4a8ddf17624.
\end{center}
It uses exact rational arithmetic and symbolic polynomials.  It does not
evaluate a Wilson path integral, locate a selected root, or manufacture a
marginal response.

\begingroup\small
\begin{longtable}{>{\raggedright\arraybackslash}p{0.27\textwidth}
                  >{\raggedright\arraybackslash}p{0.20\textwidth}
                  >{\raggedright\arraybackslash}p{0.45\textwidth}}
\caption{Version V13 proof-status ledger.}\label{tab:YMA13-ledger}\\
\toprule
\textbf{Row}&\textbf{Status}&\textbf{Exact scope}\\
\midrule
\endfirsthead
\toprule
\textbf{Row}&\textbf{Status}&\textbf{Exact scope}\\
\midrule
\endhead
Root cofactor identity&Proved exactly
&Every action response is one derivative of the V10 logarithm-free root
function; common partition factors cancel at the root.\\[0.3em]
Marginal coefficient&Reduced exactly
&\(b^{\rm marg}=-\mathfrak c_u[\Psi]/\mathfrak c_u[P]\), with no four-loop
or partition denominator.\\[0.3em]
Creutz response source&Constructed exactly
&One centered writer \(\mathfrak s_{\rm Cr}\) gives every action response as
minus a same-law covariance.\\[0.3em]
Response Gram&Positive finite card
&One three-writer covariance Gram separates response-source collapse,
plaquette invisibility, response orthogonality, and regular incidence.\\[0.3em]
Simple-root implication&Sharpened
&A simple selector root forces response-source and selector stock, but not a
nonzero plaquette angle.\\[0.3em]
Two-pair unnormalized writer&Proved exactly
&Two pair-marked cylinders assemble the cofactor before division; no
pointwise positivity of real Wilson loops is assumed.\\[0.3em]
Connected plaquette kernel&Proved exactly
&\(d_P\) is the total mass of the connected loop--plaquette kernel
\(\kappa_P\).\\[0.3em]
Boundary factorization&Proved exactly
&Distant kernel entries factor through the actual Creutz and plaquette
boundary messages under the finite-range Markov blanket.\\[0.3em]
Connected tail&Conditional explicit bound
&Exponential loop--plaquette covariance decay gives the displayed finite
head and cutoff-uniform tail.  The decay estimate itself remains physical.\\[0.3em]
Haar pair execution&Reduced exactly
&Two-copy coefficients are binomial convolutions of the one-copy V10 Haar
coefficients, with an explicit factorial remainder.\\[0.3em]
Tangent-incidence error&Positive bound
&The scalar incidence defect is bounded by the Creutz-source stock times the
literal tangent-residual variance.\\[0.3em]
Marginal ramification&Proved in cofactor coordinates
&The first nonzero higher \(P\)-cofactor has the same order and opposite sign
as the first nonzero marginal response derivative.\\[0.3em]
Selected physical cofactor&Open value
&No attached cylinder supplies the pair-marked Wilson integrals or the common
response Gram on the first selected root.\\[0.3em]
Selected response-matched symbol&\stopstatus
&The complete V11 mixed Hessian still awaits the marginal coefficient and its
same-final-law jets.\\[0.3em]
Local-vacuum contraction and mass&Not inferred
&No finite response cofactor is promoted to an OS-time or continuum mass
gap.\\
\bottomrule
\end{longtable}
\endgroup

\begin{theorem}[Version V13 root-cofactor reduction]
\label{thm:YMA13-terminal}
Versions~V1--V12 are retained.  Version~V13 additionally proves:
\begin{enumerate}[label=\textup{(E13.\arabic*)},leftmargin=5.5em]
\item the independent plaquette response is one directional derivative of the
      already existing logarithm-free root function;
\item its response-matching coefficient is a ratio of two root cofactors and
      is invariant under every common vacuum or partition factor;
\item one centered Creutz response source represents every action response by
      a same-law covariance and yields a positive target-minimal Gram;
\item the marginal zero branch splits into source collapse, plaquette
      invisibility, or nonzero-stock response orthogonality;
\item a simple weak-facing root proves response-source visibility but does not
      prove marginal regularity;
\item two pair-marked unnormalized cylinders acquire the cofactor before
      division, while the exact Haar coefficients reduce to one-copy moments;
\item the plaquette denominator is the total mass of a connected
      loop--plaquette kernel whose distant terms factor through the actual
      Markov boundary messages;
\item a source-connected exponential covariance estimate yields a finite
      head/tail certificate for \(d_P\);
\item outward cofactor intervals give the exact regular/obstructed/unresolved
      marginal-chart alternative and a sign-safe coefficient interval;
\item higher root cofactors decide the marginally ramified branch without
      changing the physical marginal writer.
\end{enumerate}
No item evaluates the first selected-root cofactor, proves the connected decay
hypothesis, supplies the final response-matched mixed field jet, proves
weak-facing-root escape, or establishes local-vacuum contraction or a
continuum Yang--Mills mass gap.
\end{theorem}

\begin{proof}
Items~\textup{(E13.1)--(E13.2)} are
\Cref{thm:YMA13-cofactor-identity,cor:YMA13-selector-cofactor,prop:YMA13-common-factor}.
Items~\textup{(E13.3)--(E13.5)} are
\Cref{thm:YMA13-Creutz-covariance,thm:YMA13-response-Gram,cor:YMA13-simple-root-source-floor}.
Items~\textup{(E13.6)--(E13.8)} are
\Cref{prop:YMA13-two-pair-acquisition,thm:YMA13-pair-convolution,thm:YMA13-connected-P-kernel,prop:YMA13-Markov-blanket,cor:YMA13-connected-tail}.
Items~\textup{(E13.9)--(E13.10)} are
\Cref{thm:YMA13-cofactor-certificate,prop:YMA13-higher-cofactor}.
\end{proof}

\begin{assessment}[V13 termination]
\label{ass:YMA13-termination}
The cofactor reduction, Creutz-source Gram, connected plaquette kernel,
Markov-boundary factorization, exact Haar convolution, interval certificate,
and marginal-ramification identity are \(\pass\) at their stated finite or
conditional scopes.  This continuation removes one more circular layer: the
first missing response row is not four marked quotient cards but the single
root cofactor \(\mathfrak c_u[P]\), equivalently the cross covariance
\(-\operatorname{Cov}(\mathfrak s_{\rm Cr},P)\).  Its physical value remains
unpopulated, so the selected symbol remains \(\stopstatus\).  No negative
physical mode has been produced.
\end{assessment}

\section*{Version V13 append-only continuation record}
\addcontentsline{toc}{section}{Version V13 append-only continuation record}
The complete Version~V12 source above is preserved byte for byte up to its
sole terminal document-closing command.  Its SHA--256 digest is
\begin{center}\scriptsize\ttfamily
40919e033e6e162d13732123da49176013c491b2eb085a8aee7bc860cc58e22d.
\end{center}
For Version~V14, preserve this accumulated source, append new mathematical
material immediately before the final document-closing command, and increment
the filename to \texttt{\_V14.tex}.  Keep the physical response target,
weak-facing root, loop family, full unnormalized action, actual selector
tangent, plaquette marginal writer, generated interaction, rooted Hodge
carrier, selected axial mode, colour normalization, and local/topological
convention fixed before reading any cofactor, covariance, head, tail,
coefficient, or field-jet value.

% =============================================================================
% END APPENDED VERSION V13 MATERIAL
% =============================================================================


% =============================================================================
% START APPENDED VERSION V14 MATERIAL
% =============================================================================

\clearpage
\hypersetup{
 pdftitle={Renewal Yang--Mills Reduced-Fibre Wilson Shell Symbol and Local-Vacuum Programme V14},
 pdfsubject={Projective response coordinates, invariant marginal efficiency, and a cofactor-native physical mode card}
}
\pdfinfo{/Title (Renewal Yang--Mills Reduced-Fibre Wilson Shell Symbol and Local-Vacuum Programme V14)
 /Subject (Projective response coordinates, invariant marginal efficiency, and a cofactor-native physical mode card)}

\section{Version V14: projective response coordinates and a cofactor-native physical mode card}
\label{sec:YMA14-entry}

Version~V13 is preserved above.  It reduced the first independent marginal
row to the single root cofactor
\(\mathfrak c_u[P]\), equivalently to the same-law covariance
\(-\operatorname{Cov}(\mathfrak s_{\rm Cr},P)\).  Continuing to introduce
new quotient writers around that scalar would be circular.  There is,
however, one upstream correction to the order of acquisition.  The numerical
coefficient
\begin{equation}
 b^{\rm marg}=-\frac{\mathfrak c_u[\Psi]}{\mathfrak c_u[P]}
 \label{eq:YMA14-b-recall}
\end{equation}
is not itself invariant under a harmless rescaling of the declared plaquette
coordinate.  The invariant objects are the response-null interaction ray and
the completely assembled physical Hessian.  A divergent coefficient can be
created by changing units while leaving both objects unchanged.

This continuation therefore replaces the order
\[
 \mathfrak c_u[P]\longrightarrow b^{\rm marg}
 \longrightarrow\text{six separate mixed-jet blocks}
\]
by the target-native order
\begin{equation}
 \boxed{
 \begin{gathered}
 (\mathfrak c_P,\mathfrak c_\Psi)
 \longrightarrow
 \text{one response-null projective interaction}\
 \Downarrow\\[-0.2em]
 \text{one cofactor-weighted deterministic writer}
 \ -\ 
 \text{one positive two-copy covariance writer}\
 \Downarrow\\[-0.2em]
 \text{physical axial margin or an attained mode witness}.
 \end{gathered}}
 \label{eq:YMA14-new-order}
\end{equation}
Here and below
\begin{equation}
 \mathfrak c_P:=\mathfrak c_u[P],
 \qquad
 \mathfrak c_\Psi:=\mathfrak c_u[\Psi].
 \label{eq:YMA14-two-cofactors}
\end{equation}
The full mixed panel of V11 remains useful when many marginal coefficients
will be reused.  It is not required for the one already selected physical
mode.

\begin{firewall}[Projective compression is not a change of marginal scheme]
\label{fire:YMA14-no-scheme-change}
The plaquette writer \(P\), generated insertion \(\Psi\), weak-facing root,
response target, unnormalized cylinder, rooted Hodge carrier, and local versus
topological split remain frozen.  The transformations below identify only
changes of coordinate which leave the final action unchanged up to a
field-independent constant.  Replacing \(P\) by \(P+K\) for a response-null
but field-active interaction \(K\), without transforming the rest of the
action so that the same final writer is recovered, is a different
renormalization scheme and is not licensed here.
\end{firewall}

% =============================================================================
\section{Affine recoding and the invariant response-null interaction}
\label{sec:YMA14-projective-interaction}
% =============================================================================

Work on one selected finite action law and let
\begin{equation}
 d:=\mathscr R(P),
 \qquad
 n:=\mathscr R(\Psi).
 \label{eq:YMA14-dn}
\end{equation}
Constants independent of the configuration and of the coarse field are
annihilated by \(\mathscr R\) and have zero field jets.

\begin{theorem}[Exact affine recoding covariance]
\label{thm:YMA14-affine-recoding}
Let \(a\in\mathbb R\setminus\{0\}\), \(\lambda,\alpha,\gamma\in\mathbb R\),
and define a recoded interaction chart by
\begin{equation}
 P'=aP+\alpha\one,
 \qquad
 \Psi'=\Psi+\lambda P+\gamma\one.
 \label{eq:YMA14-chart-recoding}
\end{equation}
Then
\begin{equation}
 \boxed{
 d'=ad,
 \qquad
 n'=n+\lambda d.}
 \label{eq:YMA14-response-recoding}
\end{equation}
On the regular branch \(d\ne0\), if
\begin{equation}
 b=-n/d,
 \qquad
 b'=-n'/d',
 \label{eq:YMA14-two-b}
\end{equation}
then
\begin{equation}
 \boxed{b'=\frac{b-\lambda}{a}}
 \label{eq:YMA14-b-transform}
\end{equation}
and
\begin{equation}
 \boxed{
 \Psi'+b'P'
 =\Psi+bP+(\gamma+b'\alpha)\one.}
 \label{eq:YMA14-matched-interaction-invariance}
\end{equation}
Thus the normalized law and every coarse-field jet are unchanged.

Define the homogeneous response-null interaction
\begin{equation}
 \boxed{
 \mathcal I_{P,\Psi}^{\wedge}
 :=d\Psi-nP.}
 \label{eq:YMA14-response-wedge}
\end{equation}
It obeys
\begin{equation}
 \mathscr R(\mathcal I_{P,\Psi}^{\wedge})=0,
 \qquad
 \Psi+bP=d^{-1}\mathcal I_{P,\Psi}^{\wedge},
 \label{eq:YMA14-wedge-null-and-normalized}
\end{equation}
and under \eqref{eq:YMA14-chart-recoding},
\begin{equation}
 \boxed{
 \mathcal I_{P',\Psi'}^{\wedge}
 =a\mathcal I_{P,\Psi}^{\wedge}
 +\bigl(ad\gamma-(n+\lambda d)\alpha\bigr)\one.}
 \label{eq:YMA14-wedge-transform}
\end{equation}
Hence its nonconstant ray and all its field jets are chart invariant.
\end{theorem}

\begin{proof}
Linearity of the response functional and its annihilation of constants give
\eqref{eq:YMA14-response-recoding}.  Substitution gives
\[
 b'=-\frac{n+\lambda d}{ad}
    =\frac{-n/d-\lambda}{a}
    =\frac{b-\lambda}{a},
\]
and then
\[
 \Psi'+b'P'
 =\Psi+\lambda P+\gamma\one
  +(b-\lambda)P+b'\alpha\one,
\]
which is \eqref{eq:YMA14-matched-interaction-invariance}.  The response of
\(d\Psi-nP\) is \(dn-nd=0\), and division by \(d\) gives the normalized
matched interaction.  Finally,
\begin{align*}
 d'\Psi'-n'P'
 & =ad(\Psi+\lambda P+\gamma\one)
   -(n+\lambda d)(aP+\alpha\one)\\
 & =a(d\Psi-nP)
   +\bigl(ad\gamma-(n+\lambda d)\alpha\bigr)\one.
\end{align*}
A field-independent constant changes only the absolute partition
normalization, not the normalized law or a coarse-field Hessian.
\end{proof}

\begin{corollary}[Cofactor-native projective interaction]
\label{cor:YMA14-cofactor-wedge}
At the selected root, V13 gives a positive scalar
\begin{equation}
 g_*=e^{-u}B_*>0
 \label{eq:YMA14-gstar}
\end{equation}
such that
\begin{equation}
 \mathfrak c_P=-g_*d,
 \qquad
 \mathfrak c_\Psi=-g_*n.
 \label{eq:YMA14-cofactor-response-proportionality}
\end{equation}
Consequently
\begin{equation}
 \boxed{
 \mathcal I_u^{\rm cof}
 :=\mathfrak c_P\Psi-\mathfrak c_\Psi P
 =-g_*\mathcal I_{P,\Psi}^{\wedge},}
 \label{eq:YMA14-cofactor-interaction}
\end{equation}
and, whenever \(\mathfrak c_P\ne0\),
\begin{equation}
 \boxed{
 \Psi+b^{\rm marg}P
 =\frac{\mathcal I_u^{\rm cof}}{\mathfrak c_P},
 \qquad
 b^{\rm marg}=-\frac{\mathfrak c_\Psi}{\mathfrak c_P}.}
 \label{eq:YMA14-cofactor-normalization}
\end{equation}
The common factor \(g_*\), including the product of loop numerators, is absent
from the projective interaction.
\end{corollary}

\begin{proof}
Use the V13 identity
\(\mathfrak c_u[X]=-e^{-u}B_*\mathscr R(X)\) for
\(X=P,\Psi\).  The remaining formulas are direct substitutions.
\end{proof}

\begin{countertheorem}[The coefficient can diverge while the physical action is fixed]
\label{cth:YMA14-b-not-physical}
Fix a regular card with \(b\ne0\).  For any sequence
\(a_r>0\) with \(a_r\to0\), set
\begin{equation}
 P_r=a_rP,
 \qquad
 d_r=a_rd,
 \qquad
 b_r=b/a_r.
 \label{eq:YMA14-rescaled-P}
\end{equation}
Then
\begin{equation}
 |b_r|\longrightarrow\infty,
 \qquad
 \boxed{\Psi+b_rP_r=\Psi+bP.}
 \label{eq:YMA14-b-diverges-action-fixed}
\end{equation}
The final law, every deterministic field jet, every score covariance, and the
physical shell Hessian are identical for all \(r\).  Hence divergence of the
coordinate \(b_r\) alone is not a physical instability.
\end{countertheorem}

\begin{proof}
The first assertion is immediate.  The second is
\(b_rP_r=(b/a_r)(a_rP)=bP\).  Every claimed equality follows because the
complete final action is the same writer, not merely a response-equivalent
writer.
\end{proof}

\begin{assessment}[Qualification of the V11 amplification language]
\label{ass:YMA14-V11-qualification}
The quantitative theorem \Cref{thm:YMA11-b-blowup} remains correct under its
stated uniform hypotheses.  Those hypotheses include a cutoff-independent
lower bound on the plaquette-score covariance.  Under the harmless recoding
\eqref{eq:YMA14-rescaled-P}, that covariance scales as \(a_r^2\) and the
hypothesis fails.  What is withdrawn is only the interpretation of
\(|b_r|\to\infty\) by itself as physical amplification.  The invariant
quantities are \(b_rP_r\), \(b_rH_{P_r}\),
\(b_r^2C_{P_rP_r}\), or equivalently the homogeneous card constructed below.
\end{assessment}

% =============================================================================
\section{Response--Riesz geometry and the exact price of the plaquette chart}
\label{sec:YMA14-response-Riesz}
% =============================================================================

Retain the centered Creutz source \(\mathfrak s_{\rm Cr}\) of V13 and write
\begin{equation}
 s:=\mathfrak s_{\rm Cr},
 \qquad
 v_s:=\operatorname{Var}(s)=\norm{s}_2^2.
 \label{eq:YMA14-source-stock}
\end{equation}
For any real writer \(X\), write \(X^\circ=X-\mathbb EX\).  Then
\begin{equation}
 \mathscr R(X)=-\ip{s}{X^\circ}.
 \label{eq:YMA14-response-Riesz}
\end{equation}
Let
\begin{equation}
 v_P:=\operatorname{Var}(P),
 \qquad
 \eta_P:=\frac{d^2}{v_sv_P}
 \label{eq:YMA14-eta}
\end{equation}
whenever \(v_sv_P>0\).  Positivity of the V13 response Gram gives
\(0\le\eta_P\le1\).  This number is unchanged by
\(P\mapsto aP+\alpha\one\).

\begin{theorem}[Least-variance response correction and plaquette null excess]
\label{thm:YMA14-minimum-correction}
Assume \(v_s>0\).  Among all centered action corrections \(C\) satisfying
\begin{equation}
 \mathscr R(C)=-n,
 \label{eq:YMA14-correction-constraint}
\end{equation}
the unique minimum-variance correction is
\begin{equation}
 \boxed{C_{\min}=\frac{n}{v_s}s,}
 \label{eq:YMA14-Cmin}
\end{equation}
with
\begin{equation}
 \boxed{\operatorname{Var}(C_{\min})=\frac{n^2}{v_s}.}
 \label{eq:YMA14-Cmin-cost}
\end{equation}

Suppose now that \(d\ne0\).  Define the source-orthogonal part of the
plaquette writer by
\begin{equation}
 P_\perp
 :=P^\circ+\frac{d}{v_s}s.
 \label{eq:YMA14-Pperp}
\end{equation}
Then
\begin{equation}
 \ip{s}{P_\perp}=0,
 \qquad
 P^\circ=-\frac{d}{v_s}s+P_\perp.
 \label{eq:YMA14-Psplit}
\end{equation}
For the physical plaquette correction \(C_P=bP^\circ\), with \(b=-n/d\),
\begin{equation}
 \boxed{
 C_P=C_{\min}+bP_\perp.}
 \label{eq:YMA14-correction-split}
\end{equation}
The two terms on the right are orthogonal, and
\begin{align}
 \boxed{
 \operatorname{Var}(C_P)
 =\frac{n^2}{v_s\eta_P},}
 \label{eq:YMA14-plaquette-correction-cost}\\
 \boxed{
 \operatorname{Var}(C_P-C_{\min})
 =\frac{n^2}{v_s}
  \left(\eta_P^{-1}-1\right).}
 \label{eq:YMA14-null-excess-cost}
\end{align}
The second quantity is the exact response-null action stock introduced by
using the one-dimensional plaquette chart rather than the Riesz-optimal
correction.
\end{theorem}

\begin{proof}
Condition \eqref{eq:YMA14-correction-constraint} is
\(\ip{s}{C}=n\).  Every such correction has the unique orthogonal
representation
\[
 C=\frac{n}{v_s}s+q,
 \qquad q\perp s.
\]
Therefore
\[
 \norm C_2^2=\frac{n^2}{v_s}+\norm q_2^2,
\]
which proves uniqueness and \eqref{eq:YMA14-Cmin-cost}.

Since \(d=-\ip{s}{P^\circ}\), the writer in
\eqref{eq:YMA14-Pperp} is orthogonal to \(s\).  Multiplying
\eqref{eq:YMA14-Psplit} by \(b=-n/d\) gives
\eqref{eq:YMA14-correction-split}.  Orthogonality yields
\[
 \operatorname{Var}(C_P)=b^2v_P
 =\frac{n^2v_P}{d^2}
 =\frac{n^2}{v_s\eta_P}.
\]
Subtracting the minimum cost proves
\eqref{eq:YMA14-null-excess-cost}.
\end{proof}

\begin{corollary}[Three-factor marginal regularity and invariant conditioning]
\label{cor:YMA14-three-factor-regularity}
Whenever \(v_sv_P>0\),
\begin{equation}
 \boxed{d^2=v_sv_P\eta_P.}
 \label{eq:YMA14-d-factorization}
\end{equation}
Consequently a uniform denominator floor follows from three independent
floors:
\begin{equation}
 v_s\ge s_*>0,
 \qquad
 v_P\ge p_*>0,
 \qquad
 \eta_P\ge\eta_*>0
 \quad\Longrightarrow\quad
 |d|\ge\sqrt{s_*p_*\eta_*}.
 \label{eq:YMA14-three-floor}
\end{equation}
The scale-invariant response condition number of the plaquette line is
\begin{equation}
 \boxed{\kappa_P^{\rm resp}:=\eta_P^{-1/2}.}
 \label{eq:YMA14-response-condition-number}
\end{equation}
Thus collapse of \(d\) is typed as response-source collapse, plaquette-stock
collapse, or source--plaquette angle collapse; a change of units for \(P\)
does not alter the angle branch.
\end{corollary}

\begin{proof}
The identity is the definition of \(\eta_P\).  The floor follows by taking
square roots.  Invariance under affine rescaling of \(P\) follows because
both \(d^2\) and \(v_P\) acquire the same factor \(a^2\).
\end{proof}

\begin{corollary}[Response cost times field influence]
\label{cor:YMA14-response-field-factorization}
Let \(j_P[v]\) be the coarse-field score of the plaquette interaction on a
fixed physical polarization \(v\), and put
\begin{equation}
 C_{PP}[v]:=\operatorname{Var}(j_P[v]).
 \label{eq:YMA14-CPPv}
\end{equation}
When \(v_P>0\), define the action-to-field influence ratio
\begin{equation}
 \Lambda_P[v]:=\frac{C_{PP}[v]}{v_P}\ge0.
 \label{eq:YMA14-P-influence}
\end{equation}
Then the quadratic plaquette-score debit in the matched Hessian factorizes as
\begin{equation}
 \boxed{
 b^2C_{PP}[v]
 =\frac{n^2}{v_s\eta_P}\,\Lambda_P[v].}
 \label{eq:YMA14-invariant-covariance-debit}
\end{equation}
Both factors on the right are invariant under
\(P\mapsto aP+\alpha\one\).  Thus an actual covariance amplification is a
joint failure of response efficiency and field influence, not a statement
about the coordinate \(b\) alone.
\end{corollary}

\begin{proof}
Use \(b^2=n^2/d^2\),
\(d^2=v_sv_P\eta_P\), and
\(C_{PP}[v]=v_P\Lambda_P[v]\).  Under rescaling of \(P\), both numerator and
denominator in \(\Lambda_P\) acquire the same factor, while \(\eta_P\) is
unchanged.
\end{proof}

\begin{firewall}[The Riesz minimizer is diagnostic, not a substituted action]
\label{fire:YMA14-Riesz-not-scheme}
The writer \(C_{\min}\) need not belong to the declared local interaction
algebra, need not have the plaquette support or reflection properties required
by the renormalization scheme, and is not inserted into the Wilson action.
The theorem measures the exact response-null excess carried by the already
frozen plaquette chart.  A large excess is a chart-efficiency or physical
interaction issue to be evaluated; it does not authorize a post hoc change of
marginal writer.
\end{firewall}

% =============================================================================
\section{The cofactor-native signed Wilson mode card}
\label{sec:YMA14-cofactor-card}
% =============================================================================

Return to the V11 field notation.  Let \(F_0\) be the complete pre-retuning
action on the pulled reduced fibre, let \(F_P\) be the same-chart plaquette
interaction, and suppose the regular cofactor branch
\begin{equation}
 c:=\mathfrak c_P\ne0,
 \qquad
 e:=\mathfrak c_\Psi,
 \qquad
 b=-e/c
 \label{eq:YMA14-ceb}
\end{equation}
has been loaded.  All expectations below are taken in the one exact final law
with action
\begin{equation}
 F_b=F_0+bF_P.
 \label{eq:YMA14-final-law}
\end{equation}
For a retained physical bank synthesis \(V\), write
\begin{align}
 \mathbf j_0&:=V^*D_VF_0,
 &\mathbf j_P&:=V^*D_VF_P,
 \label{eq:YMA14-two-bank-scores}\\
 \mathbf K_0&:=V^*D_V^2F_0V,
 &\mathbf K_P&:=V^*D_V^2F_PV.
 \label{eq:YMA14-two-bank-curvatures}
\end{align}
Define the homogeneous writers
\begin{align}
 \boxed{
 \mathbf j^{\wedge}_{c,e}
 :=c\mathbf j_0-e\mathbf j_P,}
 \label{eq:YMA14-homogeneous-score}\\
 \boxed{
 \mathbf K^{\wedge}_{c,e}
 :=c^2\mathbf K_0-ce\mathbf K_P.}
 \label{eq:YMA14-homogeneous-curvature}
\end{align}

\begin{theorem}[Exact cofactor-native Hessian identity]
\label{thm:YMA14-cofactor-Hessian}
Let \(H_b^{\rm RG}\) be the exact physical bank Hessian of the final law
\eqref{eq:YMA14-final-law}.  Then
\begin{equation}
 \boxed{
 \mathfrak H_{c,e}^{\wedge}
 :=\mathbb E_b\mathbf K^{\wedge}_{c,e}
   -\operatorname{Cov}_b(\mathbf j^{\wedge}_{c,e})
 =c^2H_b^{\rm RG}.}
 \label{eq:YMA14-cofactor-Hessian-identity}
\end{equation}
If \(T\) is the already typed topological Hessian and \(S\succ0\) is the
selected flat reference, then for every \(\delta\ge0\),
\begin{equation}
 \boxed{
 \mathfrak M_{c,e}^{\wedge}(\delta)
 :=\mathfrak H_{c,e}^{\wedge}-c^2(T+\delta S)
 =c^2\bigl(H_b^{\rm RG}-T-\delta S\bigr).}
 \label{eq:YMA14-cofactor-margin}
\end{equation}
Since \(c^2>0\), the physical margin and the cofactor-native margin have the
same inertia and the same negative witness vectors.
\end{theorem}

\begin{proof}
The final score and deterministic curvature are
\[
 \mathbf j_b=\mathbf j_0-\frac ec\mathbf j_P
             =c^{-1}\mathbf j^{\wedge}_{c,e},
 \qquad
 \mathbf K_b=\mathbf K_0-\frac ec\mathbf K_P
             =c^{-2}\mathbf K^{\wedge}_{c,e}.
\]
Insert them into the exact score--Hessian identity
\(H_b^{\rm RG}=\mathbb E_b\mathbf K_b-
\operatorname{Cov}_b(\mathbf j_b)\) and multiply by \(c^2\).  Subtracting
\(c^2(T+\delta S)\) gives \eqref{eq:YMA14-cofactor-margin}.  Multiplication
by a positive scalar preserves inertia and witness directions.
\end{proof}

\begin{corollary}[One deterministic Read and one positive two-copy Read]
\label{cor:YMA14-two-read-card}
Let \(z_1,z_2\) be independent samples from the final law.  Then
\begin{equation}
 \boxed{
 \operatorname{Cov}_b(\mathbf j^{\wedge}_{c,e})
 =\frac12\mathbb E_b^{\otimes2}
 \left[
  (\mathbf j^{\wedge}_{c,e}(z_1)
   -\mathbf j^{\wedge}_{c,e}(z_2))
  (\mathbf j^{\wedge}_{c,e}(z_1)
   -\mathbf j^{\wedge}_{c,e}(z_2))^*
 \right].}
 \label{eq:YMA14-two-copy-covariance}
\end{equation}
Hence the one selected bank requires only
\begin{equation}
 \boxed{
 \mathbb E_b\mathbf K^{\wedge}_{c,e}
 \quad\text{and}\quad
 \frac12\mathbb E_b^{\otimes2}
 [\Delta\mathbf j^{\wedge}_{c,e}
  (\Delta\mathbf j^{\wedge}_{c,e})^*],}
 \label{eq:YMA14-two-physical-reads}
\end{equation}
together with the already separated \(T,S\).  The second writer is positive
semidefinite sample by sample.  The four covariance blocks of V11 are not
separately required unless they are intended for reuse at another coefficient.
\end{corollary}

\begin{proof}
For any square-integrable random vector \(X\), independence gives
\[
 \frac12\mathbb E[(X_1-X_2)(X_1-X_2)^*]
 =\mathbb E[XX^*]-\mathbb EX(\mathbb EX)^*
 =\operatorname{Cov}(X).
\]
Apply this with \(X=\mathbf j^{\wedge}_{c,e}\).
\end{proof}

\begin{theorem}[Chart covariance of the physical mode card]
\label{thm:YMA14-card-covariance}
Let \(a\ne0\) and \(\lambda\in\mathbb R\).  Recode the same final action by
\begin{equation}
 F_0'=F_0+\lambda F_P+C_0,
 \qquad
 F_P'=aF_P+C_P,
 \label{eq:YMA14-field-chart-recoding}
\end{equation}
where \(C_0,C_P\) are field- and fibre-independent constants, and set
\begin{equation}
 c'=ac,
 \qquad
 e'=e+\lambda c.
 \label{eq:YMA14-ce-transform}
\end{equation}
Then the recoded coefficient is \(b'=(b-\lambda)/a\), the final action agrees
with \(F_b\) up to a constant, and
\begin{align}
 \boxed{
 \mathbf j_{c',e'}^{\wedge\prime}
 =a\mathbf j_{c,e}^{\wedge},}
 \label{eq:YMA14-score-scale}\\
 \boxed{
 \mathbf K_{c',e'}^{\wedge\prime}
 =a^2\mathbf K_{c,e}^{\wedge},}
 \label{eq:YMA14-curvature-scale}\\
 \boxed{
 \mathfrak M_{c',e'}^{\wedge\prime}(\delta)
 =a^2\mathfrak M_{c,e}^{\wedge}(\delta).}
 \label{eq:YMA14-margin-scale}
\end{align}
The physical verdict is therefore independent of the affine interaction
coordinates.
\end{theorem}

\begin{proof}
The coefficient and final-action identities are
\Cref{thm:YMA14-affine-recoding}.  Field differentiation gives
\[
 \mathbf j_0'=\mathbf j_0+\lambda\mathbf j_P,
 \qquad
 \mathbf j_P'=a\mathbf j_P,
\]
and the same formulas for \(\mathbf K_0',\mathbf K_P'\).  Hence
\begin{align*}
 c'\mathbf j_0'-e'\mathbf j_P'
 &=ac(\mathbf j_0+\lambda\mathbf j_P)
   -(e+\lambda c)a\mathbf j_P
 =a(c\mathbf j_0-e\mathbf j_P),\\
 (c')^2\mathbf K_0'-c'e'\mathbf K_P'
 &=a^2c^2(\mathbf K_0+\lambda\mathbf K_P)
   -a^2c(e+\lambda c)\mathbf K_P\\
 &=a^2(c^2\mathbf K_0-ce\mathbf K_P).
\end{align*}
Expectation and covariance preserve these scalar factors.  The reference and
topological subtraction acquire \((c')^2=a^2c^2\), proving the last identity.
\end{proof}

\begin{firewall}[The final law must still be the physical law]
\label{fire:YMA14-projective-not-law}
The homogeneous card removes a numerically unstable division from the
\emph{sign test}.  It does not manufacture the response-matched probability
law.  The expectation in \eqref{eq:YMA14-cofactor-Hessian-identity} must be
taken in the actual law with coefficient \(b=-e/c\), or in an independently
proved equivalent unnormalized cylinder.  Evaluating the homogeneous writers
in the pre-retuning law would be a different card.
\end{firewall}

% =============================================================================
\section{Denominator-free finite certification}
\label{sec:YMA14-finite-certificate}
% =============================================================================

The V11 expansion
\begin{equation}
 H_b^{\rm RG}-T-\delta S
 =A+bL-b^2C,
 \qquad C\succeq0,
 \label{eq:YMA14-V11-polynomial}
\end{equation}
uses
\(A=H_{\rm pre}-T-\delta S\) and the same-law Hermitian blocks \(L,C\).
Substitution of \(b=-e/c\) gives the homogeneous polynomial
\begin{equation}
 \boxed{
 \mathfrak M^{\rm poly}_{c,e}
 :=c^2A-ceL-e^2C
 =c^2(H_b^{\rm RG}-T-\delta S).}
 \label{eq:YMA14-homogeneous-polynomial}
\end{equation}
This is the block-expanded version of
\eqref{eq:YMA14-cofactor-margin}.

\begin{theorem}[Robust cofactor-native spectral interval]
\label{thm:YMA14-robust-cofactor-interval}
Let \(A,L,C\) be Hermitian operators on the finite supported physical bank,
with \(C\succeq0\).  Suppose
\begin{equation}
 |c-\widehat c|\le r_c,
 \qquad
 |e-\widehat e|\le r_e,
 \label{eq:YMA14-ce-errors}
\end{equation}
and
\begin{equation}
 \norm{A-\widehat A}\le\epsilon_A,
 \qquad
 \norm{L-\widehat L}\le\epsilon_L,
 \qquad
 \norm{C-\widehat C}\le\epsilon_C.
 \label{eq:YMA14-ALC-errors}
\end{equation}
Put
\begin{align}
 c_+&:=|\widehat c|+r_c,
 &e_+&:=|\widehat e|+r_e,
 \label{eq:YMA14-ce-plus}\\
 \Delta_{c^2}&:=r_c(2|\widehat c|+r_c),
 &\Delta_{e^2}&:=r_e(2|\widehat e|+r_e),
 \label{eq:YMA14-square-errors}\\
 \Delta_{ce}&:=|\widehat c|r_e+|\widehat e|r_c+r_cr_e,
 \label{eq:YMA14-product-error}
\end{align}
and define
\begin{equation}
 \widehat{\mathfrak M}
 :=\widehat c^{\,2}\widehat A
  -\widehat c\widehat e\widehat L
  -\widehat e^{\,2}\widehat C.
 \label{eq:YMA14-Mhat}
\end{equation}
Then
\begin{equation}
 \boxed{
 \norm{\mathfrak M^{\rm poly}_{c,e}-\widehat{\mathfrak M}}
 \le\epsilon_{\wedge},}
 \label{eq:YMA14-wedge-error}
\end{equation}
where
\begin{equation}
 \boxed{
 \begin{aligned}
 \epsilon_{\wedge}:={}&
 c_+^2\epsilon_A+c_+e_+\epsilon_L+e_+^2\epsilon_C\\
 &+\Delta_{c^2}\norm{\widehat A}
  +\Delta_{ce}\norm{\widehat L}
  +\Delta_{e^2}\norm{\widehat C}.
 \end{aligned}}
 \label{eq:YMA14-wedge-radius}
\end{equation}
If \(|\widehat c|>r_c\), the regular physical card obeys the terminating
alternative
\begin{equation}
 \boxed{
 \begin{array}{rcl}
 \lambda_{\min}(\widehat{\mathfrak M})-\epsilon_{\wedge}>0
 &\Longrightarrow&\pass,\\[0.25em]
 \lambda_{\min}(\widehat{\mathfrak M})+\epsilon_{\wedge}<0
 &\Longrightarrow&\obstruction,\\[0.25em]
 \text{otherwise}&\Longrightarrow&\stopstatus.
 \end{array}}
 \label{eq:YMA14-projective-verdict}
\end{equation}
On the obstruction branch a unit minimum eigenvector of
\(\widehat{\mathfrak M}\) is an attained physical witness.
\end{theorem}

\begin{proof}
Insert and subtract the three terms with exact coefficients and approximate
operators.  The coefficient errors satisfy
\[
 |c^2-\widehat c^{\,2}|\le\Delta_{c^2},
 \qquad
 |e^2-\widehat e^{\,2}|\le\Delta_{e^2},
 \qquad
 |ce-\widehat c\widehat e|\le\Delta_{ce}.
\]
Also \(|c|\le c_+\) and \(|e|\le e_+\).  The triangle inequality gives
\eqref{eq:YMA14-wedge-radius}.  Weyl's inequality gives the positive branch.
If \(v\) is a unit minimum eigenvector of
\(\widehat{\mathfrak M}\), then
\[
 \ip{v}{\mathfrak M^{\rm poly}_{c,e}v}
 \le\lambda_{\min}(\widehat{\mathfrak M})+\epsilon_{\wedge}<0.
\]
Since the interval for \(c\) excludes zero,
\eqref{eq:YMA14-homogeneous-polynomial} transfers the same sign and witness to
the physical margin.  The remaining case has no strict sign certificate.
\end{proof}

\begin{corollary}[Quantitative physical floor without coefficient division]
\label{cor:YMA14-quantitative-floor}
Assume \(|\widehat c|>r_c\).  If
\begin{equation}
 \widehat{\mathfrak M}\succeq
 (\epsilon_{\wedge}+\gamma)I
 \qquad(\gamma>0),
 \label{eq:YMA14-positive-numerator-floor}
\end{equation}
then
\begin{equation}
 \boxed{
 H_b^{\rm RG}-T-\delta S
 \succeq\frac{\gamma}{c_+^2}I.}
 \label{eq:YMA14-physical-floor}
\end{equation}
If a unit vector \(v\) satisfies
\begin{equation}
 \ip{v}{\widehat{\mathfrak M}v}
 \le-(\epsilon_{\wedge}+\gamma),
 \label{eq:YMA14-negative-numerator-witness}
\end{equation}
then
\begin{equation}
 \boxed{
 \ip{v}{(H_b^{\rm RG}-T-\delta S)v}
 \le-\frac{\gamma}{c_+^2}.}
 \label{eq:YMA14-physical-negative-floor}
\end{equation}
\end{corollary}

\begin{proof}
The preceding theorem gives
\(\mathfrak M^{\rm poly}_{c,e}\succeq\gamma I\) in the first case and
\(\mathfrak M^{\rm poly}_{c,e}[v]\le-\gamma\) in the second.  Divide by the
positive scalar \(c^2\le c_+^2\), taking account of the sign.
\end{proof}

\begin{corollary}[Exact marginal-chart branches in projective coordinates]
\label{cor:YMA14-projective-branches}
The response and physical cards terminate in the following order.
\begin{enumerate}[label=\textup{(P14.\arabic*)},leftmargin=5.2em]
\item If \(\mathfrak c_P=0\) and \(\mathfrak c_\Psi\ne0\), no finite
      one-dimensional plaquette coefficient exists.  This is the typed
      plaquette response-chart obstruction of V13.
\item If \(\mathfrak c_P=\mathfrak c_\Psi=0\), acquire the first nonzero
      higher cofactor.  No projective first-order interaction has been selected.
\item If \(\mathfrak c_P\ne0\), form the cofactor-native card directly.
      Failure to separate the sign in
      \eqref{eq:YMA14-projective-verdict} is \(\stopstatus\), not an
      obstruction.
\item Only a negative assembled numerator on an explicit supported transverse
      vector is a physical pure-mode obstruction.
\end{enumerate}
\end{corollary}

\begin{proof}
The first two items are the V13 marginal alternatives.  The third and fourth
are \Cref{thm:YMA14-cofactor-Hessian,thm:YMA14-robust-cofactor-interval}.
\end{proof}

% =============================================================================
\section{Invariant amplification and the corrected cofinal target}
\label{sec:YMA14-cofinal}
% =============================================================================

The response-angle factorization gives a coordinate-free version of the
near-ramification mechanism.  On a unit polarization \(v\), write the V11
physical margin as
\begin{equation}
 M_b[v]=A[v]+bL[v]-b^2C[v].
 \label{eq:YMA14-scalar-margin}
\end{equation}
The covariance term is the invariant debit
\begin{equation}
 b^2C[v]
 =\frac{n^2}{v_s\eta_P}\Lambda_P[v]
 \label{eq:YMA14-invariant-debit-repeat}
\end{equation}
whenever \(C[v]=C_{PP}[v]\) and \(v_P>0\).

\begin{theorem}[Angle-collapse amplification criterion]
\label{thm:YMA14-angle-amplification}
Let \(v_r\) be unit supported physical polarizations on a selected family.
Assume
\begin{equation}
 A_r[v_r]+b_rL_r[v_r]\le H_*<\infty,
 \label{eq:YMA14-affine-upper}
\end{equation}
and, for common constants \(n_*,V_*,\Lambda_*>0\),
\begin{equation}
 |n_r|\ge n_* ,
 \qquad
 0<v_{s,r}\le V_* ,
 \qquad
 \Lambda_{P,r}[v_r]\ge\Lambda_*.
 \label{eq:YMA14-angle-data}
\end{equation}
If \(\eta_{P,r}\to0\), then
\begin{equation}
 \boxed{M_{b_r}[v_r]\longrightarrow-\infty.}
 \label{eq:YMA14-angle-negative}
\end{equation}
In particular the selected physical mode is eventually an obstruction.
This conclusion is invariant under every affine rescaling of the plaquette
coordinate.
\end{theorem}

\begin{proof}
By \eqref{eq:YMA14-invariant-debit-repeat},
\[
 b_r^2C_r[v_r]
 \ge\frac{n_*^2\Lambda_*}{V_*\eta_{P,r}}
 \longrightarrow\infty.
\]
Subtract this quantity from the uniformly bounded affine part
\eqref{eq:YMA14-affine-upper}.  Every factor in the lower bound is invariant
under \(P_r\mapsto a_rP_r+\alpha_r\one\), so the conclusion is coordinate
independent.
\end{proof}

\begin{assessment}[What V14 removes from the acquisition]
\label{ass:YMA14-reduction}
Once \((\mathfrak c_P,\mathfrak c_\Psi)\) and the final law are loaded, V14
requires neither the numerical coefficient as a primary output nor the full
six-entry deterministic/covariance decomposition.  The physical mode is one
cofactor-weighted deterministic expectation minus one positive two-copy
covariance, with the already typed topological and reference rows.  The
response Gram remains useful only to classify failure of the plaquette chart
into stock collapse, angle collapse, or regular incidence.
\end{assessment}

\begin{openproblem}[V15: execute the cofactor-native physical axial card]
\label{op:YMA14-V15}
Preserve every frozen choice and proceed in the following order.
\begin{enumerate}[label=\textup{(V15.\arabic*)},leftmargin=5.2em]
\item Acquire outward intervals for
      \(\mathfrak c_P\) and \(\mathfrak c_\Psi\) from the already specified
      unnormalized pair cylinder, or from the exactly equivalent common-law
      response Gram.  Do not report \(b\) without its physical plaquette
      normalization.
\item If \(\mathfrak c_P\ne0\), load the exact final law and acquire directly
      the two writers in \eqref{eq:YMA14-two-physical-reads} on the first
      predeclared axial polarization.  Do not separately populate the complete
      V11 block unless reuse at another coefficient is intended.
\item Attach the existing topological interval and exact reference, then use
      \Cref{thm:YMA14-robust-cofactor-interval} to return
      \(\pass/\obstruction/\stopstatus\).
\item If the cofactor interval approaches zero along a cofinal family, acquire
      \((v_s,v_P,\eta_P)\) and the field-influence ratio
      \(\Lambda_P\).  Classify scale collapse separately from angle collapse;
      use \Cref{thm:YMA14-angle-amplification} only on its stated physical
      incidence branch.
\item Open the boundary/Witten ancestry and differentiated tail only after the
      assembled direct card has passed, failed, or produced a collapsing
      invariant margin requiring attribution.
\end{enumerate}
If the actual selected-root cylinder remains unavailable, the correct output
is the same named physical \(\stopstatus\).  A new root, response, current,
or boundary compiler does not substitute for the two cofactor-native Reads.
\end{openproblem}

% =============================================================================
\section{Version V14 proof ledger and append-only boundary}
\label{sec:YMA14-ledger}
% =============================================================================

The accompanying exact replay
\begin{center}
\path{YMA_V14_projective_response_replay.py}
\end{center}
checks a finite rational response--Riesz card, affine interaction recoding,
the response-null wedge, the cofactor-native Hessian, its positive two-copy
covariance, homogeneous chart covariance, coefficient blow-up with fixed
physical action, the homogeneous V11 polynomial, and one exact sample of the
robust error radius.  Its script and output digests are
\begin{center}\scriptsize\ttfamily
1c223831f43e8a9c2c125f4f4f7036dc4c8282df0a9f4392f93dc10ca752c824,\\[-0.15em]
389bb36eb36228e660684beeab43a98cff6168d4a043f709f0a6a1d1cca34925.
\end{center}
It uses exact rational arithmetic.  It does not evaluate a Wilson path
integral, a weak-facing root, a marginal cofactor, or a physical mode.

\begingroup\small
\begin{longtable}{>{\raggedright\arraybackslash}p{0.27\textwidth}
                  >{\raggedright\arraybackslash}p{0.20\textwidth}
                  >{\raggedright\arraybackslash}p{0.45\textwidth}}
\caption{Version V14 proof-status ledger.}\label{tab:YMA14-ledger}\\
\toprule
\textbf{Row}&\textbf{Status}&\textbf{Exact scope}\\
\midrule
\endfirsthead
\toprule
\textbf{Row}&\textbf{Status}&\textbf{Exact scope}\\
\midrule
\endhead
Affine interaction recoding&Proved exactly
&The coefficient transforms, while the final interaction and all field jets
remain unchanged up to a harmless constant.\\[0.3em]
Response-null interaction ray&Proved exactly
&\(d\Psi-nP\), equivalently
\(\mathfrak c_P\Psi-\mathfrak c_\Psi P\), is the chart-covariant physical
null direction.\\[0.3em]
Coefficient blow-up&Typed as coordinate dependent
&Rescaling \(P\) can make \(|b|\to\infty\) with the complete final action and
Hessian fixed.\\[0.3em]
Response--Riesz correction&Proved sharply
&The minimum correction has variance \(n^2/v_s\); the plaquette line pays the
exact factor \(\eta_P^{-1}\).\\[0.3em]
Plaquette null excess&Proved exactly
&The extra response-null correction stock is
\(n^2v_s^{-1}(\eta_P^{-1}-1)\).\\[0.3em]
Action-to-field amplification&Factorized exactly
&The plaquette covariance debit is response cost times the invariant field
influence \(\Lambda_P\).\\[0.3em]
Cofactor-native Hessian&Proved exactly
&One homogeneous deterministic writer minus one homogeneous covariance equals
\(\mathfrak c_P^2\) times the physical Hessian.\\[0.3em]
Target-native acquisition&Reduced
&One deterministic Read and one positive two-copy covariance Read replace the
full V11 mixed block for the single selected coefficient.\\[0.3em]
Chart covariance of the mode card&Proved exactly
&Every harmless affine recoding scales the homogeneous margin by a positive
square and leaves its verdict unchanged.\\[0.3em]
Robust projective interval&Proved
&Explicit cofactor and operator error radii give a direct
\(\pass/\obstruction/\stopstatus\) test without dividing by the marginal
cofactor.\\[0.3em]
Angle-collapse obstruction&Conditional physical theorem
&With noncollapsed generated response and field influence, collapse of the
source--plaquette angle forces a negative mode.\\[0.3em]
Selected cofactors and final-law Reads&Open physical values
&No supplied cylinder populates \(\mathfrak c_P,\mathfrak c_\Psi\) or the two
cofactor-weighted mode writers.\\[0.3em]
Selected response-matched symbol&\stopstatus
&The exact direct card remains unexecuted; no negative physical mode is
claimed.\\[0.3em]
Local-vacuum contraction and mass&Not inferred
&No projective response identity is promoted to an OS-time or continuum mass
gap.\\
\bottomrule
\end{longtable}
\endgroup

\begin{theorem}[Version V14 projective reduction]
\label{thm:YMA14-terminal}
Versions~V1--V13 are retained.  Version~V14 additionally proves:
\begin{enumerate}[label=\textup{(E14.\arabic*)},leftmargin=5.5em]
\item the marginal coefficient is an interaction-chart coordinate, while the
      response-null action ray is invariant under every harmless affine
      recoding;
\item coefficient divergence alone can be a pure change-of-units artifact;
\item the Creutz source gives a sharp Riesz theorem for response correction,
      and the plaquette scheme has exact efficiency \(\eta_P\);
\item the physical plaquette-score debit factorizes into response correction
      stock and an invariant action-to-field influence;
\item the root cofactors define one homogeneous score and one homogeneous
      deterministic curvature whose signed difference is exactly
      \(\mathfrak c_P^2\) times the physical Hessian;
\item the complete physical margin is therefore acquired by one deterministic
      expectation and one positive two-copy covariance, not by a separately
      normalized six-block panel;
\item affine chart changes scale the homogeneous card by a positive square;
\item explicit outward error radii give a denominator-free finite spectral
      verdict and an attained physical witness on the negative branch;
\item source--plaquette angle collapse yields a coordinate-invariant physical
      amplification theorem once the generated response and field influence
      remain visible.
\end{enumerate}
No item populates the selected cofactors, constructs the final selected law,
evaluates the two physical Reads, proves weak-facing-root escape, or establishes
local-vacuum contraction or a continuum Yang--Mills mass gap.
\end{theorem}

\begin{proof}
Items~\textup{(E14.1)--(E14.2)} are
\Cref{thm:YMA14-affine-recoding,cor:YMA14-cofactor-wedge,cth:YMA14-b-not-physical}.
Items~\textup{(E14.3)--(E14.4)} are
\Cref{thm:YMA14-minimum-correction,cor:YMA14-three-factor-regularity,cor:YMA14-response-field-factorization}.
Items~\textup{(E14.5)--(E14.7)} are
\Cref{thm:YMA14-cofactor-Hessian,cor:YMA14-two-read-card,thm:YMA14-card-covariance}.
Items~\textup{(E14.8)--(E14.9)} are
\Cref{thm:YMA14-robust-cofactor-interval,cor:YMA14-quantitative-floor,thm:YMA14-angle-amplification}.
\end{proof}

\begin{assessment}[V14 termination]
\label{ass:YMA14-termination}
The projective interaction, response-efficiency decomposition, cofactor-native
Hessian, chart covariance, two-read reduction, and robust finite verdict are
\(\pass\) at their exact or conditional scopes.  This is the upstream
noncircular reduction: the programme should no longer treat the numerical
coefficient or the full mixed block as the first target.  It must populate the
two root cofactors and then the one assembled homogeneous Wilson mode card.
Those physical values remain absent, so the selected symbol remains
\(\stopstatus\).  No negative physical mode has been produced.
\end{assessment}

\section*{Version V14 append-only continuation record}
\addcontentsline{toc}{section}{Version V14 append-only continuation record}
The complete Version~V13 source above is preserved byte for byte up to its
sole terminal document-closing command.  Its SHA--256 digest is
\begin{center}\scriptsize\ttfamily
f8d436402e430c409385b4138b1b7b53b0cc5f40b42282ad4a8a78a94c324ee4.
\end{center}
For Version~V15, preserve this accumulated source, append new mathematical
material immediately before the final document-closing command, and increment
the filename to \texttt{\_V15.tex}.  Keep the physical response target,
weak-facing root, loop family, unnormalized action, selector tangent,
plaquette marginal writer and its physical normalization, generated
interaction, rooted Hodge carrier, selected axial polarization, colour
normalization, and local/topological split fixed before reading any cofactor,
response angle, final-law writer, covariance, margin, ancestry, or transport
value.

% =============================================================================
% END APPENDED VERSION V14 MATERIAL
% =============================================================================

% =============================================================================
% START APPENDED VERSION V15 MATERIAL
% =============================================================================

\clearpage
\hypersetup{
 pdftitle={Renewal Yang--Mills Reduced-Fibre Wilson Shell Symbol and Local-Vacuum Programme V15},
 pdfsubject={Finite response-level continuation, tangent-endpoint separation, and an endpoint-safe Wilson mode card}
}
\pdfinfo{/Title (Renewal Yang--Mills Reduced-Fibre Wilson Shell Symbol and Local-Vacuum Programme V15)
 /Subject (Finite response-level continuation, tangent-endpoint separation, and an endpoint-safe Wilson mode card)}

\providecommand{\Cum}{\operatorname{Cum}}

\section{Version V15: response-level continuation and the endpoint Wilson card}
\label{sec:YMA15-entry}

Version~V14 is preserved above.  Its projective response coordinate removes a
spurious dependence on the units used for the plaquette writer, and its
cofactor-weighted Hessian is an exact algebraic identity for any independently
loaded affine law.  The first point of the present continuation is upstream of
that identity.  The two cofactors at one selected root determine the
\emph{tangent} to the response level.  They do not, by themselves, determine
the endpoint of a finite generated-action insertion.

Version~V11 proved by the implicit-function theorem that
\(b^{\rm marg}\) is the first derivative of a response-level graph.  Its own
statement already distinguished that first jet from the second response jet
needed for a finite deformation.  Versions~V12--V14 sharpened the writer,
coordinate, and projective meanings of this derivative, but subsequently used
its value as the coefficient of an affine endpoint law.  That substitution is
licensed only on one of two branches:
\begin{enumerate}[label=\textup{(S15.\arabic*)},leftmargin=5.2em]
\item the submitted \(\Psi\) is an infinitesimal tangent and the intended
      output is a derivative at the selected root; or
\item an independent finite continuation theorem proves that the affine
      endpoint lies on the same response level.
\end{enumerate}
For a finite generated shell increment neither branch is automatic.

The corrected acquisition order is
\begin{equation}
 \boxed{
 \begin{gathered}
 \text{physical two-parameter unnormalized action family and generated amplitude}
 \\
 \Downarrow\\[-0.2em]
 \text{connected component of the logarithm-free response level}
 \\
 \Downarrow\\[-0.2em]
 \text{finite endpoint coefficient or an attained fold/singularity}
 \\
 \Downarrow\\[-0.2em]
 \text{one unnormalized two-copy endpoint Hessian card}
 \\
 \Downarrow\\[-0.2em]
 \text{physical axial margin or an attained transverse witness}.
 \end{gathered}}
 \label{eq:YMA15-corrected-order}
\end{equation}
The root cofactors remain essential, but their correct geometric role is as the
velocity field of the response-level curve.  The endpoint law is obtained by
integrating that field or, equivalently, by validated root continuation at the
physical generated amplitude.

\begin{firewall}[No retrospective change of the submitted action]
\label{fire:YMA15-no-action-change}
The weak-facing response, loop family, physical plaquette normalization,
generated writer, unnormalized action, rooted Hodge carrier, selected axial
polarization, colour normalization, and local/topological convention remain
fixed.  The continuation below does not replace the response level, choose a
new marginal direction, or fit a path after observing the mode sign.  It only
distinguishes the local tangent of the already declared level from its finite
endpoint.
\end{firewall}

% =============================================================================
\section{The cofactor wedge is the velocity field of the response level}
\label{sec:YMA15-level-flow}
% =============================================================================

Let \(\alpha\) denote the coefficient of the physically normalized plaquette
writer and let \(s\) denote the strength of the submitted generated insertion.
On one fixed finite reduced fibre consider the affine action plane
\begin{equation}
 \Phi_{\alpha,s}=\Phi_*+\alpha P+s\Psi,
 \qquad (\alpha,s)\in U\subset\mathbb R^2,
 \label{eq:YMA15-action-plane}
\end{equation}
where \(U\) is an open neighborhood of the selected root \((0,0)\).  Let
\begin{equation}
 D(\alpha,s):=D_u[\Phi_{\alpha,s}]
 \label{eq:YMA15-root-function}
\end{equation}
be the logarithm-free weak-facing root function of Version~V10.  Thus
\(D(0,0)=0\).  Write
\begin{equation}
 c_P(\alpha,s):=\partial_\alpha D(\alpha,s),
 \qquad
 c_\Psi(\alpha,s):=\partial_s D(\alpha,s).
 \label{eq:YMA15-moving-cofactors}
\end{equation}
At the origin these are precisely the two cofactors
\(\mathfrak c_P,\mathfrak c_\Psi\) of Versions~V13--V14.

\begin{theorem}[Hamiltonian cofactor flow of the exact response level]
\label{thm:YMA15-Hamiltonian-level-flow}
Assume \(D\in C^2(U)\) and define the planar vector field
\begin{equation}
 \boxed{
 \mathcal X_D(\alpha,s)
 :=\bigl(-c_\Psi(\alpha,s),\ c_P(\alpha,s)\bigr).}
 \label{eq:YMA15-Hamiltonian-field}
\end{equation}
Then:
\begin{enumerate}[label=\textup{(H\arabic*)},leftmargin=3.8em]
\item every integral curve \(z(\tau)=(\alpha(\tau),s(\tau))\) of
      \(\mathcal X_D\) satisfies
      \begin{equation}
       \boxed{D(z(\tau))=\text{constant};}
       \label{eq:YMA15-D-constant}
      \end{equation}
\item at every regular point \(\nabla D\ne0\), the vector
      \(\mathcal X_D\ne0\) spans the tangent line of the level set of \(D\);
\item along such a curve, the action velocity is exactly the V14 cofactor
      interaction,
      \begin{equation}
       \boxed{
       \frac d{d\tau}\Phi_{\alpha(\tau),s(\tau)}
       =c_P\Psi-c_\Psi P
       =:\mathcal I_{\alpha,s}^{\rm cof};}
       \label{eq:YMA15-action-velocity}
      \end{equation}
\item on a patch where \(c_P\ne0\), the generated amplitude itself may be
      used as parameter and the exact response-preserving graph obeys
      \begin{equation}
       \boxed{
       \frac{d\alpha}{ds}
       =-\frac{c_\Psi(\alpha(s),s)}{c_P(\alpha(s),s)}.}
       \label{eq:YMA15-exact-coefficient-ODE}
      \end{equation}
      In particular the V14 coefficient is only the initial velocity
      \begin{equation}
       \boxed{
       b_0^{\rm marg}
       =\alpha'(0)
       =-\frac{\mathfrak c_\Psi}{\mathfrak c_P}.}
       \label{eq:YMA15-initial-velocity}
      \end{equation}
\end{enumerate}
\end{theorem}

\begin{proof}
The chain rule gives
\[
 \frac d{d\tau}D(z(\tau))
 =D_\alpha(-D_s)+D_sD_\alpha=0,
\]
which proves the first item.  At a regular point the tangent space of the level
set is the one-dimensional orthogonal complement of \(\nabla D\).  The vector
\((-D_s,D_\alpha)\) is nonzero and belongs to that complement, proving the
second item.  Differentiating \eqref{eq:YMA15-action-plane} along the same
curve gives
\[
 \dot\alpha P+\dot s\Psi=-D_sP+D_\alpha\Psi,
\]
which is \eqref{eq:YMA15-action-velocity}.  If \(D_\alpha\ne0\), then
\(\dot s=D_\alpha\ne0\) and division of \(\dot\alpha=-D_s\) by
\(\dot s=D_\alpha\) gives \eqref{eq:YMA15-exact-coefficient-ODE}.  Evaluation
at the selected root proves \eqref{eq:YMA15-initial-velocity}.
\end{proof}

\begin{corollary}[The finite marginal coefficient is a path integral]
\label{cor:YMA15-coefficient-path-integral}
Suppose the response-level graph exists on \(0\le s\le\sigma\), with
\(\alpha(0)=0\) and \(c_P(\alpha(s),s)\ne0\).  Then the finite endpoint
coefficient is
\begin{equation}
 \boxed{
 \alpha(\sigma)
 =-\int_0^\sigma
   \frac{c_\Psi(\alpha(t),t)}{c_P(\alpha(t),t)}\,dt.}
 \label{eq:YMA15-endpoint-integral}
\end{equation}
Consequently
\begin{equation}
 \boxed{
 \alpha(\sigma)-\sigma b_0^{\rm marg}
 =\int_0^\sigma
   \left[
    -\frac{c_\Psi(\alpha(t),t)}{c_P(\alpha(t),t)}
    -b_0^{\rm marg}
   \right]dt.}
 \label{eq:YMA15-endpoint-minus-tangent}
\end{equation}
The equality \(\alpha(\sigma)=\sigma b_0^{\rm marg}\) is therefore an
additional finite-level statement, not a consequence of the initial
cofactors.
\end{corollary}

\begin{proof}
Integrate \eqref{eq:YMA15-exact-coefficient-ODE} and subtract the integral of
the constant initial slope.
\end{proof}

\begin{assessment}[Correct interpretation of the V14 projective interaction]
\label{ass:YMA15-V14-wedge-qualification}
The interaction
\(\mathfrak c_P\Psi-\mathfrak c_\Psi P\) constructed in Version~V14 is the
unnormalized velocity of the response level at the selected root.  It is not,
without \eqref{eq:YMA15-endpoint-integral}, the finite action increment from
\(s=0\) to a unit generated insertion.  The affine-covariance and
cofactor-Hessian identities of Version~V14 remain algebraically exact for any
independently loaded affine law.  What is withdrawn is only the inference that
the initial cofactor ratio itself loads the finite response-matched endpoint.
\end{assessment}

% =============================================================================
\section{The first nonlinear obstruction to the affine tangent}
\label{sec:YMA15-linear-residual}
% =============================================================================

Assume \(\mathfrak c_P\ne0\) and put
\begin{equation}
 b_0=-\frac{\mathfrak c_\Psi}{\mathfrak c_P},
 \qquad
 v_0=(b_0,1),
 \qquad
 \widetilde\alpha(s)=b_0s.
 \label{eq:YMA15-linear-predictor}
\end{equation}
The response defect of the affine tangent is
\begin{equation}
 \mathscr E_{\rm lin}(s):=D(b_0s,s).
 \label{eq:YMA15-linear-response-defect}
\end{equation}

\begin{theorem}[Exact curvature residual of the affine tangent]
\label{thm:YMA15-linear-residual}
For every \(s\) for which the line segment
\(\{(b_0t,t):0\le t\le s\}\) lies in \(U\),
\begin{align}
 \mathscr E_{\rm lin}(0)&=0,
 &\mathscr E_{\rm lin}'(0)&=0,
 \label{eq:YMA15-linear-first-zero}\\
 \mathscr E_{\rm lin}(s)
 &=\boxed{\int_0^s(s-t)\,
 v_0^*\bigl(\operatorname{Hess}D\bigr)(b_0t,t)v_0\,dt.}
 \label{eq:YMA15-linear-Peano}
\end{align}
In particular,
\begin{equation}
 \boxed{
 \mathscr E_{\rm lin}''(0)
 =D_{ss}+2b_0D_{\alpha s}+b_0^2D_{\alpha\alpha}
 \quad\text{at }(0,0).}
 \label{eq:YMA15-level-curvature}
\end{equation}
Let
\begin{equation}
 X_0:=\mathfrak c_P\Psi-\mathfrak c_\Psi P.
 \label{eq:YMA15-X0}
\end{equation}
Then the homogeneous second cofactor is
\begin{equation}
 \boxed{
 \left.\frac{d^2}{d\tau^2}
 D_u[\Phi_*+\tau X_0]\right|_{\tau=0}
 =\mathfrak c_P^2\mathscr E_{\rm lin}''(0).}
 \label{eq:YMA15-homogeneous-curvature}
\end{equation}
If this number is nonzero, the affine tangent leaves the response level at
quadratic order.

If \(D\) is real analytic, the affine line is contained in the response level
in a neighborhood of the origin if and only if
\begin{equation}
 \boxed{
 (v_0\cdot\nabla)^mD(0,0)=0
 \quad\text{for every }m\ge2.}
 \label{eq:YMA15-straight-line-criterion}
\end{equation}
For a polynomial root function of degree \(r\), only the finitely many orders
\(2\le m\le r\) are required.
\end{theorem}

\begin{proof}
The first identity is the root condition.  The derivative at zero is
\[
 \mathfrak c_Pb_0+\mathfrak c_\Psi=0.
\]
Twice differentiating the one-variable function
\(t\mapsto D(b_0t,t)\) gives the integrand in
\eqref{eq:YMA15-linear-Peano}; integrating twice with zero initial value and
zero initial derivative proves that formula and
\eqref{eq:YMA15-level-curvature}.  The parameter direction associated with
\(X_0\) is
\((-\mathfrak c_\Psi,\mathfrak c_P)=\mathfrak c_Pv_0\), proving
\eqref{eq:YMA15-homogeneous-curvature}.

For analytic \(D\), the restriction \(s\mapsto D(b_0s,s)\) is analytic.  It
vanishes on a neighborhood exactly when all of its Taylor coefficients
vanish.  The orders zero and one already vanish, giving
\eqref{eq:YMA15-straight-line-criterion}.  A polynomial restriction has no
coefficients above its degree.
\end{proof}

The level-curvature scalar is itself an exact unnormalized insertion card.
For a loop label \(\ell\in\{11,00,10,01\}\), write
\begin{equation}
 N_{\ell;X^r}
 :=\int\mathcal O_\ell X^r e^{-\Phi_*}\,dm_{\rm red},
 \qquad r=0,1,2,
 \label{eq:YMA15-marked-moments}
\end{equation}
with \(N_{\ell;X^0}=N_\ell\), and define
\begin{equation}
 \mathfrak K_{\ell_1,\ell_2}[X]
 :=N_{\ell_1;X^2}N_{\ell_2}
   +2N_{\ell_1;X}N_{\ell_2;X}
   +N_{\ell_1}N_{\ell_2;X^2}.
 \label{eq:YMA15-pair-second-writer}
\end{equation}

\begin{proposition}[Exact pair-marked second cofactor]
\label{prop:YMA15-second-cofactor}
For every fixed insertion \(X\),
\begin{equation}
 \boxed{
 \left.\frac{d^2}{dt^2}D_u[\Phi_*+tX]\right|_{t=0}
 =\mathfrak K_{11,00}[X]
  -e^{-u}\mathfrak K_{10,01}[X].}
 \label{eq:YMA15-second-cofactor}
\end{equation}
Thus the first nonlinear test of the V14 affine tangent is obtained from the
same finite Haar grammar as the first cofactors, with at most two copies and
two marked insertions.
\end{proposition}

\begin{proof}
For an affine action perturbation,
\[
 \frac d{dt}N_\ell(t)=-N_{\ell;X}(t),
 \qquad
 \frac{d^2}{dt^2}N_\ell(t)=N_{\ell;X^2}(t).
\]
The product rule therefore gives
\[
 \frac{d^2}{dt^2}(N_{\ell_1}N_{\ell_2})
 =N_{\ell_1;X^2}N_{\ell_2}
  +2N_{\ell_1;X}N_{\ell_2;X}
  +N_{\ell_1}N_{\ell_2;X^2}.
\]
Apply this to
\(D_u=N_{11}N_{00}-e^{-u}N_{10}N_{01}\).
\end{proof}

\begin{countertheorem}[Initial cofactors do not determine a finite endpoint or its sign]
\label{cth:YMA15-initial-cofactors-no-endpoint}
For \(\lambda\in\mathbb R\), consider the positive analytic loop-numerator
germ
\begin{equation}
 N_{11}^{(\lambda)}(\alpha,s)
 =\exp[-\chi_\lambda(\alpha,s)],
 \qquad
 N_{00}^{(\lambda)}=N_{10}^{(\lambda)}=N_{01}^{(\lambda)}=1,
 \label{eq:YMA15-counter-numerators}
\end{equation}
where
\begin{equation}
 \chi_\lambda(\alpha,s)=u+\alpha+s+\lambda s^2.
 \label{eq:YMA15-counter-response}
\end{equation}
Every member has the same selected root, the same two initial cofactors up to
the common nonzero factor \(-e^{-u}\), and the same initial marginal slope
\(b_0=-1\).  Its exact response-level graph is, however,
\begin{equation}
 \boxed{\alpha_\lambda(s)=-s-\lambda s^2.}
 \label{eq:YMA15-counter-graph}
\end{equation}
At unit generated amplitude the endpoint is \(-1-\lambda\).

Attach the coarse-only local curvature writer
\begin{equation}
 C_\alpha(w)=\frac12\left(\alpha+\frac32\right)w^2,
 \label{eq:YMA15-counter-curvature}
\end{equation}
which vanishes at the response background \(w=0\) and hence does not alter the
loop root.  Its field Hessian at the endpoint is
\begin{equation}
 M_\lambda=\frac12-\lambda.
 \label{eq:YMA15-counter-margin}
\end{equation}
Thus \(\lambda=-1\) gives a strictly positive endpoint and \(\lambda=1\)
gives a strictly negative endpoint although all initial cofactor data agree.
\end{countertheorem}

\begin{proof}
The four numerator germs are positive, and their Creutz response is exactly
\(\chi_\lambda\).  Setting it equal to \(u\) gives
\eqref{eq:YMA15-counter-graph}.  At the origin both first partial derivatives
of \(\chi_\lambda\) equal one, independently of \(\lambda\), so the root
cofactor ratio is \(-1\).  Substitution of the unit-amplitude endpoint into the
second derivative of \eqref{eq:YMA15-counter-curvature} gives
\[
 \alpha_\lambda(1)+\frac32
 =-1-\lambda+\frac32
 =\frac12-\lambda.
\]
The coarse-only term is zero at \(w=0\), so it leaves all four response
numerators at that background unchanged.  This is exactly the absolute
curvature freedom isolated in Version~V8.
\end{proof}

% =============================================================================
\section{Finite continuation, validated endpoint transport, and folds}
\label{sec:YMA15-finite-continuation}
% =============================================================================

The preceding counterexample shows that a finite endpoint must be followed on
the actual response level.  The next theorem gives a simple global certificate
which is stronger than a local implicit-function statement and is directly
suited to outward interval arithmetic.

\begin{theorem}[Monotone response tube and unique finite graph]
\label{thm:YMA15-monotone-tube}
Let
\begin{equation}
 \mathcal R=[a_-,a_+]\times[0,\sigma]
 \subset U
 \label{eq:YMA15-response-tube}
\end{equation}
be compact.  Suppose, after replacing \(D\) by \(-D\) if necessary, that
\begin{equation}
 \boxed{
 \partial_\alpha D(\alpha,s)\ge m>0
 \quad\text{on }\mathcal R,}
 \label{eq:YMA15-Da-floor}
\end{equation}
and that for every \(s\in[0,\sigma]\),
\begin{equation}
 D(a_-,s)\le0\le D(a_+,s).
 \label{eq:YMA15-boundary-bracket}
\end{equation}
Then there is a unique function
\(\alpha:[0,\sigma]\to[a_-,a_+]\) satisfying
\begin{equation}
 D(\alpha(s),s)=0.
 \label{eq:YMA15-global-root-graph}
\end{equation}
It is continuous, is \(C^1\) whenever \(D\in C^1\), and obeys the exact ODE
\eqref{eq:YMA15-exact-coefficient-ODE}.  If the boundary inequalities are
strict and \(D\) is analytic on a neighborhood of \(\mathcal R\), the graph is
analytic.  It is the unique response-level branch in \(\mathcal R\) issuing
from the selected root.
\end{theorem}

\begin{proof}
For fixed \(s\), the boundary bracket and the intermediate value theorem give
at least one zero.  The derivative floor makes
\(\alpha\mapsto D(\alpha,s)\) strictly increasing, so that zero is unique.
To prove continuity, let \(s_n\to s\).  Every subsequence of
\(\alpha(s_n)\) has a convergent subsubsequence in the compact interval.
Continuity of \(D\) makes its limit a zero at \(s\), and uniqueness forces the
limit to be \(\alpha(s)\).  Hence the whole sequence converges.  The implicit-
function theorem and the nonzero derivative give local \(C^1\), respectively
analytic, graphs; uniqueness glues them to the global graph.  Differentiating
\eqref{eq:YMA15-global-root-graph} gives the ODE.  Since every zero in the tube
is on this graph, no second branch in the tube can issue from the initial root.
\end{proof}

\begin{corollary}[Explicit quadratic error of the initial tangent]
\label{cor:YMA15-tangent-error}
Under the hypotheses of \Cref{thm:YMA15-monotone-tube}, assume additionally
that on \(\mathcal R\)
\begin{equation}
 |D_s|\le n,
 \qquad
 |D_{\alpha\alpha}|,
 |D_{\alpha s}|,
 |D_{ss}|\le M.
 \label{eq:YMA15-second-derivative-bounds}
\end{equation}
Put
\begin{equation}
 B:=\frac nm,
 \qquad
 L:=\frac Mm+\frac{nM}{m^2},
 \qquad
 b_0=-\frac{D_s(0,0)}{D_\alpha(0,0)}.
 \label{eq:YMA15-BL}
\end{equation}
If \(\alpha(0)=0\), then for \(0\le s\le\sigma\),
\begin{equation}
 \boxed{
 |\alpha(s)-b_0s|
 \le\frac12L(B+1)s^2.}
 \label{eq:YMA15-tangent-error}
\end{equation}
\end{corollary}

\begin{proof}
Let
\(q(\alpha,s)=-D_s/D_\alpha\).  The derivative floor and
\eqref{eq:YMA15-second-derivative-bounds} give
\[
 |q|\le B,
 \qquad
 |\partial_\alpha q|,|\partial_sq|
 \le \frac Mm+\frac{nM}{m^2}=L.
\]
The graph equation gives \(|\alpha(t)|\le Bt\).  Hence
\[
 |q(\alpha(t),t)-q(0,0)|
 \le L(|\alpha(t)|+t)
 \le L(B+1)t.
\]
Integrating
\(\alpha(s)-b_0s=\int_0^s[q(\alpha(t),t)-q(0,0)]dt\)
proves the claim.
\end{proof}

\begin{proposition}[Endpoint residual correction]
\label{prop:YMA15-endpoint-residual}
Fix \(s=\sigma\) and suppose
\begin{equation}
 |D_\alpha(\alpha,\sigma)|\ge m>0
 \label{eq:YMA15-endpoint-slope-floor}
\end{equation}
with constant sign on an interval \(I\) containing both the unique root
\(\alpha_*\) and a candidate \(\widetilde\alpha\).  Then
\begin{equation}
 \boxed{
 |\alpha_*-\widetilde\alpha|
 \le\frac{|D(\widetilde\alpha,\sigma)|}{m}.}
 \label{eq:YMA15-residual-root-error}
\end{equation}
Thus a certified residual interval
\(|D(\widetilde\alpha,\sigma)|\le\varepsilon_D\) and the slope floor give an
endpoint interval of radius \(\varepsilon_D/m\).  The tangent predictor
\(\widetilde\alpha=\sigma b_0\) is admissible only after this held-out response
residual is read.
\end{proposition}

\begin{proof}
The one-variable mean-value theorem gives
\[
 D(\widetilde\alpha,\sigma)-D(\alpha_*,\sigma)
 =D_\alpha(\xi,\sigma)(\widetilde\alpha-\alpha_*)
\]
for some \(\xi\) between the two points.  The root term vanishes and the
absolute derivative is at least \(m\).
\end{proof}

\begin{theorem}[Generic finite marginal fold]
\label{thm:YMA15-fold}
Let \((\alpha_0,s_0)\) satisfy
\begin{equation}
 D(\alpha_0,s_0)=0,
 \qquad
 D_\alpha(\alpha_0,s_0)=0,
 \qquad
 D_s(\alpha_0,s_0)\ne0.
 \label{eq:YMA15-fold-hypotheses}
\end{equation}
Then the response level is locally a unique graph \(s=g(\alpha)\), with
\begin{equation}
 g'(\alpha_0)=0.
 \label{eq:YMA15-fold-first}
\end{equation}
If in addition \(D_{\alpha\alpha}(\alpha_0,s_0)\ne0\), then
\begin{equation}
 \boxed{
 g(\alpha)-s_0
 =-\frac{D_{\alpha\alpha}}{2D_s}
  (\alpha-\alpha_0)^2
 +O(|\alpha-\alpha_0|^3).}
 \label{eq:YMA15-fold-normal-form}
\end{equation}
Consequently no \(C^1\) plaquette-coefficient graph \(\alpha=\alpha(s)\)
passes through the fold.  On the side where roots exist, the two local
branches have square-root separation.  This is a finite marginal-chart fold,
not a negative field mode.
\end{theorem}

\begin{proof}
Since \(D_s\ne0\), the implicit-function theorem solves the equation as
\(s=g(\alpha)\).  Differentiating
\(D(\alpha,g(\alpha))=0\) gives
\[
 D_\alpha+D_sg'=0,
\]
so \(g'(\alpha_0)=0\).  A second derivative and the vanishing of \(g'\) at
the point give
\[
 D_{\alpha\alpha}+D_sg''=0.
\]
Taylor's theorem yields \eqref{eq:YMA15-fold-normal-form}.  When the quadratic
coefficient is nonzero, inversion has the usual square-root form and cannot be
\(C^1\) through \(s_0\).
\end{proof}

\begin{firewall}[Endpoint root sets do not select a physical branch]
\label{fire:YMA15-branch-selection}
If several response-level roots occur at the final generated amplitude, the
set of endpoints does not determine which one continues the selected root.
The physical branch is the connected continuation fixed before the final mode
is read.  A fold, exit from the certified tube, or singular point
\(D_\alpha=D_s=0\) is retained as the corresponding response-transport
failure.  It is not repaired by choosing whichever endpoint gives the desired
Hessian sign.
\end{firewall}

% =============================================================================
\section{Exact Hessian transport along the response level}
\label{sec:YMA15-Hessian-flow}
% =============================================================================

The V4 generated-action flow is exact for its straight interpolation.  Along
the actual response-level graph the interaction tangent changes with \(s\).
The same differentiation identity yields the following endpoint-safe flow.

Let
\begin{equation}
 \Phi_s:=\Phi_{\alpha(s),s},
 \qquad
 A_s:=\dot\Phi_s=\Psi+\alpha'(s)P.
 \label{eq:YMA15-moving-action-tangent}
\end{equation}
For the fixed retained field bank, let \(j_s\) and \(K_s\) be the complete
first and second field jets of \(\Phi_s\), and put
\begin{equation}
 j_{A,s}:=D_VA_s,
 \qquad
 K_{A,s}:=D_V^2A_s.
 \label{eq:YMA15-moving-tangent-jets}
\end{equation}
All expectations and cumulants in this section are taken in the same law
\(d\mu_s=Z_s^{-1}e^{-\Phi_s}dm_{\rm red}\).  On the nonzero axial bank the
translation and colour symmetries give
\(\mathbb E_sj_s=\mathbb E_sj_{A,s}=0\).

\begin{theorem}[Response-preserving Hessian flow]
\label{thm:YMA15-response-Hessian-flow}
The exact physical Hessian
\begin{equation}
 H_s=\mathbb E_sK_s-\operatorname{Cov}_s(j_s)
 \label{eq:YMA15-Hs}
\end{equation}
satisfies
\begin{equation}
 \boxed{
 \begin{aligned}
 \dot H_s={}&\mathbb E_sK_{A,s}
  -\operatorname{Cov}_s(K_s,A_s)\\
 &-\operatorname{Cov}_s(j_{A,s},j_s)
  -\operatorname{Cov}_s(j_s,j_{A,s})\\
 &+\Cum_s(j_s,j_s^*,A_s).
 \end{aligned}}
 \label{eq:YMA15-Hessian-flow-card}
\end{equation}
Here
\begin{equation}
 \Cum_s(j_s,j_s^*,A_s)
 :=\mathbb E_s[j_sj_s^*(A_s-\mathbb E_sA_s)]
 \label{eq:YMA15-third-cumulant}
\end{equation}
on the zero-mean bank.  Consequently
\begin{equation}
 \boxed{
 H_\sigma-H_0
 =\int_0^\sigma\dot H_s\,ds.}
 \label{eq:YMA15-Hessian-endpoint-integral}
\end{equation}
If the topological and reference forms vary along the response level, their
actual endpoint change is subtracted by integrating
\(\dot T_s+\delta\dot S_s\) on the same path.
\end{theorem}

\begin{proof}
For any differentiable observable \(X_s\) on a finite compact cylinder,
\begin{equation}
 \frac d{ds}\mathbb E_sX_s
 =\mathbb E_s\dot X_s
  -\operatorname{Cov}_s(X_s,A_s).
 \label{eq:YMA15-expectation-derivative}
\end{equation}
Apply this to \(K_s\), noting \(\dot K_s=K_{A,s}\).  For the covariance,
differentiate
\(\mathbb E_s[j_sj_s^*]\).  The explicit derivative gives the two mixed
covariances, while the change of law gives
\(-\operatorname{Cov}_s(j_sj_s^*,A_s)\).  Since both score means vanish,
this last covariance is the third cumulant in
\eqref{eq:YMA15-third-cumulant}.  Subtracting the covariance derivative from
the deterministic derivative proves \eqref{eq:YMA15-Hessian-flow-card}.
Integration gives \eqref{eq:YMA15-Hessian-endpoint-integral}.
\end{proof}

\begin{corollary}[Cofactor parameterization without division]
\label{cor:YMA15-cofactor-Hessian-flow}
Parameterize the response level by the Hamiltonian time \(\tau\) of
\Cref{thm:YMA15-Hamiltonian-level-flow}.  Then
\begin{equation}
 A_\tau=\mathcal I_{\alpha(\tau),s(\tau)}^{\rm cof}
 =c_P\Psi-c_\Psi P.
 \label{eq:YMA15-cofactor-moving-tangent}
\end{equation}
Inserting this moving cofactor interaction into
\eqref{eq:YMA15-Hessian-flow-card} gives an exact denominator-free transport
of the Hessian along the response level.  The V14 projective interaction is
the initial value of this moving tangent.
\end{corollary}

\begin{proof}
This is \eqref{eq:YMA15-action-velocity} substituted into
\Cref{thm:YMA15-response-Hessian-flow}.
\end{proof}

\begin{corollary}[One two-replica writer for the moving Hessian flow]
\label{cor:YMA15-Hessian-flow-replica}
Let \(z_1,z_2\) be independent \(\mu_s\)-samples and write
\(\Delta X=X(z_1)-X(z_2)\).  Then the matrix in
\eqref{eq:YMA15-Hessian-flow-card} is the expectation of the single signed
writer
\begin{equation}
 \boxed{
 \begin{aligned}
 \mathscr G_s={}&
 \frac12(K_{A,s}^{(1)}+K_{A,s}^{(2)})
 -\frac12\Delta K_s\,\Delta A_s\\
 &-\frac12\left(
   \Delta j_{A,s}\Delta j_s^*
  +\Delta j_s\Delta j_{A,s}^*\right)\\
 &+\frac12\Delta(j_sj_s^*)\,\Delta A_s.
 \end{aligned}}
 \label{eq:YMA15-moving-replica-writer}
\end{equation}
That is,
\begin{equation}
 \boxed{\dot H_s=\mathbb E_s^{\otimes2}\mathscr G_s.}
 \label{eq:YMA15-moving-replica-expectation}
\end{equation}
All cancellations are retained on one same-law pair cylinder before an
interval is taken.
\end{corollary}

\begin{proof}
For arbitrary square-integrable variables,
\[
 \frac12\mathbb E^{\otimes2}[\Delta X\,\Delta Y]
 =\operatorname{Cov}(X,Y).
\]
Apply this identity to the second and third lines of
\eqref{eq:YMA15-moving-replica-writer}.  The final line equals
\(\operatorname{Cov}(j_sj_s^*,A_s)\), which is the third cumulant because
\(\mathbb E_sj_s=0\).  The first line gives \(\mathbb E_sK_{A,s}\).
\end{proof}

\begin{assessment}[Scope of the V4 straight interpolation]
\label{ass:YMA15-V4-qualification}
The V4 identity for
\(\Phi_\theta=\beta V+\theta\Psi\) remains exact for the two endpoints of
that declared straight interpolation.  It does not assert that the physical
response-matched trajectory has fixed plaquette coefficient.  When response
matching is part of the shell definition, its endpoint change is given by
\Cref{thm:YMA15-response-Hessian-flow} with
\(A_s=\Psi+\alpha'(s)P\).  The two flows coincide only when the independently
loaded physical path is the V4 straight path or when the plaquette correction
vanishes identically.
\end{assessment}

% =============================================================================
\section{One unnormalized endpoint cylinder decides the physical mode}
\label{sec:YMA15-endpoint-cylinder}
% =============================================================================

Once the finite response-level endpoint has been certified, no normalized law
or separate deterministic/covariance acquisition is required for the sign of
the selected mode.  Let
\begin{equation}
 \Phi_{\rm end}=\Phi_*+\sigma\Psi+\alpha(\sigma)P,
 \qquad
 d\Lambda(z)=e^{-\Phi_{\rm end}(z)}dm_{\rm red}(z),
 \qquad
 Z=\Lambda(Y)>0.
 \label{eq:YMA15-endpoint-measure}
\end{equation}
Let \(j\) and \(K\) be the complete first and second field jets of this same
action on the selected supported bank.  Let \(T\) be the endpoint topological
form and \(S\) the endpoint reference form.

\begin{theorem}[Single unnormalized two-copy physical margin]
\label{thm:YMA15-unnormalized-endpoint-card}
For \(z_1,z_2\in Y\), define the Hermitian matrix writer
\begin{equation}
 \boxed{
 \mathscr M_\delta(z_1,z_2)
 :=\frac12(K(z_1)+K(z_2))
  -\frac12\Delta j\,\Delta j^*
  -(T+\delta S).}
 \label{eq:YMA15-endpoint-writer}
\end{equation}
Then
\begin{equation}
 \boxed{
 \mathbb N_\delta
 :=\int_{Y^2}\mathscr M_\delta(z_1,z_2)
    \,d\Lambda(z_1)d\Lambda(z_2)
 =Z^2(H_{\rm end}-T-\delta S).}
 \label{eq:YMA15-unnormalized-margin-identity}
\end{equation}
Consequently \(\mathbb N_\delta\) and the physical endpoint margin have the
same inertia, kernel, and witness vectors.  The sign is decided by one
unnormalized same-history two-copy cylinder, including the direct second
action insertion and the complete score covariance.
\end{theorem}

\begin{proof}
Division of the product measure by \(Z^2\) gives two independent samples from
the normalized endpoint law.  The first term in
\eqref{eq:YMA15-endpoint-writer} has expectation \(\mathbb E K\).  The pair-
difference identity gives
\[
 \frac12\mathbb E^{\otimes2}[\Delta j\Delta j^*]
 =\operatorname{Cov}(j).
\]
Thus the normalized expectation of the writer is
\(H_{\rm end}-T-\delta S\).  Multiplying by \(Z^2>0\) proves the identity and
the inertia statement.
\end{proof}

\begin{corollary}[Endpoint PASS, obstruction, and quantitative floor]
\label{cor:YMA15-endpoint-verdict}
Suppose a Hermitian approximation \(\widehat{\mathbb N}_\delta\) satisfies
\begin{equation}
 \norm{\mathbb N_\delta-\widehat{\mathbb N}_\delta}\le\varepsilon_N.
 \label{eq:YMA15-N-error}
\end{equation}
Then
\begin{equation}
 \boxed{
 \begin{array}{rcl}
 \lambda_{\min}(\widehat{\mathbb N}_\delta)-\varepsilon_N>0
 &\Longrightarrow&\pass,\\[0.25em]
 \lambda_{\min}(\widehat{\mathbb N}_\delta)+\varepsilon_N<0
 &\Longrightarrow&\obstruction,\\[0.25em]
 \text{otherwise}&\Longrightarrow&\stopstatus.
 \end{array}}
 \label{eq:YMA15-endpoint-verdict}
\end{equation}
On the negative branch, a unit minimum eigenvector of
\(\widehat{\mathbb N}_\delta\) is an attained physical witness.  If additionally
\(Z\le Z_+\) and
\(\mathbb N_\delta\succeq\gamma I\), then
\begin{equation}
 \boxed{H_{\rm end}-T-\delta S\succeq\gamma Z_+^{-2}I.}
 \label{eq:YMA15-endpoint-quantitative-floor}
\end{equation}
\end{corollary}

\begin{proof}
Weyl's inequality gives the sign alternatives.  On a minimum eigenvector of
the approximation, the upper error bound transfers strict negativity to the
exact numerator.  The preceding theorem transfers the same witness to the
physical margin.  For the quantitative floor, divide
\(\mathbb N_\delta=Z^2M\succeq\gamma I\) by
\(Z^2\le Z_+^2\).
\end{proof}

The endpoint coefficient itself may be enclosed rather than known exactly.
For fixed \(s=\sigma\), write
\begin{equation}
 M(\alpha)
 :=H_{\alpha,\sigma}-T_{\alpha,\sigma}
   -\delta S_{\alpha,\sigma}.
 \label{eq:YMA15-Malpha}
\end{equation}
The derivative of the Hessian with respect to \(\alpha\) is the special case
of \eqref{eq:YMA15-Hessian-flow-card} with action tangent \(P\).  Denote that
matrix by \(\mathfrak G_P(\alpha)\).

\begin{theorem}[Response residual transported into the physical margin]
\label{thm:YMA15-residual-to-margin}
Suppose on an interval \(I\) containing the exact endpoint \(\alpha_*\) and a
candidate \(\widetilde\alpha\) one has
\begin{equation}
 |D_\alpha(\alpha,\sigma)|\ge m>0
 \label{eq:YMA15-residual-margin-slope}
\end{equation}
and
\begin{equation}
 \norm{\mathfrak G_P(\alpha)
       -\partial_\alpha T_{\alpha,\sigma}
       -\delta\partial_\alpha S_{\alpha,\sigma}}
 \le L_M.
 \label{eq:YMA15-margin-Lipschitz}
\end{equation}
If
\(|D(\widetilde\alpha,\sigma)|\le\varepsilon_D\), then
\begin{equation}
 \boxed{
 \norm{M(\alpha_*)-M(\widetilde\alpha)}
 \le L_M\frac{\varepsilon_D}{m}.}
 \label{eq:YMA15-response-to-margin-error}
\end{equation}
Thus an endpoint-card numerical radius \(\varepsilon_{\rm card}\) at
\(\widetilde\alpha\) becomes the rigorous physical radius
\begin{equation}
 \boxed{
 \varepsilon_{\rm tot}
 =\varepsilon_{\rm card}
  +L_M\frac{\varepsilon_D}{m}.}
 \label{eq:YMA15-total-endpoint-radius}
\end{equation}
No separately fitted coefficient uncertainty is required.
\end{theorem}

\begin{proof}
\Cref{prop:YMA15-endpoint-residual} gives
\(|\alpha_*-\widetilde\alpha|\le\varepsilon_D/m\).  The fundamental theorem
of calculus for finite-dimensional Hermitian matrices gives
\[
 M(\alpha_*)-M(\widetilde\alpha)
 =\int_{\widetilde\alpha}^{\alpha_*}
   \partial_\alpha M(t)\,dt.
\]
The derivative is exactly the matrix in
\eqref{eq:YMA15-margin-Lipschitz}; its norm bound proves
\eqref{eq:YMA15-response-to-margin-error}.  The final radius is the triangle
inequality.
\end{proof}

\begin{firewall}[Absolute endpoint cylinder before downstream attribution]
\label{fire:YMA15-endpoint-before-ancestry}
The unnormalized numerator in
\eqref{eq:YMA15-unnormalized-margin-identity} is the first physical quadratic
verdict.  A Witten current, boundary message, response angle, connected tail,
or source-ancestry short may explain or transport a small endpoint margin, but
none selects the endpoint branch or reconstructs its missing absolute
second-action term.  Those rows remain downstream of the response-level
continuation and complete endpoint assembly.
\end{firewall}

% =============================================================================
\section{Tangent semantics, finite-shell semantics, and the exact alternative}
\label{sec:YMA15-semantics}
% =============================================================================

\begin{theorem}[The tangent and finite endpoint are different physical targets]
\label{thm:YMA15-two-semantics}
Let \(b_0=-\mathfrak c_\Psi/\mathfrak c_P\) on the regular initial branch.
Exactly one of the following interpretations must be declared before the
cofactor is used.
\begin{enumerate}[label=\textup{(T\arabic*)},leftmargin=3.8em]
\item \textbf{Infinitesimal shell tangent.}  The writer \(\Psi\) denotes a
      tangent at \(s=0\).  Then
      \begin{equation}
       A_0=\Psi+b_0P
       \label{eq:YMA15-infinitesimal-tangent}
      \end{equation}
      is the exact response-null action velocity, and the physical quadratic
      output is \(\dot H_0\) from
      \eqref{eq:YMA15-Hessian-flow-card}.  The finite affine law
      \(\Phi_*+\Psi+b_0P\) is not inferred.
\item \textbf{Finite generated shell.}  A physical amplitude
      \(\sigma>0\) is part of the action loading.  Then the exact endpoint is
      \begin{equation}
       \boxed{
       \Phi_{\rm end}
       =\Phi_*+\sigma\Psi+\alpha(\sigma)P,}
       \label{eq:YMA15-finite-endpoint-action}
      \end{equation}
      where \(\alpha(\sigma)\) is obtained from the connected response-level
      branch.  Replacing it by \(\sigma b_0\) is licensed only after the
      straight-level residual \eqref{eq:YMA15-linear-response-defect} has
      been proved to vanish on the whole segment, or after an endpoint
      residual/error certificate has been included.
\end{enumerate}
The V14 cofactor-Hessian identity remains exact for the affine law it names,
but the first cofactor ratio alone does not decide which of these two physical
targets that law represents.
\end{theorem}

\begin{proof}
The first branch is the differential statement
\eqref{eq:YMA15-initial-velocity} combined with
\Cref{thm:YMA15-response-Hessian-flow}.  The second branch is
\Cref{cor:YMA15-coefficient-path-integral}.  The straight-line criterion is
\Cref{thm:YMA15-linear-residual}, and the certified endpoint correction is
\Cref{prop:YMA15-endpoint-residual,thm:YMA15-residual-to-margin}.  These two
outputs are mathematically different unless the additional straightness or
endpoint residual row passes.
\end{proof}

\begin{corollary}[Finite response-level termination alternative]
\label{cor:YMA15-finite-alternative}
For a predeclared finite selected shell with a finite generated amplitude,
the ordered response/mode audit has the following exhaustive outcomes.
\begin{enumerate}[label=\textup{(V15.\arabic*)},leftmargin=5.2em]
\item \textbf{Rejected input.}  The action plane, generated amplitude,
      loop-root function, endpoint action, support, Hermiticity, Ward row, or
      local/topological split fails its defining identity.  This is not a
      theorem verdict.
\item \textbf{Regular finite continuation.}  A monotone response tube or an
      equivalent validated continuation supplies the unique endpoint
      \(\alpha(\sigma)\).  Apply
      \Cref{thm:YMA15-unnormalized-endpoint-card} and return
      \(\pass/\obstruction/\stopstatus\).
\item \textbf{Finite marginal fold.}  The connected branch reaches
      \(D_\alpha=0\), \(D_s\ne0\).  The coefficient ceases to be a regular
      function of the generated amplitude.  Return the fold and its local
      square-root branches; do not choose an endpoint by its field sign.
\item \textbf{Singular response point.}  The branch reaches
      \(D_\alpha=D_s=0\).  The least nonzero response jet or an independently
      specified continuation is required.  The current status is
      \(\stopstatus\).
\item \textbf{Unselected multiple endpoint.}  Several final roots are
      present but no connected continuation from the selected root has been
      certified.  The endpoint action is unpopulated and the status is
      \(\stopstatus\).
\end{enumerate}
Only the negative branch of the completely assembled endpoint numerator is a
Yang--Mills pure-mode obstruction.
\end{corollary}

\begin{proof}
At every regular point the level is locally a unique one-dimensional curve by
\Cref{thm:YMA15-Hamiltonian-level-flow}.  As the connected curve is continued
to the physical amplitude, either it remains in a regular graph tube, reaches
a point where only the plaquette graph fails, reaches a singular point of the
level itself, exits the validated region, or arrives in a final root set whose
physical component has not been selected.  The generic plaquette failure is
\Cref{thm:YMA15-fold}.  On the regular endpoint the sign alternatives are
\Cref{cor:YMA15-endpoint-verdict}.  These branches exhaust the finite
continuation possibilities under the declared audit.
\end{proof}

\begin{assessment}[Version V15 exact response-level replay]
\label{ass:YMA15-replay}
The standalone verifier and its exact output are
\begin{center}
\path{YMA_V15_response_endpoint_replay.py},\\
\path{YMA_V15_response_endpoint_replay.json}.
\end{center}
Their SHA--256 digests are
\begin{center}\scriptsize\ttfamily
878ba68141657c7983e6cdc2d2b345e05a941253660482321dfd5feca3981c73,\\[-0.15em]
1dcafafeb441c3cc3021b3c1af3762aa8b29b5737b801b1951df1ad1ca5ecd29.
\end{center}
The replay uses exact rational and symbolic arithmetic to verify the
Hamiltonian cofactor flow, the affine-tangent Peano residual, the pair-marked
second cofactor, the fold normal form, the complete connected Hessian-flow
formula and its two-copy writer, the two-by-two unnormalized endpoint-card
identity, and the endpoint-residual-to-margin error.  It does not evaluate a
Wilson path integral, choose a physical response semantics, locate a
weak-facing root, or certify a physical endpoint margin.
\end{assessment}

% =============================================================================
\section{Version V15 proof ledger and next literal acquisition}
\label{sec:YMA15-ledger}
% =============================================================================

\begingroup\small
\begin{longtable}{>{\raggedright\arraybackslash}p{0.27\textwidth}
                  >{\raggedright\arraybackslash}p{0.20\textwidth}
                  >{\raggedright\arraybackslash}p{0.45\textwidth}}
\caption{Version V15 proof-status ledger.}\label{tab:YMA15-ledger}\\
\toprule
\textbf{Row}&\textbf{Status}&\textbf{Exact scope}\\
\midrule
\endfirsthead
\toprule
\textbf{Row}&\textbf{Status}&\textbf{Exact scope}\\
\midrule
\endhead
Cofactor wedge&Reinterpreted exactly
&It is the Hamiltonian velocity of the logarithm-free response level, not an
unproved finite endpoint displacement.\\[0.3em]
Initial marginal ratio&Qualified
&It is \(\alpha'(0)\).  The finite coefficient is the integral of the moving
cofactor ratio along the connected level branch.\\[0.3em]
Affine-tangent residual&Proved exactly
&One second cofactor is the first nonlinear obstruction; all directional jets
characterize a straight analytic response level.\\[0.3em]
Second cofactor acquisition&Reduced exactly
&Two marked insertions in the existing pair-numerator grammar evaluate the
level curvature.\\[0.3em]
Initial-cofactor endpoint inference&Disproved
&Positive analytic numerator germs with identical initial cofactors have
opposite finite endpoint margins.\\[0.3em]
Finite response continuation&Proved conditionally
&A derivative floor and boundary bracket give a unique global response graph,
its exact ODE, and an explicit tangent error.\\[0.3em]
Finite marginal fold&Classified exactly
&A generic zero of the plaquette cofactor has the displayed square-root normal
form and is not a field-mode obstruction.\\[0.3em]
Response-preserving Hessian flow&Proved exactly
&The moving action tangent gives one four-term connected flow and one same-law
two-replica writer.\\[0.3em]
V4 straight interpolation&Qualified
&It remains exact for its declared endpoints; a response-matched endpoint uses
the moving plaquette correction.\\[0.3em]
Endpoint physical card&Reduced exactly
&One unnormalized two-copy writer equals \(Z^2\) times the complete physical
margin and preserves its inertia.\\[0.3em]
Root-to-margin error&Proved
&A held-out endpoint response residual, response-slope floor, and one Hessian
flow bound give an explicit physical error radius.\\[0.3em]
Finite endpoint continuation&Open physical value
&No supplied cylinder gives the physical generated amplitude together with a
certified response-level path and endpoint coefficient.\\[0.3em]
Selected endpoint symbol&\stopstatus
&The final unnormalized two-copy margin has not been populated; no negative
physical mode is claimed.\\[0.3em]
Local-vacuum contraction and mass&Not inferred
&No response-level continuation or finite symbol is promoted to an OS-time or
continuum mass gap.\\
\bottomrule
\end{longtable}
\endgroup

\begin{theorem}[Version V15 tangent--endpoint correction]
\label{thm:YMA15-terminal}
Versions~V1--V14 are retained.  Version~V15 additionally proves:
\begin{enumerate}[label=\textup{(E15.\arabic*)},leftmargin=5.5em]
\item the moving root cofactors generate the exact Hamiltonian flow of the
      response-level curve;
\item the V14 cofactor interaction is its initial action velocity, not a
      finite endpoint increment;
\item the first nonlinear failure of the affine tangent is an exact
      pair-marked second cofactor;
\item initial cofactors do not determine a finite response endpoint or its
      field-sign verdict;
\item a monotone response tube gives a unique global coefficient graph,
      explicit endpoint integration, and a quadratic tangent error;
\item a generic finite zero of the plaquette cofactor is a square-root fold,
      while a zero of the full response gradient is a distinct higher
      singularity;
\item the exact field Hessian has a moving response-preserving cumulant flow
      and one two-replica implementation;
\item after the endpoint is loaded, one unnormalized two-copy cylinder equals
      a positive scalar times the complete physical margin;
\item an endpoint response residual and slope floor propagate to an explicit
      physical margin error;
\item infinitesimal-shell and finite-shell interpretations are separated and
      have different target cards.
\end{enumerate}
No item evaluates the finite response-level endpoint, populates the endpoint
Wilson cylinder, proves weak-facing-root escape, supplies the differentiated
connected cofinal tail, or establishes local-vacuum contraction or a continuum
Yang--Mills mass gap.
\end{theorem}

\begin{proof}
Items~\textup{(E15.1)--(E15.2)} are
\Cref{thm:YMA15-Hamiltonian-level-flow,cor:YMA15-coefficient-path-integral,ass:YMA15-V14-wedge-qualification}.
Items~\textup{(E15.3)--(E15.4)} are
\Cref{thm:YMA15-linear-residual,prop:YMA15-second-cofactor,cth:YMA15-initial-cofactors-no-endpoint}.
Items~\textup{(E15.5)--(E15.6)} are
\Cref{thm:YMA15-monotone-tube,cor:YMA15-tangent-error,prop:YMA15-endpoint-residual,thm:YMA15-fold}.
Items~\textup{(E15.7)--(E15.9)} are
\Cref{thm:YMA15-response-Hessian-flow,cor:YMA15-Hessian-flow-replica,thm:YMA15-unnormalized-endpoint-card,thm:YMA15-residual-to-margin}.
Item~\textup{(E15.10)} is \Cref{thm:YMA15-two-semantics}.
\end{proof}

\begin{openproblem}[V16: continue the actual response level and read the endpoint mode]
\label{op:YMA15-V16}
Preserve every frozen physical choice.  Do not append another response,
cofactor, current, or boundary compiler.  Execute the following card.
\begin{enumerate}[label=\textup{(V16.\arabic*)},leftmargin=5.2em]
\item Declare whether the submitted \(\Psi\) is an infinitesimal tangent or a
      finite generated action.  On the finite branch, freeze its physical
      amplitude \(\sigma\).
\item Populate the same logarithm-free root function
      \(D(\alpha,s)\) on a connected tube issuing from the selected root.
      Certify a plaquette-cofactor floor and boundary bracket, or return the
      first fold, singular point, or tube escape.
\item Obtain the endpoint \(\alpha(\sigma)\) by the exact cofactor ODE or by
      interval root continuation.  Evaluate the held-out residual of the
      initial affine predictor; do not identify \(\alpha(\sigma)\) with
      \(\sigma b_0\) without this row.
\item Load the complete endpoint action
      \(\Phi_*+\sigma\Psi+\alpha(\sigma)P\), including the endpoint
      local/topological split and reference normalization.
\item Acquire the single unnormalized two-copy writer
      \eqref{eq:YMA15-endpoint-writer} on the first predeclared axial
      polarization and return
      \(\pass/\obstruction/\stopstatus\) by
      \eqref{eq:YMA15-endpoint-verdict}.
\item Open the V6--V7 Witten/boundary ancestry and the differentiated cofinal
      tail only after this endpoint verdict requires attribution or transport.
\end{enumerate}
If the generated amplitude or connected response-level branch is not part of
the supplied physical packet, that loading defect is now the first
\(\stopstatus\).  Initial cofactors and a later Hessian card cannot repair it.
\end{openproblem}

\begin{assessment}[V15 termination]
\label{ass:YMA15-termination}
The root-flow geometry, tangent--endpoint separation, second-cofactor test,
finite continuation theorem, fold classification, moving Hessian flow,
unnormalized endpoint card, and response-to-margin error budget are \(\pass\)
at their exact or conditional scopes.  The prior cofactor machinery is not
repeated; it is placed at its correct differential stage.  The selected finite
symbol remains \(\stopstatus\) because the physical generated amplitude,
connected response-level endpoint, and endpoint two-copy Wilson cylinder have
not been populated.  No negative physical mode has been produced.
\end{assessment}

\section*{Version V15 append-only continuation record}
\addcontentsline{toc}{section}{Version V15 append-only continuation record}
The complete Version~V14 source above is preserved byte for byte up to its
sole terminal document-closing command.  Its SHA--256 digest is
\begin{center}\scriptsize\ttfamily
512ca1d0c6ff6fd337f92f7e651ec6537f68e1eedcf27806a3c93d4f4c09edf5.
\end{center}
For Version~V16, preserve this accumulated source, append new mathematical
material immediately before the final document-closing command, and increment
the filename to \texttt{\_V16.tex}.  Keep the response target, selected root,
loop family, two-parameter unnormalized action family, physical generated
amplitude, plaquette normalization, rooted Hodge carrier, selected axial
polarization, colour normalization, and local/topological convention fixed
before reading any continuation path, endpoint, Hessian, covariance, margin,
ancestry, or transport value.

% =============================================================================
% END APPENDED VERSION V15 MATERIAL
% =============================================================================



% =============================================================================
% START APPENDED VERSION V16 MATERIAL
% =============================================================================

\clearpage
\hypersetup{
  pdftitle={Renewal Yang--Mills Reduced-Fibre Wilson Shell Symbol and Local-Vacuum Programme V16},
  pdfsubject={Restoration of the frozen tangent-matching scheme and one base-reference likelihood card for the physical Wilson symbol}
}

\section{Version V16: restoration of the frozen tangent scheme and a
base-reference likelihood card}
\label{sec:YMA16-direction}
% =============================================================================

The preceding versions deliberately went upstream whenever a later Gram,
current, or endpoint construction depended on a physical datum which had not
been loaded.  That discipline reveals one further correction.  The inherited
Yang--Mills response-matching rule does not start from a freely scalable
``generated action'' and then demand an exactly response-preserving endpoint at
unit amplitude.  Its native object is an already loaded exact shell action,
its frozen response differential, one declared plaquette marginal direction,
and the response-null \emph{interaction tangent}
\begin{equation}
 \Psi^{\rm RG}
 =\Psi^{\rm cut}+b^{\rm marg}P,
 \qquad
 b^{\rm marg}
 =-\frac{\mathscr R_{\Phi}(\Psi^{\rm cut})}
         {\mathscr R_{\Phi}(P)}.
 \label{eq:YMA16-inherited-tangent}
\end{equation}
The complete pre-retuning shell action is already finite.  The physical
matched action is obtained by adding the single frozen marginal term to that
action.  No extra amplitude parameter is part of this definition.

Version~V15 remains mathematically correct for a stronger, independently
specified experiment in which a finite generated amplitude is varied and the
response is required to remain exactly constant along a connected nonlinear
level set.  That experiment is optional.  Making it a prerequisite for the
inherited tangent-matched shell would inflate the target and place an
artificial loading row before the actual reduced-fibre Wilson symbol.

The corrected order is therefore
\begin{equation}
 \boxed{
 \begin{gathered}
 \text{loaded pre-retuning shell action and frozen response law}
 \\
 \longrightarrow
 (\mathfrak c_P,\mathfrak c_{\Psi})
 \longrightarrow
 b^{\rm marg}=-\mathfrak c_{\Psi}/\mathfrak c_P
 \\
 \longrightarrow
 \text{the exact finite action }F_b=F_0+b^{\rm marg}F_P
 \\
 \longrightarrow
 \text{one base-reference likelihood-weighted two-copy Hessian card}
 \\
 \longrightarrow
 \pass\;/\;\obstruction\;/\;\stopstatus.
 \end{gathered}}
 \label{eq:YMA16-corrected-order}
\end{equation}
No response-level path is inserted between the coefficient and the physical
field Hessian on this native branch.

% =============================================================================
\section{The native shell is tangent matched, not endpoint matched}
\label{sec:YMA16-native-semantics}
% =============================================================================

\begin{definition}[Frozen tangent-matched shell packet]
\label{def:YMA16-frozen-packet}
A finite tangent-matched shell packet consists of the following data on one
fixed reduced fibre and one fixed physical field chart.
\begin{enumerate}[label=\textup{(T\arabic*)},leftmargin=4.4em]
\item A complete real action $F_0(z,v)$ before plaquette marginal retuning.
      Every exactly generated local interaction, support birth, Jacobian, and
      declared coarse-only term is already included in $F_0$.
\item A declared plaquette marginal action $F_P(z,v)$ and its background
      writer
      \begin{equation}
       P(z):=F_P(z,0).
       \label{eq:YMA16-background-P}
      \end{equation}
\item A frozen response law and differential $\mathscr R_{\Phi}$, together
      with the exact cutoff interaction $\Psi^{\rm cut}$ whose marginal
      projection is being formed.
\item The two response cofactors
      \begin{equation}
       c:=\mathfrak c_u[P],
       \qquad
       e:=\mathfrak c_u[\Psi^{\rm cut}],
       \label{eq:YMA16-two-cofactors}
      \end{equation}
      acquired on that frozen response card.
\item On the regular branch $c\ne0$, the coefficient and final action
      \begin{equation}
       \boxed{
       b=-\frac ec,
       \qquad
       F_b=F_0+bF_P.}
       \label{eq:YMA16-final-action}
      \end{equation}
\end{enumerate}
If the response law used in item~\textup{(T3)} is not the law named by the
physical shell prescription, their occurrence or transport relation is an
input row; it is not repaired by the field Hessian.  No scalar amplitude
multiplying $\Psi^{\rm cut}$ is included unless the physical regulator
explicitly supplies such a family.
\end{definition}

\begin{theorem}[Restoration of the inherited tangent semantics]
\label{thm:YMA16-semantic-restoration}
For the packet of \Cref{def:YMA16-frozen-packet}, the response-null statement
is exactly
\begin{equation}
 \boxed{
 \mathscr R_{\Phi}
 \left(\Psi^{\rm cut}+bP\right)=0.}
 \label{eq:YMA16-response-null}
\end{equation}
The action $F_b$ in \eqref{eq:YMA16-final-action} is already one complete
finite physical action.  Its exact field Hessian is
\begin{equation}
 \boxed{
 H_b
 =\mathbb E_b\!\bigl[D_v^2F_b\bigr]
  -\Cov_b\!\bigl(D_vF_b,D_vF_b\bigr),}
 \label{eq:YMA16-native-Hessian}
\end{equation}
where the expectation is in the normalized law proportional to
$e^{-F_b(z,0)}\,\dd m(z)$.  Neither \eqref{eq:YMA16-response-null} nor
\eqref{eq:YMA16-native-Hessian} requires a finite path
$s\mapsto F_0+s\Psi^{\rm cut}+\alpha(s)F_P$ or the endpoint condition
$\chi(\alpha(s),s)=\chi(0,0)$.

The finite-level construction of Version~V15 agrees with the present packet
at first order when its initial tangent is
$\Psi^{\rm cut}+bP$.  Equality of its finite endpoint with $F_b$ is an
additional statement about the physical renormalization scheme and does not
follow from the frozen response differential.
\end{theorem}

\begin{proof}
By \eqref{eq:YMA13-cofactor-response}, the cofactors are a common nonzero
scalar multiple of the corresponding response derivatives.  Hence
$b=-e/c$ is the same coefficient as in
\eqref{eq:YMA16-inherited-tangent}, and linearity of $\mathscr R_{\Phi}$
gives \eqref{eq:YMA16-response-null}.  The action $F_0$ and marginal writer
$F_P$ are finite functions on the already fixed compact reduced fibre, so
\eqref{eq:YMA16-final-action} is a finite action without introducing a new
parameter.  Differentiating its partition function twice in the physical
field coordinate gives \eqref{eq:YMA16-native-Hessian} by the standard
score--curvature identity proved in Version~V1.

For the optional finite-level family, Version~V15 proves that the initial
cofactor ratio is the velocity of the exact response level.  Its own
counterexample \Cref{cth:YMA15-initial-cofactors-no-endpoint} shows that the
finite endpoint is not determined by that velocity.  Therefore endpoint
matching is strictly stronger than the native frozen tangent condition.
\end{proof}

\begin{firewall}[The V15 endpoint programme is optional, not false]
\label{fire:YMA16-V15-optional}
The Hamiltonian cofactor flow, monotone response tube, fold alternative,
response-preserving Hessian flow, and endpoint numerator of Version~V15 remain
valid for any physical scheme which explicitly declares a finite generated
amplitude and exact response preservation along the whole path.  Version~V16
withdraws only their use as prerequisites for the inherited tangent-matched
YM--A shell.  It does not identify an infinitesimal insertion with a finite
level endpoint, nor does it deny that a different renormalization prescription
may require the V15 packet.
\end{firewall}

\begin{corollary}[The former amplitude loading stop is removed]
\label{cor:YMA16-amplitude-stop-removed}
On the native tangent-matched branch, absence of a separately declared
physical generated amplitude is not a loading defect.  The first unresolved
rows are instead
\begin{equation}
 \boxed{
 c,\quad e,\quad
 \mathbb E_b[D_v^2F_b],\quad
 \Cov_b(D_vF_b),\quad
 T_b,\quad S_b,}
 \label{eq:YMA16-actual-open-rows}
\end{equation}
all on the already frozen selected card.  Version~V15's finite continuation
is opened only if the regulator declares endpoint-matched rather than
frozen-tangent semantics.
\end{corollary}

\begin{proof}
The coefficient and finite action are fixed by
\eqref{eq:YMA16-two-cofactors}--\eqref{eq:YMA16-final-action}.  The physical
margin is the Hessian \eqref{eq:YMA16-native-Hessian} after the declared
local/topological subtraction and comparison with the reference.  These are
exactly the rows in \eqref{eq:YMA16-actual-open-rows}; no further scalar path
coordinate enters their definition.
\end{proof}

% =============================================================================
\section{One base-reference likelihood cylinder for the final physical card}
\label{sec:YMA16-likelihood-card}
% =============================================================================

The final law may be sampled directly.  For exact finite evaluation it is also
useful to express it as one likelihood deformation of the already loaded
pre-retuning law.  This removes any need to reconstruct a response-level path
and keeps the deterministic and covariance terms on one common cylinder.

Let
\begin{equation}
 Z_0=\int_Y e^{-F_0(z,0)}\,\dd m(z),
 \qquad
 \dd\mu_0(z)=Z_0^{-1}e^{-F_0(z,0)}\,\dd m(z).
 \label{eq:YMA16-base-law}
\end{equation}
For real $b$, put
\begin{equation}
 w_b(z)=e^{-bP(z)},
 \qquad
 z_b=\mathbb E_0w_b>0,
 \qquad
 \dd\mu_b=z_b^{-1}w_b\,\dd\mu_0.
 \label{eq:YMA16-likelihood}
\end{equation}
Then the physical partition function is $Z_b=Z_0z_b$.

On the retained finite field bank $E$, define the complete field jets
\begin{align}
 J_0(z)&=D_vF_0(z,0),
 &J_P(z)&=D_vF_P(z,0),
 \label{eq:YMA16-score-jets}\\
 K_0(z)&=D_v^2F_0(z,0),
 &K_P(z)&=D_v^2F_P(z,0),
 \label{eq:YMA16-curvature-jets}
\end{align}
and
\begin{equation}
 J_b=J_0+bJ_P,
 \qquad
 K_b=K_0+bK_P.
 \label{eq:YMA16-total-jets}
\end{equation}
Let
\begin{equation}
 \mathcal L_{\delta,b}=T_b+\delta S_b
 \label{eq:YMA16-loss-form}
\end{equation}
be the complete deterministic topological/reference subtraction on the same
field bank.  No assumption that it is independent of $b$ is made below.

For two copies $z_1,z_2$, write
\begin{equation}
 \Delta J_b=J_b(z_1)-J_b(z_2),
 \qquad
 Q=P(z_1)+P(z_2),
 \label{eq:YMA16-pair-notation}
\end{equation}
and define the Hermitian writer
\begin{equation}
 \boxed{
 \mathscr M_{\delta,b}(z_1,z_2)
 =\frac12\bigl(K_b(z_1)+K_b(z_2)\bigr)
  -\frac12\Delta J_b\Delta J_b^*
  -\mathcal L_{\delta,b}.}
 \label{eq:YMA16-pair-writer}
\end{equation}

\begin{theorem}[Exact base-reference likelihood card]
\label{thm:YMA16-likelihood-card}
The single base-law two-copy expectation
\begin{equation}
 \boxed{
 \mathbb W_{\delta}(b)
 :=\mathbb E_0^{\otimes2}
 \left[e^{-bQ}\mathscr M_{\delta,b}\right]}
 \label{eq:YMA16-weighted-card}
\end{equation}
satisfies
\begin{equation}
 \boxed{
 \mathbb W_{\delta}(b)
 =z_b^2\left(H_b-T_b-\delta S_b\right).}
 \label{eq:YMA16-weighted-identity}
\end{equation}
Consequently $\mathbb W_{\delta}(b)$ and the complete physical margin have
identical inertia, kernel, and witness vectors.  In particular, one negative
vector of the weighted card is already a negative vector of the physical
response-matched mode.
\end{theorem}

\begin{proof}
The first term in \eqref{eq:YMA16-pair-writer} gives
\[
 \mathbb E_0^{\otimes2}
 \left[e^{-bQ}\frac{K_b(z_1)+K_b(z_2)}2\right]
 =z_b^2\mathbb E_bK_b.
\]
For independent copies from $\mu_b$,
\[
 \frac12\mathbb E_b^{\otimes2}
 \left[\Delta J_b\Delta J_b^*\right]
 =\mathbb E_b[J_bJ_b^*]
  -(\mathbb E_bJ_b)(\mathbb E_bJ_b)^*
 =\Cov_b(J_b).
\]
Multiplying by $z_b^2$ converts this identity to the corresponding weighted
$\mu_0^{\otimes2}$ expectation.  The deterministic form
$\mathcal L_{\delta,b}$ contributes $z_b^2\mathcal L_{\delta,b}$.
Substitution of \eqref{eq:YMA16-native-Hessian} proves
\eqref{eq:YMA16-weighted-identity}.  Since $z_b^2>0$, multiplication by this
scalar preserves inertia, kernel, and all signs on fixed vectors.
\end{proof}

\begin{corollary}[The first axial calibration is one scalar likelihood Read]
\label{cor:YMA16-first-axial}
On the frozen $M=4$, $q=\pi/2$ real transverse colour polarization used in
Versions~V3--V5, the reference value is
\begin{equation}
 s_0(\pi/2)=\frac{833}{1818}.
 \label{eq:YMA16-first-reference}
\end{equation}
Hence the direct scalar writer
\begin{equation}
 \mathbb W_{\delta,4}(b;\pi/2)
 =\mathbb E_0^{\otimes2}
 \left[e^{-bQ}\mathscr M_{\delta,b}(\pi/2)\right]
 \label{eq:YMA16-first-scalar-card}
\end{equation}
has the sign of
\begin{equation}
 h_{b,4}^{\rm loc}(\pi/2)
 -\delta\beta\frac{833}{1818}.
 \label{eq:YMA16-first-scalar-margin}
\end{equation}
No finite response-level continuation or boundary-source incidence is needed
to read this first quadratic verdict.
\end{corollary}

\begin{proof}
Use \eqref{eq:YMA2-reference-checks} in
\eqref{eq:YMA16-weighted-identity} and restrict to the already fixed scalar
polarization.
\end{proof}

\begin{proposition}[Gauge held-out survives likelihood recoding]
\label{prop:YMA16-gauge-heldout}
Suppose the complete final action, topological split, and reference are gauge
invariant on a coarse gauge direction $g$, so that
\begin{equation}
 (H_b-T_b)g=0,
 \qquad
 S_bg=0.
 \label{eq:YMA16-gauge-premise}
\end{equation}
Then
\begin{equation}
 \boxed{\mathbb W_{\delta}(b)g=0.}
 \label{eq:YMA16-gauge-zero}
\end{equation}
Thus a nonzero gauge row in the base-reference implementation is an input or
numerical error, not a physical mode.
\end{proposition}

\begin{proof}
Apply \eqref{eq:YMA16-weighted-identity} to $g$ and use
\eqref{eq:YMA16-gauge-premise}.
\end{proof}

% =============================================================================
\section{Exact coefficient flow and a rigorous interval in the marginal coefficient}
\label{sec:YMA16-b-flow}
% =============================================================================

The coefficient $b$ is obtained from response cofactors and will generally be
known by an outward interval.  The next theorem transports that interval to
the complete physical card without evaluating a family of normalized final
laws.

\begin{theorem}[Exact first and second coefficient derivatives]
\label{thm:YMA16-b-derivatives}
Assume $b\mapsto\mathcal L_{\delta,b}$ is twice differentiable on a compact
real interval.  Put
\begin{align}
 \dot{\mathscr M}_{\delta,b}
 ={}&\frac12\bigl(K_P(z_1)+K_P(z_2)\bigr)
 \notag\\
 &-\frac12\left(
   \Delta J_P\Delta J_b^*
   +\Delta J_b\Delta J_P^*\right)
 -\dot{\mathcal L}_{\delta,b},
 \label{eq:YMA16-M-first}\\
 \ddot{\mathscr M}_{\delta,b}
 ={}&-\Delta J_P\Delta J_P^*
 -\ddot{\mathcal L}_{\delta,b}.
 \label{eq:YMA16-M-second}
\end{align}
Then
\begin{equation}
 \boxed{
 \mathbb W_{\delta}'(b)
 =\mathbb E_0^{\otimes2}
 \left[e^{-bQ}
 \left(\dot{\mathscr M}_{\delta,b}
       -Q\mathscr M_{\delta,b}\right)\right],}
 \label{eq:YMA16-W-first}
\end{equation}
\begin{equation}
 \boxed{
 \mathbb W_{\delta}''(b)
 =\mathbb E_0^{\otimes2}
 \left[e^{-bQ}
 \left(\ddot{\mathscr M}_{\delta,b}
       -2Q\dot{\mathscr M}_{\delta,b}
       +Q^2\mathscr M_{\delta,b}\right)\right].}
 \label{eq:YMA16-W-second}
\end{equation}
Both derivatives are same-base-law two-copy Reads.  No score covariance is
charged separately from the writer in \eqref{eq:YMA16-pair-writer}.
\end{theorem}

\begin{proof}
The fibre is compact and every displayed writer is continuous on the finite
card, so differentiation under the integral is legitimate.  Differentiate
$J_b=J_0+bJ_P$ and $K_b=K_0+bK_P$ in
\eqref{eq:YMA16-pair-writer}; this gives
\eqref{eq:YMA16-M-first} and \eqref{eq:YMA16-M-second}.  Since $Q$ is
independent of $b$,
\[
 \frac{d}{db}\left(e^{-bQ}\mathscr M_b\right)
 =e^{-bQ}(\dot{\mathscr M}_b-Q\mathscr M_b),
\]
which proves \eqref{eq:YMA16-W-first}.  Differentiate once more to obtain
\eqref{eq:YMA16-W-second}.
\end{proof}

\begin{proposition}[Explicit finite-card derivative bounds]
\label{prop:YMA16-derivative-bounds}
Let $I\subset[-B,B]$ and suppose, pointwise on the reduced fibre,
\begin{equation}
 \begin{gathered}
 |P|\le p,
 \qquad
 \norm{J_0}\le r_0,
 \qquad
 \norm{J_P}\le r_P,\\
 \norm{K_0}\le k_0,
 \qquad
 \norm{K_P}\le k_P,
 \qquad
 \sup_{b\in I}\norm{\mathcal L_{\delta,b}^{(j)}}\le\ell_j
 \quad(j=0,1,2).
 \end{gathered}
 \label{eq:YMA16-uniform-stocks}
\end{equation}
Define
\begin{align}
 M_0&=k_0+Bk_P+2(r_0+Br_P)^2+\ell_0,
 \label{eq:YMA16-M0}\\
 M_1&=k_P+4r_P(r_0+Br_P)+\ell_1,
 \label{eq:YMA16-M1}\\
 M_2&=4r_P^2+\ell_2,
 \label{eq:YMA16-M2}
\end{align}
and
\begin{align}
 L_1&=e^{2Bp}(M_1+2pM_0),
 \label{eq:YMA16-L1}\\
 L_2&=e^{2Bp}(M_2+4pM_1+4p^2M_0).
 \label{eq:YMA16-L2}
\end{align}
Then, for every $b\in I$,
\begin{equation}
 \boxed{
 \norm{\mathbb W_{\delta}'(b)}\le L_1,
 \qquad
 \norm{\mathbb W_{\delta}''(b)}\le L_2.}
 \label{eq:YMA16-derivative-floor}
\end{equation}
All constants are finite and explicit on each finite Wilson card.
\end{proposition}

\begin{proof}
For $|b|\le B$,
\[
 \norm{\Delta J_b}
 \le2(r_0+Br_P),
 \qquad
 \norm{\Delta J_P}\le2r_P.
\]
Consequently \eqref{eq:YMA16-pair-writer},
\eqref{eq:YMA16-M-first}, and \eqref{eq:YMA16-M-second} give
\[
 \norm{\mathscr M_{\delta,b}}\le M_0,
 \qquad
 \norm{\dot{\mathscr M}_{\delta,b}}\le M_1,
 \qquad
 \norm{\ddot{\mathscr M}_{\delta,b}}\le M_2.
\]
Indeed, the rank-one covariance writer has norm at most
$\frac12\norm{\Delta J_b}^2$, while the mixed derivative has norm at most
$\norm{\Delta J_P}\norm{\Delta J_b}$.  Moreover
$|Q|\le2p$ and $e^{-bQ}\le e^{2Bp}$.  Insert these bounds into
\eqref{eq:YMA16-W-first}--\eqref{eq:YMA16-W-second} and use that
$\mu_0^{\otimes2}$ is a probability measure.
\end{proof}

\begin{corollary}[One centre card controls a coefficient interval]
\label{cor:YMA16-centre-card}
Let the sign-safe cofactor quotient give
\begin{equation}
 b\in I_b=[\bar b-r,\bar b+r]\subset[-B,B].
 \label{eq:YMA16-b-interval}
\end{equation}
Suppose a Hermitian matrix $\widehat W$ satisfies
\begin{equation}
 \norm{\mathbb W_{\delta}(\bar b)-\widehat W}\le\epsilon_0.
 \label{eq:YMA16-centre-error}
\end{equation}
Then the actual physical coefficient obeys
\begin{equation}
 \boxed{
 \norm{\mathbb W_{\delta}(b)-\widehat W}
 \le\epsilon_0+L_1r.}
 \label{eq:YMA16-first-order-radius}
\end{equation}

If, in addition, a Hermitian derivative approximation satisfies
\begin{equation}
 \norm{\mathbb W_{\delta}'(\bar b)-\widehat W'}\le\epsilon_1,
 \label{eq:YMA16-derivative-error}
\end{equation}
then
\begin{equation}
 \boxed{
 \norm{\mathbb W_{\delta}(b)-\widehat W}
 \le
 \epsilon_0+r(\norm{\widehat W'}+\epsilon_1)
 +\frac12L_2r^2.}
 \label{eq:YMA16-second-order-radius}
\end{equation}
The smaller of the two displayed radii may be used in the final spectral
certificate.
\end{corollary}

\begin{proof}
The first estimate is the fundamental theorem of calculus together with
\eqref{eq:YMA16-derivative-floor}.  For the second, Taylor's theorem with
integral remainder gives
\[
 \mathbb W(b)
 =\mathbb W(\bar b)
 +(b-\bar b)\mathbb W'(\bar b)+R_2,
 \qquad
 \norm{R_2}\le\frac12L_2r^2.
\]
Insert \eqref{eq:YMA16-centre-error} and
\eqref{eq:YMA16-derivative-error}, and use $|b-\bar b|\le r$.
\end{proof}

\begin{remark}[A one-sided Wilson improvement]
\label{rem:YMA16-one-sided-weight}
For the standard Wilson plaquette interaction,
$0\le P\le\frac32|\mathscr P_P|$.  If the certified coefficient interval is
nonnegative, then $e^{-bQ}\le1$ and the factors $e^{2Bp}$ in
\eqref{eq:YMA16-L1}--\eqref{eq:YMA16-L2} may be replaced by one.  If the
interaction is normalized by volume or by a selected local support, the
corresponding bound for $p$ must use that actual normalization.
\end{remark}

% =============================================================================
\section{Conditioning of the likelihood gauge}
\label{sec:YMA16-likelihood-conditioning}
% =============================================================================

The base-reference formula is an exact identity, not a claim that importance
sampling from the pre-retuning law is uniformly efficient.  The distinction
is measurable on the same finite card.

\begin{proposition}[Exact likelihood-conditioning scalar]
\label{prop:YMA16-likelihood-conditioning}
Let
\begin{equation}
 L_b(z)=\frac{e^{-bP(z)}}{z_b}
 =\frac{\dd\mu_b}{\dd\mu_0}(z).
 \label{eq:YMA16-normalized-likelihood}
\end{equation}
Then
\begin{equation}
 \boxed{
 \kappa_{\rm rw}(b)
 :=\mathbb E_0L_b^2
 =\frac{z_{2b}}{z_b^2}\ge1.}
 \label{eq:YMA16-rw-condition}
\end{equation}
For every bounded scalar entry $A$ of the final-law card,
\begin{equation}
 \Var_0(L_bA)
 \le\kappa_{\rm rw}(b)\norm{A}_{\infty}^2.
 \label{eq:YMA16-rw-variance}
\end{equation}
Thus a uniform bound for $\kappa_{\rm rw}$ is sufficient for stable
pre-retuning-law sampling.  Its failure does not change the exact identity
\eqref{eq:YMA16-weighted-identity}; it directs the acquisition to the final
law or to a source-connected interpolation instead.
\end{proposition}

\begin{proof}
The first equality follows from
\[
 \mathbb E_0L_b^2
 =z_b^{-2}\mathbb E_0e^{-2bP}
 =z_{2b}/z_b^2.
\]
Cauchy--Schwarz gives $z_b^2\le z_{2b}$ and hence the lower bound one.  Finally
\[
 \Var_0(L_bA)
 \le\mathbb E_0|L_bA|^2
 \le\norm{A}_{\infty}^2\mathbb E_0L_b^2.
\]
\end{proof}

\begin{firewall}[Exact recoding is not cofinal control]
\label{fire:YMA16-reweighting-not-uniform}
On one compact finite card, \eqref{eq:YMA16-weighted-card} is a complete
mathematical reduction.  Along a growing family, an extensive $P$ may make
$\kappa_{\rm rw}$ or the crude factor $e^{2Bp}$ grow rapidly.  That is a
sampling or uniform-estimate obstruction, not a negative Wilson mode.  A
cofinal theorem must instead provide a bounded final-law Read, a local
shell-factor likelihood, or the connected head/tail estimate already isolated
in Versions~V1, V5, and V6.  No global importance-sampling bound is inferred
from finite-card compactness.
\end{firewall}

% =============================================================================
\section{The restored finite certificate}
\label{sec:YMA16-certificate}
% =============================================================================

\begin{theorem}[Tangent-matched finite Wilson verdict]
\label{thm:YMA16-finite-verdict}
Fix one finite selected shell, one retained Hermitian field bank, and a target
margin $\delta\ge0$.  Assume the following outward data are supplied.
\begin{enumerate}[label=\textup{(C\arabic*)},leftmargin=4.4em]
\item Cofactor intervals $I_P$ and $I_{\Psi}$ for $c$ and $e$, with
      $0\notin I_P$.
\item The sign-safe quotient interval
      \begin{equation}
       I_b=-\mathcal Q(I_{\Psi},I_P)
       =[\bar b-r,\bar b+r].
       \label{eq:YMA16-quotient-interval}
      \end{equation}
\item A centre approximation $\widehat W$ to
      $\mathbb W_{\delta}(\bar b)$ with error $\epsilon_0$; optionally a
      derivative approximation $\widehat W'$ with error $\epsilon_1$.
\item The finite writer bounds in
      \eqref{eq:YMA16-uniform-stocks}, or another rigorous bound producing an
      operator radius $\epsilon_b$ for all $b\in I_b$.
\item Exact support, Hermiticity, Ward, response-law, and
      local/topological/reference loading checks.
\end{enumerate}
Let
\begin{equation}
 \epsilon_{\rm card}
 =\epsilon_0+
 \min\left\{
 L_1r,
 r(\norm{\widehat W'}+\epsilon_1)+\frac12L_2r^2
 \right\},
 \label{eq:YMA16-card-radius}
\end{equation}
where the second entry is omitted when no derivative card is supplied.  Then
\begin{equation}
 \boxed{
 \begin{array}{rcl}
 \lambda_{\min}(\widehat W)-\epsilon_{\rm card}>0
 &\Longrightarrow&\pass,\\[0.3em]
 \lambda_{\min}(\widehat W)+\epsilon_{\rm card}<0
 &\Longrightarrow&\obstruction,\\[0.3em]
 \text{otherwise}
 &\Longrightarrow&\stopstatus.
 \end{array}}
 \label{eq:YMA16-verdict}
\end{equation}
On the obstruction branch, a unit minimum eigenvector $v$ of $\widehat W$
satisfies
\begin{equation}
 v^*\mathbb W_{\delta}(b)v<0,
 \qquad
 v^*(H_b-T_b-\delta S_b)v<0,
 \label{eq:YMA16-explicit-witness}
\end{equation}
and is therefore an explicit physical polarization.

If $c=0$ is proved exactly while $e\ne0$, the output is the typed plaquette
response-chart obstruction of Version~V13.  If $c=e=0$, the least nonzero
higher marginal cofactor is required.  An interval $I_P$ meeting zero without
an exact branch decision returns \stopstatus, not a field obstruction.
\end{theorem}

\begin{proof}
The cofactor quotient in \textup{(C2)} contains the actual frozen marginal
coefficient by Version~V13.  The bounds of
\Cref{cor:YMA16-centre-card} give
\[
 \norm{\mathbb W_{\delta}(b)-\widehat W}
 \le\epsilon_{\rm card}.
\]
Weyl's eigenvalue perturbation inequality yields the positive branch of
\eqref{eq:YMA16-verdict}.  If the upper endpoint is negative and $v$ is a
minimum eigenvector of $\widehat W$, then
\[
 v^*\mathbb W_{\delta}(b)v
 \le\lambda_{\min}(\widehat W)+\epsilon_{\rm card}<0.
\]
The second inequality in \eqref{eq:YMA16-explicit-witness} follows from the
positive factor in \eqref{eq:YMA16-weighted-identity}.  The remaining cases
are the exact cofactor alternatives already proved in Version~V13.
\end{proof}

\begin{corollary}[Target-minimal common-cylinder acquisition]
\label{cor:YMA16-common-cylinder}
When the frozen response law is the pre-retuning law $\mu_0$, the complete
finite tangent-matched decision can be acquired from the following data on
one common base cylinder:
\begin{equation}
 \boxed{
 \left(
 \mathfrak c_u[P],
 \mathfrak c_u[\Psi^{\rm cut}],
 \mathbb W_{\delta}(\bar b),
 \mathbb W_{\delta}'(\bar b)
 \right).}
 \label{eq:YMA16-minimal-read}
\end{equation}
The derivative entry is optional.  The first two entries use the unnormalized
root-cofactor writer of Version~V13; the last two use
\eqref{eq:YMA16-weighted-card} and \eqref{eq:YMA16-W-first}.  They are not
four independent source prices: they are four target functionals of the same
loaded shell law and its declared likelihood deformation.
\end{corollary}

\begin{proof}
The cofactor identity of Version~V13 and the two likelihood identities of the
present version are all expectations on products of the same finite base law.
Independent copies may be supplied by one product cylinder.  The quotient,
interval transport, and verdict are then
\Cref{thm:YMA16-finite-verdict}.
\end{proof}

\begin{assessment}[What V16 changes in the proof ledger]
\label{ass:YMA16-qualification}
The finite response-level results of Version~V15 are retained but moved to an
optional endpoint-matched scheme.  On the native YM--A shell:
\begin{enumerate}[label=\textup{(Q\arabic*)},leftmargin=4.4em]
\item the phrase ``physical generated amplitude'' is removed from the default
      loading list;
\item the cofactor ratio is restored to its inherited role as the frozen
      marginal coefficient of the exact action $F_b$;
\item the final normalized law need not be separately sampled, because one
      exact likelihood-weighted base-law card has the same inertia;
\item coefficient uncertainty is transported by the exact first/second
      $b$-derivative writers rather than by a nonlinear response-level path;
\item the selected mode remains unresolved only because the cofactors and the
      weighted physical card have not been populated.
\end{enumerate}
No item turns the tangent condition into exact finite response preservation,
or suppresses a generated, coarse-only, topological, or support-birth term
already present in $F_0$.
\end{assessment}

% =============================================================================
\section{Finite execution, cofinal scope, and the next literal card}
\label{sec:YMA16-frontier}
% =============================================================================

The fixed-card algebra is now direct.  The V10 Haar engine applies to every
bounded term in \eqref{eq:YMA16-minimal-read}: the likelihood factor merely
changes the finite action from $F_0$ to $F_0+bF_P$, while the response
cofactors remain the pair-marked loop numerators of Version~V13.  The present
version does not repeat that invariant-tensor expansion or claim that its
global degree bound is computationally efficient.

For cofinal use, the meaningful quantitative alternatives are:
\begin{enumerate}[label=\textup{(U\arabic*)},leftmargin=4.4em]
\item a uniformly bounded coefficient interval, a direct final-law acquisition,
      and a Ward-relative connected displacement tail;
\item a source-local likelihood factor with uniformly bounded
      $\kappa_{\rm rw}$ and finite head/tail stocks;
\item failure of either mechanism, returned as coefficient escape,
      likelihood-conditioning collapse, local connected-tail growth, or a
      genuine harmonic/topological survivor after the assembled mode card is
      read.
\end{enumerate}
Only the last alternative with a negative assembled field vector is a
Yang--Mills mode obstruction.  The others are named acquisition or transport
failures.

\begin{assessment}[Version V16 exact replay]
\label{ass:YMA16-replay}
The accompanying verifier and exact output are
\begin{center}
\path{YMA_V16_tangent_likelihood_replay.py},\\
\path{YMA_V16_tangent_likelihood_replay.json}.
\end{center}
Their SHA--256 digests are
\begin{center}\scriptsize\ttfamily
b9a264fd233cd825c39b5485e2a8ebbafa95d453b0f14b26b0ec6a8d9ac446a7,\\[-0.15em]
fb02e4bc27d18394f184174a7ce70c94deb1da33fb627a6f07c446b67f6c6839.
\end{center}
The replay verifies on a symbolic finite law:
\begin{enumerate}[label=\textup{(R\arabic*)},leftmargin=4.4em]
\item the distinction between a frozen tangent coefficient and a nonlinear
      finite response endpoint;
\item the likelihood identity
      \eqref{eq:YMA16-weighted-identity};
\item the exact first and second derivative writers
      \eqref{eq:YMA16-W-first}--\eqref{eq:YMA16-W-second};
\item the positive pair-difference covariance identity;
\item the likelihood-conditioning scalar
      \eqref{eq:YMA16-rw-condition};
\item one rational robust spectral instance of
      \eqref{eq:YMA16-verdict}.
\end{enumerate}
It does not evaluate a Wilson path integral, a weak-facing root, the physical
response cofactors, or the selected symbol.
\end{assessment}

% =============================================================================
\section{Version V16 proof ledger and termination}
\label{sec:YMA16-ledger}
% =============================================================================

\begingroup\small
\begin{longtable}{>{\raggedright\arraybackslash}p{0.27\textwidth}
                  >{\raggedright\arraybackslash}p{0.20\textwidth}
                  >{\raggedright\arraybackslash}p{0.45\textwidth}}
\caption{Version V16 proof-status ledger.}\label{tab:YMA16-ledger}\\
\toprule
\textbf{Row}&\textbf{Status}&\textbf{Exact scope}\\
\midrule
\endfirsthead
\toprule
\textbf{Row}&\textbf{Status}&\textbf{Exact scope}\\
\midrule
\endhead
Native response semantics&Restored exactly
&The inherited coefficient is a frozen response-null tangent coefficient;
no free generated amplitude is part of the packet.\\[0.3em]
V15 finite endpoint scheme&Retained conditionally
&Valid only for a regulator which explicitly requires exact finite response
preservation along a generated-amplitude path.\\[0.3em]
Final tangent-matched action&Loaded by definition
&$F_b=F_0+bF_P$ is the complete finite action once the two cofactors and
all terms of $F_0,F_P$ are supplied.\\[0.3em]
Base-reference physical card&Proved exactly
&One likelihood-weighted two-copy base-law expectation equals a positive
scalar times the complete final physical margin.\\[0.3em]
Coefficient derivatives&Proved exactly
&First and second $b$ derivatives are explicit same-base-law two-copy
writers.\\[0.3em]
Coefficient interval transport&Proved
&One centre card, optionally its derivative, and explicit finite stocks give
a rigorous operator radius over the cofactor quotient interval.\\[0.3em]
Likelihood conditioning&Separated exactly
&$\kappa_{\rm rw}=z_{2b}/z_b^2$ tests importance-sampling conditioning without
changing the physical sign identity.\\[0.3em]
Finite mode verdict&Terminating conditional theorem
&The loaded cofactor and centre-card intervals return
$\pass/\obstruction/\stopstatus$ with an explicit negative polarization.\\[0.3em]
First physical values&Open
&Neither $\mathfrak c_u[P]$, $\mathfrak c_u[\Psi^{\rm cut}]$, nor the first
weighted axial card has been evaluated on the submitted shell.\\[0.3em]
Cofinal connected tail&Open physical estimate
&Finite likelihood recoding does not replace a uniform Ward-relative
connected head/tail theorem.\\[0.3em]
Selected shell symbol&\stopstatus
&No negative physical mode is obtained; the direct card remains unpopulated.\\[0.3em]
Local-vacuum contraction and mass&Not inferred
&No finite symbol identity is promoted to an OS-time or continuum mass gap.\\
\bottomrule
\end{longtable}
\endgroup

\begin{theorem}[Version V16 tangent-scheme correction]
\label{thm:YMA16-terminal}
Versions~V1--V15 are retained.  Version~V16 additionally proves:
\begin{enumerate}[label=\textup{(E16.\arabic*)},leftmargin=5.6em]
\item the inherited response-matched shell has frozen tangent semantics and
      requires no independently declared generated amplitude;
\item the V15 nonlinear endpoint programme is a valid optional stronger
      scheme, not the default YM--A loading;
\item the finite action $F_b=F_0+bF_P$ has one exact base-reference
      likelihood-weighted two-copy margin card;
\item the weighted card has the same inertia, kernel, and witnesses as the
      physical final-law margin;
\item its first and second marginal-coefficient derivatives are explicit
      same-base-law writers;
\item finite pointwise stocks give explicit first- and second-order interval
      radii in the cofactor quotient;
\item the likelihood second moment separates mathematical exactness from
      sampling conditioning;
\item one cofactor pair and one centre Hessian card give a terminating finite
      symbol certificate.
\end{enumerate}
No item supplies the numerical cofactors, evaluates the first axial weighted
card, proves weak-facing-root escape, controls the cofinal connected tail, or
establishes local-vacuum contraction or a continuum Yang--Mills mass gap.
\end{theorem}

\begin{proof}
Items~\textup{(E16.1)--(E16.2)} are
\Cref{thm:YMA16-semantic-restoration,fire:YMA16-V15-optional,cor:YMA16-amplitude-stop-removed}.
Items~\textup{(E16.3)--(E16.4)} are
\Cref{thm:YMA16-likelihood-card,cor:YMA16-first-axial}.
Items~\textup{(E16.5)--(E16.6)} are
\Cref{thm:YMA16-b-derivatives,prop:YMA16-derivative-bounds,cor:YMA16-centre-card}.
Item~\textup{(E16.7)} is
\Cref{prop:YMA16-likelihood-conditioning,fire:YMA16-reweighting-not-uniform}.
Item~\textup{(E16.8)} is
\Cref{thm:YMA16-finite-verdict,cor:YMA16-common-cylinder}.
\end{proof}

\begin{openproblem}[V17: execute the restored tangent-matched mode card]
\label{op:YMA16-V17}
Preserve every frozen physical choice and do not reopen a finite generated-
amplitude continuation unless the regulator explicitly declares endpoint-
matched semantics.  Execute the following native card.
\begin{enumerate}[label=\textup{(V17.\arabic*)},leftmargin=5.3em]
\item On the actual frozen response law, acquire outward intervals for
      $\mathfrak c_u[P]$ and
      $\mathfrak c_u[\Psi^{\rm cut}]$ using the V13 unnormalized pair
      cofactors.  Return the exact chart branch if the plaquette cofactor
      vanishes.
\item Form the sign-safe interval for
      $b=-\mathfrak c_u[\Psi^{\rm cut}]/\mathfrak c_u[P]$ and load the exact
      finite action $F_b=F_0+bF_P$.
\item On the first predeclared $M=4$, $q=\pi/2$ axial polarization, acquire
      $\mathbb W_{\delta}(\bar b)$ from
      \eqref{eq:YMA16-weighted-card}.  Acquire
      $\mathbb W_{\delta}'(\bar b)$ only if it reduces the coefficient
      uncertainty materially.
\item Run the gauge held-out and verify that every term uses the same loaded
      $F_0,F_P$, response law, topological split, and reference normalization.
\item Apply \Cref{thm:YMA16-finite-verdict}.  Only after a direct physical
      interval is small or negative should the V6--V7 Witten/boundary
      ancestry or the differentiated cofinal tail be reopened.
\end{enumerate}
If the action writers or the two response cofactors are not supplied, the
first status is the corresponding loading or scalar-acquisition
\stopstatus.  A missing finite generated amplitude is no longer reported on
the native tangent branch.
\end{openproblem}

\begin{assessment}[V16 termination]
\label{ass:YMA16-termination}
The upstream semantic issue is resolved: the native programme does not await
a generated amplitude or a nonlinear endpoint path.  It awaits two actual
response cofactors and one likelihood-weighted physical mode card.  The
restored finite certificate is exact and terminating once those values are
populated.  The selected symbol remains \stopstatus because none of those
three physical values is present in the supplied cylinder.  No negative
physical mode has been produced.
\end{assessment}

\section*{Version V16 append-only continuation record}
\addcontentsline{toc}{section}{Version V16 append-only continuation record}
The complete Version~V15 source above is preserved byte for byte up to its
sole terminal document-closing command.  Its SHA--256 digest is
\begin{center}\scriptsize\ttfamily
e84a77425493adca90f704469eeea09f34818fa2b4e1da046c00f26b6b44ce61.
\end{center}
For Version~V17, preserve this accumulated source, append new mathematical
material immediately before the final document-closing command, and increment
the filename to \texttt{\_V17.tex}.  Keep the frozen response law, exact
pre-retuning action, exact cutoff interaction, plaquette marginal writer,
response target, root selector, rooted Hodge carrier, selected axial
polarization, colour normalization, local/topological split, and reference
normalization fixed before reading any cofactor, coefficient, likelihood,
Hessian, covariance, margin, ancestry, or transport value.

% =============================================================================
% END APPENDED VERSION V16 MATERIAL
% =============================================================================


% =============================================================================
% BEGIN APPENDED VERSION V17 MATERIAL
% =============================================================================

\clearpage
\hypersetup{
  pdftitle={Renewal Yang--Mills Reduced-Fibre Wilson Shell Symbol and Local-Vacuum Programme V17},
  pdfsubject={Physical-momentum indexing, fixed-index stiffness certification, and the trajectory-first cofinal card}
}
\section*{Version V17: physical-momentum indexing, the cofinal axial stiffness card,
and the trajectory-first stop}
\addcontentsline{toc}{section}{Version V17: physical-momentum indexing and the cofinal axial stiffness card}

\begin{abstract}
Version~V16 restored the native frozen-tangent response scheme and reduced one
finite response-matched mode to two response cofactors and one likelihood-
weighted two-copy Hessian card.  Its first calibration remained the
$M=4$, $q=\pi/2$ axial polarization.  The present continuation goes one layer
upstream before that finite number is promoted along a cutoff sequence.

On a fixed physical box, a Fourier angle is not itself a physical momentum.
At lattice side $M$, the axial character of integer index $n$ has angle
$q_{M,n}=2\pi n/M$ and physical lattice momentum
$2M\sin(\pi n/M)/L$.  Hence the $M=4$, $q=\pi/2$ card is the first-index
card only at the base cutoff.  Its finite-physical-momentum continuation keeps
$n=1$ and replaces the angle by $q_{M,1}=2\pi/M$.  Holding $q=\pi/2$ while
$M$ grows instead follows $n=M/4$ and is an ultraviolet sequence.

This distinction leads to an exact cofinal card.  The already evaluated flat
reference satisfies
\[
 \frac{s_0(q)}{\omega(q)}
 =\frac{(t+8)(t+24)}{36(t^2+18t+16)},
 \qquad t=\sin^2(q/2),
\]
and its deviation from the continuum stiffness $1/3$ is explicitly positive
and of order $t$.  The Version~V1 fourth-displacement theorem therefore gives,
on the fixed-index mode $q_{M,n}$, a two-sided physical margin interval whose
error is $O(n^2M^{-2})$ times the fused fourth displacement stock.  A uniform
normalized fourth stock makes the sign of every fixed physical mode eventually
equal to the sign of the zero-angle stiffness margin
$\kappa_M/\beta_M-\delta/3$.  Conversely, disagreement with a stiffness gap
forces a fourth-moment stock of order at least $M^2$ and returns a typed
spatial-delocalization row rather than another unidentified low mode.

The weak-facing trajectory dichotomy is therefore placed before the symbol
calculation.  A bounded cofinal root subsequence remains on the inherited
rough-Haar/contact branch.  On an escaping branch, the first physical axial
card is the fixed-index sequence, not a fixed-angle sequence.  Version~V17
does not populate the response cofactors, stiffness, fourth moment, or direct
mode value.  It corrects the cofinal target, proves the exact finite-index
certificate and its failure alternative, and leaves the selected physical
symbol at \stopstatus.
\end{abstract}

% =============================================================================
\section{Upstream correction: an axial angle is not a physical momentum}
\label{sec:YMA17-direction}
% =============================================================================

The finite $M=4$, $q=\pi/2$ calculation of Versions~V3 and V16 is a valid
calibration.  A cofinal conclusion, however, requires a rule identifying the
same physical momentum across changing lattices.  Repeating the same numerical
angle does not provide that rule.

Fix a physical periodic box of axial length $L>0$.  At lattice side $M$, put
\begin{equation}
 a_M:=\frac{L}{M},
 \qquad
 q_{M,n}:=\frac{2\pi n}{M},
 \qquad
 0\le n\le\left\lfloor\frac M2\right\rfloor .
 \label{eq:YMA17-angle-index}
\end{equation}
The associated physical lattice momentum is
\begin{equation}
 \boxed{
 k_{M,n}^{\rm phys}
 :=\frac{2}{a_M}\sin\frac{q_{M,n}}2
 =\frac{2M}{L}\sin\frac{\pi n}{M}.}
 \label{eq:YMA17-physical-momentum}
\end{equation}
The index $n=0$ is the harmonic/gauge character and is not part of the
nonzero axial bank.

\begin{theorem}[Exhaustive axial momentum classification]
\label{thm:YMA17-momentum-classification}
Let $M_j\to\infty$ and let
$1\le n_j\le M_j/2$.  After passage to a subsequence, exactly one of the
following alternatives holds.
\begin{enumerate}[label=\textup{(P17.\arabic*)},leftmargin=5.2em]
\item \emph{Finite physical momentum.}  The integers $n_j$ are constant,
      say $n_j=n$, and
      \begin{equation}
       k_{M_j,n}^{\rm phys}\longrightarrow\frac{2\pi n}{L}.
       \label{eq:YMA17-fixed-index-limit}
      \end{equation}
\item \emph{Mesoscopic ultraviolet escape.}  One has
      $n_j\to\infty$ and $n_j/M_j\to0$; then
      \begin{equation}
       q_{M_j,n_j}\longrightarrow0,
       \qquad
       k_{M_j,n_j}^{\rm phys}\sim\frac{2\pi n_j}{L}
       \longrightarrow\infty.
       \label{eq:YMA17-mesoscopic-UV}
      \end{equation}
\item \emph{Lattice-angle ultraviolet escape.}  There is an
      $\varepsilon>0$ with $n_j/M_j\ge\varepsilon$; then
      \begin{equation}
       k_{M_j,n_j}^{\rm phys}
       \ge\frac{2M_j}{L}\sin(\pi\varepsilon)
       \longrightarrow\infty.
       \label{eq:YMA17-angle-UV}
      \end{equation}
\end{enumerate}
Thus a cofinal sequence has finite nonzero physical momentum precisely when
its Fourier index is eventually a fixed positive integer.
\end{theorem}

\begin{proof}
If $(n_j)$ is bounded, an integer subsequence is constant.  With
$x_j=\pi n/M_j$, the limit
$M_j\sin x_j\to\pi n$ proves
\eqref{eq:YMA17-fixed-index-limit}.

Suppose $(n_j)$ is unbounded and pass to a subsequence on which $n_j\to\infty$.
If $n_j/M_j\to0$, then
$\sin(\pi n_j/M_j)\sim\pi n_j/M_j$, proving
\eqref{eq:YMA17-mesoscopic-UV}.  Otherwise the ratio has a positive limsup;
a further subsequence satisfies $n_j/M_j\ge\varepsilon>0$.  Since the ratio
is at most $1/2$, the sine is bounded below by
$\sin(\pi\varepsilon)>0$, which proves
\eqref{eq:YMA17-angle-UV}.  The alternatives are disjoint by construction.
\end{proof}

\begin{corollary}[The $M=4$, $q=\pi/2$ calibration has two different continuations]
\label{cor:YMA17-calibration-two-continuations}
At $M=4$ the angle $q=\pi/2$ is the index-one character.  Under dyadic
refinement $M_r=4\,2^r$, its finite-physical-momentum continuation is
\begin{equation}
 \boxed{n_r=1,
 \qquad q_r=\frac{2\pi}{M_r}=\frac{\pi}{2^{r+1}}.}
 \label{eq:YMA17-first-index-branch}
\end{equation}
By contrast, keeping $q_r=\pi/2$ gives $n_r=M_r/4$ and belongs to
\textup{(P17.3)}.  It is an ultraviolet card.
\end{corollary}

\begin{proof}
The relation $q_{M,n}=2\pi n/M$ gives both assertions.  The physical-momentum
classification follows from \Cref{thm:YMA17-momentum-classification}.
\end{proof}

\begin{firewall}[A fixed-angle pass is not a finite-physical-momentum pass]
\label{fire:YMA17-angle-not-physical}
A positive interval at $q=\pi/2$ is a legitimate finite-card statement and
may be useful for ultraviolet control.  It does not certify the fixed-index
low-angle branch on larger lattices.  Conversely, failure of a fixed-index
card is not inferred from the fixed-angle calibration.  The two branches can
be related only through an independently proved displacement-kernel or
adjacent-cutoff transport estimate.
\end{firewall}

% =============================================================================
\section{Exact dyadic transport of the physical axial label}
\label{sec:YMA17-dyadic-label}
% =============================================================================

\begin{proposition}[Dyadic angle and momentum transport]
\label{prop:YMA17-dyadic-transport}
For a fixed index $n$ and $M\ge2n$,
\begin{equation}
 q_{2M,n}=\frac12q_{M,n},
 \label{eq:YMA17-angle-halving}
\end{equation}
and
\begin{equation}
 \boxed{
 \frac{k_{2M,n}^{\rm phys}}{k_{M,n}^{\rm phys}}
 =\sec\frac{\pi n}{2M}.}
 \label{eq:YMA17-momentum-ratio}
\end{equation}
In particular,
\begin{equation}
 0\le\frac{2\pi n}{L}-k_{M,n}^{\rm phys}
 \le\frac{\pi^3n^3}{3LM^2}.
 \label{eq:YMA17-momentum-rate}
\end{equation}
For the Ward factor $\omega(q)=2(1-\cos q)$,
\begin{equation}
 \boxed{
 \frac{(2M)^2\omega(q_{2M,n})}
      {M^2\omega(q_{M,n})}
 =\sec^2\frac{\pi n}{2M}.}
 \label{eq:YMA17-Ward-scale-ratio}
\end{equation}
Thus the fixed-index physical normalization is asymptotically invariant under
one dyadic refinement, while a fixed numerical angle is not.
\end{proposition}

\begin{proof}
The angle identity is immediate.  Put $x=\pi n/M$.  Then
\[
 \frac{k_{2M,n}^{\rm phys}}{k_{M,n}^{\rm phys}}
 =\frac{4M\sin(x/2)}{2M\sin x}
 =\frac1{\cos(x/2)}.
\]
The elementary estimate $0\le x-\sin x\le x^3/6$ gives
\eqref{eq:YMA17-momentum-rate}.  Finally
$\omega(q)=4\sin^2(q/2)$ and the double-angle identity give
\eqref{eq:YMA17-Ward-scale-ratio}.
\end{proof}

% =============================================================================
\section{The exact fixed-index scaling of the Wilson reference}
\label{sec:YMA17-reference}
% =============================================================================

Recall the exact reference obtained in Version~V2:
\begin{equation}
 s_0(q)
 =\frac{t(t+8)(t+24)}{9(t^2+18t+16)},
 \qquad
 t=\sin^2(q/2),
 \qquad
 \omega(q)=4t.
 \label{eq:YMA17-reference-recall}
\end{equation}

\begin{theorem}[Reference quotient and its sharp small-angle sign]
\label{thm:YMA17-reference-quotient}
For every $q\ne0$,
\begin{equation}
 \boxed{
 r_0(q):=\frac{s_0(q)}{\omega(q)}
 =\frac{(t+8)(t+24)}{36(t^2+18t+16)}.}
 \label{eq:YMA17-reference-quotient}
\end{equation}
Moreover,
\begin{equation}
 \boxed{
 \frac13-r_0(q)
 =\frac{t(11t+184)}{36(t^2+18t+16)}\ge0,}
 \label{eq:YMA17-reference-defect-exact}
\end{equation}
and, for $0\le t\le1$,
\begin{equation}
 \boxed{
 0\le\frac13-r_0(q)\le\frac{23}{72}t.}
 \label{eq:YMA17-reference-defect-bound}
\end{equation}
The coefficient $23/72$ is the limiting ratio of the two sides as
$t\downarrow0$.
\end{theorem}

\begin{proof}
Divide \eqref{eq:YMA17-reference-recall} by $4t$.  Subtracting the result from
$1/3$ gives
\eqref{eq:YMA17-reference-defect-exact}.  To prove the upper bound, observe
that
\[
 \frac{23}{72}
 -\frac{11t+184}{36(t^2+18t+16)}
 =\frac{t(23t+392)}{72(t^2+18t+16)}\ge0.
\]
Multiplication by $t$ proves \eqref{eq:YMA17-reference-defect-bound}.  The
limit at $t=0$ is immediate.
\end{proof}

\begin{corollary}[Noncollapsed fixed-index reference]
\label{cor:YMA17-reference-scaling}
For $1\le n\le M/2$,
\begin{equation}
 \frac{20}{7}n^2
 \le M^2s_0(q_{M,n})
 \le\frac{4\pi^2}{3}n^2.
 \label{eq:YMA17-reference-window}
\end{equation}
For each fixed $n$,
\begin{equation}
 \boxed{
 M^2s_0(q_{M,n})
 \longrightarrow\frac{4\pi^2n^2}{3}.}
 \label{eq:YMA17-reference-limit}
\end{equation}
More quantitatively,
\begin{equation}
 \boxed{
 \left|
  M^2s_0(q_{M,n})-\frac{4\pi^2n^2}{3}
 \right|
 \le\frac{31\pi^4n^4}{18M^2}.}
 \label{eq:YMA17-reference-rate}
\end{equation}
\end{corollary}

\begin{proof}
Put $x=\pi n/M\in[0,\pi/2]$.  The inequalities
$2x/\pi\le\sin x\le x$ give
\[
 16n^2\le M^2\omega(q_{M,n})\le4\pi^2n^2.
\]
Combine these with the inherited global bound
$5/28\le s_0/\omega\le1/3$ to obtain
\eqref{eq:YMA17-reference-window}.

For the rate, $0\le x^2-\sin^2x\le x^4/3$ follows from
$\sin x\ge x-x^3/6$.  Hence
\[
 0\le4\pi^2n^2-M^2\omega(q_{M,n})
 \le\frac{4\pi^4n^4}{3M^2}.
\]
By \eqref{eq:YMA17-reference-defect-bound} and
$t=\sin^2x\le x^2$,
\[
 0\le
 \frac13M^2\omega(q_{M,n})-M^2s_0(q_{M,n})
 \le\frac{23\pi^4n^4}{18M^2}.
\]
Adding one third of the first error to the second gives
$(4/9+23/18)\pi^4n^4/M^2=31\pi^4n^4/(18M^2)$.
\end{proof}

\begin{assessment}[What the reference scaling fixes]
\label{ass:YMA17-reference-meaning}
The dimensionless symbol itself vanishes like $M^{-2}$ on a fixed-index
branch.  This is the Ward zero and not a collapse of the physical reference.
The scale-correct object is either the quotient by $\omega(q_{M,n})$ or the
rescaled form $M^2s_0(q_{M,n})$.  A cutoff-independent absolute eigenvalue
floor at fixed index would be the wrong target.
\end{assessment}

% =============================================================================
\section{The cofinal fixed-index physical card}
\label{sec:YMA17-physical-card}
% =============================================================================

Let $M_r=2^rM_0$ be the selected lattice sequence on the fixed box.  On an
escaping weak-facing branch, let $\beta_r$ be the selected coupling, let
$b_r$ be the native response-matched marginal coefficient of Version~V16,
and let
\begin{equation}
 h_r^{\rm loc}(q)
 \label{eq:YMA17-local-symbol}
\end{equation}
be the completely assembled central-axial local symbol after the declared
topological subtraction.  The topological or harmonic packet is retained
separately whenever it cannot be represented by the local displacement
kernel.

For a fixed positive integer $n$ define
\begin{equation}
 \boxed{
 \mathfrak m_{r,n}(\delta)
 :=h_r^{\rm loc}(q_{M_r,n})
   -\delta\beta_rs_0(q_{M_r,n}).}
 \label{eq:YMA17-fixed-index-margin}
\end{equation}

\begin{proposition}[V16 likelihood card on the correct cofinal label]
\label{prop:YMA17-likelihood-fixed-index}
For every finite $r,n$, the Version~V16 base-reference card evaluated at
$q_{M_r,n}$ satisfies
\begin{equation}
 \boxed{
 \mathbb W_{\delta,r,n}(b_r)
 =z_{r,b_r}^2\,\mathfrak m_{r,n}(\delta).}
 \label{eq:YMA17-likelihood-sign}
\end{equation}
Consequently it has the same sign and the same one-dimensional transverse
witness as the physical fixed-index margin.  When $M_0=4$ and $n=1$, its
$r=0$ member is the V16 $q=\pi/2$ calibration; its next member is at
$q=\pi/4$, not at $q=\pi/2$.
\end{proposition}

\begin{proof}
Apply the exact V16 likelihood identity at the mode
$q=q_{M_r,n}$.  The likelihood normalizer is strictly positive, so
multiplication by its square preserves sign and witnesses.  The final
statement is \eqref{eq:YMA17-first-index-branch}.
\end{proof}

% =============================================================================
\section{A finite fixed-index stiffness certificate}
\label{sec:YMA17-stiffness}
% =============================================================================

Let the contact-repaired scalar local kernel at cutoff $r$ have aggregated
axial coefficients $\alpha_{r,m}$ in the Version~V1 convention
\begin{equation}
 h_r^{\rm loc}(q)
 =2\sum_{m\ge1}\alpha_{r,m}(\cos(mq)-1).
 \label{eq:YMA17-kernel-expansion}
\end{equation}
Its Ward quotient, stiffness, and fourth stock are
\begin{equation}
 g_r(q):=\frac{h_r^{\rm loc}(q)}{\omega(q)},
 \qquad
 \kappa_r:=-\sum_{m\ge1}m^2\alpha_{r,m},
 \qquad
 \mathcal M_{4,r}:=\sum_{m\ge1}m^4\abs{\alpha_{r,m}}.
 \label{eq:YMA17-stiffness-stocks}
\end{equation}
When only a finite displacement head is acquired, denote its stiffness and
fourth stock by $\kappa_{r,R}$ and $\mathcal M_{4,r,R}$, and denote the
Version~V1 Ward-relative omitted-tail stock by $\tau_{r,R}^{(2)}$.

\begin{theorem}[Exact fixed-index pass/obstruction interval]
\label{thm:YMA17-fixed-index-certificate}
Fix $n\ge1$ and assume $M_r\ge2\pi n$, so that
$q_{M_r,n}\le1$.  Put
\begin{equation}
 E_{r,n,R}
 :=\tau_{r,R}^{(2)}
   +\frac{\pi^2n^2}{M_r^2}\mathcal M_{4,r,R}.
 \label{eq:YMA17-kernel-error}
\end{equation}
Then
\begin{equation}
 \boxed{
 \begin{aligned}
 \frac{\mathfrak m_{r,n}(\delta)}
      {\omega(q_{M_r,n})}
 &\ge
 \kappa_{r,R}-\frac{\delta\beta_r}{3}-E_{r,n,R},\\
 \frac{\mathfrak m_{r,n}(\delta)}
      {\omega(q_{M_r,n})}
 &\le
 \kappa_{r,R}-\frac{\delta\beta_r}{3}+E_{r,n,R}
 +\frac{23\delta\beta_r\pi^2n^2}{72M_r^2}.
 \end{aligned}}
 \label{eq:YMA17-margin-interval}
\end{equation}
Therefore:
\begin{enumerate}[label=\textup{(C17.\arabic*)},leftmargin=5.2em]
\item if the first endpoint in \eqref{eq:YMA17-margin-interval} is positive,
      the physical fixed-index card is \pass;
\item if the second endpoint is negative, the actual axial index-$n$
      polarization is a \obstruction;
\item if the interval meets zero, the finite card is \stopstatus.
\end{enumerate}
No fixed-angle value enters this certificate.
\end{theorem}

\begin{proof}
The Version~V1 low-angle interval gives
\[
 \kappa_{r,R}-E_{r,n,R}
 \le g_r(q_{M_r,n})
 \le\kappa_{r,R}+E_{r,n,R},
\]
because $q_{M_r,n}^2/4=\pi^2n^2/M_r^2$.
Now
\[
 \frac{\mathfrak m_{r,n}(\delta)}{\omega(q_{M_r,n})}
 =g_r(q_{M_r,n})-\delta\beta_rr_0(q_{M_r,n}).
\]
The upper inequality $r_0\le1/3$ gives the lower endpoint.  The bound
\[
 r_0(q_{M_r,n})
 \ge\frac13-\frac{23}{72}\sin^2\frac{\pi n}{M_r}
 \ge\frac13-\frac{23\pi^2n^2}{72M_r^2}
\]
gives the upper endpoint.  Since the Ward factor is positive at nonzero
index, its removal does not change sign.  The three verdicts are ordinary
one-sided interval inclusion.
\end{proof}

\begin{corollary}[Full-kernel form]
\label{cor:YMA17-full-kernel-certificate}
If $\mathcal M_{4,r}<\infty$, take $R=\infty$ in
\Cref{thm:YMA17-fixed-index-certificate}.  Then
\begin{equation}
 E_{r,n}
 =\frac{\pi^2n^2}{M_r^2}\mathcal M_{4,r}.
 \label{eq:YMA17-full-kernel-error}
\end{equation}
Thus the fixed-index mode is determined by the zero-angle stiffness up to an
explicit fourth-moment correction of order $n^2/M_r^2$.
\end{corollary}

\begin{firewall}[The stiffness is the fused physical stiffness]
\label{fire:YMA17-fused-stiffness}
The coefficient $\kappa_r$ in this section is the signed second moment of the
already assembled physical local kernel: deterministic curvature, connected
score structure, generated marginal correction, and every local support birth
have first been combined with their actual signs.  A deterministic stiffness,
a covariance stiffness, or a pointwise Bochner floor cannot be substituted
for it separately.  A nonlocal harmonic or topological row remains outside the
local kernel and is carried by its own finite card.
\end{firewall}

% =============================================================================
\section{Cofinal stiffness alternative on every finite physical bank}
\label{sec:YMA17-cofinal-stiffness}
% =============================================================================

\begin{theorem}[Fixed physical bank from one normalized stiffness limit]
\label{thm:YMA17-cofinal-bank}
Let $n_0\ge1$ be fixed.  Assume along an escaping selected branch that
\begin{equation}
 \sup_r\frac{\mathcal M_{4,r}}{\beta_r}<\infty.
 \label{eq:YMA17-normalized-W4}
\end{equation}
Then the following alternatives hold.
\begin{enumerate}[label=\textup{(B17.\arabic*)},leftmargin=5.2em]
\item If, for some $\eta>0$,
      \begin{equation}
       \liminf_{r\to\infty}
       \left(\frac{\kappa_r}{\beta_r}-\frac\delta3\right)
       \ge\eta,
       \label{eq:YMA17-stiffness-positive-gap}
      \end{equation}
      then every fixed index $1\le n\le n_0$ is eventually \pass.
\item If, for some $\eta>0$,
      \begin{equation}
       \limsup_{r\to\infty}
       \left(\frac{\kappa_r}{\beta_r}-\frac\delta3\right)
       \le-\eta,
       \label{eq:YMA17-stiffness-negative-gap}
      \end{equation}
      then every fixed index $1\le n\le n_0$ is eventually a
      \obstruction.
\item If neither separated stiffness branch holds, the stiffness itself
      approaches the target threshold on a subsequence.  This is the exact
      low-physical-momentum borderline row.
\end{enumerate}
The theorem concerns every predeclared finite bank of physical Fourier
indices simultaneously; it does not assert a full ultraviolet symbol margin.
\end{theorem}

\begin{proof}
Divide the full-kernel endpoints of
\Cref{cor:YMA17-full-kernel-certificate} by $\beta_r$.  Under
\eqref{eq:YMA17-normalized-W4}, the fourth-moment correction is bounded by
$C\pi^2n_0^2/M_r^2$ and tends to zero uniformly for
$1\le n\le n_0$.  The reference correction in the upper endpoint is bounded
by $23\delta\pi^2n_0^2/(72M_r^2)$ and also tends to zero.  The positive and
negative stiffness gaps therefore dominate the corresponding endpoints for
all sufficiently large $r$.  If neither one-sided gap exists, elementary
subsequence selection gives a sequence on which
$\kappa_r/\beta_r-\delta/3$ approaches zero or changes sign arbitrarily close
to zero.
\end{proof}

\begin{corollary}[Continuum stiffness value]
\label{cor:YMA17-continuum-stiffness}
If
\begin{equation}
 \frac{\kappa_r}{\beta_r}\longrightarrow\kappa_\infty,
 \qquad
 \sup_r\frac{\mathcal M_{4,r}}{\beta_r}<\infty,
 \label{eq:YMA17-stiffness-limit-assumption}
\end{equation}
then, for every fixed $n$,
\begin{equation}
 \boxed{
 \frac{h_r^{\rm loc}(q_{M_r,n})}
      {\beta_r\omega(q_{M_r,n})}
 \longrightarrow\kappa_\infty,}
 \label{eq:YMA17-physical-stiffness-limit}
\end{equation}
and
\begin{equation}
 \boxed{
 \frac{h_r^{\rm loc}(q_{M_r,n})}
      {\beta_rs_0(q_{M_r,n})}
 \longrightarrow3\kappa_\infty.}
 \label{eq:YMA17-reference-relative-limit}
\end{equation}
Thus a positive target margin $\delta$ on every fixed physical index is
equivalent, away from the borderline, to
$\kappa_\infty>\delta/3$.
\end{corollary}

\begin{proof}
The first limit is the V1 low-angle estimate divided by $\beta_r$ and the
bounded normalized fourth stock.  The second uses
$r_0(q_{M_r,n})\to1/3$ from
\Cref{thm:YMA17-reference-quotient}.
\end{proof}

% =============================================================================
\section{A failed stiffness prediction forces displacement delocalization}
\label{sec:YMA17-delocalization}
% =============================================================================

\begin{theorem}[Fourth-moment price of disagreement with a stiffness gap]
\label{thm:YMA17-W4-obstruction}
Fix $n\ge1$ and use the full kernel.  Suppose
\begin{equation}
 \frac{\kappa_r}{\beta_r}-\frac\delta3\ge\eta>0.
 \label{eq:YMA17-positive-stiffness-card}
\end{equation}
If the actual index-$n$ margin is nonpositive, then
\begin{equation}
 \boxed{
 \frac{\mathcal M_{4,r}}{\beta_r}
 \ge\frac{\eta M_r^2}{\pi^2n^2}.}
 \label{eq:YMA17-W4-lower-positive}
\end{equation}

Conversely, suppose
\begin{equation}
 \frac{\kappa_r}{\beta_r}-\frac\delta3\le-\eta<0.
 \label{eq:YMA17-negative-stiffness-card}
\end{equation}
If the actual index-$n$ margin is nonnegative, then
\begin{equation}
 \boxed{
 \frac{\mathcal M_{4,r}}{\beta_r}
 \ge
 \frac{\eta M_r^2}{\pi^2n^2}-\frac{23\delta}{72}.}
 \label{eq:YMA17-W4-lower-negative}
\end{equation}
In particular, for fixed $n$ and large $M_r$, disagreement between the
physical mode and a separated stiffness sign forces a normalized fourth
moment of order $M_r^2$.
\end{theorem}

\begin{proof}
Under \eqref{eq:YMA17-positive-stiffness-card}, the lower endpoint in
\eqref{eq:YMA17-margin-interval} with the full kernel is
\[
 \beta_r\eta
 -\frac{\pi^2n^2}{M_r^2}\mathcal M_{4,r}.
\]
If the actual margin is nonpositive, this lower bound cannot be positive,
which proves \eqref{eq:YMA17-W4-lower-positive}.

Under \eqref{eq:YMA17-negative-stiffness-card}, the upper endpoint is at
most
\[
 -\beta_r\eta
 +\frac{\pi^2n^2}{M_r^2}\mathcal M_{4,r}
 +\frac{23\delta\beta_r\pi^2n^2}{72M_r^2}.
\]
If the actual margin is nonnegative, this upper bound cannot be negative.
Rearranging proves \eqref{eq:YMA17-W4-lower-negative}.
\end{proof}

\begin{assessment}[Typed meaning of fourth-moment growth]
\label{ass:YMA17-W4-meaning}
The lower bounds above do not prove a negative Yang--Mills symbol.  They say
that a mode which disagrees with its separated zero-angle stiffness must draw
its sign from displacements whose fourth moment reaches the physical box
scale.  After the V1 contact repair, this is a typed spatial-delocalization or
positive-range row.  It must be located by the physical boundary/tail packet;
it is not repaired by recomputing the same local stiffness.
\end{assessment}

% =============================================================================
\section{Trajectory-first termination and the corrected next acquisition}
\label{sec:YMA17-trajectory}
% =============================================================================

\begin{theorem}[Trajectory gate before cofinal symbol promotion]
\label{thm:YMA17-trajectory-gate}
Apply the inherited weak-facing alternative before using any finite card as a
cofinal local-vacuum statement.
\begin{enumerate}[label=\textup{(T17.\arabic*)},leftmargin=5.2em]
\item If the selected roots have a bounded cofinal subsequence, the inherited
      cylindrical-Haar/contact theorem returns the rough-contact trajectory
      obstruction.  A fixed-$M$ symbol calibration does not reopen that raw
      smooth local-vacuum branch.
\item If the selected roots escape, then for each fixed physical Fourier
      index $n$ the correct axial mode is $q_{M_r,n}=2\pi n/M_r$.  The V16
      response cofactors and likelihood card are evaluated at that label.
\item A fixed numerical angle on the escaping branch is an ultraviolet card
      and is classified separately by
      \Cref{thm:YMA17-momentum-classification}.
\end{enumerate}
No fourth branch called ``repeat the $M=4$ angle at every cutoff'' is a
finite-physical-momentum continuation.
\end{theorem}

\begin{proof}
The first item is the inherited trajectory dichotomy and bounded-branch
landing.  The second and third items are
\Cref{thm:YMA17-momentum-classification,cor:YMA17-calibration-two-continuations}.
They are exhaustive for the selected axial labels.
\end{proof}

\begin{corollary}[Target-minimal escaping-branch acquisition]
\label{cor:YMA17-target-minimal-card}
On an escaping branch, a fixed finite physical axial bank
$1\le n\le n_0$ can be closed by either of the following nonduplicating
packets.
\begin{enumerate}[label=\textup{(A17.\arabic*)},leftmargin=5.2em]
\item Directly acquire the V16 likelihood-weighted card
      $\mathbb W_{\delta,r,n}(b_r)$ at each selected fixed index.
\item Acquire the fused stiffness $\kappa_r$, a fourth-moment or finite
      head/tail certificate, and apply
      \Cref{thm:YMA17-fixed-index-certificate}.
\end{enumerate}
The two packets evaluate the same physical mode and are alternative
implementations, not two source prices.  Boundary/Witten ancestry is opened
only after a direct failure or a fourth-moment delocalization has been frozen.
\end{corollary}

\begin{firewall}[No mass conclusion from physical-momentum bookkeeping]
\label{fire:YMA17-no-mass}
Correct fixed-index transport prevents a UV card from being mislabeled as a
finite-physical-momentum card.  It does not populate the response cofactors,
the fused stiffness, the fourth displacement stock, the topological packet,
a full centre-neutral network gap, an Osterwalder--Schrader time scale, or a
continuum Yang--Mills mass.  Those remain separate physical rows.
\end{firewall}

% =============================================================================
\section{Version V17 proof ledger and exact next attack}
\label{sec:YMA17-ledger}
% =============================================================================

\begingroup\small
\begin{longtable}{>{\raggedright\arraybackslash}p{0.28\textwidth}
                  >{\raggedright\arraybackslash}p{0.19\textwidth}
                  >{\raggedright\arraybackslash}p{0.45\textwidth}}
\caption{Version V17 proof-status ledger.}\label{tab:YMA17-ledger}\\
\toprule
\textbf{Row}&\textbf{Status}&\textbf{Exact scope}\\
\midrule
\endfirsthead
\toprule
\textbf{Row}&\textbf{Status}&\textbf{Exact scope}\\
\midrule
\endhead
Axial momentum classification&Proved exactly
&Fixed index, mesoscopic UV, and fixed-angle UV are exhaustive after
subsequence extraction.\\[0.3em]
Base calibration continuation&Corrected exactly
&$M=4,q=\pi/2$ is the $n=1$ base card; fixed physical momentum halves the
angle under every dyadic refinement.\\[0.3em]
Physical momentum transport&Proved exactly
&The dyadic ratios are the displayed secant and squared-secant factors, with
an explicit $M^{-2}$ convergence rate.\\[0.3em]
Wilson reference scaling&Proved exactly
&$s_0/\omega$ tends to $1/3$ at the Ward zero with the exact defect and an
explicit fixed-index $M^{-2}$ rate.\\[0.3em]
Fixed-index symbol interval&Proved exactly
&The V1 stiffness/fourth-moment theorem gives the two endpoints in
\eqref{eq:YMA17-margin-interval}.\\[0.3em]
Finite physical bank&Conditional cofinal theorem
&A normalized fourth-stock bound makes the zero-angle stiffness decide every
fixed finite index bank away from the threshold.\\[0.3em]
Disagreement with stiffness&Typed quantitative alternative
&A contrary physical sign forces normalized fourth-moment growth of order
$M^2$.\\[0.3em]
Trajectory ordering&Corrected
&Bounded roots stop on the inherited contact branch; escaping roots use the
fixed-index card; fixed angle is UV.\\[0.3em]
Response cofactors&Open physical values
&No interval for $\mathfrak c_u[P]$ or
$\mathfrak c_u[\Psi^{\rm cut}]$ is supplied.\\[0.3em]
Fused stiffness and fourth stock&Open physical values
&Neither the complete second displacement moment nor its fourth stock is
populated on the escaping selected law.\\[0.3em]
Selected cofinal symbol&\stopstatus
&The correct cofinal label is now fixed, but no direct or stiffness packet is
numerically acquired.\\[0.3em]
Local-vacuum contraction and mass&Not inferred
&No fixed-index bookkeeping theorem is promoted to a temporal or continuum
gap.\\
\bottomrule
\end{longtable}
\endgroup

\begin{theorem}[Version V17 physical-index correction]
\label{thm:YMA17-terminal}
Versions~V1--V16 are retained.  Version~V17 additionally proves:
\begin{enumerate}[label=\textup{(E17.\arabic*)},leftmargin=5.6em]
\item finite physical axial momentum is represented by fixed Fourier index,
      not fixed lattice angle;
\item the $M=4,q=\pi/2$ card is the base member of the $n=1$ branch, while
      its fixed-angle continuation is ultraviolet;
\item the physical momentum, Ward normalization, and Wilson reference have
      exact dyadic transport formulas and explicit $M^{-2}$ errors;
\item the complete fixed-index margin has the finite stiffness interval
      \eqref{eq:YMA17-margin-interval};
\item a uniform normalized fourth stock makes the fused zero-angle stiffness
      decide every fixed finite physical bank away from the threshold;
\item disagreement with a separated stiffness sign forces a box-scale fourth
      displacement stock;
\item the weak-facing trajectory alternative precedes every cofinal symbol
      promotion.
\end{enumerate}
No item supplies the response cofactors, selects the escaping branch,
populates the fused physical stiffness or fourth moment, evaluates the direct
fixed-index likelihood card, or proves local-vacuum contraction or a continuum
Yang--Mills mass gap.
\end{theorem}

\begin{proof}
Items~\textup{(E17.1)--(E17.2)} are
\Cref{thm:YMA17-momentum-classification,cor:YMA17-calibration-two-continuations}.
Item~\textup{(E17.3)} is
\Cref{prop:YMA17-dyadic-transport,thm:YMA17-reference-quotient,cor:YMA17-reference-scaling}.
Item~\textup{(E17.4)} is
\Cref{thm:YMA17-fixed-index-certificate}.
Item~\textup{(E17.5)} is
\Cref{thm:YMA17-cofinal-bank,cor:YMA17-continuum-stiffness}.
Item~\textup{(E17.6)} is
\Cref{thm:YMA17-W4-obstruction}.
Item~\textup{(E17.7)} is
\Cref{thm:YMA17-trajectory-gate}.
\end{proof}

\begin{openproblem}[V18: execute the first fixed-physical-momentum card]
\label{op:YMA17-V18}
Preserve every frozen physical choice.  Do not repeat the numerical angle
$q=\pi/2$ at larger cutoffs and call it the first physical mode.  Proceed in
this order.
\begin{enumerate}[label=\textup{(V18.\arabic*)},leftmargin=5.3em]
\item Apply the weak-facing trajectory gate.  On a bounded cofinal
      subsequence, export the inherited rough-Haar/contact obstruction and
      stop the raw smooth local-vacuum branch.
\item On an escaping selected shell sequence $M_r=2^rM_0$, freeze the first
      physical axial label $n=1$, hence $q_r=2\pi/M_r$.
\item Acquire the actual response cofactors
      $\mathfrak c_{u,r}[P_r]$ and
      $\mathfrak c_{u,r}[\Psi_r^{\rm cut}]$ and load the native
      response-matched law at each selected shell.
\item Either acquire the direct V16 card
      $\mathbb W_{\delta,r,1}(b_r)$, or acquire the fused stiffness and a
      certified fourth-moment/head-tail packet and apply
      \Cref{thm:YMA17-fixed-index-certificate}.
\item Run the gauge held-out and retain the topological/harmonic card
      separately before assigning a sign.
\item If the stiffness is separated but the physical sign disagrees, export
      the attained fourth-moment delocalization bound and open the physical
      boundary/tail router on that same mode.
\item Only after the $n=1$ card closes should the bank be enlarged to fixed
      $n=2,3,\ldots$ or to the ultraviolet fixed-angle sector.
\end{enumerate}
If no escaping selected law or no cofactor/direct/stiffness packet is supplied,
the exact verdict remains \stopstatus{} at the first missing row.  Another
fixed-angle calibration is not a substitute.
\end{openproblem}

\begin{assessment}[V17 termination]
\label{ass:YMA17-termination}
The circularity in the preceding finite target has been removed.  The
$M=4,q=\pi/2$ value is retained as the base $n=1$ calibration, but the
cofinal physical card follows $q=2\pi/M$.  Its reference is noncollapsed in
Ward units, and its sign is reduced to one fused stiffness and one fourth
moment with explicit errors.  The selected symbol remains \stopstatus{} because
the escaping trajectory, response cofactors, and direct or stiffness values
are not populated.  No negative physical mode has been produced.
\end{assessment}

\section*{Version V17 append-only continuation record}
\addcontentsline{toc}{section}{Version V17 append-only continuation record}
The complete Version~V16 source above is preserved byte for byte up to its
sole terminal document-closing command.  Its SHA--256 digest is
\begin{center}\scriptsize\ttfamily
e12e18360d59297b58122a9363fb69acf148066112ac816a03ceb7030a293f57.
\end{center}
For Version~V18, preserve this accumulated source, append new mathematical
material immediately before the final document-closing command, and increment
the filename to \texttt{\_V18.tex}.  Keep the fixed physical box, dyadic
lattice sequence, weak-facing selector, response law, exact pre-retuning
action, cutoff interaction, plaquette writer, rooted Hodge carrier, local and
topological support convention, colour normalization, and reference
normalization fixed.  At each cutoff, identify physical axial modes by their
integer Fourier index before reading any cofactor, stiffness, fourth moment,
likelihood card, symbol, ancestry, tail, or transport value.

% =============================================================================
% END APPENDED VERSION V17 MATERIAL
% =============================================================================


% =============================================================================
% BEGIN APPENDED VERSION V18 MATERIAL
% =============================================================================

\clearpage
\hypersetup{
  pdftitle={Renewal Yang--Mills Reduced-Fibre Wilson Shell Symbol and Local-Vacuum Programme V18},
  pdfsubject={Carrier-first physical indexing, resonance-centred toron labels, and a two-window certificate for the actual buffered high shell}
}
\section*{Version V18: carrier-first physical indexing, resonance-centred toron
labels, and the actual buffered-high axial card}
\addcontentsline{toc}{section}{Version V18: carrier-first physical indexing and the actual buffered-high axial card}

\begin{abstract}
Version~V17 correctly distinguished fixed physical momentum from a fixed
lattice angle on a fixed physical box.  It then promoted the index-$1$ branch
to the first cofinal YM--A card.  The inherited selected-shell programme has
an earlier carrier gate: before the reduced symbol is formed, the uniformly
finite physical low-loss space and its buffered transition band are retained,
and the direct Wilson symbol is tested on the rooted-Hodge buffered high
projection.  The first nonzero Fourier index need not survive that projection.

This continuation restores the carrier-first order.  At the central identity
background the exact first hard-high axial index is
\[
 n_{+,M}^{\rm hi}
 =1+\left\lfloor
    \frac{M}{\pi}
    \arcsin\!\left(\frac{L\Lambda_+}{2M}\right)
   \right\rfloor,
\]
whenever the outer physical threshold lies below the lattice Nyquist edge.
It stabilizes to
$1+\lfloor L\Lambda_+/(2\pi)\rfloor$.  Thus the Version~V17 index-$1$
sequence is the first high-carrier sequence only when the physical threshold
lies below the first box momentum; otherwise it is an exceptional low-loss
card and cannot decide the YM--A high-shell margin.

For a commuting toron channel, fixed Fourier index is not the invariant
physical label.  The finite-momentum modes are exactly the fixed integer
offsets from the nearest toron-resonance index.  After a stable subsequence,
the corresponding physical momenta converge to
$|2\pi j+\rho_\alpha|/L$, where $\rho_\alpha\in[-\pi,\pi]$ is the scaled
residual toron phase.  The hard-high/transition/low assignment stabilizes
away from the physical thresholds; equality with a threshold remains a
buffered transition row rather than being forced into either carrier.

The Version~V17 stiffness estimate is extended to every resonance-centred
bounded-physical-momentum sequence.  A separate finite-range counterexample
proves that passing any fixed finite physical-index bank, even with positive
stiffness and a finite fourth displacement moment, does not imply the uniform
high-carrier margin: a fixed-angle ultraviolet mode may still fail.  The
correct finite all-angle packet therefore has two windows: a low-angle
stiffness/fourth-moment certificate and the existing finite-head
Fej\'er--Riesz certificate away from the Ward zero.

No response cofactor, selected Wilson expectation, signed symbol coefficient,
or high-carrier margin is invented.  The first physical acquisition is now
the least mode which actually survives the buffered high projection, followed
by the away-from-zero head certificate.  The selected shell remains at
\stopstatus.
\end{abstract}

% =============================================================================
\section{The carrier gate precedes physical-momentum ordering}
\label{sec:YMA18-carrier-gate}
% =============================================================================

Retain the fixed physical box and the inherited coarse spacing
\begin{equation}
 a_{c,r}=\frac{L}{M_r}.
 \label{eq:YMA18-coarse-spacing}
\end{equation}
At a central background, the physical momentum of the axial Fourier index
$n$ is the Version~V17 quantity
\begin{equation}
 p_{M,n}
 :=\frac{2M}{L}\sin\frac{\pi n}{M},
 \qquad 0\le n\le\left\lfloor\frac M2\right\rfloor.
 \label{eq:YMA18-central-momentum}
\end{equation}

The inherited weak-shell construction fixes two physical thresholds
\begin{equation}
 0<\Lambda_-<\Lambda_+<\infty
 \label{eq:YMA18-thresholds}
\end{equation}
before any symbol value is read.  The hard outer low-loss projection is the
spectral projection onto physical Hodge energies at most $\Lambda_+^2$; its
hard high complement contains only modes with physical momentum strictly
larger than $\Lambda_+$.  The finite transition band between $\Lambda_-$ and
$\Lambda_+$ is retained by the buffered frame.  The final symbol is tested on
the still smaller rooted-Hodge, toron-resonance, and buffered-high projection
$P_{r,\sigma}^{\rm hi}(k)$.

\begin{theorem}[Carrier-first mode ordering]
\label{prin:YMA18-carrier-first}
A raw Fourier--colour character is ordered as a YM--A physical shell card only
after the following data have been frozen:
\begin{equation}
 \boxed{
 (a_{c,r},M_r,T_r,\sigma,k,
   E_{r,T_r}^-,F_{r,T_r},Q_{r,T_r}^+,
   P_{r,\sigma}^{\rm hi}(k),
   \widehat S_{0,r,\sigma}^{\rm hi}(k)).}
 \label{eq:YMA18-carrier-passport}
\end{equation}
In particular, a value on a character with
$P_{r,\sigma}^{\rm hi}(k)=0$ is not an entry of the global high-carrier loss
edge.  It may be a valid exceptional low-loss or transition-band card, but it
cannot be promoted to a high-shell obstruction or pass.
\end{theorem}

\begin{proof}
The retained high space is the orthogonal direct sum of the nonzero ranges of
$P_{r,\sigma}^{\rm hi}(k)$.  Compression by a zero projection contributes no
summand to either the physical Hessian or the reference.  The global normalized
loss edge is therefore the maximum only over nonzero retained blocks.  The
physical scale, toron phase, and reference support determine which raw
characters represent those blocks and must be fixed before their numerical
values are ordered.
\end{proof}

\begin{firewall}[Kinematic Fourier modes versus the YM--A carrier]
\label{fire:YMA18-kinematic-not-carrier}
Version~V17's fixed-index formulas are correct kinematic statements on the
full periodic lattice.  They do not by themselves prove that the index lies
in the inherited buffered high carrier.  Carrier membership is an earlier
physical-loading row.  A later stiffness, likelihood, Witten, boundary, or
tail theorem cannot repair a mode which was evaluated on the wrong side of
the low/high short.
\end{firewall}

% =============================================================================
\section{Exact central axial entrance index}
\label{sec:YMA18-central-threshold}
% =============================================================================

\begin{lemma}[Monotonicity of the central axial momentum]
\label{lem:YMA18-central-monotone}
For fixed $M$ and $L$, the map
$n\mapsto p_{M,n}$ is strictly increasing on
$0\le n\le M/2$.
\end{lemma}

\begin{proof}
The angle $\pi n/M$ belongs to $[0,\pi/2]$, where the sine is strictly
increasing.  Multiplication by the positive factor $2M/L$ preserves the
order.
\end{proof}

For a physical threshold $\Lambda>0$ with $L\Lambda<2M$, put
\begin{equation}
 x_{\Lambda,M}
 :=\frac{M}{\pi}
   \arcsin\!\left(\frac{L\Lambda}{2M}\right).
 \label{eq:YMA18-threshold-coordinate}
\end{equation}

\begin{theorem}[Exact first hard-high axial index]
\label{thm:YMA18-exact-first-high}
Assume $L\Lambda_+<2M$.  The smallest nonnegative axial index whose central
physical momentum lies strictly above the outer threshold is
\begin{equation}
 \boxed{
 n_{+,M}^{\rm hi}
 =\left\lfloor x_{\Lambda_+,M}\right\rfloor+1.}
 \label{eq:YMA18-exact-high-index}
\end{equation}
Equivalently,
\begin{equation}
 p_{M,n}\le\Lambda_+
 \quad(0\le n<n_{+,M}^{\rm hi}),
 \qquad
 p_{M,n_{+,M}^{\rm hi}}>\Lambda_+.
 \label{eq:YMA18-high-index-order}
\end{equation}
The indices in the hard buffered transition strip are exactly
\begin{equation}
 \boxed{
 \left\lfloor x_{\Lambda_-,M}\right\rfloor+1
 \le n\le
 \left\lfloor x_{\Lambda_+,M}\right\rfloor,}
 \label{eq:YMA18-transition-indices}
\end{equation}
with the interval empty when its lower endpoint exceeds its upper endpoint.
\end{theorem}

\begin{proof}
By \Cref{lem:YMA18-central-monotone}, the inequality
$p_{M,n}>\Lambda$ is equivalent to
\[
 \sin\frac{\pi n}{M}>\frac{L\Lambda}{2M}.
\]
Both arguments lie in $[0,\pi/2]$, so this is equivalent to
$n>x_{\Lambda,M}$.  The least integer satisfying a strict inequality
$n>x$ is $\lfloor x\rfloor+1$, which proves
\eqref{eq:YMA18-exact-high-index}--\eqref{eq:YMA18-high-index-order}.
The transition strip consists of the simultaneous inequalities
$p_{M,n}>\Lambda_-$ and $p_{M,n}\le\Lambda_+$; applying the same integer
calculus to the two endpoints gives \eqref{eq:YMA18-transition-indices}.
\end{proof}

\begin{theorem}[Eventual continuum entrance index with an explicit cutoff]
\label{thm:YMA18-eventual-first-high}
Define
\begin{equation}
 n_+
 :=1+\left\lfloor\frac{L\Lambda_+}{2\pi}\right\rfloor,
 \qquad
 \Delta_+
 :=\frac{2\pi n_+}{L}-\Lambda_+>0.
 \label{eq:YMA18-continuum-index-gap}
\end{equation}
If
\begin{equation}
 M\ge2n_+,
 \qquad
 M^2>
 \frac{\pi^3n_+^3}{3L\Delta_+},
 \label{eq:YMA18-index-stability-cutoff}
\end{equation}
then
\begin{equation}
 \boxed{n_{+,M}^{\rm hi}=n_+.}
 \label{eq:YMA18-eventual-index}
\end{equation}
In particular, the first hard-high index is eventually constant along every
fixed-box dyadic refinement.
\end{theorem}

\begin{proof}
Version~V17 proved
\begin{equation}
 0\le\frac{2\pi n}{L}-p_{M,n}
 \le\frac{\pi^3n^3}{3LM^2}
 \label{eq:YMA18-momentum-error-recall}
\end{equation}
for $2n\le M$.  Every $n<n_+$ obeys
$2\pi n/L\le\Lambda_+$ by the definition of $n_+$, and hence
$p_{M,n}\le\Lambda_+$.  At $n=n_+$, conditions
\eqref{eq:YMA18-index-stability-cutoff} and
\eqref{eq:YMA18-momentum-error-recall} give
\[
 p_{M,n_+}
 >\frac{2\pi n_+}{L}-\Delta_+
 =\Lambda_+.
\]
Monotonicity now proves \eqref{eq:YMA18-eventual-index}.
\end{proof}

\begin{corollary}[Exact scope of the Version V17 index-$1$ card]
\label{cor:YMA18-V17-scope}
The fixed-index branch $n=1$ is eventually the first central hard-high branch
if and only if
\begin{equation}
 \boxed{\Lambda_+<\frac{2\pi}{L}.}
 \label{eq:YMA18-index-one-condition}
\end{equation}
If $\Lambda_+\ge2\pi/L$, the index-$1$ mode belongs to the outer physical
low-loss carrier for every sufficiently fine lattice.  At the equality
$\Lambda_+=2\pi/L$, one has the strict finite-lattice inequality
$p_{M,1}<\Lambda_+$ for every $M>2$.
\end{corollary}

\begin{proof}
The first statement is $n_+=1$ in
\eqref{eq:YMA18-continuum-index-gap}.  At equality,
$\sin(\pi/M)<\pi/M$ for $M>2$, so
$p_{M,1}<2\pi/L=\Lambda_+$.
\end{proof}

\begin{assessment}[Correction to the V17 acquisition order]
\label{ass:YMA18-V17-correction}
Version~V17 is retained as a theorem about fixed physical Fourier indices.
Its designation of $n=1$ as the first YM--A cofinal card is conditional on
\eqref{eq:YMA18-index-one-condition} and on survival of every further rooted-
Hodge and stabilizer projection.  Without those checks, $n=1$ is an
exceptional low-loss calibration, not the first high-shell mode.  The actual
central hard-high scan begins at \eqref{eq:YMA18-exact-high-index}; if an
additional inherited projection kills that block, it proceeds to the least
later index with nonzero retained range.
\end{assessment}

% =============================================================================
\section{Resonance-centred physical labels on a commuting toron}
\label{sec:YMA18-toron-labels}
% =============================================================================

A fixed Fourier index is intrinsic only in a Cartan channel at the identity
background.  In a root channel of a commuting toron, retain the inherited
coarse translation phase
\begin{equation}
 \vartheta_{r,\alpha}\in[-\pi,\pi)
 \label{eq:YMA18-toron-phase}
\end{equation}
in the selected axial direction.  Let $m_{r,\alpha}\in\Z/M_r\Z$ be the
least index in the predeclared ordering minimizing
\begin{equation}
 \left|
  \left(
   \frac{2\pi m}{M_r}+\vartheta_{r,\alpha}
  \right)_{2\pi}
 \right|.
 \label{eq:YMA18-resonance-minimizer}
\end{equation}
Define
\begin{equation}
 \delta_{r,\alpha}
 :=\left(
   \frac{2\pi m_{r,\alpha}}{M_r}
   +\vartheta_{r,\alpha}
  \right)_{2\pi},
 \qquad
 \rho_{r,\alpha}:=M_r\delta_{r,\alpha}.
 \label{eq:YMA18-residual-phase}
\end{equation}
Nearest-grid-point geometry gives
\begin{equation}
 \boxed{
 |\delta_{r,\alpha}|\le\frac{\pi}{M_r},
 \qquad
 |\rho_{r,\alpha}|\le\pi.}
 \label{eq:YMA18-residual-phase-bound}
\end{equation}

For a fixed integer offset $j$, put
\begin{equation}
 n_{r,\alpha}(j)=m_{r,\alpha}+j\pmod{M_r}
 \label{eq:YMA18-offset-index}
\end{equation}
and let $q_{r,\alpha}(j)$ be the corresponding wrapped covariant angle.

\begin{theorem}[Resonance-centred finite-momentum theorem]
\label{thm:YMA18-resonance-centred}
Fix a root channel on a stable commuting-toron patch.
\begin{enumerate}[label=\textup{(R18.\arabic*)},leftmargin=5.5em]
\item For every fixed $j$, all sufficiently large $r$ satisfy
      \begin{equation}
       \boxed{
       M_rq_{r,\alpha}(j)
       =\rho_{r,\alpha}+2\pi j.}
       \label{eq:YMA18-resonance-angle}
      \end{equation}
      The physical momentum obeys
      \begin{equation}
       \boxed{
       p_{r,\alpha}(j)
       =\frac{2M_r}{L}
        \left|
         \sin\frac{\rho_{r,\alpha}+2\pi j}{2M_r}
        \right|.}
       \label{eq:YMA18-resonance-momentum}
      \end{equation}
\item Its continuum approximation has the explicit error
      \begin{equation}
       \boxed{
       0\le
       \frac{|\rho_{r,\alpha}+2\pi j|}{L}
       -p_{r,\alpha}(j)
       \le
       \frac{|\rho_{r,\alpha}+2\pi j|^3}
            {24LM_r^2}.}
       \label{eq:YMA18-resonance-error}
      \end{equation}
\item After passage to a subsequence on which
      $\rho_{r,\alpha}\to\rho_\alpha\in[-\pi,\pi]$,
      \begin{equation}
       \boxed{
       p_{r,\alpha}(j)
       \longrightarrow
       p_{\alpha,j}^{\infty}
       :=\frac{|\rho_\alpha+2\pi j|}{L}.}
       \label{eq:YMA18-resonance-limit}
      \end{equation}
\item Conversely, if a sequence of root-channel indices has bounded physical
      momentum, then its offsets from $m_{r,\alpha}$ are bounded.  After a
      further subsequence, it is one fixed offset $j$ and therefore has the
      form above.
\end{enumerate}
Thus finite physical momentum in a toron root channel is equivalent, after
subsequence extraction, to a fixed resonance-centred offset.  Fixed raw
Fourier index is the special central/Cartan case $m_{r,\alpha}=0$.
\end{theorem}

\begin{proof}
By \eqref{eq:YMA18-residual-phase-bound}, for fixed $j$ one has
$|\rho_{r,\alpha}+2\pi j|<\pi M_r$ for all sufficiently large $r$.  Hence no
additional $2\pi$ wrap occurs after adding $2\pi j/M_r$, and
\eqref{eq:YMA18-resonance-angle} follows.  Substitution into the inherited
physical covariant momentum gives \eqref{eq:YMA18-resonance-momentum}.

For $x\ge0$, $0\le x-\sin x\le x^3/6$.  Apply this with
$x=|\rho_{r,\alpha}+2\pi j|/(2M_r)$ and multiply by $2M_r/L$ to obtain
\eqref{eq:YMA18-resonance-error}.  Compactness of $[-\pi,\pi]$ gives a
convergent residual-phase subsequence and proves
\eqref{eq:YMA18-resonance-limit}.

For the converse, let $q_r\in[-\pi,\pi]$ be the wrapped covariant angle of a
mode with $p_r\le P$.  Since
$\sin(|q_r|/2)\ge|q_r|/\pi$,
\[
 M_r|q_r|\le\frac{\pi LP}{2}.
\]
The nearest-resonance angle satisfies
$M_r|\delta_{r,\alpha}|\le\pi$.  Choose the integer offset representative
$j_r$ for which
$q_r-\delta_{r,\alpha}=2\pi j_r/M_r$ once both angles are taken in their
small representatives.  Then
\[
 |j_r|
 \le\frac{M_r|q_r|+M_r|\delta_{r,\alpha}|}{2\pi}
 \le\frac{LP}{4}+\frac12.
\]
Only finitely many offsets occur, so one is constant after a further
subsequence.
\end{proof}

\begin{corollary}[Stable hard-high membership away from a threshold]
\label{cor:YMA18-toron-membership}
On a subsequence of \Cref{thm:YMA18-resonance-centred}, fix an offset $j$ and
suppose
\begin{equation}
 p_{\alpha,j}^{\infty}\ne\Lambda_+.
 \label{eq:YMA18-threshold-separation}
\end{equation}
Then the hard high/low membership of that mode is eventually constant and is
the sign of
$p_{\alpha,j}^{\infty}-\Lambda_+$.  If equality holds, the mode remains a
threshold-transition row: a hard assignment is not transported through an
arbitrary toron or cutoff perturbation, and the buffered frame must be used.
\end{corollary}

\begin{proof}
The convergence \eqref{eq:YMA18-resonance-limit} preserves every strict
inequality for sufficiently large $r$.  At equality, arbitrarily small changes
of the residual toron phase can move the physical momentum to either side of
the threshold, which is precisely the reason for retaining the buffered
transition band.
\end{proof}

\begin{definition}[First actual retained physical block]
\label{def:YMA18-first-actual-block}
On a stable patch, order the resonance-centred pairs $(\sigma,j)$ first by
their physical momentum and then by the predeclared finite tie-breaker.  The
\emph{first actual retained block} is the least pair for which
\begin{equation}
 p_{r,\sigma}(j)>\Lambda_+,
 \qquad
 E_{r,\sigma}^{\rm hi}(j)\ne\{0\}.
 \label{eq:YMA18-first-actual-condition}
\end{equation}
The first inequality is the hard physical carrier gate; the second retains
all further rooted-Hodge, resonance, stabilizer, and buffered support data.
\end{definition}

\begin{assessment}[Central reduction]
\label{ass:YMA18-central-reduction}
At the identity central representative, every adjoint root phase is zero, so
$m_{r,\alpha}=0$ and $\rho_{r,\alpha}=0$.  The resonance-centred labels reduce
to the ordinary fixed indices of Version~V17.  The new information is not a
change of Fourier convention: it is the requirement that the index first
survive the already frozen physical low/high and rooted-Hodge projections.
\end{assessment}

% =============================================================================
\section{A bounded-physical-angle stiffness certificate}
\label{sec:YMA18-bounded-angle}
% =============================================================================

The Version~V17 proof used only that the selected angle was small.  It can be
stated directly for every resonance-centred finite-momentum sequence.
Retain the fused kernel, stiffness, finite head, fourth stock, and omitted
Ward-relative tail from
\eqref{eq:YMA17-kernel-expansion}--\eqref{eq:YMA17-stiffness-stocks}.

\begin{theorem}[Exact certificate at an arbitrary small covariant angle]
\label{thm:YMA18-small-angle-certificate}
Let $q\in[-1,1]\setminus\{0\}$ be an actual retained covariant axial angle and
put
\begin{equation}
 E_{r,R}(q)
 :=\tau_{r,R}^{(2)}
   +\frac{q^2}{4}\mathcal M_{4,r,R}.
 \label{eq:YMA18-small-angle-error}
\end{equation}
Then the assembled physical margin
\begin{equation}
 \mathfrak m_r(q;\delta)
 :=h_r^{\rm loc}(q)-\delta\beta_rs_0(q)
 \label{eq:YMA18-general-margin}
\end{equation}
satisfies
\begin{equation}
 \boxed{
 \begin{aligned}
 \frac{\mathfrak m_r(q;\delta)}{\omega(q)}
 &\ge
 \kappa_{r,R}-\frac{\delta\beta_r}{3}-E_{r,R}(q),\\
 \frac{\mathfrak m_r(q;\delta)}{\omega(q)}
 &\le
 \kappa_{r,R}-\frac{\delta\beta_r}{3}+E_{r,R}(q)
 +\frac{23\delta\beta_r}{288}q^2.
 \end{aligned}}
 \label{eq:YMA18-general-margin-interval}
\end{equation}
A positive lower endpoint is a \pass{} for that actual retained block; a
negative upper endpoint is a \obstruction{} on its already frozen
momentum--colour polarization; otherwise the block is \stopstatus.
\end{theorem}

\begin{proof}
The Version~V1 fourth-moment estimate gives
\[
 \left|
  \frac{h_r^{\rm loc}(q)}{\omega(q)}-\kappa_{r,R}
 \right|
 \le\tau_{r,R}^{(2)}+\frac{q^2}{4}\mathcal M_{4,r,R}.
\]
Version~V17 proved
\[
 \frac13-\frac{23}{72}\sin^2\frac q2
 \le\frac{s_0(q)}{\omega(q)}\le\frac13.
\]
Since $\sin^2(q/2)\le q^2/4$, subtraction of the reference gives the two
endpoints in \eqref{eq:YMA18-general-margin-interval}.  The Ward factor is
strictly positive for $q\ne0$, so division preserves the sign.
\end{proof}

\begin{corollary}[Resonance-centred fixed-offset form]
\label{cor:YMA18-offset-certificate}
For the offset $j$ of \Cref{thm:YMA18-resonance-centred}, put
\begin{equation}
 \xi_{r,\alpha,j}:=\rho_{r,\alpha}+2\pi j.
 \label{eq:YMA18-scaled-angle}
\end{equation}
For all sufficiently large $r$,
$q_{r,\alpha}(j)=\xi_{r,\alpha,j}/M_r$ and
\begin{equation}
 \boxed{
 E_{r,R}(q_{r,\alpha}(j))
 =\tau_{r,R}^{(2)}
  +\frac{\xi_{r,\alpha,j}^{2}}{4M_r^2}
   \mathcal M_{4,r,R}.}
 \label{eq:YMA18-offset-error}
\end{equation}
If the offsets range over a fixed finite bank, the residual phases stay in
$[-\pi,\pi]$, and
$\sup_r\mathcal M_{4,r}/\beta_r<\infty$, then the fourth-moment part of the
error divided by $\beta_r$ vanishes uniformly on that bank.  The complete
error has the same property either for the full kernel, where
$\tau_{r,\infty}^{(2)}=0$, or for a head sequence $R_r$ satisfying
$\tau_{r,R_r}^{(2)}/\beta_r\to0$.  Under that additional tail exhaustion and
away from the exact threshold $\kappa_r/\beta_r=\delta/3$, the fused
stiffness decides every such actual finite-physical-momentum block.
\end{corollary}

\begin{proof}
Equation \eqref{eq:YMA18-resonance-angle} gives the first identity.  For a
finite offset bank, $|\xi_{r,\alpha,j}|$ has one uniform bound.  The second
term in \eqref{eq:YMA18-offset-error}, divided by $\beta_r$, is therefore
$O(M_r^{-2})$ under the stated fourth-stock hypothesis.  The tail term vanishes
in the full-kernel form or by the explicitly assumed finite-head exhaustion.
Apply \Cref{thm:YMA18-small-angle-certificate}.
\end{proof}

\begin{firewall}[A high threshold is not the Ward zero]
\label{fire:YMA18-threshold-not-zero}
The hard-high carrier begins at a fixed positive physical momentum
$\Lambda_+$, but its lattice angles still tend to zero like $a_{c,r}$ on a
fixed box.  The Ward-relative stiffness theorem is therefore relevant to the
near-threshold high modes.  This does not make the omitted physical low-loss
carrier disappear; it remains a separately retained finite-rank exceptional
space.
\end{firewall}

% =============================================================================
\section{A fixed physical bank cannot certify the complete high carrier}
\label{sec:YMA18-no-finite-bank}
% =============================================================================

The actual YM--A target is a floor on every retained high block.  Passing the
first retained finite-physical-momentum block is necessary, but it is not
sufficient.

\begin{countertheorem}[Positive stiffness and every fixed index do not control a
fixed-angle mode]
\label{cth:YMA18-fixed-bank-not-global}
Fix numbers $\delta\beta\ge0$ and
\begin{equation}
 \kappa>\frac{\delta\beta}{3}.
 \label{eq:YMA18-counter-stiffness}
\end{equation}
Choose $C>\kappa/4$ and define the real even finite-range Ward symbol
\begin{equation}
 \boxed{
 h_C(q)=\kappa\omega(q)-C\omega(q)^2.}
 \label{eq:YMA18-counter-symbol}
\end{equation}
Then:
\begin{enumerate}[label=\textup{(N18.\arabic*)},leftmargin=5.4em]
\item $h_C(0)=0$, its stiffness is exactly $\kappa$, and its fourth
      displacement stock is finite;
\item for every fixed finite index bank $1\le n\le n_0$,
      \begin{equation}
       h_C(q_{M,n})-\delta\beta s_0(q_{M,n})>0
       \label{eq:YMA18-counter-fixed-pass}
      \end{equation}
      for every sufficiently large $M$;
\item at the fixed angle $q=\pi$,
      \begin{equation}
       \boxed{
       h_C(\pi)-\delta\beta s_0(\pi)<0.}
       \label{eq:YMA18-counter-UV-fail}
      \end{equation}
\end{enumerate}
Thus no theorem based only on the stiffness and finitely many fixed physical
indices can imply the uniform high-carrier symbol margin, even within the
class of finite-range zero-row-sum kernels.
\end{countertheorem}

\begin{proof}
Using
$\omega(q)=2(1-\cos q)$ and
$\omega(q)^2=6-8\cos q+2\cos2q$, one may write
\[
 h_C(q)
 =2\alpha_1(\cos q-1)+2\alpha_2(\cos2q-1),
 \qquad
 \alpha_1=-\kappa+4C,
 \quad
 \alpha_2=-C.
\]
Hence the kernel has range two and finite fourth stock
$|\alpha_1|+16|\alpha_2|$.  Its stiffness is
\[
 -\alpha_1-4\alpha_2=\kappa.
\]
Moreover,
\[
 \frac{h_C(q)}{\omega(q)}=\kappa-C\omega(q).
\]
For every fixed $n$, $q_{M,n}\to0$ and
$s_0(q_{M,n})/\omega(q_{M,n})\to1/3$.  Condition
\eqref{eq:YMA18-counter-stiffness} therefore proves
\eqref{eq:YMA18-counter-fixed-pass}, uniformly on a fixed finite bank.  At
$q=\pi$, $\omega(\pi)=4$ and
$h_C(\pi)=4\kappa-16C<0$; subtracting the nonnegative reference gives
\eqref{eq:YMA18-counter-UV-fail}.
\end{proof}

\begin{assessment}[What the countertheorem means]
\label{ass:YMA18-counter-meaning}
The countertheorem is not a claim that the Wilson cylinder realizes the
particular kernel \eqref{eq:YMA18-counter-symbol}.  It proves a logical
boundary: the near-Ward-zero stiffness packet and the first physical-index
card cannot, by themselves, close a margin whose quantified carrier also
contains fixed-angle ultraviolet blocks.  The away-from-zero finite head/tail
certificate remains independently necessary unless a Wilson-specific theorem
controls that sector directly.
\end{assessment}

% =============================================================================
\section{A target-minimal two-window certificate for the central axial carrier}
\label{sec:YMA18-two-window}
% =============================================================================

The existing programme already contains both ingredients required by the
countertheorem: the low-angle fourth-moment estimate and the finite
Fej\'er--Riesz head certificate away from the Ward zero.  V18 assembles them
only after the actual carrier gate.

Fix $0<q_0\le1$ and a displacement head radius $R$.  Let
\begin{equation}
 g_{r,R}(q)
 :=\frac{h_{r,R}^{\rm loc}(q)}{\omega(q)}
 \label{eq:YMA18-head-quotient}
\end{equation}
be the contact-repaired finite-head quotient, and retain the exact Wilson
reference quotient
$r_0(q)=s_0(q)/\omega(q)$.  Define
\begin{align}
 \mathfrak l_{r,R}^{\rm near}(q_0)
 &:={}
 \kappa_{r,R}-\frac{\delta\beta_r}{3}
 -\tau_{r,R}^{(2)}
 -\frac{q_0^2}{4}\mathcal M_{4,r,R},
 \label{eq:YMA18-near-margin}\\
 \mathfrak l_{r,R}^{\rm far}(q_0)
 &:={}
 \inf_{q_0\le|q|\le\pi}
 \left[g_{r,R}(q)-\delta\beta_rr_0(q)\right]
 -\tau_{r,R}^{(2)}.
 \label{eq:YMA18-far-margin}
\end{align}
The finite infimum in \eqref{eq:YMA18-far-margin} is a trigonometric-polynomial
positivity problem and can be certified by the Version~V1 Gram/interval
procedure.

\begin{theorem}[Two-window all-angle axial pass]
\label{thm:YMA18-two-window-pass}
Assume the complete central axial local kernel has the declared zero-row-sum
split and the omitted tail satisfies the Version~V1 Ward-relative bound.  If
\begin{equation}
 \boxed{
 \mathfrak l_{r,R}^{\rm near}(q_0)>0,
 \qquad
 \mathfrak l_{r,R}^{\rm far}(q_0)>0,}
 \label{eq:YMA18-two-window-condition}
\end{equation}
then
\begin{equation}
 \boxed{
 h_r^{\rm loc}(q)
 >\delta\beta_rs_0(q)
 \qquad(0<|q|\le\pi).}
 \label{eq:YMA18-all-angle-pass}
\end{equation}
In particular every nonzero retained central axial high block passes.  The
near-window condition may be weakened to the actual hard-high angles only;
the displayed form is a convenient stronger certificate.
\end{theorem}

\begin{proof}
For $0<|q|\le q_0$, the Version~V1 fourth-moment estimate and
$r_0(q)\le1/3$ give
\[
 \frac{h_r^{\rm loc}(q)-\delta\beta_rs_0(q)}{\omega(q)}
 \ge\mathfrak l_{r,R}^{\rm near}(q_0)>0.
\]
For $q_0\le|q|\le\pi$, the complete quotient differs from the finite-head
quotient by at most $\tau_{r,R}^{(2)}$.  Definition
\eqref{eq:YMA18-far-margin} therefore gives the same strict positivity.
Multiplication by $\omega(q)>0$ proves
\eqref{eq:YMA18-all-angle-pass}.
\end{proof}

\begin{corollary}[Carrier-restricted three-way verdict]
\label{cor:YMA18-carrier-verdict}
On one selected central axial card, exactly one of the following mathematical
outcomes is certified by the populated data.
\begin{enumerate}[label=\textup{(C18.\arabic*)},leftmargin=5.4em]
\item \pass: the two-window lower bounds are positive, or direct outward
      intervals are positive on every nonzero retained block.
\item \obstruction: one actual block with
      $P_{r,\sigma}^{\rm hi}(k)\ne0$ has a negative upper margin endpoint;
      its eigenvector is the physical witness.
\item \stopstatus: a carrier passport, response cofactor, finite head,
      Ward-relative tail, transition-band assignment, or direct mode interval
      remains unresolved.
\end{enumerate}
Failure of either sufficient inequality in
\eqref{eq:YMA18-two-window-condition} is only \stopstatus{} unless a negative
direct interval is supplied.
\end{corollary}

\begin{proof}
The passing branch is \Cref{thm:YMA18-two-window-pass}.  A negative upper
endpoint on a nonzero retained block is a direct finite-dimensional Rayleigh
witness.  The remaining case lacks either a domain input or a one-sided sign
certificate, and therefore does not imply the opposite sign.
\end{proof}

% =============================================================================
\section{Revised execution order and proof ledger}
\label{sec:YMA18-ledger}
% =============================================================================

The noncircular execution order is now
\begin{equation}
 \boxed{
 \begin{gathered}
 \text{weak-facing trajectory branch}
 \\
 \Downarrow
 \\
 \text{buffered low/high carrier passport}
 \\
 \Downarrow
 \\
 \text{central or resonance-centred physical label}
 \\
 \Downarrow
 \\
 \text{first actual retained block and response-matched law}
 \\
 \Downarrow
 \\
 \text{near-threshold card plus away-from-zero head/tail}
 \\
 \Downarrow
 \\
 \text{uniform high-carrier margin or an attained pure mode}.
 \end{gathered}}
 \label{eq:YMA18-execution-order}
\end{equation}

\begingroup\small
\begin{longtable}{>{\raggedright\arraybackslash}p{0.28\textwidth}
                  >{\raggedright\arraybackslash}p{0.19\textwidth}
                  >{\raggedright\arraybackslash}p{0.45\textwidth}}
\caption{Version V18 proof-status ledger.}\label{tab:YMA18-ledger}\\
\toprule
\textbf{Row}&\textbf{Status}&\textbf{Exact scope}\\
\midrule
\endfirsthead
\toprule
\textbf{Row}&\textbf{Status}&\textbf{Exact scope}\\
\midrule
\endhead
Carrier-before-value rule&Proved exactly
&Only nonzero rooted-Hodge buffered-high symbol blocks enter the global
YM--A loss edge.\\[0.3em]
Central hard-high entrance&Proved exactly
&The finite-$M$ entrance index is
$1+\lfloor(M/\pi)\arcsin(L\Lambda_+/(2M))\rfloor$.\\[0.3em]
Continuum entrance index&Proved with explicit cutoff
&The first central hard-high index stabilizes to
$1+\lfloor L\Lambda_+/(2\pi)\rfloor$.\\[0.3em]
Version V17 index-$1$ status&Qualified exactly
&It is the first high card only below the first physical threshold and after
all further retained projections pass.\\[0.3em]
Toron physical indexing&Proved exactly
&Finite physical momentum is a fixed offset from the nearest resonance
index; fixed raw index is only the central/Cartan case.\\[0.3em]
Threshold transport&Proved away from equality
&Hard membership stabilizes off the threshold; equality remains in the
buffered transition row.\\[0.3em]
Small-angle mode interval&Proved exactly
&The stiffness/fourth-moment endpoints hold for every actual small covariant
angle, including toron-centred offsets.\\[0.3em]
Finite physical bank versus global carrier&Separated by countertheorem
&A range-two Ward kernel can pass every fixed index and fail at $q=\pi$.\\[0.3em]
Two-window axial certificate&Proved conditionally on populated heads
&Near-zero stiffness plus away-from-zero Fej\'er--Riesz head/tail positivity
closes the complete central axial angle bank.\\[0.3em]
Weak-facing branch&Open physical value
&No supplied card decides bounded roots versus an escaping selected sequence.\\[0.3em]
Response-matched selected law&Open physical value
&The response cofactors and actual final shell law remain unpopulated.\\[0.3em]
Near-threshold and UV symbol Reads&Open physical values
&Neither the first actual retained mode nor the away-from-zero signed head is
evaluated on the selected Wilson law.\\[0.3em]
Selected high-shell margin&\stopstatus
&The carrier and correct ordering are fixed, but the Wilson-specific values
are absent.\\[0.3em]
Local-vacuum contraction and mass&Not inferred
&No carrier-index theorem supplies higher jets, temporal transfer, continuum
landing, or a Yang--Mills mass gap.\\
\bottomrule
\end{longtable}
\endgroup

\begin{theorem}[Version V18 carrier-first correction]
\label{thm:YMA18-terminal}
Versions~V1--V17 are retained at their stated scopes.  Version~V18 additionally
proves:
\begin{enumerate}[label=\textup{(E18.\arabic*)},leftmargin=5.7em]
\item the physical low/high and rooted-Hodge carrier gate precedes every
      momentum ordering and every symbol value;
\item the exact central first hard-high index is
      \eqref{eq:YMA18-exact-high-index} and stabilizes to
      \eqref{eq:YMA18-eventual-index};
\item the Version~V17 index-$1$ card is the first high card only under
      \eqref{eq:YMA18-index-one-condition} and the remaining support checks;
\item on a commuting toron, finite physical momentum is exactly a fixed
      resonance-centred offset after subsequence extraction;
\item the stiffness/fourth-moment certificate extends to every such actual
      small-angle block;
\item a finite physical-index bank cannot replace the away-from-zero
      high-carrier certificate;
\item the two-window condition
      \eqref{eq:YMA18-two-window-condition} is a sufficient all-angle central
      axial margin theorem.
\end{enumerate}
No item decides the weak-facing branch, evaluates a response cofactor,
populates the selected Wilson symbol, or proves local-vacuum contraction or a
continuum Yang--Mills mass gap.
\end{theorem}

\begin{proof}
Items~\textup{(E18.1)--(E18.3)} are
\Cref{prin:YMA18-carrier-first,thm:YMA18-exact-first-high,thm:YMA18-eventual-first-high,cor:YMA18-V17-scope}.
Item~\textup{(E18.4)} is
\Cref{thm:YMA18-resonance-centred,cor:YMA18-toron-membership}.
Item~\textup{(E18.5)} is
\Cref{thm:YMA18-small-angle-certificate,cor:YMA18-offset-certificate}.
Item~\textup{(E18.6)} is
\Cref{cth:YMA18-fixed-bank-not-global}.
Item~\textup{(E18.7)} is
\Cref{thm:YMA18-two-window-pass}.
\end{proof}

\begin{openproblem}[V19: execute the first actual buffered-high card]
\label{op:YMA18-V19}
Preserve every frozen physical choice.  Do not automatically identify the
first nonzero Fourier index with the first retained high-shell block.  Proceed
in this order.
\begin{enumerate}[label=\textup{(V19.\arabic*)},leftmargin=5.4em]
\item Apply the weak-facing trajectory gate.  A bounded cofinal subsequence
      remains on the inherited rough-Haar/contact branch.
\item On an escaping shell, load the outer threshold $\Lambda_+$, buffered
      transition frame, rooted-Hodge projection, toron patch, and stabilizer
      before reading any mode value.
\item At the central identity, compute the exact entrance index
      $n_{+,M}^{\rm hi}$ from \eqref{eq:YMA18-exact-high-index}.  On a toron
      channel, use the resonance-centred offsets of
      \Cref{thm:YMA18-resonance-centred}.
\item Retain the least candidate only when its complete projection
      $P_{r,\sigma}^{\rm hi}(k)$ is nonzero.  If it is killed, continue in the
      predeclared physical-momentum ordering.
\item Acquire the two response cofactors and load the native response-matched
      law on that actual block.
\item Acquire the direct V16 likelihood card, or the fused stiffness and
      fourth-moment interval, on the first actual near-threshold block.
\item Independently populate the finite signed head and Ward-relative tail on
      $q_0\le|q|\le\pi$ and apply
      \Cref{thm:YMA18-two-window-pass}; a finite physical-index pass is not a
      substitute for this ultraviolet row.
\item Run the exact gauge held-out and keep the topological/harmonic packet
      separate before assigning a sign.
\item On a negative actual block, return its full carrier passport and
      polarization.  Open boundary/Witten ancestry only for attribution or
      cofinal transport of that frozen mode.
\end{enumerate}
If the carrier passport, response-matched law, first actual block, or
away-from-zero head/tail value is absent, the exact verdict remains
\stopstatus{} at that row.  Another theorem about a raw index outside the
retained carrier is not an advance.
\end{openproblem}

\begin{assessment}[V18 termination]
\label{ass:YMA18-termination}
The circularity in Version~V17's next attack has been removed.  Fixed-index
kinematics remain valid, but the physical low/high short determines which
index is a YM--A high-shell card.  The exact central entrance index and the
toron resonance-centred labels are now known.  The first retained finite-
physical-momentum block is only one part of the target: the finite-range
counterexample proves that the away-from-zero ultraviolet head must also be
controlled.  The selected Wilson margin remains \stopstatus{} because none of
those Wilson-specific values is supplied.  No negative physical mode has
been produced.
\end{assessment}

\section*{Version V18 append-only continuation record}
\addcontentsline{toc}{section}{Version V18 append-only continuation record}
The complete Version~V17 source above is preserved byte for byte up to its
sole terminal document-closing command.  Its SHA--256 digest is
\begin{center}\scriptsize\ttfamily
94fa145e1ec9f2139ad87bf58d4cd387662a94394b84e2342ef9e5b096959405.
\end{center}
For Version~V19, preserve this accumulated source, append new mathematical
material immediately before the final document-closing command, and increment
the filename to \texttt{\_V19.tex}.  Keep the fixed physical box, coarse
spacing, dyadic lattice sequence, weak-facing selector, buffered physical
thresholds, toron/stabilizer patch, rooted-Hodge high projection, response law,
pre-retuning action, cutoff interaction, plaquette writer, local/topological
support split, colour normalization, and reference normalization fixed.
Before reading any cofactor, likelihood, stiffness, head, tail, symbol,
ancestry, or transport value, verify that the proposed momentum--colour block
survives the complete carrier passport \eqref{eq:YMA18-carrier-passport}.

% =============================================================================
% END APPENDED VERSION V18 MATERIAL
% =============================================================================


% =============================================================================
% BEGIN APPENDED VERSION V19 MATERIAL
% =============================================================================

\clearpage
\hypersetup{
  pdftitle={Renewal Yang--Mills Reduced-Fibre Wilson Shell Symbol and Local-Vacuum Programme V19},
  pdfsubject={Volume-route separation, thermodynamic carrier grazing, and a guarded physical-annulus symbol certificate}
}
\section*{Version V19: volume-route separation, thermodynamic carrier grazing,
and the guarded physical-annulus card}
\addcontentsline{toc}{section}{Version V19: volume-route separation and the guarded physical-annulus card}

\begin{abstract}
Version~V18 correctly restored the buffered low/high carrier before ordering
Fourier--colour values.  Its quantitative continuation nevertheless retained
the Version~V17 hypothesis of one fixed physical box length $L$.  That is one
legitimate target: the shell-renormalized finite-volume toron operator.  It is
not the only target left by the inherited Yang--Mills frontier, which also
allows a uniformly local centre-neutral algebra followed by an
infinite-volume vacuum core.  On that branch the physical volume $L_r=M_ra_r$
must grow, and the fixed-index and fixed-resonance-offset conclusions of
Versions~V17--V18 no longer identify a nonzero physical momentum.

This continuation inserts the volume route into the carrier passport.  For a
central axial index $n$ it distinguishes the continuum wave number
$\kappa_{r,n}=2\pi n/L_r$ from the lattice momentum
$p_{r,n}=2a_r^{-1}\sin(a_r\kappa_{r,n}/2)$ and proves the uniform comparison
\[
 \frac2\pi\kappa_{r,n}\le p_{r,n}\le\kappa_{r,n},
 \qquad
 0\le\kappa_{r,n}-p_{r,n}
 \le\frac{a_r^2\kappa_{r,n}^3}{24}.
\]
Finite nonzero physical momentum is a fixed index only on a finite-volume
route.  On a thermodynamic route it requires
$n_r/L_r\to p/(2\pi)$ and hence $n_r\to\infty$, while the lattice angle still
tends to zero.  The same correction holds in a toron root channel: a fixed
offset from the nearest resonance gives a nonzero finite momentum only at
fixed volume; at growing volume the offset itself must grow linearly with
$L_r$.

The first hard-high central index is evaluated on an arbitrary volume route.
It obeys
\[
 n_{+,r}^{\rm hi}
 =1+\left\lfloor\frac{M_r}{\pi}
   \arcsin\!\left(\frac{a_r\Lambda_+}{2}\right)\right\rfloor,
 \qquad
 \frac{n_{+,r}^{\rm hi}}{L_r}\longrightarrow
 \frac{\Lambda_+}{2\pi}
\]
on an escaping thermodynamic branch.  Its physical overshoot is at most
$2\pi/L_r$ and therefore vanishes.  The first hard-high grid mode consequently
grazes the threshold and is not a uniformly separated carrier block.  A
fixed physical guard $\Delta>0$, or the full buffered transition projection,
is required before adjacent-cutoff mode transport.

A finite spectral-passport theorem makes this statement executable: a
Hermitian selector approximation with error $\varepsilon$ fixes every high
rank whose eigenvalue interval avoids $\Lambda_+^2$, while a crossing interval
is retained as a transition-band \stopstatus.  Carrier selection is held out
from the Wilson margin; a countertheorem shows that selecting the projection
after seeing the Hessian can manufacture either sign.

Finally, the small-angle stiffness theorem is rewritten uniformly on a
physical annulus.  For every central or toron mode with
$\Lambda_++\Delta\le p\le P$ and $a_rP\le2\sin(1/2)$, one stiffness, fourth-displacement
stock, and Ward-relative tail control the entire growing mode bank.  The error
is bounded by $a_r^2P^2\mathcal M_{4,r}$ rather than by an index-dependent
quantity.  Failure against a separated stiffness sign forces a normalized
fourth moment of order $a_r^{-2}$ or a nonvanishing tail.  The fixed-volume,
thermodynamic-local, and unspecified-volume branches are thereby separated
before any response cofactor or Wilson symbol is read.

No weak-facing root, volume route, response cofactor, selected-law kernel, or
symbol value is manufactured.  The physical margin remains at \stopstatus.
\end{abstract}

% =============================================================================
\section{The volume route is part of the physical carrier passport}
\label{sec:YMA19-volume-route}
% =============================================================================

Let $M_r\ge4$ be the coarse periodic side and let $a_r>0$ be its calibrated
coarse spacing.  The physical axial period is
\begin{equation}
 \boxed{L_r:=M_ra_r.}
 \label{eq:YMA19-volume-route}
\end{equation}
Versions~V17--V18 specialized to $L_r=L$.  In this continuation no asymptotic
condition on $L_r$ is imposed until the target branch is declared.

For a central axial index
$0\le n\le\lfloor M_r/2\rfloor$, put
\begin{equation}
 q_{r,n}:=\frac{2\pi n}{M_r},
 \qquad
 \kappa_{r,n}:=\frac{2\pi n}{L_r}=\frac{q_{r,n}}{a_r},
 \qquad
 p_{r,n}:=\frac2{a_r}\sin\frac{q_{r,n}}2.
 \label{eq:YMA19-three-momenta}
\end{equation}
Here $q_{r,n}$ is a dimensionless lattice angle,
$\kappa_{r,n}$ is the unwrapped continuum wave number attached to that
integer label, and $p_{r,n}$ is the physical lattice Hodge momentum.

\begin{theorem}[Exact lattice-to-physical momentum comparison]
\label{thm:YMA19-momentum-comparison}
For every index in \eqref{eq:YMA19-three-momenta},
\begin{equation}
 \boxed{
 \frac2\pi\,\kappa_{r,n}
 \le p_{r,n}\le\kappa_{r,n}.}
 \label{eq:YMA19-momentum-two-sided}
\end{equation}
Moreover,
\begin{equation}
 \boxed{
 0\le\kappa_{r,n}-p_{r,n}
 \le\frac{a_r^2\kappa_{r,n}^3}{24}.}
 \label{eq:YMA19-momentum-cubic-error}
\end{equation}
Consequently a family $(p_{r,n_r})$ is bounded if and only if
$(\kappa_{r,n_r})$ is bounded.  On every bank
$0\le\kappa_{r,n}\le P$,
\begin{equation}
 \sup_{\kappa_{r,n}\le P}
 |p_{r,n}-\kappa_{r,n}|
 \le\frac{a_r^2P^3}{24}.
 \label{eq:YMA19-uniform-momentum-error}
\end{equation}
\end{theorem}

\begin{proof}
Put $x=a_r\kappa_{r,n}/2=\pi n/M_r\in[0,\pi/2]$.  On this interval,
\[
 \frac{2x}{\pi}\le\sin x\le x,
 \qquad
 0\le x-\sin x\le\frac{x^3}{6}.
\]
Multiplying the first pair by $2/a_r$ gives
\eqref{eq:YMA19-momentum-two-sided}.  Multiplying the second inequality by
$2/a_r$ gives
\[
 0\le\kappa_{r,n}-p_{r,n}
 \le\frac2{a_r}\frac1{6}
       \left(\frac{a_r\kappa_{r,n}}2\right)^3
 =\frac{a_r^2\kappa_{r,n}^3}{24},
\]
which is \eqref{eq:YMA19-momentum-cubic-error}.  The boundedness equivalence
follows from the first pair; taking the supremum proves
\eqref{eq:YMA19-uniform-momentum-error}.
\end{proof}

\begin{definition}[Three geometric volume routes]
\label{def:YMA19-volume-branches}
After passage to a subsequence, every positive sequence $(L_r)$ belongs to
one of the following branches:
\begin{enumerate}[label=\textup{(V\arabic*)},leftmargin=3.8em]
\item $L_r\to L\in(0,\infty)$: the \emph{finite-volume route};
\item $L_r\to\infty$: the \emph{thermodynamic route};
\item $L_r\to0$: the \emph{collapsing-volume route}.
\end{enumerate}
A sequence with several cluster behaviours is classified separately on each
cofinal subsequence.  The route is an absolute length-calibration row, not a
property inferred from the dimensionless lattice symbols.
\end{definition}

\begin{theorem}[Finite physical momentum on the three volume routes]
\label{thm:YMA19-volume-classification}
Assume $a_r\to0$ and let $n_r\ge1$.
\begin{enumerate}[label=\textup{(C19.\arabic*)},leftmargin=5.3em]
\item On a finite-volume route $L_r\to L\in(0,\infty)$, bounded physical
      momentum forces $(n_r)$ to be bounded.  After a further subsequence,
      $n_r=n$ and
      \begin{equation}
       p_{r,n}\longrightarrow\frac{2\pi n}{L}.
       \label{eq:YMA19-fixed-volume-limit}
      \end{equation}
\item On a thermodynamic route $L_r\to\infty$, if
      $p_{r,n_r}\to p\in(0,\infty)$, then
      \begin{equation}
       \boxed{
       n_r\longrightarrow\infty,
       \qquad
       \frac{n_r}{L_r}\longrightarrow\frac{p}{2\pi},
       \qquad q_{r,n_r}\longrightarrow0.}
       \label{eq:YMA19-thermodynamic-index-density}
      \end{equation}
      Conversely, $n_r/L_r\to p/(2\pi)$ implies
      $p_{r,n_r}\to p$.
\item On a collapsing-volume route, every nonzero mode escapes:
      \begin{equation}
       \boxed{p_{r,n_r}\ge\frac4{L_r}\longrightarrow\infty.}
       \label{eq:YMA19-collapsing-volume-escape}
      \end{equation}
\end{enumerate}
Thus ``finite physical momentum equals fixed Fourier index'' is exact only on
a finite-volume route.
\end{theorem}

\begin{proof}
By \eqref{eq:YMA19-momentum-two-sided}, bounded $p_{r,n_r}$ implies bounded
$\kappa_{r,n_r}=2\pi n_r/L_r$.  If $L_r$ stays bounded above, the integers
$n_r$ are bounded; a constant subsequence and
\eqref{eq:YMA19-momentum-cubic-error} give
\eqref{eq:YMA19-fixed-volume-limit}.

On the thermodynamic route, the same comparison and
\eqref{eq:YMA19-momentum-cubic-error} give
$\kappa_{r,n_r}\to p$.  Hence
$n_r/L_r\to p/(2\pi)$.  Since $p>0$ and $L_r\to\infty$, one has
$n_r\to\infty$.  Also
$q_{r,n_r}=a_r\kappa_{r,n_r}\to0$.  The converse follows directly from
\eqref{eq:YMA19-momentum-cubic-error}.

Finally, $n_r\ge1$ gives
$\kappa_{r,n_r}\ge2\pi/L_r$.  The lower bound in
\eqref{eq:YMA19-momentum-two-sided} yields
$p_{r,n_r}\ge4/L_r$, proving the last branch.
\end{proof}

\begin{corollary}[Exact scope of the Version V17 fixed-index theorem]
\label{cor:YMA19-V17-scope}
The fixed-index continuation of Version~V17 and the eventual constant central
entrance index of Version~V18 are finite-volume results.  On a thermodynamic
route, every fixed positive index satisfies
\begin{equation}
 p_{r,n}\longrightarrow0
 \label{eq:YMA19-fixed-index-zero}
\end{equation}
and eventually belongs below every fixed positive hard-high threshold.
\end{corollary}

\begin{proof}
For fixed $n$, \eqref{eq:YMA19-momentum-two-sided} gives
$0\le p_{r,n}\le2\pi n/L_r$.  The right side tends to zero when
$L_r\to\infty$.
\end{proof}

\begin{firewall}[The volume route cannot be supplied by a later mode theorem]
\label{fire:YMA19-volume-loading}
A fixed-volume toron operator and a uniformly local infinite-volume vacuum
core are different physical targets.  The dimensionless symbol, response
cofactor, Witten current, boundary message, or adjacent-cutoff Gram does not
determine which limit of $L_r=M_ra_r$ is intended.  If the volume route is
absent, the carrier passport is incomplete before any mode sign is read.
\end{firewall}

% =============================================================================
\section{The first hard-high grid mode grazes a thermodynamic threshold}
\label{sec:YMA19-threshold-grazing}
% =============================================================================

Fix a physical hard-high threshold $\Lambda_+>0$.  Whenever
$a_r\Lambda_+<2$, put
\begin{equation}
 x_{+,r}
 :=\frac{M_r}{\pi}
   \arcsin\!\left(\frac{a_r\Lambda_+}{2}\right),
 \qquad
 n_{+,r}^{\rm hi}:=\lfloor x_{+,r}\rfloor+1.
 \label{eq:YMA19-general-first-index}
\end{equation}
This is the Version~V18 formula written without the fixed-volume
specialization $a_r=L/M_r$.

\begin{lemma}[A useful arcsine remainder]
\label{lem:YMA19-arcsine-remainder}
For $0\le y\le1/2$,
\begin{equation}
 \boxed{0\le\arcsin y-y\le\frac{y^3}{3}.}
 \label{eq:YMA19-arcsine-bound}
\end{equation}
\end{lemma}

\begin{proof}
The lower bound is immediate.  Moreover,
\[
 \arcsin y-y
 =\int_0^y
   \frac{t^2}
   {\sqrt{1-t^2}\,[1+\sqrt{1-t^2}]}
   \,\dd t.
\]
For $t\le1/2$ the denominator exceeds one, so the integrand is at most
$t^2$.  Integration gives the upper bound.
\end{proof}

\begin{theorem}[Hard-high index count and threshold overshoot]
\label{thm:YMA19-threshold-grazing}
Assume $a_r\Lambda_+\le1$.  Then
\begin{equation}
 \boxed{
 0< n_{+,r}^{\rm hi}
     -\frac{\Lambda_+L_r}{2\pi}
 \le
 1+\frac{\Lambda_+^3a_r^2L_r}{24\pi}.}
 \label{eq:YMA19-index-count-bound}
\end{equation}
The first hard-high physical momentum satisfies the exact overshoot bound
\begin{equation}
 \boxed{
 0<p_{r,n_{+,r}^{\rm hi}}-\Lambda_+
 \le\frac{2\pi}{L_r}.}
 \label{eq:YMA19-threshold-overshoot}
\end{equation}
Consequently, if $a_r\to0$ and $L_r\to\infty$,
\begin{equation}
 \boxed{
 \frac{n_{+,r}^{\rm hi}}{L_r}
 \longrightarrow\frac{\Lambda_+}{2\pi},
 \qquad
 n_{+,r}^{\rm hi}\longrightarrow\infty,
 \qquad
 p_{r,n_{+,r}^{\rm hi}}\longrightarrow\Lambda_+
 \text{ from above}.}
 \label{eq:YMA19-thermodynamic-first-high}
\end{equation}
\end{theorem}

\begin{proof}
Put $y=a_r\Lambda_+/2$.  Since $L_r=M_ra_r$,
\[
 x_{+,r}-\frac{\Lambda_+L_r}{2\pi}
 =\frac{M_r}{\pi}(\arcsin y-y).
\]
Apply \Cref{lem:YMA19-arcsine-remainder} and use
$M_ry^3/(3\pi)=\Lambda_+^3a_r^2L_r/(24\pi)$.  Since
$0<n_{+,r}^{\rm hi}-x_{+,r}\le1$, this proves
\eqref{eq:YMA19-index-count-bound}.

By the defining strict inequality,
$p_{r,n_{+,r}^{\rm hi}}>\Lambda_+$ while
$p_{r,n_{+,r}^{\rm hi}-1}\le\Lambda_+$.  The sine is one-Lipschitz, hence
\[
 p_{r,n+1}-p_{r,n}
 \le\frac2{a_r}\frac{\pi}{M_r}
 =\frac{2\pi}{L_r}.
\]
This proves \eqref{eq:YMA19-threshold-overshoot}.  Divide
\eqref{eq:YMA19-index-count-bound} by $L_r$ and use
$a_r\to0$, $L_r\to\infty$.  The remaining conclusions follow from the
positive threshold and the overshoot bound.
\end{proof}

\begin{corollary}[The threshold collar contains macroscopically many modes]
\label{cor:YMA19-collar-count}
Fix $\Delta>0$ and assume $a_r\to0$, $L_r\to\infty$, and
$a_r(\Lambda_++\Delta)\le1$.  Let
$N_r^{\rm col}(\Delta)$ be the number of central axial indices satisfying
\begin{equation}
 \Lambda_+<p_{r,n}\le\Lambda_++\Delta.
 \label{eq:YMA19-collar-definition}
\end{equation}
On a thermodynamic route,
\begin{equation}
 \boxed{
 \frac{N_r^{\rm col}(\Delta)}{L_r}
 \longrightarrow\frac{\Delta}{2\pi}.}
 \label{eq:YMA19-collar-density}
\end{equation}
In particular, the physical buffer is not a bounded list of exceptional
indices in the infinite-volume limit.
\end{corollary}

\begin{proof}
The exact count is
\[
 N_r^{\rm col}(\Delta)
 =\left\lfloor\frac{M_r}{\pi}
   \arcsin\!\left(\frac{a_r(\Lambda_++\Delta)}2\right)\right\rfloor
  -\left\lfloor\frac{M_r}{\pi}
   \arcsin\!\left(\frac{a_r\Lambda_+}2\right)\right\rfloor.
\]
Apply \eqref{eq:YMA19-index-count-bound} separately at the two thresholds,
divide by $L_r$, and use that a floor changes a real number by less than one.
\end{proof}

\begin{assessment}[Why the first high grid mode is not a stable cofinal card]
\label{ass:YMA19-first-mode-grazes}
At fixed volume, a threshold separated from the discrete continuum spectrum
can select one eventually constant first high index.  At growing volume,
\eqref{eq:YMA19-threshold-overshoot} forces the first high grid mode to approach
the carrier boundary.  An $o(1)$ perturbation of the toron phase, physical
length calibration, or selector matrix can change its membership.  The
cofinal thermodynamic card must therefore use either the complete buffered
transition projection or a fixed guard away from the threshold; it cannot
transport the identity of the first strict grid point as though it had a
uniform spectral gap.
\end{assessment}

% =============================================================================
\section{A finite held-out carrier-passport certificate}
\label{sec:YMA19-passport-certificate}
% =============================================================================

Let $\mathcal E$ be one finite momentum--colour block after every frozen
algebraic support, rooted-gauge, Hodge, and stabilizer restriction except the
final physical low/high split.  Let
\begin{equation}
 A=A^*=A_{r,\sigma}(k)
 \label{eq:YMA19-selector}
\end{equation}
be the independently loaded physical selector on $\mathcal E$.  The hard-high
projection is
\begin{equation}
 P^{\rm hi}:=\one_{(\Lambda_+^2,\infty)}(A).
 \label{eq:YMA19-true-high-projection}
\end{equation}
At the central identity, $A=p_{r,n}^2I$ on the irreducible axial block; on a
toron or a smaller stabilizer block it may be matrix valued.

\begin{theorem}[Outward spectral-passport theorem]
\label{thm:YMA19-spectral-passport}
Suppose a Hermitian approximation $\widehat A$ satisfies
\begin{equation}
 \norm{A-\widehat A}\le\varepsilon_A.
 \label{eq:YMA19-selector-error}
\end{equation}
Let $\widehat\lambda_1\le\cdots\le\widehat\lambda_d$ be the ordered
eigenvalues of $\widehat A$.
\begin{enumerate}[label=\textup{(S19.\arabic*)},leftmargin=5.3em]
\item If
      $\widehat\lambda_j-\varepsilon_A>\Lambda_+^2$, then the corresponding
      ordered true eigenvalue is hard high.
\item If
      $\widehat\lambda_j+\varepsilon_A\le\Lambda_+^2$, then the corresponding
      ordered true eigenvalue is hard low.
\item If every interval
      $[\widehat\lambda_j-\varepsilon_A,
        \widehat\lambda_j+\varepsilon_A]$
      lies strictly on one side of $\Lambda_+^2$, the rank of
      $P^{\rm hi}$ is exactly the number of intervals above the threshold.
\item If some interval meets the threshold, that direction is a buffered
      transition row.  Its hard membership is \stopstatus{} at the submitted
      precision; it is not assigned by the sign of the Wilson Hessian.
\end{enumerate}
\end{theorem}

\begin{proof}
Weyl's Hermitian perturbation inequality gives
$|\lambda_j(A)-\widehat\lambda_j|\le\varepsilon_A$ for every ordered
eigenvalue.  The four conclusions follow by interval separation and the
definition \eqref{eq:YMA19-true-high-projection}.
\end{proof}

\begin{proposition}[Stable projection orientation under a selector gap]
\label{prop:YMA19-projection-angle}
Assume
\begin{equation}
 \widehat g
 :=\dist(\Lambda_+^2,\spec\widehat A)
 >\varepsilon_A.
 \label{eq:YMA19-approx-selector-gap}
\end{equation}
Let
$\widehat P^{\rm hi}=\one_{(\Lambda_+^2,\infty)}(\widehat A)$.  Then
$P^{\rm hi}$ and $\widehat P^{\rm hi}$ have the same rank and
\begin{equation}
 \boxed{
 \norm{P^{\rm hi}-\widehat P^{\rm hi}}
 \le\frac{\varepsilon_A}{2\widehat g-\varepsilon_A}<1.}
 \label{eq:YMA19-projection-angle-bound}
\end{equation}
\end{proposition}

\begin{proof}
Rank equality is \Cref{thm:YMA19-spectral-passport}.  Put
$P=P^{\rm hi}$ and $Q=\widehat P^{\rm hi}$.  Weyl's inequality places the
true low spectrum below
$\Lambda_+^2-(\widehat g-\varepsilon_A)$ and the approximate high spectrum
above $\Lambda_+^2+\widehat g$.  On the cross operator
$X=(I-P)Q$, the identity $\widehat A=A+(\widehat A-A)$ gives the Sylvester
equation
\[
 A_-X-X\widehat A_+
 =-(I-P)(\widehat A-A)Q,
\]
where the two spectra are separated by at least
$2\widehat g-\varepsilon_A$.  The spectral theorem, or the standard integral
formula for the inverse Sylvester map, gives
\[
 \norm{(I-P)Q}
 \le\frac{\varepsilon_A}{2\widehat g-\varepsilon_A}.
\]
The same argument with the high and low spaces interchanged gives the
identical bound for $P(I-Q)$.  For two orthogonal projections of equal finite
rank,
$\norm{P-Q}$ is the larger of these two cross norms.  This proves
\eqref{eq:YMA19-projection-angle-bound}.
\end{proof}

\begin{definition}[Guarded high carrier]
\label{def:YMA19-guarded-carrier}
For a fixed physical guard $\Delta>0$, the guarded high projection is
\begin{equation}
 P^{\rm ghi}
 :=\one_{((\Lambda_++\Delta)^2,\infty)}(A).
 \label{eq:YMA19-guarded-projection}
\end{equation}
The band between $\Lambda_+$ and $\Lambda_++\Delta$ remains in the declared
buffer.  The guard is fixed before any Wilson margin value is inspected.
\end{definition}

\begin{corollary}[A guard converts calibration error into stable membership]
\label{cor:YMA19-guard-stability}
If the selector is the physical Hodge energy and the error in the physical
momentum of a candidate is at most $\Delta/2$, then every mode with certified
momentum at least $\Lambda_++\Delta$ remains hard high under that error.
Equivalently, for a matrix selector, every approximate spectral interval
lying above $(\Lambda_++\Delta)^2$ and separated from $\Lambda_+^2$ by its
outward error has stable high rank under
\Cref{thm:YMA19-spectral-passport}.
\end{corollary}

\begin{countertheorem}[Postselecting the carrier can manufacture either sign]
\label{cth:YMA19-no-postselection}
Let $M=M^*$ have at least one positive and one negative eigenvalue.  If the
``retained carrier'' is allowed to be chosen after $M$ is observed, there are
nonzero orthogonal projections $P_+$ and $P_-$ such that
\begin{equation}
 P_+MP_+\succ0,
 \qquad
 P_-MP_-\prec0.
 \label{eq:YMA19-postselected-signs}
\end{equation}
Hence neither a pass nor an obstruction is physically meaningful unless the
selector, threshold, buffer, and support projection are held out from the
Wilson margin card.
\end{countertheorem}

\begin{proof}
Take $P_+$ and $P_-$ to be the spectral projections onto any positive and any
negative eigenspace of $M$, respectively.  Their compressions have the stated
strict signs.  The construction uses the target matrix itself and therefore
has no held-out carrier content.
\end{proof}

\begin{firewall}[Rank stability is not yet the Wilson margin]
\label{fire:YMA19-passport-not-margin}
\Cref{thm:YMA19-spectral-passport,prop:YMA19-projection-angle} certify the
physical subspace on which the signed Wilson card must be evaluated.  They do
not assign the sign of direct action curvature minus score covariance.  A
projection interval which crosses the threshold is a carrier \stopstatus,
not a negative Yang--Mills mode.
\end{firewall}

% =============================================================================
\section{Toron momentum labels on a growing physical volume}
\label{sec:YMA19-toron-volume}
% =============================================================================

Retain the Version~V18 residual phase
$\rho_{r,\alpha}\in[-\pi,\pi]$ in a stable root channel.  For an integer
offset $j_r$ whose representative satisfies
$|\rho_{r,\alpha}+2\pi j_r|\le\pi M_r$, put
\begin{equation}
 q_{r,\alpha}(j_r)
 :=\frac{\rho_{r,\alpha}+2\pi j_r}{M_r},
 \qquad
 \kappa_{r,\alpha}(j_r)
 :=\frac{|\rho_{r,\alpha}+2\pi j_r|}{L_r},
 \label{eq:YMA19-toron-kappa}
\end{equation}
\begin{equation}
 p_{r,\alpha}(j_r)
 :=\frac2{a_r}
   \left|\sin\frac{q_{r,\alpha}(j_r)}2\right|.
 \label{eq:YMA19-toron-p}
\end{equation}

\begin{theorem}[Volume-relative toron label theorem]
\label{thm:YMA19-toron-volume-label}
Assume $a_r\to0$.  For every admissible offset,
\begin{equation}
 \boxed{
 \frac2\pi\kappa_{r,\alpha}(j_r)
 \le p_{r,\alpha}(j_r)
 \le\kappa_{r,\alpha}(j_r),}
 \label{eq:YMA19-toron-comparison}
\end{equation}
\begin{equation}
 \boxed{
 0\le
 \kappa_{r,\alpha}(j_r)-p_{r,\alpha}(j_r)
 \le
 \frac{a_r^2\kappa_{r,\alpha}(j_r)^3}{24}.}
 \label{eq:YMA19-toron-error}
\end{equation}
Consequently:
\begin{enumerate}[label=\textup{(T19.\arabic*)},leftmargin=5.3em]
\item on a finite-volume route, bounded physical momentum forces bounded
      offsets, and Version~V18's fixed-offset subsequence theorem applies;
\item on a thermodynamic route, if
      $p_{r,\alpha}(j_r)\to p\in(0,\infty)$, then
      \begin{equation}
       \boxed{
       |j_r|\longrightarrow\infty,
       \qquad
       \frac{|j_r|}{L_r}\longrightarrow\frac{p}{2\pi},
       \qquad
       q_{r,\alpha}(j_r)\longrightarrow0.}
       \label{eq:YMA19-growing-toron-offset}
      \end{equation}
      After a sign-stable subsequence,
      $j_r/L_r\to\pm p/(2\pi)$.
\item every fixed resonance-centred offset satisfies
      $p_{r,\alpha}(j)\to0$ when $L_r\to\infty$.
\end{enumerate}
Thus ``finite toron momentum equals fixed resonance offset'' is also a
finite-volume statement.
\end{theorem}

\begin{proof}
Equations \eqref{eq:YMA19-toron-comparison}--
\eqref{eq:YMA19-toron-error} are the proof of
\Cref{thm:YMA19-momentum-comparison} applied to
$x=|q_{r,\alpha}(j_r)|/2$.

If $L_r$ is bounded above, bounded $p$ gives bounded $\kappa$ and hence
\[
 |j_r|
 \le\frac{L_r\kappa_{r,\alpha}(j_r)+\pi}{2\pi},
\]
so only finitely many offsets occur.  On a thermodynamic route,
\eqref{eq:YMA19-toron-error} gives $\kappa_{r,\alpha}(j_r)\to p$.  Since
$|\rho_{r,\alpha}|/L_r\to0$,
\[
 \frac{2\pi|j_r|}{L_r}\longrightarrow p.
\]
The positive limit forces $|j_r|\to\infty$.  Also
$|q_{r,\alpha}(j_r)|=a_r\kappa_{r,\alpha}(j_r)\to0$.
A subsequence fixes the sign of $j_r$.  For fixed $j$, the numerator in
\eqref{eq:YMA19-toron-kappa} is bounded while $L_r\to\infty$, proving the
last claim.
\end{proof}

\begin{assessment}[Correct toron transport variable]
\label{ass:YMA19-toron-density}
On a finite physical torus, a fixed resonance-centred offset is the natural
cofinal label.  On a thermodynamic route the natural label is instead the
physical offset density $j_r/L_r$, or equivalently the unwrapped covariant
wave number $\kappa_{r,\alpha}$.  The residual phase remains an order-one
boundary correction and disappears after division by $L_r$ on a nonzero
finite-momentum branch.
\end{assessment}

% =============================================================================
\section{One stiffness packet controls a growing bounded physical annulus}
\label{sec:YMA19-annulus}
% =============================================================================

Let $h_r^{\rm loc}(q)$ be the fully assembled local scalar axial symbol on an
actual retained block, with all deterministic, connected-score, generated,
and local support-birth terms fused before the displacement kernel is formed.
Retain the Version~V18 finite-head quantities
\begin{equation}
 \kappa_{r,R},
 \qquad
 \mathcal M_{4,r,R},
 \qquad
 \tau_{r,R}^{(2)},
 \label{eq:YMA19-head-stocks}
\end{equation}
and the exact Wilson reference $s_0(q)=\omega(q)r_0(q)$.

\begin{lemma}[A bounded physical momentum has uniformly small lattice angle]
\label{lem:YMA19-angle-from-p}
If $0\le p\le P$ and $a_rP\le2\sin(1/2)$, then the small representative of the
corresponding central or toron covariant angle satisfies
\begin{equation}
 \boxed{|q|\le1,\qquad |q|\le2a_rP.}
 \label{eq:YMA19-angle-bound}
\end{equation}
\end{lemma}

\begin{proof}
The exact dispersion relation is
$|q|=2\arcsin(a_rp/2)$.  The hypothesis gives
$a_rp/2\le\sin(1/2)$ and therefore $|q|\le1$.  Also
$\sin(1/2)<1/2$ and, for $0\le y\le1/2$,
$\arcsin y\le2y$; this follows, for example, because its derivative is at
most $2$ on that interval and both functions vanish at zero.  Substitution
gives the second inequality in \eqref{eq:YMA19-angle-bound}.
\end{proof}

\begin{theorem}[Uniform physical-annulus margin interval]
\label{thm:YMA19-annulus-certificate}
Fix $P>\Lambda_++\Delta>0$ and assume
$a_rP\le2\sin(1/2)$.  For every actual retained
central or toron axial block whose physical momentum satisfies
\begin{equation}
 \Lambda_++\Delta\le p\le P,
 \label{eq:YMA19-guarded-annulus}
\end{equation}
the normalized target margin obeys
\begin{equation}
 \boxed{
 \frac{h_r^{\rm loc}(q)-\delta\beta_rs_0(q)}{\omega(q)}
 \ge
 \kappa_{r,R}-\frac{\delta\beta_r}{3}
 -\tau_{r,R}^{(2)}
 -a_r^2P^2\mathcal M_{4,r,R}.}
 \label{eq:YMA19-annulus-lower}
\end{equation}
It also obeys
\begin{equation}
 \boxed{
 \begin{aligned}
 \frac{h_r^{\rm loc}(q)-\delta\beta_rs_0(q)}{\omega(q)}
 \le{}&
 \kappa_{r,R}-\frac{\delta\beta_r}{3}
 +\tau_{r,R}^{(2)}
 +a_r^2P^2\mathcal M_{4,r,R}\\
 &+\frac{23\delta\beta_r}{72}a_r^2P^2.
 \end{aligned}}
 \label{eq:YMA19-annulus-upper}
\end{equation}
The bounds are uniform over the entire annulus and do not depend on the
number of lattice indices or toron offsets it contains.
\end{theorem}

\begin{proof}
Version~V18 proved, for every actual small representative $|q|\le1$,
\[
 \frac{h_r^{\rm loc}(q)-\delta\beta_rs_0(q)}{\omega(q)}
 \ge
 \kappa_{r,R}-\frac{\delta\beta_r}{3}
 -\tau_{r,R}^{(2)}-\frac{q^2}{4}\mathcal M_{4,r,R},
\]
and the matching upper bound with the additional term
$23\delta\beta_rq^2/288$.  The hypothesis and
\Cref{lem:YMA19-angle-from-p} give $|q|\le1$, so the Version~V18
estimate applies, and also $q^2/4\le a_r^2P^2$.  The reference correction satisfies
$23\delta\beta_rq^2/288
 \le23\delta\beta_ra_r^2P^2/72$.
Substitution gives the two displayed bounds.  No counting estimate is used.
\end{proof}

\begin{corollary}[Thermodynamic bounded-momentum alternative]
\label{cor:YMA19-thermodynamic-annulus}
Assume $a_r\to0$, $L_r\to\infty$, and choose fixed
$0<\Delta$ and $P>\Lambda_++\Delta$.  Suppose heads $R_r$ are chosen so that
\begin{equation}
 \frac{\tau_{r,R_r}^{(2)}}{\beta_r}\longrightarrow0,
 \qquad
 a_r^2\frac{\mathcal M_{4,r,R_r}}{\beta_r}
 \longrightarrow0.
 \label{eq:YMA19-annulus-errors-vanish}
\end{equation}
Then:
\begin{enumerate}[label=\textup{(A19.\arabic*)},leftmargin=5.3em]
\item if
      \begin{equation}
       \liminf_{r\to\infty}
       \left(\frac{\kappa_{r,R_r}}{\beta_r}-\frac\delta3\right)>0,
       \label{eq:YMA19-positive-stiffness-gap}
      \end{equation}
      every actual retained mode in
      $[\Lambda_++\Delta,P]$ is eventually \pass, uniformly over that
      growing bank;
\item if
      \begin{equation}
       \limsup_{r\to\infty}
       \left(\frac{\kappa_{r,R_r}}{\beta_r}-\frac\delta3\right)<0,
       \label{eq:YMA19-negative-stiffness-gap}
      \end{equation}
      every such mode is eventually a \obstruction{} for the target margin;
\item on the remaining branch, the normalized stiffness approaches the
      target, or one of the two error stocks fails to vanish.
\end{enumerate}
A uniform bound on
$\mathcal M_{4,r,R_r}/\beta_r$ is a sufficient stronger condition for the
second limit in \eqref{eq:YMA19-annulus-errors-vanish}.
\end{corollary}

\begin{proof}
Divide \eqref{eq:YMA19-annulus-lower} and
\eqref{eq:YMA19-annulus-upper} by $\beta_r$.  The two error assumptions and
$a_r^2\to0$ remove every displayed debit uniformly on the annulus.  A strict
positive or negative limiting stiffness gap therefore fixes the sign for all
sufficiently large $r$.  The final statement is immediate.
\end{proof}

\begin{corollary}[Disagreement forces an $a_r^{-2}$ fourth-moment scale]
\label{cor:YMA19-annulus-disagreement}
Suppose at one cutoff
\begin{equation}
 \frac{\kappa_{r,R}}{\beta_r}-\frac\delta3\ge\eta>0
 \label{eq:YMA19-positive-finite-gap}
\end{equation}
and an actual retained mode with $p\le P$ has nonpositive target margin.
Then
\begin{equation}
 \boxed{
 a_r^2P^2\frac{\mathcal M_{4,r,R}}{\beta_r}
 +\frac{\tau_{r,R}^{(2)}}{\beta_r}
 \ge\eta.}
 \label{eq:YMA19-annulus-disagreement-lower}
\end{equation}
In particular, if
$\tau_{r,R}^{(2)}/\beta_r\le\eta/2$, then
\begin{equation}
 \boxed{
 \frac{\mathcal M_{4,r,R}}{\beta_r}
 \ge\frac{\eta}{2a_r^2P^2}.}
 \label{eq:YMA19-a-minus-two-growth}
\end{equation}
Thus a bounded-physical-momentum sign which disagrees with a separated
stiffness floor must be carried by a physical displacement stock reaching the
cutoff scale, or by an unresolved Ward-relative tail.
\end{corollary}

\begin{proof}
If the margin is nonpositive, the lower endpoint in
\eqref{eq:YMA19-annulus-lower} cannot be strictly positive.  Divide by
$\beta_r$ and rearrange to obtain
\eqref{eq:YMA19-annulus-disagreement-lower}.  The second assertion follows by
subtracting the assumed tail bound.
\end{proof}

\begin{firewall}[The bounded annulus is not the ultraviolet completion]
\label{fire:YMA19-annulus-not-UV}
\Cref{thm:YMA19-annulus-certificate} controls every mode in one fixed physical
annulus, even though the number of indices grows with $L_r$.  It does not
control modes whose physical momentum tends to infinity.  The mesoscopic and
fixed-angle ultraviolet branches still require the away-from-zero signed
head/tail or another Wilson-specific uniform estimate.  A bounded-annulus
pass is not promoted to the complete high-carrier margin.
\end{firewall}

% =============================================================================
\section{Corrected finite-volume versus local-infinite-volume target split}
\label{sec:YMA19-target-split}
% =============================================================================

\begin{theorem}[Volume-route-first YM--A alternative]
\label{thm:YMA19-volume-route-alternative}
For every cofinal selected-shell family, apply the following ordered
alternative before reading a Wilson symbol.
\begin{enumerate}[label=\textup{(R19.\arabic*)},leftmargin=5.3em]
\item \emph{Unloaded volume route.}  If $a_r$, $M_r$, and the intended
      asymptotics of $L_r=M_ra_r$ are not jointly calibrated, the physical
      momentum carrier is \stopstatus.  Neither fixed indices nor fixed
      offsets are assigned a cofinal meaning.
\item \emph{Finite-volume route.}  If $L_r\to L\in(0,\infty)$,
      bounded physical modes reduce after subsequence extraction to fixed
      indices or fixed resonance offsets.  When $L_r\equiv L$, the formulas
      of Versions~V17--V18 apply verbatim; for a varying convergent length,
      use \Cref{thm:YMA19-momentum-comparison} together with the held-out
      spectral passport of \Cref{thm:YMA19-spectral-passport}.
\item \emph{Thermodynamic local route.}  If $L_r\to\infty$, label bounded
      physical momentum by $n_r/L_r$ or the toron offset density
      $j_r/L_r$.  The first strict high grid mode is threshold-grazing and is
      not transported as an isolated uniformly separated block.  Use the
      complete buffer or a fixed guard $\Delta>0$, and apply
      \Cref{thm:YMA19-annulus-certificate} on every bounded physical annulus.
\item \emph{Ultraviolet completion.}  Independently of the first three
      branches, modes with unbounded physical momentum require the far-window
      signed head/tail certificate of Version~V18 or another directly
      populated Wilson estimate.
\end{enumerate}
Only a negative assembled interval on a nonzero block of the already certified
carrier is a Yang--Mills pure-mode obstruction.
\end{theorem}

\begin{proof}
The first three branches are
\Cref{thm:YMA19-volume-classification,thm:YMA19-threshold-grazing,thm:YMA19-toron-volume-label,thm:YMA19-spectral-passport,thm:YMA19-annulus-certificate}.
The ultraviolet separation is
\Cref{fire:YMA19-annulus-not-UV} together with the Version~V18 finite-bank
counterexample.  The final sign rule is the carrier-first principle already
proved in Version~V18.
\end{proof}

\begin{assessment}[What has moved upstream]
\label{ass:YMA19-upstream-progress}
The first missing datum is no longer described merely as ``the first retained
mode.''  Before such a mode exists cofinally, the programme must select one of
the two inherited geometric targets:
\[
 \boxed{
 \begin{gathered}
 \text{shell-renormalized finite-volume toron operator}\
 \text{or}\
 \text{uniformly local thermodynamic vacuum core}
 \end{gathered}}
\]
On the first branch, after the finite physical length and any threshold
separation have been calibrated, the fixed-index passport is the correct
finite card.  On the second, the first high index diverges and grazes the threshold;
the correct target is a guarded physical annulus plus a separate ultraviolet
completion.  This is an absolute geometry/loading distinction and cannot be
repaired by another response, current, boundary, or symbol theorem.
\end{assessment}

% =============================================================================
\section{Version V19 proof ledger and next literal acquisition}
\label{sec:YMA19-ledger}
% =============================================================================

\begingroup\small
\begin{longtable}{>{\raggedright\arraybackslash}p{0.28\textwidth}
                  >{\raggedright\arraybackslash}p{0.19\textwidth}
                  >{\raggedright\arraybackslash}p{0.45\textwidth}}
\caption{Version V19 proof-status ledger.}\label{tab:YMA19-ledger}\\
\toprule
\textbf{Row}&\textbf{Status}&\textbf{Exact scope}\\
\midrule
\endfirsthead
\toprule
\textbf{Row}&\textbf{Status}&\textbf{Exact scope}\\
\midrule
\endhead
Volume route in carrier passport&Proved necessary
&The physical momentum labels depend on the absolute route
$L_r=M_ra_r$; dimensionless symbols do not choose it.\\[0.3em]
Lattice/continuum momentum comparison&Proved exactly
&$2\kappa/\pi\le p\le\kappa$ and
$0\le\kappa-p\le a^2\kappa^3/24$.\\[0.3em]
Finite-volume labels&Proved exactly
&Bounded physical momentum reduces, after subsequence extraction, to fixed
indices or fixed toron resonance offsets.\\[0.3em]
Thermodynamic labels&Proved exactly
&Nonzero finite momentum requires index or toron-offset density
$n_r/L_r$ or $j_r/L_r$, not a fixed integer.\\[0.3em]
First hard-high thermodynamic index&Proved exactly
&It grows as $\Lambda_+L_r/(2\pi)$ and its physical overshoot is at most
$2\pi/L_r$.\\[0.3em]
Threshold collar density&Proved exactly
&A fixed physical collar contains asymptotic mode density
$\Delta/(2\pi)$ per unit axial length.\\[0.3em]
Carrier spectral passport&Finite theorem
&Weyl intervals decide hard ranks away from the threshold; crossing intervals
remain buffered \stopstatus.\\[0.3em]
Projection orientation stability&Proved quantitatively
&A held-out selector gap gives the bound
$\|P-\widehat P\|\le\varepsilon/(2\widehat g-\varepsilon)$.\\[0.3em]
Carrier postselection&Excluded by countertheorem
&Choosing the projection after seeing the margin can manufacture either sign.\\[0.3em]
Thermodynamic bounded annulus&Proved conditionally on physical stocks
&One fused stiffness, fourth moment, and Ward-relative tail control every
retained mode in a fixed guarded physical annulus.\\[0.3em]
Annulus/stiffness disagreement&Typed quantitatively
&A disagreement forces fourth-moment growth of order $a_r^{-2}$ or an
unresolved relative tail.\\[0.3em]
Weak-facing branch&Open physical value
&No supplied card decides bounded roots versus an escaping selected sequence.\\[0.3em]
Absolute volume route&Open physical loading
&The supplied YM--A packet does not choose finite-volume versus thermodynamic
local-vacuum semantics.\\[0.3em]
Response-matched Wilson law&Open physical value
&The response cofactors and the actual final shell law remain unpopulated.\\[0.3em]
Guarded annulus and ultraviolet Reads&Open physical values
&Neither the bounded-physical-momentum fused kernel nor the far-window signed
head/tail has been evaluated on the selected law.\\[0.3em]
Selected high-shell margin&\stopstatus
&The corrected geometric alternatives are exact, but the Wilson-specific
values are absent.\\[0.3em]
Local-vacuum contraction and mass&Not inferred
&No volume-label theorem supplies the quadratic margin, higher jets, temporal
transfer, continuum landing, or a Yang--Mills mass gap.\\
\bottomrule
\end{longtable}
\endgroup

\begin{theorem}[Version V19 volume-route correction]
\label{thm:YMA19-terminal}
Versions~V1--V18 are retained at their stated finite or conditional scopes.
Version~V19 additionally proves:
\begin{enumerate}[label=\textup{(E19.\arabic*)},leftmargin=5.7em]
\item the absolute volume route $L_r=M_ra_r$ is part of the carrier passport;
\item fixed Fourier indices and fixed toron offsets label nonzero finite
      physical momentum only on a finite-volume route;
\item on a thermodynamic route, physical labels are index or resonance-offset
      densities and the first strict hard-high mode grazes the threshold;
\item a fixed guard or the complete buffer is required for stable carrier
      transport near that threshold;
\item a finite selector interval and gap certify the high rank and projection
      before the Wilson margin is inspected;
\item one stiffness/fourth-moment/tail packet controls every actual retained
      mode in a fixed bounded physical annulus, even when its mode count grows;
\item disagreement with a separated stiffness sign forces an
      $a_r^{-2}$-scale fourth moment or a nonvanishing relative tail;
\item the ultraviolet completion remains an independent Wilson-specific
      acquisition.
\end{enumerate}
No item decides the weak-facing branch, chooses the volume route, evaluates a
response cofactor or Wilson symbol, or proves local-vacuum contraction or a
continuum Yang--Mills mass gap.
\end{theorem}

\begin{proof}
Items~\textup{(E19.1)--(E19.3)} are
\Cref{thm:YMA19-momentum-comparison,thm:YMA19-volume-classification,thm:YMA19-threshold-grazing,thm:YMA19-toron-volume-label}.
Items~\textup{(E19.4)--(E19.5)} are
\Cref{def:YMA19-guarded-carrier,thm:YMA19-spectral-passport,prop:YMA19-projection-angle,cth:YMA19-no-postselection}.
Items~\textup{(E19.6)--(E19.7)} are
\Cref{thm:YMA19-annulus-certificate,cor:YMA19-thermodynamic-annulus,cor:YMA19-annulus-disagreement}.
Item~\textup{(E19.8)} is \Cref{fire:YMA19-annulus-not-UV}.
\end{proof}

\begin{openproblem}[V20: choose the volume route and execute its first Wilson card]
\label{op:YMA19-V20}
Preserve every frozen source, response, gauge, topological, and reference
choice.  Do not append another index or carrier theorem before the following
physical loading is supplied.
\begin{enumerate}[label=\textup{(V20.\arabic*)},leftmargin=5.4em]
\item Apply the weak-facing trajectory gate.
\item Load the absolute geometric sequence
      \[
       (a_r,M_r,L_r),\qquad L_r=M_ra_r,
      \]
      and declare either the shell-renormalized finite-volume target or the
      uniformly local thermodynamic target.
\item Load the independent selector matrix, thresholds, buffer, rooted-Hodge
      support, toron patch, and stabilizer; certify the carrier with
      \Cref{thm:YMA19-spectral-passport}.
\item On the finite-volume branch, execute the first actual Version~V18 block.
      On the thermodynamic branch, fix a guard $\Delta>0$ and a first bounded
      physical annulus $[\Lambda_++\Delta,P]$; do not transport the first
      threshold-grazing grid index as an isolated block.
\item Acquire the two response cofactors and load the native response-matched
      Wilson law on the selected carrier.
\item Acquire the fused stiffness, fourth moment, and Ward-relative tail, or
      the direct V16 likelihood card, on that same physical packet.
\item Independently populate the unbounded-momentum/far-window signed head and
      tail required by Version~V18.
\item Run the gauge held-out and keep the topological/harmonic packet separate.
      Return \pass, an attained \obstruction, or \stopstatus{} at the first
      unresolved physical row.
\end{enumerate}
If the volume route is not loaded, the programme stops before mode indexing.
If the carrier is loaded but the Wilson values are absent, it stops at the
literal signed mode card.  Boundary/Witten ancestry is opened only after a
physical mode has been frozen for attribution or transport.
\end{openproblem}

\begin{assessment}[V19 termination]
\label{ass:YMA19-termination}
The fixed-box assumption hidden in the previous cofinal labels has been made
explicit and separated from the thermodynamic local-vacuum target.  On the
latter branch, a single first index is not a stable physical object: it grows
with volume and approaches the carrier threshold.  The guarded annulus theorem
replaces the growing finite list by one uniform stiffness packet, while the
ultraviolet sector remains separate.  The selected Wilson margin is still
\stopstatus{} because the geometric route and Wilson-specific expectations
have not been supplied.  No negative physical mode has been produced.
\end{assessment}

\section*{Version V19 append-only continuation record}
\addcontentsline{toc}{section}{Version V19 append-only continuation record}
The complete Version~V18 source above is preserved byte for byte up to its
sole terminal document-closing command.  Its SHA--256 digest is
\begin{center}\scriptsize\ttfamily
40ae6f6627cc0f7da0b01f9e2579a76e0218fbb60d923a067078f8f431b43803.
\end{center}
For Version~V20, preserve this accumulated source, append new mathematical
material immediately before the final document-closing command, and increment
the filename to \texttt{\_V20.tex}.  Keep the weak-facing selector, response
law, pre-retuning action, cutoff interaction, plaquette writer, local/topological
split, colour normalization, reference normalization, and carrier thresholds
fixed.  In addition, freeze the absolute volume route $(a_r,M_r,L_r)$ and the
finite-volume versus thermodynamic target before assigning a cofinal Fourier
or toron label.  On a thermodynamic branch, use a guarded physical annulus or
the complete buffer rather than the first threshold-grazing grid point.

% =============================================================================
% END APPENDED VERSION V19 MATERIAL
% =============================================================================


% =============================================================================
% BEGIN APPENDED VERSION V20 MATERIAL
% =============================================================================

\clearpage
\hypersetup{
  pdftitle={Renewal Yang--Mills Reduced-Fibre Wilson Shell Symbol and Local-Vacuum Programme V20},
  pdfsubject={Two-regulator separation, selected local-vacuum state, continuous Bloch fibres, and finite-periodic realization of the Wilson symbol}
}
\providecommand{\Wthree}{\mathsf W_3}

\section*{Version V20: the selected local vacuum and continuous Bloch card precede a diagonal volume route}
\addcontentsline{toc}{section}{Version V20: selected local vacuum and continuous Bloch card before a diagonal volume route}

\begin{abstract}
Version~V19 correctly observed that a cofinal Fourier label has no physical
meaning until the lattice spacing and physical volume are specified.  It
nevertheless used one integer $M_r$ simultaneously as the ultraviolet shell
index and the infrared box size.  For the local-vacuum target this is still
one step too late and one regulator too few.  A diagonal sequence
$(a_r,M_r,L_r)$ can converge while failing to represent the thermodynamic
local state at any fixed shell; conversely it can miss a perfectly existing
local state.  The physical reduced-fibre symbol depends first on the selected
response-matched local vacuum law and its connected displacement kernel, not
on a global Fourier grid.

This continuation separates the shell/cutoff index $r$ from an independent
infrared side $N$.  At fixed $r$, the finite law is
$\mu_{r,N}^{\mathfrak b}$ on a box of physical side
$L_{r,N}=Na_r$, with boundary, toron, and stabilizer label
$\mathfrak b$ fixed before the limit.  A selected local vacuum is a local
weak limit as $N\to\infty$.  A pair of explicit bounded arrays proves that
an arbitrary diagonal $N=N_r$ does not determine this iterated limit.  A
second finite countermodel proves that the same local action and geometry can
produce opposite curvature signs in two selected sectors; volume data do not
choose the state.

For the actual axial Wilson card, local convergence determines every finite
head coefficient.  A uniform Ward-relative second-displacement tail then
upgrades head convergence to uniform convergence of the complete Ward
quotient.  The limiting convolution kernel has a continuous Bloch symbol on
all quasi-momenta $q\in\mathbb T$, not only on periodic grid points.  Every
$q$ is realized exactly by a finite twisted boundary condition, while an
ordinary periodic grid is only a sampling device.  Exact periodization agrees
with the infinite symbol at allowed grid points.

The quotient sampling error is made explicit.  For
$D_m(q)=(1-\cos mq)/(1-\cos q)$,
\[
 D_m(q)=m+2\sum_{j=1}^{m-1}(m-j)\cos(jq),
 \qquad
 \|D_m'\|_\infty\le\frac{m^3-m}{3}.
\]
Hence a fused kernel with finite third displacement stock has a Lipschitz
Ward quotient.  The exact Wilson reference quotient
$r_0=s_0/\omega$ has Lipschitz constant $23/144$.  If $q_N$ is the nearest
periodic grid point to a predeclared Bloch momentum $q_*$, the total
margin-quotient error is bounded by
\[
 \varepsilon_{r,N}^{\rm loc}
 +\frac{\pi}{N}
  \left(\frac{\mathsf W_{3,r}}{3}
        +\frac{23\delta\beta_r}{144}\right).
\]
A separated positive or negative local Bloch margin therefore produces the
same verdict on all sufficiently large periodic realizations.  No fixed
finite-volume route is needed to define or evaluate that first local card.
The V19 volume-route theorems remain valid for global spectra, mode counting,
and finite-torus obstruction sequences; they are no longer treated as an
upstream prerequisite for the local symbol itself.

No local-vacuum state, response cofactor, fused Wilson kernel, or Bloch-symbol
value is manufactured.  The selected shell remains at \stopstatus.  The first
physical row is now the selected response-matched local state, followed by
one finite head and one certified connected tail on a predeclared continuous
Bloch block.
\end{abstract}

% =============================================================================
\section{The ultraviolet shell and infrared volume are independent regulators}
\label{sec:YMA20-two-regulators}
% =============================================================================

Let $r$ denote the already frozen Wilson shell/cutoff index.  It carries the
coarse spacing $a_r>0$, response-matched local action $\Phi_r$, colour and
field-tangent normalization, local/topological split, and the physical
carrier thresholds.  Introduce a second integer $N\ge4$ for the infrared
periodic side and put
\begin{equation}
 \boxed{L_{r,N}:=Na_r.}
 \label{eq:YMA20-two-scale-volume}
\end{equation}
Let
\begin{equation}
 \mu_{r,N}^{\mathfrak b}
 \label{eq:YMA20-finite-law}
\end{equation}
be the complete finite response-matched reduced-fibre law with a
predeclared boundary, topological, toron, and stabilizer label $\mathfrak b$.  The
label is part of the law and is not chosen after a symbol value is read.

\begin{definition}[Selected local-vacuum packet]
\label{def:YMA20-local-vacuum}
At fixed $r$, a \emph{selected local vacuum} is a probability law
$\mu_r^{\rm vac}$ on the infinite covering lattice together with a declared
finite-volume family $\mu_{r,N}^{\mathfrak b}$ such that, for every finite
cylinder support $\Lambda$ and every bounded writer $F$ measurable on
$\Lambda$,
\begin{equation}
 \boxed{
 \mathbb E_{\mu_{r,N}^{\mathfrak b}}F
 \longrightarrow
 \mathbb E_{\mu_r^{\rm vac}}F
 \qquad(N\to\infty).}
 \label{eq:YMA20-local-convergence}
\end{equation}
When quantitative use is intended, write
\begin{equation}
 \varepsilon_{r,N}(\Lambda)
 :=\left\|
   (\mu_{r,N}^{\mathfrak b})_\Lambda
   -(\mu_r^{\rm vac})_\Lambda
  \right\|_{\rm TV,*},
 \label{eq:YMA20-local-TV}
\end{equation}
where
\begin{equation}
 \|\rho-\nu\|_{\rm TV,*}
 :=\sup_{\|f\|_\infty\le1}
   \left|\int f\,\dd\rho-\int f\,\dd\nu\right|.
 \label{eq:YMA20-TV-convention}
\end{equation}
The state, boundary label, and convergence mode are physical loading data;
they are not consequences of the geometric identity $L_{r,N}=Na_r$.
\end{definition}

\begin{countertheorem}[A diagonal regulator does not determine the local-vacuum limit]
\label{cth:YMA20-diagonal-no-local-limit}
There are bounded finite-state expectation tables $x_{r,N}$ and $y_{r,N}$
such that every fixed-$r$ infrared limit exists, while the diagonal
$N_r=r$ gives the opposite value:
\begin{equation}
 \boxed{
 x_{r,N}:=\mathbf1_{\{N=r\}},
 \qquad
 \lim_{N\to\infty}x_{r,N}=0,
 \qquad
 x_{r,r}=1,}
 \label{eq:YMA20-diagonal-x}
\end{equation}
whereas
\begin{equation}
 \boxed{
 y_{r,N}:=\mathbf1_{\{N>r\}},
 \qquad
 \lim_{N\to\infty}y_{r,N}=1,
 \qquad
 y_{r,r}=0.}
 \label{eq:YMA20-diagonal-y}
\end{equation}
Both arrays are realized as expectations of one binary writer under positive
probability laws.  Consequently a single diagonal triple
$(a_r,N_r,L_{r,N_r})$ neither proves nor identifies the fixed-shell local
vacuum in \eqref{eq:YMA20-local-convergence}.  A diagonal passage is valid
only after an infrared convergence modulus, uniform estimate, or explicit
selected-state theorem has been supplied.
\end{countertheorem}

\begin{proof}
The displayed fixed-$r$ and diagonal limits are immediate.  To realize either
array by positive laws, let $\Omega=\{0,1\}$, let $X(\omega)=\omega$, and set
\[
 \mu_{r,N}\{1\}=x_{r,N}
 \quad\hbox{or}\quad
 \mu_{r,N}\{1\}=y_{r,N}.
\]
Then $\mathbb E_{\mu_{r,N}}X$ is the desired table.  Thus the failure is
already present for finite positive cylinders and is not an artefact of an
unbounded field or a sign-indefinite measure.
\end{proof}

\begin{countertheorem}[Geometry and the local action formula do not select a vacuum sector]
\label{cth:YMA20-geometry-no-state}
Let the local configuration carry one sector bit $\sigma\in\{-1,+1\}$ and
let the same action deformation be
\begin{equation}
 F_t(\sigma)=t^2\sigma.
 \label{eq:YMA20-sector-action}
\end{equation}
Conditioning the same action on the two predeclared sectors gives
$\mu^+=\delta_{+1}$ and $\mu^-=\delta_{-1}$.  Their absolute partition
curvatures are
\begin{equation}
 \boxed{
 \left.\frac{\dd^2}{\dd t^2}[-\log e^{-F_t(+1)}]\right|_{t=0}=2,
 \qquad
 \left.\frac{\dd^2}{\dd t^2}[-\log e^{-F_t(-1)}]\right|_{t=0}=-2.}
 \label{eq:YMA20-sector-opposite-curvature}
\end{equation}
Thus the same geometry, deformation formula, and volume route can have
opposite loaded curvature signs when the selected state/sector is changed.
This is a logical countermodel to state-free inference; it is not a claim
that the Wilson vacuum has these two sectors or these values.
\end{countertheorem}

\begin{proof}
On the conditioned one-point state, the partition function is
$Z_\pm(t)=e^{\mp t^2}$ and the effective action is
$-\log Z_\pm(t)=\pm t^2$.  Differentiation gives
\eqref{eq:YMA20-sector-opposite-curvature}.
\end{proof}

\begin{firewall}[What Version V19 still proves]
\label{fire:YMA20-V19-scope}
The kinematic classifications in Version~V19 remain exact for a declared
finite periodic spectral target: the meaning of a periodic Fourier index,
the threshold-grazing first hard-high mode, the transition-collar density,
and the finite carrier passport all depend on $(a_r,N,L_{r,N})$.  The present
correction concerns the \emph{local-vacuum symbol}.  Its first input is the
selected local state and connected kernel.  The global volume route enters
later when that continuous symbol is represented by periodic eigenmodes or
used in a full finite-volume spectral statement.
\end{firewall}

% =============================================================================
\section{Local convergence determines the finite Wilson head}
\label{sec:YMA20-local-head}
% =============================================================================

The direct action curvature at a fixed displacement is a bounded cylinder
writer on the compact Wilson fibre.  A finite-range score covariance is a
joint cylinder writer on the union of the two insertion collars.  Generated
quasilocal terms are handled by a finite head plus their already declared
tail; no infinite support is inserted into the next lemma.

\begin{lemma}[Local total variation controls a connected coefficient]
\label{lem:YMA20-TV-covariance}
Let $A$ and $B$ be real bounded writers measurable on a finite support
$\Lambda$, and let $\rho,\nu$ be two laws on that support.  With the convention
\eqref{eq:YMA20-TV-convention},
\begin{equation}
 \boxed{
 \left|\operatorname{Cov}_\rho(A,B)
       -\operatorname{Cov}_\nu(A,B)\right|
 \le
 3\|A\|_\infty\|B\|_\infty
 \|\rho-\nu\|_{\rm TV,*}.}
 \label{eq:YMA20-TV-covariance}
\end{equation}
Moreover, for every bounded $K$ on the same support,
\begin{equation}
 \boxed{
 |\mathbb E_\rho K-\mathbb E_\nu K|
 \le\|K\|_\infty\|\rho-\nu\|_{\rm TV,*}.}
 \label{eq:YMA20-TV-mean}
\end{equation}
\end{lemma}

\begin{proof}
The mean estimate is the definition of the dual total-variation norm after
scaling by $\|K\|_\infty$.  For the covariance, write
\[
 \operatorname{Cov}_\rho(A,B)-\operatorname{Cov}_\nu(A,B)
 =\bigl(\mathbb E_\rho AB-\mathbb E_\nu AB\bigr)
  -\bigl(\mathbb E_\rho A\mathbb E_\rho B
         -\mathbb E_\nu A\mathbb E_\nu B\bigr).
\]
The first difference is at most
$\|A\|_\infty\|B\|_\infty\|\rho-\nu\|_{\rm TV,*}$.
For the product of means, add and subtract
$\mathbb E_\nu A\mathbb E_\rho B$.  Each of the resulting two terms has the
same upper bound.  Summing gives \eqref{eq:YMA20-TV-covariance}.
\end{proof}

\begin{corollary}[Convergence of every finite fused displacement head]
\label{cor:YMA20-head-convergence}
Fix $r$ and a real axial transverse colour polarization.  Suppose that for
each displacement $m$ the deterministic insertion $K_{r,m}$ and the two
score insertions $J_{r,0},J_{r,m}$ are bounded cylinder writers on a finite
collar $\Lambda_{r,m}$.  Define
\begin{equation}
 \alpha_{r,N}(m)
 :=\mathbb E_{\mu_{r,N}^{\mathfrak b}}K_{r,m}
   -\operatorname{Cov}_{\mu_{r,N}^{\mathfrak b}}
      (J_{r,0},J_{r,m})
   -t_{r,N}(m),
 \label{eq:YMA20-fused-coefficient}
\end{equation}
where the local topological term $t_{r,N}(m)$ is retained on its actual
support.  If the selected-state convergence
\eqref{eq:YMA20-local-convergence} holds on every $\Lambda_{r,m}$ and the
corresponding topological writer also converges, then for every finite $R$,
\begin{equation}
 \boxed{
 \sum_{m=1}^{R}m^2
 |\alpha_{r,N}(m)-\alpha_r(m)|\longrightarrow0.}
 \label{eq:YMA20-weighted-head-convergence}
\end{equation}
Here $\alpha_r(m)$ is the same expectation--covariance--topology expression
in the selected local vacuum $\mu_r^{\rm vac}$.
\end{corollary}

\begin{proof}
Apply \Cref{lem:YMA20-TV-covariance} to the score pair and
\eqref{eq:YMA20-TV-mean} to each deterministic or topological writer.  Each
fixed coefficient converges.  The displayed head contains finitely many
coefficients, so the weighted sum converges.
\end{proof}

\begin{remark}[Why local convergence alone is not enough]
\label{rem:YMA20-head-not-tail}
Local convergence controls every fixed displacement, but the Ward quotient
contains the weighted sum of all displacements.  Mass may escape to distances
which grow with $N$ while every fixed head converges.  The next theorem keeps
that possibility as the physical connected-tail row rather than hiding it in
a diagonal limit.
\end{remark}

% =============================================================================
\section{Head convergence plus a Ward-relative tail gives the local symbol}
\label{sec:YMA20-head-tail}
% =============================================================================

For a real reflection-even axial kernel write, in the convention of
Version~V1,
\begin{equation}
 h(q)=2\sum_{m\ge1}\alpha(m)(\cos(mq)-1),
 \qquad
 \omega(q)=2(1-\cos q).
 \label{eq:YMA20-axial-symbol}
\end{equation}
Its continuously extended Ward quotient is
\begin{equation}
 g(q)=\frac{h(q)}{\omega(q)}
 =-\sum_{m\ge1}\alpha(m)D_m(q),
 \qquad
 D_m(q):=\frac{1-\cos(mq)}{1-\cos q},
 \label{eq:YMA20-Ward-quotient}
\end{equation}
with $D_m(0)=m^2$.

For a periodic side $N$, lift each shortest axial displacement to
$1\le m\le\lfloor N/2\rfloor$.  Let $z_{r,N}$ be the unrepaired zero symbol
and let $\alpha_{r,N}^{\circ}$ be the contact-repaired coefficient family of
Version~V4.  Put
\begin{equation}
 \begin{aligned}
 E_{r,N}^{\rm head}(R)
 &:=\sum_{m=1}^{R}m^2
       |\alpha_{r,N}^{\circ}(m)-\alpha_r(m)|,\\
 T_{r,N}^{\rm per}(R)
 &:=\sum_{R<m\le N/2}m^2
       |\alpha_{r,N}^{\circ}(m)|,\\
 T_r^{\infty}(R)
 &:=\sum_{m>R}m^2|\alpha_r(m)|,\\
 H_{r,N}^{\rm harm}
 &:=\frac{N^2}{16}|z_{r,N}|.
 \end{aligned}
 \label{eq:YMA20-head-tail-stocks}
\end{equation}

\begin{theorem}[Local-state head--tail completion of the Ward quotient]
\label{thm:YMA20-head-tail-completion}
Assume $N>2R$ and let $g_{r,N}$ be the complete finite periodic quotient on
its nonzero axial grid.  Then
\begin{equation}
 \boxed{
 \max_{q\in(2\pi/N)\Z\setminus\{0\}}
 |g_{r,N}(q)-g_r(q)|
 \le
 E_{r,N}^{\rm head}(R)
 +T_{r,N}^{\rm per}(R)
 +T_r^{\infty}(R)
 +H_{r,N}^{\rm harm}.}
 \label{eq:YMA20-head-tail-uniform-bound}
\end{equation}
If the complete physical finite kernel has exact Ward zero, then
$H_{r,N}^{\rm harm}=0$.  Consequently, if
\begin{equation}
 \lim_{R\to\infty}\sup_{N>2R}T_{r,N}^{\rm per}(R)=0,
 \qquad
 \lim_{R\to\infty}T_r^{\infty}(R)=0,
 \label{eq:YMA20-uniform-tail-hyp}
\end{equation}
\begin{equation}
 \lim_{N\to\infty}E_{r,N}^{\rm head}(R)=0
 \quad\hbox{for every fixed }R,
 \qquad
 H_{r,N}^{\rm harm}\longrightarrow0,
 \label{eq:YMA20-head-harm-hyp}
\end{equation}
then $g_{r,N}$ converges uniformly on the nonzero periodic grids to the
continuous local-vacuum quotient $g_r$.
\end{theorem}

\begin{proof}
For the contact-repaired finite kernel and the limiting zero-row-sum kernel,
\Cref{thm:YMA1-W2-control} bounds the quotient of their coefficient
difference by the second displacement moment.  Split that difference into
the head $m\le R$, the finite periodic tail $R<m\le N/2$, and the limiting
tail $m>R$.  This gives the first three terms in
\eqref{eq:YMA20-head-tail-uniform-bound}.  The unrepaired harmonic scalar is
handled by \Cref{lem:YMA4-harmonic-Ward-router}, which contributes exactly
$N^2|z_{r,N}|/16$ on the nonzero grid.  The limiting statement follows by
first choosing $R$ so that both tail terms are small and then choosing $N$
so that the head and harmonic terms are small.
\end{proof}

\begin{corollary}[Matrix-valued rooted-Hodge blocks]
\label{cor:YMA20-matrix-head-tail}
The same theorem holds for a finite-dimensional Hermitian momentum--colour
block when absolute values are replaced by operator norms and the coefficient
families satisfy
$A(-m)=A(m)^*$.  In particular, every ordered eigenvalue of the finite block
converges uniformly to the corresponding ordered eigenvalue of the local
Bloch block.
\end{corollary}

\begin{proof}
The Dirichlet factors $D_m(q)$ are nonnegative scalars bounded by $m^2$.
The triangle inequality in operator norm gives the same estimate, and Weyl's
Hermitian eigenvalue inequality transfers the norm convergence to every
ordered eigenvalue.
\end{proof}

\begin{corollary}[Failure of the uniform tail is a typed positive-range row]
\label{cor:YMA20-tail-survivor}
Suppose every fixed finite head converges as in
\eqref{eq:YMA20-weighted-head-convergence}, but
\eqref{eq:YMA20-uniform-tail-hyp} fails.  Then there are
$\epsilon>0$, radii $R_j\to\infty$, and periodic sizes $N_j>2R_j$ such that
\begin{equation}
 \boxed{
 \sum_{R_j<m\le N_j/2}m^2
 |\alpha_{r,N_j}^{\circ}(m)|\ge\epsilon.}
 \label{eq:YMA20-positive-range-witness}
\end{equation}
This is an attained displacement-tail or positive-range survivor on the
selected state family.  It is not repaired by choosing a different diagonal
Fourier index.
\end{corollary}

\begin{proof}
The displayed sequence is exactly the negation of the first limit in
\eqref{eq:YMA20-uniform-tail-hyp}.  Since the fixed heads converge, the stock
cannot be attributed to any fixed displacement range.
\end{proof}

% =============================================================================
\section{The local symbol is a continuous Bloch fibre, not a periodic-grid object}
\label{sec:YMA20-Bloch}
% =============================================================================

The preceding scalar quotient is the one-polarization reduction of a
translation-covariant matrix kernel.  It is useful to state the exact
finite-cell representation before returning to the axial scalar.

\begin{theorem}[Continuous Bloch symbol and exact periodization]
\label{thm:YMA20-Bloch-periodization}
Let $E$ be a finite-dimensional Hermitian fibre and let
$K:\Z^4\to\operatorname{End}(E)$ satisfy
\begin{equation}
 K(-x)=K(x)^*,
 \qquad
 \sum_{x\in\Z^4}\|K(x)\|_{\rm op}<\infty.
 \label{eq:YMA20-kernel-l1}
\end{equation}
The convolution operator
\begin{equation}
 (\mathcal Kf)(y)=\sum_{x\in\Z^4}K(x)f(y-x)
 \label{eq:YMA20-convolution}
\end{equation}
has the continuous Hermitian Bloch symbol
\begin{equation}
 \boxed{
 \widehat K(q)=\sum_{x\in\Z^4}K(x)e^{-\mathrm i q\cdot x},
 \qquad q\in[-\pi,\pi]^4.}
 \label{eq:YMA20-continuous-symbol}
\end{equation}
For every $v\in E$, the generalized plane wave
$f_{q,v}(y)=e^{\mathrm iq\cdot y}v$ obeys
\begin{equation}
 \mathcal Kf_{q,v}(y)=e^{\mathrm iq\cdot y}\widehat K(q)v.
 \label{eq:YMA20-plane-wave}
\end{equation}

For a periodic side $N$, define the exact periodization
\begin{equation}
 K_N^{\rm per}(\bar x)
 :=\sum_{n\in\Z^4}K(x+Nn).
 \label{eq:YMA20-periodized-kernel}
\end{equation}
At every ordinary periodic grid momentum
$q\in(2\pi/N)\Z^4$,
\begin{equation}
 \boxed{\widehat K_N^{\rm per}(q)=\widehat K(q).}
 \label{eq:YMA20-periodization-symbol}
\end{equation}
For arbitrary $q$, let $\mathcal H_{N,q}$ be the finite space of
$q$-quasiperiodic functions satisfying
$f(y+Ne_j)=e^{\mathrm iNq_j}f(y)$.  The restriction of
\eqref{eq:YMA20-convolution} to this finite Bloch sector is represented on
one fundamental cell by
\begin{equation}
 K_{N,q}^{\rm tw}(x,y)
 =\sum_{n\in\Z^4}
   K(x-y+Nn)e^{-\mathrm iNn\cdot q},
 \label{eq:YMA20-twisted-kernel}
\end{equation}
and it contains the exact finite-dimensional block
$\widehat K(q)$.  Thus a continuous quasi-momentum is a literal finite twisted
card and need not be approximated by a periodic grid.
\end{theorem}

\begin{proof}
Absolute summability makes \eqref{eq:YMA20-continuous-symbol} uniformly
convergent and therefore continuous.  Hermiticity follows by replacing
$x$ with $-x$.  Substitution of $f_{q,v}$ in
\eqref{eq:YMA20-convolution} gives \eqref{eq:YMA20-plane-wave}.

For an allowed periodic momentum, $e^{-\mathrm iq\cdot Nn}=1$.  Hence
\[
 \sum_{\bar x}K_N^{\rm per}(\bar x)e^{-\mathrm iq\cdot x}
 =\sum_{\bar x}\sum_nK(x+Nn)e^{-\mathrm iq\cdot(x+Nn)}
 =\sum_{z\in\Z^4}K(z)e^{-\mathrm iq\cdot z},
\]
which proves \eqref{eq:YMA20-periodization-symbol}.

For the twisted sector, every $z\in\Z^4$ has a unique representation
$z=x-y+Nn$ once $x$ and the residue class of $y$ are fixed.  Acting with
\eqref{eq:YMA20-twisted-kernel} on $e^{\mathrm iq\cdot y}v$ gives
\[
 \sum_{y,n}K(x-y+Nn)e^{-\mathrm iNn\cdot q}
             e^{\mathrm iq\cdot y}v
 =e^{\mathrm iq\cdot x}\sum_zK(z)e^{-\mathrm iq\cdot z}v.
\]
This is the asserted finite Bloch block.
\end{proof}

\begin{assessment}[The local-vacuum card is volume independent after state selection]
\label{ass:YMA20-volume-downstream}
Once $\mu_r^{\rm vac}$ and its summable fused kernel are fixed, the matrix
\eqref{eq:YMA20-continuous-symbol} is the fixed-shell local Wilson symbol.
Changing the total periodic volume changes only the ordinary grid on which it
is sampled and the finite-volume defect from the selected local state.  It
does not change the definition of the local Bloch block.  The absolute volume
route of Version~V19 remains indispensable for full finite-torus spectra,
mode densities, and global obstruction sequences, but it is downstream of
this local object.
\end{assessment}

% =============================================================================
\section{A sharp sampling bound for the Ward quotient}
\label{sec:YMA20-sampling}
% =============================================================================

\begin{lemma}[Exact Dirichlet quotient polynomial and derivative bound]
\label{lem:YMA20-Dirichlet-Lipschitz}
For every integer $m\ge1$,
\begin{equation}
 \boxed{
 D_m(q)
 =m+2\sum_{j=1}^{m-1}(m-j)\cos(jq).}
 \label{eq:YMA20-Dirichlet-polynomial}
\end{equation}
Consequently
\begin{equation}
 \boxed{
 \|D_m'\|_{L^\infty(\T)}
 \le2\sum_{j=1}^{m-1}j(m-j)
 =\frac{m^3-m}{3}.}
 \label{eq:YMA20-Dirichlet-derivative}
\end{equation}
\end{lemma}

\begin{proof}
The geometric sum gives
\[
 D_m(q)
 =\left|\sum_{k=0}^{m-1}e^{\mathrm ikq}\right|^2.
\]
Grouping the pairs with difference $j$ yields
\eqref{eq:YMA20-Dirichlet-polynomial}, including its continuous value at
$q=0$.  Differentiate the finite trigonometric polynomial and use
$|\sin(jq)|\le1$.  Finally,
\[
 2\sum_{j=1}^{m-1}j(m-j)
 =2\left(m\frac{m(m-1)}2
       -\frac{m(m-1)(2m-1)}6\right)
 =\frac{m^3-m}{3}.
\]
\end{proof}

For a scalar reflected kernel put
\begin{equation}
 \|\alpha\|_{\Wthree}
 :=\sum_{m\ge1}m^3|\alpha(m)|.
 \label{eq:YMA20-W3}
\end{equation}

\begin{theorem}[Lipschitz Ward quotient]
\label{thm:YMA20-W3-Lipschitz}
If $\|\alpha\|_{\Wthree}<\infty$, then the quotient in
\eqref{eq:YMA20-Ward-quotient} is globally Lipschitz and
\begin{equation}
 \boxed{
 |g(q)-g(q')|
 \le\frac13\|\alpha\|_{\Wthree}|q-q'|_{\T}.}
 \label{eq:YMA20-g-Lipschitz}
\end{equation}
The same estimate holds for Hermitian matrix coefficients with the operator
norm inside $\Wthree$.
\end{theorem}

\begin{proof}
Use
$g=-\sum_m\alpha(m)D_m$ and
\Cref{lem:YMA20-Dirichlet-Lipschitz}.  Since
$m^3-m\le m^3$, the derivative series converges uniformly and has norm at
most $\|\alpha\|_{\Wthree}/3$.  The matrix statement follows from the
operator-norm triangle inequality.
\end{proof}

\begin{lemma}[Exact Lipschitz constant for the Wilson reference quotient]
\label{lem:YMA20-reference-Lipschitz}
Let
\begin{equation}
 r_0(q)=\frac{s_0(q)}{\omega(q)}
 =\frac{(t+8)(t+24)}{36(t^2+18t+16)},
 \qquad t=\sin^2(q/2).
 \label{eq:YMA20-r0}
\end{equation}
Then
\begin{equation}
 \boxed{
 |r_0(q)-r_0(q')|
 \le\frac{23}{144}|q-q'|_{\T}.}
 \label{eq:YMA20-r0-Lipschitz}
\end{equation}
\end{lemma}

\begin{proof}
Differentiation in $t$ gives
\begin{equation}
 \frac{\dd r_0}{\dd t}
 =-\frac{7t^2+176t+1472}
         {18(t^2+18t+16)^2}.
 \label{eq:YMA20-r0-t-derivative}
\end{equation}
The derivative of its absolute value is
\begin{equation}
 -\frac{7t^3+264t^2+4416t+25088}
        {9(t^2+18t+16)^3}<0
 \qquad(0\le t\le1).
 \label{eq:YMA20-r0-monotone-derivative}
\end{equation}
Thus the maximum occurs at $t=0$ and equals $23/72$.  Since
$|\dd t/\dd q|=|\sin q|/2\le1/2$, the chain rule gives
\eqref{eq:YMA20-r0-Lipschitz}.
\end{proof}

\begin{lemma}[Nearest periodic sample]
\label{lem:YMA20-nearest-grid}
For every $q_*\in\T$ and every $N\ge1$, there is
$q_N\in(2\pi/N)\Z/2\pi\Z$ such that
\begin{equation}
 \boxed{|q_N-q_*|_{\T}\le\frac{\pi}{N}.}
 \label{eq:YMA20-nearest-grid}
\end{equation}
\end{lemma}

\begin{proof}
Round $Nq_*/(2\pi)$ to the nearest integer.  The rounding error is at most
$1/2$, which gives the displayed circular-distance bound.
\end{proof}

% =============================================================================
\section{From one continuous local Bloch card to an actual periodic witness}
\label{sec:YMA20-realization}
% =============================================================================

Fix $r$ and let $g_r^{\rm loc}$ be the selected local-vacuum Ward quotient of
the complete fused axial Hessian.  Define the target margin quotient
\begin{equation}
 \mathfrak m_r(q;\delta)
 :=g_r^{\rm loc}(q)-\delta\beta_rr_0(q).
 \label{eq:YMA20-local-margin-quotient}
\end{equation}
Let $g_{r,N}^{\rm loc}$ be the complete finite-periodic quotient on the same
response-matched law and carrier convention.  Suppose the local-state
head/tail certificate gives
\begin{equation}
 \max_{q\in(2\pi/N)\Z\setminus\{0\}}
 |g_{r,N}^{\rm loc}(q)-g_r^{\rm loc}(q)|
 \le\varepsilon_{r,N}^{\rm loc}.
 \label{eq:YMA20-localization-defect}
\end{equation}
Put
\begin{equation}
 W_{3,r}:=\|\alpha_r^{\rm loc}\|_{\Wthree}.
 \label{eq:YMA20-W3-r}
\end{equation}

\begin{theorem}[Bloch-to-periodic sign transfer]
\label{thm:YMA20-Bloch-periodic-transfer}
Let $q_*\ne0$ be a predeclared continuous Bloch momentum and let $q_N$ be a
nearest periodic sample from \Cref{lem:YMA20-nearest-grid}.  Then
\begin{equation}
 \boxed{
 \begin{aligned}
 &\left|
  \bigl[g_{r,N}^{\rm loc}(q_N)-\delta\beta_rr_0(q_N)\bigr]
  -\mathfrak m_r(q_*;\delta)
 \right|\\
 &\qquad\le
 \varepsilon_{r,N}^{\rm loc}
 +\frac{\pi}{N}
  \left(\frac{W_{3,r}}3
        +\frac{23\delta\beta_r}{144}\right)
 =:\mathcal E_{r,N}(q_*;\delta).
 \end{aligned}}
 \label{eq:YMA20-Bloch-periodic-error}
\end{equation}
Consequently:
\begin{align}
 \mathfrak m_r(q_*;\delta)>
 \mathcal E_{r,N}(q_*;\delta)
 &\Longrightarrow
 g_{r,N}^{\rm loc}(q_N)>
 \delta\beta_rr_0(q_N),
 \label{eq:YMA20-periodic-pass}\\
 \mathfrak m_r(q_*;\delta)<
 -\mathcal E_{r,N}(q_*;\delta)
 &\Longrightarrow
 g_{r,N}^{\rm loc}(q_N)<
 \delta\beta_rr_0(q_N).
 \label{eq:YMA20-periodic-obstruction}
\end{align}
On the second branch, the ordinary periodic mode $q_N$ is an attained
finite-volume pure-mode witness.  If twisted boundary conditions are admitted,
the exact $q_*$ card of \Cref{thm:YMA20-Bloch-periodization} removes the grid
term entirely.
\end{theorem}

\begin{proof}
Insert and subtract $g_r^{\rm loc}(q_N)$ and
$\delta\beta_rr_0(q_*)$.  The finite/local difference is bounded by
\eqref{eq:YMA20-localization-defect}.  The two continuous-symbol differences
are bounded by \Cref{thm:YMA20-W3-Lipschitz} and
\Cref{lem:YMA20-reference-Lipschitz}; use
\eqref{eq:YMA20-nearest-grid}.  This proves
\eqref{eq:YMA20-Bloch-periodic-error}.  The sign conclusions are immediate.
\end{proof}

\begin{corollary}[Physical-momentum form]
\label{cor:YMA20-physical-momentum}
Fix a physical shell scale $a_r$ and a physical momentum $p_*$ with
$0<a_rp_*<\pi$.  Put
\begin{equation}
 q_*=a_rp_*.
 \label{eq:YMA20-q-from-p}
\end{equation}
Then \Cref{thm:YMA20-Bloch-periodic-transfer} applies without first choosing
a total physical volume.  On a central branch this is the continuous Bloch
card at physical momentum $p_*$.  On a toron root channel, replace $q_*$ by
the loaded covariant quasi-momentum relative to the selected resonance.  The
toron stabilizer and rooted-Hodge support remain part of the carrier passport.
\end{corollary}

\begin{proof}
The relation $q_*=a_rp_*$ uses only the fixed shell spacing.  The continuous
Bloch symbol is defined by \Cref{thm:YMA20-Bloch-periodization}.  A periodic
side is needed only if one elects to replace the twisted kernel realization by an ordinary
periodic eigenmode.
\end{proof}

\begin{corollary}[Retained matrix block]
\label{cor:YMA20-matrix-sign-transfer}
Let $\mathfrak M_r(q)$ be a Hermitian rooted-Hodge momentum--colour margin
block and assume an operator-norm version of
\eqref{eq:YMA20-Bloch-periodic-error} with radius $\mathcal E_{r,N}$.  Then
\begin{align}
 \lambda_{\min}(\mathfrak M_r(q_*))>\mathcal E_{r,N}
 &\Longrightarrow\pass,\\
 \lambda_{\min}(\mathfrak M_r(q_*))<-\mathcal E_{r,N}
 &\Longrightarrow\obstruction,
\end{align}
for the corresponding periodic block.  On the obstruction branch, a minimum
eigenvector of the local block, transported through the spectral passport,
provides an explicit periodic witness after including the same perturbation
radius.
\end{corollary}

\begin{proof}
Apply Weyl's Hermitian perturbation inequality to the block difference.
\end{proof}

% =============================================================================
\section{The correct two-stage cofinal order}
\label{sec:YMA20-two-stage}
% =============================================================================

The local-vacuum programme has two regulator passages:
\begin{equation}
 \boxed{
 \text{infrared local-state limit }N\to\infty\text{ at fixed }r
 \quad\Longrightarrow\quad
 \text{shell/cutoff transport }r\to\infty.}
 \label{eq:YMA20-two-stage-order}
\end{equation}
An arbitrary diagonal may be used only after its infrared error is controlled.

\begin{theorem}[Quantitative diagonal extraction for a fixed physical Bloch bank]
\label{thm:YMA20-diagonal-extraction}
Let $\mathcal P\subset(0,\infty)$ be a finite predeclared bank of physical
momenta and put $q_r(p)=a_rp$.  Assume $q_r(p)\in(0,\pi)$ for all sufficiently
large $r$.  Suppose:
\begin{enumerate}[label=\textup{(D\arabic*)},leftmargin=3.8em]
\item for every $r$ there is a selected local vacuum and a complete fused
      quotient $g_r^{\rm loc}$;
\item the local head/tail completion supplies numbers
      $\varepsilon_{r,N}^{\rm loc}\to0$ as $N\to\infty$;
\item $W_{3,r}<\infty$;
\item the local margin is separated on the bank by
      \begin{equation}
       \eta_r
       :=\min_{p\in\mathcal P}
       |\mathfrak m_r(q_r(p);\delta)|>0.
       \label{eq:YMA20-eta-r}
      \end{equation}
\end{enumerate}
Then one can choose $N_r\to\infty$ such that, simultaneously for every
$p\in\mathcal P$,
\begin{equation}
 \boxed{
 \mathcal E_{r,N_r}(q_r(p);\delta)<\frac12\eta_r.}
 \label{eq:YMA20-diagonal-choice}
\end{equation}
Every periodic realization on this chosen diagonal has the same sign as its
local Bloch card.

If instead there are uniform constants
\begin{equation}
 \inf_r\frac{\eta_r}{\beta_r}\ge\eta>0,
 \qquad
 \sup_r\frac{W_{3,r}}{\beta_r}\le C_3,
 \qquad
 \frac{\varepsilon_{r,N_r}^{\rm loc}}{\beta_r}\longrightarrow0,
 \qquad
 N_r\longrightarrow\infty,
 \label{eq:YMA20-uniform-diagonal-data}
\end{equation}
then the normalized periodic and local margin quotients differ by $o(1)$
uniformly on $\mathcal P$, and every separated sign is eventually preserved.
\end{theorem}

\begin{proof}
For each fixed $r$, first choose $N$ so large that
$\varepsilon_{r,N}^{\rm loc}<\eta_r/4$.  Increase it further until
\[
 \frac{\pi}{N}
 \left(\frac{W_{3,r}}3+\frac{23\delta\beta_r}{144}\right)
 <\frac{\eta_r}{4}.
\]
The maximum of these finitely many requirements over $p\in\mathcal P$ is the
same because the error bound is independent of $p$.  Choosing the sizes
strictly increasing gives $N_r\to\infty$ and
\eqref{eq:YMA20-diagonal-choice}.  The sign statement is
\Cref{thm:YMA20-Bloch-periodic-transfer}.

Under \eqref{eq:YMA20-uniform-diagonal-data}, divide
\eqref{eq:YMA20-Bloch-periodic-error} by $\beta_r$.  The first term tends to
zero, and the grid term is bounded by
\[
 \frac{\pi}{N_r}
 \left(\frac{C_3}{3}+\frac{23\delta}{144}\right),
\]
which also tends to zero.
\end{proof}

\begin{countertheorem}[A local symbol does not select a global volume route]
\label{cth:YMA20-local-no-volume-route}
Fix any finite-range Hermitian convolution kernel $K$ on $\Z^4$.  Its
continuous Bloch matrix $\widehat K(q)$ is the same local object when it is
represented on:
\begin{enumerate}[label=\textup{(\roman*)},leftmargin=2.8em]
\item a fixed physical torus through twisted boundary conditions;
\item a sequence of growing periodic boxes;
\item any other sequence on which the exact periodization is used.
\end{enumerate}
These representations can have different periodic mode sets and different
physical volumes while agreeing on every common Bloch block.  Therefore the
local symbol cannot determine the finite-volume versus thermodynamic target;
that is an independent global loading choice.
\end{countertheorem}

\begin{proof}
Finite range implies \eqref{eq:YMA20-kernel-l1}.  Apply
\Cref{thm:YMA20-Bloch-periodization}.  The formula
\eqref{eq:YMA20-continuous-symbol} contains no total-volume parameter, whereas
the allowed ordinary periodic grids do.  Hence the same local matrix has all
three representations.
\end{proof}

\begin{proposition}[Consistent field-tangent rescaling is a congruence, not a new sign]
\label{prop:YMA20-field-calibration}
Let the frozen coarse-link path be reparametrized by
$t\mapsto ct$ with $c\ne0$, while the physical Hessian and its reference are
both differentiated in that same coordinate.  Then
\begin{equation}
 \boxed{
 H^{(c)}(q)=c^2H(q),
 \qquad
 S^{(c)}(q)=c^2S(q),
 \qquad
 \mathfrak M^{(c)}(q;\delta)=c^2\mathfrak M(q;\delta).}
 \label{eq:YMA20-field-congruence}
\end{equation}
Every displacement coefficient and every $\Wtwo$, $\Wthree$, or $\Wfour$
stock scales by $c^2$.  The sign, kernel, and normalized margin are unchanged.
Absolute partition-ratio step sizes and derivative bounds do depend on the
chosen tangent units and must use one common calibration.
\end{proposition}

\begin{proof}
The first field derivative scales by $c$, the second field derivative by
$c^2$, and the covariance of two first derivatives by $c^2$.  Hence the
assembled Hessian and reference scale by $c^2$.  The remaining assertions are
linear consequences for their displacement kernels and positive congruence by
the scalar $c^2$.
\end{proof}

% =============================================================================
\section{Corrected termination order and the next literal Wilson acquisition}
\label{sec:YMA20-frontier}
% =============================================================================

\begin{theorem}[Local-vacuum-first YM--A alternative]
\label{thm:YMA20-local-vacuum-alternative}
For the local-vacuum target, proceed in the following order.
\begin{enumerate}[label=\textup{(L20.\arabic*)},leftmargin=5.2em]
\item \emph{State loading.}  Freeze the response-matched action, boundary or
      toron sector, and a selected local-vacuum law
      $\mu_r^{\rm vac}$, or provide a uniqueness theorem which makes the
      choice independent of boundary conditions.
\item \emph{Local assembly.}  Acquire the complete deterministic insertion,
      connected score covariance, generated support births, and declared
      local/topological split in that state.
\item \emph{Head--tail completion.}  Populate a finite displacement head and
      the Ward-relative tail in
      \Cref{thm:YMA20-head-tail-completion}.  Failure of the tail gives the
      positive-range witness \eqref{eq:YMA20-positive-range-witness}.
\item \emph{Continuous physical block.}  Freeze one carrier-supported
      physical momentum--colour block and evaluate its continuous Bloch
      matrix, using a twisted kernel card if the momentum is not on a periodic grid.
\item \emph{Finite realization.}  Only when a periodic witness or global
      finite-torus spectrum is intended, choose the volume route and apply
      \Cref{thm:YMA20-Bloch-periodic-transfer} together with the independent
      carrier passport.
\end{enumerate}
The terminating outcomes are:
\begin{enumerate}[label=\textup{(T20.\arabic*)},leftmargin=5.2em]
\item \pass: a positive lower interval for the assembled local Bloch margin,
      and, when requested, a periodic realization whose transfer error is
      smaller than that interval;
\item \obstruction: a negative upper interval on an explicit retained Bloch
      polarization, with an ordinary periodic witness supplied by
      \eqref{eq:YMA20-periodic-obstruction} when finite-volume realization is
      part of the target;
\item \stopstatus: the selected state, response-matched law, finite head,
      harmonic row, connected tail, carrier support, or strict interval is
      unresolved.
\end{enumerate}
A failure to select a vacuum state is an upstream local-vacuum loading defect,
not a negative Yang--Mills mode.  A failure of the uniform tail is a typed
positive-range row but becomes a symbol obstruction only when an assembled
retained Bloch block is negative.
\end{theorem}

\begin{proof}
State loading is necessary by
\Cref{cth:YMA20-diagonal-no-local-limit,cth:YMA20-geometry-no-state}.  Local
assembly and head convergence are
\Cref{lem:YMA20-TV-covariance,cor:YMA20-head-convergence}.  The tail and
harmonic alternatives are
\Cref{thm:YMA20-head-tail-completion,cor:YMA20-tail-survivor}.  The continuous
block and finite realization are
\Cref{thm:YMA20-Bloch-periodization,thm:YMA20-Bloch-periodic-transfer}.  The
three sign outcomes follow from the direct Hermitian interval test and Weyl's
inequality.
\end{proof}

\begin{assessment}[Exact qualification of Versions V18--V19]
\label{ass:YMA20-qualification}
Versions~V18--V19 remain exact for their stated finite periodic carrier and
volume-route questions.  Their global conclusions are not retracted.  The
statement that the absolute volume route must be chosen before \emph{any}
local symbol card is, however, too strong for the local-vacuum branch.  After
a selected local state and summable connected kernel have been loaded, an
arbitrary continuous Bloch block is already a finite twisted card and the
global volume is only a realization coordinate.  Conversely, a volume route
without local-state control does not define the local symbol.  The two rows
are independent and are now placed in the correct order.
\end{assessment}

\begin{firewall}[This is not a proof of vacuum uniqueness or the Wilson margin]
\label{fire:YMA20-no-physical-value}
No theorem above constructs an infinite-volume Yang--Mills Gibbs state,
proves independence of boundary conditions, evaluates a response cofactor,
controls the full nonlinear connected tail, or assigns a numerical sign to
the selected Wilson symbol.  The head--tail and Bloch theorems say exactly how
those physical values determine the local card once supplied.  They do not
replace them by geometry, a diagonal limit, or a twisted boundary condition.
\end{firewall}

% =============================================================================
\section{Version V20 proof ledger and replay boundary}
\label{sec:YMA20-ledger}
% =============================================================================

\begingroup\small
\begin{longtable}{>{\raggedright\arraybackslash}p{0.28\textwidth}
                  >{\raggedright\arraybackslash}p{0.19\textwidth}
                  >{\raggedright\arraybackslash}p{0.45\textwidth}}
\caption{Version V20 proof-status ledger.}\label{tab:YMA20-ledger}\\
\toprule
\textbf{Row}&\textbf{Status}&\textbf{Exact scope}\\
\midrule
\endfirsthead
\toprule
\textbf{Row}&\textbf{Status}&\textbf{Exact scope}\\
\midrule
\endhead
UV/IR regulator separation&Proved necessary
&The shell index $r$ and infrared side $N$ are independent; one arbitrary
 diagonal can reverse the fixed-shell infrared limit.\\[0.3em]
Selected local-vacuum law&Identified as upstream loading
&Boundary, toron, and stabilizer sector must be fixed, or uniqueness proved,
 before local expectations define the fused kernel.\\[0.3em]
Geometry-only state selection&Excluded by countertheorem
&The same geometry and action deformation can have opposite curvature signs
 in different selected sectors.\\[0.3em]
Local connected coefficients&Proved quantitatively
&Dual local total variation controls means and covariances, giving every
 finite fused head.\\[0.3em]
Ward head--tail completion&Proved
&Finite-head convergence plus a uniform $\Wtwo$ tail and harmonic control
 gives uniform quotient convergence.\\[0.3em]
Tail failure&Typed exactly
&Failure returns a nonvanishing second-displacement stock outside every fixed
 head.\\[0.3em]
Continuous Bloch symbol&Proved exactly
&Every absolutely summable Hermitian kernel has a continuous finite matrix
 on arbitrary quasi-momentum.\\[0.3em]
Periodization and twist&Proved exactly
&Allowed periodic momenta reproduce the infinite symbol; arbitrary momenta
 are exact finite twisted kernel cards.\\[0.3em]
Ward-quotient sampling&Proved sharply at the displayed level
&$\operatorname{Lip}g\le\Wthree/3$ and
 $\operatorname{Lip}r_0\le23/144$.\\[0.3em]
Bloch-to-periodic transfer&Proved quantitatively
&The error is the local-state defect plus the explicit nearest-grid term in
 \eqref{eq:YMA20-Bloch-periodic-error}.\\[0.3em]
Quantitative diagonal extraction&Proved
&After infrared moduli are known, one can choose a diagonal preserving every
 separated sign on a finite physical bank.\\[0.3em]
V19 global volume theorems&Retained
&They remain necessary for ordinary periodic spectra, mode counting, and
 global finite-volume obstruction sequences.\\[0.3em]
Selected response-matched local state&Open physical loading
&No supplied packet identifies the infinite-volume state or proves boundary
 independence.\\[0.3em]
Fused Wilson head and connected tail&Open physical values
&No selected-state coefficient head or uniform nonlinear tail is populated.\\[0.3em]
Continuous retained Bloch margin&\stopstatus
&The executable reduction is complete, but the Wilson-specific value is
 absent.\\[0.3em]
Local-vacuum contraction and mass&Not inferred
&No state-selection or Bloch-realization theorem supplies higher jets,
 temporal transfer, a reflection-positive continuum, or a mass gap.\\
\bottomrule
\end{longtable}
\endgroup

\begin{theorem}[Version V20 two-regulator correction]
\label{thm:YMA20-terminal}
Versions~V1--V19 are retained at their stated exact, finite, conditional, or
open scopes.  Version~V20 additionally proves:
\begin{enumerate}[label=\textup{(E20.\arabic*)},leftmargin=5.7em]
\item the UV shell/cutoff index and the IR volume side are independent
      regulators;
\item an arbitrary diagonal does not determine the selected fixed-shell local
      vacuum;
\item local-state convergence determines every finite deterministic and
      connected covariance head;
\item a uniform Ward-relative tail upgrades that head convergence to the
      complete continuous local symbol;
\item the local symbol is a continuous Bloch matrix and arbitrary momentum is
      represented exactly by a finite twisted kernel card;
\item ordinary periodic samples approximate it with the explicit error
      \eqref{eq:YMA20-Bloch-periodic-error};
\item a separated local sign yields the same periodic sign after a
      quantitative diagonal is chosen;
\item the V19 volume route remains an independent downstream choice for global
      spectra and mode counting.
\end{enumerate}
No item selects the Yang--Mills vacuum, evaluates the response-matched kernel,
or proves a positive or negative physical margin.
\end{theorem}

\begin{proof}
Items~\textup{(E20.1)--(E20.2)} are
\Cref{def:YMA20-local-vacuum,cth:YMA20-diagonal-no-local-limit,cth:YMA20-geometry-no-state}.
Items~\textup{(E20.3)--(E20.4)} are
\Cref{cor:YMA20-head-convergence,thm:YMA20-head-tail-completion}.
Items~\textup{(E20.5)--(E20.7)} are
\Cref{thm:YMA20-Bloch-periodization,thm:YMA20-W3-Lipschitz,lem:YMA20-reference-Lipschitz,thm:YMA20-Bloch-periodic-transfer,thm:YMA20-diagonal-extraction}.
Item~\textup{(E20.8)} is
\Cref{fire:YMA20-V19-scope,cth:YMA20-local-no-volume-route}.
\end{proof}

\begin{openproblem}[V21: load one selected local vacuum and execute its first continuous Bloch card]
\label{op:YMA20-V21}
Preserve the weak-facing selector, response normalization, pre-retuning
action, cutoff interaction, plaquette writer, local/topological split, colour
metric, field-tangent calibration, carrier thresholds, toron/stabilizer patch,
and flat reference.  Do not append another global index, volume, or carrier
theorem before the following local packet is attempted.
\begin{enumerate}[label=\textup{(V21.\arabic*)},leftmargin=5.4em]
\item Separate the shell index $r$ from an independent infrared side $N$.
      Declare the boundary/topological sector used to approach the local
      vacuum.
\item Prove local convergence or provide an explicit selected DLR/vacuum law
      for the response-matched shell.  If several cluster states remain,
      retain them as distinct branches rather than averaging them after the
      symbol is read.
\item Acquire the two response cofactors and load the complete final local
      law on each surviving branch.
\item Freeze one carrier-supported physical momentum--colour block in a
      guarded annulus and set its continuous covariant quasi-momentum before
      reading the Hessian.
\item Populate one finite fused displacement head from the direct curvature,
      complete same-history score covariance, generated support births, and
      declared local/topological term.
\item Prove or acquire the $\Wtwo$ tail and, for periodic realization, the
      $\Wthree$ sampling stock.  A failed tail returns
      \eqref{eq:YMA20-positive-range-witness}.
\item Evaluate the continuous Bloch margin directly, preferably with the
      twisted finite card of \Cref{thm:YMA20-Bloch-periodization}.  Return
      \pass, an attained \obstruction, or \stopstatus{}.
\item Use a global volume route and the V19 carrier passport only when an
      ordinary periodic witness, full finite-torus spectrum, or thermodynamic
      mode-density statement is additionally intended.
\end{enumerate}
Boundary/Witten ancestry, action-flow conductance, and differentiated cutoff
tails are reopened only after the local mode and state branch have been
frozen for attribution or transport.
\end{openproblem}

\begin{assessment}[V20 termination]
\label{ass:YMA20-termination}
The previous continuation was beginning to organize an ever more detailed
Fourier grid before the local vacuum law itself had been selected.  The two
regulator axes are now separated.  A local response-matched state plus one
finite head and one Ward-relative tail produces a continuous Bloch symbol;
a periodic box is an exact twist or a controlled sample of that symbol, not
its definition.  The first unclosed physical arrow is therefore state
selection and the actual fused connected kernel.  Since neither is supplied,
the selected Wilson margin remains \stopstatus{} and no negative physical
mode has been produced.
\end{assessment}

The accompanying replay verifies only the finite algebra used in this
continuation: the two diagonal counterarrays, the Dirichlet quotient identity
and derivative sum, the Wilson-reference derivative bound, exact Laurent
periodization modulo $z^N-1$, the nearest-grid estimate, and a rational
head--tail/sampling example.  It does not construct a Gibbs state, evaluate a
Wilson path integral, or certify a physical margin.

\section*{Version V20 append-only continuation record}
\addcontentsline{toc}{section}{Version V20 append-only continuation record}
The complete Version~V19 source above is preserved byte for byte up to its
sole terminal document-closing command.  Its SHA--256 digest is
\begin{center}\scriptsize\ttfamily
4d6645759c619e50cc4391bb963e0368b0739462d3a47e06cfc16922d018ebad.
\end{center}
For Version~V21, preserve this accumulated source, append new mathematical
material immediately before the final document-closing command, and increment
the filename to \texttt{\_V21.tex}.  Keep all previously frozen physical
action, response, gauge, topology, colour, tangent, carrier, and reference
choices.  In addition, keep the shell index and infrared volume side distinct,
freeze the selected local-vacuum branch before the fused kernel is formed,
and use the continuous Bloch card before choosing an ordinary periodic
realization.

% =============================================================================
% END APPENDED VERSION V20 MATERIAL
% =============================================================================


% =============================================================================
% BEGIN APPENDED VERSION V21 MATERIAL
% =============================================================================

\clearpage
\hypersetup{
  pdftitle={Renewal Yang--Mills Reduced-Fibre Wilson Shell Symbol and Local-Vacuum Programme V21},
  pdfsubject={Phase selection before response matching, the exact mixture debit, and a deformation-stable local Wilson vacuum}
}
\providecommand{\DLR}{\mathrm{DLR}}
\providecommand{\Ph}{\mathrm{Ph}}

\section*{Version V21: phase selection precedes response matching and the local Bloch Hessian}
\addcontentsline{toc}{section}{Version V21: phase selection before response matching and the local Bloch Hessian}

\begin{abstract}
Version~V20 placed the selected infinite-volume local law upstream of a
periodic Fourier realization.  That correction is necessary, but its terminal
order still froze the response-matched action before deciding whether the
pre-retuning local law has one phase or several.  At phase coexistence those
operations need not commute.  A coefficient which is response null in a
mixture need not be response null in any constituent phase, and the Hessian
of the mixture is not the average of the phase Hessians: it contains an
additional negative between-phase score Gram.

This continuation first proves a fixed-shell existence theorem for compact
reduced-fibre Gibbs specifications with an absolutely summable Wilsonian
interaction.  The set of translation-invariant DLR states is nonempty,
compact, and convex, so extreme state branches exist whenever the submitted
generated tail has the required summability.  Existence does not select a
branch or make its field deformation differentiable.

For an arbitrary state decomposition
$\mu=\int\mu_\theta\,\pi(\dd\theta)$, the complete score--Hessian card obeys
\[
 H_\mu=\int H_\theta\,\pi(\dd\theta)
       -\operatorname{Cov}_\pi(m_\theta),
 \qquad m_\theta=\mathbb E_{\mu_\theta}J.
\]
The final term is positive and is subtracted.  On translation-invariant
phases it appears as a constant connected covariance at every separation.
Unless the phase score means agree, the mixed covariance has no summable
second- or third-displacement tail and the continuous Ward quotient of
Version~V20 is not the ordinary local kernel of the mixture.  On a periodic
torus the same term is supported exactly at zero momentum, so it is a typed
phase/harmonic survivor rather than an ordinary nonzero local-curvature mode.

The colour content is explicit.  If a gauge-fixed phase has nonzero adjoint
score mean $m\in\mathfrak{su}(3)$, averaging its global-conjugation orbit
produces zero mean but the exact between-orbit Gram
\[
 \frac{\|m\|^2}{8}I_{\mathfrak{su}(3)}.
\]
Thus restoring the residual colour symmetry can create a nondecaying
zero-momentum covariance even when every orbit component has a summable
connected kernel.  No claim is made that such a broken-colour phase occurs;
the theorem identifies its exact consequence if it does.

A separate two-phase cylinder shows that phase-conditioned curvatures may
all vanish while the mixed finite-volume curvature diverges and the limiting
pressure has a cusp.  Weak local convergence at the base action is therefore
insufficient.  The needed object is a \emph{differentiably selected phase
chart}: one fixed boundary/sector branch, locally uniform convergence under
the physical deformation, and a uniformly summable susceptibility kernel.
Under these conditions the finite-volume linear-response identity passes to
the selected DLR state, the response cofactors are phase specific, and the
physical direct-curvature-minus-covariance kernel is the derivative of that
same branch.  Only then may the V20 head--tail and continuous Bloch theorems
be invoked.

No selected Wilson phase, response cofactor, connected tail, or Bloch margin
is numerically populated here.  The advance is the exact upstream order and
the identification of the positive phase-mixture debit.  The present
physical verdict remains \stopstatus{}.
\end{abstract}

% =============================================================================
\section{The fixed-shell Gibbs specification and its nonempty phase set}
\label{sec:YMA21-DLR}
% =============================================================================

The first task is not to choose a preferred infinite-volume state by name.
It is to state the minimum finite-shell hypotheses under which at least one
local Gibbs state exists.  The theorem below is applied only to the already
reduced, gauge-fixed shell coordinates.  It does not reconstruct that
reduction or the response-matched interaction.

\begin{definition}[Compact reduced-cell interaction]
\label{def:YMA21-interaction}
Fix a shell $r$.  Let $\mathsf X_r$ be a compact metric cell space carrying a
full-support probability measure $m_r$; for the literal parity-tree fibre it
is a finite product of copies of $\SU(3)$ with product Haar measure.  Put
\[
 \Omega_r=\mathsf X_r^{\mathbb Z^4}
\]
with its product topology and shift action $(\tau_x)_{x\in\mathbb Z^4}$.

A translation-covariant interaction is a family
$\boldsymbol\Phi=(\Phi_A)$ indexed by finite nonempty
$A\Subset\mathbb Z^4$, where each $\Phi_A$ is a real continuous cylinder
function on $\Omega_r$ and
\[
 \Phi_{A+x}=\Phi_A\circ\tau_{-x}.
\]
It is \emph{uniformly absolutely summable} when
\begin{equation}
 \boxed{
 \|\boldsymbol\Phi\|_{\rm abs}
 :=\sum_{A\ni0}\|\Phi_A\|_\infty<\infty.}
 \label{eq:YMA21-abs-interaction}
\end{equation}
For finite $\Lambda\Subset\mathbb Z^4$ and exterior configuration $\eta$,
define
\begin{equation}
 H_\Lambda^{\boldsymbol\Phi}(z_\Lambda\mid\eta)
 :=\sum_{A:\,A\cap\Lambda\ne\varnothing}
   \Phi_A(z_\Lambda\eta_{\Lambda^c}),
 \label{eq:YMA21-finite-Hamiltonian}
\end{equation}
and the conditional kernel
\begin{equation}
 \gamma_\Lambda^{\boldsymbol\Phi}(\dd z\mid\eta)
 =\frac{
   e^{-H_\Lambda^{\boldsymbol\Phi}(z_\Lambda\mid\eta)}
   m_r^{\otimes\Lambda}(\dd z_\Lambda)
   \delta_{\eta_{\Lambda^c}}(\dd z_{\Lambda^c})}
  {Z_\Lambda^{\boldsymbol\Phi}(\eta)}.
 \label{eq:YMA21-specification}
\end{equation}
A probability measure $\mu$ on $\Omega_r$ is a DLR state when
$\mu\gamma_\Lambda^{\boldsymbol\Phi}=\mu$ for every finite $\Lambda$.
\end{definition}

The sum in \eqref{eq:YMA21-finite-Hamiltonian} is absolutely and uniformly
convergent.  Indeed, every interaction set meeting $\Lambda$ contains at
least one $x\in\Lambda$, and translation covariance gives
\begin{equation}
 \sum_{A:\,A\cap\Lambda\ne\varnothing}\|\Phi_A\|_\infty
 \le |\Lambda|\|\boldsymbol\Phi\|_{\rm abs}.
 \label{eq:YMA21-Hamiltonian-bound}
\end{equation}
Terms independent of $z_\Lambda$ may be removed without changing the kernel.

\begin{theorem}[Existence and compactness of fixed-shell DLR phases]
\label{thm:YMA21-DLR-existence}
For every interaction of \Cref{def:YMA21-interaction}:
\begin{enumerate}[label=\textup{(D\arabic*)},leftmargin=4.0em]
\item $(\gamma_\Lambda^{\boldsymbol\Phi})_\Lambda$ is a proper, consistent,
      quasilocal specification;
\item the DLR set $\mathcal G_r(\boldsymbol\Phi)$ is nonempty, compact, and
      convex in the weak topology;
\item the translation-invariant DLR set
      $\mathcal G_r^{\rm tr}(\boldsymbol\Phi)$ is nonempty, compact, and
      convex;
\item $\mathcal G_r^{\rm tr}(\boldsymbol\Phi)$ has extreme points.
\end{enumerate}
Accordingly, the pure Wilson finite-range shell has at least one
translation-invariant DLR state.  The same conclusion holds for the complete
response-matched shell whenever its generated interaction has a
representation satisfying \eqref{eq:YMA21-abs-interaction}.
\end{theorem}

\begin{proof}
Uniform absolute convergence in \eqref{eq:YMA21-Hamiltonian-bound} makes the
finite Hamiltonian continuous in both the interior variables and the product
boundary topology.  Its exponential is strictly positive and continuous on
a compact space; hence the normalizer is positive and continuous.  This
proves quasilocality.  Properness is immediate because the kernel fixes the
exterior.  If $\Delta\subseteq\Lambda$, Fubini's theorem and cancellation of
all terms which do not depend on $z_\Delta$ show
$\gamma_\Lambda\gamma_\Delta=\gamma_\Lambda$; this is the ordinary Gibbs
conditionalization identity and proves consistency.

Choose increasing boxes $\Lambda_n\uparrow\mathbb Z^4$ and one exterior
configuration $\eta$.  Extend
$\gamma_{\Lambda_n}^{\boldsymbol\Phi}(\cdot\mid\eta)$ to a probability
measure $\mu_n$ on $\Omega_r$ by retaining $\eta$ outside $\Lambda_n$.
The countable product of compact metric spaces is compact and metrizable, so
its probability measures are weakly compact.  Let $\mu$ be a cluster point.
For a finite $\Delta$ and continuous local $f$, consistency gives
\[
 \mu_n(\gamma_\Delta f)=\mu_n(f)
\]
whenever $\Delta\subseteq\Lambda_n$.  The function $\gamma_\Delta f$ is
continuous by quasilocality.  Passing to the cluster point yields
$\mu(\gamma_\Delta f)=\mu(f)$.  Continuous local functions determine
probability measures, so $\mu$ is DLR.  This proves nonemptiness.

The DLR equations are affine and weakly closed because every
$\gamma_\Delta f$ is continuous.  Hence the DLR set is compact and convex.
Translations preserve it.  Starting from any DLR state, average its
translates over boxes:
\[
 \overline\mu_n
 =|\Lambda_n|^{-1}\sum_{x\in\Lambda_n}\tau_x\mu.
\]
Every average is DLR.  Any weak cluster point is DLR, and the boundary-to-
volume ratio shows that it is translation invariant.  This proves the third
claim.  Finally a nonempty compact convex subset of the locally convex space
of signed measures has an extreme point by the Krein--Milman theorem.
\end{proof}

\begin{firewall}[Existence is not selection or differentiability]
\label{fire:YMA21-DLR-scope}
The theorem proves nonemptiness only after the complete generated action has
been shown to satisfy the displayed summability condition.  It does not prove
uniqueness, clustering, a response derivative, stability under the marginal
cofactor deformation, or equality of different boundary-condition limits.
An arbitrary extreme point of the compact convex set is not automatically the
physical vacuum selected by the inherited boundary, toron, stabilizer, and
reflection packet.
\end{firewall}

% =============================================================================
\section{Phase selection and response matching do not commute}
\label{sec:YMA21-response-mixture}
% =============================================================================

Let $(\Theta,\pi)$ be a probability space and let
\begin{equation}
 \mu=\int_\Theta\mu_\theta\,\pi(\dd\theta)
 \label{eq:YMA21-state-decomposition}
\end{equation}
be any barycentric decomposition of one state.  No uniqueness of the
decomposition is assumed in the algebraic identities below.

\begin{theorem}[Exact response decomposition across phases]
\label{thm:YMA21-response-decomposition}
Let $S,P,\Psi\in L^2(\mu;\mathbb R)$.  Put
\[
 s_\theta=\mathbb E_\theta S,
 \qquad p_\theta=\mathbb E_\theta P,
 \qquad \psi_\theta=\mathbb E_\theta\Psi,
\]
and define the phase responses
\[
 d_\theta=-\Cov_\theta(S,P),
 \qquad
 n_\theta=-\Cov_\theta(S,\Psi).
\]
Then the responses in the mixed state are
\begin{equation}
 \boxed{
 \begin{aligned}
 d_\mu
 &:=-\Cov_\mu(S,P)
   =\int d_\theta\,\pi(\dd\theta)
    -\Cov_\pi(s_\theta,p_\theta),\\
 n_\mu
 &:=-\Cov_\mu(S,\Psi)
   =\int n_\theta\,\pi(\dd\theta)
    -\Cov_\pi(s_\theta,\psi_\theta).
 \end{aligned}}
 \label{eq:YMA21-response-total-covariance}
\end{equation}
For $I_b=\Psi+bP$ one has the exact identity
\begin{equation}
 \boxed{
 \mathscr R_\mu(I_b)
 =\int\bigl(n_\theta+bd_\theta\bigr)\,\pi(\dd\theta)
  -\Cov_\pi\bigl(s_\theta,\psi_\theta+bp_\theta\bigr).}
 \label{eq:YMA21-mixed-response-null}
\end{equation}
Thus a coefficient which is response null in $\mu$ need not be response null
in any component $\mu_\theta$.
\end{theorem}

\begin{proof}
The law of total covariance gives
\[
 \Cov_\mu(S,P)
 =\int\Cov_\theta(S,P)\,\pi(\dd\theta)
  +\Cov_\pi(s_\theta,p_\theta).
\]
Multiplying by minus one proves the first formula.  The second is identical.
Linearity in the insertion $I_b$ gives
\eqref{eq:YMA21-mixed-response-null}.
\end{proof}

\begin{countertheorem}[A mixture-matched coefficient need not match either phase]
\label{cth:YMA21-response-noncommutation}
There is a four-point positive cylinder with two equally weighted phases for
which
\[
 d_+=d_-=1,
 \qquad
 n_+=1,
 \qquad
 n_-=-1,
\]
so the two phase coefficients are
\[
 b_+=-1,
 \qquad
 b_-=1,
\]
whereas the mixed coefficient is $b_\mu=0$.
\end{countertheorem}

\begin{proof}
Let $\theta\in\{-1,+1\}$ and $X\in\{-1,+1\}$ be independent uniform
variables.  Condition on the phase $\theta$ and set
\[
 S=X,
 \qquad P=-X,
 \qquad \Psi=-\theta X.
\]
All component means vanish.  Hence
\[
 -\Cov_\theta(S,P)=1,
 \qquad
 -\Cov_\theta(S,\Psi)=\theta.
\]
The phase coefficients are $-n_\theta/d_\theta=-\theta$.  In the equal
mixture, $d_\mu=1$ and $n_\mu=0$, so $b_\mu=0$.  The mixed interaction is
response null only after the two opposite phase responses have cancelled.
\end{proof}

\begin{corollary}[Exact commutation criterion for a common coefficient]
\label{cor:YMA21-response-commutation}
A number $b$ which is response null in every component is response null in
the mixture if and only if
\begin{equation}
 \boxed{
 \Cov_\pi\bigl(s_\theta,\psi_\theta+bp_\theta\bigr)=0.}
 \label{eq:YMA21-between-phase-response-zero}
\end{equation}
Conversely, mixed response nullity alone implies only that the integral of the
component responses equals the between-phase covariance in
\eqref{eq:YMA21-mixed-response-null}; it does not imply phasewise nullity.
\end{corollary}

\begin{proof}
If $n_\theta+bd_\theta=0$ almost surely, the first term on the right of
\eqref{eq:YMA21-mixed-response-null} vanishes, so the remaining condition is
exactly \eqref{eq:YMA21-between-phase-response-zero}.  The converse statement
is the same identity read without setting either term separately to zero.
\end{proof}

\begin{assessment}[Correct order of the marginal card]
\label{ass:YMA21-response-order}
At coexistence, the admissible order is
\[
 \boxed{
 \begin{gathered}
 \text{select the physical phase branch}
 \longrightarrow
 (\mathfrak c_{P,\theta},\mathfrak c_{\Psi,\theta})\\
 \longrightarrow b_\theta
 \longrightarrow
 \text{continue the same phase under the final action}
 \end{gathered}}
\]
One may reverse the first two arrows only after proving that the response
cofactor ratio is phase independent and that the between-phase term
\eqref{eq:YMA21-between-phase-response-zero} vanishes.  The finite periodic
cofactor of Versions~V13--V16 remains exact for the finite law on which it is
computed; it is not silently promoted to every infinite-volume phase.
\end{assessment}

% =============================================================================
\section{The exact between-phase debit in the physical Hessian}
\label{sec:YMA21-mixture-Hessian}
% =============================================================================

\begin{theorem}[Total-covariance decomposition of the signed Wilson card]
\label{thm:YMA21-Hessian-mixture}
Let $J\in L^2(\mu;\mathbb C^d)$ be the complete field-score column and let
$K\in L^1(\mu;\operatorname{Herm}_d)$ be the complete deterministic second
writer, both already assembled in one physical action chart.  Define
\[
 m_\theta=\mathbb E_\theta J,
 \qquad
 \overline m=\int m_\theta\,\pi(\dd\theta),
\]
\[
 H_\theta=\mathbb E_\theta K-\Cov_\theta(J),
 \qquad
 H_\mu=\mathbb E_\mu K-\Cov_\mu(J).
\]
Then
\begin{equation}
 \boxed{
 H_\mu
 =\int H_\theta\,\pi(\dd\theta)-\mathbb B_\pi,}
 \qquad
 \boxed{
 \mathbb B_\pi
 =\int(m_\theta-\overline m)(m_\theta-\overline m)^*
        \,\pi(\dd\theta)\succeq0.}
 \label{eq:YMA21-Hessian-mixture}
\end{equation}
Moreover,
\begin{equation}
 \boxed{
 \mathbb B_\pi=0
 \iff
 m_\theta=\overline m
 \quad\text{for $\pi$-almost every $\theta$.}}
 \label{eq:YMA21-phase-mean-zero-criterion}
\end{equation}
\end{theorem}

\begin{proof}
For a vector-valued random variable, total covariance is
\[
 \Cov_\mu(J)
 =\int\Cov_\theta(J)\,\pi(\dd\theta)
  +\mathbb B_\pi.
\]
The deterministic expectation is affine in the state.  Subtracting proves
\eqref{eq:YMA21-Hessian-mixture}.  Finally
\[
 \Tr\mathbb B_\pi
 =\int\|m_\theta-\overline m\|^2\,\pi(\dd\theta),
\]
so positivity gives the equivalence in
\eqref{eq:YMA21-phase-mean-zero-criterion}.
\end{proof}

\begin{corollary}[Translation-invariant kernel and the harmonic phase term]
\label{cor:YMA21-kernel-mixture}
Assume every $\mu_\theta$ is translation invariant and let
$J_x=J_0\circ\tau_x$.  Put
\[
 C_\theta(x)
 =\Cov_\theta(J_0,J_x),
 \qquad
 C_\mu(x)=\Cov_\mu(J_0,J_x).
\]
Then for every $x\in\mathbb Z^4$,
\begin{equation}
 \boxed{
 C_\mu(x)
 =\int C_\theta(x)\,\pi(\dd\theta)+\mathbb B_\pi.}
 \label{eq:YMA21-kernel-mixture}
\end{equation}
If
\begin{equation}
 \sum_{x\in\mathbb Z^4}(1+\|x\|_1^k)
 \int\|C_\theta(x)\|_{\rm op}\,\pi(\dd\theta)<\infty
 \label{eq:YMA21-component-tail}
\end{equation}
for some $k\ge0$, then $\mathbb B_\pi\ne0$ implies that the mixed covariance
has no absolutely summable $k$th displacement stock.

On a periodic torus $\Lambda_N$, using the unnormalized Fourier convention
\[
 \widehat C_N(q)=\sum_{x\in\Lambda_N}e^{-\mathrm iq\cdot x}C_N(x),
\]
the constant phase term contributes exactly
\begin{equation}
 \boxed{
 \sum_{x\in\Lambda_N}e^{-\mathrm iq\cdot x}\mathbb B_\pi
 =N^4\mathbb B_\pi\,\one_{\{q=0\}}.}
 \label{eq:YMA21-phase-zero-mode}
\end{equation}
Thus the between-phase debit is a zero-momentum/harmonic survivor on each
finite periodic realization, while it prevents an ordinary continuous
summable-kernel extension through $q=0$.
\end{corollary}

\begin{proof}
Translation invariance makes the component mean of $J_x$ equal to
$m_\theta$ for every $x$.  Applying total covariance to $(J_0,J_x)$ gives
\eqref{eq:YMA21-kernel-mixture}.  Under
\eqref{eq:YMA21-component-tail}, the averaged connected term tends to zero as
$\|x\|_1\to\infty$.  If $\mathbb B_\pi\ne0$, the mixed kernel therefore tends
to a nonzero matrix and cannot have a finite weighted absolute sum.
Equation \eqref{eq:YMA21-phase-zero-mode} is the character orthogonality
identity on the finite abelian group $\Lambda_N$.
\end{proof}

\begin{firewall}[A phase debit is not automatically a nonzero pure-mode failure]
\label{fire:YMA21-phase-not-mode}
The constant term in \eqref{eq:YMA21-kernel-mixture} is killed by every
nonzero ordinary periodic character.  Its presence obstructs the summable
local-vacuum kernel and creates a macroscopic zero-mode susceptibility; it
does not by itself make a retained nonzero momentum--colour block negative.
A pure-mode obstruction still requires the assembled direct interval on that
block.
\end{firewall}

% =============================================================================
\section{Residual colour restoration and its exact adjoint debit}
\label{sec:YMA21-colour-orbit}
% =============================================================================

The reduced parity-tree section retains a global conjugation action at a
central background.  The following theorem identifies exactly what happens
if one first chooses a phase which is not invariant under that action and
then restores the symmetry by orbit averaging.

\begin{theorem}[Exact $\SU(3)$ orbit-mixture Gram]
\label{thm:YMA21-SU3-orbit}
Let
\[
 V=\mathfrak{su}(3)
\]
with real inner product $\ip{X}{Y}=-\Tr(XY)$, so $\dim_\mathbb R V=8$.  Let
$\mu$ be a translation-invariant state and suppose an adjoint-equivariant
score field has one-point mean
\[
 m=\mathbb E_\mu J_0\in V.
\]
For $g\in\SU(3)$ let $\mu_g$ be the globally conjugated state, and let
\[
 \overline\mu=\int_{\SU(3)}\mu_g\,\dd g
\]
be its Haar orbit average.  Then
\begin{equation}
 \boxed{
 \mathbb E_{\overline\mu}J_0=0,}
 \qquad
 \boxed{
 \mathbb B_{\rm orb}
 =\int_{\SU(3)}
   (\Ad_gm)(\Ad_gm)^*\,\dd g
 =\frac{\|m\|^2}{8}I_V.}
 \label{eq:YMA21-SU3-orbit-Gram}
\end{equation}
Consequently, if the orbit-component connected kernels have an averaged
summable displacement tail, the colour-restored state has the constant
long-range covariance floor $\|m\|^2I_V/8$ and hence the zero-momentum phase
survivor of \Cref{cor:YMA21-kernel-mixture}.
\end{theorem}

\begin{proof}
The Haar average
$\int\Ad_gm\,\dd g$ is invariant under every adjoint transformation.  Since
$\mathfrak{su}(3)$ is simple, its adjoint representation has no nonzero
invariant vector, so the average is zero.

Let
\[
 T=\int(\Ad_gm)(\Ad_gm)^*\,\dd g.
\]
Haar invariance gives $\Ad_hT\Ad_h^*=T$ for every $h$.  The self-adjoint
eigenspaces of $T$ are therefore adjoint invariant.  Any real adjoint-
invariant subspace of $\mathfrak{su}(3)$ is an ideal, hence is either zero or
all of $V$.  Thus $T=cI_V$.  Taking traces gives
\[
 8c=\Tr T=\int\|\Ad_gm\|^2\,\dd g=\|m\|^2,
\]
which proves \eqref{eq:YMA21-SU3-orbit-Gram}.  The final assertion is
\Cref{cor:YMA21-kernel-mixture}.
\end{proof}

\begin{remark}[No spontaneous-breaking assertion]
\label{rem:YMA21-no-breaking-claim}
The theorem is conditional on a component with $m\ne0$.  It neither proves
nor assumes that the physical Yang--Mills vacuum breaks residual colour
symmetry.  On a globally conjugation-invariant phase with $m=0$, this
particular orbit debit vanishes.  The finite card must test that fact in the
same gauge-fixed and sector-selected law used by the Wilson Hessian.
\end{remark}

% =============================================================================
\section{A phase-coexistence cylinder with divergent mixed curvature}
\label{sec:YMA21-coexistence}
% =============================================================================

\begin{countertheorem}[Phasewise flat curvature and a divergent mixed Hessian]
\label{cth:YMA21-coexistence-cusp}
For every integer volume $V\ge1$ and $a>0$, consider two one-point phase
cylinders with actions
\[
 F_{+,V}(t)=-Vat,
 \qquad
 F_{-,V}(t)=Vat.
\]
Each phase effective action has zero second derivative.  The unselected
mixture has partition function and effective action
\[
 Z_V(t)=e^{Vat}+e^{-Vat}=2\cosh(Vat),
 \qquad
 \Phi_V(t)=-\log Z_V(t),
\]
and satisfies
\begin{equation}
 \boxed{
 \Phi_V''(0)=-V^2a^2.}
 \label{eq:YMA21-divergent-curvature}
\end{equation}
Moreover,
\[
 -V^{-1}\log Z_V(t)\longrightarrow-|at|,
\]
which is not differentiable at $t=0$.
\end{countertheorem}

\begin{proof}
The two phase effective actions are $F_{+,V}$ and $F_{-,V}$, both affine in
$t$.  Direct differentiation gives
\[
 \Phi_V''(0)
 =-\left.\frac{\dd^2}{\dd t^2}\log\cosh(Vat)\right|_{t=0}
 =-V^2a^2.
\]
Finally
\[
 \frac1V\log(2\cosh(Vat))\longrightarrow|at|
\]
by the elementary log-sum-exp limit.  The negative of this function has a
cusp at zero.
\end{proof}

\begin{assessment}[Meaning of the coexistence model]
\label{ass:YMA21-coexistence-meaning}
The countermodel is not a claim about the phase diagram of four-dimensional
Wilson theory.  It proves a logical point needed by the programme: local
convergence at the undeformed action and existence of phase-conditioned
Hessians do not determine the curvature of an unselected mixture.  The
between-phase first-derivative variance can be volume squared and can destroy
twice differentiable pressure.  A selected local-vacuum symbol therefore
requires a deformation-stable phase branch, not merely one weak cluster state
at the base parameter.
\end{assessment}

% =============================================================================
\section{Differentiably selected phase charts and the physical Kubo row}
\label{sec:YMA21-phase-chart}
% =============================================================================

\begin{definition}[Differentiably selected local-vacuum chart]
\label{def:YMA21-differentiable-phase}
Fix a shell $r$; fix boundary, topological, toron, and stabilizer data
$\mathfrak b$; and let $t\in[-t_0,t_0]$.  Let
$\mu_{N,t}^{\mathfrak b}$ be the corresponding finite-volume Gibbs law and
let $\mu_t^{\Ph}$ be a translation-invariant DLR state.  The local writers
$A_t$ in the target bank $\mathcal A$ are continuously differentiable in
$t$, uniformly bounded, and have fixed finite support.  The family is a
\emph{differentiably selected phase chart} when:
\begin{enumerate}[label=\textup{(P\arabic*)},leftmargin=4.0em]
\item on every finite cylinder, the restrictions of
      $\mu_{N,t}^{\mathfrak b}$ converge to those of $\mu_t^{\Ph}$ in dual
      total variation, locally uniformly in $t$;
\item the finite-volume action derivative has the decomposition
      \begin{equation}
       \dot F_{N,t}
       =\sum_{x\in\Lambda_N}J_{x,t}+R_{N,t},
       \label{eq:YMA21-source-decomposition}
      \end{equation}
      where $J_{x,t}=J_{0,t}\circ\tau_x$ is a uniformly bounded local writer;
\item for every $A_t\in\mathcal A$,
      \begin{equation}
       \sup_{N,\,|t|\le t_0}
       \sum_{x\in\Lambda_N}
       \left|\Cov_{\mu_{N,t}^{\mathfrak b}}(A_t,J_{x,t})\right|<\infty,
       \label{eq:YMA21-susceptibility-sum}
      \end{equation}
      and the corresponding tails vanish uniformly as the distance from
      $\operatorname{supp}A_t$ tends to infinity;
\item the boundary remainder is asymptotically invisible to the bank:
      \begin{equation}
       \sup_{|t|\le t_0}
       \left|\Cov_{\mu_{N,t}^{\mathfrak b}}(A_t,R_{N,t})\right|
       \longrightarrow0.
       \label{eq:YMA21-boundary-response-vanish}
      \end{equation}
\end{enumerate}
All writers and the boundary branch are fixed before their expectations are
inspected.
\end{definition}

\begin{theorem}[Selected-phase linear response]
\label{thm:YMA21-selected-Kubo}
On a differentiably selected phase chart, every $A_t\in\mathcal A$ is
differentiable in expectation and
\begin{equation}
 \boxed{
 \frac{\dd}{\dd t}\mathbb E_{\mu_t^{\Ph}}A_t
 =\mathbb E_{\mu_t^{\Ph}}\dot A_t
  -\sum_{x\in\mathbb Z^4}
   \Cov_{\mu_t^{\Ph}}(A_t,J_{x,t}).}
 \label{eq:YMA21-selected-Kubo}
\end{equation}
The series is absolutely and locally uniformly convergent in $t$.
\end{theorem}

\begin{proof}
At finite volume, differentiation under the compact integral gives the exact
identity
\[
 \frac{\dd}{\dd t}\mathbb E_{\mu_{N,t}^{\mathfrak b}}A_t
 =\mathbb E_{\mu_{N,t}^{\mathfrak b}}\dot A_t
  -\Cov_{\mu_{N,t}^{\mathfrak b}}(A_t,\dot F_{N,t}).
\]
Insert \eqref{eq:YMA21-source-decomposition}.  Local convergence passes every
fixed covariance to the selected DLR state.  The uniform tail in
\eqref{eq:YMA21-susceptibility-sum} permits passage through the infinite sum,
and \eqref{eq:YMA21-boundary-response-vanish} removes the boundary remainder.
The resulting finite-volume derivatives converge locally uniformly to the
right side of \eqref{eq:YMA21-selected-Kubo}.  Integrating those derivatives
between two parameter values and passing to the limit proves that the limiting
expectation is differentiable with the displayed derivative.
\end{proof}

\begin{corollary}[Phase-specific score--Hessian identity]
\label{cor:YMA21-phase-Hessian}
Assume the phase chart contains the complete first and second field jets of a
finite supported coefficient bank.  Let $J_v$ and $K_{v,w}$ denote the first
and second action derivatives in coefficient directions $v,w$.  Then the
selected phase curvature is
\begin{equation}
 \boxed{
 H_\Ph[v,w]
 =\mathbb E_{\mu_0^\Ph}K_{v,w}
  -\Cov_{\mu_0^\Ph}(J_v,J_w).}
 \label{eq:YMA21-phase-score-Hessian}
\end{equation}
If the translation kernel of the two terms has the $\Wtwo$ and $\Wthree$
stocks required in Version~V20, its continuous Bloch symbol is the derivative
of this same selected phase branch and all V20 head--tail and sampling
certificates apply.
\end{corollary}

\begin{proof}
Apply \Cref{thm:YMA21-selected-Kubo} to the first field derivative in the
second coefficient direction.  The direct derivative is $K_{v,w}$ and the
change of the law is the covariance with $J_w$.  This gives
\eqref{eq:YMA21-phase-score-Hessian}.  The Bloch assertion is exactly the
head--tail and continuous-symbol theorem of Version~V20 applied to this
branch-specific kernel.
\end{proof}

\begin{firewall}[Weak DLR convergence at one parameter is insufficient]
\label{fire:YMA21-weak-not-differentiable}
A DLR state $\mu_0$ obtained at the undeformed action does not determine
$\frac{\dd}{\dd t}\mu_t|_{t=0}$.  Without a fixed phase continuation and the
susceptibility tail of \Cref{def:YMA21-differentiable-phase}, one may not
identify a limit of finite-volume partition curvatures with
\eqref{eq:YMA21-phase-score-Hessian}.  The cusp in
\Cref{cth:YMA21-coexistence-cusp} is an explicit failure.
\end{firewall}

% =============================================================================
\section{Phase-conditioned response matching and final-law assembly}
\label{sec:YMA21-phase-matching}
% =============================================================================

\begin{theorem}[Phase-conditioned marginal coefficient]
\label{thm:YMA21-phase-marginal}
Let $\mu_t^\Ph$ be a differentiably selected pre-retuning phase chart carrying
the physical Creutz response source $S_{\rm Cr}$, plaquette writer $P$, and
cutoff interaction $\Psi^{\rm cut}$.  Define
\[
 d_P^\Ph=-\Cov_{\mu_0^\Ph}(S_{\rm Cr},P),
 \qquad
 n_\Psi^\Ph=-\Cov_{\mu_0^\Ph}(S_{\rm Cr},\Psi^{\rm cut}).
\]
On the regular branch $d_P^\Ph\ne0$, the phase-specific native coefficient is
\begin{equation}
 \boxed{
 b_\Ph=-\frac{n_\Psi^\Ph}{d_P^\Ph}.}
 \label{eq:YMA21-phase-b}
\end{equation}
Assume that $b_\Ph$ lies in a certified deformation interval on which the
same boundary-selected phase chart persists.  The physical final law is then
the continuation of that \emph{same} phase under
\[
 F_{b_\Ph}=F_0+b_\Ph F_P.
\]
Its selected shell Hessian is
\begin{equation}
 \boxed{
 H_{b_\Ph}^{\Ph}
 =\mathbb E_{\mu_{b_\Ph}^{\Ph}}K_{b_\Ph}
  -\Cov_{\mu_{b_\Ph}^{\Ph}}(J_{b_\Ph}).}
 \label{eq:YMA21-phase-final-Hessian}
\end{equation}
A coefficient obtained from a mixed finite-volume law may replace
$b_\Ph$ only after the commutation criterion of
\Cref{cor:YMA21-response-commutation} and deformation stability of the phase
have been verified.
\end{theorem}

\begin{proof}
Equation \eqref{eq:YMA21-phase-b} is the unique scalar which makes
\[
 -\Cov_{\mu_0^\Ph}
 \bigl(S_{\rm Cr},\Psi^{\rm cut}+bP\bigr)=0.
\]
The selected phase chart specifies which continuation of the local Gibbs law
is used after changing the action.  Applying
\Cref{cor:YMA21-phase-Hessian} in that chart gives
\eqref{eq:YMA21-phase-final-Hessian}.  The final qualification follows from
\Cref{thm:YMA21-response-decomposition}.
\end{proof}

\begin{corollary}[When response matching and phase averaging do commute]
\label{cor:YMA21-full-commutation}
Suppose there is one number $b$ such that, for $\pi$-almost every phase,
\[
 n_\theta+bd_\theta=0,
 \qquad
 \mathbb E_\theta J_b=m_b
\]
with the same vector $m_b$, and suppose
\eqref{eq:YMA21-between-phase-response-zero} holds.  Then the common
coefficient is also response null in the mixture and
\[
 H_{b,\mu}=\int H_{b,\theta}\,\pi(\dd\theta).
\]
If any one of these conditions fails, the mixture and phase-conditioned cards
must remain distinct.
\end{corollary}

\begin{proof}
The response statement is
\Cref{cor:YMA21-response-commutation}.  Equality of the score means makes the
between-phase Hessian Gram in \Cref{thm:YMA21-Hessian-mixture} vanish.  The
remaining Hessian is therefore the phase average.
\end{proof}

% =============================================================================
\section{Finite phase card, robust mixture enclosure, and termination}
\label{sec:YMA21-finite-card}
% =============================================================================

Even when a mixed state is the declared physical target, the additional debit
can be acquired without constructing a new source atlas.

\begin{proposition}[Positive two-copy acquisition of the phase Gram]
\label{prop:YMA21-phase-Gram-replica}
Let $\Theta$ denote a phase label in a finite or countable decomposition and
let
\[
 m_\Theta=\mathbb E[J\mid\Theta].
\]
For independent phase labels $\Theta_1,\Theta_2$ with law $\pi$,
\begin{equation}
 \boxed{
 \mathbb B_\pi
 =\frac12\mathbb E
  \left[(m_{\Theta_1}-m_{\Theta_2})
        (m_{\Theta_1}-m_{\Theta_2})^*\right].}
 \label{eq:YMA21-phase-Gram-replica}
\end{equation}
The writer is positive semidefinite sample by sample.
\end{proposition}

\begin{proof}
Expand the right side, use independence, and collect
$\mathbb E[m_\Theta m_\Theta^*]-\overline m\,\overline m^*$.
\end{proof}

\begin{proposition}[Outward error for a finite phase mixture]
\label{prop:YMA21-mixture-error}
Let $\pi_1,\ldots,\pi_s$ be exact phase weights.  Suppose
\[
 \|H_i-\widehat H_i\|\le\varepsilon_{H,i},
 \qquad
 \|m_i-\widehat m_i\|\le\varepsilon_m
\]
for every phase, and put
\[
 \widehat{\overline m}=\sum_i\pi_i\widehat m_i,
 \qquad
 \widehat B
 =\sum_i\pi_i\widehat m_i\widehat m_i^*
  -\widehat{\overline m}\widehat{\overline m}^*,
\]
\[
 \widehat H_{\rm mix}
 =\sum_i\pi_i\widehat H_i-\widehat B,
 \qquad
 M_m=\max_i\|\widehat m_i\|.
\]
Then
\begin{equation}
 \boxed{
 \|H_{\rm mix}-\widehat H_{\rm mix}\|
 \le
 \sum_i\pi_i\varepsilon_{H,i}
 +4M_m\varepsilon_m+2\varepsilon_m^2.}
 \label{eq:YMA21-mixture-error}
\end{equation}
\end{proposition}

\begin{proof}
For each $i$,
\[
 \|m_im_i^*-\widehat m_i\widehat m_i^*\|
 \le\varepsilon_m(2M_m+\varepsilon_m).
\]
The same estimate holds for the outer products of the exact and approximate
mixture means because
$\|\overline m-\widehat{\overline m}\|\le\varepsilon_m$ and
$\|\widehat{\overline m}\|\le M_m$.  Adding the two covariance errors gives
$4M_m\varepsilon_m+2\varepsilon_m^2$.  The phase-Hessian errors add with
weights $\pi_i$.
\end{proof}

\begin{theorem}[Phase-first YM--A terminating alternative]
\label{thm:YMA21-termination}
At a fixed shell, freeze the physical boundary, topological, toron,
stabilizer, reflection, response, field-tangent, and carrier packet before
reading the symbol.  Then proceed in the following order.
\begin{enumerate}[label=\textup{(A21.\arabic*)},leftmargin=5.7em]
\item Verify that the complete response-matched interaction defines a compact
      DLR specification, for example by
      \eqref{eq:YMA21-abs-interaction}.  Failure to load the generated tail is
      \stopstatus{}.
\item Either prove uniqueness of the relevant local state or select one
      boundary/sector phase branch.  If the declared target is a mixture,
      acquire its phase-mean Gram $\mathbb B_\pi$.
\item Establish a differentiably selected phase chart.  A proved pressure
      cusp, divergent phase susceptibility, or nonzero mixture Gram against a
      demanded summable local kernel is a typed phase/local-vacuum
      \obstruction; an unresolved chart is \stopstatus{}.
\item Acquire the response cofactors in that phase, load the final law by
      \eqref{eq:YMA21-phase-b}, and populate the fused finite head and
      $\Wtwo/\Wthree$ tail of the same phase.
\item Evaluate one predeclared carrier-supported continuous Bloch block.
      A positive lower interval is \pass.  A negative upper interval on an
      explicit retained polarization is a pure-mode \obstruction.  Otherwise
      the result is \stopstatus{}.
\item If a mixed state rather than a phase is the physical target, subtract
      the phase Gram according to \eqref{eq:YMA21-Hessian-mixture}; retain the
      zero-momentum term separately according to
      \eqref{eq:YMA21-phase-zero-mode}.
\end{enumerate}
There is no branch in which a coefficient fitted in an unselected mixture is
silently inserted into a constituent local vacuum.
\end{theorem}

\begin{proof}
DLR nonemptiness is \Cref{thm:YMA21-DLR-existence}.  The need to select a
branch before response matching is
\Cref{thm:YMA21-response-decomposition,cth:YMA21-response-noncommutation}.
The phase-mixture debit and its zero-mode typing are
\Cref{thm:YMA21-Hessian-mixture,cor:YMA21-kernel-mixture}.  Differentiability
and the physical branch Hessian are
\Cref{thm:YMA21-selected-Kubo,cor:YMA21-phase-Hessian,thm:YMA21-phase-marginal}.
After these upstream rows pass, the head--tail, continuous Bloch, and direct
Hermitian interval alternatives are the theorems of Version~V20.
\end{proof}

\begin{assessment}[Qualification of Version V20]
\label{ass:YMA21-V20-qualification}
Version~V20 correctly put local-state selection before a global periodic
realization.  Its phrase ``freeze the response-matched action, then select a
local vacuum'' is valid on a unique-phase branch or after phase-independence
of the response chart has been proved.  At coexistence the order must be
reversed: select the physical pre-retuning phase, acquire its response
cofactors, and continue that same phase under the final action.  All V20
head--tail, Bloch, twist, sampling, and diagonal-extraction theorems remain
unchanged once applied to that phase-specific fused kernel.
\end{assessment}

\begin{firewall}[No Wilson phase or margin is manufactured]
\label{fire:YMA21-no-Wilson-value}
The DLR existence theorem is conditional on an absolutely summable submitted
interaction and does not prove uniqueness or clustering.  The mixture and
orbit theorems are exact identities and alternatives, not evidence that the
physical Wilson vacuum has multiple phases or broken residual colour.  No
response cofactor, selected-phase head, nonlinear connected tail, continuous
Bloch eigenvalue, or local-vacuum contraction is evaluated in this version.
\end{firewall}

% =============================================================================
\section{Version V21 proof ledger and next literal acquisition}
\label{sec:YMA21-ledger}
% =============================================================================

\begingroup\small
\begin{longtable}{>{\raggedright\arraybackslash}p{0.28\textwidth}
                  >{\raggedright\arraybackslash}p{0.19\textwidth}
                  >{\raggedright\arraybackslash}p{0.45\textwidth}}
\caption{Version V21 proof-status ledger.}\label{tab:YMA21-ledger}\\
\toprule
\textbf{Row}&\textbf{Status}&\textbf{Exact scope}\\
\midrule
\endfirsthead
\toprule
\textbf{Row}&\textbf{Status}&\textbf{Exact scope}\\
\midrule
\endhead
Compact fixed-shell DLR specification&Proved conditionally
&Every continuous uniformly absolutely summable reduced-cell interaction has
 a nonempty compact convex DLR set and a translation-invariant extreme
 branch.\\[0.3em]
Actual generated interaction summability&Open physical loading
&The pure Wilson head is finite range; the complete submitted generated tail
 has not been populated in the required interaction norm.\\[0.3em]
Response matching versus phase selection&Separated exactly
&Total covariance gives the two between-phase response terms in
 \eqref{eq:YMA21-response-total-covariance}; a finite counterexample proves
 noncommutation.\\[0.3em]
Phase-mixture Hessian&Proved exactly
&The mixed signed card is the phase average minus the positive score-mean
 Gram $\mathbb B_\pi$.\\[0.3em]
Mixture tail and zero mode&Proved exactly
&A nonzero phase Gram is constant in displacement, fails every summable
 moment stock, and lands only at the periodic zero character.\\[0.3em]
Residual-colour orbit average&Proved exactly
&An adjoint mean $m$ produces the isotropic debit $\|m\|^2I_8/8$ after global
 $\SU(3)$ restoration.\\[0.3em]
Coexistence pressure cusp&Explicit countertheorem
&Two affine phase branches have zero individual curvature while the mixed
 curvature is $-V^2a^2$ and the limiting action density is nondifferentiable.\\[0.3em]
Differentiably selected phase&Exact conditional theorem
&Uniform local convergence, susceptibility summability, and boundary
 invisibility pass the finite-volume Kubo identity to one DLR branch.\\[0.3em]
Phase-specific final Hessian&Proved conditionally
&After phase-specific response cofactors are acquired, the direct curvature
 minus covariance is the derivative of the same continued phase.\\[0.3em]
Finite phase-Gram acquisition&Proved exactly
&One positive two-copy writer yields the complete between-phase Gram.\\[0.3em]
Robust finite-mixture interval&Proved
&Equation \eqref{eq:YMA21-mixture-error} gives a direct operator-norm error
 radius.\\[0.3em]
Selected Wilson phase chart&\stopstatus
&No supplied packet proves uniqueness or differentiable boundary-selected
 continuation for the response-matched shell.\\[0.3em]
Phase response cofactors and fused kernel&Open physical values
&No phase-specific cofactor, deterministic/covariance head, or nonlinear
 $\Wtwo/\Wthree$ tail has been populated.\\[0.3em]
Continuous retained Bloch margin&\stopstatus
&The phase-first logical reduction is complete, but the Wilson-specific value
 remains absent.\\[0.3em]
Local-vacuum contraction and mass&Not inferred
&Neither DLR existence nor a phase-mixture identity supplies higher jets,
 temporal transfer, a reflection-positive continuum, or a Yang--Mills mass
 gap.\\
\bottomrule
\end{longtable}
\endgroup

\begin{theorem}[Version V21 phase-first correction]
\label{thm:YMA21-terminal}
Versions~V1--V20 are retained at their stated exact, finite, conditional, and
open scopes.  Version~V21 additionally proves:
\begin{enumerate}[label=\textup{(E21.\arabic*)},leftmargin=5.8em]
\item existence and compactness of fixed-shell DLR states under an explicit
      summable-interaction hypothesis;
\item exact noncommutation of phase averaging and response matching;
\item the positive between-phase debit in the complete signed Hessian;
\item its constant-displacement and periodic zero-mode realization;
\item the exact isotropic $\SU(3)$ orbit debit $\|m\|^2I_8/8$;
\item a finite coexistence model where mixed curvature diverges although every
      phase curvature vanishes;
\item a sufficient differentiable-phase packet which passes the Kubo response
      and the score--Hessian identity to the selected local vacuum;
\item a phase-specific response-matched final-law card and a finite robust
      mixture enclosure.
\end{enumerate}
No item selects the actual Wilson phase or evaluates its symbol.
\end{theorem}

\begin{proof}
Items~\textup{(E21.1)--(E21.2)} are
\Cref{thm:YMA21-DLR-existence,thm:YMA21-response-decomposition,cth:YMA21-response-noncommutation}.
Items~\textup{(E21.3)--(E21.5)} are
\Cref{thm:YMA21-Hessian-mixture,cor:YMA21-kernel-mixture,thm:YMA21-SU3-orbit}.
Item~\textup{(E21.6)} is
\Cref{cth:YMA21-coexistence-cusp}.  Items~\textup{(E21.7)--(E21.8)} are
\Cref{thm:YMA21-selected-Kubo,cor:YMA21-phase-Hessian,thm:YMA21-phase-marginal,prop:YMA21-mixture-error}.
\end{proof}

\begin{openproblem}[V22: select one deformation-stable phase and read its first Wilson block]
\label{op:YMA21-V22}
Preserve every frozen action, response, gauge, topology, colour, tangent,
carrier, reference, shell, and infrared-regulator convention.  Do not return
to a mixed periodic cofactor or another Fourier-volume compiler before the
following phase card is attempted.
\begin{enumerate}[label=\textup{(V22.\arabic*)},leftmargin=5.5em]
\item Exhibit the complete response-matched interaction as a finite-range plus
      absolutely summable potential, or return the first unsupported generated
      tail coefficient.
\item Fix the physical boundary/topological/toron/stabilizer branch.  Prove
      uniqueness on the relevant local algebra or retain each DLR cluster
      phase separately.
\item On every retained phase, acquire the one-point complete score mean and
      the phase-mixture Gram.  Route any nonzero constant term to the harmonic
      phase-survivor row before a Ward quotient is formed.
\item Prove the differentiable phase chart for the plaquette and cutoff
      insertions, including the source-specific susceptibility tail.
\item Acquire the phase-specific response cofactors, load the final law by
      \eqref{eq:YMA21-phase-b}, and verify that the same phase branch persists.
\item Populate one finite fused head and the $\Wtwo/\Wthree$ connected tail in
      that final phase.
\item Evaluate the first guarded continuous Bloch momentum--colour block and
      return \pass, an attained pure-mode \obstruction, or \stopstatus{}.
\end{enumerate}
If a uniqueness or cluster-expansion theorem is unavailable, the first useful
finite calculation is not a phase average.  It is the same boundary-conditioned
head, score mean, and susceptibility tail on two increasing infrared boxes,
with a held-out boundary comparison.
\end{openproblem}

\begin{assessment}[V21 termination]
\label{ass:YMA21-termination}
Version~V20 had reached the correct local-vacuum level but still allowed the
response-matched coefficient to be frozen in an unselected finite-volume law.
The exact total-covariance identities show that this can mix distinct
response charts and subtract an additional phase Gram.  The first unclosed
physical arrow is therefore not another Bloch or carrier theorem.  It is a
\emph{deformation-stable selected phase} for the pre-retuning law, followed by
that phase's response cofactors and fused connected kernel.  Since none of
those values is supplied, the selected Wilson margin remains \stopstatus{}
and no negative physical mode has been produced.
\end{assessment}

The accompanying replay checks only the finite algebra in this continuation:
total covariance, response-matching noncommutation, the mixture Hessian debit,
the positive two-copy phase Gram, the phase-coexistence curvature, an exact
finite specification consistency example, and the robust mixture error
bound.  It does not prove a Wilson cluster expansion, construct the physical
DLR phase, or evaluate a Yang--Mills symbol.

The standalone exact verifier and its result are
\begin{center}
\path{YMA_V21_phase_selection_replay.py},\qquad
\path{YMA_V21_phase_selection_replay.json}.
\end{center}
Their SHA--256 digests, followed by the replay-package digest, are
\begin{center}\scriptsize\ttfamily
3ea3343607f765245d48a88066e054c88521db02dd5cedaf44f36277ab04be2b,\\
6c5a75e450290f286bd56cc1ab625abddca0c6ca4c175d361872c367a0acf42c,\\
8f64015fc85d32708ee33ed515836c52954f2451052798caf171072c4cf5040d.
\end{center}
All eight declared replay checks pass.  The finite signed-permutation twirl
in the replay verifies the irreducible-orbit second-moment algebra; the
$\SU(3)$ adjoint coefficient $1/8$ is proved in the text from simplicity and
Haar invariance rather than inferred from that finite analogue.

\section*{Version V21 append-only continuation record}
\addcontentsline{toc}{section}{Version V21 append-only continuation record}
The complete Version~V20 source above is preserved byte for byte up to its
sole terminal document-closing command.  Its SHA--256 digest is
\begin{center}\scriptsize\ttfamily
c094d724b7f95b19207aeed679e683b4b2c4f518c7f40aa865776b789cd7fb07.
\end{center}
For Version~V22, preserve this accumulated source, append new mathematical
material immediately before the final document-closing command, and increment
the filename to \texttt{\_V22.tex}.  Keep the shell and infrared regulators
distinct.  In addition, freeze the selected pre-retuning DLR phase before its
response cofactors are formed, continue that same phase under the final
response-matched action, and remove every between-phase score Gram before a
summable local Bloch kernel is claimed.

% =============================================================================
% END APPENDED VERSION V21 MATERIAL
% =============================================================================



% =============================================================================
% START APPENDED VERSION V22 MATERIAL
% =============================================================================

\clearpage
\hypersetup{
  pdftitle={Renewal Yang--Mills Reduced-Fibre Wilson Shell Symbol and Local-Vacuum Programme V22},
  pdfsubject={Fixed-action phase decomposition, specification-first quasilocality, and an obstruction-first nonzero Wilson mode card}
}

\section{Version V22: the fixed action precedes phase decomposition, and a
nonzero mode may be tested before phase uniqueness}
\label{sec:YMA22-direction}
% =============================================================================

Version~V21 correctly proved that phase averaging can add a positive
between-phase score Gram to a partition curvature and that a coefficient
fitted in a mixture need not be response null in any constituent phase.  Its
next-step order nevertheless imposed a stronger semantics than the native
YM--A shell.  The inherited shell prescription was already fixed in
Version~V16: the two response cofactors determine one finite coefficient
\[
 b_r=-\frac{\mathfrak c_{\Psi,r}}{\mathfrak c_{P,r}},
\]
and hence one finite response-matched action
\[
 F_{r,b_r}=F_{r,0}+b_rF_{r,P}
\]
before an infrared DLR decomposition is taken.  Recomputing a different
coefficient in each limiting phase generally produces a different family of
actions.  It is an optional phase-retuned scheme, not an automatic
interpretation of the frozen shell.

The upstream order for the native programme is therefore
\begin{equation}
 \boxed{
 \begin{gathered}
 \text{finite-shell response cofactors and the fixed action }F_{r,b_r}
 \\
 \longrightarrow
 \text{the exact quasilocal specification of that action}
 \\
 \longrightarrow
 \text{its translation-invariant DLR phase set}
 \\
 \longrightarrow
 \text{a nonzero-character fixed-action symbol card}
 \\
 \longrightarrow
 \begin{matrix}
 \text{an extreme-phase obstruction, or a phase-uniform pass}\\[-0.15em]
 \text{after a selected phase, boundary-pair floor,}\\[-0.15em]
 \text{or phase-dispersion certificate}
 \end{matrix}.
 \end{gathered}}
 \label{eq:YMA22-upstream-order}
\end{equation}

This correction does not discard any total-covariance identity from
Version~V21.  It changes what those identities are used for.  Phasewise
response residuals become diagnostics of the fixed action or of an optional
phase-retuned experiment; they do not silently redefine the Wilsonian shell.

% =============================================================================
\section{The native fixed-action phase packet}
\label{sec:YMA22-fixed-action}
% =============================================================================

\begin{definition}[Fixed-action and phase-retuned semantics]
\label{def:YMA22-two-semantics}
Fix a shell $r$ and a predeclared finite response card.  Let
\begin{equation}
 c_r=\mathfrak c_{u,r}[P_r],
 \qquad
 e_r=\mathfrak c_{u,r}[\Psi_r^{\rm cut}],
 \qquad
 b_r=-e_r/c_r,
 \label{eq:YMA22-fixed-b}
\end{equation}
with $c_r\ne0$, and let
\begin{equation}
 F_{r,b_r}=F_{r,0}+b_rF_{r,P}
 \label{eq:YMA22-fixed-action}
\end{equation}
be the complete loaded action, including every generated term, support birth,
Jacobian, and declared local/topological convention.

The \emph{native fixed-action packet} consists of the DLR states of the one
specification generated by \eqref{eq:YMA22-fixed-action}.  The number $b_r$ is
common to every such state.

A \emph{phase-retuned packet} starts instead from pre-retuning phases
$\nu_{r,\theta}$ and assigns, when defined,
\begin{equation}
 b_{r,\theta}
 =-\frac{n_{r,\theta}}{d_{r,\theta}},
 \qquad
 d_{r,\theta}=-\Cov_{\nu_{r,\theta}}(S_{\rm Cr},P_r),
 \qquad
 n_{r,\theta}=-\Cov_{\nu_{r,\theta}}
                  (S_{\rm Cr},\Psi_r^{\rm cut}).
 \label{eq:YMA22-phase-retuned-b}
\end{equation}
This second construction is a family of phase-dependent actions and is a
strictly stronger physical prescription unless
$b_{r,\theta}=b_r$ on every retained phase.
\end{definition}

\begin{theorem}[Native action loading precedes its DLR decomposition]
\label{thm:YMA22-fixed-action-first}
For the packet of \Cref{def:YMA22-two-semantics}, the finite response-null
identity
\begin{equation}
 \mathscr R_r\bigl(\Psi_r^{\rm cut}+b_rP_r\bigr)=0
 \label{eq:YMA22-finite-response-null}
\end{equation}
loads one action \eqref{eq:YMA22-fixed-action}.  Every DLR state and every
barycentric phase of that action is evaluated with this same $b_r$.  No
phasewise response identity is required for the existence or definition of
the resulting field Hessian.

If one additionally demands that the same response target be stationary in
every pre-retuning phase, that demand is the separate family of equations
\begin{equation}
 r_{r,\theta}
 :=n_{r,\theta}+b_rd_{r,\theta}=0.
 \label{eq:YMA22-phase-response-residual}
\end{equation}
When $d_{r,\theta}\ne0$, the optional phase-retuned coefficient satisfies the
exact identity
\begin{equation}
 \boxed{
 b_{r,\theta}-b_r
 =-\frac{r_{r,\theta}}{d_{r,\theta}}.}
 \label{eq:YMA22-retuning-difference}
\end{equation}
Thus a nonzero $r_{r,\theta}$ is a response-scheme discrepancy for the fixed
action; it is not a license to replace $b_r$ inside its Wilson symbol.
\end{theorem}

\begin{proof}
The first assertion is precisely the frozen tangent semantics of
\eqref{eq:YMA16-final-action}: the cofactors are acquired on the declared
finite response law and their quotient is part of the action loading.  A DLR
decomposition is a decomposition of a state of that already fixed
specification and cannot change its coefficients.

The phase response of the fixed tangent is, by linearity,
$n_{r,\theta}+b_rd_{r,\theta}$.  If $d_{r,\theta}\ne0$, subtract
$b_r$ from \eqref{eq:YMA22-phase-retuned-b} to obtain
\[
 b_{r,\theta}-b_r
 =-\frac{n_{r,\theta}+b_rd_{r,\theta}}
         {d_{r,\theta}},
\]
which is \eqref{eq:YMA22-retuning-difference}.
\end{proof}

\begin{corollary}[Quantitative cost of optional phase retuning]
\label{cor:YMA22-retuning-cost}
Suppose the selected phase branch admits a differentiable margin matrix
$M_{r,\theta}(b)$ on the interval joining $b_r$ and $b_{r,\theta}$ and
\begin{equation}
 \sup_b\|\partial_bM_{r,\theta}(b)\|_{\rm op}
 \le L_{r,\theta}.
 \label{eq:YMA22-b-Lipschitz}
\end{equation}
Then
\begin{equation}
 \boxed{
 \|M_{r,\theta}(b_{r,\theta})
      -M_{r,\theta}(b_r)\|_{\rm op}
 \le
 L_{r,\theta}
 \frac{|r_{r,\theta}|}{|d_{r,\theta}|}.}
 \label{eq:YMA22-retuning-margin-cost}
\end{equation}
In particular, closeness of the fixed-action and phase-retuned cards requires
both a response residual bound and a phasewise denominator floor.
\end{corollary}

\begin{proof}
Combine the fundamental theorem of calculus in $b$ with
\eqref{eq:YMA22-retuning-difference}.
\end{proof}

\begin{assessment}[Qualification of the phase-specific coefficient in V21]
\label{ass:YMA22-V21-coefficient-qualification}
The total-covariance theorem, phase Gram, coexistence cusp, and selected-phase
Kubo theorem of Version~V21 remain unchanged.  Its coefficient
$b_\Ph=-n_\Psi^\Ph/d_P^\Ph$ is the correct coefficient for a separately
specified phase-retuned response experiment.  For the native YM--A shell,
however, the coefficient is the finite-shell $b_r$ of
\eqref{eq:YMA22-fixed-b}; phase selection is subsequently performed inside
the DLR set of \eqref{eq:YMA22-fixed-action}.  If the programme also demands
phasewise response stationarity, the residual
\eqref{eq:YMA22-phase-response-residual} is retained as an additional
physical row rather than absorbed by changing the action.
\end{assessment}

% =============================================================================
\section{A specification-first replacement for a chosen generated potential}
\label{sec:YMA22-specification}
% =============================================================================

The first V21 loading row asked for one finite-range plus absolutely summable
interaction representation.  Such a representation is sufficient, but the
physical object needed for phase existence is the conditional specification
itself.  The exact reduced-fibre shell may therefore be certified directly,
without first choosing a preferred decomposition of its generated tail into
interaction sets.

\begin{definition}[Exterior oscillation of an exact conditional cylinder]
\label{def:YMA22-exterior-oscillation}
Let $\Omega=\mathsf X^{\mathbb Z^4}$ with $\mathsf X$ compact.  For every
finite $\Lambda\Subset\mathbb Z^4$, suppose an exact conditional kernel has
strictly positive density
\begin{equation}
 \gamma_\Lambda(\dd z_\Lambda\mid\eta)
 =\frac{e^{-\mathcal H_\Lambda(z_\Lambda\mid\eta)}}
        {Z_\Lambda(\eta)}
   m^{\otimes\Lambda}(\dd z_\Lambda).
 \label{eq:YMA22-direct-kernel}
\end{equation}
For $R\ge0$, write $B_R(\Lambda)$ for the $\ell^1$ neighbourhood of
$\Lambda$ and define
\begin{equation}
 \begin{aligned}
 \varepsilon_\Lambda(R)
 :=\sup_{\substack{\eta,\eta':\;
        \eta=\eta'\text{ on }B_R(\Lambda)\setminus\Lambda}}
 \inf_{c\in\R}
 \sup_{z_\Lambda}
 \bigl|
 &\mathcal H_\Lambda(z_\Lambda\mid\eta)
 -\mathcal H_\Lambda(z_\Lambda\mid\eta')-c
 \bigr|.
 \end{aligned}
 \label{eq:YMA22-exterior-modulus}
\end{equation}
The scalar $c$ removes boundary-dependent normalizing terms which are
constant in the conditioned variables.
\end{definition}

\begin{lemma}[Log-density oscillation controls conditional total variation]
\label{lem:YMA22-log-density-TV}
Define
\[
 \|\rho-\nu\|_{\rm TV,*}
 :=\sup_{\|f\|_\infty\le1}|\rho(f)-\nu(f)|.
\]
If two kernels of the form \eqref{eq:YMA22-direct-kernel} have conditional
Hamiltonians $H,H'$ satisfying
\begin{equation}
 \inf_{c\in\R}\|H-H'-c\|_\infty\le\varepsilon,
 \label{eq:YMA22-log-osc-assumption}
\end{equation}
then
\begin{equation}
 \boxed{
 \|\gamma_H-\gamma_{H'}\|_{\rm TV,*}
 \le \min\{2,e^{2\varepsilon}-1\}.}
 \label{eq:YMA22-TV-bound}
\end{equation}
Consequently, $\varepsilon_\Lambda(R)\to0$ for every finite $\Lambda$
implies quasilocality of the exact specification.
\end{lemma}

\begin{proof}
Choose $c$ up to an arbitrarily small enlargement of $\varepsilon$ and put
$u=H'-H-c$.  With $\mu=\gamma_H$ and $\nu=\gamma_{H'}$,
\[
 \frac{\dd\nu}{\dd\mu}
 =\frac{e^{-u}}{\mathbb E_\mu e^{-u}}.
\]
Since $|u|\le\varepsilon$, both $e^{-u}$ and its $\mu$-mean lie in
$[e^{-\varepsilon},e^{\varepsilon}]$.  Hence the density ratio lies in
$[e^{-2\varepsilon},e^{2\varepsilon}]$, and
\[
 \|\nu-\mu\|_{\rm TV,*}
 =\mathbb E_\mu\left|
   \frac{\dd\nu}{\dd\mu}-1\right|
 \le e^{2\varepsilon}-1.
\]
The norm of a difference of probability measures is at most two.  Applying
this estimate to boundary conditions agreeing on growing neighbourhoods
proves quasilocality.
\end{proof}

\begin{theorem}[Specification-first fixed-shell phase existence]
\label{thm:YMA22-specification-existence}
Suppose the kernels $(\gamma_\Lambda)_\Lambda$ of
\eqref{eq:YMA22-direct-kernel} are proper and consistent and satisfy
\begin{equation}
 \varepsilon_\Lambda(R)\longrightarrow0
 \qquad(R\to\infty)
 \label{eq:YMA22-direct-quasilocality}
\end{equation}
for every finite $\Lambda$.  Then the DLR set of the specification is
nonempty, compact, and convex.  If the specification is translation
covariant, it has a nonempty compact convex translation-invariant DLR subset
and at least one translation-invariant extreme point.

Thus the first V21 loading row may be passed in either of two ways:
\begin{enumerate}[label=\textup{(S\arabic*)},leftmargin=4.0em]
\item exhibit an absolutely summable potential as in
      \eqref{eq:YMA21-abs-interaction}; or
\item verify consistency of the exact reduced-fibre conditional cylinders
      and the direct modulus \eqref{eq:YMA22-direct-quasilocality}.
\end{enumerate}
The second route does not select a phase or prove a Wilson symbol margin.
\end{theorem}

\begin{proof}
By \Cref{lem:YMA22-log-density-TV}, every $\gamma_\Lambda f$ is continuous
for continuous local $f$.  Choose increasing boxes and any exterior
condition.  Compactness of $\Omega$ gives a weak cluster point of the
finite-volume conditional laws.  Properness and consistency imply
$\mu_n\gamma_\Delta=\mu_n$ once $\Delta$ lies inside the volume.
Quasilocality permits passage to the cluster point, giving the DLR equations.
They are affine and weakly closed, so the DLR set is compact and convex.
Translation averaging over Følner boxes preserves the DLR equations and
produces a translation-invariant cluster point.  An extreme point exists by
Krein--Milman.  This is the V21 compactness proof with the potential
representation replaced by the direct quasilocality hypothesis.
\end{proof}

\begin{corollary}[The V21 potential norm implies the direct modulus]
\label{cor:YMA22-potential-implies-modulus}
For the interaction of \Cref{def:YMA21-interaction},
\begin{equation}
 \boxed{
 \varepsilon_\Lambda(R)
 \le
 2\sum_{\substack{A:\,A\cap\Lambda\ne\varnothing,\\
                   A\not\subseteq B_R(\Lambda)}}
   \|\Phi_A\|_\infty.}
 \label{eq:YMA22-potential-tail-modulus}
\end{equation}
For every fixed $\Lambda$, the right side tends to zero under uniform
absolute summability.
\end{corollary}

\begin{proof}
Terms contained in $B_R(\Lambda)$ are identical under the two exterior
conditions.  Every remaining term changes by at most twice its supremum norm.
Terms independent of $z_\Lambda$ may be absorbed in the scalar $c$ of
\eqref{eq:YMA22-exterior-modulus}.  Summing the remaining bounds proves the
claim.  The tail tends to zero because $\Lambda$ is finite and the
translation-covariant interaction norm is summable.
\end{proof}

\begin{firewall}[Specification existence is not an interaction-tail estimate]
\label{fire:YMA22-specification-scope}
The direct modulus avoids a representation-dependent potential ledger, but
it is still a physical estimate on the exact response-matched shell.  No
value of $\varepsilon_\Lambda(R)$ is supplied here for the nonlinear Wilson
generated tail.  Exact finite-volume consistency alone does not imply
quasilocality, clustering, a differentiable phase, or a summable connected
Hessian kernel.
\end{firewall}

% =============================================================================
\section{Zero-total field packets remove the between-phase score mean exactly}
\label{sec:YMA22-zero-total}
% =============================================================================

The phase Gram of Version~V21 is a genuine zero-character obstruction.  For
nonzero momentum there is a stronger exact statement which requires no phase
selection and no spatial decay estimate.

\begin{theorem}[Fixed-action zero-total phase barycentre]
\label{thm:YMA22-zero-total-barycentre}
Let
\begin{equation}
 \mu=\int_\Theta\mu_\theta\,\pi(\dd\theta)
 \label{eq:YMA22-fixed-action-decomposition}
\end{equation}
be a barycentric decomposition of translation-invariant states of one fixed
action.  Let $J_x\in L^2(\mu;\C^d)$ be a translation-covariant complete field
score with
\[
 \mathbb E_\theta J_x=m_\theta
\]
independent of $x$.  For a finitely supported coefficient field
$a=(a_x)\subset\C^d$, put
\[
 J[a]=\sum_x a_x^*J_x
\]
and let $K[a]$ be the complete deterministic second-action writer in the
same field direction.  Define
\[
 H_\nu[a]=\mathbb E_\nu K[a]-\Var_\nu(J[a]).
\]
For fixed deterministic topological and reference forms, put
\begin{equation}
 M_\nu[a;\delta]
 :=H_\nu[a]-T[a]-\delta S[a].
 \label{eq:YMA22-finite-probe-margin}
\end{equation}
If
\begin{equation}
 \sum_xa_x=0,
 \label{eq:YMA22-zero-total-condition}
\end{equation}
then
\begin{equation}
 \boxed{
 H_\mu[a]
 =\int_\Theta H_{\mu_\theta}[a]\,\pi(\dd\theta).}
 \label{eq:YMA22-zero-total-Hessian}
\end{equation}
The same identity holds after subtracting any fixed deterministic reference,
topological, or carrier form.
\end{theorem}

\begin{proof}
Translation invariance gives
\[
 \mathbb E_\theta J[a]
 =\left(\sum_xa_x\right)^*m_\theta=0.
\]
The law of total variance therefore has no between-phase term:
\[
 \Var_\mu(J[a])
 =\int\Var_{\mu_\theta}(J[a])\,\pi(\dd\theta)
  +\Var_\pi(\mathbb E_\theta J[a])
 =\int\Var_{\mu_\theta}(J[a])\,\pi(\dd\theta).
\]
The deterministic expectation is affine in the state.  Subtraction gives
\eqref{eq:YMA22-zero-total-Hessian}.  A fixed quadratic form is unchanged by
barycentric integration.
\end{proof}

\begin{corollary}[Every nonzero periodic character is phase-mean free]
\label{cor:YMA22-nonzero-character}
On a periodic box $\Lambda_N$, let $q$ be a nonzero character and let
\begin{equation}
 a_x=|\Lambda_N|^{-1/2}e^{\mathrm iq\cdot x}v,
 \qquad v\in\C^d.
 \label{eq:YMA22-plane-wave-packet}
\end{equation}
Then $\sum_xa_x=0$ and
\begin{equation}
 \boxed{
 v^*H_{\mu,N}(q)v
 =\int v^*H_{\theta,N}(q)v\,\pi(\dd\theta).}
 \label{eq:YMA22-periodic-symbol-average}
\end{equation}
Polarization gives the corresponding matrix identity.
\end{corollary}

\begin{proof}
Character orthogonality gives
$\sum_{x\in\Lambda_N}e^{\mathrm iq\cdot x}=0$.  Apply
\Cref{thm:YMA22-zero-total-barycentre} and polarize in $v$.
\end{proof}

\begin{corollary}[Regular continuous Bloch symbol away from the phase atom]
\label{cor:YMA22-continuous-symbol-average}
Suppose the phase-specific assembled kernels $A_\theta(x)$ satisfy
\begin{equation}
 \int_\Theta
 \sum_{x\in\mathbb Z^4}
 \|A_\theta(x)\|_{\rm op}
 \,\pi(\dd\theta)<\infty.
 \label{eq:YMA22-averaged-kernel-summability}
\end{equation}
Then the regular nonzero-momentum symbol is
\begin{equation}
 \boxed{
 \overline H(q)
 =\int_\Theta H_\theta(q)\,\pi(\dd\theta),
 \qquad q\ne0.}
 \label{eq:YMA22-continuous-symbol-barycentre}
\end{equation}
The only additional distributional term in the uncentered mixture covariance
is the phase Gram of Version~V21 at $q=0$.
\end{corollary}

\begin{proof}
Condition \eqref{eq:YMA22-averaged-kernel-summability} permits Fubini's
theorem in the Fourier series.  The connected phase-average kernel therefore
has symbol equal to the displayed integral.  The total-covariance term is
constant in displacement and hence is supported at the zero character, as
proved in \eqref{eq:YMA21-phase-zero-mode}.
\end{proof}

\begin{assessment}[What phase selection is and is not needed for]
\label{ass:YMA22-phase-needed}
For the fixed action, phase selection is not needed to remove the
between-phase score mean from a predeclared nonzero-character test: the
zero-total identity does so exactly.  Phase selection remains necessary to
claim positivity for one particular constituent phase from a merely positive
phase average, to define a branch-specific response derivative, and to
control the zero-momentum or nonsummable phase survivor.
\end{assessment}

% =============================================================================
\section{Obstruction-first nonzero-mode logic}
\label{sec:YMA22-obstruction-first}
% =============================================================================

Let $S(q)\succeq0$ be the already frozen reference on one supported
momentum--colour block and put
\begin{equation}
 M_\theta(q;\delta)
 =H_\theta^{\rm loc}(q)-\delta S(q),
 \qquad
 \overline M(q;\delta)
 =\int M_\theta(q;\delta)\,\pi(\dd\theta).
 \label{eq:YMA22-phase-margin-matrices}
\end{equation}

\begin{theorem}[A negative fixed-action mixture exposes a bad phase]
\label{thm:YMA22-negative-average-phase}
Fix $q\ne0$ and a supported vector $v$.  If
\begin{equation}
 v^*\overline M(q;\delta)v<0,
 \label{eq:YMA22-negative-average}
\end{equation}
then
\begin{equation}
 \pi\left(
  \left\{\theta:
   v^*M_\theta(q;\delta)v<0
  \right\}
 \right)>0.
 \label{eq:YMA22-positive-bad-phase-measure}
\end{equation}
In particular, at least one component state has the same explicit
momentum--colour polarization as a negative pure-mode witness.  An extreme
phase witness is supplied by the envelope theorem below.

Conversely, if
$M_\theta(q;\delta)\succeq0$ for almost every phase, then
$\overline M(q;\delta)\succeq0$.
\end{theorem}

\begin{proof}
By \Cref{thm:YMA22-zero-total-barycentre}, the scalar in
\eqref{eq:YMA22-negative-average} is the integral of
$v^*M_\theta v$.  An integrable function with negative integral is negative
on a set of positive measure.  The converse follows because the positive
semidefinite cone is convex and closed under barycentric integration.
\end{proof}

\begin{theorem}[Attained extreme translation-invariant phase envelope]
\label{thm:YMA22-extreme-envelope}
Let $\mathcal G_{r,b}^{\rm ti}$ be the compact convex set of
translation-invariant DLR states of the fixed action, and let $a$ be one
finite zero-total field packet.  Define
\begin{equation}
 \underline m_{r,b}[a;\delta]
 :=\min_{\nu\in\mathcal G_{r,b}^{\rm ti}}
   \lambda_{\min}\bigl(M_\nu[a;\delta]\bigr),
 \label{eq:YMA22-extreme-envelope}
\end{equation}
where $M_\nu[a;\delta]$ is the complete fixed-action Hessian after the frozen
topological and reference subtractions.  Then
\begin{equation}
 \boxed{
 \underline m_{r,b}[a;\delta]
 =\min_{\theta\in\operatorname{Ext}(\mathcal G_{r,b}^{\rm ti})}
   \lambda_{\min}\bigl(M_\theta[a;\delta]\bigr).}
 \label{eq:YMA22-extreme-envelope-equality}
\end{equation}
The minimum is attained by an extreme point of
$\mathcal G_{r,b}^{\rm ti}$ and a unit polarization.  Such a state is a
translation-ergodic DLR phase.  Consequently a negative
translation-invariant mixed card always has an extreme
translation-invariant phase witness for the same fixed action.
\end{theorem}

\begin{proof}
By \Cref{thm:YMA22-zero-total-barycentre}, the map
$\nu\mapsto M_\nu[a;\delta]$ is affine.  It is weakly continuous because the
packet is finite and its writers are bounded local functions.  The function
\[
 (\nu,v)\longmapsto v^*M_\nu[a;\delta]v
\]
is therefore continuous on the compact product of
$\mathcal G_{r,b}^{\rm ti}$ and the unit sphere.  Choose a minimizing pair
$(\nu_0,v_0)$.  For fixed $v_0$, the set of states minimizing the affine
functional $\nu\mapsto v_0^*M_\nu v_0$ is a nonempty compact face.  By
Krein--Milman it contains an extreme point $\theta_0$, and an extreme point
of a face is extreme in the ambient DLR set.  Hence
\[
 \lambda_{\min}(M_{\theta_0})
 \le v_0^*M_{\theta_0}v_0
 =v_0^*M_{\nu_0}v_0
 =\underline m_{r,b}[a;\delta].
\]
The reverse inequality is immediate because the extreme points are states in
the full set.  This proves equality and attainment.  No assertion about
non-translation-invariant extreme DLR states is needed.
\end{proof}

\begin{countertheorem}[A positive mixture can hide a negative phase]
\label{cth:YMA22-positive-average-not-phasewise}
A positive nonzero-character phase average does not imply phasewise
positivity.  With two equally weighted one-dimensional phase cards
\begin{equation}
 M_+=2,
 \qquad
 M_-=-1,
 \label{eq:YMA22-masking-example}
\end{equation}
one has $\overline M=1/2>0$ although the minus phase fails.
Consequently a positive mixed card is a pass only for the mixed target.  It
is not a pass for an independently selected phase without additional phase
control.
\end{countertheorem}

\begin{proof}
The average is $(2-1)/2=1/2$, while $M_-<0$.
\end{proof}

\begin{corollary}[Obstruction-first use of an unselected infrared limit]
\label{cor:YMA22-obstruction-first-use}
Suppose a boundary-independent or mixed fixed-action infrared limit supplies
an outward interval for $\overline M(q;\delta)$ at $q\ne0$.  A negative upper
interval on an explicit vector is already a typed
\emph{extremal-phase pure-mode obstruction}: at least one DLR phase fails on
that vector.  If the declared target is the mixed state or uniform positivity
of every phase, this is terminal.  If the target is the existence of one
unspecified good phase, it narrows the phase set but does not by itself
exclude another passing phase.
\end{corollary}

% =============================================================================
\section{A positive phase-dispersion card upgrades an average pass}
\label{sec:YMA22-phase-dispersion}
% =============================================================================

The missing information in
\Cref{cth:YMA22-positive-average-not-phasewise} is not another action
coefficient.  It is the spread of the already fixed-action symbol across the
retained phase set.

\begin{theorem}[Finite phase-dispersion identity and uniform pass]
\label{thm:YMA22-phase-dispersion}
Let $M_1,\ldots,M_s$ be Hermitian $d\times d$ margin matrices of one fixed
action, with weights $\pi_i>0$, $\sum_i\pi_i=1$.  Put
\begin{equation}
 \overline M=\sum_i\pi_iM_i
 \label{eq:YMA22-finite-average}
\end{equation}
and define the positive scalar dispersion
\begin{equation}
 \boxed{
 \mathcal D_\pi
 :=\frac12\sum_{i,j}\pi_i\pi_j
       \|M_i-M_j\|_{\mathrm{HS}}^2.}
 \label{eq:YMA22-phase-dispersion}
\end{equation}
Then
\begin{equation}
 \boxed{
 \mathcal D_\pi
 =\sum_i\pi_i\|M_i-\overline M\|_{\mathrm{HS}}^2.}
 \label{eq:YMA22-dispersion-Pythagoras}
\end{equation}
If $\pi_*:=\min_i\pi_i$, then every phase satisfies
\begin{equation}
 \boxed{
 \lambda_{\min}(M_i)
 \ge
 \lambda_{\min}(\overline M)
 -\sqrt{\mathcal D_\pi/\pi_*}.}
 \label{eq:YMA22-phasewise-lower}
\end{equation}
Consequently
\begin{equation}
 \boxed{
 \lambda_{\min}(\overline M)
 >\sqrt{\mathcal D_\pi/\pi_*}}
 \label{eq:YMA22-uniform-phase-pass}
\end{equation}
implies $M_i\succ0$ for every retained phase.
\end{theorem}

\begin{proof}
Expand the squared Hilbert--Schmidt distance in
\eqref{eq:YMA22-phase-dispersion}.  The cross terms give
$\|\overline M\|_{\mathrm{HS}}^2$, yielding
\eqref{eq:YMA22-dispersion-Pythagoras}.  For each $i$,
\[
 \pi_i\|M_i-\overline M\|_{\mathrm{HS}}^2
 \le\mathcal D_\pi.
\]
The operator norm is bounded by the Hilbert--Schmidt norm, and Weyl's
inequality gives
\[
 \lambda_{\min}(M_i)
 \ge\lambda_{\min}(\overline M)
     -\|M_i-\overline M\|_{\rm op}.
\]
Insert the preceding estimate and $\pi_i\ge\pi_*$.
\end{proof}

\begin{corollary}[Positive two-phase acquisition]
\label{cor:YMA22-dispersion-replica}
For independent phase labels $\Theta_1,\Theta_2$ with law $\pi$,
\begin{equation}
 \boxed{
 \mathcal D_\pi
 =\frac12\mathbb E
   \|M_{\Theta_1}-M_{\Theta_2}\|_{\mathrm{HS}}^2.}
 \label{eq:YMA22-dispersion-two-copy}
\end{equation}
The writer is nonnegative sample by sample.  Moreover,
$\mathcal D_\pi=0$ if and only if every positive-weight phase has the same
margin matrix.
\end{corollary}

\begin{proof}
The first identity is the definition written as an expectation.  In
\eqref{eq:YMA22-dispersion-Pythagoras}, a sum of nonnegative terms vanishes
exactly when every $M_i=\overline M$.
\end{proof}

\begin{proposition}[Infinite phase bank: a quantitative exceptional set]
\label{prop:YMA22-infinite-phase-dispersion}
For a measurable Hermitian phase family with
\begin{equation}
 \mathcal D_\pi
 :=\int\|M_\theta-\overline M\|_{\mathrm{HS}}^2
       \,\pi(\dd\theta)<\infty,
 \label{eq:YMA22-infinite-dispersion}
\end{equation}
one has, for every $t>0$,
\begin{equation}
 \boxed{
 \pi\bigl(
  \|M_\theta-\overline M\|_{\rm op}\ge t
 \bigr)
 \le\mathcal D_\pi/t^2.}
 \label{eq:YMA22-phase-Chebyshev}
\end{equation}
If $m:=\lambda_{\min}(\overline M)>0$, then all but phase measure at most
$\mathcal D_\pi/m^2$ have a positive margin.  This does not identify an
independently boundary-selected phase inside the exceptional set.
\end{proposition}

\begin{proof}
The operator norm is bounded by the Hilbert--Schmidt norm, so Markov's
inequality applied to the squared norm gives
\eqref{eq:YMA22-phase-Chebyshev}.  A phase with
$\|M_\theta-\overline M\|_{\rm op}<m$ is positive by Weyl's inequality.
\end{proof}

\begin{firewall}[Dispersion is not phase occurrence]
\label{fire:YMA22-dispersion-not-selection}
A small phase-dispersion card can upgrade a positive average to a uniform
finite phase-bank pass.  For an infinite phase family it controls only the
measure of an exceptional set.  It does not prove that a predeclared boundary
condition selects a phase outside that set.  Boundary-phase incidence remains
necessary when the target names one specific vacuum rather than the mixed
state or all phases.
\end{firewall}

% =============================================================================
\section{A boundary-pair cylinder gives a uniqueness-free all-phase pass}
\label{sec:YMA22-boundary-pair}
% =============================================================================

Let $a$ be one finite zero-total packet whose score $J[a]$, deterministic
writer $K[a]$, and fixed subtraction $L[a]=T[a]+\delta S[a]$ are supported
strictly inside a finite region $\Lambda$.  For an exterior configuration
$\eta$, abbreviate the exact conditional law by
$\nu_\eta=\gamma_\Lambda(\cdot\mid\eta)$ and put
\begin{equation}
 m_\eta:=\mathbb E_{\nu_\eta}J[a].
 \label{eq:YMA22-boundary-mean}
\end{equation}
For two independent interiors with possibly different exteriors define
\begin{equation}
 \boxed{
 \begin{aligned}
 Q_\Lambda^{\eta_1,\eta_2}[a;\delta]
 :={}&\frac12\bigl(
    \mathbb E_{\nu_{\eta_1}}K[a]
   +\mathbb E_{\nu_{\eta_2}}K[a]\bigr)-L[a]\\
 &-\frac12\mathbb E_{\nu_{\eta_1}\otimes\nu_{\eta_2}}
   \left[(J[a]^{(1)}-J[a]^{(2)})
         (J[a]^{(1)}-J[a]^{(2)})^*\right].
 \end{aligned}}
 \label{eq:YMA22-boundary-pair-card}
\end{equation}
This is the same direct-curvature minus same-history covariance card, written
before an infrared phase has been selected.

\begin{theorem}[Exact DLR boundary-pair disintegration]
\label{thm:YMA22-boundary-pair-disintegration}
For every DLR state $\mu$ of the fixed action,
\begin{equation}
 \boxed{
 M_\mu[a;\delta]
 =\iint Q_\Lambda^{\eta_1,\eta_2}[a;\delta]
       \,\mu(\dd\eta_1)\mu(\dd\eta_2).}
 \label{eq:YMA22-DLR-boundary-average}
\end{equation}
Therefore a common lower bound for all exterior pairs is a lower bound for
every DLR state, with no uniqueness assumption.
\end{theorem}

\begin{proof}
Apply the DLR equation to the local writers $K[a]$ and $J[a]$ in each of two
independent copies.  The first line of
\eqref{eq:YMA22-boundary-pair-card} averages to $\mathbb E_\mu K[a]-L[a]$,
while
\[
 \frac12\mathbb E_{\mu\otimes\mu}
 [(J^{(1)}-J^{(2)})(J^{(1)}-J^{(2)})^*]
 =\operatorname{Cov}_\mu(J[a]).
\]
This is exactly $M_\mu[a;\delta]$.
\end{proof}

\begin{theorem}[Same-boundary floor minus boundary-mean spread]
\label{thm:YMA22-boundary-spread}
For every pair of exteriors,
\begin{equation}
 \boxed{
 \begin{aligned}
 Q_\Lambda^{\eta_1,\eta_2}
 ={}&\frac12Q_\Lambda^{\eta_1,\eta_1}
    +\frac12Q_\Lambda^{\eta_2,\eta_2}\\
 &-\frac12(m_{\eta_1}-m_{\eta_2})
              (m_{\eta_1}-m_{\eta_2})^*.
 \end{aligned}}
 \label{eq:YMA22-boundary-pair-split}
\end{equation}
Put
\begin{equation}
 \ell_{\Lambda,a}
 :=\inf_\eta\lambda_{\min}
   \bigl(Q_\Lambda^{\eta,\eta}[a;\delta]\bigr),
 \qquad
 \rho_{\Lambda,a}
 :=\frac12\sup_{\eta_1,\eta_2}
   \|m_{\eta_1}-m_{\eta_2}\|^2.
 \label{eq:YMA22-boundary-two-stocks}
\end{equation}
Then every DLR state obeys
\begin{equation}
 \boxed{
 M_\mu[a;\delta]
 \succeq(\ell_{\Lambda,a}-\rho_{\Lambda,a})I.}
 \label{eq:YMA22-boundary-uniform-floor}
\end{equation}
In particular,
\begin{equation}
 \boxed{\ell_{\Lambda,a}>\rho_{\Lambda,a}}
 \label{eq:YMA22-boundary-all-phase-pass}
\end{equation}
is a finite all-phase $\pass$ certificate.
\end{theorem}

\begin{proof}
For arbitrary laws $\nu_1,\nu_2$ with means $m_1,m_2$,
\[
 \frac12\mathbb E_{\nu_1\otimes\nu_2}
 [(J_1-J_2)(J_1-J_2)^*]
 =\frac12\Cov_{\nu_1}(J)+\frac12\Cov_{\nu_2}(J)
  +\frac12(m_1-m_2)(m_1-m_2)^*.
\]
Substitution proves \eqref{eq:YMA22-boundary-pair-split}.  The final rank-one
operator is bounded above by $\rho_{\Lambda,a}I$.  Use
\Cref{thm:YMA22-boundary-pair-disintegration}.
\end{proof}

\begin{theorem}[Negative same-boundary exhaustion extracts a DLR obstruction]
\label{thm:YMA22-negative-boundary-extraction}
Let $\Lambda_n\uparrow\mathbb Z^4$ and choose exteriors $\eta_n$.  Suppose a
fixed finite zero-total packet and unit vectors $v_n$ satisfy
\begin{equation}
 v_n^*Q_{\Lambda_n}^{\eta_n,\eta_n}[a;\delta]v_n
 \le-\epsilon
 \label{eq:YMA22-negative-boundary-sequence}
\end{equation}
for one $\epsilon>0$.  Then a subsequence of the same-boundary Gibbs states
has a DLR cluster state $\mu_\infty$ and a unit limit vector $v$ with
\begin{equation}
 \boxed{v^*M_{\mu_\infty}[a;\delta]v\le-\epsilon.}
 \label{eq:YMA22-negative-DLR-limit}
\end{equation}
If the cluster state is translation invariant, the extreme-phase envelope
contains an attained negative phase.  The same conclusion holds for a Bloch
packet whose finite heads converge and whose phase-uniform connected tail
tends to zero as in Version~V20.
\end{theorem}

\begin{proof}
Compactness of the configuration space gives a weakly convergent subsequence.
For every fixed finite $\Delta$ and all sufficiently large $n$, consistency
of the specification gives $\nu_n\gamma_\Delta=\nu_n$.  Quasilocality permits
passage to the limit, so the limit is DLR.  The unit sphere is compact, and
all writers in the fixed finite card are bounded and local; the inequality
therefore passes to the limit.  For a translation-invariant limit use
\Cref{thm:YMA22-extreme-envelope}.  The Bloch statement follows by the
head--tail triangle inequality already proved in Version~V20.
\end{proof}

\begin{firewall}[Cross-boundary negativity is not a same-phase mode]
\label{fire:YMA22-cross-boundary}
A negative card with $\eta_1\ne\eta_2$ may be caused entirely by the positive
mean-difference Gram in \eqref{eq:YMA22-boundary-pair-split}.  It is then a
boundary/phase-separation witness.  A pure-mode obstruction requires a
negative same-state card, a negative translation-invariant mixture at
nonzero character, or the same-boundary exhaustion above.
\end{firewall}

% =============================================================================
\section{Validated fixed-action nonzero-mode certificate}
\label{sec:YMA22-certificate}
% =============================================================================

\begin{theorem}[Mixture-first obstruction and dispersion-certified pass]
\label{thm:YMA22-validated-certificate}
Fix one shell, one already loaded coefficient $b_r$, one nonzero supported
momentum--colour block, and one target margin $\delta$.  Let
$\overline M$ be the fixed-action phase-average margin matrix and suppose
\begin{equation}
 \|\overline M-\widehat M\|_{\rm op}\le\varepsilon_M.
 \label{eq:YMA22-average-error}
\end{equation}
For a finite retained phase bank, suppose additionally
\begin{equation}
 |\mathcal D_\pi-\widehat{\mathcal D}|
 \le\varepsilon_D,
 \qquad
 \pi_i\ge\pi_*>0,
 \label{eq:YMA22-dispersion-error}
\end{equation}
and put
\[
 D_+=\max\{0,\widehat{\mathcal D}+\varepsilon_D\}.
\]
Then the following one-sided conclusions are rigorous.
\begin{enumerate}[label=\textup{(C22.\arabic*)},leftmargin=5.8em]
\item If
      \begin{equation}
       \lambda_{\min}(\widehat M)+\varepsilon_M<0,
       \label{eq:YMA22-negative-mixed-certificate}
      \end{equation}
      a unit minimum eigenvector $v$ of $\widehat M$ satisfies
      $v^*\overline Mv<0$.  Hence at least one retained phase has
      $v^*M_iv<0$.  This is a typed extremal-phase pure-mode obstruction on
      the same momentum and polarization.
\item If
      \begin{equation}
       \lambda_{\min}(\widehat M)-\varepsilon_M
       -\sqrt{D_+/\pi_*}>0,
       \label{eq:YMA22-all-phase-certificate}
      \end{equation}
      every phase in the bank passes the target margin.
\item If
      \begin{equation}
       \lambda_{\min}(\widehat M)-\varepsilon_M>0
       \label{eq:YMA22-mixture-only-pass}
      \end{equation}
      but \eqref{eq:YMA22-all-phase-certificate} is unavailable, the mixed
      target passes but a particular constituent phase remains
      $\stopstatus$.
\item If none of the displayed strict inequalities holds, the card is
      $\stopstatus$.  A failed sufficient dispersion bound is not a negative
      Wilson mode.
\end{enumerate}
At $q=0$, the phase Gram of Version~V21 must be retained explicitly and this
nonzero-character certificate is inapplicable.
\end{theorem}

\begin{proof}
Let $v$ be a minimum eigenvector of $\widehat M$.  From
\eqref{eq:YMA22-average-error},
\[
 v^*\overline Mv
 \le\lambda_{\min}(\widehat M)+\varepsilon_M.
\]
Thus \eqref{eq:YMA22-negative-mixed-certificate} and
\Cref{thm:YMA22-negative-average-phase} prove item~\textup{(C22.1)}.

Weyl's inequality gives
\[
 \lambda_{\min}(\overline M)
 \ge\lambda_{\min}(\widehat M)-\varepsilon_M.
\]
Apply \eqref{eq:YMA22-phasewise-lower} with
$\mathcal D_\pi\le D_+$ to obtain item~\textup{(C22.2)}.  Item~\textup{(C22.3)}
is the ordinary positive interval for the average together with
\Cref{cth:YMA22-positive-average-not-phasewise}.  The remaining branches are
unresolved by the stated one-sided data.
\end{proof}

\begin{corollary}[Boundary-selected phase as a special zero-dispersion card]
\label{cor:YMA22-selected-phase-specialization}
If a boundary theorem proves that every retained infrared cluster state has
the same nonzero-block margin matrix, then $\mathcal D_\pi=0$ and the mixed,
boundary-selected, and extremal-phase cards coincide on that block.  In this
case \eqref{eq:YMA22-mixture-only-pass} is already a phase-uniform pass.
\end{corollary}

\begin{proof}
Equality of the matrices is equivalent to zero dispersion by
\Cref{cor:YMA22-dispersion-replica}; the common matrix equals its barycentre.
\end{proof}

% =============================================================================
\section{The first literal V23 acquisition after the upstream repair}
\label{sec:YMA22-frontier}
% =============================================================================

The new order removes two circular loops.

First, the programme no longer waits for a preferred generated-potential
ledger when the exact conditional shell kernels and their exterior
oscillation can be tested directly.  Second, it no longer refits the already
loaded marginal coefficient separately in every phase before asking whether
a nonzero physical mode is negative.

\begin{theorem}[Fixed-action obstruction-first termination]
\label{thm:YMA22-termination}
For the native YM--A shell, freeze the response card, coefficient $b_r$,
complete action, gauge/topological packet, and retained nonzero block before
reading any symbol value.  Then proceed as follows.
\begin{enumerate}[label=\textup{(A22.\arabic*)},leftmargin=5.8em]
\item Verify a fixed-shell Gibbs specification either by an absolutely
      summable interaction or by the direct exterior modulus
      \eqref{eq:YMA22-direct-quasilocality}.  A missing exact conditional
      cylinder or unresolved quasilocality row is $\stopstatus$.
\item Form the infrared limit or cluster set of this one fixed action.  Do not
      replace $b_r$ by $b_{r,\theta}$ unless a phase-retuned scheme is
      independently declared.
\item Evaluate the first predeclared retained nonzero-character Wilson margin
      in any available fixed-action mixed or boundary-conditioned limit.
      A negative interval returns the extremal-phase witness of
      \Cref{thm:YMA22-negative-average-phase} before uniqueness is known.
\item If the mixed interval is positive, either identify the selected phase
      by a boundary theorem or acquire the phase-dispersion card.  Apply
      \eqref{eq:YMA22-all-phase-certificate} for a phase-uniform pass.
\item Only after the fixed-action nonzero block is decided should a phasewise
      response residual be used to test the optional stronger demand of
      phasewise response stationarity.
\item Route every constant between-phase score Gram to the harmonic
      zero-character row.  It is not inserted into a Ward quotient at
      nonzero momentum.
\end{enumerate}
The possible outputs are
\[
 \pass,\qquad \obstruction,\qquad \stopstatus.
\]
There is no fourth branch which changes the action coefficient to make the
chosen phase pass.
\end{theorem}

\begin{proof}
Specification existence is
\Cref{thm:YMA22-specification-existence}.  The common action coefficient and
optional retuning discrepancy are
\Cref{thm:YMA22-fixed-action-first,cor:YMA22-retuning-cost}.  Exact removal of
the phase-mean debit at nonzero character is
\Cref{thm:YMA22-zero-total-barycentre,cor:YMA22-nonzero-character}.  The
obstruction and positive alternatives are
\Cref{thm:YMA22-negative-average-phase,thm:YMA22-phase-dispersion,thm:YMA22-validated-certificate}.
The zero-character qualification is the phase-Gram theorem of Version~V21.
\end{proof}

\begin{assessment}[What is now genuinely upstream]
\label{ass:YMA22-genuine-frontier}
The first missing Wilson-specific object is no longer described as
``choose a phase, refit a coefficient, and then return to the symbol.''  It is
one of two concrete fixed-action reads:
\begin{equation}
 \boxed{
 \begin{aligned}
 &\text{the exterior conditional-cylinder modulus }
   \varepsilon_\Lambda(R),\\
 &\text{the first nonzero fixed-action margin matrix }
   \overline M_r(q;\delta)
   \text{ with an outward error interval.}
 \end{aligned}}
 \label{eq:YMA22-two-first-reads}
\end{equation}
A negative second read already returns a bad extreme phase.  A positive
second read requires a boundary-selected phase, the boundary-pair floor, or
the positive phase-dispersion read.  This is strictly closer to the actual Wilson
coefficient than another DLR, Bloch, current, or carrier compiler.
\end{assessment}

\begin{firewall}[No phase or symbol value is manufactured]
\label{fire:YMA22-no-value}
No direct exterior modulus, response cofactor, fixed-action mixed symbol,
phase-dispersion value, selected-phase connected tail, or retained Bloch
eigenvalue is supplied by the present data.  The specification and
barycentric theorems classify what those values would prove.  They do not
establish uniqueness of the four-dimensional Wilson vacuum, positivity of a
selected shell, local-vacuum contraction, an Osterwalder--Schrader time gap,
or a continuum Yang--Mills mass gap.
\end{firewall}

% =============================================================================
\section{Version V22 proof ledger and next exact card}
\label{sec:YMA22-ledger}
% =============================================================================

\begingroup\small
\begin{longtable}{>{\raggedright\arraybackslash}p{0.28\textwidth}
                  >{\raggedright\arraybackslash}p{0.19\textwidth}
                  >{\raggedright\arraybackslash}p{0.45\textwidth}}
\caption{Version V22 proof-status ledger.}\label{tab:YMA22-ledger}\\
\toprule
\textbf{Row}&\textbf{Status}&\textbf{Exact scope}\\
\midrule
\endfirsthead
\toprule
\textbf{Row}&\textbf{Status}&\textbf{Exact scope}\\
\midrule
\endhead
Native response coefficient before phase decomposition
&Corrected exactly
&The finite cofactors load one common action $F_{r,b_r}$; phasewise refitting
 is an optional different scheme.\\[0.3em]
Phasewise response residual
&Proved exactly
&$b_{r,\theta}-b_r=-r_{r,\theta}/d_{r,\theta}$ on a regular phase, with the
 explicit Hessian perturbation bound
 \eqref{eq:YMA22-retuning-margin-cost}.\\[0.3em]
Specification-first DLR entrance
&Proved conditionally
&Proper, consistent exact conditional kernels with
 $\varepsilon_\Lambda(R)\to0$ possess compact nonempty DLR phase sets.\\[0.3em]
Conditional density comparison
&Proved exactly
&The boundary Hamiltonian oscillation gives the direct total-variation bound
 \eqref{eq:YMA22-TV-bound}.\\[0.3em]
Generated Wilson exterior modulus
&Open physical estimate
&No bound on $\varepsilon_\Lambda(R)$ is populated for the complete nonlinear
 response-matched shell.\\[0.3em]
Zero-total phase barycentre
&Proved exactly
&Every finite zero-total field packet has Hessian equal to the average of its
 fixed-action phase Hessians.\\[0.3em]
Nonzero periodic character
&Proved exactly
&The between-phase score-mean Gram cancels by character orthogonality.\\[0.3em]
Negative mixed nonzero card
&Proved to expose a phase
&A negative vector for the fixed-action average is negative in a positive
 measure set of constituent phases.\\[0.3em]
Positive mixed nonzero card
&Insufficient phasewise
&An explicit two-phase counterexample shows that a positive average can mask
 a negative phase.\\[0.3em]
Boundary-pair DLR card
&Proved exactly
&The same-boundary floor minus one positive boundary-mean-spread stock gives
 a uniqueness-free lower bound for every DLR state.\\[0.3em]
Negative boundary exhaustion
&Proved conditionally
&A separated same-boundary card has a DLR cluster obstruction; Bloch passage
 uses the already stated uniform connected tail.\\[0.3em]
Finite phase dispersion
&Proved exactly
&One positive pair-difference scalar plus a phase-weight floor upgrades the
 average according to \eqref{eq:YMA22-uniform-phase-pass}.\\[0.3em]
First fixed-action Wilson margin
&\stopstatus
&No outward matrix interval is supplied on the first actual retained nonzero
 momentum--colour block.\\[0.3em]
Local-vacuum contraction and mass
&Not inferred
&The phase/specification corrections do not populate higher jets, temporal
 transfer, reflection-positive continuum, or a mass gap.\\
\bottomrule
\end{longtable}
\endgroup

\begin{theorem}[Version V22 fixed-action correction]
\label{thm:YMA22-terminal}
Versions~V1--V21 are retained at their stated exact, finite, conditional, and
open scopes.  Version~V22 additionally proves:
\begin{enumerate}[label=\textup{(E22.\arabic*)},leftmargin=5.8em]
\item the native frozen response coefficient loads one action before its DLR
      phase decomposition;
\item phasewise refitting differs from that action by the exact response
      residual quotient;
\item a direct exterior conditional-cylinder modulus is sufficient for DLR
      existence and subsumes the absolutely summable-potential route;
\item every zero-total field packet, hence every nonzero periodic character,
      removes the between-phase score mean exactly;
\item a negative fixed-action average exposes a bad extreme phase on the
      same momentum--colour polarization;
\item a positive average is upgraded to a phase-uniform pass by either the
      boundary-pair floor or the positive phase-dispersion certificate;
\item the finite outward conditions in
      \Cref{thm:YMA22-validated-certificate} give a terminating
      obstruction/pass/stop card.
\end{enumerate}
No item evaluates the required Wilson-specific physical value.
\end{theorem}

\begin{proof}
Items~\textup{(E22.1)--(E22.2)} are
\Cref{thm:YMA22-fixed-action-first,cor:YMA22-retuning-cost}.  Item~\textup{(E22.3)}
is \Cref{lem:YMA22-log-density-TV,thm:YMA22-specification-existence,cor:YMA22-potential-implies-modulus}.
Items~\textup{(E22.4)--(E22.5)} are
\Cref{thm:YMA22-zero-total-barycentre,cor:YMA22-nonzero-character,thm:YMA22-negative-average-phase}.
Items~\textup{(E22.6)--(E22.7)} are
\Cref{thm:YMA22-boundary-pair-disintegration,thm:YMA22-boundary-spread,%
thm:YMA22-phase-dispersion,thm:YMA22-validated-certificate}.
\end{proof}

\begin{openproblem}[V23: execute the fixed-action nonzero card before another phase compiler]
\label{op:YMA22-V23}
Preserve every frozen response, action, gauge, topology, colour, tangent,
carrier, reference, shell, and infrared-regulator convention.  Do not refit
the marginal coefficient phase by phase unless the programme is explicitly
changed to the phase-retuned scheme.
\begin{enumerate}[label=\textup{(V23.\arabic*)},leftmargin=5.6em]
\item Load the two finite response cofactors and the native coefficient $b_r$.
\item For the resulting exact conditional shell, acquire either the direct
      exterior modulus $\varepsilon_\Lambda(R)$ or the first boundary pair
      on which it fails to decay.
\item On one predeclared retained $q\ne0$ block, evaluate the fixed-action
      mixed or boundary-conditioned likelihood/Hessian card with an outward
      Hermitian error radius.
\item If its upper interval is negative, return the explicit
      extremal-phase pure-mode witness and stop that branch without proving
      uniqueness.
\item If its lower interval is positive, acquire either the selected-boundary
      phase card, the boundary-pair stocks
      \eqref{eq:YMA22-boundary-two-stocks}, or the positive phase-dispersion
      scalar.  Apply \eqref{eq:YMA22-boundary-all-phase-pass} or
      \eqref{eq:YMA22-all-phase-certificate} before declaring a phase-uniform
      pass.
\item Only after that block is decided, populate the selected-phase
      $\Wtwo/\Wthree$ connected tail, higher jets, and adjacent-shell
      transport.
\end{enumerate}
If none of the two physical reads in \eqref{eq:YMA22-two-first-reads} is
available, the exact output is $\stopstatus$, not another abstract phase or
Bloch theorem.
\end{openproblem}

\begin{assessment}[V22 termination]
\label{ass:YMA22-termination}
Version~V21 identified a real phase-mixture debit, but its phase-first
response refit would change the native action already frozen in Version~V16.
The present correction keeps that action fixed, replaces a preferred
potential ledger by an exact conditional-cylinder modulus, and proves that
nonzero-character testing may proceed obstruction first.  A negative mixed
card would already expose a bad extreme phase; a positive card still needs a
selected phase, a boundary-pair floor, or phase-dispersion control.  Since neither physical value has
been supplied, the selected Wilson margin remains $\stopstatus$ and no
negative mode has been produced.
\end{assessment}

The accompanying replay checks the finite algebra of this continuation:
zero-total phase affinity, zero-character support of the phase Gram, the
extreme-envelope minimum on a finite phase polytope, the negative- and
positive-mixture alternatives, the boundary-pair decomposition, the exact
DLR pair-average algebra, and one sign-safe boundary-floor interval.  It does
not prove a Wilson quasilocality modulus, construct a DLR phase, or evaluate
a Yang--Mills symbol.

The standalone verifier and its result are
\begin{center}
\path{YMA_V22_phase_envelope_replay.py},\qquad
\path{YMA_V22_phase_envelope_replay.json}.
\end{center}
Their SHA--256 digests, followed by the replay-package digest, are
\begin{center}\scriptsize\ttfamily
345437a0988d5123b42635b2c5f32676d182ca642fbef08301b97d0be3b21154,\\
71c2997eb4b9dfd68adce05e0fa81e36eec54b88fe755ff48d013f10f96b96ff,\\
1b88891263985f6f90465f938a66761b008eeab74dd751a13e0f7807852f972a.
\end{center}
All seven declared replay groups pass.

\section*{Version V22 append-only continuation record}
\addcontentsline{toc}{section}{Version V22 append-only continuation record}
The complete Version~V21 source above is preserved byte for byte up to its
sole terminal document-closing command.  Its SHA--256 digest is
\begin{center}\scriptsize\ttfamily
6248af751ff5099ed38b04c49bc22bf92df667d18d71498e1ba1b49304819971.
\end{center}
For Version~V23, preserve this accumulated source, append new material
immediately before the final document-closing command, and increment the
filename to \texttt{\_V23.tex}.  Keep the native finite-shell action
coefficient fixed before the DLR decomposition.  Test a nonzero-character
fixed-action block for an obstruction before demanding phase uniqueness; a
positive phase claim still requires boundary selection or a phase-dispersion
certificate.

% =============================================================================
% END APPENDED VERSION V22 MATERIAL
% =============================================================================



% =============================================================================
% START APPENDED VERSION V23 MATERIAL
% =============================================================================

\clearpage
\hypersetup{
  pdftitle={Renewal Yang--Mills Reduced-Fibre Wilson Shell Symbol and Local-Vacuum Programme V23},
  pdfsubject={Target-quotient half-density transport, boundary-condition gain, and a noninflated fixed-action Wilson card}
}

\section*{Version V23: target-quotient half-density transport, boundary-condition gain, and one noninflated Wilson card}
\addcontentsline{toc}{section}{Version V23: target-quotient half-density transport and one noninflated Wilson card}

Version~V22 placed one fixed response-matched action before its infrared phase
decomposition and showed that a nonzero-character obstruction can be read before
vacuum uniqueness is known.  Its first proposed entrance still asked for either a
quasilocality modulus of the complete conditional shell or a phase-level boundary
packet.  The related boundary-sector continuation identifies a smaller object for a
smaller target.  When two laws act on the same multiplication-writer algebra, the
unique order-positive Hilbert incidence is multiplication by the square root of the
Radon--Nikodym derivative.  One nonnegative half-density difference then measures
the endpoint source displacement.  No stochastic coupling or path law is required
unless chronology is itself part of the claim.

The present continuation imports that module-incidence theorem without reproving
it and applies it to the \emph{target quotient} generated by the complete one-copy
Wilson jet needed by one predeclared momentum--colour block.  The full conditional
configuration law is not reconstructed merely to compare one matrix expectation.
At the same time, the target writer norm is retained.  This is essential: the
related fixed-datum and terminal-source programmes show in different settings that
a small source residual before the physical route need not imply a small response
after a large source-to-output gain.  Here the corresponding gain is an explicit
operator oscillation radius of the assembled Hessian writer.

The revised order is
\begin{equation}
 \boxed{
 \begin{gathered}
 \text{one fixed response-matched Wilson action and one retained }q\ne0\text{ block}
 \\
 \Downarrow\\
 \text{one-copy target record }(J_{q,R},K_{q,R})
 \text{ and its boundary-conditioned pushforward laws}
 \\
 \Downarrow\\
 \text{one positive target-module distance and one reference boundary card}
 \\
 \Downarrow\\
 \text{target-writer gain plus the already required Ward-relative tail}
 \\
 \Downarrow\\
 \text{all-boundary pass, all-boundary obstruction, or a typed boundary/gain stop.}
 \end{gathered}}
 \label{eq:YMA23-order}
\end{equation}
This is a target-relative reduction of the Version~V22 specification entrance.  It
does not replace the selected local vacuum when higher jets, products, a complete
local algebra, or local-vacuum contraction are later claimed.

% =============================================================================
\section{The one-copy Wilson target record and the assembled pair writer}
\label{sec:YMA23-target-record}
% =============================================================================

Fix throughout this continuation:
\begin{enumerate}[label=\textup{(F23.\arabic*)},leftmargin=5.2em]
\item one shell index and the native response-matched action already loaded before
      any infrared decomposition;
\item one nonzero covariant momentum, toron/stabilizer block, rooted-Hodge support,
      and finite retained coefficient space $E$;
\item one finite displacement radius $R$ and the complete deterministic,
      generated, support-birth, and local topological conventions on that head;
\item the same flat reference and target margin $\delta$ used in Versions~V1--V22.
\end{enumerate}
Write $d=\dim E$.  On one interior configuration $z$, let
\begin{equation}
 J_{q,R}(z)\in\C^d,
 \qquad
 K_{q,R}(z)=K_{q,R}(z)^*\in\operatorname{Herm}_d
 \label{eq:YMA23-one-copy-jets}
\end{equation}
be the complete first and second field jets of the finite head.  All terms in
\eqref{eq:YMA23-one-copy-jets} are evaluated from the same action and on the same
history.  Put
\begin{equation}
 L_{q,R}^{\delta}=T_{q,R}+\delta S_{q,R},
 \label{eq:YMA23-deterministic-subtraction}
\end{equation}
where $T_{q,R}$ is the declared deterministic local/topological subtraction and
$S_{q,R}$ is the flat reference on the same supported block.

\begin{definition}[One-copy target record]
\label{def:YMA23-target-record}
The finite Wilson target record is
\begin{equation}
 \boxed{
 \mathfrak T_{q,R}(z)
 =\bigl(J_{q,R}(z),K_{q,R}(z)\bigr)
 \in\mathcal Y_{q,R}\subseteq
 \C^d\times\operatorname{Herm}_d.}
 \label{eq:YMA23-target-map}
\end{equation}
The target space is given the Borel structure inherited from the finite-dimensional
ambient space.  On the compact reduced fibre its image has compact closure.
\end{definition}

For $y_i=(J_i,K_i)\in\mathcal Y_{q,R}$ define the Hermitian pair writer
\begin{equation}
 \boxed{
 \mathscr M_{q,R}^{\delta}(y_1,y_2)
 =\frac12(K_1+K_2)
  -\frac12(J_1-J_2)(J_1-J_2)^*
  -L_{q,R}^{\delta}.}
 \label{eq:YMA23-pair-writer}
\end{equation}
When a Ward quotient is desired, put
\begin{equation}
 \mathscr G_{q,R}^{\delta}
 =\omega(q)^{-1}\mathscr M_{q,R}^{\delta},
 \qquad
 \omega(q)>0.
 \label{eq:YMA23-Ward-pair-writer}
\end{equation}
No division is made at the zero character.

Let $\lambda$ be any probability law on $\mathcal Y_{q,R}$.  Define
\begin{equation}
 \mathbb Q_{q,R}^{\delta}[\lambda_1,\lambda_2]
 :=\mathbb E_{\lambda_1\otimes\lambda_2}
       \mathscr G_{q,R}^{\delta},
 \label{eq:YMA23-cross-card}
\end{equation}
and abbreviate
$\mathbb Q_{q,R}^{\delta}[\lambda]
 =\mathbb Q_{q,R}^{\delta}[\lambda,\lambda]$.

\begin{lemma}[The target record reconstructs the complete finite-head card]
\label{lem:YMA23-target-reconstructs-card}
For every probability law $\lambda$ on the target record,
\begin{equation}
 \boxed{
 \mathbb Q_{q,R}^{\delta}[\lambda]
 =\frac1{\omega(q)}
 \left(
   \mathbb E_\lambda K_{q,R}
   -\operatorname{Cov}_\lambda(J_{q,R})
   -T_{q,R}-\delta S_{q,R}
 \right).}
 \label{eq:YMA23-card-identity}
\end{equation}
For two different laws,
\begin{equation}
 \boxed{
 \begin{aligned}
 \mathbb Q_{q,R}^{\delta}[\lambda_1,\lambda_2]
 =\frac1{\omega(q)}\bigg[{}
 &\frac12\bigl(\mathbb E_{\lambda_1}K_{q,R}
                +\mathbb E_{\lambda_2}K_{q,R}\bigr)
 -L_{q,R}^{\delta}\\
 &-\frac12\mathbb E_{\lambda_1\otimes\lambda_2}
   (J_1-J_2)(J_1-J_2)^*
 \bigg].
 \end{aligned}}
 \label{eq:YMA23-cross-card-identity}
\end{equation}
\end{lemma}

\begin{proof}
For independent copies $J_1,J_2$ with common law $\lambda$,
\[
 \frac12\mathbb E(J_1-J_2)(J_1-J_2)^*
 =\mathbb E(JJ^*)-(\mathbb EJ)(\mathbb EJ)^*
 =\operatorname{Cov}_\lambda(J).
\]
Substitution in \eqref{eq:YMA23-pair-writer} proves
\eqref{eq:YMA23-card-identity}.  The second identity is the definition with no
further simplification.
\end{proof}

\begin{firewall}[The record is joint and action specific]
\label{fire:YMA23-joint-record}
The entries of $\mathfrak T_{q,R}$ are the field jets of one already fixed action
on one history.  A $K$ marginal from one action, a score marginal from another
law, or separately fitted boundary panels do not populate
\eqref{eq:YMA23-pair-writer}.  The target quotient reduces irrelevant
configuration information; it does not manufacture the same-history incidence of
the retained writers.
\end{firewall}

\begin{firewall}[Common value space versus a common target map]
\label{fire:YMA23-common-target-map}
The direct module distance compares the laws of the actual jet values on the
common finite-dimensional value space $\mathcal Y_{q,R}$; the boundary-dependent
maps which produce those values need not be identified.  By contrast,
\Cref{lem:YMA23-data-processing,prop:YMA23-exterior-implies-module} may be used to
push a full conditional-law estimate down to the target only when the same
measurable map $T$ is used for both laws.  In the Wilson application this holds
when the conditioned region contains the complete support of the finite-head jet,
so that the target map is independent of the exterior condition.  Otherwise the
target distance must be acquired directly on $\mathcal Y_{q,R}$; no data-processing
claim is made from the full conditional laws.
\end{firewall}

% =============================================================================
\section{Target-quotient half-density distance}
\label{sec:YMA23-module-distance}
% =============================================================================

Let $\lambda$ and $\nu$ be probability laws on a common measurable target space.
Choose any dominating probability measure $m$, with densities $p$ and $q$.

\begin{definition}[Target-module distance and affinity]
\label{def:YMA23-module-distance}
Set
\begin{equation}
 \boxed{
 \mathsf D(\lambda,\nu)
 :=\int(\sqrt p-\sqrt q)^2\,\dd m
 =2-2\mathsf A(\lambda,\nu),}
 \label{eq:YMA23-module-distance}
\end{equation}
where
\begin{equation}
 \mathsf A(\lambda,\nu)
 :=\int\sqrt{pq}\,\dd m.
 \label{eq:YMA23-affinity}
\end{equation}
Then $0\le\mathsf D\le2$.  Both quantities are independent of $m$.
\end{definition}

The boundary-sector Version~V22 theorem identifies multiplication by
$\sqrt{\dd\lambda/\dd\nu}$ as the unique order-positive unitary intertwining the
complete common target-writer algebra.  We use only that imported incidence and
its distance; no stochastic coupling between the two laws is introduced.

\begin{proposition}[One bounded positive Read]
\label{prop:YMA23-positive-distance-read}
Let $\overline\lambda=(\lambda+\nu)/2$.  On $p+q>0$ define
\begin{equation}
 \boxed{
 \mathscr D_{\lambda/\nu}(y)
 =\frac{2(\sqrt p-\sqrt q)^2}{p+q},}
 \label{eq:YMA23-positive-distance-writer}
\end{equation}
and set it to zero on $p+q=0$.  Then
\begin{equation}
 \boxed{
 0\le\mathscr D_{\lambda/\nu}\le2,
 \qquad
 \mathbb E_{\overline\lambda}\mathscr D_{\lambda/\nu}
 =\mathsf D(\lambda,\nu).}
 \label{eq:YMA23-positive-distance-expectation}
\end{equation}
Thus one nonnegative scalar on the fifty-fifty target-law mixture acquires the
complete module distance.
\end{proposition}

\begin{proof}
The density of $\overline\lambda$ relative to $m$ is $(p+q)/2$, so the factors
cancel in the expectation.  The pointwise estimate
$(\sqrt p-\sqrt q)^2\le p+q$ gives the upper bound two.
\end{proof}

\begin{lemma}[Target pushforward contracts module distance]
\label{lem:YMA23-data-processing}
Let $T:\Omega\to\mathcal Y$ be measurable and let $\mu,\nu$ be probability laws
on $\Omega$.  Then
\begin{equation}
 \boxed{
 \mathsf D(T_*\mu,T_*\nu)
 \le\mathsf D(\mu,\nu).}
 \label{eq:YMA23-data-processing}
\end{equation}
Equality is not required for the target card.
\end{lemma}

\begin{proof}
Take $m=(\mu+\nu)/2$ and write $p=\dd\mu/\dd m$,
$q=\dd\nu/\dd m$.  Relative to $T_*m$, the two pushforward densities are the
conditional expectations
\[
 p_T\circ T=\mathbb E_m[p\mid\sigma(T)],
 \qquad
 q_T\circ T=\mathbb E_m[q\mid\sigma(T)].
\]
Conditional Cauchy--Schwarz gives
\[
 \mathbb E_m[\sqrt{pq}\mid\sigma(T)]
 \le\sqrt{p_Tq_T}\circ T.
\]
After integration,
$\mathsf A(T_*\mu,T_*\nu)\ge\mathsf A(\mu,\nu)$.  Use
\eqref{eq:YMA23-module-distance}.
\end{proof}

\begin{assessment}[Target-relative meaning]
\label{ass:YMA23-target-relative-meaning}
A full conditional-shell Hellinger or total-variation estimate is a sufficient
route to \eqref{eq:YMA23-module-distance}, but it is not necessary for one Wilson
block.  The target record discards every exterior distinction which has no effect
on the actual finite-head score and curvature writers.  Conversely, a small target
module distance is only a law-comparison statement.  Its conversion into a margin
comparison still carries the target-writer gain derived below.
\end{assessment}

% =============================================================================
\section{Product amplitudes and the target-writer transport bound}
\label{sec:YMA23-product-transport}
% =============================================================================

\begin{lemma}[Exact product-law module distance]
\label{lem:YMA23-product-distance}
Let $\lambda_0,\lambda_1,\lambda_2$ be probability laws on the same target space
and put
\begin{equation}
 d_i=\mathsf D(\lambda_i,\lambda_0),
 \qquad i=1,2.
 \label{eq:YMA23-di}
\end{equation}
Then
\begin{equation}
 \boxed{
 \mathsf D(\lambda_1\otimes\lambda_2,
           \lambda_0\otimes\lambda_0)
 =d_1+d_2-\frac12d_1d_2.}
 \label{eq:YMA23-product-distance}
\end{equation}
In particular, if $d_1,d_2\le d_*$, then
\begin{equation}
 \boxed{
 \mathsf D(\lambda_1\otimes\lambda_2,
           \lambda_0\otimes\lambda_0)
 \le d_*^{(2)}:=2d_*-\frac12d_*^2.}
 \label{eq:YMA23-product-distance-uniform}
\end{equation}
\end{lemma}

\begin{proof}
The Hellinger affinity is multiplicative under products.  Since
$\mathsf A(\lambda_i,\lambda_0)=1-d_i/2$,
\[
 \mathsf D(\lambda_1\otimes\lambda_2,
           \lambda_0\otimes\lambda_0)
 =2-2(1-d_1/2)(1-d_2/2),
\]
which is \eqref{eq:YMA23-product-distance}.  The right-hand side is increasing
in each $d_i\in[0,2]$, proving the uniform bound.
\end{proof}

\begin{definition}[Operator oscillation radius of the assembled writer]
\label{def:YMA23-writer-radius}
For a bounded Hermitian matrix writer $\mathscr A$ on a probability space, put
\begin{equation}
 \boxed{
 \operatorname{rad}_{\op}(\mathscr A)
 :=\inf_{C=C^*}
   \operatorname*{ess\,sup}\norm{\mathscr A-C}_{\op}.}
 \label{eq:YMA23-writer-radius}
\end{equation}
For the Ward-normalized finite-head writer define
\begin{equation}
 \mathfrak b_{q,R}^{\delta}
 :=\operatorname{rad}_{\op}(\mathscr G_{q,R}^{\delta}).
 \label{eq:YMA23-Wilson-gain}
\end{equation}
The deterministic constant $L_{q,R}^{\delta}/\omega(q)$ does not affect this
radius.
\end{definition}

\begin{theorem}[Half-density transport of a Hermitian target]
\label{thm:YMA23-matrix-transport}
Let $\rho$ and $\sigma$ be probability laws on one space and let
$\mathscr A=\mathscr A^*$ be a bounded $d\times d$ matrix writer.  Put
$D=\mathsf D(\rho,\sigma)$.  Then
\begin{equation}
 \boxed{
 \norm{\mathbb E_\rho\mathscr A-
       \mathbb E_\sigma\mathscr A}_{\op}
 \le
 \operatorname{rad}_{\op}(\mathscr A)
 \sqrt{D(4-D)}.}
 \label{eq:YMA23-matrix-transport}
\end{equation}
Consequently,
\begin{equation}
 \boxed{
 \begin{aligned}
 &\norm{
 \mathbb Q_{q,R}^{\delta}[\lambda_1,\lambda_2]
 -\mathbb Q_{q,R}^{\delta}[\lambda_0,\lambda_0]
 }_{\op}\\
 &\quad\le
 \mathfrak b_{q,R}^{\delta}
 \sqrt{d_{12}(4-d_{12})},
 \qquad
 d_{12}=d_1+d_2-\frac12d_1d_2.
 \end{aligned}}
 \label{eq:YMA23-cross-boundary-transport}
\end{equation}
\end{theorem}

\begin{proof}
Choose a common dominating measure with densities $r,s$.  For any Hermitian
constant $C$ and any unit vector $v$, put
$f_v=v^*(\mathscr A-C)v$.  Since constants have the same expectation under both
laws,
\[
 v^*(\mathbb E_\rho\mathscr A-\mathbb E_\sigma\mathscr A)v
 =\int f_v(r-s)\,\dd m.
\]
Factor $r-s=(\sqrt r-\sqrt s)(\sqrt r+\sqrt s)$ and apply Cauchy--Schwarz:
\[
 \left|\int f_v(r-s)\,\dd m\right|
 \le\norm{\mathscr A-C}_{L^\infty(\op)}
      \norm{\sqrt r-\sqrt s}_2
      \norm{\sqrt r+\sqrt s}_2.
\]
The squared norms of the last two factors are $D$ and $4-D$.  Take the supremum
in $v$, then the infimum in $C$.  Apply the result to the two product laws and use
\Cref{lem:YMA23-product-distance}.
\end{proof}

\begin{corollary}[Uniform reference-boundary radius]
\label{cor:YMA23-uniform-boundary-radius}
Let $\lambda_\eta$ be the target-record law under exterior condition $\eta$, and
fix a reference exterior $\eta_0$.  Assume
\begin{equation}
 \sup_\eta\mathsf D(\lambda_\eta,\lambda_{\eta_0})\le d_*\le2.
 \label{eq:YMA23-uniform-target-distance}
\end{equation}
Then every boundary pair satisfies
\begin{equation}
 \boxed{
 \norm{
 \mathbb Q_{q,R}^{\delta}[\lambda_{\eta_1},\lambda_{\eta_2}]
 -\mathbb Q_{q,R}^{\delta}[\lambda_{\eta_0}]
 }_{\op}
 \le\varepsilon_{\rm mod},}
 \label{eq:YMA23-uniform-boundary-card}
\end{equation}
where
\begin{equation}
 \boxed{
 \varepsilon_{\rm mod}
 =\mathfrak b_{q,R}^{\delta}
   \sqrt{d_*^{(2)}(4-d_*^{(2)})},
 \qquad
 d_*^{(2)}=2d_*-\frac12d_*^2.}
 \label{eq:YMA23-module-error}
\end{equation}
\end{corollary}

\begin{proof}
Combine \Cref{lem:YMA23-product-distance,thm:YMA23-matrix-transport}.
\end{proof}

% =============================================================================
\section{A uniqueness-free all-boundary Wilson certificate}
\label{sec:YMA23-all-boundary}
% =============================================================================

Let the complete boundary-pair Ward quotient be decomposed as
\begin{equation}
 \mathbb M_{q}^{\delta}[\eta_1,\eta_2]
 =\mathbb Q_{q,R}^{\delta}[\lambda_{\eta_1},\lambda_{\eta_2}]
  +\mathbb R_{q,R}^{\delta}[\eta_1,\eta_2],
 \label{eq:YMA23-head-tail-boundary}
\end{equation}
where the remainder contains every omitted differentiated connected displacement,
nonlocal support birth, and separately declared topological contribution.  Assume
the sign-symmetric outward bound
\begin{equation}
 \boxed{
 \sup_{\eta_1,\eta_2}
 \norm{\mathbb R_{q,R}^{\delta}[\eta_1,\eta_2]}_{\op}
 \le\tau_{q,R}.}
 \label{eq:YMA23-uniform-tail}
\end{equation}
This is the boundary-uniform form of the existing Ward-relative tail row; it is not
inferred from the module distance.

\begin{theorem}[Reference boundary plus one target-module radius controls every DLR phase]
\label{thm:YMA23-all-DLR-certificate}
Assume \eqref{eq:YMA23-uniform-target-distance} and
\eqref{eq:YMA23-uniform-tail}.  Suppose a validated Hermitian matrix
$\widehat Q_0$ satisfies
\begin{equation}
 \norm{\mathbb Q_{q,R}^{\delta}[\lambda_{\eta_0}]-\widehat Q_0}_{\op}
 \le\varepsilon_0.
 \label{eq:YMA23-reference-error}
\end{equation}
Put
\begin{equation}
 \boxed{
 \varepsilon_{\rm tot}
 =\varepsilon_0+\varepsilon_{\rm mod}+\tau_{q,R}.}
 \label{eq:YMA23-total-radius}
\end{equation}
Then:
\begin{enumerate}[label=\textup{(B23.\arabic*)},leftmargin=5.2em]
\item If
      \begin{equation}
       \boxed{
       \lambda_{\min}(\widehat Q_0)-\varepsilon_{\rm tot}>0,}
       \label{eq:YMA23-all-boundary-pass}
      \end{equation}
      every boundary-pair card is positive definite.  Every DLR state of the
      same fixed action passes the selected finite-head-plus-tail block.
\item If
      \begin{equation}
       \boxed{
       \lambda_{\min}(\widehat Q_0)+\varepsilon_{\rm tot}<0,}
       \label{eq:YMA23-all-boundary-obstruction}
      \end{equation}
      a minimum eigenvector $v_0$ of $\widehat Q_0$ is negative for every
      boundary-pair card and therefore for every DLR state of the same fixed
      action.
\item If neither strict inequality holds, the stated data return $\stopstatus$.
\end{enumerate}
No vacuum-uniqueness theorem is used.
\end{theorem}

\begin{proof}
By \Cref{cor:YMA23-uniform-boundary-radius} and the two error bounds,
\begin{equation}
 \norm{\mathbb M_q^{\delta}[\eta_1,\eta_2]-\widehat Q_0}_{\op}
 \le\varepsilon_{\rm tot}
 \label{eq:YMA23-all-boundary-error}
\end{equation}
for every exterior pair.  Weyl's inequality proves the first item.  For the second,
\[
 v_0^*\mathbb M_q^{\delta}[\eta_1,\eta_2]v_0
 \le\lambda_{\min}(\widehat Q_0)+\varepsilon_{\rm tot}<0.
\]
The boundary-pair DLR disintegration proved in Version~V22 expresses the margin of
any DLR state as the barycentre of these cross-boundary cards under two independent
exterior draws.  Positivity or negativity on the common vector therefore passes
through the integral.
\end{proof}

\begin{corollary}[Selected boundary branch]
\label{cor:YMA23-selected-boundary}
Let $\eta_N$ be a predeclared boundary sequence and suppose its target laws satisfy
\begin{equation}
 \mathsf D(\lambda_{\eta_N},\lambda_*)\longrightarrow0
 \label{eq:YMA23-selected-target-convergence}
\end{equation}
for one target law $\lambda_*$, while the target-writer radius is finite and the
Ward-relative tail converges.  Then the selected boundary cards converge in
operator norm to the card of $\lambda_*$.  A separated sign at the target law is
eventually attained by the finite selected cards.
\end{corollary}

\begin{proof}
Apply \Cref{lem:YMA23-product-distance,thm:YMA23-matrix-transport} with both
endpoint laws equal to $\lambda_{\eta_N}$ and the reference equal to $\lambda_*$.
The product distance tends to zero.  Add the assumed tail convergence and use
Weyl's inequality.
\end{proof}

\begin{theorem}[Target-card limit without reconstruction of a complete state]
\label{thm:YMA23-target-card-Cauchy}
Fix one finite target record and one bounded pair writer $\mathscr G$.  Let
$\lambda_n$ be any sequence of target laws and put
$d_n=\mathsf D(\lambda_{n+1},\lambda_n)$.  If
\begin{equation}
 \boxed{
 \sum_{n=1}^{\infty}
 \operatorname{rad}_{\op}(\mathscr G)
 \sqrt{d_n^{(2)}(4-d_n^{(2)})}<\infty,
 \qquad
 d_n^{(2)}=2d_n-\frac12d_n^2,}
 \label{eq:YMA23-target-Cauchy-condition}
\end{equation}
then the Hermitian matrices
$\mathbb E_{\lambda_n^{\otimes2}}\mathscr G$ converge in operator norm.
This conclusion constructs the selected target value only.  It does not construct a
DLR state, a local algebra, products outside the target record, or higher field
jets.
\end{theorem}

\begin{proof}
The adjacent matrix difference is bounded by the $n$th summand in
\eqref{eq:YMA23-target-Cauchy-condition}, by
\Cref{lem:YMA23-product-distance,thm:YMA23-matrix-transport}.  The matrices are
therefore Cauchy in the complete finite-dimensional operator norm.
\end{proof}

% =============================================================================
\section{Relation to the Version V22 exterior modulus}
\label{sec:YMA23-exterior-modulus}
% =============================================================================

\begin{proposition}[Exterior Hamiltonian oscillation implies target-module control]
\label{prop:YMA23-exterior-implies-module}
Let $\gamma_\Lambda^\eta$ and $\gamma_\Lambda^{\eta'}$ be two finite conditional
laws with densities proportional to $e^{-H_\eta}$ and $e^{-H_{\eta'}}$ relative to
the same compact product-Haar measure.  Assume that for some scalar $c$,
\begin{equation}
 \sup_z|H_\eta(z)-H_{\eta'}(z)-c|\le\varepsilon.
 \label{eq:YMA23-Hamiltonian-oscillation}
\end{equation}
Then
\begin{equation}
 \boxed{
 \mathsf D(\gamma_\Lambda^\eta,
           \gamma_\Lambda^{\eta'})
 \le\min\{2,(e^\varepsilon-1)^2\}.}
 \label{eq:YMA23-Hellinger-from-oscillation}
\end{equation}
For every target map $T$,
\begin{equation}
 \boxed{
 \mathsf D(T_*\gamma_\Lambda^\eta,
           T_*\gamma_\Lambda^{\eta'})
 \le\min\{2,(e^\varepsilon-1)^2\}.}
 \label{eq:YMA23-target-Hellinger-from-oscillation}
\end{equation}
Thus the direct exterior modulus of Version~V22 is a sufficient, but generally
stronger, route to the present target-module card.
\end{proposition}

\begin{proof}
The assumption gives
\[
 e^{-c-\varepsilon}
 \le e^{-(H_\eta-H_{\eta'})}
 \le e^{-c+\varepsilon}.
\]
After integration, the corresponding partition-function ratio lies between the
same two endpoints.  Hence the normalized likelihood ratio $r$ satisfies
\[
 e^{-2\varepsilon}\le r\le e^{2\varepsilon}.
\]
Taking square roots gives
$|\sqrt r-1|\le e^\varepsilon-1$.  Since
$\mathsf D=\mathbb E_{\gamma^{\eta'}}(\sqrt r-1)^2$, the first bound follows,
with the universal cap two from \Cref{def:YMA23-module-distance}.  Apply
\Cref{lem:YMA23-data-processing} for the target statement.
\end{proof}

\begin{assessment}[What is no longer first]
\label{ass:YMA23-no-full-spec-first}
To evaluate one finite retained Wilson block, it is unnecessary to prove
quasilocality of every conditional writer or uniqueness of the complete Gibbs
specification before reading the target.  One reference card, one target-module
distance, the target-writer gain, and the existing signed tail are sufficient.
The full exterior modulus remains necessary when the target is a complete local
state, arbitrary local observables, products, or an all-cylinder DLR theorem.
\end{assessment}

% =============================================================================
\section{A target-Fisher bridge as an alternative positive route}
\label{sec:YMA23-target-Fisher}
% =============================================================================

Let $t\mapsto\mu_t$ be a differentiable path of probability laws on a fixed finite
configuration cylinder, with density $r_t$ relative to a common measure and centered
score.  Assume differentiation under the integral is valid and that the pushforward
laws below admit a common dominating target measure.  Write
\begin{equation}
 s_t=\partial_t\log r_t,
 \qquad
 \mathbb E_{\mu_t}s_t=0.
 \label{eq:YMA23-full-score}
\end{equation}
Let $T$ be the fixed one-copy target map and put $\lambda_t=T_*\mu_t$.

\begin{theorem}[Target-quotient Fisher speed]
\label{thm:YMA23-target-Fisher-speed}
The score of the target law is
\begin{equation}
 \boxed{
 \bar s_t(T)
 =\mathbb E_{\mu_t}[s_t\mid\sigma(T)].}
 \label{eq:YMA23-target-score}
\end{equation}
If $p_t$ is the density of $\lambda_t$ relative to a common target measure, then
\begin{equation}
 \boxed{
 \left\|\partial_t\sqrt{p_t}\right\|_2^2
 =\frac14\mathbb E_{\lambda_t}\bar s_t^2.}
 \label{eq:YMA23-target-speed}
\end{equation}
Consequently,
\begin{equation}
 \boxed{
 \sqrt{\mathsf D(\lambda_0,\lambda_1)}
 \le\frac12\int_0^1
     \sqrt{\mathbb E_{\lambda_t}\bar s_t^2}\,\dd t
 \le\frac12\int_0^1
     \sqrt{\mathbb E_{\mu_t}s_t^2}\,\dd t.}
 \label{eq:YMA23-target-length}
\end{equation}
For an action bridge, $s_t=-\dot F_t+\mathbb E_t\dot F_t$, so the first integrand
is the variance of the action score after projection to the actual Wilson target
record.
\end{theorem}

\begin{proof}
For every bounded target function $f$,
\[
 \frac{\dd}{\dd t}\mathbb E_{\lambda_t}f
 =\frac{\dd}{\dd t}\mathbb E_{\mu_t}f(T)
 =\mathbb E_{\mu_t}[f(T)s_t]
 =\mathbb E_{\lambda_t}[f\bar s_t].
\]
This identifies the target score and proves
\eqref{eq:YMA23-target-score}.  Differentiating the target amplitude gives
$\partial_t\sqrt{p_t}=\bar s_t\sqrt{p_t}/2$, proving
\eqref{eq:YMA23-target-speed}.  The chord between the two unit amplitudes is no
longer than the path length.  Conditional expectation is an $L^2$ contraction,
which proves the second inequality.
\end{proof}

\begin{corollary}[Finite target-speed quadrature]
\label{cor:YMA23-target-speed-quadrature}
For a partition $0=t_0<\cdots<t_J=1$ fixed before evaluation, suppose
\begin{equation}
 \mathbb E_{\lambda_t}\bar s_t^2\le E_j
 \qquad(t_{j-1}\le t\le t_j).
 \label{eq:YMA23-target-speed-panels}
\end{equation}
Then
\begin{equation}
 \boxed{
 \mathsf D(\lambda_0,\lambda_1)
 \le\frac14
 \left[\sum_{j=1}^J(t_j-t_{j-1})\sqrt{E_j}\right]^2.}
 \label{eq:YMA23-target-speed-distance}
\end{equation}
This value may be inserted into the module-error certificate.
\end{corollary}

\begin{proof}
Apply \eqref{eq:YMA23-target-length} and the piecewise upper bounds, then square.
\end{proof}

\begin{firewall}[A path upper bound is not an obstruction]
\label{fire:YMA23-path-not-obstruction}
The direct positive module distance is the target-native endpoint value.  The
Fisher length is a sufficient upper bound and may be loose because the target
amplitude path may turn.  Failure to make
\eqref{eq:YMA23-target-speed-distance} small is $\stopstatus$, not a negative
Wilson mode and not a boundary-transport obstruction without a direct endpoint
lower bound.
\end{firewall}

% =============================================================================
\section{The boundary-distance versus target-gain alternative}
\label{sec:YMA23-gain}
% =============================================================================

\begin{theorem}[Small target-law distance transports the card only after the gain is paid]
\label{thm:YMA23-gain-alternative}
Let $\rho_n>0$ be any predeclared physical normalization, let
$d_n=\mathsf D(\lambda_n,\lambda_0)$, and let
$\mathfrak b_n$ be the operator oscillation radius of the corresponding assembled
pair writer.  If
\begin{equation}
 \norm{\mathbb Q[\lambda_n]-\mathbb Q[\lambda_0]}_{\op}
 \ge\eta\rho_n
 \label{eq:YMA23-persistent-output}
\end{equation}
for some $\eta>0$, then
\begin{equation}
 \boxed{
 \frac{\mathfrak b_n}{\rho_n}
 \sqrt{d_n^{(2)}(4-d_n^{(2)})}
 \ge\eta,
 \qquad
 d_n^{(2)}=2d_n-\frac12d_n^2.}
 \label{eq:YMA23-gain-lower}
\end{equation}
Consequently, after passage to a subsequence, either
\begin{enumerate}[label=\textup{(G23.\arabic*)},leftmargin=5.2em]
\item $d_n\ge d_*>0$, a target-record boundary survivor; or
\item $d_n\to0$ and
      \begin{equation}
       \boxed{
       \frac{\mathfrak b_n}{\rho_n}\longrightarrow\infty,}
       \label{eq:YMA23-writer-gain-diverges}
      \end{equation}
      a target-writer amplification packet.
\end{enumerate}
Neither branch is by itself a negative Wilson mode; the sign is assigned only by
the assembled margin matrix.
\end{theorem}

\begin{proof}
Apply \Cref{thm:YMA23-matrix-transport} to the two product laws.  If $d_n$ has no
positive lower bound, pass to a subsequence with $d_n\to0$.  Then
$d_n^{(2)}\to0$, and \eqref{eq:YMA23-gain-lower} forces
\eqref{eq:YMA23-writer-gain-diverges}.
\end{proof}

\begin{countertheorem}[Vanishing module distance alone does not control an amplified target]
\label{cth:YMA23-distance-alone}
There are two-point probability laws $\mu_\varepsilon,\nu$ and bounded scalar
writers $A_\varepsilon$ such that
\begin{equation}
 \mathsf D(\mu_\varepsilon,\nu)\longrightarrow0,
 \qquad
 \mathbb E_{\mu_\varepsilon}A_\varepsilon
 -\mathbb E_\nu A_\varepsilon=1.
 \label{eq:YMA23-amplified-counterexample}
\end{equation}
Indeed, for $0<\varepsilon<1/2$ take
\begin{equation}
 \nu=(1/2,1/2),
 \qquad
 \mu_\varepsilon=(1/2+\varepsilon,1/2-\varepsilon),
 \qquad
 A_\varepsilon=\frac1{2\varepsilon}(1,-1).
 \label{eq:YMA23-counterexample-data}
\end{equation}
Then
\begin{equation}
 \boxed{
 \mathsf D(\mu_\varepsilon,\nu)
 =2-\sqrt{1+2\varepsilon}-\sqrt{1-2\varepsilon}
 =\varepsilon^2+O(\varepsilon^4),}
 \label{eq:YMA23-counterexample-distance}
\end{equation}
while $\norm{A_\varepsilon}_\infty=(2\varepsilon)^{-1}$.
\end{countertheorem}

\begin{proof}
The expectation under $\nu$ is zero, whereas
\[
 \mathbb E_{\mu_\varepsilon}A_\varepsilon
 =\frac{(1/2+\varepsilon)-(1/2-\varepsilon)}{2\varepsilon}=1.
\]
The affinity is
$\frac12(\sqrt{1+2\varepsilon}+\sqrt{1-2\varepsilon})$, which gives the exact
distance and its Taylor expansion.
\end{proof}

\begin{assessment}[Interpretation for YM--A]
\label{ass:YMA23-gain-meaning}
For a fixed compact finite card the writer radius is finite, so target-module
convergence transports the finite head.  Along a cofinal family, however, the
radius may grow with the response coupling, field normalization, head radius, or
support-birth stock.  The correct small parameter is the product in
\eqref{eq:YMA23-gain-lower}, normalized by the physical reference.  This is the
precise analogue of a source-to-state gain: the source distance and the target
amplification are not interchangeable rows.
\end{assessment}

% =============================================================================
\section{The complete finite-head, module, and Ward-tail certificate}
\label{sec:YMA23-complete-certificate}
% =============================================================================

The preceding statements are now specialized to the first actual retained Wilson
block.  Let $\widehat Q_0$ be a validated reference-boundary approximation to the
Ward-normalized finite-head margin, let $d_*$ be an outward upper bound for every
required target-module distance, and let $\tau_{q,R}$ be the uniform
Ward-relative remainder in \eqref{eq:YMA23-uniform-tail}.

\begin{theorem}[Target-native finite Wilson alternative]
\label{thm:YMA23-finite-alternative}
Assume the action, carrier, target map, and boundary family pass their exact domain
checks.  Define $d_*^{(2)}$, $\varepsilon_{\rm mod}$, and
$\varepsilon_{\rm tot}$ by
\eqref{eq:YMA23-module-error}--\eqref{eq:YMA23-total-radius}.  Then the following implications are valid.  A terminating audit returns the
first applicable branch in the displayed priority order; before that priority is
imposed, the structural survivor rows need not be mutually exclusive.
\begin{enumerate}[label=\textup{(C23.\arabic*)},leftmargin=5.2em]
\item \emph{All-boundary pass.}
      If \eqref{eq:YMA23-all-boundary-pass} holds, every DLR phase of the fixed
      action passes this nonzero block.
\item \emph{All-boundary pure-mode obstruction.}
      If \eqref{eq:YMA23-all-boundary-obstruction} holds, the minimum eigenvector
      of $\widehat Q_0$ is an explicit negative momentum--colour polarization in
      every DLR phase.
\item \emph{Target-record boundary survivor.}
      If a positive lower interval proves
      $\mathsf D(\lambda_{\eta_n},\lambda_{\eta_0})\ge d_*>0$ along an exhausting
      boundary sequence, the selected finite target record remains
      boundary-sensitive.  This is an auxiliary local-vacuum survivor, not yet a
      signed mode obstruction.
\item \emph{Target-writer amplification.}
      If the target distance tends to zero but a normalized card discrepancy
      persists, \eqref{eq:YMA23-writer-gain-diverges} returns the amplified
      finite-head writer stock.
\item \emph{Ward-tail survivor.}
      If the finite target card is controlled but no vanishing or finite outward
      bound is obtained for $\tau_{q,R}$, the programme returns the already typed
      connected positive-range or topological tail row.
\item \emph{Unresolved.}
      Every remaining valid overlap is $\stopstatus$.
\end{enumerate}
Malformed target records, action mismatches, a zero-character division, or failed
Hermiticity are input errors rather than theorem verdicts.
\end{theorem}

\begin{proof}
The first two items are \Cref{thm:YMA23-all-DLR-certificate}.  The third is the
definition of a nonvanishing target-module boundary difference, acquired by the
positive writer in \Cref{prop:YMA23-positive-distance-read}.  The fourth is
\Cref{thm:YMA23-gain-alternative}.  The fifth is the failure of the independent
remainder hypothesis \eqref{eq:YMA23-uniform-tail}; Versions~V1, V17--V20 already
classify its displacement and harmonic alternatives.  The last branch records the
failure of every strict one-sided interval without inventing a dual witness.
\end{proof}

\begin{corollary}[A direct selected-boundary obstruction needs no phase compiler]
\label{cor:YMA23-direct-selected-obstruction}
Suppose one boundary-selected exhaustion has a target-card limit in the sense of
\Cref{thm:YMA23-target-card-Cauchy}, its Ward tail converges, and the limiting
margin has a negative upper interval on a unit vector $v$.  Then $v$ is a negative
pure-mode witness for that selected local target.  No phase dispersion, complete
DLR simplex, or stochastic boundary coupling is needed for this selected-branch
claim.
\end{corollary}

\begin{proof}
Operator-norm convergence passes the strict negative inequality to the limit and,
eventually, to the selected finite cards.  The target-card theorem supplies exactly
the matrix expectation required for the stated branch and no stronger state claim.
\end{proof}

\begin{firewall}[No target inflation and no target deflation]
\label{fire:YMA23-target-scope}
The module card does not require a stochastic cutoff route, a complete DLR phase
simplex, or full conditional-shell quasilocality when the declared output is one
fixed finite Wilson matrix.  Conversely, it does not establish products, higher
jets, a local algebra, vacuum uniqueness, boundary chronology, adjacent-shell
transport, or local-vacuum contraction.  Those stronger targets reopen their own
native rows after the first Wilson block has been decided.
\end{firewall}

% =============================================================================
\section{Version V23 integrated conclusion and revised acquisition order}
\label{sec:YMA23-conclusion}
% =============================================================================

\begin{theorem}[Version V23 target-quotient module theorem]
\label{thm:YMA23-terminal}
Versions~V1--V22 are retained at their stated exact, finite, conditional, and open
scopes.  Version~V23 additionally proves:
\begin{enumerate}[label=\textup{(E23.\arabic*)},leftmargin=5.6em]
\item the complete finite-head Wilson Hessian is determined by the one-copy target
      record $(J_{q,R},K_{q,R})$ and one Hermitian two-copy writer;
\item the target pushforward contracts the canonical half-density/module distance;
\item one bounded nonnegative target-mixture writer acquires that distance;
\item product-law affinity gives the exact boundary-pair distance
      $d_1+d_2-d_1d_2/2$;
\item the operator oscillation radius gives the explicit factor
      $\sqrt{D(4-D)}$ transporting every Hermitian target expectation;
\item one reference boundary card, one target-module radius, and one uniform
      Ward-tail interval give a uniqueness-free all-DLR pass or an all-DLR
      pure-mode obstruction;
\item a summable adjacent target-module card constructs one selected matrix limit
      without reconstructing a complete state;
\item the Version~V22 exterior Hamiltonian modulus implies the target distance but
      is not required for this weaker target;
\item target-projected Fisher speed gives an alternative positive bridge-length
      upper bound;
\item a persistent card discrepancy yields the exact boundary-survival versus
      target-writer-amplification alternative.
\end{enumerate}
No item populates the required Wilson-specific numerical values.
\end{theorem}

\begin{proof}
Item~\textup{(E23.1)} is \Cref{lem:YMA23-target-reconstructs-card}.  Items
\textup{(E23.2)--(E23.3)} are
\Cref{prop:YMA23-positive-distance-read,lem:YMA23-data-processing}.  Items
\textup{(E23.4)--(E23.5)} are
\Cref{lem:YMA23-product-distance,thm:YMA23-matrix-transport}.  Item
\textup{(E23.6)} is \Cref{thm:YMA23-all-DLR-certificate}.  Item
\textup{(E23.7)} is \Cref{thm:YMA23-target-card-Cauchy}.  Item
\textup{(E23.8)} is \Cref{prop:YMA23-exterior-implies-module}.  Item
\textup{(E23.9)} is \Cref{thm:YMA23-target-Fisher-speed}.  Item
\textup{(E23.10)} is \Cref{thm:YMA23-gain-alternative}.
\end{proof}

\begin{assessment}[Effect of the related continuations]
\label{ass:YMA23-related-effect}
The boundary-sector continuation removes one target inflation: same-observable
Hilbert incidence is canonical and does not require a stochastic re-equilibration
kernel.  The fixed-datum and terminal-source continuations show why a source
distance cannot be passed through an unpriced response gain.  The
Einstein--Standard--Model continuation independently warns that positive control
marginals do not identify a same-history mixed physical gate.  The present target
record incorporates all three lessons: it uses the canonical positive module on
the exact fixed-action writer algebra, retains the target gain, and evaluates the
assembled pair writer rather than combining separately fitted marginals.
\end{assessment}

\begingroup\small
\begin{longtable}{>{\raggedright\arraybackslash}p{0.28\textwidth}
                  >{\raggedright\arraybackslash}p{0.19\textwidth}
                  >{\raggedright\arraybackslash}p{0.45\textwidth}}
\caption{Version V23 proof-status ledger.}\label{tab:YMA23-ledger}\\
\toprule
\textbf{Row}&\textbf{Status}&\textbf{Exact scope}\\
\midrule
\endfirsthead
\toprule
\textbf{Row}&\textbf{Status}&\textbf{Exact scope}\\
\midrule
\endhead
V1--V22 prefix
&Preserved exactly
&No inherited source, action, phase, carrier, or margin theorem is rewritten.\\[0.3em]
Same-observable module incidence
&Imported, not reproved
&Boundary-sector V22 supplies the unique positive half-density unitary; no
 stochastic coupling is charged for the target Gram.\\[0.3em]
Finite Wilson target record
&Proved exact
&The one-copy pair $(J,K)$ and one two-copy Hermitian writer reconstruct the
 complete finite-head direct-minus-covariance card.\\[0.3em]
Target data processing
&Proved exact
&Hellinger/module distance decreases under the physical target pushforward.\\[0.3em]
Positive target-distance Read
&Proved exact
&One writer in $[0,2]$ on the endpoint mixture has expectation equal to the
 target distance.\\[0.3em]
Product boundary distance
&Proved exact
&The two-copy distance is $d_1+d_2-d_1d_2/2$.\\[0.3em]
Target-writer transport
&Proved exact
&Operator expectation change is bounded by the oscillation radius times
 $\sqrt{D(4-D)}$.\\[0.3em]
All-boundary nonzero card
&Conditional terminating theorem
&A reference card, uniform target distance, and Ward tail give an all-DLR pass or
 common-vector obstruction.\\[0.3em]
Target-card limit
&Proved exact
&Summable adjacent module errors give a matrix limit without reconstructing a
 complete state.\\[0.3em]
Full exterior modulus
&Retained as sufficient stronger row
&Its Hamiltonian oscillation implies the target-module estimate; it is not the
 first requirement for one finite block.\\[0.3em]
Target-Fisher bridge
&Proved exact
&The projected action score is the target-law score and its amplitude length
 bounds the endpoint distance.\\[0.3em]
Boundary/gain alternative
&Proved quantitative
&A persistent output under vanishing target distance forces divergence of the
 normalized target-writer radius.\\[0.3em]
Selected Wilson value
&\stopstatus
&No target-module distance, reference boundary Hessian, response cofactor, or
 Ward-tail interval is numerically populated.\\[0.3em]
Higher jets, contraction, and mass
&Not inferred
&The target quotient does not create a complete local state, products, temporal
 transfer, or continuum Yang--Mills mass.\\
\bottomrule
\end{longtable}
\endgroup

\begin{openproblem}[V24: execute one target-module Wilson card]
\label{op:YMA23-V24}
Preserve every fixed response, action, gauge, topology, colour, carrier, reference,
and regulator convention.  Do not reopen a full coupling process, phase simplex,
or global Fourier-volume compiler before the following finite values are read.
\begin{enumerate}[label=\textup{(V24.\arabic*)},leftmargin=5.6em]
\item Load the two native response cofactors and the resulting fixed action.
\item Freeze one actual retained $q\ne0$ block and one finite head radius $R$.
      Construct the literal target record $(J_{q,R},K_{q,R})$ from that action.
\item Choose one predeclared reference boundary and evaluate its assembled
      Ward-normalized pair card with an outward Hermitian radius.
\item Acquire the positive target-module distance to the required boundary family,
      or populate a target-score bridge-length upper bound.  Retain the target-writer
      oscillation radius; do not report law closeness without its gain.
\item Acquire the uniform Ward-relative omitted tail on the same boundary family.
\item Apply \Cref{thm:YMA23-all-DLR-certificate}.  Return an all-boundary pass, the
      explicit common-vector pure-mode obstruction, or $\stopstatus$ at the first
      unresolved target distance, gain, or tail.
\item If only one selected boundary branch is intended, use
      \Cref{thm:YMA23-target-card-Cauchy,cor:YMA23-direct-selected-obstruction}
      and state that weaker semantics explicitly.
\item Open the complete DLR/local-algebra, Witten/boundary ancestry, higher-jet,
      adjacent-shell, and local-vacuum contraction rows only after this first block
      closes.
\end{enumerate}
A stochastic boundary coupling is required only when a common re-equilibration
history or chronology is itself claimed.
\end{openproblem}

\begin{assessment}[V23 termination]
\label{ass:YMA23-termination}
The first missing object has been reduced from a complete infrared phase compiler to
one positive target-law distance, one target-writer gain, one reference Wilson card,
and the already required signed tail.  This is strictly closer to the actual
Wilson-specific physical value and does not discard the stronger phase and
specification theorems when their stronger outputs are later needed.  None of the
four required numerical values is supplied by the present packet.  The selected
Wilson margin therefore remains $\stopstatus$, and no negative physical mode has
been produced.
\end{assessment}

\section*{Version V23 append-only continuation record}
\addcontentsline{toc}{section}{Version V23 append-only continuation record}
The complete Version~V22 source above is preserved byte for byte up to the point
immediately preceding its sole terminal document-closing command.  Its SHA--256
digest is
\begin{center}\scriptsize\ttfamily
e3fd23524e8b27233fd711ebe7d0bebf518eb8395792aa5510f1590be36f4789.
\end{center}
For Version~V24, preserve this accumulated source, append new material immediately
before the final document-closing command, and increment the filename to
\texttt{\_V24.tex}.  The target quotient is used only for the first finite Wilson
matrix.  Reopen the complete local-vacuum state and its products before claiming
higher jets or local-vacuum contraction.

% =============================================================================
% END APPENDED VERSION V23 MATERIAL
% =============================================================================


% =============================================================================
% START APPENDED VERSION V24 MATERIAL
% =============================================================================

\clearpage
\hypersetup{
  pdftitle={Renewal Yang--Mills Reduced-Fibre Wilson Shell Symbol and Local-Vacuum Programme V24},
  pdfsubject={Score-module sufficiency, direct curvature Reads, and a target-inflation correction}
}

\section*{Version V24: score-module sufficiency, direct curvature Reads, and a target-inflation correction}
\addcontentsline{toc}{section}{Version V24: score-module sufficiency and direct curvature Reads}

Version~V23 compressed the boundary problem to the pushforward law of the joint
finite-head record
\(
 (J_{q,R},K_{q,R})
\).
That law is sufficient for the Wilson card, but it is not target minimal.  The
physical Hessian uses the deterministic curvature only through
\(\mathbb E K_{q,R}\), and it uses the score only through its mean and covariance.
No mixed moment between $J_{q,R}$ and $K_{q,R}$ occurs.  A joint target-law distance
can therefore be maximal even when the complete finite-head Hessian is identical.

The related boundary programme supplies the correct smaller positive geometry: two
centered source families living under different laws are compared in the canonical
common half-density amplitude space.  One positive amplitude-difference writer
measures their score-module displacement; a stochastic coupling is not required
unless a common re-equilibration history is itself claimed.  The related
source-to-return programmes add the necessary caution: a small source displacement
must be multiplied by the actual response gain before a downstream value is called
small.

Version~V24 therefore replaces the joint target-law entrance by the exact packet
\begin{equation}
 \boxed{
 \left(
  \overline K_\eta,
  m_\eta,
  C_\eta
 \right)
 =\left(
  \mathbb E_\eta K_{q,R},
  \mathbb E_\eta J_{q,R},
  \operatorname{Cov}_\eta(J_{q,R})
 \right).}
 \label{eq:YMA24-minimal-moment-packet}
\end{equation}
For one selected boundary, these three finite quantities and the Ward-relative tail
already decide the physical block.  For transport from a reference boundary, the
curvature mean is acquired by one direct matrix Read, while the covariance is
transported by the positive score-module difference.  For an all-boundary claim,
the boundary score-mean diameter is retained as the additional negative dispersion
term.  The corrected order is
\begin{equation}
 \boxed{
 \begin{gathered}
 \text{one fixed response-matched action and one retained }q\ne0\text{ block}
 \\
 \Downarrow\\
 \text{direct curvature mean and centred score amplitude}
 \\
 \Downarrow\\
 \text{one selected-boundary Wilson card}
 \\
 \Downarrow\quad\text{only for boundary transport}\\
 \text{positive score-module difference, curvature shift, and score-mean spread}
 \\
 \Downarrow\\
 \text{Ward tail, then }\pass\;|\;\obstruction\;|\;\stopstatus.
 \end{gathered}}
 \label{eq:YMA24-order}
\end{equation}
The complete joint record remains admissible when another target actually uses its
mixed $J$--$K$ incidence.  It is no longer charged merely to evaluate the Wilson
Hessian.

% =============================================================================
\section{A maximal joint-target displacement with an unchanged Wilson card}
\label{sec:YMA24-target-inflation}
% =============================================================================

\begin{countertheorem}[Joint target-law distance can be maximally inflated]
\label{cth:YMA24-joint-target-inflation}
Let the target value space be $\{-1,+1\}^2$, with first coordinate $J$ and second
coordinate $K$.  Define
\begin{align}
 \lambda_+
 &=\frac12\delta_{(1,1)}+\frac12\delta_{(-1,-1)},
 \label{eq:YMA24-lambda-plus}\\
 \lambda_-
 &=\frac12\delta_{(1,-1)}+\frac12\delta_{(-1,1)}.
 \label{eq:YMA24-lambda-minus}
\end{align}
Then the two joint laws have disjoint supports, so their squared Hellinger distance
in the Version~V23 normalization is
\begin{equation}
 \boxed{\mathsf D(\lambda_+,\lambda_-)=2.}
 \label{eq:YMA24-max-Hellinger}
\end{equation}
Nevertheless, their score marginals agree, their deterministic curvature means
agree, and their physical scalar Hessians agree:
\begin{equation}
 \boxed{
 \mathbb E_+K=\mathbb E_-K=0,
 \qquad
 \operatorname{Var}_+(J)=\operatorname{Var}_-(J)=1,
 \qquad
 H_+=H_-=-1.}
 \label{eq:YMA24-same-card}
\end{equation}
The canonical centred score-module distance is exactly zero.
\end{countertheorem}

\begin{proof}
The supports displayed in \eqref{eq:YMA24-lambda-plus}--
\eqref{eq:YMA24-lambda-minus} are disjoint, which gives the maximal Hellinger
value.  Under either law, $J$ is uniform on $\{-1,+1\}$, hence has mean zero and
variance one.  Under either law, $K$ is also uniform on $\{-1,+1\}$ and therefore
has mean zero.  Thus
$H=\mathbb E K-\operatorname{Var}(J)=-1$ in both cases.  The two centred score
amplitudes are identical because the score marginal laws are identical.
\end{proof}

\begin{assessment}[Exact qualification of Version V23]
\label{ass:YMA24-V23-qualification}
The Version~V23 joint target quotient is a valid sufficient packet.  What is
withdrawn is its priority as the smallest target.  A distance on $(J,K)$ may pay for
correlation information which the Hessian never reads.  For the quadratic Wilson
card, the direct curvature mean and the centred score module are sufficient and, in
the sense of the preceding counterexample, strictly less inflated.
\end{assessment}

% =============================================================================
\section{The exact moment factorization of the Wilson pair card}
\label{sec:YMA24-moment-factorization}
% =============================================================================

Fix the response-matched action, the nonzero retained block, and the finite head
radius $R$ exactly as in Version~V23.  For every exterior or boundary condition
$\eta$, let
\begin{equation}
 J_\eta:=J_{q,R},
 \qquad
 K_\eta:=K_{q,R},
 \qquad
 \overline K_\eta:=\mathbb E_\eta K_\eta,
 \qquad
 m_\eta:=\mathbb E_\eta J_\eta,
 \label{eq:YMA24-boundary-moments}
\end{equation}
where the subscript on $J_\eta,K_\eta$ records the law and, when relevant, the
boundary-dependent local realization of the same fixed action.  Put
\begin{equation}
 C_\eta
 :=\mathbb E_\eta
 \bigl[(J_\eta-m_\eta)(J_\eta-m_\eta)^*\bigr]
 \succeq0.
 \label{eq:YMA24-score-covariance}
\end{equation}
Retain
\begin{equation}
 L_{q,R}^{\delta}=T_{q,R}+\delta S_{q,R}
 \label{eq:YMA24-deterministic-subtraction}
\end{equation}
from Version~V23.

\begin{theorem}[Target-minimal same-boundary and cross-boundary formulas]
\label{thm:YMA24-moment-factorization}
The same-boundary Ward-normalized finite-head margin is
\begin{equation}
 \boxed{
 \mathbb Q_\eta
 =\frac{
  \overline K_\eta-C_\eta-L_{q,R}^{\delta}}
  {\omega(q)}.}
 \label{eq:YMA24-same-boundary-formula}
\end{equation}
For two different boundary conditions,
\begin{equation}
 \boxed{
 \mathbb Q_{\eta_1,\eta_2}
 =\frac12\bigl(\mathbb Q_{\eta_1}+\mathbb Q_{\eta_2}\bigr)
 -\frac{
  (m_{\eta_1}-m_{\eta_2})
  (m_{\eta_1}-m_{\eta_2})^*}
  {2\omega(q)}.}
 \label{eq:YMA24-cross-boundary-factorization}
\end{equation}
Consequently, the complete finite-head boundary-pair card depends only on the
three moment packets \eqref{eq:YMA24-minimal-moment-packet}.  No mixed $J$--$K$
moment and no joint law of $(J,K)$ occurs.
\end{theorem}

\begin{proof}
The first identity is the same-law pair-difference identity of Version~V23.  For
independent copies with laws $\eta_1,\eta_2$,
\begin{align*}
 &\mathbb E
 \bigl[(J_{\eta_1}^{(1)}-J_{\eta_2}^{(2)})
       (J_{\eta_1}^{(1)}-J_{\eta_2}^{(2)})^*\bigr]\\
 &\qquad={}
 C_{\eta_1}+C_{\eta_2}
 +(m_{\eta_1}-m_{\eta_2})
  (m_{\eta_1}-m_{\eta_2})^*.
\end{align*}
Insert this into the Version~V23 cross-boundary pair writer and group the first two
terms as the average of \eqref{eq:YMA24-same-boundary-formula}.
\end{proof}

\begin{firewall}[Separate moment Reads still use one fixed action]
\label{fire:YMA24-same-action}
The factorization does not permit arbitrary marginals to be spliced together.  The
curvature, score mean, and score covariance must all come from the same fixed
action, field tangent, boundary condition, support convention, and topological
split.  Their joint statistical correlation is unnecessary because it is absent
from the displayed formula; their common physical loading remains mandatory.
\end{firewall}

% =============================================================================
\section{The canonical centred score module}
\label{sec:YMA24-score-module}
% =============================================================================

Consider two score-value laws $\lambda_0,\lambda_1$ on $\C^d$.  They may be the
pushforwards of two boundary-conditioned Wilson laws by their complete finite-head
score maps.  Choose any common dominating probability measure $m$ and write
\begin{equation}
 p_i=\frac{\dd\lambda_i}{\dd m},
 \qquad
 m_i=\int_{\C^d}x\,p_i(x)\,\dd m(x),
 \qquad i=0,1.
 \label{eq:YMA24-score-densities}
\end{equation}
Define the centred half-density amplitudes
\begin{equation}
 \Phi_i(x):=(x-m_i)\sqrt{p_i(x)}\in\C^d.
 \label{eq:YMA24-score-amplitudes}
\end{equation}
Their covariance matrices are
\begin{equation}
 \boxed{
 C_i=\int\Phi_i(x)\Phi_i(x)^*\,\dd m(x).}
 \label{eq:YMA24-covariance-amplitude}
\end{equation}

\begin{proposition}[Dominating-measure independence]
\label{prop:YMA24-dominating-independence}
The amplitude inner products, covariance matrices, amplitude differences, and all
Grams defined below are independent of the common dominating measure, up to the
canonical density-square-root unitary between the two amplitude spaces.
\end{proposition}

\begin{proof}
If $m'=r m$ with $r>0$, then $p_i'=p_i/r$ and
$\Phi_i'=r^{-1/2}\Phi_i$.  Multiplication by $r^{-1/2}$ is the canonical unitary
from $L^2(m)$ to $L^2(m')$ and preserves every amplitude inner product.
\end{proof}

Put
\begin{equation}
 E:=\Phi_1-\Phi_0,
 \qquad
 D_{1/0}:=\int EE^*\,\dd m\succeq0,
 \qquad
 B_{1/0}:=\int\Phi_0E^*\,\dd m.
 \label{eq:YMA24-module-blocks}
\end{equation}
The matrix $D_{1/0}$ is the unnormalized centred score-module distance.  For one
unit field polarization $v$,
\begin{equation}
 v^*D_{1/0}v
 =\int
 \left|
  v^*(x-m_1)\sqrt{p_1(x)}
  -v^*(x-m_0)\sqrt{p_0(x)}
 \right|^2\dd m(x).
 \label{eq:YMA24-directional-module-distance}
\end{equation}

\begin{theorem}[Positive module Gram and exact covariance transport]
\label{thm:YMA24-score-module-Gram}
The complete score-module Gram is
\begin{equation}
 \boxed{
 \mathbb G_{1/0}^{\rm sc}
 =\begin{pmatrix}
   C_0&B_{1/0}\\
   B_{1/0}^*&D_{1/0}
  \end{pmatrix}
 =\int
  \binom{\Phi_0}{E}
  \binom{\Phi_0}{E}^{\!*}\dd m
 \succeq0.}
 \label{eq:YMA24-score-module-Gram}
\end{equation}
The endpoint covariance satisfies the exact identity
\begin{equation}
 \boxed{
 C_1-C_0
 =B_{1/0}+B_{1/0}^*+D_{1/0}.}
 \label{eq:YMA24-covariance-transport}
\end{equation}
For every unit $v$,
\begin{equation}
 \boxed{
 \left|v^*(C_1-C_0)v\right|
 \le
 2\sqrt{
  (v^*C_0v)(v^*D_{1/0}v)}
 +v^*D_{1/0}v.}
 \label{eq:YMA24-directional-covariance-bound}
\end{equation}
In operator norm,
\begin{equation}
 \boxed{
 \|C_1-C_0\|_{\mathrm{op}}
 \le
 2\sqrt{
  \|C_0\|_{\mathrm{op}}
  \|D_{1/0}\|_{\mathrm{op}}}
 +\|D_{1/0}\|_{\mathrm{op}}.}
 \label{eq:YMA24-operator-covariance-bound}
\end{equation}
If the mixed block $B_{1/0}$ is directly acquired, the covariance transport is
exact.  If only the positive module distance is acquired, the last two displays
are the target-native sufficient bounds.
\end{theorem}

\begin{proof}
The first matrix is the Gram of the two amplitude columns and is therefore
positive.  Since $\Phi_1=\Phi_0+E$,
\[
 C_1
 =\int(\Phi_0+E)(\Phi_0+E)^*\,\dd m
 =C_0+B_{1/0}+B_{1/0}^*+D_{1/0}.
\]
For a unit vector $v$, Cauchy--Schwarz gives
\[
 |v^*B_{1/0}v|
 \le
 \left(\int|v^*\Phi_0|^2\right)^{1/2}
 \left(\int|v^*E|^2\right)^{1/2}
 =\sqrt{(v^*C_0v)(v^*D_{1/0}v)}.
\]
This proves \eqref{eq:YMA24-directional-covariance-bound}.  The operator estimate
follows from
$\|B_{1/0}\|_{\mathrm{op}}
 \le\|C_0\|_{\mathrm{op}}^{1/2}
      \|D_{1/0}\|_{\mathrm{op}}^{1/2}$.
\end{proof}

\begin{corollary}[Directional covariance support]
\label{cor:YMA24-covariance-support}
If $D_{1/0}v=0$, then $B_{1/0}v=0$ and
\begin{equation}
 \boxed{v^*C_1v=v^*C_0v.}
 \label{eq:YMA24-covariance-support}
\end{equation}
Thus a field polarization invisible to the positive score-module displacement has
exactly unchanged covariance \emph{quadratic form}.  No claim that
$C_1v=C_0v$ is made without the additional condition $B_{1/0}^*v=0$.
\end{corollary}

\begin{proof}
The condition $D_{1/0}v=0$ means that the amplitude-difference synthesis annihilates
$v$, so $B_{1/0}v=0$.  Hence
$v^*(B_{1/0}+B_{1/0}^*)v=0$.  Use
\eqref{eq:YMA24-covariance-transport} and $v^*D_{1/0}v=0$.
\end{proof}

% =============================================================================
\section{One positive writer acquires the score-module distance}
\label{sec:YMA24-positive-module-writer}
% =============================================================================

Let
\begin{equation}
 \overline\lambda:=\frac12(\lambda_0+\lambda_1)
 \label{eq:YMA24-score-mixture}
\end{equation}
be the score-value mixture.  On the set $p_0+p_1>0$, define the positive matrix
writer
\begin{equation}
 \boxed{
 \mathscr D_{1/0}^{\rm sc}(x)
 :=\frac{2E(x)E(x)^*}{p_0(x)+p_1(x)}
 \succeq0,}
 \label{eq:YMA24-positive-module-writer}
\end{equation}
with value zero on the null set where the denominator vanishes.

\begin{theorem}[Positive score-module acquisition]
\label{thm:YMA24-positive-module-writer}
The writer in \eqref{eq:YMA24-positive-module-writer} satisfies
\begin{equation}
 \boxed{
 \mathbb E_{\overline\lambda}
 \mathscr D_{1/0}^{\rm sc}
 =D_{1/0}.}
 \label{eq:YMA24-positive-module-expectation}
\end{equation}
It is independent of the common dominating measure.  Hence one positive
matrix-valued target-mixture Read populates the complete score-module distance.
For a scalarized axial block or one frozen polarization, it is one bounded
nonnegative scalar writer whenever the score is bounded.
\end{theorem}

\begin{proof}
The density of $\overline\lambda$ relative to $m$ is $(p_0+p_1)/2$.  Multiplication
by \eqref{eq:YMA24-positive-module-writer} cancels the denominator and leaves
$EE^*$.  Dominating-measure independence is
\Cref{prop:YMA24-dominating-independence}.
\end{proof}

\begin{firewall}[Module incidence is not a stochastic boundary route]
\label{fire:YMA24-module-not-Markov}
The common half-density amplitude space compares the same score-value observable
under two endpoint laws.  It does not assert a Markov transition, common
re-equilibration history, or pathwise coupling between the endpoints.  Such a
process is charged only when chronology or an occurring boundary evolution is part
of the conclusion.
\end{firewall}

% =============================================================================
\section{The curvature mean is a separate direct target}
\label{sec:YMA24-curvature-read}
% =============================================================================

The deterministic part of the Wilson head requires only
\begin{equation}
 \overline K_i=\mathbb E_iK_i.
 \label{eq:YMA24-curvature-means}
\end{equation}
It does not require the distribution of $K_i$ jointly with the score.

\begin{proposition}[Canonical direct curvature contrast]
\label{prop:YMA24-curvature-contrast}
Let $\kappa_0,\kappa_1$ be the laws of the Hermitian curvature values on their
common finite-dimensional value space.  Put
\begin{equation}
 \overline\kappa=\frac12(\kappa_0+\kappa_1),
 \qquad
 \chi_K
 =\frac12\left(
  \frac{\dd\kappa_1}{\dd\overline\kappa}
  -\frac{\dd\kappa_0}{\dd\overline\kappa}
 \right).
 \label{eq:YMA24-curvature-contrast}
\end{equation}
Then $|\chi_K|\le1$, $\mathbb E_{\overline\kappa}\chi_K=0$, and
\begin{equation}
 \boxed{
 \overline K_1-\overline K_0
 =2\mathbb E_{\overline\kappa}(\chi_KK).}
 \label{eq:YMA24-curvature-contrast-identity}
\end{equation}
This one mixed target Read, a direct pair of curvature expectations, or a certified
matrix-order sandwich is sufficient.  No score-module or joint $J$--$K$ law is
needed for the curvature row.
\end{proposition}

\begin{proof}
The mixture dominates both endpoint laws and their densities add to two.  Subtract
the two expectation formulas entrywise.
\end{proof}

\begin{assessment}[The two target coordinates are not two source births]
\label{ass:YMA24-two-coordinates}
The curvature mean and score covariance are two terms of the same physical Hessian.
Their separate target-minimal acquisition does not assign them independent
physical occurrence prices.  It only avoids measuring a $J$--$K$ correlation which
the Hessian formula discards.
\end{assessment}

% =============================================================================
\section{Exact and bounded reference-to-boundary Wilson transport}
\label{sec:YMA24-card-transport}
% =============================================================================

Let boundary $0$ be a predeclared reference and boundary $1$ another target.  Put
\begin{equation}
 \Delta K_{1/0}:=\overline K_1-\overline K_0.
 \label{eq:YMA24-Delta-K}
\end{equation}
Let $B_{1/0},D_{1/0}$ be the score-module blocks of
\eqref{eq:YMA24-module-blocks}.

\begin{theorem}[Exact target-minimal Wilson transport]
\label{thm:YMA24-exact-card-transport}
The same-boundary finite-head cards satisfy
\begin{equation}
 \boxed{
 \mathbb Q_1-\mathbb Q_0
 =\frac{
  \Delta K_{1/0}
  -B_{1/0}-B_{1/0}^*-D_{1/0}}
  {\omega(q)}.}
 \label{eq:YMA24-exact-card-transport}
\end{equation}
Thus one curvature difference and the complete score-module Gram give the exact
transport.  If only the positive distance is acquired and
\begin{equation}
 \|\Delta K_{1/0}\|_{\mathrm{op}}\le k,
 \qquad
 \|C_0\|_{\mathrm{op}}\le c,
 \qquad
 \|D_{1/0}\|_{\mathrm{op}}\le d,
 \label{eq:YMA24-transport-inputs}
\end{equation}
then
\begin{equation}
 \boxed{
 \|\mathbb Q_1-\mathbb Q_0\|_{\mathrm{op}}
 \le
 \frac{k+2\sqrt{cd}+d}{\omega(q)}.}
 \label{eq:YMA24-card-transport-bound}
\end{equation}
For a fixed unit polarization $v$, the corresponding directional interval uses
\begin{equation}
 \boxed{
 \left|
  v^*(C_1-C_0)v-v^*D_{1/0}v
 \right|
 \le
 2\sqrt{
  (v^*C_0v)(v^*D_{1/0}v)}.}
 \label{eq:YMA24-directional-transport-interval}
\end{equation}
\end{theorem}

\begin{proof}
Subtract \eqref{eq:YMA24-same-boundary-formula} at the two endpoints and use
\eqref{eq:YMA24-covariance-transport}.  The norm and directional bounds are
\Cref{thm:YMA24-score-module-Gram}.
\end{proof}

\begin{countertheorem}[The positive distance alone does not fix the covariance sign]
\label{cth:YMA24-distance-no-orientation}
There are scalar amplitude pairs with the same reference stock $C_0$ and the same
positive distance $D_{1/0}$ for which $C_1-C_0$ takes the two extremal values
\begin{equation}
 \boxed{
 \pm2\sqrt{C_0D_{1/0}}+D_{1/0}.}
 \label{eq:YMA24-distance-extrema}
\end{equation}
Hence the mixed block $B_{1/0}$ is independent orientation data.  The bound
\eqref{eq:YMA24-card-transport-bound} is sharp at the scalar Hilbert-space level.
\end{countertheorem}

\begin{proof}
Take nonzero collinear amplitude vectors $\Phi_0$ and
$E=\pm\sqrt{D_{1/0}/C_0}\,\Phi_0$.  Then
$B_{1/0}=\pm\sqrt{C_0D_{1/0}}$ and use
\eqref{eq:YMA24-covariance-transport}.
\end{proof}

% =============================================================================
\section{Selected-boundary and all-boundary finite certificates}
\label{sec:YMA24-certificates}
% =============================================================================

For one selected boundary $\eta_*$, let an outward head approximation and the
remaining Ward-relative tail satisfy
\begin{equation}
 \|\mathbb Q_{\eta_*}-\widehat Q_*\|_{\mathrm{op}}
 \le\varepsilon_*,
 \qquad
 \|\mathbb M_{\eta_*}-\mathbb Q_{\eta_*}\|_{\mathrm{op}}
 \le\tau_*.
 \label{eq:YMA24-selected-errors}
\end{equation}

\begin{theorem}[Direct selected-boundary verdict]
\label{thm:YMA24-selected-verdict}
With $r_*=\varepsilon_*+\tau_*$,
\begin{equation}
 \boxed{
 \lambda_{\min}(\widehat Q_*)-r_*>0
 \quad\Longrightarrow\quad\pass,}
 \label{eq:YMA24-selected-pass}
\end{equation}
whereas
\begin{equation}
 \boxed{
 \lambda_{\min}(\widehat Q_*)+r_*<0
 \quad\Longrightarrow\quad\obstruction.}
 \label{eq:YMA24-selected-obstruction}
\end{equation}
On the second branch, a minimum eigenvector of $\widehat Q_*$ is an explicit
negative physical polarization for the selected boundary.  Every remaining valid
interval is $\stopstatus$.  No boundary comparison or score-module distance is
required.
\end{theorem}

\begin{proof}
Apply Weyl's inequality to
$\|\mathbb M_{\eta_*}-\widehat Q_*\|_{\mathrm{op}}\le r_*$.  Evaluation on a
minimum eigenvector gives the obstruction witness.
\end{proof}

Now suppose an all-boundary target is intended.  Fix a reference boundary $0$ and
assume uniformly that
\begin{equation}
 \|\Delta K_{\eta/0}\|_{\mathrm{op}}\le k_*,
 \qquad
 \|D_{\eta/0}\|_{\mathrm{op}}\le d_*,
 \qquad
 \|C_0\|_{\mathrm{op}}\le c_*.
 \label{eq:YMA24-uniform-module-inputs}
\end{equation}
Put
\begin{equation}
 \varepsilon_{\rm sb}
 :=\frac{k_*+2\sqrt{c_*d_*}+d_*}{\omega(q)}.
 \label{eq:YMA24-same-boundary-radius}
\end{equation}
Retain the score-mean spread
\begin{equation}
 \rho_m
 :=\frac1{2\omega(q)}
 \sup_{\eta_1,\eta_2}
 \|m_{\eta_1}-m_{\eta_2}\|^2.
 \label{eq:YMA24-mean-spread}
\end{equation}
Let
\begin{equation}
 \|\mathbb Q_0-\widehat Q_0\|_{\mathrm{op}}\le\varepsilon_0,
 \qquad
 \sup_{\eta_1,\eta_2}
 \|\mathbb M_{\eta_1,\eta_2}-\mathbb Q_{\eta_1,\eta_2}\|_{\mathrm{op}}
 \le\tau_{q,R}.
 \label{eq:YMA24-reference-tail}
\end{equation}

\begin{theorem}[All-boundary score-module certificate]
\label{thm:YMA24-all-boundary-certificate}
If
\begin{equation}
 \boxed{
 \lambda_{\min}(\widehat Q_0)
 -\varepsilon_0-\varepsilon_{\rm sb}-\rho_m-\tau_{q,R}>0,}
 \label{eq:YMA24-all-boundary-pass}
\end{equation}
then every boundary-pair card and every DLR-state card of the fixed action passes
the retained nonzero block.

If
\begin{equation}
 \boxed{
 \lambda_{\min}(\widehat Q_0)
 +\varepsilon_0+\varepsilon_{\rm sb}+\tau_{q,R}<0,}
 \label{eq:YMA24-all-boundary-obstruction}
\end{equation}
then a minimum eigenvector $v_0$ of $\widehat Q_0$ is negative for every
boundary-pair card and hence for every DLR state.  The mean-spread term is not paid
on this branch because it enters \eqref{eq:YMA24-cross-boundary-factorization} with
a negative sign.
\end{theorem}

\begin{proof}
The transport bound gives
$\|\mathbb Q_\eta-\mathbb Q_0\|_{\mathrm{op}}
 \le\varepsilon_{\rm sb}$.  By
\eqref{eq:YMA24-cross-boundary-factorization},
\[
 \lambda_{\min}(\mathbb Q_{\eta_1,\eta_2})
 \ge
 \lambda_{\min}(\mathbb Q_0)-\varepsilon_{\rm sb}-\rho_m,
\]
which proves the passing branch after the reference and tail errors are included.
For the obstruction branch, evaluation on $v_0$ gives
\[
 v_0^*\mathbb Q_{\eta_1,\eta_2}v_0
 \le
 \lambda_{\min}(\widehat Q_0)
 +\varepsilon_0+\varepsilon_{\rm sb},
\]
because the boundary-mean term is negative semidefinite.  Add the tail radius.
The DLR statement is the exact boundary-pair disintegration of Version~V22.
\end{proof}

\begin{firewall}[A positive average and an all-boundary pass remain different]
\label{fire:YMA24-average-versus-all}
The score-module certificate is a uniform boundary theorem.  A positive card in one
mixed or periodic state does not imply \eqref{eq:YMA24-all-boundary-pass} unless the
curvature shift, score-module distance, mean spread, and Ward tail have the stated
uniform bounds.  Conversely, a negative selected-boundary card already obstructs
that selected target and is not postponed until an all-boundary theorem is proved.
\end{firewall}

% =============================================================================
\section{Source-specific bridge speed for the score module}
\label{sec:YMA24-score-speed}
% =============================================================================

The positive module distance can be read directly.  When an actual path of score
laws is already present, it can alternatively be bounded by its target-native
speed.

Let $t\mapsto\lambda_t$ be a differentiable path of score-value laws on a common
value space, with density $p_t$ relative to $m$.  Let
\begin{equation}
 F_t\in L^2(\lambda_t;\C^d),
 \qquad
 \mathbb E_{\lambda_t}F_t=0
 \label{eq:YMA24-centred-score-path}
\end{equation}
be the centred physical score family and put
\begin{equation}
 \Phi_t=F_t\sqrt{p_t},
 \qquad
 s_t=\partial_t\log p_t,
 \qquad
 G_t=\dot F_t+\frac12s_tF_t.
 \label{eq:YMA24-score-speed-data}
\end{equation}

\begin{theorem}[Exact score-amplitude speed]
\label{thm:YMA24-score-speed}
One has
\begin{equation}
 \boxed{
 \dot\Phi_t=G_t\sqrt{p_t},
 \qquad
 \|\dot\Phi_t\|_{L^2(m;\C^d)}^2
 =\mathbb E_{\lambda_t}\|G_t\|^2.}
 \label{eq:YMA24-score-speed}
\end{equation}
If $D_{1/0}$ is the endpoint module distance, then
\begin{equation}
 \boxed{
 \|D_{1/0}\|_{\mathrm{op}}^{1/2}
 \le
 \int_0^1
 \left(
  \mathbb E_{\lambda_t}\|G_t\|^2
 \right)^{1/2}\dd t.}
 \label{eq:YMA24-score-length}
\end{equation}
A predeclared finite quadrature of upper speed panels therefore supplies a
sufficient module-distance bound.  Failure of the path-length bound to be small is
$\stopstatus$, not an obstruction, because the source path may turn.
\end{theorem}

\begin{proof}
Differentiate $F_t\sqrt{p_t}$ and use
$\partial_t\sqrt{p_t}=s_t\sqrt{p_t}/2$.  This gives the speed identity.  Regard
$\Phi_t$ as the synthesis operator
$v\mapsto v^*\Phi_t$ from $\C^d$ to $L^2(m)$.  Its endpoint difference has Gram
$D_{1/0}$, hence operator norm $\|D_{1/0}\|_{\mathrm{op}}^{1/2}$.  The fundamental
theorem of calculus, the operator-norm triangle inequality, and the
Hilbert--Schmidt upper bound prove \eqref{eq:YMA24-score-length}.
\end{proof}

\begin{assessment}[Why this is smaller than the Version V23 target-Fisher path]
\label{ass:YMA24-score-speed-scope}
Only the centred score amplitude is transported.  Curvature values, unused
$J$--$K$ correlations, and complete conditional configurations are absent.  The
path is still optional: one positive endpoint mixture writer
\eqref{eq:YMA24-positive-module-writer} is the target-minimal acquisition.
\end{assessment}

% =============================================================================
\section{Adjacent-cutoff transport on a fixed physical score screen}
\label{sec:YMA24-cofinal}
% =============================================================================

Suppose adjacent response-matched shells have already been transported to one fixed
physical momentum--colour screen with common Ward normalization.  Let
\begin{equation}
 k_n\ge
 \|\overline K_{n+1}-\overline K_n\|_{\mathrm{op}},
 \qquad
 d_n\ge\|D_{n+1/n}\|_{\mathrm{op}},
 \qquad
 c_n\ge\|C_n\|_{\mathrm{op}},
 \label{eq:YMA24-adjacent-stocks}
\end{equation}
let $r_n$ bound the transported deterministic reference/topological increment, and
let $t_n$ bound the differentiated Ward-tail increment.

\begin{corollary}[Score-module Cauchy transport]
\label{cor:YMA24-score-module-Cauchy}
If
\begin{equation}
 \boxed{
 \sum_{n=0}^{\infty}
 \left[
  \frac{k_n+2\sqrt{c_nd_n}+d_n}{\omega_n}
  +r_n+t_n
 \right]
 <\infty,}
 \label{eq:YMA24-Cauchy-sum}
\end{equation}
then the transported Wilson margin matrices are Cauchy in operator norm on the
fixed screen.  A uniform positive lower margin larger than the remaining tail gives
a cofinal pass; a persistent negative upper interval gives a transported pure-mode
obstruction.  The theorem transports one selected physical matrix and does not
construct a complete local-vacuum state or path law.
\end{corollary}

\begin{proof}
Apply \eqref{eq:YMA24-card-transport-bound} at each adjacent step, add the
reference/topological and tail increments, and use the Cauchy criterion in the
finite-dimensional operator space.
\end{proof}

% =============================================================================
\section{Exact finite replay}
\label{sec:YMA24-replay}
% =============================================================================

The accompanying exact replay verifies eight independent finite layers.

\begin{enumerate}[label=\textup{(R24.\arabic*)},leftmargin=5.5em]
\item The disjoint joint target laws
      \eqref{eq:YMA24-lambda-plus}--\eqref{eq:YMA24-lambda-minus} have Hellinger
      distance two, score-module distance zero, and identical Hessian $-1$.
\item On a three-point score-value space, one explicit endpoint-mixture writer
      reproduces the exact centred score-module distance
      \begin{equation}
       \boxed{
       D_{1/0}=\frac{25}{48}-\frac{\sqrt2}{3}.}
       \label{eq:YMA24-replay-module-distance}
      \end{equation}
\item A rational $3\times2$ amplitude perturbation verifies
      \eqref{eq:YMA24-covariance-transport} exactly and checks the directional
      bound with a strict positive margin.
\item A scalar two-boundary packet verifies the exact mean/covariance decomposition
      \eqref{eq:YMA24-cross-boundary-factorization}, including the negative
      score-mean dispersion term.
\item A two-point Hermitian-curvature value law verifies the canonical contrast
      identity \eqref{eq:YMA24-curvature-contrast-identity} exactly.
\item For the two-point exponential family, the direct amplitude speed and the
      formula in \eqref{eq:YMA24-score-speed} both equal
      \begin{equation}
       \boxed{
       \frac{\cosh(3t)}{4\cosh(t)^5},}
       \label{eq:YMA24-replay-speed}
      \end{equation}
      with value $1/4$ at $t=0$.
\item A rational finite certificate has total radius $49/100$ and returns strict
      positive and negative endpoints $151/100$ and $-151/100$.
\item The model adjacent sequence $c_n=9$, $d_n=16^{-n}$ has exact summed
      covariance-increment majorant
      \begin{equation}
       \boxed{
       \sum_{n\ge1}
       \left(2\sqrt{c_nd_n}+d_n\right)
       =\frac{31}{15}.}
       \label{eq:YMA24-replay-Cauchy-sum}
      \end{equation}
\end{enumerate}

Every replay statement uses exact rational or symbolic-radical arithmetic.  The
replay does not evaluate a Wilson path integral, populate a selected root, or prove
a nonlinear Ward tail.

% =============================================================================
\section{Integrated V24 theorem and revised stopping boundary}
\label{sec:YMA24-integrated}
% =============================================================================

\begin{theorem}[Version V24 score-module reduction theorem]
\label{thm:YMA24-integrated}
Preserving Versions~V1--V23, Version~V24 proves the following.
\begin{enumerate}[label=\textup{(V24.\arabic*)},leftmargin=5.7em]
\item The joint law of $(J,K)$ is not target minimal for the Wilson Hessian; it can
      be maximally displaced while the physical card is unchanged.
\item The complete finite-head boundary-pair card is determined exactly by the
      curvature mean, score mean, and score covariance, with no $J$--$K$ mixed
      moment.
\item The centred score covariance is the Gram of a canonical half-density
      amplitude, independent of the dominating measure.
\item One positive score-module difference writer and one mixed module block give
      exact covariance transport; the positive distance alone gives the sharp
      square-root gain bound.
\item The deterministic curvature shift is one separate direct target Read and
      does not require a score module or joint target law.
\item One direct selected-boundary card plus its Ward tail gives a terminating
      physical verdict without any boundary comparison.
\item An all-boundary pass is certified by the reference card, curvature shift,
      score-module distance, boundary score-mean spread, and Ward tail.  A
      common-vector all-boundary obstruction does not pay the negative mean-spread
      debit.
\item An actual score-law bridge supplies an optional source-specific speed bound,
      and summable adjacent score-module increments transport one fixed physical
      matrix cofinally.
\end{enumerate}
No item populates the two native response cofactors, an actual selected Wilson
curvature mean, score covariance, score-module difference, or differentiated tail.
\end{theorem}

\begin{proof}
Item~\textup{(V24.1)} is
\Cref{cth:YMA24-joint-target-inflation}.  Item~\textup{(V24.2)} is
\Cref{thm:YMA24-moment-factorization}.  Items~\textup{(V24.3)--(V24.4)} are
\Cref{prop:YMA24-dominating-independence,thm:YMA24-score-module-Gram,thm:YMA24-positive-module-writer}.
Item~\textup{(V24.5)} is \Cref{prop:YMA24-curvature-contrast}.  Items
\textup{(V24.6)--(V24.7)} are
\Cref{thm:YMA24-selected-verdict,thm:YMA24-all-boundary-certificate}.  The final
item is \Cref{thm:YMA24-score-speed,cor:YMA24-score-module-Cauchy}.
\end{proof}

\begingroup\small
\begin{longtable}{>{\raggedright\arraybackslash}p{0.27\textwidth}
                  >{\raggedright\arraybackslash}p{0.18\textwidth}
                  >{\raggedright\arraybackslash}p{0.47\textwidth}}
\caption{Version V24 proof-status ledger.}\label{tab:YMA24-ledger}\\
\toprule
\textbf{Row}&\textbf{Status}&\textbf{Exact scope}\\
\midrule
\endfirsthead
\toprule
\textbf{Row}&\textbf{Status}&\textbf{Exact scope}\\
\midrule
\endhead
Joint-target inflation
& \textbf{excluded}
& Explicit maximal-distance counterexample with unchanged score covariance,
  curvature mean, and Hessian.\\[0.35em]
Moment factorization
& \textbf{proved} exact
& The finite head uses only $\overline K$, $m$, and $C$; the cross-boundary mean
  spread appears as one negative Gram.\\[0.35em]
Centred score module
& \textbf{proved} exact
& Canonical half-density amplitude, positive module Gram, exact covariance
  transport, and support theorem.\\[0.35em]
Positive module writer
& \textbf{proved} exact
& One endpoint-mixture matrix square has expectation $D_{1/0}$.\\[0.35em]
Curvature Read
& \textbf{proved} target minimal
& One direct mean difference, target-contrast Read, or order sandwich supplies the
  deterministic shift.\\[0.35em]
Selected-boundary verdict
& \textbf{proved} finite compiler
& One direct head and one Ward-tail radius return pass, an attained selected
  pure mode, or stop.\\[0.35em]
All-boundary verdict
& \textbf{proved} conditional compiler
& Reference card, curvature shift, score-module distance, mean spread, and tail
  give a uniform decision.\\[0.35em]
Score bridge and cutoff transport
& \textbf{proved} conditional
& Target-native speed panels and summable adjacent increments transport the
  selected covariance/card.\\[0.35em]
Native response cofactors
& \textbf{open}
& The selected-root plaquette and generated-action cofactors remain unpopulated.\\[0.35em]
Actual Wilson head
& \textbf{open}
& No supplied cylinder evaluates the direct curvature mean and complete score
  covariance on one retained physical block.\\[0.35em]
Differentiated Ward tail
& \textbf{open}
& The complete nonlinear connected, support-birth, and topological omitted tail
  remains a physical value.\\[0.35em]
Local-vacuum contraction and mass
& Not inferred
& Require higher jets, cutoff transport, conductance, reflection-positive
  continuum, and the remaining independent physical certificates.\\
\bottomrule
\end{longtable}
\endgroup

\begin{openproblem}[V25: execute the direct moment card]
\label{op:YMA24-V25}
Preserve every response, action, gauge, topology, carrier, colour, and reference
convention.  Do not reconstruct a joint target law unless another declared target
uses its $J$--$K$ correlation.  Proceed in this order.
\begin{enumerate}[label=\textup{(V25.\arabic*)},leftmargin=5.6em]
\item Acquire the two native response cofactors and load the one fixed
      response-matched action.  Terminate on the already classified response-chart
      branch if the plaquette cofactor vanishes.
\item Freeze one actual retained $q\ne0$ block, one selected physical boundary or
      local-vacuum branch, and one finite head radius $R$ before reading any sign.
\item Acquire directly, under that same action and branch,
      \(
       \overline K_{\eta_*},m_{\eta_*},C_{\eta_*}
      \),
      assemble \eqref{eq:YMA24-same-boundary-formula}, and populate the
      Ward-relative omitted tail.
\item Apply \Cref{thm:YMA24-selected-verdict}.  No module distance, reference law,
      boundary process, or phase simplex is charged for this selected card.
\item Only if all-boundary semantics are intended, choose one reference boundary,
      acquire the curvature shift and the positive score-module writer
      \eqref{eq:YMA24-positive-module-writer}, retain the boundary score-mean
      spread, and apply \Cref{thm:YMA24-all-boundary-certificate}.
\item Acquire the mixed module block only when the sufficient distance bound is too
      wide; do not infer its orientation from the diagonal stocks.
\item Return \pass, the explicit \obstruction, or \stopstatus
after the first unpopulated cofactor, moment, module, mean-spread, or tail row.
\item Reopen Witten ancestry and boundary chronology only after the first direct
      matrix block closes.  DLR reconstruction, higher jets, adjacent-shell
      transport, and local-vacuum contraction remain later targets.
\end{enumerate}
\end{openproblem}

\begin{assessment}[Version V24 termination]
\label{ass:YMA24-termination}
The first target has been reduced below the Version~V23 joint law.  It is now one
direct curvature expectation, one centred score covariance, and the already
required Ward-relative tail, all under one fixed response-matched action.  The
positive score-module distance is charged only for transport to another boundary.
No supplied note populates these Wilson-specific values.  The selected physical
symbol therefore remains $\stopstatus$, and no negative physical mode has been
produced.
\end{assessment}

\section*{Version V24 append-only continuation record}
\addcontentsline{toc}{section}{Version V24 append-only continuation record}
The complete Version~V23 source above is preserved byte for byte up to the point
immediately preceding its terminal document-closing command.  Its SHA--256 digest
is
\begin{center}\scriptsize\ttfamily
8a3e6ed20e5b1733fe8e3019d18cf4c5e9a20c1ecc93875becaede641a92183f.
\end{center}
For Version~V25, preserve this accumulated source, append new material immediately
before the final document-closing command, and increment the filename to
\texttt{\_V25.tex}.  Execute the direct moment card before reconstructing a joint
target law; use the positive score-module distance only for an actual boundary or
cutoff transport claim.

% =============================================================================
% END APPENDED VERSION V24 MATERIAL
% =============================================================================


% =============================================================================
% START APPENDED VERSION V25 MATERIAL
% =============================================================================

\clearpage
\hypersetup{
 pdftitle={Renewal Yang--Mills Reduced-Fibre Wilson Shell Symbol and Local-Vacuum Programme V25},
 pdfsubject={The exact strong-coupling Wilson surface coefficient, a source-connected polymer radius, and the finite-card margin crossing}
}

\pdfinfo{
 /Title (Renewal Yang--Mills Reduced-Fibre Wilson Shell Symbol and Local-Vacuum Programme V25)
 /Subject (The exact strong-coupling Wilson surface coefficient, a source-connected polymer radius, and the finite-card margin crossing)
}

\section{Version V25: the exact strong-coupling Wilson anchor, a source-connected polymer radius, and the finite-card margin crossing}
\label{sec:YMA25-direction}

Version~V24 is preserved above.  It reduced a boundary comparison to the
curvature mean and the centred score covariance.  That reduction is correct,
but another continuation through boundary modules would again postpone the
first literal Wilson coefficient.  The programme was opened to evaluate the
actual reduced-fibre Wilson symbol, not to build an indefinitely refined
comparison language.

The present continuation therefore goes upstream to the canonical point at
which the pure Wilson cylinder is exactly soluble coefficient by coefficient:
the strong-coupling anchor $\beta=0$.  It proves the following new facts.
\begin{enumerate}[label=\textup{(U25.\arabic*)},leftmargin=5.5em]
\item The first nonconstant term of the straight dyadic blocked Wilson action
      occurs at order four in the fine coupling.  Its coefficient is computed
      exactly from the four-plaquette fundamental surface.
\item In the normalization used throughout this programme, the induced coarse
      Wilson coefficient is
      \[
       \beta_{\rm c}(\beta)=\frac{\beta^4}{5832}+O(\beta^5),
      \]
      and the central axial transverse symbol is
      \[
       h_{M,\beta}^{\rm W}(q)
       =\frac{\beta^4}{17496}\,\omega(q)+O(\beta^5\omega(q)).
      \]
\item A source-connected polymer expansion converges uniformly in the volume
      for an explicit nonzero interval of $\beta$.  It gives an exponentially
      summable differentiated kernel and a volume-independent analytic
      remainder for the Ward quotient.
\item The pure Wilson Hessian is strictly positive on every nonzero central
      axial mode for sufficiently small positive coupling, but it is much
      softer than every fixed positive fraction of the flat reference
      $\beta s_0(q)$.
\item Combining this exact strong-coupling endpoint with the fixed-card
      large-$\beta$ asymptotic proved in Version~V3 gives an actual finite-card
      crossing of every target ratio $\delta\in(0,1)$.  No monotonicity or
      uniqueness is inferred.
\end{enumerate}

This is a Wilson-specific physical calculation.  It does not populate the
response-matched generated action, the selected weak-facing root, or the
nonlinear large-$\beta$ remainder.  It does, however, replace the statement
``the first Wilson value is absent'' by an exact coefficient, a convergent
local branch, and a finite-card crossing theorem.

\begin{firewall}[The strong-coupling anchor is not the weak-coupling conclusion]
\label{fire:YMA25-strong-not-weak}
The estimates below concern the pure Wilson shell before an independently
loaded generated or marginal interaction is added.  They prove a local
strong-coupling branch and determine its first coefficient.  A generated term
whose Taylor order is below five can alter the displayed remainder or even the
leading order.  No result below is promoted to the response-matched selected
root until that action is loaded on the same cylinder.  Likewise, a positive
Euclidean symbol at small $\beta$ is not an Osterwalder--Schrader mass gap.
\end{firewall}

% =============================================================================
\section{Positive character expansion of the Wilson plaquette weight}
\label{sec:YMA25-character}
% =============================================================================

Let $\chi_{\mathbf3}$ and $\chi_{\overline{\mathbf3}}$ be the fundamental and
antifundamental characters of $\SU(3)$.  In the normalization of the inherited
Wilson action, the field-dependent one-plaquette weight is
\begin{equation}
 w_\beta(U)
 :=\exp\left[
   \frac{\beta}{6}
   \bigl(\chi_{\mathbf3}(U)+\chi_{\overline{\mathbf3}}(U)\bigr)
  \right].
 \label{eq:YMA25-plaquette-weight}
\end{equation}
The field-independent factor $e^{-\beta}$ is kept outside this weight and has
no effect on a coarse-field Hessian.

\begin{theorem}[Entire positive character expansion]
\label{thm:YMA25-positive-character}
There are entire functions $a_\lambda(\beta)$, indexed by irreducible
$\SU(3)$ representations, such that
\begin{equation}
 \boxed{
 w_\beta(U)=\sum_\lambda a_\lambda(\beta)\chi_\lambda(U).}
 \label{eq:YMA25-character-expansion}
\end{equation}
For every real $\beta\ge0$,
\begin{equation}
 \boxed{a_\lambda(\beta)\ge0.}
 \label{eq:YMA25-character-positive}
\end{equation}
The two fundamental coefficients agree and satisfy
\begin{equation}
 \boxed{
 a_{\mathbf3}(\beta)=a_{\overline{\mathbf3}}(\beta)
 =\frac{\beta}{6}+\frac{\beta^2}{72}+O(\beta^3).}
 \label{eq:YMA25-fundamental-character-coefficient}
\end{equation}
\end{theorem}

\begin{proof}
Expand the exponential in \eqref{eq:YMA25-plaquette-weight}:
\[
 w_\beta
 =\sum_{m,n\ge0}
   \frac{(\beta/6)^{m+n}}{m!n!}
   \chi_{\mathbf3}^{\,m}\chi_{\overline{\mathbf3}}^{\,n}.
\]
A product of characters is the character of the corresponding tensor product.
Decomposition into irreducibles has nonnegative integral multiplicities.
Collecting each irreducible character therefore gives an entire coefficient
whose power-series coefficients are nonnegative, proving
\eqref{eq:YMA25-character-expansion}--\eqref{eq:YMA25-character-positive}.

At first order, the coefficient of either fundamental character is $\beta/6$.
At second order use
\[
 \mathbf3\otimes\mathbf3=\mathbf6\oplus\overline{\mathbf3},
 \qquad
 \overline{\mathbf3}\otimes\overline{\mathbf3}
 =\overline{\mathbf6}\oplus\mathbf3,
 \qquad
 \mathbf3\otimes\overline{\mathbf3}=\mathbf1\oplus\mathbf8.
\]
The prefactor of the quadratic term is $\beta^2/72$.  Exactly one copy of
$\mathbf3$ occurs in
$\overline{\mathbf3}\otimes\overline{\mathbf3}$ and exactly one copy of
$\overline{\mathbf3}$ occurs in $\mathbf3\otimes\mathbf3$.  This proves
\eqref{eq:YMA25-fundamental-character-coefficient}.
\end{proof}

\begin{remark}[What positivity does and does not prove]
\label{rem:YMA25-character-scope}
Positivity of the one-plaquette character coefficients is the correct entrance
to the surface expansion.  It does not by itself make the full reduced-fibre
field Hessian positive at every coupling: different connected surfaces and
score terms are assembled by the logarithm.  The theorem below extracts the
first coefficient before any such higher-order competition occurs.
\end{remark}

% =============================================================================
\section{The exact four-plaquette disk coefficient}
\label{sec:YMA25-four-surface}
% =============================================================================

We first record the only representation-theoretic integration identity needed
for the leading dyadic surface.  If $r$ is irreducible of dimension $d_r$,
Schur orthogonality gives
\begin{equation}
 \boxed{
 \int_{\SU(3)}
 \chi_r(AU)\chi_r(U^{-1}B)\,\dd U
 =\frac1{d_r}\chi_r(AB).}
 \label{eq:YMA25-character-convolution}
\end{equation}

\begin{lemma}[Oriented disk gluing]
\label{lem:YMA25-disk-gluing}
Let $\mathscr D$ be a simply connected oriented two-dimensional plaquette disk
with $A$ elementary plaquettes and fixed ordered boundary holonomy $W_{\partial
\mathscr D}$.  Put the same irreducible character $\chi_r$ on every plaquette,
with the orientation induced by the disk, and integrate all interior links with
normalized Haar measure.  Then
\begin{equation}
 \boxed{
 \int_{\mathrm{int}(\mathscr D)}
 \prod_{p\subset\mathscr D}\chi_r(U_p)
 \prod_{e\in\mathrm{int}(\mathscr D)}\dd U_e
 =d_r^{\,1-A}\chi_r(W_{\partial\mathscr D}).}
 \label{eq:YMA25-disk-formula}
\end{equation}
In particular, for the $2\times2$ fine disk tiling one coarse plaquette and
$r=\mathbf3$,
\begin{equation}
 \boxed{
 \int\prod_{p\subset P}\chi_{\mathbf3}(U_p)\,\dd U_{\rm int}
 =\frac1{27}\chi_{\mathbf3}(W_P).}
 \label{eq:YMA25-four-plaquette-gluing}
\end{equation}
The conjugate orientation gives the corresponding
$\overline{\mathbf3}$ formula.
\end{lemma}

\begin{proof}
Fix an axial gauge on the disk: choose a rooted spanning tree in its one-skeleton
which contains every interior vertex and all boundary vertices except for the
single boundary holonomy variable.  Haar invariance makes the change from tree
links to plaquette holonomies measure preserving.  The remaining disk integral
is an iterated convolution of $A$ copies of $\chi_r$, with one fixed boundary
product.  Applying \eqref{eq:YMA25-character-convolution} successively removes
$A-1$ independent plaquette holonomies and contributes one factor $d_r^{-1}$
at each step.  The last character is evaluated on the ordered boundary
holonomy.  This gives $d_r^{1-A}\chi_r(W_{\partial\mathscr D})$.
For $A=4$ and $d_{\mathbf3}=3$ the factor is $3^{-3}=1/27$.
\end{proof}

\begin{lemma}[Minimum dyadic coarse surface]
\label{lem:YMA25-minimum-surface}
For the straight two-link block map, a nonconstant contractible coarse Wilson
word requires at least four fine plaquette factors in the strong-coupling
expansion.  At total degree four, the only one-sheet contribution to an
oriented elementary coarse plaquette is the $2\times2$ tiled disk of
\Cref{lem:YMA25-disk-gluing}, with one fundamental character on each tile and
one common orientation.
\end{lemma}

\begin{proof}
A term in the character expansion is a finite fine-plaquette flux complex.
After every integrated fine link is contracted, a nonconstant dependence on
the retained coarse links must have a nonempty closed coarse boundary.  Under
the straight block map the smallest nontrivial contractible coarse boundary is
an elementary coarse plaquette, and its minimum fine filling has area four.
Thus total degree below four cannot carry nonconstant coarse boundary flux.

At degree four, any character of tensor order at least two on one tile leaves
at most three occupied tiles and cannot fill the coarse boundary.  Hence each
of the four tiles carries a first-order fundamental or antifundamental
character.  An internal edge is incident with two tiles.  Haar orthogonality
annihilates an inconsistent orientation, while a consistent orientation glues
the four tiles into the single-sheet disk.  Epsilon contractions at this total
degree either close a determinant inside the fine complex, giving a
coarse-field-independent scalar, or leave nonzero boundary triality and hence
vanish.  The only nonconstant one-sheet terms are therefore the two conjugate
surfaces stated in the lemma.
\end{proof}

% =============================================================================
\section{The first exact blocked Wilson coefficient}
\label{sec:YMA25-blocked-coefficient}
% =============================================================================

Let $Z_{M,\beta}^{\rm W}(W)$ be the exact reduced-fibre partition function for
one straight dyadic pure Wilson shell, with normalized product Haar fibre, and
let
\begin{equation}
 S_{M,\beta}^{\rm W}(W)=-\log Z_{M,\beta}^{\rm W}(W).
 \label{eq:YMA25-effective-action}
\end{equation}
Actions which differ by a $W$-independent scalar are identified in the next
theorem, since such a scalar has zero field jet.

\begin{theorem}[Exact order-four blocked Wilson action]
\label{thm:YMA25-order-four-action}
For every finite periodic dyadic card with no coarse plaquette identified with
itself by the period,
\begin{equation}
 \boxed{
 S_{M,\beta}^{\rm W}(W)
 =C_M(\beta)
 +\frac{\beta^4}{5832}
  \sum_{P\in\mathscr P_c}
  \left(1-\frac13\operatorname{Re}\Tr W_P\right)
 +O_M(\beta^5),}
 \label{eq:YMA25-blocked-action-leading}
\end{equation}
where $C_M$ is independent of $W$.  Equivalently, the induced fundamental
coarse Wilson coefficient is
\begin{equation}
 \boxed{
 \beta_{\rm c}(\beta)
 =\frac{\beta^4}{5832}+O(\beta^5).}
 \label{eq:YMA25-beta-coarse}
\end{equation}
The coefficient $1/5832$ is exact.
\end{theorem}

\begin{proof}
By \Cref{thm:YMA25-positive-character}, the first-order fundamental coefficient
on one fine plaquette is $\beta/6$.  By
\Cref{lem:YMA25-minimum-surface}, the first nonconstant coarse-field term uses
one such coefficient on each of the four tiles of one elementary coarse
plaquette.  The disk integration in
\eqref{eq:YMA25-four-plaquette-gluing} therefore gives
\begin{equation}
 \frac{1}{27}\left(\frac{\beta}{6}\right)^4
 \left(
  \chi_{\mathbf3}(W_P)+\chi_{\overline{\mathbf3}}(W_P)
 \right)
 =\frac{\beta^4}{17496}\operatorname{Re}\Tr W_P
 \label{eq:YMA25-partition-leading}
\end{equation}
inside the partition function.  Terms of lower degree are independent of the
coarse field.  Consequently taking $-\log$ does not mix any lower-order
nonconstant term into degree four.  The degree-four field-dependent part of
the effective action is the negative of
\eqref{eq:YMA25-partition-leading}.

Since
\[
 \frac{\beta^4}{5832}
 \left(1-\frac13\operatorname{Re}\Tr W_P\right)
 =\frac{\beta^4}{5832}
  -\frac{\beta^4}{17496}\operatorname{Re}\Tr W_P,
\]
the first term is absorbed into $C_M(\beta)$ and the second is exactly the
field-dependent coefficient just obtained.  Summing the translated elementary
surfaces proves \eqref{eq:YMA25-blocked-action-leading}.
\end{proof}

\begin{corollary}[Exact leading central axial symbol]
\label{cor:YMA25-axial-leading}
Let $h_{M,\beta}^{\rm W}(q)$ be the scalar central axial transverse Hessian on
a normalized nonzero coarse momentum $q$.  Then
\begin{equation}
 \boxed{
 h_{M,\beta}^{\rm W}(q)
 =\frac{\beta^4}{17496}\,\omega(q)+O_M(\beta^5),
 \qquad
 \omega(q)=2(1-\cos q).}
 \label{eq:YMA25-axial-leading}
\end{equation}
Equivalently, the Ward quotient has the expansion
\begin{equation}
 \boxed{
 g_{M,\beta}^{\rm W}(q)
 :=\frac{h_{M,\beta}^{\rm W}(q)}{\omega(q)}
 =\frac{\beta^4}{17496}+O_M(\beta^5).}
 \label{eq:YMA25-Ward-leading}
\end{equation}
\end{corollary}

\begin{proof}
Version~V5 fixed the Wilson normalization
\[
 D^2\left[
  \beta_c\sum_P\left(1-\frac13\operatorname{Re}\Tr W_P\right)
 \right]_{W=I}
 =\frac{\beta_c}{3}\,d_1^*d_1.
\]
On an axial transverse unit mode, $d_1^*d_1$ has eigenvalue $\omega(q)$.
Substitute $\beta_c=\beta^4/5832+O(\beta^5)$ from
\eqref{eq:YMA25-beta-coarse}.  Division by the nonzero Ward factor gives the
second formula.
\end{proof}

\begin{corollary}[Strong-coupling scaling relative to the flat shell reference]
\label{cor:YMA25-reference-scaling}
For every fixed nonzero axial momentum,
\begin{equation}
 \boxed{
 \frac{h_{M,\beta}^{\rm W}(q)}{\beta s_0(q)}
 =\frac{\beta^3}{17496\,r_0(q)}+O_M(\beta^4),
 \qquad
 r_0(q)=\frac{s_0(q)}{\omega(q)}.}
 \label{eq:YMA25-reference-ratio}
\end{equation}
In particular,
\begin{equation}
 \boxed{
 \lim_{\beta\downarrow0}
 \frac{h_{M,\beta}^{\rm W}(q)}{\beta s_0(q)}=0.}
 \label{eq:YMA25-reference-zero}
\end{equation}
The pure Wilson shell is therefore positive but sub-reference at the
strong-coupling endpoint.
\end{corollary}

\begin{proof}
Use $s_0(q)=\omega(q)r_0(q)$ and
\eqref{eq:YMA25-Ward-leading}.  Version~V2 proves $r_0(q)>0$ on every nonzero
mode, so division is legitimate.
\end{proof}

% =============================================================================
\section{A volume-uniform source-connected polymer branch}
\label{sec:YMA25-polymer}
% =============================================================================

The preceding coefficient is finite-volume exact.  We now prove that it is the
first term of one volume-uniform local branch rather than an isolated formal
coefficient.

Let $\mathscr G_{\rm plaq}$ be the graph whose vertices are fine plaquettes and
whose edges join plaquettes sharing an integrated fine link.  In four
dimensions every link belongs to six plaquettes, so
\begin{equation}
 \boxed{\Delta_{\rm plaq}\le4(6-1)=20.}
 \label{eq:YMA25-plaquette-degree}
\end{equation}
Put
\begin{equation}
 u(\beta)=e^{|\beta|}-1,
 \qquad
 \beta_{\rm KP}
 :=\log\left(1+\frac{1}{2e^3\Delta_{\rm plaq}}\right).
 \label{eq:YMA25-KP-radius}
\end{equation}
For $\Delta_{\rm plaq}=20$ this gives
\(
 \beta_{\rm KP}=0.0012439027\ldots
\).

\begin{lemma}[Rooted plaquette-polymer bound]
\label{lem:YMA25-rooted-polymer}
Write
\begin{equation}
 e^{(\beta/3)\operatorname{Re}\Tr U_p}=1+f_{p,\beta}(U_p).
 \label{eq:YMA25-activity-factor}
\end{equation}
For complex $\beta$,
\begin{equation}
 \abs{f_{p,\beta}}\le u(\beta).
 \label{eq:YMA25-activity-bound}
\end{equation}
If $\gamma$ is a connected plaquette set, its activity obeys
\begin{equation}
 \abs{z_\gamma}\le u(\beta)^{|\gamma|}.
 \label{eq:YMA25-polymer-activity}
\end{equation}
The number of connected plaquette sets of cardinality $n$ containing one fixed
plaquette is at most
\begin{equation}
 \boxed{(e\Delta_{\rm plaq})^{n-1}.}
 \label{eq:YMA25-animal-count}
\end{equation}
For $|\beta|<\beta_{\rm KP}$ the hard-core polymer expansion of the logarithm
converges absolutely and uniformly in the volume and in the retained coarse
links.
\end{lemma}

\begin{proof}
Since $\abs{\operatorname{Re}\Tr U/3}\le1$,
\(
 |e^{\beta x}-1|\le e^{|\beta|}-1
\)
for the real number $x=\operatorname{Re}\Tr U/3$, proving
\eqref{eq:YMA25-activity-bound}.  A polymer activity is an integral of a
product of the corresponding $f_p$'s against normalized Haar measure, so
\eqref{eq:YMA25-polymer-activity} follows.

A connected set containing a root admits a rooted spanning tree.  The standard
exploration which records, for every newly discovered vertex, one of at most
$\Delta_{\rm plaq}$ parent incidences and one of the rooted tree shapes gives
the elementary lattice-animal bound
\eqref{eq:YMA25-animal-count}; the factor $e^{n-1}$ bounds the rooted-tree
shape count.

The subset expansion of the plaquette product factors over connected
components in $\mathscr G_{\rm plaq}$ and is therefore a hard-core polymer gas.
Two distinct connected components are incompatible when they overlap or are
adjacent in $\mathscr G_{\rm plaq}$.  The closed plaquette neighbourhood of a
polymer $\gamma$ has at most
$(\Delta_{\rm plaq}+1)|\gamma|$ vertices.  Hence the sum of activities
incompatible with $\gamma$, weighted by $e^{2|\gamma'|}$, is at most
\begin{align*}
 &(\Delta_{\rm plaq}+1)|\gamma|
 \sum_{n\ge1}
 (e\Delta_{\rm plaq})^{n-1}
 u(\beta)^n e^{2n}\\
 &\qquad\le
 |\gamma|\,
 \frac{(\Delta_{\rm plaq}+1)u(\beta)e^2}
      {1-e^3\Delta_{\rm plaq}u(\beta)}.
\end{align*}
When $|\beta|<\beta_{\rm KP}$, the denominator exceeds $1/2$ and the last
coefficient is bounded by
\(
 (\Delta_{\rm plaq}+1)/(e\Delta_{\rm plaq})<1
\).
The rooted tree expansion of the Mayer coefficients is therefore dominated by
a convergent geometric tree sum with weight $e^{2|\gamma|}$.  This proves absolute convergence of the logarithm, uniformly
in the number of plaquettes.  The activity bound is uniform in the fixed
coarse links, so the same is true of the boundary parameters.
\end{proof}

\begin{theorem}[Uniform analytic differentiated Wilson kernel]
\label{thm:YMA25-uniform-analytic-kernel}
On the pure Wilson branch and for $|\beta|<\beta_{\rm KP}$, the thermodynamic
or selected local-vacuum effective action has a unique connected expansion on
every finite coarse cylinder.  Its central axial Hessian has a translation
kernel $a_\beta(x)$ satisfying, for some explicit $C_\beta<\infty$ and
$c_\beta>0$ independent of the volume,
\begin{equation}
 \boxed{
 \norm{a_\beta(x)}_{\rm op}
 \le C_\beta e^{-c_\beta\,\dist_{\rm plaq}(0,x)}.}
 \label{eq:YMA25-exponential-kernel}
\end{equation}
Consequently every weighted displacement stock is finite, the Ward quotient
extends continuously to $q=0$, and
\begin{equation}
 g_\beta^{\rm W}(q)
 =\frac{\beta^4}{17496}+R_\beta(q)
 \label{eq:YMA25-uniform-Ward-expansion}
\end{equation}
with a volume-independent analytic remainder.

More explicitly, put
\begin{equation}
 r_*:=\frac12\beta_{\rm KP},
 \qquad
 \vartheta:=e^3\Delta_{\rm plaq}(e^{r_*}-1)<\frac12,
 \qquad b_\partial:=12,
 \label{eq:YMA25-marked-constants}
\end{equation}
where $b_\partial$ bounds the number of fine plaquettes incident with the two
fine constituents of one coarse link.  One admissible marked-cluster bound is
\begin{equation}
 B_*:=64b_\partial^2(\Delta_{\rm plaq}+1)^2
 e^{2r_*}(1+r_*)^2
 \frac{1+11\vartheta+11\vartheta^2+\vartheta^3}
      {(1-\vartheta)^5}.
 \label{eq:YMA25-Bstar}
\end{equation}
For real $0\le\beta\le r_*/2$,
\begin{equation}
 \boxed{
 \sup_{q\in[-\pi,\pi]}
 \abs{R_\beta(q)}
 \le C_*\beta^5,
 \qquad
 C_*:=\frac{2B_*}{r_*^5}.}
 \label{eq:YMA25-uniform-remainder}
\end{equation}
The constants are deliberately conservative and are not proposed as a useful
numerical strong-coupling threshold.
\end{theorem}

\begin{proof}
Differentiate the absolutely convergent connected polymer expansion with
respect to two retained coarse-link coordinates.  A first derivative can mark
only a plaquette incident with one of the two fine constituents of the coarse
link; there are at most $b_\partial=12$ such plaquettes.  A second derivative
either places a second mark or differentiates the first marked factor.  For
$|\beta|\le r_*$, the marked plaquette factors are bounded by
$e^{r_*}(r_*+r_*^2)$, while the unmarked factors retain the activity bound of
\Cref{lem:YMA25-rooted-polymer}.

A connected cluster contributing to displacement $x$ must contain a plaquette
path joining the two marked incidence sets.  The rooted tree estimate therefore
adds a geometric factor
\(
 \vartheta^{\dist_{\rm plaq}(0,x)}
\), proving \eqref{eq:YMA25-exponential-kernel}.  Weighting by the axial second
moment contributes at most a quadratic factor in the cluster cardinality.
The two mark choices and the second-derivative product rule contribute another
quadratic factor.  Hence the complete marked Ward-quotient sum is bounded by
\[
 64b_\partial^2(\Delta_{\rm plaq}+1)^2
 e^{2r_*}(1+r_*)^2
 \sum_{n\ge1}n^4\vartheta^{n-1}.
\]
The exact generating function
\[
 \sum_{n\ge1}n^4\vartheta^{n-1}
 =\frac{1+11\vartheta+11\vartheta^2+\vartheta^3}
       {(1-\vartheta)^5}
\]
gives \eqref{eq:YMA25-Bstar}.  The inherited Ward identity and the weighted
second-moment theorem of Version~V1 then extend the quotient through $q=0$.
The convergence is uniform on the complex disk $|\beta|\le r_*$, so the
quotient is analytic there and has modulus at most $B_*$.

By \Cref{cor:YMA25-axial-leading}, its coefficients of orders zero through
three vanish and its order-four coefficient is $1/17496$.  Cauchy's estimate
bounds every later coefficient by $B_*r_*^{-n}$.  For
$0\le\beta\le r_*/2$,
\[
 \sum_{n\ge5}B_*r_*^{-n}\beta^n
 \le\frac{B_*(\beta/r_*)^5}{1-\beta/r_*}
 \le\frac{2B_*}{r_*^5}\beta^5,
\]
which is \eqref{eq:YMA25-uniform-remainder}.
\end{proof}

\begin{corollary}[Uniform pure-Wilson strong-coupling positivity]
\label{cor:YMA25-uniform-positive}
Put
\begin{equation}
 c_0:=\frac1{17496},
 \qquad
 \beta_{\rm pos}
 :=\min\left\{
  \frac{r_*}{2},\frac{c_0}{2C_*}
 \right\}.
 \label{eq:YMA25-beta-positive}
\end{equation}
Then for every $0<\beta\le\beta_{\rm pos}$ and every nonzero central axial
momentum,
\begin{equation}
 \boxed{
 h_\beta^{\rm W}(q)
 \ge\frac{\beta^4}{34992}\,\omega(q)>0.}
 \label{eq:YMA25-small-beta-positive}
\end{equation}
The statement is uniform in the periodic volume and passes to the
strong-coupling local-vacuum limit.
\end{corollary}

\begin{proof}
Equation \eqref{eq:YMA25-uniform-remainder} and the definition of
$\beta_{\rm pos}$ give
\[
 g_\beta^{\rm W}(q)
 \ge c_0\beta^4-C_*\beta^5
 \ge\frac12c_0\beta^4.
\]
Multiply by $\omega(q)>0$.
\end{proof}

\begin{corollary}[Every fixed positive reference fraction fails near
strong coupling]
\label{cor:YMA25-positive-delta-fails}
Fix $\delta>0$.  Define
\begin{equation}
 \beta_{\rm soft}(\delta)
 :=\min\left\{
  \beta_{\rm pos},
  \left(\frac{5\delta}{42c_0}\right)^{1/3}
 \right\}.
 \label{eq:YMA25-beta-soft}
\end{equation}
For every $0<\beta<\beta_{\rm soft}(\delta)$ and every nonzero central axial
momentum,
\begin{equation}
 \boxed{
 h_\beta^{\rm W}(q)-\delta\beta s_0(q)<0.}
 \label{eq:YMA25-delta-obstruction}
\end{equation}
This is an attained obstruction to the declared $\delta$-fraction reference
floor.  It is not a negative Wilson Hessian, because
\eqref{eq:YMA25-small-beta-positive} still holds.
\end{corollary}

\begin{proof}
On the range in question,
\(
 g_\beta^{\rm W}(q)\le(3c_0/2)\beta^4
\).
Version~V2 proves the exact reference floor
\(
 r_0(q)=s_0(q)/\omega(q)\ge5/28
\).
Therefore
\[
 \frac{h_\beta^{\rm W}(q)-\delta\beta s_0(q)}{\omega(q)}
 \le\frac32c_0\beta^4-\frac{5\delta}{28}\beta.
\]
The second bound in \eqref{eq:YMA25-beta-soft} makes the right-hand side
strictly negative.
\end{proof}

% =============================================================================
\section{The finite-card margin must cross}
\label{sec:YMA25-crossing}
% =============================================================================

The strong-coupling coefficient can be combined with the fixed-card
large-coupling theorem already proved in Version~V3.  No new asymptotic input is
needed.

\begin{theorem}[Existence of a pure-Wilson finite-card margin crossing]
\label{thm:YMA25-finite-crossing}
Fix a finite periodic card $M$, one retained nonzero central axial momentum
$q$, and a target $\delta\in(0,1)$.  Define
\begin{equation}
 F_{M,q,\delta}(\beta)
 :=h_{M,\beta}^{\rm W}(q)-\delta\beta s_0(q).
 \label{eq:YMA25-margin-function}
\end{equation}
Then:
\begin{enumerate}[label=\textup{(C\arabic*)},leftmargin=4em]
\item $F_{M,q,\delta}$ is real analytic for every finite real $\beta>0$ and
      continuous at $\beta=0$;
\item $F_{M,q,\delta}(\beta)<0$ for all sufficiently small positive $\beta$;
\item $F_{M,q,\delta}(\beta)>0$ for all sufficiently large $\beta$;
\item consequently there is at least one
      \begin{equation}
       \boxed{\beta_{M,q,\delta}^{\rm cross}\in(0,\infty)}
       \label{eq:YMA25-crossing-exists}
      \end{equation}
      such that
      \(
       F_{M,q,\delta}(\beta_{M,q,\delta}^{\rm cross})=0
      \).
\end{enumerate}
No monotonicity, simplicity, or uniqueness of this crossing is asserted.
\end{theorem}

\begin{proof}
At finite volume the reduced fibre is compact and the Wilson integrand is an
entire function of $\beta$.  Its real partition function is strictly positive,
so differentiation and division give a real-analytic effective Hessian for
$\beta>0$.  Continuity at zero follows from dominated convergence or from the
convergent expansion above.  Item~\textup{(C2)} is
\Cref{cor:YMA25-positive-delta-fails}.

Version~V3 proves on every fixed finite card that
\[
 \frac{h_{M,\beta}^{\rm W}(q)}{\beta s_0(q)}\longrightarrow1
 \qquad(\beta\to\infty).
\]
Since $\delta<1$, the ratio is eventually larger than $\delta$, proving
item~\textup{(C3)}.  The intermediate value theorem gives
\eqref{eq:YMA25-crossing-exists}.
\end{proof}

\begin{corollary}[Finite axial-bank crossing]
\label{cor:YMA25-bank-crossing}
Fix one finite card $M$ and let $\mathcal Q_M$ be any finite nonzero central
axial bank.  Put
\begin{equation}
 \rho_M(\beta)
 :=\min_{q\in\mathcal Q_M}
 \frac{h_{M,\beta}^{\rm W}(q)}{\beta s_0(q)}.
 \label{eq:YMA25-bank-ratio}
\end{equation}
Then $\rho_M(\beta)\to0$ as $\beta\downarrow0$ and
$\rho_M(\beta)\to1$ as $\beta\to\infty$.  For every $\delta\in(0,1)$ the
bank therefore has at least one transition value at which
\(
 \rho_M(\beta)=\delta
\).
\end{corollary}

\begin{proof}
The bank is finite.  The strong-coupling expansion and its small-$\beta$
remainder are uniform over the bank, so the minimum ratio tends to zero.  The
fixed-card large-$\beta$ expansion of Version~V3 is also uniform over a finite
bank, so the minimum tends to one.  Continuity of a finite minimum of
continuous functions and the intermediate value theorem finish the proof.
\end{proof}

\begin{assessment}[What the crossing does not select]
\label{ass:YMA25-crossing-scope}
The crossing in \Cref{thm:YMA25-finite-crossing} is a crossing of the pure
Wilson shell margin relative to the flat Schur reference.  It is not silently
identified with the weak-facing Creutz root used to select the bare trajectory.
The two scalar equations may be compared only after their common action and
response incidence are physically loaded.  The theorem nevertheless proves
that the Wilson mode ratio itself is a nontrivial finite-coupling quantity and
cannot be replaced by either endpoint asymptotic.
\end{assessment}

% =============================================================================
\section{Fixed-loop strong-coupling response anchor}
\label{sec:YMA25-Creutz-anchor}
% =============================================================================

For completeness, the same surface coefficient gives a direct upstream check
on the response selector, while preserving the selector-versus-marginal
distinction of Version~V12.

\begin{theorem}[Fixed contractible loop and Creutz leading law]
\label{thm:YMA25-loop-leading}
Let $C$ be a fixed simply connected fine Wilson loop with minimum plaquette
area $A(C)$, and use the normalized fundamental writer
\(
 \mathcal O_C=\frac13\operatorname{Re}\Tr U_C
\).
On the pure Wilson strong-coupling branch,
\begin{equation}
 \boxed{
 \mathbb E_\beta\mathcal O_C
 =\left(\frac{\beta}{18}\right)^{A(C)}
  \bigl(1+O_C(\beta)\bigr).}
 \label{eq:YMA25-loop-leading}
\end{equation}
For a fixed Creutz quartet whose signed area second difference is one,
\begin{equation}
 \boxed{
 \chi(\beta)
 =-\log\frac{\beta}{18}+O(\beta),
 \qquad
 \chi'(\beta)=-\frac1\beta+O(1).}
 \label{eq:YMA25-Creutz-leading}
\end{equation}
Hence the bare-selector response is regular and strictly decreasing for all
sufficiently small positive $\beta$.
\end{theorem}

\begin{proof}
To obtain a nonzero numerator containing $\mathcal O_C$, the strong-coupling
expansion must supply an oppositely oriented fundamental sheet filling $C$.
The minimum sheet has $A(C)$ plaquettes.  Applying
\Cref{lem:YMA25-disk-gluing} and then the final character orthogonality on the
loop gives one factor $1/3$ per occupied plaquette.  Together with the
one-plaquette coefficient $\beta/6$, the leading value is
$(\beta/18)^{A(C)}$.  The connected polymer expansion controls the relative
remainder for a fixed loop.

Insert this formula into the signed four-loop logarithm.  The perimeter and
vacuum factors cancel, while the signed area second difference contributes one
copy of $-\log(\beta/18)$.  Differentiation of the convergent fixed-loop
expansion proves the derivative formula.
\end{proof}

\begin{corollary}[Large-target strong-coupling selector root]
\label{cor:YMA25-large-target-root}
For every sufficiently large response target $u$, the equation
\(
 \chi(\beta)=u
\)
has one strong-coupling solution satisfying
\begin{equation}
 \boxed{
 \beta(u)=18e^{-u}\bigl(1+O(e^{-u})\bigr),
 \qquad
 \chi'(\beta(u))\ne0.}
 \label{eq:YMA25-large-u-root}
\end{equation}
This conclusion concerns the bare selector derivative
$d_{\rm sel}$.  It does not identify the independent plaquette marginal
response $d_P$ of Version~V12.
\end{corollary}

\begin{proof}
Set $y=-\log(\beta/18)$.  Equation
\eqref{eq:YMA25-Creutz-leading} becomes $u=y+O(e^{-y})$.  For large $u$ the
map is a contraction on a fixed neighbourhood of $u$, giving
$y=u+O(e^{-u})$ and hence the displayed $\beta(u)$.  The derivative formula in
\eqref{eq:YMA25-Creutz-leading} is nonzero there.
\end{proof}

% =============================================================================
\section{Exact replay and proof-status ledger}
\label{sec:YMA25-replay}
% =============================================================================

The accompanying replay checks the new exact finite algebra rather than a
Wilson path-integral sample.  It verifies:
\begin{enumerate}[label=\textup{(R25.\arabic*)},leftmargin=5.5em]
\item the fundamental tensor-product decompositions through second order and
      the coefficient
      \(
       a_{\mathbf3}=\beta/6+\beta^2/72+O(\beta^3)
      \);
\item the disk-gluing sequence $3^{1-A}$ and the $2\times2$ factor $1/27$;
\item the exact induced coefficients
      \(
       \beta_c^{(4)}=1/5832
      \)
      and
      \(
       g^{(4)}=1/17496
      \);
\item the fixed-loop leading factor $(\beta/18)^A$ and the unit Creutz area
      second difference;
\item the plaquette dependency degree bound $20$, the displayed
      Koteck\'y--Preiss radius, and the marked generating function
      \[
       \sum_{n\ge1}n^4x^{n-1}
       =\frac{1+11x+11x^2+x^3}{(1-x)^5};
      \]
\item the algebra of the small-$\beta$ positive and positive-$\delta$ soft
      branches;
\item one finite continuous model with a ratio tending from zero to one and an
      attained intermediate crossing.
\end{enumerate}
The replay does not locate the physical weak-facing root, evaluate a generated
interaction, or prove a large-$\beta$ uniform remainder.

\begin{theorem}[Version V25 strong-coupling anchor theorem]
\label{thm:YMA25-integrated}
Preserving Versions~V1--V24, Version~V25 proves:
\begin{enumerate}[label=\textup{(V25.\arabic*)},leftmargin=5.8em]
\item the Wilson plaquette weight has an entire nonnegative character
      expansion for nonnegative coupling;
\item the fundamental four-plaquette disk integrates to
      $3^{-3}\chi_{\mathbf3}(W_P)$;
\item the first nonconstant straight-dyadic blocked action is the coarse
      Wilson plaquette action with exact coefficient $\beta^4/5832$;
\item the first central axial Ward quotient is
      $\beta^4/17496+O(\beta^5)$;
\item a volume-uniform source-connected polymer branch gives exponential
      differentiated tails and a constructive analytic remainder;
\item the pure Wilson symbol is positive on that branch but fails every fixed
      positive fraction of the $\beta s_0$ reference near $\beta=0$;
\item every fixed finite pure-Wilson card crosses every target ratio
      $\delta\in(0,1)$ between strong and weak coupling;
\item fixed-loop strong coupling gives the selector anchor
      $\chi(\beta)=-\log(\beta/18)+O(\beta)$ without identifying the separate
      plaquette marginal response.
\end{enumerate}
No item evaluates the response-matched selected-root card or yields a
continuum mass gap.
\end{theorem}

\begin{proof}
The first two items are
\Cref{thm:YMA25-positive-character,lem:YMA25-disk-gluing}.  The blocked action
and symbol are
\Cref{thm:YMA25-order-four-action,cor:YMA25-axial-leading}.  The uniform branch
is
\Cref{lem:YMA25-rooted-polymer,thm:YMA25-uniform-analytic-kernel}.  Its two
sign consequences are
\Cref{cor:YMA25-uniform-positive,cor:YMA25-positive-delta-fails}.  The crossing
is \Cref{thm:YMA25-finite-crossing,cor:YMA25-bank-crossing}, and the selector
anchor is \Cref{thm:YMA25-loop-leading,cor:YMA25-large-target-root}.
\end{proof}

\begingroup\small
\begin{longtable}{>{\raggedright\arraybackslash}p{0.27\textwidth}
                  >{\raggedright\arraybackslash}p{0.18\textwidth}
                  >{\raggedright\arraybackslash}p{0.47\textwidth}}
\caption{Version V25 proof-status ledger.}\label{tab:YMA25-ledger}\\
\toprule
\textbf{Row}&\textbf{Status}&\textbf{Exact scope}\\
\midrule
\endfirsthead
\toprule
\textbf{Row}&\textbf{Status}&\textbf{Exact scope}\\
\midrule
\endhead
Positive Wilson character coefficients
& \textbf{proved} exact
& Entire character expansion; every coefficient is nonnegative for real
  $\beta\ge0$.\\[0.35em]
Four-plaquette surface
& \textbf{proved} exact
& Fundamental disk factor $1/27$ and conjugate orientation.\\[0.35em]
Blocked coarse coefficient
& \textbf{populated} exact leading value
& $\beta_c=\beta^4/5832+O(\beta^5)$ for the straight dyadic pure-Wilson shell.\\[0.35em]
Axial Ward quotient
& \textbf{populated} exact leading value
& $g_\beta^{\rm W}=\beta^4/17496+O(\beta^5)$.\\[0.35em]
Source-connected strong-coupling branch
& \textbf{proved} quantitative
& Explicit polymer radius, exponential differentiated kernel, and uniform
  analytic remainder.\\[0.35em]
Zero-floor symbol
& \pass
& Strictly positive on every nonzero central axial mode for sufficiently small
  pure-Wilson coupling.\\[0.35em]
Positive reference fraction
& \obstruction
& Every fixed $\delta>0$ fails near $\beta=0$ on an attained mode, although the
  Hessian itself remains positive.\\[0.35em]
Finite-card $\delta$ crossing
& \textbf{proved} existence
& Strong-coupling softness and the V3 large-$\beta$ limit force at least one
  crossing for every $\delta\in(0,1)$; uniqueness is open.\\[0.35em]
Bare Creutz selector anchor
& \textbf{proved} fixed-loop asymptotic
& $\chi=-\log(\beta/18)+O(\beta)$ and nonzero selector slope near zero.\\[0.35em]
Response-matched selected-root symbol
& \textbf{open}
& Generated/marginal action and the physical selected root are not populated by
  the pure-Wilson anchor.\\[0.35em]
Weak-coupling uniform remainder
& \textbf{open}
& Version~V3 remains fixed-card asymptotic; a cutoff-uniform nonlinear remainder
  is not proved here.\\[0.35em]
Local-vacuum contraction and mass
& Not inferred
& Require the response-matched physical branch, higher jets, transport,
  conductance, reflection positivity, and continuum certificates.\\
\bottomrule
\end{longtable}
\endgroup

% =============================================================================
\section{Revised V26 acquisition}
\label{sec:YMA25-next}
% =============================================================================

\begin{openproblem}[V26: bracket the physical branch against the exact Wilson anchor]
\label{op:YMA25-V26}
Preserve the complete present source.  Do not append another boundary, phase,
module, or target-law reduction before using the coefficient now available.
Proceed in the following order.
\begin{enumerate}[label=\textup{(V26.\arabic*)},leftmargin=5.7em]
\item Freeze whether the intended card is the pure Wilson shell or the native
      response-matched shell.  On the latter branch, load the two response
      cofactors and the complete added action before using the pure-Wilson
      coefficient as a comparison.
\item On one literal finite card, compute the selected response root or bracket
      it against the strong-coupling selector asymptotic
      \eqref{eq:YMA25-Creutz-leading}.  Keep $d_{\rm sel}$ and the independent
      plaquette marginal response $d_P$ distinct.
\item Evaluate the first nontrivial correction to
      $g_\beta^{\rm W}=\beta^4/17496+O(\beta^5)$, or acquire one rigorous
      partition-ratio/direct-moment interval at a coupling between the
      strong-coupling and V3 large-$\beta$ regimes.
\item Use the crossing theorem to predeclare a bracket for one target
      $\delta\in(0,1)$ and isolate every crossing in that bracket.  Return its
      multiplicity rather than assuming monotonicity.
\item If the response-matched addition is present, acquire its direct action
      curvature and complete same-history score covariance on the same block;
      do not transport the pure-Wilson sign through normalized-law data alone.
\item Only after one physical crossing or selected-root block is populated,
      reopen the differentiated Ward tail and adjacent-shell transport.
\item Return \pass, an attained \obstruction, or \stopstatus at the first
      unpopulated action, root, coefficient, or remainder row.
\end{enumerate}
The exact order-four coefficient is now the upstream anchor.  A new abstract
comparison that does not sharpen the $O(\beta^5)$ term, locate a crossing, or
populate the selected-root card is not an advance of YM--A.
\end{openproblem}

\begin{assessment}[Version V25 termination]
\label{ass:YMA25-termination}
The pure Wilson shell now has one literal physical coefficient and one
volume-uniform convergent local branch.  Its zero-floor symbol passes, every
fixed positive reference fraction fails near strong coupling, and every finite
card crosses to the V3 weak-coupling regime.  The response-matched selected
root remains unpopulated, so the programme-level verdict is still
$\stopstatus$.  The next progress must be a correction coefficient, a root
bracket, or a direct selected-card interval—not another enlargement of the
comparison carrier.
\end{assessment}

\section*{Version V25 append-only continuation record}
\addcontentsline{toc}{section}{Version V25 append-only continuation record}
The complete Version~V24 source above is preserved byte for byte up to the
point immediately preceding its terminal document-closing command.  Its
SHA--256 digest is
\begin{center}\scriptsize\ttfamily
4553a46fb4a86359b14a6277f8212d672feda6eafa746b7c58f427118e392484.
\end{center}
For Version~V26, preserve this accumulated source, append new material
immediately before the final document-closing command, and increment the
filename to \texttt{\_V26.tex}.  Use the exact strong-coupling coefficient to
sharpen one physical Wilson interval before reopening boundary or ancestry
machinery.

% =============================================================================
% END APPENDED VERSION V25 MATERIAL
% =============================================================================


% =============================================================================
% START APPENDED VERSION V26 MATERIAL
% =============================================================================

\clearpage
\hypersetup{
 pdftitle={Renewal Yang--Mills Reduced-Fibre Wilson Shell Symbol and Local-Vacuum Programme V26},
 pdfsubject={The vacuum-normalized fundamental activity, exact fifth-order Wilson coefficient, and the unique strong-coupling margin branch}
}

\pdfinfo{
 /Title (Renewal Yang--Mills Reduced-Fibre Wilson Shell Symbol and Local-Vacuum Programme V26)
 /Subject (The vacuum-normalized fundamental activity, exact fifth-order Wilson coefficient, and the unique strong-coupling margin branch)
}

\section{Version V26: vacuum-normalized activity, the exact fifth-order Wilson coefficient, and the first unique margin branch}
\label{sec:YMA26-direction}

Version~V25 is preserved above.  It computed the first nonzero pure-Wilson
coarse coefficient and directed the next continuation either to the complete
order-five correction or to a rigorous crossing interval.  The present
continuation does both at the first scale where they can be done without
introducing another boundary, phase, or source-comparison carrier.

The upstream point is that the raw coefficient
$a_{\mathbf3}(\beta)$ is not itself the surface activity.  Every plaquette
weight contains a nonzero trivial-character coefficient.  The physically
relevant character activity after removal of the common vacuum factor is
\[
 u_{\mathbf3}(\beta)
 =\frac{a_{\mathbf3}(\beta)}{a_{\mathbf1}(\beta)}.
\]
At fifth order the vacuum factor cannot mix into the coarse boundary term, but
making this division first reveals an exact cubic cancellation and identifies
the first order at which vacuum normalization can matter.  A separate
triality-filling lemma then excludes every five-plaquette competing coarse
surface.  The result is the exact expansion
\[
 \boxed{
 g_{M,\beta}^{\rm W}(q)
 =\frac{\beta^4}{17496}
  +\frac{\beta^5}{52488}
  +O(\beta^6)}
\]
for the Ward quotient of every nonzero central axial block.

Combining that coefficient with the volume-uniform analytic disk of
Version~V25 gives a volume-independent sixth-order remainder, strict
monotonicity of the relative Wilson margin on one explicit strong-coupling
interval, and a unique simple crossing for every sufficiently small target
fraction.  This is a Wilson-specific physical coefficient and crossing
statement.  It does not populate the response-matched generated action or
identify the crossing with a weak-facing selector root.

\begin{firewall}[No promotion of a pure-Wilson Taylor coefficient]
\label{fire:YMA26-pure-scope}
Every result below is for the pure Wilson shell and the already fixed straight
dyadic block map.  A generated or marginal interaction can contribute below
order six and must be inserted before these coefficients are used on the
response-matched branch.  The exact fifth-order coefficient is an anchor for
that comparison, not a substitute for it.  Positivity of the Euclidean
quadratic symbol is not an Osterwalder--Schrader mass statement.
\end{firewall}

% =============================================================================
\section{Vacuum-normalized fundamental activity}
\label{sec:YMA26-normalized-activity}
% =============================================================================

Put
\begin{equation}
 x:=\frac{\beta}{6},
 \qquad
 V:=\mathbf3\oplus\overline{\mathbf3}.
 \label{eq:YMA26-x-V}
\end{equation}
The one-plaquette weight of Version~V25 is
\(
 w_\beta=\exp(x\chi_V)
\).
For an irreducible representation $r$, let
\(
 m_n(r)=[r:V^{\otimes n}]
\)
be its tensor multiplicity.  Then
\begin{equation}
 a_r(\beta)=\sum_{n\ge0}\frac{x^n}{n!}m_n(r).
 \label{eq:YMA26-character-multiplicity-series}
\end{equation}

\begin{lemma}[Low-order tensor multiplicities]
\label{lem:YMA26-low-multiplicities}
The trivial and fundamental multiplicities through tensor degree three are
\begin{equation}
 \boxed{
 \begin{array}{c|rrrr}
 n&0&1&2&3\\
 \hline
 m_n(\mathbf1)&1&0&2&2\\
 m_n(\mathbf3)&0&1&1&6\\
 m_n(\overline{\mathbf3})&0&1&1&6.
 \end{array}}
 \label{eq:YMA26-multiplicity-table}
\end{equation}
Consequently,
\begin{align}
 \boxed{
 a_{\mathbf1}(\beta)
 =1+x^2+\frac{x^3}{3}+O(x^4),}
 \label{eq:YMA26-a1-series}\\
 \boxed{
 a_{\mathbf3}(\beta)
 =a_{\overline{\mathbf3}}(\beta)
 =x+\frac{x^2}{2}+x^3+O(x^4).}
 \label{eq:YMA26-a3-series}
\end{align}
\end{lemma}

\begin{proof}
The degree-two decompositions are
\[
 \mathbf3\otimes\mathbf3
 =\mathbf6\oplus\overline{\mathbf3},
 \qquad
 \mathbf3\otimes\overline{\mathbf3}
 =\mathbf1\oplus\mathbf8,
 \qquad
 \overline{\mathbf3}\otimes\overline{\mathbf3}
 =\overline{\mathbf6}\oplus\mathbf3.
\]
Since both ordered mixed products occur in $V^{\otimes2}$, the trivial
multiplicity is two; the fundamental and antifundamental multiplicities are
one.

At degree three, one trivial representation occurs in
$\mathbf3^{\otimes3}$ and one in
$\overline{\mathbf3}^{\otimes3}$, giving $m_3(\mathbf1)=2$.
A copy of $\mathbf3$ can occur only in a word containing two fundamental and
one antifundamental factors.  For one fixed ordering,
\[
 (\mathbf3\otimes\mathbf3)\otimes\overline{\mathbf3}
 =(\mathbf6\oplus\overline{\mathbf3})
   \otimes\overline{\mathbf3}
\]
contains two copies of $\mathbf3$: one in
$\mathbf6\otimes\overline{\mathbf3}$ and one in
$\overline{\mathbf3}\otimes\overline{\mathbf3}$.  There are three positions
for the antifundamental factor, so $m_3(\mathbf3)=6$.  Conjugation gives the
last entry.  Substitution in
\eqref{eq:YMA26-character-multiplicity-series} proves the displayed series.
\end{proof}

\begin{definition}[Vacuum-normalized character activities]
\label{def:YMA26-normalized-activity}
Factor the trivial coefficient from every plaquette weight:
\begin{equation}
 w_\beta(U)
 =a_{\mathbf1}(\beta)
  \left(
   1+\sum_{r\ne\mathbf1}u_r(\beta)\chi_r(U)
  \right),
 \qquad
 u_r(\beta):=\frac{a_r(\beta)}{a_{\mathbf1}(\beta)}.
 \label{eq:YMA26-normalized-characters}
\end{equation}
The product of the factors $a_{\mathbf1}$ is independent of every retained
coarse link and is therefore absorbed into the scalar part of the effective
action.
\end{definition}

\begin{theorem}[Exact cubic cancellation in the fundamental activity]
\label{thm:YMA26-activity-cancellation}
The vacuum-normalized fundamental activity satisfies
\begin{equation}
 \boxed{
 u_{\mathbf3}(\beta)
 =u_{\overline{\mathbf3}}(\beta)
 =x+\frac{x^2}{2}+0\,x^3+O(x^4).}
 \label{eq:YMA26-u3-series}
\end{equation}
In particular,
\begin{equation}
 \boxed{
 u_{\mathbf3}(\beta)^4
 =x^4+2x^5+\frac32x^6+O(x^7).}
 \label{eq:YMA26-u3-fourth}
\end{equation}
The coefficient $3/2$ in the second formula is the contribution of the
minimal four-plaquette sheet only; it is not yet asserted to be the complete
order-six blocked coefficient.
\end{theorem}

\begin{proof}
From \eqref{eq:YMA26-a1-series},
\[
 a_{\mathbf1}^{-1}
 =1-x^2-\frac{x^3}{3}+O(x^4).
\]
Multiplication by \eqref{eq:YMA26-a3-series} gives
\[
 \left(x+\frac{x^2}{2}+x^3+O(x^4)\right)
 \left(1-x^2-\frac{x^3}{3}+O(x^4)\right)
 =x+\frac{x^2}{2}+O(x^4).
\]
The raw cubic fundamental coefficient is cancelled exactly by the first
vacuum-normalization correction.  Raising the resulting series to the fourth
power gives \eqref{eq:YMA26-u3-fourth}; the order-six coefficient is
$\binom42(1/2)^2=3/2$.
\end{proof}

\begin{assessment}[Why the normalization first matters]
\label{ass:YMA26-normalization-first}
The order-five coarse coefficient is unchanged if one works with the raw
$a_{\mathbf3}$ because $a_{\mathbf1}=1+O(\beta^2)$.  At order six that shortcut
fails.  The cancellation in \eqref{eq:YMA26-u3-series} must be performed
before local sheet terms are compared with genuinely connected order-six
spin-foam decorations.  Vacuum factors and connected source polymers are not
separate physical sources and are not charged twice.
\end{assessment}

% =============================================================================
\section{No competing five-plaquette triality filling}
\label{sec:YMA26-five-face}
% =============================================================================

The coefficient of $x^5$ could fail to be local to the four-plaquette disk
only if a distinct five-factor flux complex carried the same coarse boundary.
The next lemma rules this out before any numerical enumeration.

Let $D$ be the positively oriented $2\times2$ disk in the $(1,2)$ coordinate
plane whose boundary is one elementary coarse plaquette.  A first-order
fundamental or antifundamental plaquette factor is assigned coefficient
$+1$ or $-1$ in $\mathbb Z_3$.  Invariance of Haar integration under the
centre of $\SU(3)$ forces the signed incident trialities at every integrated
link to sum to zero modulo three.  Thus every nonzero first-order monomial
carries a $\mathbb Z_3$ plaquette chain with the prescribed retained boundary.

\begin{lemma}[Five-face triality filling gap]
\label{lem:YMA26-five-face-gap}
Assume the $2\times2$ coarse-plaquette neighbourhood is embedded in the
periodic card.  Let $C$ be a finite $\mathbb Z_3$ plaquette chain whose
nonzero coefficients are represented by $\pm1$, whose support has at most five
distinct elementary plaquettes, and whose boundary equals $\partial D$ in
$\mathbb Z_3$.  Then
\begin{equation}
 \boxed{C=D.}
 \label{eq:YMA26-five-face-unique}
\end{equation}
In particular, no chain with five distinct plaquettes carries that boundary.
\end{lemma}

\begin{proof}
Call a plaquette horizontal if it is parallel to the $(1,2)$ plane and
transverse otherwise.  Suppose first that the support contains a transverse
plaquette $p$.  The plaquette $p$ has two opposite edges parallel to at least
one transverse coordinate direction.  A horizontal plaquette contains neither
of these edges.  At each of the two edges, cancellation of the nonzero
coefficient of $p$ modulo three requires at least one further transverse
plaquette.  No distinct elementary plaquette shares both opposite edges of
$p$.  Hence the support contains at least three transverse plaquettes.

Every one of the eight fine edges of $\partial D$ must occur in the support.
A transverse plaquette contains at most one of these boundary edges, while a
horizontal elementary plaquette contains at most two.  With at least three
transverse plaquettes and at most two horizontal plaquettes, at most
$3+2\cdot2=7$ boundary edges can be covered.  This contradicts
$\partial C=\partial D$.

The support is therefore entirely horizontal.  Decompose it by transverse
coordinate layer.  In every layer different from the layer of $D$, the chain
has zero boundary.  The boundary map on finite-support two-chains in the
square lattice is injective: choose an occupied square extremal in the first
coordinate and then in the second; one of its outer edges cannot be cancelled.
Thus all off-plane layers vanish.  In the plane of $D$, the finite chain
$C-D$ has zero boundary, so the same extremal-square argument gives $C-D=0$.
This proves \eqref{eq:YMA26-five-face-unique}.
\end{proof}

\begin{theorem}[Complete degree-five boundary classification]
\label{thm:YMA26-degree-five-classification}
For one embedded elementary coarse plaquette, every nonconstant pure-Wilson
term of total Taylor degree five is obtained from the flat four-plaquette disk
by placing the quadratic fundamental coefficient on exactly one of its four
tiles and the linear fundamental coefficient on the other three.  The four
choices add to
\begin{equation}
 \boxed{
 4\left(\frac{x^2}{2}\right)x^3=2x^5.}
 \label{eq:YMA26-degree-five-weight}
\end{equation}
There is no other coarse-boundary word at degree five.
\end{theorem}

\begin{proof}
Version~V25 proves that a nonconstant coarse boundary requires at least four
occupied fine plaquettes and that the unique degree-four support is $D$ with a
common fundamental orientation.  At total degree five there are two
possibilities.

If five distinct plaquettes carry first-order characters, centre invariance
produces the triality chain of
\Cref{lem:YMA26-five-face-gap}, which must equal $D$ and therefore cannot have
five distinct support plaquettes.  Such a term vanishes.

Otherwise exactly four plaquettes are occupied.  Their support must be $D$,
and the degree partition is $(2,1,1,1)$.  Integration on each interior disk
edge forces the same irreducible boundary label on its two incident tiles.
To produce the fundamental coarse character, every tile must therefore lie in
the fundamental channel.  The degree-two contribution to that channel is
$x^2/2$ by \eqref{eq:YMA26-a3-series}; the three remaining tiles contribute
$x$.  There are four possible locations of the quadratic coefficient, giving
\eqref{eq:YMA26-degree-five-weight}.  Every other coarse loop has minimum fine
area greater than five, while a product of two nonconstant coarse words starts
at degree eight.
\end{proof}

\begin{firewall}[What the five-face lemma does not decide]
\label{fire:YMA26-six-open}
The theorem closes total degree five only.  At degree six, trivial-character
normalization, two quadratic insertions, cubic local tensor multiplicities,
and genuinely connected six-degree spin-foam decorations must be assembled.
Equation \eqref{eq:YMA26-u3-fourth} gives the minimal-sheet part of that row but
not the complete coefficient.
\end{firewall}

% =============================================================================
\section{The exact fifth-order blocked action and axial symbol}
\label{sec:YMA26-fifth-coefficient}
% =============================================================================

\begin{theorem}[Exact fifth-order induced Wilson coefficient]
\label{thm:YMA26-beta-coarse-fifth}
For every finite periodic straight-dyadic card with the local coarse-plaquette
neighbourhood embedded, the pure-Wilson effective action has the expansion
\begin{equation}
 \boxed{
 \begin{aligned}
 S_{M,\beta}^{\rm W}(W)
 ={}&C_M(\beta)\\
 &+\left(
   \frac{\beta^4}{5832}
   +\frac{\beta^5}{17496}
  \right)
  \sum_{P\in\mathscr P_c}
  \left(1-\frac13\operatorname{Re}\Tr W_P\right)
  +O_M(\beta^6),
 \end{aligned}}
 \label{eq:YMA26-blocked-action-fifth}
\end{equation}
where $C_M$ is independent of the retained coarse field.  Equivalently,
\begin{equation}
 \boxed{
 \beta_{\rm c}(\beta)
 =\frac{\beta^4}{5832}
  +\frac{\beta^5}{17496}
  +O_M(\beta^6).}
 \label{eq:YMA26-beta-coarse-fifth}
\end{equation}
Both displayed coefficients are exact.
\end{theorem}

\begin{proof}
Factor $a_{\mathbf1}(\beta)$ from every fine plaquette as in
\eqref{eq:YMA26-normalized-characters}.  The product of these factors is a
coarse-field-independent scalar.  The two oriented minimal disks contribute
\begin{equation}
 \frac{u_{\mathbf3}(\beta)^4}{27}
 \left(
  \chi_{\mathbf3}(W_P)
  +\chi_{\overline{\mathbf3}}(W_P)
 \right)
 \label{eq:YMA26-normalized-disk-density}
\end{equation}
inside the normalized reduced partition function.  By
\Cref{thm:YMA26-degree-five-classification}, this is the complete nonconstant
contribution through degree five.  Equation \eqref{eq:YMA26-u3-fourth} gives
\[
 \frac{x^4+2x^5}{27}
 \left(
  \chi_{\mathbf3}(W_P)
  +\chi_{\overline{\mathbf3}}(W_P)
 \right)
 +O(\beta^6).
\]
The first nonconstant partition term has degree four, so nonlinear terms in
$-\log Z$ begin at degree eight.  Hence the effective action has the negative
of the displayed field-dependent coefficient through degree five.

Since
\[
 \chi_{\mathbf3}(W)+\chi_{\overline{\mathbf3}}(W)
 =2\operatorname{Re}\Tr W
\]
and
\[
 \beta_{\rm c}
 \left(1-\frac13\operatorname{Re}\Tr W\right)
 =\beta_{\rm c}-\frac{\beta_{\rm c}}6
  \left(\chi_{\mathbf3}(W)+\chi_{\overline{\mathbf3}}(W)\right),
\]
the induced coefficient is
\[
 \beta_{\rm c}
 =\frac{6}{27}u_{\mathbf3}^4
 =\frac29\left(x^4+2x^5\right)+O(\beta^6).
\]
Substituting $x=\beta/6$ gives
\[
 \frac29x^4=\frac{\beta^4}{5832},
 \qquad
 \frac49x^5=\frac{\beta^5}{17496},
\]
which proves the theorem.
\end{proof}

\begin{corollary}[Exact fifth-order axial Ward quotient]
\label{cor:YMA26-axial-fifth}
For every nonzero central axial transverse mode,
\begin{align}
 \boxed{
 h_{M,\beta}^{\rm W}(q)
 =\left(
   \frac{\beta^4}{17496}
   +\frac{\beta^5}{52488}
  \right)\omega(q)
  +O_M(\beta^6\omega(q)),}
 \label{eq:YMA26-h-fifth}\\
 \boxed{
 g_{M,\beta}^{\rm W}(q)
 :=\frac{h_{M,\beta}^{\rm W}(q)}{\omega(q)}
 =\frac{\beta^4}{17496}
  +\frac{\beta^5}{52488}
  +O_M(\beta^6).}
 \label{eq:YMA26-g-fifth}
\end{align}
The fifth-order correction is positive and equals one third of the leading
coefficient after one power of $\beta$ is removed:
\begin{equation}
 \boxed{
 g_{M,\beta}^{\rm W}(q)
 =\frac{\beta^4}{17496}
  \left(1+\frac{\beta}{3}+O_M(\beta^2)\right).}
 \label{eq:YMA26-g-factorized}
\end{equation}
\end{corollary}

\begin{proof}
The correctly normalized coarse Wilson Hessian is
$\frac13\beta_{\rm c}d_1^*d_1$.  On the axial transverse block,
$d_1^*d_1$ has eigenvalue $\omega(q)$.  Substitute
\eqref{eq:YMA26-beta-coarse-fifth} and divide by the nonzero Ward factor.
\end{proof}

\begin{assessment}[Minimal-sheet order-six preview]
\label{ass:YMA26-order-six-preview}
The minimal disk alone would contribute
\begin{equation}
 \boxed{
 \beta_{\rm c}^{\rm disk}(\beta)
 =\frac{\beta^4}{5832}
  +\frac{\beta^5}{17496}
  +\frac{\beta^6}{139968}
  +O(\beta^7),}
 \label{eq:YMA26-beta-disk-six}
\end{equation}
so its axial Ward quotient has order-six coefficient $1/419904$.  This number
is a rigorously computed sector coefficient, not yet the complete order-six
Wilson coefficient.  The latter must add every connected six-degree source
foam before a sign or crossing correction is reported.
\end{assessment}

% =============================================================================
\section{A volume-uniform sixth-order remainder}
\label{sec:YMA26-uniform-remainder}
% =============================================================================

Retain the volume-independent analytic disk and constants $r_*,B_*$ from
Version~V25.  In particular, the pure-Wilson Ward quotient is analytic on
$|\beta|\le r_*$ and its Taylor coefficients are bounded in modulus by
$B_*r_*^{-n}$, uniformly in the finite periodic volume and the central axial
momentum.

Put
\begin{equation}
 c_4:=\frac1{17496},
 \qquad
 c_5:=\frac1{52488}=\frac{c_4}{3},
 \qquad
 C_6:=\frac{2B_*}{r_*^6}.
 \label{eq:YMA26-uniform-constants}
\end{equation}

\begin{theorem}[Uniform remainder after the exact fifth coefficient]
\label{thm:YMA26-uniform-sixth}
For every finite periodic volume, every nonzero central axial momentum, and
$0\le\beta\le r_*/2$,
\begin{equation}
 \boxed{
 \left|
  g_{M,\beta}^{\rm W}(q)
  -c_4\beta^4-c_5\beta^5
 \right|
 \le C_6\beta^6.}
 \label{eq:YMA26-uniform-sixth}
\end{equation}
Moreover,
\begin{equation}
 \boxed{
 \left|
 \frac{\dd}{\dd\beta}
 \left[
  \frac{
   g_{M,\beta}^{\rm W}(q)
   -c_4\beta^4-c_5\beta^5
  }{\beta}
 \right]
 \right|
 \le6C_6\beta^4.}
 \label{eq:YMA26-ratio-derivative-tail}
\end{equation}
Both estimates pass to the strong-coupling local-vacuum limit.
\end{theorem}

\begin{proof}
Write
\[
 g_{M,\beta}^{\rm W}(q)
 =\sum_{n\ge4}c_{n,M,q}\beta^n.
\]
The preceding section proves
$c_{4,M,q}=c_4$ and $c_{5,M,q}=c_5$ for every card.  Cauchy's estimate from
Version~V25 gives
$|c_{n,M,q}|\le B_*r_*^{-n}$ for $n\ge6$.  With
$x=\beta/r_*\le1/2$,
\[
 \sum_{n\ge6}|c_{n,M,q}|\beta^n
 \le B_*\sum_{n\ge6}x^n
 \le2B_*x^6=C_6\beta^6.
\]
For the derivative, termwise differentiation gives
\[
 \left|
 \sum_{n\ge6}(n-1)c_{n,M,q}\beta^{n-2}
 \right|
 \le\frac{B_*}{r_*^2}
 \sum_{n\ge6}(n-1)x^{n-2}.
\]
The exact sum is
\[
 \sum_{n\ge6}(n-1)x^{n-2}
 =\frac{x^4(5-4x)}{(1-x)^2}
 \le12x^4
 \qquad(0\le x\le1/2).
\]
This equals $6C_6\beta^4$.  Uniform convergence permits passage to the local
limit.
\end{proof}

\begin{corollary}[Sharpened uniform strong-coupling sandwich]
\label{cor:YMA26-sharpened-sandwich}
On the same interval,
\begin{equation}
 \boxed{
 \beta^4\left(c_4+c_5\beta-C_6\beta^2\right)
 \le g_{M,\beta}^{\rm W}(q)
 \le
 \beta^4\left(c_4+c_5\beta+C_6\beta^2\right).}
 \label{eq:YMA26-sharpened-sandwich}
\end{equation}
Hence the first correction strengthens, rather than weakens, the pure-Wilson
zero-floor on the strong-coupling branch.
\end{corollary}

% =============================================================================
\section{The unique small-target margin crossing}
\label{sec:YMA26-unique-crossing}
% =============================================================================

For $q\ne0$, define the relative Wilson ratio
\begin{equation}
 \mathscr R_{M,q}(\beta)
 :=\frac{h_{M,\beta}^{\rm W}(q)}{\beta s_0(q)}
 =\frac{g_{M,\beta}^{\rm W}(q)}{\beta r_0(q)},
 \qquad
 \mathscr R_{M,q}(0):=0.
 \label{eq:YMA26-relative-ratio}
\end{equation}
Recall the exact reference window
\begin{equation}
 \frac5{28}\le r_0(q)\le\frac13.
 \label{eq:YMA26-reference-window}
\end{equation}

Define
\begin{equation}
 \boxed{
 \beta_{\uparrow}
 :=\min\left\{
  \frac{r_*}{2},
  \sqrt{\frac{c_4}{4C_6}}
 \right\},
 \qquad
 \delta_{\uparrow}
 :=\frac94c_4\beta_{\uparrow}^3.}
 \label{eq:YMA26-monotone-window}
\end{equation}

\begin{theorem}[Uniform monotone strong-coupling branch]
\label{thm:YMA26-monotone-branch}
For every finite periodic volume and every nonzero central axial momentum,
\begin{equation}
 \boxed{
 \frac{\dd}{\dd\beta}
 \mathscr R_{M,q}(\beta)>0
 \qquad(0<\beta\le\beta_{\uparrow}).}
 \label{eq:YMA26-ratio-increasing}
\end{equation}
Moreover,
\begin{equation}
 \boxed{
 \mathscr R_{M,q}(\beta_{\uparrow})
 \ge\delta_{\uparrow}>0.}
 \label{eq:YMA26-ratio-end-floor}
\end{equation}
Consequently, for every
\begin{equation}
 0<\delta<\delta_{\uparrow},
 \label{eq:YMA26-small-delta-range}
\end{equation}
there is a unique
\begin{equation}
 \boxed{
 \beta_{M,q}^{\rm sc}(\delta)
 \in(0,\beta_{\uparrow})}
 \label{eq:YMA26-unique-beta}
\end{equation}
with
\begin{equation}
 h_{M,\beta}^{\rm W}(q)
 =\delta\beta s_0(q).
 \label{eq:YMA26-crossing-equation}
\end{equation}
This crossing is simple and upward:
\begin{equation}
 \boxed{
 \left.
 \frac{\dd}{\dd\beta}
 \bigl(
  h_{M,\beta}^{\rm W}(q)-\delta\beta s_0(q)
 \bigr)
 \right|_{\beta=\beta_{M,q}^{\rm sc}(\delta)}
 >0.}
 \label{eq:YMA26-crossing-upward}
\end{equation}
\end{theorem}

\begin{proof}
Write the remainder in \eqref{eq:YMA26-uniform-sixth} as $E(\beta)$.  Then
\[
 \mathscr R_{M,q}'(\beta)
 =\frac1{r_0(q)}
 \left[
  3c_4\beta^2+4c_5\beta^3
  +\frac{\dd}{\dd\beta}\frac{E(\beta)}{\beta}
 \right].
\]
By \eqref{eq:YMA26-ratio-derivative-tail} and the definition of
$\beta_{\uparrow}$,
\[
 6C_6\beta^4
 \le\frac32c_4\beta^2.
\]
Hence
\[
 \mathscr R_{M,q}'(\beta)
 \ge\frac1{r_0(q)}
 \left(
  \frac32c_4\beta^2+4c_5\beta^3
 \right)>0.
\]
At $\beta=\beta_{\uparrow}$, the same definition gives
$C_6\beta_{\uparrow}^2\le c_4/4$.  Therefore
\[
 \mathscr R_{M,q}(\beta_{\uparrow})
 \ge\frac1{r_0(q)}
 \left(
  \frac34c_4\beta_{\uparrow}^3
  +c_5\beta_{\uparrow}^4
 \right)
 \ge3\cdot\frac34c_4\beta_{\uparrow}^3
 =\delta_{\uparrow}.
\]
Strict monotonicity, continuity, and $\mathscr R_{M,q}(0)=0$ give the unique
solution for every \eqref{eq:YMA26-small-delta-range}.

Finally,
\[
 h_{M,\beta}^{\rm W}(q)-\delta\beta s_0(q)
 =\beta s_0(q)
  \left(\mathscr R_{M,q}(\beta)-\delta\right).
\]
At the crossing the derivative of the first factor multiplies zero, while the
remaining term is
$\beta s_0(q)\mathscr R_{M,q}'(\beta)>0$.
\end{proof}

\begin{theorem}[Two-term location of the strong-coupling crossing]
\label{thm:YMA26-crossing-asymptotic}
Put
\begin{equation}
 t_q(\delta)
 :=\left(17496\,r_0(q)\,\delta\right)^{1/3}.
 \label{eq:YMA26-t-delta}
\end{equation}
As $\delta\downarrow0$, the unique crossing in
\Cref{thm:YMA26-monotone-branch} satisfies
\begin{equation}
 \boxed{
 \beta_{M,q}^{\rm sc}(\delta)
 =t_q(\delta)
  -\frac19t_q(\delta)^2
  +O\bigl(t_q(\delta)^3\bigr).}
 \label{eq:YMA26-crossing-expansion}
\end{equation}
The remainder is uniform in the finite periodic volume and in every central
axial momentum because $r_0(q)$ lies in the compact interval
\eqref{eq:YMA26-reference-window} and the analytic remainder is uniform.
\end{theorem}

\begin{proof}
Equation \eqref{eq:YMA26-g-factorized} gives
\begin{equation}
 \mathscr R_{M,q}(\beta)
 =\frac{c_4}{r_0(q)}\beta^3
  \left(1+\frac{\beta}{3}+O(\beta^2)\right).
 \label{eq:YMA26-ratio-local-form}
\end{equation}
By definition of $t=t_q(\delta)$,
\(
 \delta=(c_4/r_0(q))t^3
\).
Write the solution as $\beta=ty(t)$.  After division by the nonzero factor
$(c_4/r_0)t^3$, the crossing equation becomes
\[
 y^3\left(1+\frac{ty}{3}+O(t^2)\right)=1.
\]
At $t=0$ the positive solution is $y=1$, and the derivative with respect to
$y$ is three.  The analytic implicit-function theorem gives a unique analytic
positive branch.  Set $y=1+at+O(t^2)$.  Comparison of the coefficient of $t$
gives
\(
 3a+1/3=0
\), so $a=-1/9$.  Multiplying by $t$ proves
\eqref{eq:YMA26-crossing-expansion}.  Uniformity follows from the common
analytic disk and the fixed reference window.
\end{proof}

\begin{corollary}[A predeclared outward crossing bracket]
\label{cor:YMA26-crossing-bracket}
There are constants $C_{\rm cr}>0$ and $\delta_{\rm cr}>0$, independent of the
finite periodic volume and the central axial momentum, such that for
$0<\delta<\delta_{\rm cr}$,
\begin{equation}
 \boxed{
 t_q-\frac19t_q^2-C_{\rm cr}t_q^3
 \le
 \beta_{M,q}^{\rm sc}(\delta)
 \le
 t_q-\frac19t_q^2+C_{\rm cr}t_q^3.}
 \label{eq:YMA26-crossing-outward}
\end{equation}
Both endpoints lie inside the monotone interval when $\delta_{\rm cr}$ is
chosen sufficiently small.  Bisection or exact Haar interval evaluation on
this predeclared bracket isolates the unique strong-coupling crossing without
assuming global monotonicity.
\end{corollary}

\begin{proof}
Uniform analyticity of the implicit branch on the compact reference window
bounds its third and higher Taylor coefficients on one common neighbourhood of
$t=0$.  Taylor's theorem gives the displayed interval after reducing that
neighbourhood if necessary.  The endpoint inclusion follows from
$t_q(\delta)\to0$ uniformly as $\delta\to0$.
\end{proof}

\begin{assessment}[What has and has not been isolated]
\label{ass:YMA26-crossing-scope}
The theorem isolates exactly one crossing attached to the strong-coupling
endpoint.  Other crossings can occur at larger coupling, and no global
monotonicity is asserted.  The isolated root is a root of the pure-Wilson
relative mode margin.  It is not identified with the Creutz weak-facing root
or with a response-matched marginal coefficient.  Those comparisons require
the common unnormalized action and response cofactors already specified in
Versions~V12--V16.
\end{assessment}

% =============================================================================
\section{Exact finite replay and status ledger}
\label{sec:YMA26-replay}
% =============================================================================

The accompanying exact replay performs the following independent algebraic
checks.
\begin{enumerate}[label=\textup{(R26.\arabic*)},leftmargin=5.7em]
\item It recursively tensors every $\SU(3)$ Dynkin representation with
      $\mathbf3\oplus\overline{\mathbf3}$ and recovers the multiplicity table
      \eqref{eq:YMA26-multiplicity-table}.
\item It reconstructs $a_{\mathbf1}$ and $a_{\mathbf3}$ through degree four
      and verifies the exact disappearance of the cubic term in
      $u_{\mathbf3}=a_{\mathbf3}/a_{\mathbf1}$.
\item It expands $u_{\mathbf3}^4$ and verifies the coefficients
      $1,2,3/2$ at powers $x^4,x^5,x^6$.
\item It converts the disk activity to the exact induced coefficients
      \[
       \beta_{\rm c}^{(4)}=\frac1{5832},
       \qquad
       \beta_{\rm c}^{(5)}=\frac1{17496},
      \]
      and the exact Ward coefficients
      \[
       c_4=\frac1{17496},
       \qquad
       c_5=\frac1{52488}.
      \]
\item It symbolically solves the relative-margin equation through second order
      in $t=(17496r_0\delta)^{1/3}$ and obtains the coefficient $-1/9$.
\item It checks the derivative-tail generating function
      \[
       \sum_{n\ge6}(n-1)x^{n-2}
       =\frac{x^4(5-4x)}{(1-x)^2}
      \]
      and its bound by $12x^4$ on $0\le x\le1/2$.
\end{enumerate}
The replay does not enumerate the complete order-six spin-foam row, evaluate a
Wilson path integral at finite coupling, or identify a selected response root.

\begin{theorem}[Version V26 fifth-order and local-crossing theorem]
\label{thm:YMA26-integrated}
Preserving Versions~V1--V25, the following statements are now proved.
\begin{enumerate}[label=\textup{(V26.\arabic*)},leftmargin=5.7em]
\item The vacuum-normalized fundamental activity has no cubic term.
\item No distinct five-plaquette triality complex carries the elementary
      coarse-plaquette boundary.
\item The induced coarse Wilson coefficient through order five is
      \[
       \beta_{\rm c}(\beta)
       =\frac{\beta^4}{5832}
        +\frac{\beta^5}{17496}
        +O(\beta^6).
      \]
\item The pure-Wilson axial Ward quotient through order five is
      \[
       g_{M,\beta}^{\rm W}(q)
       =\frac{\beta^4}{17496}
        +\frac{\beta^5}{52488}
        +O(\beta^6).
      \]
\item The remainder after these two coefficients is $O(\beta^6)$ uniformly in
      the finite volume and axial momentum on the V25 analytic disk.
\item The relative mode ratio is strictly increasing on one explicit
      volume-independent strong-coupling interval.
\item Every sufficiently small target fraction has a unique simple upward
      crossing on that interval, with the two-term location
      \eqref{eq:YMA26-crossing-expansion}.
\end{enumerate}
No item evaluates the complete order-six coefficient or the native
response-matched selected-root card.
\end{theorem}

\begin{proof}
Items~\textup{(V26.1)--(V26.2)} are
\Cref{thm:YMA26-activity-cancellation,lem:YMA26-five-face-gap,thm:YMA26-degree-five-classification}.
Items~\textup{(V26.3)--(V26.4)} are
\Cref{thm:YMA26-beta-coarse-fifth,cor:YMA26-axial-fifth}.  Item~\textup{(V26.5)}
is \Cref{thm:YMA26-uniform-sixth}.  Items~\textup{(V26.6)--(V26.7)} are
\Cref{thm:YMA26-monotone-branch,thm:YMA26-crossing-asymptotic,cor:YMA26-crossing-bracket}.
\end{proof}

\begingroup\small
\begin{longtable}{>{\raggedright\arraybackslash}p{0.27\textwidth}
                  >{\raggedright\arraybackslash}p{0.17\textwidth}
                  >{\raggedright\arraybackslash}p{0.48\textwidth}}
\caption{Version V26 proof-status ledger.}\label{tab:YMA26-ledger}\\
\toprule
\textbf{Row}&\textbf{Status}&\textbf{Exact scope}\\
\midrule
\endfirsthead
\toprule
\textbf{Row}&\textbf{Status}&\textbf{Exact scope}\\
\midrule
\endhead
V1--V25 prefix
& \pass
& Complete preterminal source retained byte-for-byte; no inherited theorem is
  rewritten.\\[0.35em]
Vacuum-normalized fundamental activity
& \pass
& Exact tensor multiplicities give
  $u_{\mathbf3}=x+x^2/2+0x^3+O(x^4)$.\\[0.35em]
Five-face triality audit
& \pass
& Every supported $\mathbb Z_3$ filling with at most five plaquettes is the
  flat $2\times2$ disk; no competing fifth-order surface survives.\\[0.35em]
Pure-Wilson fifth-order coefficient
& \pass
& $\beta_{\rm c}^{(5)}=1/17496$ and
  $g^{(5)}=1/52488$ are exact on every embedded finite card.\\[0.35em]
Volume-uniform sixth-order remainder
& \pass
& V25 Cauchy control gives
  $|g-c_4\beta^4-c_5\beta^5|\le C_6\beta^6$.\\[0.35em]
Small-target crossing branch
& \pass
& One explicit common interval is strictly monotone; every sufficiently small
  $\delta$ has one simple upward crossing with a two-term bracket.\\[0.35em]
Complete order-six coefficient
& \stopstatus
& Minimal-disk contribution is known, but every connected six-degree source
  foam has not yet been enumerated and assembled.\\[0.35em]
Native response-matched action
& \stopstatus
& The two physical response cofactors and the complete added action remain
  unpopulated.\\[0.35em]
Selected-root Wilson margin
& \stopstatus
& No supplied cylinder identifies the weak-facing root with the isolated
  pure-Wilson margin branch or evaluates the same-action card there.\\[0.35em]
OS/local-vacuum mass conclusion
& Not inferred
& Requires the response-matched head, differentiated tail, conductance,
  ancestry, cutoff transport, and reflection-positive continuum rows.\\
\bottomrule
\end{longtable}
\endgroup

% =============================================================================
\section{Revised V27 acquisition}
\label{sec:YMA26-next}
% =============================================================================

\begin{openproblem}[V27: compute the first genuinely connected correction]
\label{op:YMA26-V27}
Preserve the complete present source.  The order-five row and the local crossing
are closed, so do not reopen a boundary, phase, module, carrier, or target-law
compiler.  Proceed in this order.
\begin{enumerate}[label=\textup{(V27.\arabic*)},leftmargin=5.7em]
\item Enumerate every connected total-degree-six spin-foam term incident with
      one elementary coarse plaquette, including all $\SU(3)$ epsilon and
      higher-irrep channels.
\item Subtract the already computed minimal-sheet coefficient
      $1/419904$ in the axial Ward quotient and return the exact remaining
      connected coefficient, with its sign.
\item Use the resulting sixth-order polynomial and the V25 analytic remainder
      to sharpen the outward interval
      \eqref{eq:YMA26-crossing-outward} or to isolate one finite-coupling
      crossing by exact Haar interval arithmetic.
\item On the native response-matched branch, load the two cofactors and the
      complete added action before comparing its crossing with the pure-Wilson
      anchor.
\item If the selected weak-facing root lies in the strong-coupling bracket,
      evaluate the same-action direct curvature and complete score covariance
      there.  If it does not, record the disjoint intervals and stop.
\item Reopen the differentiated Ward tail or boundary/Witten ancestry only
      after this literal coefficient or selected-root card has been populated.
\item Return \pass, an attained \obstruction, or \stopstatus at the first
      unresolved coefficient, action, root, or remainder row.
\end{enumerate}
A new abstract comparison which neither computes the degree-six connected
coefficient nor sharpens a physical root interval is not an advance of YM--A.
\end{openproblem}

\begin{assessment}[Version V26 termination]
\label{ass:YMA26-termination}
The first requested correction is no longer open.  The pure Wilson shell now
has exact order-four and order-five axial coefficients, a volume-uniform
sixth-order remainder, and a unique rigorously bracketed strong-coupling
relative-margin crossing for every sufficiently small target fraction.  The
next mathematical unknown is a finite connected order-six coefficient, not a
missing comparison architecture.  The response-matched selected-root symbol
remains \stopstatus, and no negative pure-Wilson Hessian has been found.
\end{assessment}

\section*{Version V26 append-only continuation record}
\addcontentsline{toc}{section}{Version V26 append-only continuation record}
The complete Version~V25 source above is preserved byte for byte up to the
point immediately preceding its terminal document-closing command.  Its
SHA--256 digest is
\begin{center}\scriptsize\ttfamily
 d5d09b295a9bce1e0a1ac53ef3b85751551922ed00e5a180559d9678ea582d79.
\end{center}
For Version~V27, preserve this accumulated source, append new material
immediately before the final document-closing command, and increment the
filename to \texttt{\_V27.tex}.  Compute the connected degree-six coefficient
or a literal selected-root interval before reopening any comparison machinery.

% =============================================================================
% END APPENDED VERSION V26 MATERIAL
% =============================================================================


% =============================================================================
% START APPENDED VERSION V27 MATERIAL
% =============================================================================

\clearpage
\hypersetup{
 pdftitle={Renewal Yang--Mills Reduced-Fibre Wilson Shell Symbol and Local-Vacuum Programme V27},
 pdfsubject={The seven-face support gap, exact sixth and seventh Wilson coefficients, and a four-term strong-coupling crossing}
}

\pdfinfo{
 /Title (Renewal Yang--Mills Reduced-Fibre Wilson Shell Symbol and Local-Vacuum Programme V27)
 /Subject (The seven-face support gap, exact sixth and seventh Wilson coefficients, and a four-term strong-coupling crossing)
}

\section{Version V27: the seven-face support gap, exact sixth and seventh Wilson coefficients, and the first degree-eight geometry}
\label{sec:YMA27-direction}

Version~V26 is preserved above.  Its final mandate asked for a complete
enumeration of total-degree-six spin foams before the minimal-sheet coefficient
could be promoted to the full Wilson coefficient.  The first step of the
present continuation is to check whether such a new support can occur at that
degree at all.

The answer is no.  On the embedded straight-dyadic coarse plaquette, every
nonzero character monomial carrying the fundamental coarse boundary has the
flat four-tile disk as its complete plaquette support through total degree
seven.  The proof is representation independent up to the final boundary
label: a singly incident integrated link kills the monomial, and an exact
finite completion search shows that the flat disk is the only plaquette support
with the required boundary and no dangling integrated link before eight
faces.  At eight faces exactly sixteen new supports occur, namely the
four possible elementary three-cube bumps on each of the four disk tiles.

The consequence is stronger than the Version~V26 request:
\begin{equation}
 \boxed{
 \begin{gathered}
 \text{there is no external or baryonic degree-six source foam,}\\
 \text{there is no external degree-seven source foam,}\\
 \text{and the first genuinely new geometry occurs at degree eight.}
 \end{gathered}}
 \label{eq:YMA27-main-geometric-answer}
\end{equation}
Hence the complete coefficients through degree seven are obtained from the
vacuum-normalized fundamental activity on the minimal disk.  The resulting
axial Ward quotient is
\begin{equation}
 \boxed{
 \begin{aligned}
 g_{M,\beta}^{\rm W}(q)
 ={}&\frac{\beta^4}{17496}
   +\frac{\beta^5}{52488}
   +\frac{\beta^6}{419904}
   -\frac{\beta^7}{11337408}
   +O_M(\beta^8).
 \end{aligned}}
 \label{eq:YMA27-Ward-summary}
\end{equation}
In particular, the order-six coefficient left open in Version~V26 is exactly
its already computed minimal-sheet value and is positive.  The order-seven
coefficient is also exact and is negative.

This is a correction of the stopping boundary, not a deletion of an earlier
proof.  Version~V26 correctly kept the degree-six coefficient provisional.
Version~V27 proves that the provisional sector exhausts that degree.  It also
prevents the programme from opening an unnecessary epsilon-intertwiner
compiler before the first degree at which a nonminimal support actually
exists.

\begin{firewall}[Scope of the strong-coupling continuation]
\label{fire:YMA27-scope}
All statements below concern the embedded pure-Wilson straight-dyadic shell.
They do not insert the response-matched generated action, identify a
weak-facing selected root, or evaluate the Wilson symbol at that root.  The
full-link support argument is performed before the parity-tree quotient; the
normalized gauge-volume factor is coarse-field independent, so this is an
exact computational gauge for the same gauge-invariant reduced-fibre
integral.  The resulting Euclidean Hessian coefficients are not a transfer
mass or a continuum Yang--Mills mass gap.
\end{firewall}

% =============================================================================
\section{Full-link support before the parity-tree quotient}
\label{sec:YMA27-full-link}
% =============================================================================

Let $D$ be the positively oriented $2\times2$ fine plaquette disk in the
$(1,2)$ plane whose boundary is one elementary straight-dyadic coarse
plaquette.  Its support consists of four elementary fine plaquettes and its
boundary consists of eight oriented fine links.

For the support audit it is convenient to restore the unreduced product-Haar
fibre.  Internal gauge transformations act freely up to the already separated
stabilizer, normalized Haar gauge volume is one, and the integrand is gauge
invariant.  Consequently the unreduced and parity-tree reduced integrals have
the same coarse-field-dependent coefficients.  In the unreduced
representation every nonretained fine link is integrated, so ordinary Haar
selection rules can be applied link by link.

Expand every normalized plaquette factor as in
\eqref{eq:YMA26-normalized-characters}.  A \emph{character monomial} assigns a
nontrivial irreducible character to each plaquette in a finite support
$\mathscr S$ and the trivial character elsewhere.  Its total Taylor degree is
the sum of the Taylor degrees selected from the corresponding activities.
Every occupied plaquette costs at least one Taylor degree.

\begin{lemma}[No dangling integrated link]
\label{lem:YMA27-no-dangling}
If a character monomial has an integrated fine link incident with exactly one
plaquette in its nontrivial support, its Haar integral is zero.  Consequently
every nonzero monomial satisfies
\begin{equation}
 \boxed{
 \deg_{\mathscr S}(e)\ne1
 \quad\text{for every integrated fine link }e.}
 \label{eq:YMA27-no-dangling}
\end{equation}
This conclusion is independent of the nontrivial irreducible label on the
incident plaquette.
\end{lemma}

\begin{proof}
Fix such a link $e$.  After all other link variables are held fixed, the
integrand depends on $U_e$ through one matrix coefficient of one nontrivial
irreducible representation, possibly with fixed matrices multiplied on its
left and right and possibly with $U_e$ replaced by $U_e^{-1}$.  The Haar
average of every matrix coefficient of a nontrivial irreducible representation
is zero.  Fubini's theorem therefore annihilates the complete monomial.
\end{proof}

\begin{lemma}[A closed plaquette component has at least six faces]
\label{lem:YMA27-closed-six}
Let $\mathscr C$ be a nonempty finite edge-connected plaquette set with no
distinguished boundary edge and with no edge of incidence one.  Then
\begin{equation}
 \boxed{|\mathscr C|\ge6.}
 \label{eq:YMA27-closed-six}
\end{equation}
The bound is attained by the six faces of an elementary three-cube.
\end{lemma}

\begin{proof}
Choose $p\in\mathscr C$.  Each of its four edges must be shared with another
plaquette.  Two distinct elementary plaquettes share at most one edge, so at
least four further plaquettes are required.  If there were exactly five
plaquettes in total, each of the four additional plaquettes would share one
edge with $p$.  Around $p$ their outward edges cannot all be paired: in the
only arrangement pairing the four edges transverse to $p$, the four remaining
outer edges are precisely the boundary of the missing opposite face of an
elementary three-cube and each has incidence one.  Every other arrangement
leaves a dangling edge earlier.  Hence a sixth plaquette is necessary.  The
boundary of an elementary three-cube shows sharpness.
\end{proof}

A monomial whose retained boundary coefficient contains
$\chi_{\mathbf3}(W_{\partial D})$ has a fundamental line on every edge of
$\partial D$ before boundary gauge fixing.  Therefore every boundary edge is
incident with at least one support plaquette.  Together with
\Cref{lem:YMA27-no-dangling}, this reduces the degree-six question to a finite
cubical support problem.

% =============================================================================
\section{The exact seven-face support classification}
\label{sec:YMA27-support-classification}
% =============================================================================

\begin{theorem}[Seven-face support gap and first cube bumps]
\label{thm:YMA27-seven-face-gap}
Let the fine periodic side be $N$.  For the seven-face conclusion assume
$N\ge8$; then every support of at most seven faces carrying the displayed
contractible boundary has a lift to the universal cubical cover.  For the
classification at exactly eight faces assume in addition that the support is
nonwrapping, for example $N\ge10$.  Let $\mathscr S$ be a finite set of
elementary plaquettes in such a lift such that:
\begin{enumerate}[label=\textup{(S\arabic*)},leftmargin=3.5em]
\item every one of the eight edges of $\partial D$ belongs to at least one
      plaquette of $\mathscr S$;
\item every edge outside $\partial D$ incident with $\mathscr S$ belongs to
      either zero or at least two plaquettes of $\mathscr S$.
\end{enumerate}
Then:
\begin{align}
 \boxed{
 |\mathscr S|\le7
 \quad\Longrightarrow\quad
 \mathscr S=D,}
 \label{eq:YMA27-seven-face-only-disk}\\
 \boxed{
 |\mathscr S|=8,
 \quad \mathscr S\ne D
 \quad\Longrightarrow\quad
 \mathscr S\text{ is one of exactly sixteen elementary cube bumps}.}
 \label{eq:YMA27-sixteen-bumps}
\end{align}
A cube bump is obtained by choosing one of the four tiles of $D$, one of the
two transverse coordinate directions, and one of its two sides, deleting the
chosen tile and inserting the other five faces of the adjacent elementary
three-cube.
\end{theorem}

\begin{proof}
The statement is a finite exhaustive cubical calculation.  We give the
completion algorithm and its completeness proof so that the computation is
not a black-box geometric search.

For a partial plaquette set $S$, form the edge-incidence counts.  Apply the
following deterministic recursion, choosing the lexicographically first edge
whenever a choice of defect edge is required.
\begin{enumerate}[label=\textup{(A\arabic*)},leftmargin=3.6em]
\item If a boundary edge is uncovered, branch over the six elementary
      plaquettes incident with that edge.
\item Otherwise, if an edge outside $\partial D$ has incidence one, branch
      over the five other plaquettes incident with that edge.
\item If neither defect occurs, record $S$.  If the face cutoff has not been
      reached, branch over every plaquette sharing an edge with $S$.
\item Discard every branch as soon as the face cutoff is exceeded, and
      identify repeated partial supports exactly by their integer-coordinate
      plaquette lists.
\end{enumerate}

Every admissible support occurs in this tree.  Indeed, at a defect edge any
admissible completion must contain one of the plaquettes in the corresponding
branch list.  If a defect-free partial support is a proper subset of an
edge-connected admissible support, the next plaquette may be chosen to share
an edge with the partial support.  A component not meeting the prescribed
boundary would be a closed component and would contain at least six faces by
\Cref{lem:YMA27-closed-six}; it cannot coexist with the four-face boundary
component below the eight-face cutoff.  Thus the recursion is exhaustive.
It is finite because every added plaquette lies within at most seven
edge-adjacency steps of $\partial D$.

The exact integer replay visits $126670$ distinct partial supports at cutoff
seven and returns one admissible support, of cardinality four: $D$.  At cutoff
eight it visits $628372$ distinct partial supports and returns seventeen
supports: $D$ and sixteen supports of cardinality eight.  An independent
generator forms the sets
\[
 (D\setminus\{p\})\cup
 (\partial Q\setminus\{p\}),
\]
where $p$ ranges over the four disk tiles and $Q$ ranges over the four
three-cubes incident with $p$ in the two transverse directions and their two
orientations.  Exact set comparison identifies these sixteen generated sets
with the sixteen nonplanar outputs of the exhaustive recursion.  This proves
\eqref{eq:YMA27-seven-face-only-disk}--
\eqref{eq:YMA27-sixteen-bumps}.
\end{proof}

\begin{remark}[Why epsilon channels do not enter before degree eight]
\label{rem:YMA27-epsilon-delay}
An $\SU(3)$ epsilon intertwiner can occur at a link incident with three
fundamental lines of the same orientation.  The support theorem does not
forbid such an intertwiner abstractly; it proves that no support on which it
could alter the elementary coarse-boundary coefficient exists through seven
faces.  On the unique support $D$, every internal edge has incidence two, so
Schur orthogonality, not an epsilon tensor, performs the gluing.
\end{remark}

\begin{corollary}[First nonminimal geometric contribution]
\label{cor:YMA27-first-bump}
At total Taylor degree eight, the sixteen coherently oriented fundamental cube
bumps contribute
\begin{equation}
 \boxed{
 \frac{16x^8}{3^7}
 \left(
  \chi_{\mathbf3}(W_{\partial D})
  +\chi_{\overline{\mathbf3}}(W_{\partial D})
 \right)}
 \label{eq:YMA27-bump-partition-piece}
\end{equation}
to the normalized partition function for one elementary coarse plaquette.
This is a rigorously identified part of the order-eight row, not its complete
effective-action or Hessian coefficient.
\end{corollary}

\begin{proof}
Every cube-bump support is a simply connected oriented plaquette disk of area
eight.  At total degree eight every occupied face uses a first-order
fundamental or antifundamental activity.  Each internal edge has incidence two,
so Schur orthogonality forces one common irreducible label and coherent
orientation throughout the disk.  The disk-gluing formula
\eqref{eq:YMA25-disk-formula} gives $3^{1-8}=3^{-7}$ for either orientation.
There are sixteen supports.
\end{proof}

\begin{firewall}[Why the complete order-eight symbol is still open]
\label{fire:YMA27-order-eight-open}
At degree eight the logarithm of the partition function begins to mix two
order-four coarse-boundary terms.  Adjacent coarse plaquettes can also produce
new connected multi-plaquette Hessian kernels, and the flat disk has its own
higher local activity coefficient.  Therefore
\eqref{eq:YMA27-bump-partition-piece} is not promoted to the complete
order-eight axial coefficient.  Version~V27 uses the support classification
only to close degrees six and seven, where none of these competitions can yet
occur.
\end{firewall}

% =============================================================================
\section{The normalized fundamental activity through fourth order}
\label{sec:YMA27-activity}
% =============================================================================

Retain
\(
 x=\beta/6
\)
and
\(
 V=\mathbf3\oplus\overline{\mathbf3}
\)
from Version~V26.  The next coefficient needed for the degree-seven disk is
fixed by one further finite representation-ring step.

\begin{lemma}[Fourth tensor multiplicities]
\label{lem:YMA27-fourth-multiplicities}
The fourth tensor power of $V$ has multiplicities
\begin{equation}
 \boxed{
 [\mathbf1:V^{\otimes4}]=12,
 \qquad
 [\mathbf3:V^{\otimes4}]
 =[\overline{\mathbf3}:V^{\otimes4}]=15.}
 \label{eq:YMA27-fourth-multiplicities}
\end{equation}
Consequently,
\begin{align}
 \boxed{
 a_{\mathbf1}
 =1+x^2+\frac{x^3}{3}+\frac{x^4}{2}+O(x^5),}
 \label{eq:YMA27-a1-four}\\
 \boxed{
 a_{\mathbf3}=a_{\overline{\mathbf3}}
 =x+\frac{x^2}{2}+x^3+\frac{5x^4}{8}+O(x^5).}
 \label{eq:YMA27-a3-four}
\end{align}
\end{lemma}

\begin{proof}
Use the $\SU(3)$ Pieri rules in Dynkin labels:
\begin{align*}
 (p,q)\otimes\mathbf3
 &=(p+1,q)
  \oplus(p-1,q+1)_{p\ge1}
  \oplus(p,q-1)_{q\ge1},\\
 (p,q)\otimes\overline{\mathbf3}
 &=(p,q+1)
  \oplus(p+1,q-1)_{q\ge1}
  \oplus(p-1,q)_{p\ge1}.
\end{align*}
The degree-three recursion gives the needed multiplicities
\[
 m_3(\mathbf1)=2,
 \quad m_3(\mathbf3)=m_3(\overline{\mathbf3})=6,
 \quad m_3(\mathbf6)=m_3(\overline{\mathbf6})=3,
 \quad m_3(\mathbf8)=4.
\]
A trivial representation at degree four arises from a terminal
$\mathbf3\otimes\overline{\mathbf3}$ or
$\overline{\mathbf3}\otimes\mathbf3$, giving $6+6=12$.
A terminal fundamental arises from
\[
 \mathbf1\otimes\mathbf3,
 \qquad
 \mathbf8\otimes\mathbf3,
 \qquad
 \overline{\mathbf3}\otimes\overline{\mathbf3},
 \qquad
 \mathbf6\otimes\overline{\mathbf3},
\]
with respective inherited multiplicities $2,4,6,3$, whose sum is fifteen.
Conjugation gives the antifundamental multiplicity.  Division by $4!$ in
\eqref{eq:YMA26-character-multiplicity-series} proves the two series.
\end{proof}

\begin{theorem}[Normalized activity and four-disk series through degree seven]
\label{thm:YMA27-activity-seven}
The vacuum-normalized fundamental activity satisfies
\begin{equation}
 \boxed{
 u_{\mathbf3}(x)=u_{\overline{\mathbf3}}(x)
 =x+\frac{x^2}{2}-\frac{5x^4}{24}+O(x^5).}
 \label{eq:YMA27-u3-four}
\end{equation}
Its fourth power is
\begin{equation}
 \boxed{
 u_{\mathbf3}(x)^4
 =x^4+2x^5+\frac32x^6-\frac13x^7+O(x^8).}
 \label{eq:YMA27-u3-fourth-seven}
\end{equation}
\end{theorem}

\begin{proof}
From \eqref{eq:YMA27-a1-four},
\[
 a_{\mathbf1}^{-1}
 =1-x^2-\frac{x^3}{3}+\frac{x^4}{2}+O(x^5).
\]
Multiplication by \eqref{eq:YMA27-a3-four} gives
\[
 x+\frac{x^2}{2}
 +\left(1-1\right)x^3
 +\left(\frac58-\frac12-\frac13\right)x^4
 +O(x^5),
\]
which is \eqref{eq:YMA27-u3-four}.  In the fourth power, the degree-six
coefficient is obtained by choosing two quadratic terms and equals
\(
 \binom42(1/2)^2=3/2
\).
At degree seven there are two contributions: three quadratic choices give
\(
 \binom43(1/2)^3=1/2
\), while one quartic choice gives
\(
 4(-5/24)=-5/6
\).  Their sum is $-1/3$.
\end{proof}

% =============================================================================
\section{Complete pure-Wilson coefficients through degree seven}
\label{sec:YMA27-coefficients}
% =============================================================================

\begin{theorem}[No nonminimal coarse-boundary term through degree seven]
\label{thm:YMA27-degree-seven-classification}
For one embedded elementary coarse plaquette, every pure-Wilson character
monomial of total Taylor degree at most seven which contributes to the
fundamental or antifundamental coarse boundary has support $D$.  On that
support every plaquette carries the same fundamental or antifundamental
irreducible label.  Hence the complete normalized coarse-boundary partition
coefficient through degree seven is
\begin{equation}
 \boxed{
 \frac{u_{\mathbf3}(x)^4}{27}
 \left(
  \chi_{\mathbf3}(W_P)
  +\chi_{\overline{\mathbf3}}(W_P)
 \right)
 +O(x^8).}
 \label{eq:YMA27-complete-partition-seven}
\end{equation}
No epsilon or distinct higher-irrep channel remains outside the activity
$u_{\mathbf3}$ at these degrees.
\end{theorem}

\begin{proof}
At total Taylor degree $n\le7$, the number of plaquettes carrying a nontrivial
character is at most $n$.  A nonzero monomial satisfies the no-dangling rule of
\Cref{lem:YMA27-no-dangling} and covers the prescribed boundary.  The support
classification \eqref{eq:YMA27-seven-face-only-disk} therefore forces the
support to be $D$.

Every internal link of $D$ is incident with exactly two support plaquettes.
The two-representation Haar integral is zero unless the irreducible labels
match with opposite link orientation.  Connectivity propagates one common
label over all four tiles.  The retained boundary coefficient is fundamental
or antifundamental, so the common label is respectively $\mathbf3$ or
$\overline{\mathbf3}$.  All local tensor multiplicities and trivial-character
normalizations are already contained in the exact activity $u_r$.  The disk
gluing factor is $3^{-3}=1/27$, proving the formula.
\end{proof}

\begin{theorem}[Exact induced coarse Wilson coupling through degree seven]
\label{thm:YMA27-beta-coarse-seven}
For every embedded finite periodic straight-dyadic pure-Wilson card,
\begin{equation}
 \boxed{
 \begin{aligned}
 \beta_{\rm c}(\beta)
 ={}&\frac{\beta^4}{5832}
  +\frac{\beta^5}{17496}
  +\frac{\beta^6}{139968}
  -\frac{\beta^7}{3779136}
  +O_M(\beta^8).
 \end{aligned}}
 \label{eq:YMA27-beta-coarse-seven}
\end{equation}
Equivalently, through degree seven the field-dependent effective action is
\begin{equation}
 \boxed{
 S_{M,\beta}^{\rm W}(W)
 =C_M(\beta)
 +\beta_{\rm c}(\beta)
  \sum_{P\in\mathscr P_c}
  \left(1-\frac13\operatorname{Re}\Tr W_P\right)
 +O_M(\beta^8).}
 \label{eq:YMA27-action-seven}
\end{equation}
The coefficient $1/139968$ is the complete order-six coefficient left open in
Version~V26.
\end{theorem}

\begin{proof}
By \Cref{thm:YMA27-degree-seven-classification} and
\eqref{eq:YMA27-u3-fourth-seven}, the normalized partition coefficient of
$\chi_{\mathbf3}+\chi_{\overline{\mathbf3}}$ is
\[
 \frac1{27}
 \left(x^4+2x^5+\frac32x^6-\frac13x^7\right)
 +O(x^8).
\]
The first nonconstant partition term has degree four, so nonlinear terms in
$-\log Z$ begin at degree eight.  Thus no logarithmic product changes degrees
four through seven.  As in Version~V26, conversion from the coefficient of
$\chi_{\mathbf3}+\chi_{\overline{\mathbf3}}$ to the coefficient of
$1-\frac13\operatorname{Re}\Tr W_P$ multiplies by $6$.  Hence
\[
 \beta_{\rm c}
 =\frac29
 \left(x^4+2x^5+\frac32x^6-\frac13x^7\right)
 +O(x^8).
\]
Substitution of $x=\beta/6$ gives
\[
 \frac29\frac1{6^4}=\frac1{5832},
 \qquad
 \frac49\frac1{6^5}=\frac1{17496},
\]
\[
 \frac13\frac1{6^6}=\frac1{139968},
 \qquad
 -\frac2{27}\frac1{6^7}=-\frac1{3779136},
\]
which proves the theorem.
\end{proof}

\begin{corollary}[Exact central axial Ward quotient through degree seven]
\label{cor:YMA27-Ward-seven}
For every nonzero central axial transverse mode,
\begin{align}
 \boxed{
 h_{M,\beta}^{\rm W}(q)
 =\omega(q)
 \left(
  \frac{\beta^4}{17496}
  +\frac{\beta^5}{52488}
  +\frac{\beta^6}{419904}
  -\frac{\beta^7}{11337408}
 \right)
 +O_M(\beta^8\omega(q)),}
 \label{eq:YMA27-h-seven}\\
 \boxed{
 g_{M,\beta}^{\rm W}(q)
 =\frac{\beta^4}{17496}
  +\frac{\beta^5}{52488}
  +\frac{\beta^6}{419904}
  -\frac{\beta^7}{11337408}
  +O_M(\beta^8).}
 \label{eq:YMA27-g-seven}
\end{align}
Equivalently,
\begin{equation}
 \boxed{
 g_{M,\beta}^{\rm W}(q)
 =\frac{\beta^4}{17496}
 \left(
  1+\frac\beta3+\frac{\beta^2}{24}
  -\frac{\beta^3}{648}
  +O_M(\beta^4)
 \right).}
 \label{eq:YMA27-g-factorized-seven}
\end{equation}
\end{corollary}

\begin{proof}
The coarse Wilson Hessian is
\(
 \frac13\beta_{\rm c}d_1^*d_1
\), and the axial transverse eigenvalue of $d_1^*d_1$ is $\omega(q)$.
Apply \Cref{thm:YMA27-beta-coarse-seven}.  The factorized ratios follow from
\[
 \frac{1/52488}{1/17496}=\frac13,
 \qquad
 \frac{1/419904}{1/17496}=\frac1{24},
 \qquad
 \frac{-1/11337408}{1/17496}=-\frac1{648}.
\]
\end{proof}

\begin{assessment}[Exact qualification of Version V26]
\label{ass:YMA27-V26-qualification}
The Version~V26 statement that $1/419904$ was a minimal-sheet preview was
correctly cautious.  The support theorem now proves that no additional
connected degree-six support exists, so
\begin{equation}
 \boxed{c_6^{\rm full}=\frac1{419904}>0.}
 \label{eq:YMA27-c6-full}
\end{equation}
The same argument gives the previously unrequested exact coefficient
\begin{equation}
 \boxed{c_7^{\rm full}=-\frac1{11337408}.}
 \label{eq:YMA27-c7-full}
\end{equation}
The first coefficient requiring a genuinely nonminimal geometry is $c_8$.
\end{assessment}

% =============================================================================
\section{Volume-uniform order-eight remainder and a terminating finite card}
\label{sec:YMA27-uniform-remainder}
% =============================================================================

Retain the common analytic disk $|\beta|\le r_*$ and the volume-independent
bound $B_*$ from Version~V25.  Put
\begin{equation}
 \begin{gathered}
 c_4=\frac1{17496},
 \qquad c_5=\frac1{52488},
 \qquad c_6=\frac1{419904},
 \qquad c_7=-\frac1{11337408},\\
 P_7(\beta)=c_4\beta^4+c_5\beta^5+c_6\beta^6+c_7\beta^7,
 \qquad
 C_8:=\frac{2B_*}{r_*^8}.
 \end{gathered}
 \label{eq:YMA27-C8-def}
\end{equation}

\begin{theorem}[Uniform remainder after the exact seventh coefficient]
\label{thm:YMA27-uniform-eight}
For every finite periodic volume, every central axial momentum, and
$0\le\beta\le r_*/2$,
\begin{equation}
 \boxed{
 \left|
  g_{M,\beta}^{\rm W}(q)-P_7(\beta)
 \right|
 \le C_8\beta^8.}
 \label{eq:YMA27-uniform-eight}
\end{equation}
Moreover,
\begin{equation}
 \boxed{
 \left|
 \frac{\dd}{\dd\beta}
 \left[
  \frac{g_{M,\beta}^{\rm W}(q)-P_7(\beta)}{\beta}
 \right]
 \right|
 \le8C_8\beta^6.}
 \label{eq:YMA27-derivative-eight}
\end{equation}
Both bounds pass to the strong-coupling local-vacuum limit.
\end{theorem}

\begin{proof}
Write
\(
 g_{M,\beta}^{\rm W}(q)=\sum_{n\ge4}c_{n,M,q}\beta^n
\).
The preceding section identifies $c_{n,M,q}=c_n$ for $4\le n\le7$, while
Cauchy's estimate from Version~V25 gives
\(
 |c_{n,M,q}|\le B_*r_*^{-n}
\)
for $n\ge8$.  With $x=\beta/r_*\le1/2$,
\[
 \sum_{n\ge8}|c_{n,M,q}|\beta^n
 \le B_*\sum_{n\ge8}x^n
 \le2B_*x^8=C_8\beta^8.
\]
For the derivative,
\[
 \sum_{n\ge8}(n-1)x^{n-2}
 =\frac{x^6(7-6x)}{(1-x)^2}
 \le16x^6
 \qquad(0\le x\le1/2).
\]
Multiplication by $B_*/r_*^2$ gives
\(
 16B_*\beta^6/r_*^8=8C_8\beta^6
\).
Uniform convergence gives the local-limit statement.
\end{proof}

\begin{theorem}[Degree-seven strong-coupling sign interval]
\label{cert:YMA27-finite-sign}
Fix $q\ne0$, target fraction $\delta\ge0$, and
$0<\beta\le r_*/2$.  Put
\begin{equation}
 \begin{aligned}
 L_7(\beta,q;\delta)
 &:={P_7(\beta)-\delta\beta r_0(q)-C_8\beta^8},\\
 U_7(\beta,q;\delta)
 &:={P_7(\beta)-\delta\beta r_0(q)+C_8\beta^8}.
 \end{aligned}
 \label{eq:YMA27-sign-interval}
\end{equation}
Then
\begin{equation}
 \boxed{
 L_7
 \le
 \frac{h_{M,\beta}^{\rm W}(q)-\delta\beta s_0(q)}{\omega(q)}
 \le U_7.}
 \label{eq:YMA27-sign-interval-bound}
\end{equation}
Consequently,
\begin{equation}
 \boxed{
 \begin{array}{rcl}
 L_7>0&\Longrightarrow&\pass,\\[0.2em]
 U_7<0&\Longrightarrow&\obstruction,\\[0.2em]
 0\in[L_7,U_7]&\Longrightarrow&\stopstatus.
 \end{array}}
 \label{eq:YMA27-sign-verdict}
\end{equation}
On the obstruction branch the predeclared axial transverse polarization is an
attained witness.  A degree-eight enumeration is required only when this
already rigorous interval overlaps zero at the physically supplied coupling.
\end{theorem}

\begin{proof}
Use $s_0(q)=\omega(q)r_0(q)$ and
\Cref{thm:YMA27-uniform-eight}.  Since $\omega(q)>0$, division preserves the
sign and the displayed alternatives are immediate.
\end{proof}

% =============================================================================
\section{Four-term location of the unique small-target crossing}
\label{sec:YMA27-crossing}
% =============================================================================

Retain the unique monotone strong-coupling crossing of
\Cref{thm:YMA26-monotone-branch}.  Put
\begin{equation}
 t_q(\delta)
 =\left(17496\,r_0(q)\,\delta\right)^{1/3}.
 \label{eq:YMA27-t}
\end{equation}

\begin{theorem}[Four-term crossing expansion]
\label{thm:YMA27-crossing-four-term}
As $\delta\downarrow0$, the unique strong-coupling crossing satisfies
\begin{equation}
 \boxed{
 \begin{aligned}
 \beta_{M,q}^{\rm sc}(\delta)
 ={}&t_q
 -\frac19t_q^2
 +\frac5{216}t_q^3
 -\frac{41}{8748}t_q^4
 +O(t_q^5).
 \end{aligned}}
 \label{eq:YMA27-crossing-four-term}
\end{equation}
The remainder is uniform in the finite periodic volume and in the central
axial momentum.
\end{theorem}

\begin{proof}
By \eqref{eq:YMA27-g-factorized-seven}, the crossing equation is
\begin{equation}
 \beta^3
 \left(
  1+\frac\beta3+\frac{\beta^2}{24}
  -\frac{\beta^3}{648}+O(\beta^4)
 \right)
 =t^3,
 \qquad t=t_q(\delta).
 \label{eq:YMA27-crossing-normalized}
\end{equation}
The analytic implicit branch constructed in Version~V26 may be written
\[
 \beta=t+At^2+Bt^3+Ct^4+O(t^5).
\]
Substitution in \eqref{eq:YMA27-crossing-normalized} and comparison of the
coefficients of $t,t^2,t^3$ after division by $t^3$ gives successively
\begin{align*}
 3A+\frac13&=0,\\
 3A^2+\frac43A+3B+\frac1{24}&=0,\\
 A^3+2A^2+6AB+\frac5{24}A+\frac43B+3C-\frac1{648}&=0.
\end{align*}
Solving in order yields
\[
 A=-\frac19,
 \qquad
 B=\frac5{216},
 \qquad
 C=-\frac{41}{8748}.
\]
The common analytic disk, the uniform order-eight remainder, and the compact
reference interval $5/28\le r_0(q)\le1/3$ give a common fifth-order Taylor
bound.
\end{proof}

\begin{corollary}[Uniform outward crossing bracket through fourth order]
\label{cor:YMA27-crossing-bracket}
There are constants $C_{\rm cr,5}>0$ and $\delta_{\rm cr,5}>0$, independent of
the finite volume and central axial momentum, such that for
$0<\delta<\delta_{\rm cr,5}$,
\begin{equation}
 \boxed{
 \begin{aligned}
 \beta_-^{(7)}
 &:={t_q-\frac19t_q^2+\frac5{216}t_q^3
     -\frac{41}{8748}t_q^4-C_{\rm cr,5}t_q^5},\\
 \beta_+^{(7)}
 &:={t_q-\frac19t_q^2+\frac5{216}t_q^3
     -\frac{41}{8748}t_q^4+C_{\rm cr,5}t_q^5}
 \end{aligned}}
 \label{eq:YMA27-crossing-bracket}
\end{equation}
satisfy
\begin{equation}
 \boxed{
 \beta_-^{(7)}
 \le\beta_{M,q}^{\rm sc}(\delta)
 \le\beta_+^{(7)}.}
 \label{eq:YMA27-crossing-bracket-order}
\end{equation}
After reducing $\delta_{\rm cr,5}$ if necessary, both endpoints lie inside the
strictly monotone interval of Version~V26.
\end{corollary}

\begin{proof}
Uniform analyticity of the implicit branch gives a common bound on its fifth
derivative near $t=0$.  Taylor's theorem yields the interval.  Since both
endpoints tend to zero with $t_q$, they lie inside the inherited monotone
window for sufficiently small common $\delta_{\rm cr,5}$.
\end{proof}

\begin{assessment}[Why the coefficient ladder now stops at degree seven]
\label{ass:YMA27-stop-ladder}
The strong-coupling expansion now gives a rigorous finite sign interval and a
four-term isolated crossing.  Computing degree eight is target native only if
an actual selected-root or response-matched coupling lies inside this analytic
window and the interval \eqref{eq:YMA27-sign-interval} still overlaps zero.
Without such a physical coupling, another formal coefficient would again move
away from the programme's selected-shell objective.
\end{assessment}

% =============================================================================
\section{Exact replay, integrated theorem, and revised stopping boundary}
\label{sec:YMA27-replay}
% =============================================================================

The accompanying exact replay performs two independent finite calculations.
First, the $\SU(3)$ Pieri recursion constructs the tensor multiplicities
through degree four, forms the vacuum-normalized activity, and verifies every
coefficient in \eqref{eq:YMA27-u3-fourth-seven},
\eqref{eq:YMA27-beta-coarse-seven}, and
\eqref{eq:YMA27-g-seven}.  Second, an integer-coordinate cubical completion
search verifies the support classification through eight faces and compares
its sixteen nonplanar outputs with an independently generated list of cube
bumps.

The replay reports:
\begin{equation}
 \boxed{
 \begin{array}{c|rrr}
 \text{face cutoff}&\text{visited partial supports}
 &\text{admissible supports}&\text{nonplanar supports}\\
 \hline
 7&126670&1&0\\
 8&628372&17&16.
 \end{array}}
 \label{eq:YMA27-replay-counts}
\end{equation}
The canonical serialization of the sixteen cube-bump supports has SHA--256
\begin{center}\scriptsize\ttfamily
698a9829612ea03c6d1151855790f1d7f84fc6502d0bbdbd03e8a708a45e8d33.
\end{center}
Normal and optimized Python executions produce byte-identical JSON.

\begin{theorem}[Version V27 seven-face and exact-coefficient theorem]
\label{thm:YMA27-integrated}
Preserving Versions~V1--V26, Version~V27 proves:
\begin{enumerate}[label=\textup{(V27.\arabic*)},leftmargin=5.6em]
\item every nonzero character monomial has no singly incident integrated fine
      link;
\item the flat four-plaquette disk is the only support carrying the elementary
      coarse boundary through seven faces;
\item the first nonplanar supports occur at eight faces and are exactly the
      sixteen elementary cube bumps;
\item the complete pure-Wilson degree-six and degree-seven coarse and axial
      coefficients are those in
      \eqref{eq:YMA27-beta-coarse-seven}--\eqref{eq:YMA27-g-seven};
\item the order-six axial coefficient is $1/419904>0$, while the order-seven
      coefficient is $-1/11337408$;
\item the remainder after degree seven is $O(\beta^8)$ uniformly in the finite
      periodic volume and central axial momentum on the inherited analytic
      disk;
\item the degree-seven polynomial and that remainder form a terminating finite
      margin certificate;
\item the unique small-target strong-coupling crossing has the four-term
      expansion \eqref{eq:YMA27-crossing-four-term} and the uniform outward
      bracket \eqref{eq:YMA27-crossing-bracket}.
\end{enumerate}
No item evaluates the response-matched selected root or proves local-vacuum
contraction.
\end{theorem}

\begin{proof}
Items~\textup{(V27.1)--(V27.3)} follow from
\Cref{lem:YMA27-no-dangling,lem:YMA27-closed-six} and
\Cref{thm:YMA27-seven-face-gap,cor:YMA27-first-bump}.

Items~\textup{(V27.4)--(V27.5)} follow from
\Cref{lem:YMA27-fourth-multiplicities,thm:YMA27-activity-seven},
\Cref{thm:YMA27-degree-seven-classification}, and
\Cref{thm:YMA27-beta-coarse-seven,cor:YMA27-Ward-seven}.

Items~\textup{(V27.6)--(V27.7)} are
\Cref{thm:YMA27-uniform-eight,cert:YMA27-finite-sign}.  Item~\textup{(V27.8)}
is \Cref{thm:YMA27-crossing-four-term,cor:YMA27-crossing-bracket}.
\end{proof}

\begingroup\small
\begin{longtable}{>{\raggedright\arraybackslash}p{0.27\textwidth}
                  >{\raggedright\arraybackslash}p{0.18\textwidth}
                  >{\raggedright\arraybackslash}p{0.47\textwidth}}
\caption{Version V27 proof-status ledger.}\label{tab:YMA27-ledger}\\
\toprule
\textbf{Row}&\textbf{Status}&\textbf{Exact scope}\\
\midrule
\endfirsthead
\toprule
\textbf{Row}&\textbf{Status}&\textbf{Exact scope}\\
\midrule
\endhead
V1--V26 prefix
& Preserved
& Complete preceding source retained byte for byte before the terminal command.\\[0.35em]
Seven-face support classification
& \pass
& Exact finite cubical completion proves that only the flat disk occurs through
  seven faces.\\[0.35em]
First nonminimal geometry
& \pass
& Exactly sixteen size-eight supports, all elementary three-cube bumps.\\[0.35em]
Complete degree-six Ward coefficient
& \pass
& $c_6=1/419904$; the V26 minimal-sheet preview is the full coefficient.\\[0.35em]
Complete degree-seven Ward coefficient
& \pass
& $c_7=-1/11337408$; no external support contributes.\\[0.35em]
Uniform post-seventh remainder
& \pass
& $|g-P_7|\le C_8\beta^8$ on the inherited common analytic disk.\\[0.35em]
Small-target crossing location
& \pass
& Four-term uniform expansion and an $O(t^5)$ outward bracket.\\[0.35em]
Complete degree-eight axial coefficient
& \stopstatus
& Cube bumps are classified, but logarithmic mixing and connected
  multi-coarse-plaquette Hessian terms begin at this degree.\\[0.35em]
Native response cofactors and selected-root action
& \stopstatus
& No supplied physical cylinder populates the two cofactors or the complete
  added action.\\[0.35em]
Selected-root Wilson symbol
& \stopstatus
& The exact degree-seven interval has not been evaluated on an actual
  response-matched selected-root box.\\[0.35em]
Local-vacuum or continuum mass
& Not inferred
& Requires the response-matched signed head, differentiated tail, conductance,
  ancestry, cutoff transport, and reflection-positive continuum rows.\\
\bottomrule
\end{longtable}
\endgroup

% =============================================================================
\section{Revised V28 acquisition}
\label{sec:YMA27-next}
% =============================================================================

\begin{openproblem}[V28: use the exact degree-seven card before opening degree eight]
\label{op:YMA27-V28}
Preserve the complete present source.  Proceed in the following order.
\begin{enumerate}[label=\textup{(V28.\arabic*)},leftmargin=5.7em]
\item Load one actual pure-Wilson or response-matched coupling interval and one
      retained central axial momentum inside the common analytic disk.
\item Evaluate the terminating interval
      \eqref{eq:YMA27-sign-interval} using the exact degree-seven polynomial and
      the inherited $C_8$ remainder.
\item If the interval is strictly signed, return the corresponding finite
      \pass{} or attained \obstruction{} and do not compute degree eight.
\item If the interval overlaps zero and the coupling is physically selected,
      enumerate the complete order-eight connected Hessian row, beginning with
      the sixteen cube bumps, the flat-disk activity term, and the first
      logarithmic/multi-coarse-plaquette contributions.
\item On the native selected-shell branch, populate the two response cofactors
      and the complete added action before comparing with the pure-Wilson
      crossing bracket.
\item Reopen the differentiated Ward tail, boundary incidence, Witten return,
      or cutoff transport only after the same-action selected mode has been
      evaluated.
\item Return \pass, an attained \obstruction, or \stopstatus{} at the first
      unresolved action, coupling, coefficient, or remainder row.
\end{enumerate}
A degree-eight calculation without a supplied coupling interval or a
selected-root comparison is optional strong-coupling information, not the next
mandatory YM--A acquisition.
\end{openproblem}

\begin{assessment}[Version V27 termination]
\label{ass:YMA27-termination}
The finite spin-foam question posed by Version~V26 is closed more strongly than
requested: no new support exists at degree six or seven, both complete
coefficients are exact, and the first nonminimal support is classified at
degree eight.  The next target-native step is to apply the rigorous
seventh-order interval to a physically supplied coupling or selected-root box.
The response-matched selected-shell symbol remains \stopstatus, and no negative
pure-Wilson Hessian has been found.
\end{assessment}

\section*{Version V27 append-only continuation record}
\addcontentsline{toc}{section}{Version V27 append-only continuation record}
The complete Version~V26 source preceding its terminal document-closing
command is preserved byte for byte.  The SHA--256 digest of the complete V26
source is
\begin{center}\scriptsize\ttfamily
88ab3b9e76d503413e5d7b22eb6a9c59caead8fdb2784daa46044216fe0845ec.
\end{center}
Its preterminal prefix contains $1215767$ bytes and has SHA--256
\begin{center}\scriptsize\ttfamily
ebace2f9bcbff20cec5c5ed7814da124b84e13efbea94f4a122d5bf83bb959a5.
\end{center}
For Version~V28, preserve this accumulated source, append new material
immediately before the final active document-closing command, and increment
the filename to \texttt{\_V28.tex}.

% =============================================================================
% END APPENDED VERSION V27 MATERIAL
% =============================================================================


% =============================================================================
% START APPENDED VERSION V28 MATERIAL
% =============================================================================

\clearpage
\hypersetup{
 pdftitle={Renewal Yang--Mills Reduced-Fibre Wilson Shell Symbol and Local-Vacuum Programme V28},
 pdfsubject={Character-tension calibration, the first cube-bump renormalization, and scale-matched selected-root margins}
}

\pdfinfo{
 /Title (Renewal Yang--Mills Reduced-Fibre Wilson Shell Symbol and Local-Vacuum Programme V28)
 /Subject (Character-tension calibration, the first cube-bump renormalization, and scale-matched selected-root margins)
}

\section{Version V28: character tension, first cube-bump renormalization, and scale-matched margins}
\label{sec:YMA28-direction}

Version~V27 is preserved above.  Its stopping rule correctly refuses to open a
complete degree-eight Hessian computation before an actual selected coupling
is supplied.  Continuing only by adding further Taylor coefficients would
therefore become circular.  The upstream issue is instead the scalar used to
compare two different physical scales.

Three quantities have been denoted by coupling-like symbols in the preceding
programme, but they are not the same object:
\begin{equation}
 \boxed{
 \begin{gathered}
 \beta
   =\text{incoming fine Wilson parameter},\\
 \zeta(\beta)
   =\text{vacuum-normalized fundamental character activity},\\
 \beta_{\square}(\beta)
   =\text{outgoing fundamental plaquette coefficient},\\
 \beta_{\rm H}(\beta,q):=3g_{M,\beta}^{\rm W}(q)
   =\text{coupling read from the actual axial Hessian}.
 \end{gathered}}
 \label{eq:YMA28-four-calibrations}
\end{equation}
At strong coupling the first quantity is of order \(\beta\), whereas the last
three are of order \(\beta^4\).  A margin relative to the incoming parameter
and a margin relative to an outgoing or Hessian calibration are consequently
different statements.  This distinction must be resolved before the
small-\(\beta\) relative-margin failure of Version~V25 is interpreted as a
failure of the local coarse action.

The present continuation makes this calibration explicit.  It proves four
new Wilson-specific facts.
\begin{equation}
 \boxed{
 \begin{gathered}
 \text{the normalized fundamental activity is computed one order further;}\\
 \text{the first four-dimensional correction to its dyadic area map is the}\\
 \text{sum of the sixteen cube bumps found in Version~V27;}\\
 \text{the associated B\"ottcher tension agrees with the physical planar}\\
 \text{Creutz selector through the first nonplanar relative order;}\\
 \text{incoming- and outgoing-calibrated margins have}\\
 \text{different computable strong-coupling limits.}
 \end{gathered}}
 \label{eq:YMA28-main-claims}
\end{equation}
The same calculation also proves that the scalar Wilson-action line ceases to
exhaust even the one-plaquette character sector at degree eight: nonzero
sextet and adjoint action coordinates occur.  Thus the scalar activity is an
exact calibration coordinate, but not a substitute for the complete
generated action.

\begin{firewall}[Scope of Version V28]
\label{fire:YMA28-scope}
Every coefficient below belongs to the embedded pure-Wilson straight-dyadic
shell and to the fixed planar Creutz quartet already used in Version~V25.  The
large-response root theorem concerns the strong-coupling branch of that
finite selector.  It does not supply the predeclared numerical response value,
prove that the selected cofinal roots lie on this branch, populate the two
native response cofactors, or identify the scalar character projection with
the full four-dimensional renormalization map.  No transfer mass or continuum
Yang--Mills mass is inferred.
\end{firewall}

% =============================================================================
\section{The normalized fundamental activity through fifth order}
\label{sec:YMA28-activity}
% =============================================================================

Retain
\begin{equation}
 x=\frac\beta6,
 \qquad
 z=\frac\beta{18}=\frac x3,
 \qquad
 V=\mathbf3\oplus\overline{\mathbf3}.
 \label{eq:YMA28-x-z}
\end{equation}
For the one-plaquette character expansion write
\begin{equation}
 e^{x(\chi_{\mathbf3}+\chi_{\overline{\mathbf3}})}
 =\sum_{r}a_r(x)\chi_r,
 \qquad
 u_r(x)=\frac{a_r(x)}{a_{\mathbf1}(x)}.
 \label{eq:YMA28-character-activities}
\end{equation}
The dimension-normalized fundamental activity is
\begin{equation}
 \boxed{
 \zeta(\beta)
 :=\frac{u_{\mathbf3}(\beta)}{3}
 =\frac{u_{\overline{\mathbf3}}(\beta)}{3}.}
 \label{eq:YMA28-zeta-definition}
\end{equation}
The division by three is forced by disk gluing: it is this coordinate, rather
than the raw character coefficient, which multiplies under planar area.

\begin{lemma}[Fifth tensor multiplicities]
\label{lem:YMA28-fifth-multiplicities}
The fifth tensor power of \(V\) has
\begin{equation}
 \boxed{
 [\mathbf1:V^{\otimes5}]=30,
 \qquad
 [\mathbf3:V^{\otimes5}]
 =[\overline{\mathbf3}:V^{\otimes5}]=65.}
 \label{eq:YMA28-fifth-multiplicities}
\end{equation}
Consequently,
\begin{align}
 a_{\mathbf1}(x)
 &=1+x^2+\frac{x^3}{3}+\frac{x^4}{2}+\frac{x^5}{4}
   +O(x^6),
 \label{eq:YMA28-a1-five}\\
 a_{\mathbf3}(x)=a_{\overline{\mathbf3}}(x)
 &=x+\frac{x^2}{2}+x^3+\frac{5x^4}{8}+\frac{13x^5}{24}
   +O(x^6).
 \label{eq:YMA28-a3-five}
\end{align}
\end{lemma}

\begin{proof}
Use the Pieri recursion of
\Cref{lem:YMA27-fourth-multiplicities}.  A trivial representation at degree
five is obtained only by appending \(\overline{\mathbf3}\) to a terminal
\(\mathbf3\), or \(\mathbf3\) to a terminal
\(\overline{\mathbf3}\).  The degree-four multiplicities are fifteen in
each case, giving thirty.

A terminal \(\mathbf3=(1,0)\) at degree five can arise under tensoring by
\(\mathbf3\) from \(\mathbf1=(0,0)\) or \(\mathbf8=(1,1)\), and under
tensoring by \(\overline{\mathbf3}\) from
\(\overline{\mathbf3}=(0,1)\) or \(\mathbf6=(2,0)\).  Their degree-four
multiplicities are respectively
\[
 12,\qquad24,\qquad15,\qquad14,
\]
whose sum is sixty-five.  Conjugation gives the antifundamental value.
Division by \(5!\) in the exponential series proves
\eqref{eq:YMA28-a1-five}--\eqref{eq:YMA28-a3-five}.
\end{proof}

\begin{theorem}[Fundamental character calibration through fifth order]
\label{thm:YMA28-zeta-series}
The vacuum-normalized fundamental coefficient and its dimension-normalized
activity are
\begin{align}
 \boxed{
 u_{\mathbf3}(x)
 =x+\frac{x^2}{2}-\frac{5x^4}{24}-\frac{x^5}{8}+O(x^6),}
 \label{eq:YMA28-u3-five}\\
 \boxed{
 \zeta(\beta)
 =z+\frac32z^2-\frac{45}{8}z^4-\frac{81}{8}z^5+O(z^6).}
 \label{eq:YMA28-zeta-series}
\end{align}
There is no cubic term in either series.
\end{theorem}

\begin{proof}
Series division of \eqref{eq:YMA28-a3-five} by
\eqref{eq:YMA28-a1-five} gives
\[
 x+\frac{x^2}{2}
 +0x^3-\frac{5x^4}{24}-\frac{x^5}{8}+O(x^6).
\]
Substitute \(x=3z\) and divide by three.
\end{proof}

\begin{assessment}[Why \(\zeta\) is an absolute calibration row]
\label{ass:YMA28-zeta-calibration}
The normalized law alone does not determine an arbitrary coarse-only addition
to the effective action, as proved in Version~V8.  The coefficient
\(\zeta\), however, is not inferred from the normalized reduced-fibre law: it
is a Peter--Weyl coefficient of the unnormalized plaquette weight after one
specified vacuum normalization.  It is therefore an absolute partition
coordinate of the declared Wilson cylinder.
\end{assessment}

% =============================================================================
\section{The first cubical deformation of a planar Wilson sheet}
\label{sec:YMA28-first-bump}
% =============================================================================

Let \(D_C\) be a finite simply connected plaquette disk in one coordinate
plane of the four-dimensional cubical lattice, with boundary \(C\) and area
\begin{equation}
 A(C)=|D_C|.
 \label{eq:YMA28-area}
\end{equation}
Assume the disk does not meet a periodic image of itself.  This is automatic
for every fixed loop on a sufficiently large torus.

A triality sheet is a finite \(\mathbb Z_3\)-valued plaquette two-chain whose
boundary is the triality-one chain carried by \(C\).  This is the support
shadow of every nonzero character monomial: Haar invariance under the centre
forces triality conservation at every integrated link.

\begin{lemma}[Cubical first-deformation lemma]
\label{lem:YMA28-cubical-first-deformation}
Let \(\mathscr S\) be the support of a triality sheet with boundary
\(\partial D_C\).  Then either \(\mathscr S=D_C\), or
\begin{equation}
 \boxed{|\mathscr S|\ge A(C)+4.}
 \label{eq:YMA28-area-gap}
\end{equation}
Equality holds exactly when one plaquette \(p\in D_C\) is replaced by the
other five faces of one elementary three-cube incident with \(p\) in a
transverse direction.  There are exactly
\begin{equation}
 \boxed{4A(C)}
 \label{eq:YMA28-bump-count}
\end{equation}
such first deformations.
\end{lemma}

\begin{proof}
Subtract the planar chain \(D_C\).  The difference is a compactly supported
cubical two-cycle with \(\mathbb Z_3\) coefficients.  The cubical complex of
\(\mathbb R^4\) is contractible, so the cycle is the boundary of a finite
three-chain \(B\).

First remove every cancellable three-cycle from \(B\), and among the remaining
fillings choose one with minimum number of three-cells.  A three-cell whose
coordinate directions do not contain the plane of \(D_C\) has no planar face
which can cancel a face of \(D_C\); an outermost such cell contributes a
nonempty exposed boundary and cannot improve the support count.  Hence a
minimizer for the first area increase may be taken as a union of elementary
three-cubes obtained by extruding planar plaquettes in one of the two
transverse coordinate directions, with either transverse orientation.

Fix one transverse direction and one orientation.  At each positive height,
the occupied three-cubes have a finite planar footprint \(R_j\).  The lower
and upper horizontal faces replace equal numbers of planar faces; the net new
faces are the vertical faces over the edge boundary of the footprint.  The
same conclusion holds coefficientwise for the two nonzero elements of
\(\mathbb Z_3\).  Every nonempty finite planar polyomino has at least four
boundary edges, so every nontrivial extrusion increases the support by at
least four.  Contributions in different transverse directions have different
face orientations and cannot cancel this first exposed perimeter.

Equality forces one footprint, one height layer, one connected component, and
planar perimeter four.  The footprint is therefore a single plaquette and the
three-chain is one elementary cube.  Its boundary cancels that plaquette of
\(D_C\) and inserts the other five faces, increasing the support by four.
For each of the \(A(C)\) plaquettes there are two transverse coordinate
directions and two sides, giving \(4A(C)\) choices.
\end{proof}

\begin{remark}[Relation to the V27 exhaustive search]
\label{rem:YMA28-search-relation}
For the \(2\times2\) disk, \Cref{lem:YMA28-cubical-first-deformation} gives
sixteen first bumps, exactly the sixteen supports returned by the exhaustive
V27 recursion.  The V28 replay independently repeats the no-dangling search
for the \(1\times1\), \(1\times2\), \(1\times3\), and \(2\times2\)
rectangles and obtains respectively \(4,8,12,16\) first bumps.
\end{remark}

\begin{theorem}[Planar Wilson loop through the first nonplanar relative order]
\label{thm:YMA28-loop-first-bump}
Let
\(
 O_C=\frac13\operatorname{Re}\Tr U_C
\)
be the normalized fundamental writer of a fixed embedded planar rectangle.
For real positive \(\beta\) sufficiently small,
\begin{equation}
 \boxed{
 \mathbb E_\beta O_C
 =\zeta(\beta)^{A(C)}
  +4A(C)\zeta(\beta)^{A(C)+4}
  +O_C\!\left(\zeta(\beta)^{A(C)+5}\right).}
 \label{eq:YMA28-loop-bump-expansion}
\end{equation}
Equivalently,
\begin{equation}
 \boxed{
 \mathbb E_\beta O_C
 =\zeta^{A(C)}
  \left(1+4A(C)\zeta^4+O_C(\zeta^5)\right).}
 \label{eq:YMA28-loop-bump-relative}
\end{equation}
\end{theorem}

\begin{proof}
Expand the normalized plaquette weights in irreducible characters.  On the
minimal disk every internal edge has incidence two.  Schur orthogonality
forces one common irreducible label, and the fundamental boundary insertion
selects \(\mathbf3\) or its conjugate according to orientation.  The disk
gluing identity of Version~V25 therefore gives exactly one factor
\(u_{\mathbf3}/3=\zeta\) per occupied plaquette, producing
\(\zeta^{A(C)}\).

A different nonzero monomial has a triality support with the same boundary.
By \Cref{lem:YMA28-cubical-first-deformation}, the first possible support has
area \(A(C)+4\), and equality gives exactly the \(4A(C)\) cube bumps.  Every
such support is again a simply connected disk whose internal links have
incidence two, so each contributes \(\zeta^{A(C)+4}\).  Every remaining
connected source support has total Taylor degree at least \(A(C)+5\), hence is
\(O_C(\zeta^{A(C)+5})\) because \(\zeta\) is a local analytic coordinate at
\(\beta=0\).

A disconnected closed surface is a vacuum factor and cancels between the
Wilson-loop numerator and the partition denominator.  The smallest closed
cubical surface has six plaquettes, so it cannot modify either displayed
coefficient before that cancellation.  Absolute convergence of the inherited
strong-coupling expansion controls the remainder.
\end{proof}

\begin{corollary}[Physical Creutz selector through the cube-bump order]
\label{cor:YMA28-Creutz-tension}
For the fixed planar Creutz quartet of Version~V25, whose signed area second
difference is one,
\begin{equation}
 \boxed{
 \chi(\beta)
 =-\log\zeta(\beta)-4\zeta(\beta)^4+O(\zeta(\beta)^5).}
 \label{eq:YMA28-Creutz-zeta}
\end{equation}
\end{corollary}

\begin{proof}
Taking logarithms in \eqref{eq:YMA28-loop-bump-relative} gives
\[
 \log\mathbb E_\beta O_C
 =A(C)\log\zeta+4A(C)\zeta^4+O_C(\zeta^5).
\]
Insert the four expressions into the signed Creutz logarithm.  The signed area
second difference is one, so both area-proportional terms survive with one
copy and all fixed-loop remainders remain \(O(\zeta^5)\).
\end{proof}

% =============================================================================
\section{The fundamental activity map and its exact area coordinate}
\label{sec:YMA28-Boettcher}
% =============================================================================

For one elementary straight-dyadic coarse plaquette, let
\(\mathcal F(\zeta)\) denote the dimension-normalized fundamental character
coefficient of its normalized boundary weight.  This is a scalar projection
of the complete blocked action, fixed before its value is read.

\begin{theorem}[First four-dimensional correction to the area map]
\label{thm:YMA28-activity-map}
Near \(\zeta=0\),
\begin{equation}
 \boxed{
 \mathcal F(\zeta)
 =\zeta^4+16\zeta^8+O(\zeta^9).}
 \label{eq:YMA28-activity-map}
\end{equation}
The first term is the flat \(2\times2\) disk.  The coefficient sixteen is the
complete nonplanar support correction at total degree eight.
\end{theorem}

\begin{proof}
The normalized fundamental boundary activity of an oriented disk of area
four is \(\zeta^4\) by the exact disk-gluing formula.  Version~V27 proves that
there is no other support through seven faces and exactly sixteen supports at
eight faces, all elementary cube bumps.  Each bump is an area-eight
fundamental disk and contributes \(\zeta^8\).  A different support contains at
least nine occupied plaquettes and therefore contributes \(O(\zeta^9)\).
Disconnected vacuum components first affect the normalized boundary activity
only after multiplication by the four-face source and hence do not change the
displayed degree-eight coefficient.
\end{proof}

\begin{theorem}[B\"ottcher character coordinate]
\label{thm:YMA28-Boettcher}
There is a unique real-analytic germ
\begin{equation}
 \Phi(\zeta)=\zeta+O(\zeta^2)
 \label{eq:YMA28-Phi-normalization}
\end{equation}
with
\begin{equation}
 \boxed{
 \Phi(\mathcal F(\zeta))=\Phi(\zeta)^4.}
 \label{eq:YMA28-Boettcher-equation}
\end{equation}
Its first nontrivial coefficient is
\begin{equation}
 \boxed{
 \Phi(\zeta)=\zeta+4\zeta^5+O(\zeta^6).}
 \label{eq:YMA28-Phi-series}
\end{equation}
For positive sufficiently small \(\zeta\), define the area coordinate
\begin{equation}
 \boxed{\mathcal U(\zeta):=-\log\Phi(\zeta).}
 \label{eq:YMA28-U-definition}
\end{equation}
Then
\begin{equation}
 \boxed{
 \mathcal U(\mathcal F(\zeta))=4\mathcal U(\zeta).}
 \label{eq:YMA28-U-scaling}
\end{equation}
\end{theorem}

\begin{proof}
Write
\(
 \mathcal F(\zeta)=\zeta^4a(\zeta)
\)
with \(a(0)=1\).  On a sufficiently small complex disk choose the analytic
branch of \(\log a\), and put
\begin{equation}
 h(\zeta)
 =\sum_{n=0}^{\infty}4^{-n-1}
   \log a\!\left(\mathcal F^{\circ n}(\zeta)\right).
 \label{eq:YMA28-h-series}
\end{equation}
The iterates converge to zero superexponentially, so the series converges
normally.  It satisfies
\[
 4h(\zeta)-h(\mathcal F(\zeta))=\log a(\zeta).
\]
Therefore \(\Phi(\zeta)=\zeta e^{h(\zeta)}\) obeys
\eqref{eq:YMA28-Boettcher-equation}.  If two normalized germs obey the same
equation, their quotient is invariant under \(\mathcal F\) and tends to one
under iteration to zero; hence the quotient is one.

From \eqref{eq:YMA28-activity-map},
\(
 a(\zeta)=1+16\zeta^4+O(\zeta^5)
\).
Only the \(n=0\) term of \eqref{eq:YMA28-h-series} contributes through degree
four, so \(h(\zeta)=4\zeta^4+O(\zeta^5)\).  This gives
\eqref{eq:YMA28-Phi-series}.  Taking the negative logarithm of
\eqref{eq:YMA28-Boettcher-equation} proves
\eqref{eq:YMA28-U-scaling}.
\end{proof}

\begin{corollary}[Bare-coupling series of the area coordinate]
\label{cor:YMA28-U-beta-series}
With \(z=\beta/18\),
\begin{align}
 \boxed{
 \Phi(\zeta(\beta))
 =z+\frac32z^2-\frac{45}{8}z^4-\frac{49}{8}z^5+O(z^6),}
 \label{eq:YMA28-Phi-z}\\
 \boxed{
 \mathcal U(\beta)
 =-\log z-\frac32z+\frac98z^2+\frac92z^3
  -\frac{67}{64}z^4+O(z^5).}
 \label{eq:YMA28-U-z}
\end{align}
Moreover, the physical Creutz response satisfies
\begin{equation}
 \boxed{
 \chi(\beta)=\mathcal U(\beta)+O(z^5).}
 \label{eq:YMA28-Creutz-U}
\end{equation}
\end{corollary}

\begin{proof}
Substitute \eqref{eq:YMA28-zeta-series} in
\eqref{eq:YMA28-Phi-series}.  Since
\(
 4\zeta^5=4z^5+O(z^6)
\), the fifth coefficient changes from \(-81/8\) to \(-49/8\), proving
\eqref{eq:YMA28-Phi-z}.  Expanding
\(-\log(\Phi/z)\) gives \eqref{eq:YMA28-U-z}.  Finally,
\eqref{eq:YMA28-Creutz-zeta} and
\(
 \mathcal U=-\log\zeta-4\zeta^4+O(\zeta^5)
\)
coincide through the displayed order.
\end{proof}

\begin{assessment}[What the exact fourth-power law does not say]
\label{ass:YMA28-projected-scope}
Equation~\eqref{eq:YMA28-U-scaling} is exact for the fundamental character
projection \(\mathcal F\).  It is not an autonomous recursion for the complete
four-dimensional action.  At the same degree at which the coefficient sixteen
appears, independent local character directions and potentially connected
multi-coarse-plaquette directions are generated.  Iterating
\(\mathcal F\) while discarding those directions defines a projected
one-coupling scheme, not the exact Wilson renormalization trajectory.
\end{assessment}

% =============================================================================
\section{The degree-eight local action leaves the Wilson line}
\label{sec:YMA28-order-eight-nonclosure}
% =============================================================================

The preceding activity coordinate is sufficient for area calibration.  It is
not sufficient for the complete action.  The first exact nonclosure can be
seen without enumerating every degree-eight multi-plaquette Hessian.

Put
\begin{equation}
 F=\chi_{\mathbf3}+\chi_{\overline{\mathbf3}},
 \qquad
 G=\chi_{\mathbf6}+\chi_{\overline{\mathbf6}},
 \qquad
 A=\chi_{\mathbf8}.
 \label{eq:YMA28-FGA}
\end{equation}
The representation-ring identity is
\begin{equation}
 \boxed{F^2=G+F+2+2A.}
 \label{eq:YMA28-F-square}
\end{equation}

\begin{lemma}[Leading sextet and adjoint disk activities]
\label{lem:YMA28-six-eight-activities}
The one-plaquette activities begin as
\begin{equation}
 u_{\mathbf6}=u_{\overline{\mathbf6}}
 =\frac{x^2}{2}+O(x^3),
 \qquad
 u_{\mathbf8}=x^2+O(x^3).
 \label{eq:YMA28-u6-u8}
\end{equation}
Consequently the flat four-face disk contributes
\begin{equation}
 \boxed{
 x^8\left[
  \frac1{3456}G+\frac1{512}A
 \right]}
 \label{eq:YMA28-raw-six-eight}
\end{equation}
to the normalized boundary partition function at degree eight.
\end{lemma}

\begin{proof}
The tensor products
\(
 \mathbf3\otimes\mathbf3
 =\mathbf6\oplus\overline{\mathbf3}
\)
and
\(
 \mathbf3\otimes\overline{\mathbf3}
 =\mathbf1\oplus\mathbf8
\), together with their conjugates, give one sextet at second tensor degree
and two adjoints.  Division by \(2!\) gives
\eqref{eq:YMA28-u6-u8}; vacuum normalization changes neither leading term.
Disk gluing contributes \(d_r^{1-4}\).  Hence
\[
 \frac{(x^2/2)^4}{6^3}=\frac{x^8}{3456},
 \qquad
 \frac{x^8}{8^3}=\frac{x^8}{512}.
\]
\end{proof}

\begin{theorem}[Exact local degree-eight non-Wilson coordinates]
\label{thm:YMA28-local-nonclosure}
In the one-coarse-plaquette character-support component of the effective
action, the degree-eight part orthogonal to the constant and fundamental
Wilson line is
\begin{equation}
 \boxed{
 x^8\left[
  \frac{37}{93312}
   \left(\chi_{\mathbf6}+\chi_{\overline{\mathbf6}}\right)
  -\frac{217}{373248}\chi_{\mathbf8}
 \right].}
 \label{eq:YMA28-order-eight-nonwilson}
\end{equation}
In particular, exact shell integration leaves the one-parameter Wilson action
line at degree eight.
\end{theorem}

\begin{proof}
Through degree four the normalized boundary partition function has
\[
 1+\frac{x^4}{27}F+O(x^5).
\]
At degree eight, \Cref{lem:YMA28-six-eight-activities} supplies the raw
sextet and adjoint terms.  The sixteen cube bumps have boundary triality one
and contribute only to the fundamental/antifundamental sector, so they do not
enter the present projection.  Expanding the negative logarithm contributes
\[
 \frac12\left(\frac{x^4}{27}F\right)^2
\]
in addition to the negatives of the raw degree-eight terms.  By
\eqref{eq:YMA28-F-square}, the sextet-pair coefficient is
\[
 \frac1{2\cdot27^2}-\frac1{3456}
 =\frac{37}{93312},
\]
and the adjoint coefficient is
\[
 \frac1{27^2}-\frac1{512}
 =-\frac{217}{373248}.
\]
Both are nonzero.  Since the three real class functions
\(F,G,A\) are linearly independent, no change of one fundamental Wilson
coefficient removes these two directions.
\end{proof}

\begin{corollary}[Signed nonfundamental contribution to the local Hessian]
\label{cor:YMA28-nonfund-Hessian}
Normalize \(X\in\mathfrak{su}(3)\) by
\(-\Tr_{\mathbf3}X^2=1\).  The order-eight sextet/adjoint sector contributes
\begin{equation}
 \boxed{
 -\frac{89}{186624}x^8\,\omega(q)
 =-\frac{89}{313456656384}\beta^8\,\omega(q)}
 \label{eq:YMA28-nonfund-Hessian}
\end{equation}
to the central axial Ward symbol.
\end{corollary}

\begin{proof}
The Dynkin-index ratios relative to the fundamental are five for
\(\mathbf6\) and six for \(\mathbf8\).  Hence the second derivative of
\(G\) is \(-10\) and that of \(A\) is \(-6\) in the stated normalization.
Apply these factors to \eqref{eq:YMA28-order-eight-nonwilson}:
\[
 -10\frac{37}{93312}
 -6\left(-\frac{217}{373248}\right)
 =-\frac{89}{186624}.
\]
The axial plaquette curl contributes the Ward factor \(\omega(q)\), and
\(x^8=\beta^8/6^8\).
\end{proof}

\begin{firewall}[A signed component is not the full order-eight sign]
\label{fire:YMA28-order-eight-sign}
The negative number in \eqref{eq:YMA28-nonfund-Hessian} is one exact local
component.  The complete order-eight axial coefficient must additionally
assemble the fundamental flat-disk term, all sixteen fundamental cube bumps,
the complete logarithmic fundamental term, and every connected
multi-coarse-plaquette Hessian kernel.  No sign for that assembled coefficient
is inferred here.
\end{firewall}

% =============================================================================
\section{Large-response selected roots in the area coordinate}
\label{sec:YMA28-large-root}
% =============================================================================

Let \(u>0\) be the fixed planar Creutz response value, and let
\(\beta_\chi(u)\) denote the unique strong-coupling solution for sufficiently
large \(u\) supplied by Version~V25.  Put
\begin{equation}
 y=e^{-u}.
 \label{eq:YMA28-y}
\end{equation}

\begin{theorem}[Five-term large-response root]
\label{thm:YMA28-large-root}
As \(u\to\infty\),
\begin{equation}
 \boxed{
 \begin{aligned}
 \frac{\beta_\chi(u)}{18}
 ={}&y-\frac32y^2+\frac92y^3-\frac{45}{4}y^4
     +\frac{211}{8}y^5+O(y^6).
 \end{aligned}}
 \label{eq:YMA28-large-root-series}
\end{equation}
The root is simple.  Moreover,
\begin{equation}
 \boxed{
 \zeta(\beta_\chi(u))
 =y-4y^5+O(y^6).}
 \label{eq:YMA28-zeta-at-root}
\end{equation}
\end{theorem}

\begin{proof}
First solve the exact calibration equation
\(
 \mathcal U(\beta)=u
\).
By \eqref{eq:YMA28-U-z}, substitution of
\[
 z=y+Ay^2+By^3+Cy^4+Dy^5
\]
and coefficient comparison gives successively
\[
 A=-\frac32,
 \qquad
 B=\frac92,
 \qquad
 C=-\frac{45}{4},
 \qquad
 D=\frac{211}{8}.
\]
Equivalently, this is the local inverse of
\(
 y=\Phi(\zeta(18z))
\).

By \eqref{eq:YMA28-Creutz-U}, the physical response differs from
\(\mathcal U\) by \(O(z^5)\).  Both derivatives equal
\(-z^{-1}+O(1)\), so their magnitudes are bounded below by
\(1/(2z)\) on a sufficiently small positive interval.  The mean-value theorem
therefore moves the root by at most \(O(z^6)=O(y^6)\).  This proves
\eqref{eq:YMA28-large-root-series}; simplicity is inherited from the same
derivative floor.  Substitution in \eqref{eq:YMA28-zeta-series} gives
\eqref{eq:YMA28-zeta-at-root}.
\end{proof}

\begin{corollary}[Area-covariant selector scaling]
\label{cor:YMA28-area-scaling}
On the fundamental activity projection,
\begin{equation}
 \boxed{
 \mathcal U_{\rm out}=4\mathcal U_{\rm in}.}
 \label{eq:YMA28-exact-area-scaling}
\end{equation}
For the physical planar Creutz response on the same strong-coupling branch,
\begin{equation}
 \boxed{
 \chi_{\rm out}=4\chi_{\rm in}
  +O\!\left(e^{-5\chi_{\rm in}}\right).}
 \label{eq:YMA28-Creutz-area-scaling}
\end{equation}
Thus using the same large numerical response target at two adjacent dyadic
areas is not an area-covariant root prescription.
\end{corollary}

\begin{proof}
The first identity is \eqref{eq:YMA28-U-scaling}.  At the incoming point,
\(\chi=\mathcal U+O(\zeta^5)\); at the outgoing point,
\(\zeta_{\rm out}=O(\zeta^4)\), so its corresponding error is
\(O(\zeta^{20})\).  Therefore
\[
 \chi_{\rm out}-4\chi_{\rm in}=O(\zeta^5)
 =O(e^{-5\chi_{\rm in}}).
\]
\end{proof}

\begin{assessment}[The response value is a scale-bearing calibration]
\label{ass:YMA28-response-calibration}
The finite root theorem of Version~V10 remains exact for any predeclared
computable \(u_*\).  The present result adds that, on the strong-coupling
area branch, \(u_*\) is not a scale-free number: the character-tension
coordinate multiplies by four under one \(2\times2\) area block.  A fixed
response value at every cutoff defines a legitimate scheme, but it is not the
same scheme as transporting one area density.  The intended convention must
be declared before root escape or adjacent-cutoff alignment is interpreted.
\end{assessment}

% =============================================================================
\section{Incoming versus outgoing relative margins}
\label{sec:YMA28-scale-margins}
% =============================================================================

Let \(\beta_{\square}(\beta)\) be the coefficient of the fundamental Wilson
plaquette interaction in the canonical one-plaquette support projection of the
blocked action.  Let
\begin{equation}
 \beta_{\rm H}(\beta,q)=3g_{M,\beta}^{\rm W}(q).
 \label{eq:YMA28-beta-H}
\end{equation}
Versions~V25--V27 give the same degree-seven series for these two quantities.

\begin{theorem}[Calibration triangle at strong coupling]
\label{thm:YMA28-calibration-triangle}
Uniformly in the finite periodic volume and central axial momentum,
\begin{align}
 \boxed{
 \beta_{\square}(\beta)
 =18\zeta(\beta)^4+O(\zeta(\beta)^8),}
 \label{eq:YMA28-beta-square-zeta}\\
 \boxed{
 \beta_{\rm H}(\beta,q)
 =18\zeta(\beta)^4+O(\zeta(\beta)^8),}
 \label{eq:YMA28-beta-H-zeta}\\
 \boxed{
 \frac{\beta_{\square}(\beta)}{\beta}
 =\frac{\beta^3}{5832}
  \left(1+\frac\beta3+\frac{\beta^2}{24}
        -\frac{\beta^3}{648}+O(\beta^4)\right),}
 \label{eq:YMA28-out-in-ratio}\\
 \boxed{
 \frac{\beta_{\rm H}(\beta,q)}{\beta_{\square}(\beta)}
 =1+O(\beta^4).}
 \label{eq:YMA28-H-square-ratio}
\end{align}
At the large-response root of \Cref{thm:YMA28-large-root},
\begin{equation}
 \boxed{
 \begin{aligned}
 \beta_{\square}(\beta_\chi(u))
 &=18e^{-4u}\left(1+O(e^{-4u})\right),\\
 g_{M,\beta_\chi(u)}^{\rm W}(q)
 &=6e^{-4u}\left(1+O(e^{-4u})\right).
 \end{aligned}}
 \label{eq:YMA28-root-outgoing-scales}
\end{equation}
\end{theorem}

\begin{proof}
The flat fundamental disk gives exactly
\(
 \beta_{\square}^{\rm disk}=18\zeta^4
\).
Version~V27 proves that the first additional support and the first logarithmic
competition occur at total degree eight, yielding
\eqref{eq:YMA28-beta-square-zeta}.  Its degree-seven axial calculation gives
\(
 3g=\beta_{\square}+O(\beta^8)
\), which is \eqref{eq:YMA28-beta-H-zeta} because
\(\zeta\asymp\beta\).

Factoring the exact degree-seven series
\[
 \beta_{\square}
 =\frac{\beta^4}{5832}
  \left(1+\frac\beta3+\frac{\beta^2}{24}
        -\frac{\beta^3}{648}+O(\beta^4)\right)
\]
proves \eqref{eq:YMA28-out-in-ratio}; division of the two common leading
series proves \eqref{eq:YMA28-H-square-ratio}.  Finally,
\eqref{eq:YMA28-zeta-at-root} gives
\(
 \zeta^4=e^{-4u}(1+O(e^{-4u}))
\), which proves \eqref{eq:YMA28-root-outgoing-scales}.
\end{proof}

Recall
\begin{equation}
 r_0(q)=\frac{s_0(q)}{\omega(q)},
 \qquad
 \frac5{28}\le r_0(q)\le\frac13.
 \label{eq:YMA28-r0-range}
\end{equation}
The inherited relative margin uses the incoming parameter \(\beta\).  A
scale-matched outgoing comparison would instead use
\(\beta_{\square}\).  V28 keeps these as two separately named targets.

\begin{theorem}[Strong-coupling calibration phase diagram]
\label{thm:YMA28-margin-phase-diagram}
For every fixed nonzero central axial momentum,
\begin{align}
 \boxed{
 \frac{h_{M,\beta_\chi(u)}^{\rm W}(q)}
      {\beta_\chi(u)s_0(q)}
 =\frac{e^{-3u}}{3r_0(q)}
  \left(1+O(e^{-u})\right)
 \longrightarrow0,}
 \label{eq:YMA28-incoming-root-ratio}\\
 \boxed{
 \frac{h_{M,\beta_\chi(u)}^{\rm W}(q)}
      {\beta_{\square}(\beta_\chi(u))s_0(q)}
 =\frac1{3r_0(q)}+O(e^{-4u}).}
 \label{eq:YMA28-outgoing-root-ratio}
\end{align}
The limiting outgoing ratio obeys
\begin{equation}
 \boxed{
 1\le\frac1{3r_0(q)}\le\frac{28}{15}.}
 \label{eq:YMA28-outgoing-ratio-range}
\end{equation}
Consequently:
\begin{enumerate}[label=\textup{(M\arabic*)},leftmargin=4em]
\item for every fixed incoming fraction \(\delta_{\rm in}>0\),
      \[
       h^{\rm W}-\delta_{\rm in}\beta s_0<0
      \]
      at all sufficiently large response targets;
\item for every outgoing fraction \(0\le\delta_{\rm out}<1\),
      \[
       h^{\rm W}-\delta_{\rm out}\beta_{\square}s_0>0
      \]
      at all sufficiently large response targets, uniformly over the central
      nonzero axial angles;
\item for a fixed angle, the outgoing borderline is
      \(
       \delta_{\rm out}=1/(3r_0(q))
      \); equality requires the complete order-eight and higher card.
\end{enumerate}
\end{theorem}

\begin{proof}
Use
\(
 h=\omega g
\), \eqref{eq:YMA28-large-root-series}, and
\eqref{eq:YMA28-root-outgoing-scales}.  The incoming denominator is
\[
 \beta_\chi s_0
 =18e^{-u}\omega r_0(q)(1+O(e^{-u})),
\]
while the numerator is
\(
 6e^{-4u}\omega(1+O(e^{-4u}))
\).  This proves \eqref{eq:YMA28-incoming-root-ratio}.  Dividing instead by
\(\beta_{\square}\omega r_0\) proves
\eqref{eq:YMA28-outgoing-root-ratio}.  The range
\eqref{eq:YMA28-outgoing-ratio-range} is the reciprocal image of
\eqref{eq:YMA28-r0-range}.  Its minimum is one, so every
\(\delta_{\rm out}<1\) has a uniform positive separation for sufficiently
large \(u\).  The remaining sign statements follow from the two limits.
\end{proof}

\begin{assessment}[Qualification of the V25 relative obstruction]
\label{ass:YMA28-V25-qualification}
Version~V25 correctly proves a typed obstruction to the inherited condition
\(
 h\succeq\delta\beta S_0
\)
near \(\beta=0\).  V28 does not withdraw that theorem.  It proves that the
theorem compares an outgoing blocked Hessian with the incoming bare coupling.
When the declared target is instead normalized by the outgoing fundamental
plaquette coefficient, the same strong-coupling card has an order-one positive
ratio.  Choosing between these comparisons is a physical calibration decision,
not an algebraic simplification.
\end{assessment}

% =============================================================================
\section{Finite replay, integrated theorem, and revised stopping boundary}
\label{sec:YMA28-replay}
% =============================================================================

The accompanying exact replay performs the following independent checks:
\begin{enumerate}[label=\textup{(R28.\arabic*)},leftmargin=5.5em]
\item the degree-five trivial and fundamental \(\SU(3)\) multiplicities;
\item the series \eqref{eq:YMA28-u3-five}--\eqref{eq:YMA28-zeta-series};
\item exhaustive first-bump searches for planar rectangles of areas one through
      four, obtaining \(4A\) bumps in each case;
\item the coefficient four in the B\"ottcher germ and every coefficient in
      \eqref{eq:YMA28-Phi-z}--\eqref{eq:YMA28-large-root-series};
\item the exact sextet, adjoint, and nonfundamental Hessian coefficients in
      \eqref{eq:YMA28-order-eight-nonwilson}--\eqref{eq:YMA28-nonfund-Hessian};
\item the identity \(\beta_{\square}=3g\) through degree seven and the
      cancellation of the \(e^{-5u},e^{-6u},e^{-7u}\) terms after substitution
      of the calibrated root series;
\item the exact outgoing ratio interval \([1,28/15]\).
\end{enumerate}
Normal and optimized executions produce byte-identical JSON.

\begin{theorem}[Version V28 calibration and first-nonclosure theorem]
\label{thm:YMA28-integrated}
Preserving Versions~V1--V27, Version~V28 proves:
\begin{enumerate}[label=\textup{(V28.\arabic*)},leftmargin=5.8em]
\item the dimension-normalized fundamental activity has the exact fifth-order
      series \eqref{eq:YMA28-zeta-series};
\item a nonplanar triality sheet with the same planar boundary costs at least
      four additional plaquettes, with equality exactly for one cube bump;
\item every planar Wilson rectangle has the first-bump expansion
      \eqref{eq:YMA28-loop-bump-expansion}, and the selected Creutz response is
      \(-\log\zeta-4\zeta^4+O(\zeta^5)\);
\item the elementary coarse fundamental activity map is
      \(\zeta^4+16\zeta^8+O(\zeta^9)\) and admits the unique exact area
      coordinate \(\mathcal U\);
\item the physical large-response root has the five-term expansion
      \eqref{eq:YMA28-large-root-series};
\item the local action develops nonzero sextet and adjoint coordinates at
      degree eight, so the scalar Wilson line is not closed there;
\item the incoming, outgoing, and Hessian couplings obey the calibration
      triangle \eqref{eq:YMA28-beta-square-zeta}--\eqref{eq:YMA28-H-square-ratio};
\item the incoming relative margin collapses exponentially along the
      large-response root, while the outgoing-calibrated ratio tends to the
      explicit interval \([1,28/15]\).
\end{enumerate}
No item populates the response-matched selected root at its predeclared
numerical target or proves local-vacuum contraction.
\end{theorem}

\begin{proof}
Items~\textup{(V28.1)--(V28.3)} are
\Cref{lem:YMA28-fifth-multiplicities,thm:YMA28-zeta-series,lem:YMA28-cubical-first-deformation,thm:YMA28-loop-first-bump,cor:YMA28-Creutz-tension}.  Items~\textup{(V28.4)--(V28.5)} are
\Cref{thm:YMA28-activity-map,thm:YMA28-Boettcher,cor:YMA28-U-beta-series,thm:YMA28-large-root,cor:YMA28-area-scaling}.
Item~\textup{(V28.6)} is
\Cref{lem:YMA28-six-eight-activities,thm:YMA28-local-nonclosure,cor:YMA28-nonfund-Hessian}.  Items~\textup{(V28.7)--(V28.8)} are
\Cref{thm:YMA28-calibration-triangle,thm:YMA28-margin-phase-diagram}.
\end{proof}

\begingroup\small
\begin{longtable}{>{\raggedright\arraybackslash}p{0.27\textwidth}
                  >{\raggedright\arraybackslash}p{0.18\textwidth}
                  >{\raggedright\arraybackslash}p{0.47\textwidth}}
\caption{Version V28 proof-status ledger.}\label{tab:YMA28-ledger}\\
\toprule
\textbf{Row}&\textbf{Status}&\textbf{Exact scope}\\
\midrule
\endfirsthead
\toprule
\textbf{Row}&\textbf{Status}&\textbf{Exact scope}\\
\midrule
\endhead
V1--V27 prefix
& Preserved
& Complete preceding source retained byte for byte before the terminal command.\\[0.35em]
Fundamental character calibration
& \pass
& Exact \(u_{\mathbf3}\) and \(\zeta=u_{\mathbf3}/3\) series through fifth order.\\[0.35em]
Planar first-bump theorem
& \pass
& The first nonplanar triality sheet costs four faces; exactly four cube bumps occur per planar tile.\\[0.35em]
Creutz first-bump response
& \pass
& \(\chi=-\log\zeta-4\zeta^4+O(\zeta^5)\) on the frozen planar quartet.\\[0.35em]
Character-tension coordinate
& \pass
& Unique analytic \(\Phi\) with \(\Phi\circ\mathcal F=\Phi^4\), and exact area coordinate \(\mathcal U=-\log\Phi\).\\[0.35em]
Large-response root
& \pass
& Five explicit terms with \(O(e^{-6u})\) remainder in the bare root coordinate.\\[0.35em]
Local degree-eight action closure
& \obstruction
& Nonzero sextet and adjoint coefficients prove that the one-parameter Wilson action line is not closed.\\[0.35em]
Complete order-eight axial sign
& \stopstatus
& Fundamental cube-bump, logarithmic, and multi-coarse-plaquette Hessian terms remain to be assembled.\\[0.35em]
Incoming bare-reference margin
& \obstruction
& Every fixed positive incoming fraction fails on the large-response strong-coupling root; the Hessian itself remains positive at leading order.\\[0.35em]
Outgoing-character calibration
& \pass
& Every fixed fraction below one passes asymptotically; the limiting modewise ratio lies in \([1,28/15]\).\\[0.35em]
Physical response convention
& \stopstatus
& The programme has not supplied the numerical \(u_*\) or declared whether it is transported as a fixed scheme value or an area coordinate.\\[0.35em]
Native response-matched selected-shell symbol
& \stopstatus
& The two response cofactors and complete same-action selected-root card remain unpopulated.\\[0.35em]
Local-vacuum or continuum mass
& Not inferred
& Requires the complete response-matched signed head, differentiated tail, conductance, ancestry, cutoff transport, and reflected continuum.\\
\bottomrule
\end{longtable}
\endgroup

% =============================================================================
\section{Revised V29 acquisition}
\label{sec:YMA28-next}
% =============================================================================

\begin{openproblem}[V29: freeze the calibration, then read one physical root]
\label{op:YMA28-V29}
Preserve the complete present source.  Proceed in this order.
\begin{enumerate}[label=\textup{(V29.\arabic*)},leftmargin=5.7em]
\item Declare whether the selected response value is a fixed scheme number,
      the transported character-tension coordinate \(\mathcal U\), or another
      physically calibrated area density.  Do not compare roots across cutoffs
      before this row is fixed.
\item Supply the predeclared numerical response value or interval.  If it lies
      on the large-response branch, apply
      \eqref{eq:YMA28-large-root-series} and refine it by the exact
      logarithm-free root evaluator of Version~V10.
\item Keep the incoming-bare and outgoing-character relative margins as
      separately named targets.  A pass for one is not reported as a pass for
      the other.
\item On the native response-matched branch, acquire the two physical response
      cofactors and load the complete same-action selected root before any
      pure-Wilson asymptotic is transported to it.
\item If a degree-eight calculation is required, acquire the complete generated
      coefficient vector, not only one scalar: the fundamental cube-bump and
      logarithmic terms, the sextet and adjoint coordinates, and every connected
      multi-coarse-plaquette Hessian term must be assembled together.
\item Evaluate one retained momentum--colour margin with its differentiated
      Ward tail and return \pass, an attained \obstruction, or \stopstatus{} at
      the first unresolved calibration, action, coefficient, or tail row.
\end{enumerate}
No further scalar coefficient ladder, boundary compiler, phase simplex, or
Witten-current atlas is opened before one of these physical root cards is
populated.
\end{openproblem}

\begin{assessment}[Version V28 termination]
\label{ass:YMA28-termination}
The apparent circle after Version~V27 came from treating an incoming bare
parameter, an outgoing character activity, and a physical Hessian scale as one
coupling.  V28 separates them and shows that the large-response selected root
has two different relative-margin verdicts depending on the declared
calibration.  It also identifies the exact degree at which the scalar Wilson
action ceases to close.  The next advance must therefore be a physically
calibrated root and complete same-action card, not another unlabelled Taylor
coefficient.
\end{assessment}

\section*{Version V28 append-only continuation record}
\addcontentsline{toc}{section}{Version V28 append-only continuation record}
The complete Version~V27 source preceding its terminal document-closing
command is preserved byte for byte.  The SHA--256 digest of the complete V27
source is
\begin{center}\scriptsize\ttfamily
 d7d8995901e16dd70dfe20edbe4d9f3ed07e32a5596a830a3f052d7d717ef686.
\end{center}
Its preterminal prefix contains \(1254540\) bytes and has SHA--256
\begin{center}\scriptsize\ttfamily
 b3702f5f54b848d8ea95014aed2ab4d422f174aea871e210a78b5681cb610b7f.
\end{center}
For Version~V29, preserve this accumulated source, append new material
immediately before the final active document-closing command, and increment
the filename to \texttt{\_V29.tex}.

% =============================================================================
% END APPENDED VERSION V28 MATERIAL
% =============================================================================


% =============================================================================
% START APPENDED VERSION V29 MATERIAL
% =============================================================================

\clearpage
\hypersetup{
 pdftitle={Renewal Yang--Mills Reduced-Fibre Wilson Shell Symbol and Local-Vacuum Programme V29},
 pdfsubject={Physical-rectangle calibration, exact Creutz-cell transport, and the strong-coupling/continuum separation}
}

\pdfinfo{
 /Title (Renewal Yang--Mills Reduced-Fibre Wilson Shell Symbol and Local-Vacuum Programme V29)
 /Subject (Physical-rectangle calibration, exact Creutz-cell transport, and the strong-coupling/continuum separation)
}

\section{Version V29: physical rectangles, exact Creutz-cell transport, and the strong-coupling/continuum fork}
\label{sec:YMA29-direction}

Version~V28 is preserved above.  It correctly separated the incoming bare
Wilson parameter, the normalized fundamental character activity, the outgoing
fundamental plaquette coefficient, and the coupling read from the axial
Hessian.  Its final mandate nevertheless still allowed a ``fixed scheme
number'', a transported character-tension coordinate, or a physical area
density to be chosen before the geometry of the four loop writers had been
made part of the calibration packet.  That geometry is logically earlier.

The Creutz response is a mixed finite difference in the two side lengths of a
rectangular Wilson loop.  Under an exact dyadic block, one coarse mixed
difference is not, in general, four times one fine mixed difference.  It is
the sum of the four fine Creutz cells which tile the corresponding rectangle
in loop-size space.  After division by the physical cell area, the coarse
response density is the arithmetic mean of those four fine densities.  The
scalar law
\(
 \mathcal U_{\rm out}=4\mathcal U_{\rm in}
\)
of Version~V28 is therefore the homogeneous zero-detail specialization of a
larger exact identity.

This observation changes the upstream order to
\begin{equation}
 \boxed{
 \begin{gathered}
 \text{lattice spacing and physical rectangle passport}\\
 \Downarrow\\
 \text{exact four-child Creutz response panel}\\
 \Downarrow\\
 \text{fixed lattice-cell, fixed physical-rectangle, or large-loop semantics}\\
 \Downarrow\\
 \text{scale-compatible response target and selected root}\\
 \Downarrow\\
 \text{native response cofactors and the same-action Wilson mode card}.
 \end{gathered}}
 \label{eq:YMA29-upstream-order}
\end{equation}
A fixed numerical response at a shrinking lattice cell remains a legitimate
finite-cutoff scheme.  It is not a finite physical area-density calibration:
its response per physical cell area diverges like the inverse square lattice
spacing.  Conversely, a finite physical response density requires a
dimensionless target tending to zero, and is asymptotically disjoint from the
large-response strong-coupling roots evaluated in Versions~V25--V28.

\begin{firewall}[Scope of Version V29]
\label{fire:YMA29-scope}
The results below are exact identities for rectangular loop expectations,
finite differences, exact dyadic pushforward, and the projected
character-tension coordinate.  The continuum mixed-derivative statements are
conditional on convergence of a declared renormalized loop free energy with
the displayed differentiability.  No such continuum Wilson-loop limit,
string tension, weak-facing-root escape, response-matched action, negative
mode, or mass gap is manufactured.
\end{firewall}

% =============================================================================
\section{The physical rectangle passport and its Creutz cell}
\label{sec:YMA29-rectangle-passport}
% =============================================================================

Fix a lattice spacing \(a>0\).  On a positive rectangular-loop branch, let
\begin{equation}
 W_a(m,n)>0,
 \qquad m,n\in\mathbb N_0,
 \qquad
 W_a(0,n)=W_a(m,0)=1,
 \label{eq:YMA29-loop-positive}
\end{equation}
be the normalized expectation of the oriented \(m\)-by-\(n\) rectangular
Wilson loop, and put
\begin{equation}
 F_a(m,n):=-\log W_a(m,n).
 \label{eq:YMA29-loop-free-energy}
\end{equation}
For \(m,n\ge1\), define the backward mixed difference
\begin{equation}
 \boxed{
 \chi_a(m,n)
 :=F_a(m,n)-F_a(m-1,n)-F_a(m,n-1)+F_a(m-1,n-1).}
 \label{eq:YMA29-Creutz-cell}
\end{equation}
Equivalently,
\begin{equation}
 e^{-\chi_a(m,n)}
 =\frac{W_a(m,n)W_a(m-1,n-1)}
        {W_a(m-1,n)W_a(m,n-1)}.
 \label{eq:YMA29-Creutz-ratio}
\end{equation}
The corresponding physical response density is
\begin{equation}
 \boxed{
 \Sigma_a(m,n):=\frac{\chi_a(m,n)}{a^2}.}
 \label{eq:YMA29-response-density}
\end{equation}
The denominator is the physical area of the elementary cell in the
rectangle-side parameter plane, not the area \(mn a^2\) enclosed by the
largest loop in the quartet.

\begin{definition}[Rectangle calibration passport]
\label{def:YMA29-rectangle-passport}
A rectangular response card at cutoff \(r\) carries
\begin{equation}
 \boxed{
 \mathfrak C_r^{\rm rect}
 =\bigl(a_r,m_r,n_r;\,m_ra_r,n_ra_r;\,u_r,\mathfrak s_r\bigr),}
 \label{eq:YMA29-passport}
\end{equation}
where \(a_r\) is the lattice spacing, \((m_r,n_r)\) is the largest loop in
the quartet, \(u_r\) is the declared dimensionless response target, and
\(\mathfrak s_r\) is one of the following semantic labels:
\begin{enumerate}[label=\textup{(P\arabic*)},leftmargin=3.8em]
\item \emph{fixed lattice cell}: \(m_r,n_r\) and usually \(u_r\) are fixed;
\item \emph{fixed physical rectangle}: \(m_ra_r\to\ell_1>0\) and
      \(n_ra_r\to\ell_2>0\), while \(u_r/a_r^2\) is the calibrated target;
\item \emph{large-loop area density}: the physical side lengths also tend to
      infinity and the target is the limiting mixed area coefficient.
\end{enumerate}
The semantic label and loop geometry are frozen before a root or mode value is
inspected.
\end{definition}

\begin{assessment}[Three targets which must not be identified]
\label{ass:YMA29-three-targets}
A fixed lattice quartet is a finite-cutoff ultraviolet response scheme.  A
fixed physical rectangle probes the mixed curvature of the loop free energy
at that finite physical size.  A string-tension statement additionally
requires a large-physical-rectangle limit.  Equality of these three numbers is
an additional theorem, not part of the definition of the Creutz ratio.
\end{assessment}

% =============================================================================
\section{Exact dyadic four-cell transport}
\label{sec:YMA29-four-cell}
% =============================================================================

The next theorem is purely algebraic and applies before any asymptotic or
strong-coupling approximation.

\begin{lemma}[Mixed-difference telescoping at arbitrary block factor]
\label{lem:YMA29-b-mixed-telescope}
Let \(G:\mathbb N_0^2\to\mathbb R\), and write
\begin{equation}
 (\Delta_1\Delta_2G)(k,\ell)
 :=G(k,\ell)-G(k-1,\ell)-G(k,\ell-1)+G(k-1,\ell-1).
 \label{eq:YMA29-unit-mixed}
\end{equation}
For every integer \(b\ge1\) and \(m,n\ge1\),
\begin{equation}
 \boxed{
 \begin{aligned}
 &G(bm,bn)-G(bm-b,bn)-G(bm,bn-b)+G(bm-b,bn-b)\\
 &\qquad
 =\sum_{i=0}^{b-1}\sum_{j=0}^{b-1}
   (\Delta_1\Delta_2G)(bm-i,bn-j).
 \end{aligned}}
 \label{eq:YMA29-b-telescope}
\end{equation}
\end{lemma}

\begin{proof}
For fixed \(\ell\), summing the first difference in the first variable gives
\[
 G(bm,\ell)-G(bm-b,\ell)
 =\sum_{i=0}^{b-1}
  \bigl[G(bm-i,\ell)-G(bm-i-1,\ell)\bigr].
\]
Apply this at \(\ell=bn\) and \(\ell=bn-b\), subtract, and telescope each
resulting first difference in the second variable.  The summand is precisely
\eqref{eq:YMA29-unit-mixed}.
\end{proof}

\begin{theorem}[Exact block transport of rectangular Creutz cells]
\label{thm:YMA29-exact-block-Creutz}
Assume an exact straight block of factor \(b\) maps every coarse rectangular
loop to the corresponding \(bm\)-by-\(bn\) fine loop.  Thus
\begin{equation}
 W_{ba}^{\rm out}(m,n)=W_a^{\rm in}(bm,bn).
 \label{eq:YMA29-exact-loop-pushforward}
\end{equation}
Then
\begin{equation}
 \boxed{
 \chi_{ba}^{\rm out}(m,n)
 =\sum_{i=0}^{b-1}\sum_{j=0}^{b-1}
   \chi_a^{\rm in}(bm-i,bn-j).}
 \label{eq:YMA29-exact-Creutz-pushforward}
\end{equation}
After physical-area normalization,
\begin{equation}
 \boxed{
 \Sigma_{ba}^{\rm out}(m,n)
 =\frac1{b^2}
   \sum_{i=0}^{b-1}\sum_{j=0}^{b-1}
   \Sigma_a^{\rm in}(bm-i,bn-j).}
 \label{eq:YMA29-density-average}
\end{equation}
For the dyadic shell \(b=2\), one coarse response density is exactly the
arithmetic mean of four neighboring fine response densities.
\end{theorem}

\begin{proof}
Take minus logarithms in \eqref{eq:YMA29-exact-loop-pushforward}, substitute
\(G=F_a^{\rm in}\) in \Cref{lem:YMA29-b-mixed-telescope}, and use
\eqref{eq:YMA29-Creutz-cell}.  Division by \((ba)^2=b^2a^2\) proves the
second identity.
\end{proof}

\begin{corollary}[Exact dyadic target compatibility]
\label{cor:YMA29-target-compatibility}
For one exact dyadic step, a declared parent response target and its four child
targets are area compatible only if
\begin{equation}
 \boxed{
 u_{\rm out}(m,n)
 =\sum_{i,j\in\{0,1\}}
   u_{\rm in}(2m-i,2n-j).}
 \label{eq:YMA29-target-cocycle}
\end{equation}
A homogeneous child target \(u_{\rm in}\) therefore transports as
\begin{equation}
 \boxed{u_{\rm out}=4u_{\rm in}.}
 \label{eq:YMA29-homogeneous-target}
\end{equation}
Using the same nonzero numerical target at both scales is a valid pair of
finite schemes, but it has the nonzero area-calibration defect
\begin{equation}
 \boxed{u_{\rm out}-4u_{\rm in}=-3u.}
 \label{eq:YMA29-fixed-u-defect}
\end{equation}
\end{corollary}

\begin{proof}
Equation \eqref{eq:YMA29-target-cocycle} is the root-level version of
\eqref{eq:YMA29-exact-Creutz-pushforward}.  The remaining statements are its
constant-target specialization.
\end{proof}

\begin{definition}[Four-child Creutz dispersion]
\label{def:YMA29-cell-dispersion}
For the four dyadic child densities \(\Sigma_1,\ldots,\Sigma_4\), put
\begin{equation}
 \overline\Sigma=\frac14\sum_{j=1}^4\Sigma_j,
 \qquad
 \boxed{
 \mathcal V_{\rm Cr}
 :=\frac14\sum_{j=1}^4(\Sigma_j-\overline\Sigma)^2\ge0.}
 \label{eq:YMA29-cell-variance}
\end{equation}
By \eqref{eq:YMA29-density-average}, \(\overline\Sigma\) is the parent
density.
\end{definition}

\begin{theorem}[Exact Creutz-cell Pythagoras and scalarization residual]
\label{thm:YMA29-cell-Pythagoras}
The four-child packet satisfies
\begin{equation}
 \boxed{
 \frac14\sum_{j=1}^4\Sigma_j^2
 =\overline\Sigma^{\,2}+\mathcal V_{\rm Cr}.}
 \label{eq:YMA29-cell-Pythagoras}
\end{equation}
Moreover,
\begin{equation}
 \boxed{
 |\Sigma_j-\overline\Sigma|
 \le\sqrt{3\mathcal V_{\rm Cr}}
 \qquad(1\le j\le4).}
 \label{eq:YMA29-cell-deviation}
\end{equation}
The dispersion vanishes exactly when all four child densities coincide.  On
that branch the exact four-cell transport reduces to the scalar homogeneous
law of \eqref{eq:YMA29-homogeneous-target}.
\end{theorem}

\begin{proof}
Expand the squares around the mean; the cross term vanishes because the four
deviations sum to zero.  If one deviation equals \(x\), the remaining three
sum to \(-x\), so Cauchy--Schwarz gives a sum of their squares at least
\(x^2/3\).  Hence
\[
 4\mathcal V_{\rm Cr}
 \ge x^2+\frac{x^2}{3}=\frac43x^2,
\]
which proves \eqref{eq:YMA29-cell-deviation}.  A finite sum of nonnegative
squares vanishes exactly when every summand does.
\end{proof}

\begin{countertheorem}[One fine Creutz cell does not determine its coarse parent]
\label{cth:YMA29-one-cell-no-parent}
There are two positive rectangular-loop tables with the same value on one
predeclared fine Creutz cell but different exact dyadic parent responses.
\end{countertheorem}

\begin{proof}
On a \(2\)-by-\(2\) cell array, let Card A have unit Creutz values on all four
cells.  Let Card B have value one on the predeclared upper-right cell and value
two on each of the other three.  For either array \(c_{ij}>0\), define
\begin{equation}
 F(m,n)=\sum_{1\le i\le m}\sum_{1\le j\le n}c_{ij},
 \qquad W(m,n)=e^{-F(m,n)}.
 \label{eq:YMA29-countertable}
\end{equation}
Then \(0<W\le1\), the mixed differences of \(F\) are exactly \(c_{ij}\),
and the chosen fine value is one in both cards.  The dyadic parent responses
are respectively four and seven by
\eqref{eq:YMA29-exact-Creutz-pushforward}.  This is a positive-table
calibration counterexample, not a claim that both tables are generated by the
Wilson action.
\end{proof}

\begin{corollary}[Summable dyadic cell detail gives one scalar calibration limit]
\label{cor:YMA29-summable-cell-detail}
Follow one nested chain of dyadic rectangle cells and let
\(\mathcal V_r\) be the four-child dispersion at the \(r\)-th parent.  If
\begin{equation}
 \boxed{
 \sum_r\sqrt{\mathcal V_r}<\infty,}
 \label{eq:YMA29-dispersion-summable}
\end{equation}
then the parent density and every predeclared descendant-child density differ
by a summable amount and hence have the same scalar limit whenever one of them
converges.  A persistent parent/child discrepancy forces failure of
\eqref{eq:YMA29-dispersion-summable}.
\end{corollary}

\begin{proof}
Apply \eqref{eq:YMA29-cell-deviation} at every step and sum the resulting
bounds.  The Cauchy criterion proves the conclusion.
\end{proof}

% =============================================================================
\section{Fixed physical rectangles and mixed loop-free-energy curvature}
\label{sec:YMA29-continuum-rectangle}
% =============================================================================

The exact finite-difference identity also determines what a fixed physical
rectangle would measure if a sufficiently regular continuum loop free energy
has been constructed.

\begin{theorem}[Creutz density is an exact cell average of mixed curvature]
\label{thm:YMA29-mixed-curvature-average}
Let \(\mathscr F\in C^2\) on a neighborhood of
\([x-a,x]\times[y-a,y]\), and suppose the lattice loop free energy on that
rectangle is sampled as
\begin{equation}
 F_a(m,n)=\mathscr F(ma,na),
 \qquad x=ma,
 \qquad y=na.
 \label{eq:YMA29-continuum-sampling}
\end{equation}
Then
\begin{equation}
 \boxed{
 \frac{\chi_a(m,n)}{a^2}
 =\frac1{a^2}
   \int_{x-a}^{x}\int_{y-a}^{y}
   \partial_1\partial_2\mathscr F(s,t)\,\dd t\,\dd s.}
 \label{eq:YMA29-cell-average-curvature}
\end{equation}
If \(\mathscr F\in C^4\) and
\(
 c=(x-a/2,y-a/2)
\)
is the cell center, then
\begin{equation}
 \boxed{
 \begin{aligned}
 \left|
 \frac{\chi_a(m,n)}{a^2}
 -\partial_1\partial_2\mathscr F(c)
 \right|
 \le\frac{a^2}{24}
 \bigl(&\|\partial_1^3\partial_2\mathscr F\|_{\infty,\square}\nobreak\\
      &+\|\partial_1\partial_2^3\mathscr F\|_{\infty,\square}\bigr),
 \end{aligned}}
 \label{eq:YMA29-midpoint-error}
\end{equation}
where the norms are over the displayed cell.
\end{theorem}

\begin{proof}
Apply the fundamental theorem of calculus first in the first variable and then
in the second:
\[
 \begin{aligned}
 &\mathscr F(x,y)-\mathscr F(x-a,y)
  -\mathscr F(x,y-a)+\mathscr F(x-a,y-a)\\
 &\qquad
 =\int_{x-a}^{x}\int_{y-a}^{y}
   \partial_1\partial_2\mathscr F(s,t)\,\dd t\,\dd s.
 \end{aligned}
\]
This is \eqref{eq:YMA29-cell-average-curvature}.  Put
\(f=\partial_1\partial_2\mathscr F\).  The one-dimensional midpoint rule for
an interval average has error at most \(a^2\|f''\|_\infty/24\).  Apply it in
one variable while averaging the other, then apply it in the second variable.
The two errors add and give \eqref{eq:YMA29-midpoint-error}.
\end{proof}

\begin{corollary}[Fixed lattice and fixed physical rectangles have different limits]
\label{cor:YMA29-fixed-lattice-vs-physical}
Assume the hypotheses of \Cref{thm:YMA29-mixed-curvature-average} hold along a
refining sequence \(a_r\to0\).
\begin{enumerate}[label=\textup{(R\arabic*)},leftmargin=3.8em]
\item If \(m_r=m\) and \(n_r=n\) are fixed, the physical rectangle shrinks to
      the origin and the response density tends to
      \(\partial_1\partial_2\mathscr F(0,0)\), when that derivative extends
      continuously there.
\item If \(m_ra_r\to\ell_1>0\) and \(n_ra_r\to\ell_2>0\), then
      \begin{equation}
       \boxed{
       \Sigma_{a_r}(m_r,n_r)
       \longrightarrow
       \partial_1\partial_2\mathscr F(\ell_1,\ell_2).}
       \label{eq:YMA29-fixed-rectangle-limit}
      \end{equation}
\end{enumerate}
Neither limit is, by itself, a large-loop string tension.
\end{corollary}

\begin{proof}
The cell centers have the indicated limits.  Use
\eqref{eq:YMA29-midpoint-error} and continuity of the mixed derivative.
\end{proof}

\begin{theorem}[Exact cancellation of perimeter and one-side terms]
\label{thm:YMA29-perimeter-cancellation}
Suppose a continuum rectangular-loop free energy has the decomposition
\begin{equation}
 \mathscr F(\ell_1,\ell_2)
 =\sigma\ell_1\ell_2
  +p_1(\ell_1)+p_2(\ell_2)+c
  +\mathcal R(\ell_1,\ell_2).
 \label{eq:YMA29-free-energy-decomposition}
\end{equation}
Every discrete Creutz cell cancels \(p_1,p_2,c\) exactly, and
\begin{equation}
 \boxed{
 \frac{\chi_a}{a^2}
 =\sigma
 +\frac1{a^2}\Delta_a^{(1)}\Delta_a^{(2)}\mathcal R.}
 \label{eq:YMA29-perimeter-cancelled}
\end{equation}
If
\(
 \partial_1\partial_2\mathcal R(\ell_1,\ell_2)\to0
\)
as both physical side lengths tend to infinity, and the discretization error
is controlled as in \Cref{thm:YMA29-mixed-curvature-average}, then the
large-loop Creutz density tends to \(\sigma\).
\end{theorem}

\begin{proof}
A function of one side alone has zero mixed finite difference.  The mixed
difference of \(\sigma\ell_1\ell_2\) over a cell of side \(a\) is
\(\sigma a^2\).  The final assertion follows from
\eqref{eq:YMA29-cell-average-curvature} applied to \(\mathcal R\).
\end{proof}

\begin{assessment}[What must be proved for a string-tension use]
\label{ass:YMA29-string-scope}
A fixed physical rectangle requires a continuum loop free energy on that
rectangle.  A string-tension use requires, in addition, a controlled
large-rectangle limit of its mixed curvature.  The fixed-loop
strong-coupling expansion of Version~V25 proves neither uniform statement by
itself because its remainder constant depends on the loop.
\end{assessment}

% =============================================================================
\section{The first cube bump already exposes nonuniform fixed-loop expansion}
\label{sec:YMA29-growing-loop}
% =============================================================================

Version~V28 proved for a fixed planar loop of area \(A\) that
\begin{equation}
 W_A(\beta)
 =\zeta(\beta)^A
  \left[1+4A\zeta(\beta)^4+O_A(\zeta(\beta)^5)\right].
 \label{eq:YMA29-fixed-loop-bump}
\end{equation}
The explicitly known first correction is enough to identify the missing
uniformity for a fixed physical loop.

\begin{proposition}[Fixed-loop perturbation is not uniform on growing physical rectangles]
\label{prop:YMA29-growing-loop-nonuniform}
Let \(A_r\to\infty\) be a sequence of lattice areas.  If
\begin{equation}
 A_r\zeta_r^4\not\longrightarrow0,
 \label{eq:YMA29-bump-smallness}
\end{equation}
then the explicit relative cube-bump contribution
\(4A_r\zeta_r^4\) in \eqref{eq:YMA29-fixed-loop-bump} is not a vanishing
perturbation of the flat sheet.  In particular, for a rectangle of fixed
positive physical area \(A_{\rm phys}\), where
\(
 A_r\asymp A_{\rm phys}/a_r^2
\), any trajectory with \(\zeta_r\ge\zeta_*>0\) violates
\eqref{eq:YMA29-bump-smallness} as \(a_r\to0\).
\end{proposition}

\begin{proof}
The first statement is the displayed coefficient itself.  On a fixed physical
area, \(A_r\zeta_r^4\ge c\zeta_*^4a_r^{-2}\to\infty\).
\end{proof}

\begin{assessment}[Surface resummation precedes a growing-loop root]
\label{ass:YMA29-surface-resummation}
The B\"ottcher coordinate of Version~V28 resums the homogeneous fundamental
surface activity through the first cube-bump order.  The degree-eight
nonclosure theorem in the same version proves that a complete physical
surface pressure must also retain the generated higher representations and
connected multi-plaquette terms.  A growing-loop selected root therefore
cannot be obtained by substituting \(A_r\to\infty\) into a fixed-loop
remainder symbol.
\end{assessment}

% =============================================================================
\section{Character-area density and the strong-coupling/continuum separation}
\label{sec:YMA29-strong-continuum}
% =============================================================================

Retain the normalized B\"ottcher character-tension coordinate of Version~V28,
\begin{equation}
 \mathcal U(\beta)=-\log\Phi(\zeta(\beta)),
 \qquad
 \mathcal U_{\rm out}=4\mathcal U_{\rm in}
 \label{eq:YMA29-U-recall}
\end{equation}
for one dyadic projected fundamental step.  Its physical area-density
calibration at lattice spacing \(a\) is
\begin{equation}
 \boxed{
 \Sigma_{\rm char}(a,\beta)
 :=\frac{\mathcal U(\beta)}{a^2}.}
 \label{eq:YMA29-character-density}
\end{equation}

\begin{theorem}[Exact projected area-density invariance]
\label{thm:YMA29-character-density-invariance}
If the outgoing activity and coupling satisfy
\(
 \zeta(\beta_{\rm out})=\mathcal F(\zeta(\beta_{\rm in}))
\)
for the projected map of Version~V28, and the outgoing spacing is
\(a_{\rm out}=2a_{\rm in}\), then
\begin{equation}
 \boxed{
 \Sigma_{\rm char}(a_{\rm out},\beta_{\rm out})
 =\Sigma_{\rm char}(a_{\rm in},\beta_{\rm in}).}
 \label{eq:YMA29-character-density-invariant}
\end{equation}
Thus \(\mathcal U\) is an exact scalar area coordinate for the projected
homogeneous activity map.  The exact physical Creutz density instead obeys
the four-child averaging law \eqref{eq:YMA29-density-average}; the two agree
as scalar transports only after the loop-cell detail has been controlled.
\end{theorem}

\begin{proof}
Use \(\mathcal U_{\rm out}=4\mathcal U_{\rm in}\) and
\(a_{\rm out}^2=4a_{\rm in}^2\).  The final statement is
\Cref{thm:YMA29-exact-block-Creutz,thm:YMA29-cell-Pythagoras}.
\end{proof}

\begin{theorem}[No finite area-density continuum route remains inside a fixed strong-coupling neighborhood]
\label{thm:YMA29-no-finite-density-strong}
There is a positive \(\beta_{\rm sc}\) inside the common strong-coupling
analytic disk and a constant \(u_{\rm sc}>0\) such that
\begin{equation}
 \mathcal U(\beta)\ge u_{\rm sc}
 \qquad(0<\beta\le\beta_{\rm sc}).
 \label{eq:YMA29-U-floor}
\end{equation}
Consequently, for every refining sequence \(a_r\to0\) and every
\(0<\beta_r\le\beta_{\rm sc}\),
\begin{equation}
 \boxed{
 \Sigma_{\rm char}(a_r,\beta_r)
 \ge\frac{u_{\rm sc}}{a_r^2}
 \longrightarrow\infty.}
 \label{eq:YMA29-strong-density-diverges}
\end{equation}
A trajectory with a finite nonzero limiting character-area density must
therefore leave every fixed sufficiently small-\(\beta\) neighborhood.
\end{theorem}

\begin{proof}
Version~V28 gives
\(
 \Phi(\zeta(\beta))=\beta/18+O(\beta^2)
\), so for sufficiently small positive \(\beta_{\rm sc}\),
\(
 0<\Phi(\zeta(\beta))<1
\)
and \(\mathcal U(\beta)>0\) on
\((0,\beta_{\rm sc}]\).  Moreover
\(
 \mathcal U(\beta)=-\log(\beta/18)+O(\beta)\to\infty
\)
as \(\beta\downarrow0\).  Positivity and continuity therefore give a
strict positive infimum \(u_{\rm sc}\) on that interval.  Divide by
\(a_r^2\).
\end{proof}

\begin{corollary}[Fixed and large response targets have divergent physical density]
\label{cor:YMA29-fixed-u-diverges}
Let \(u_r=\chi_{a_r}(m_r,n_r)\) be the actual selected Creutz value.
If \(a_r\to0\) and either
\begin{equation}
 u_r\ge u_0>0
 \quad\text{eventually},
 \qquad\text{or}\qquad
 u_r\to\infty,
 \label{eq:YMA29-u-positive-branches}
\end{equation}
then
\begin{equation}
 \boxed{
 \frac{u_r}{a_r^2}\longrightarrow\infty.}
 \label{eq:YMA29-u-density-diverges}
\end{equation}
In particular, the large-response roots of Versions~V25--V28 cannot be a
continuum trajectory with finite Creutz area density.
\end{corollary}

\begin{proof}
The first conclusion is immediate from \(a_r^2\to0\).  The large-response
root is the second branch.
\end{proof}

\begin{corollary}[Finite physical response density forces a small target]
\label{cor:YMA29-finite-density-small-target}
If
\begin{equation}
 \frac{u_r}{a_r^2}\longrightarrow\sigma_*
 \qquad\text{with}\qquad
 0<\sigma_*<\infty,
 \label{eq:YMA29-finite-density-target}
\end{equation}
then
\begin{equation}
 \boxed{
 u_r=\sigma_*a_r^2+o(a_r^2)\longrightarrow0.}
 \label{eq:YMA29-small-u-target}
\end{equation}
The large-response variable \(e^{-u_r}\) then tends to one rather than zero,
so the inverse-root series of Version~V28 is outside its asymptotic domain.
On the projected character branch, any corresponding root must eventually
leave the fixed strong-coupling neighborhood of
\Cref{thm:YMA29-no-finite-density-strong}.
\end{corollary}

\begin{proof}
Multiply \eqref{eq:YMA29-finite-density-target} by \(a_r^2\).  The final
statement follows from \eqref{eq:YMA29-strong-density-diverges}.
\end{proof}

\begin{assessment}[Correct status of the V28 outgoing-margin pass]
\label{ass:YMA29-outgoing-pass-scope}
The outgoing-character margin of Version~V28 is a valid finite-scale and
strong-coupling result.  On a cutoff sequence with fixed or large Creutz
response, it lies on a branch with divergent response per physical cell area.
It is therefore not, without a different calibrated root sequence, a pass for
a finite-area-density continuum local vacuum.
\end{assessment}

% =============================================================================
\section{Exact root transport and the minimum calibration packet}
\label{sec:YMA29-root-transport}
% =============================================================================

The logarithm-free root evaluator of Version~V10 remains exact after the loop
passport is supplied.  Its target must now carry the scale and rectangle
labels.

\begin{theorem}[Exact coarse-root equation on the fine loop table]
\label{thm:YMA29-root-transport}
On the positive loop branch of an exact dyadic block, the coarse selected-root
equation
\begin{equation}
 \chi_{2a}^{\rm out}(m,n;\beta)=u_{\rm out}
 \label{eq:YMA29-coarse-root}
\end{equation}
is equivalent to
\begin{equation}
 \boxed{
 \sum_{i,j\in\{0,1\}}
 \chi_a^{\rm in}(2m-i,2n-j;\beta)
 =u_{\rm out}.}
 \label{eq:YMA29-fine-four-root}
\end{equation}
Knowledge of one fine selected-root equation
\(
 \chi_a(k,\ell;\beta)=u_{\rm in}
\)
does not determine \eqref{eq:YMA29-fine-four-root} unless an additional
four-cell homogeneity or dispersion estimate is supplied.
\end{theorem}

\begin{proof}
The equivalence is \eqref{eq:YMA29-exact-Creutz-pushforward}.  The last
statement is witnessed by \Cref{cth:YMA29-one-cell-no-parent}.
\end{proof}

\begin{definition}[Scale-complete selected-root manifest]
\label{def:YMA29-root-manifest}
A selected-root card is scale complete only when it contains:
\begin{enumerate}[label=\textup{(M\arabic*)},leftmargin=4em]
\item the lattice spacing \(a_r\), loop counts \((m_r,n_r)\), and physical
      side lengths;
\item the semantic branch in \Cref{def:YMA29-rectangle-passport};
\item the target \(u_r\) with its physical scaling rule;
\item for adjacent-cutoff transport, the four-child response panel or a
      certified Creutz-cell dispersion bound;
\item the logarithm-free root interval for precisely that quartet and target;
\item only afterward, the two native response cofactors and complete
      same-action Wilson Hessian card.
\end{enumerate}
\end{definition}

\begin{theorem}[Terminating calibration alternative]
\label{thm:YMA29-calibration-alternative}
For a proposed cofinal selected-root sequence, exactly the first applicable
branch below is returned.
\begin{enumerate}[label=\textup{(C\arabic*)},leftmargin=4em]
\item If \((a_r,m_r,n_r)\) or the semantic branch is absent, the card is
      \stopstatus{} at the absolute rectangle-calibration row.
\item If the fixed lattice-cell branch is declared, the Version~V10 finite
      root theorem applies, but no fixed-physical-rectangle or string-tension
      conclusion is exported.
\item If a fixed physical rectangle is declared, the target must have the form
      \(
       u_r=a_r^2\kappa_r
      \)
      with calibrated \(\kappa_r\), and the loop counts must grow so that the
      physical side lengths converge.  The output is a finite-rectangle mixed
      loop-free-energy curvature.
\item If a large-loop area density is declared, a uniform growing-loop or
      surface-pressure theorem is additionally required.  A fixed-loop
      remainder does not populate this row.
\item A fixed nonzero target or a large-response root used as a finite
      physical area-density continuum calibration is a
      \obstruction{} to that declared calibration, by
      \Cref{cor:YMA29-fixed-u-diverges}; it is not a negative Yang--Mills mode.
\end{enumerate}
After the calibration branch passes, the physical Wilson mode still returns
\pass, an attained \obstruction, or \stopstatus{} through the same-action
Hessian and Ward-tail certificate of the preceding versions.
\end{theorem}

\begin{proof}
Items~\textup{(C1)--(C4)} are the definitions and
\Cref{cor:YMA29-fixed-lattice-vs-physical,ass:YMA29-string-scope}.  Item
\textup{(C5)} is \Cref{cor:YMA29-fixed-u-diverges}.  None of these calibration
statements changes the spectral sign rule for the fully assembled Hessian.
\end{proof}

% =============================================================================
\section{Integrated Version V29 theorem and proof-status ledger}
\label{sec:YMA29-integrated}
% =============================================================================

\begin{theorem}[Physical-rectangle and Creutz-transport theorem]
\label{thm:YMA29-integrated}
Preserving Versions~V1--V28, the following statements are now proved at their
stated scopes.
\begin{enumerate}[label=\textup{(V29.\arabic*)},leftmargin=5.7em]
\item A rectangular Creutz response is one mixed finite difference of the
      logarithmic loop expectation, and its physical density is obtained by
      dividing by the elementary response-cell area \(a^2\).
\item Exact blocking by a factor \(b\) transports one coarse Creutz response
      into the sum of the \(b^2\) fine response cells which tile it; the
      physical response density is their arithmetic mean.
\item The dyadic four-cell packet has an exact positive Pythagoras residual.
      Its vanishing is the homogeneous scalar branch; one fine cell does not
      determine its parent.
\item Under a \(C^4\) continuum loop-free-energy hypothesis, the Creutz
      density is the cell average of the mixed physical-side derivative, with
      the explicit midpoint error \eqref{eq:YMA29-midpoint-error}.
\item Fixed lattice counts, fixed physical rectangles, and the large-loop
      string-tension limit are distinct targets.  Perimeter and one-side terms
      cancel exactly, but the large-loop remainder remains an independent
      uniformity row.
\item The first cube-bump correction grows as \(4A\zeta^4\), so the fixed-loop
      strong-coupling expansion cannot simply be substituted into a growing
      physical rectangle unless a surface resummation is supplied.
\item The B\"ottcher character-area density is exactly invariant under the
      projected dyadic map, while the physical Creutz density obeys the more
      general four-child averaging identity.
\item Every trajectory remaining in a fixed strong-coupling neighborhood has
      divergent character response per physical cell area as \(a\to0\).
      Fixed and large response targets have the same calibration obstruction.
\item A finite physical response density requires a target
      \(u_r=O(a_r^2)\to0\), outside the large-response inverse-root regime of
      Version~V28.
\item The scale-complete root manifest and
      \Cref{thm:YMA29-calibration-alternative} give a terminating upstream
      branch before the response cofactors or Wilson mode card are opened.
\end{enumerate}
No item evaluates an unspecified physical response target or supplies the
native same-action selected-root Hessian.
\end{theorem}

\begin{proof}
Items~\textup{(V29.1)--(V29.3)} are
\Cref{thm:YMA29-exact-block-Creutz,thm:YMA29-cell-Pythagoras,cth:YMA29-one-cell-no-parent}.
Items~\textup{(V29.4)--(V29.5)} are
\Cref{thm:YMA29-mixed-curvature-average,cor:YMA29-fixed-lattice-vs-physical,thm:YMA29-perimeter-cancellation}.
Item~\textup{(V29.6)} is \Cref{prop:YMA29-growing-loop-nonuniform}.
Items~\textup{(V29.7)--(V29.9)} are
\Cref{thm:YMA29-character-density-invariance,thm:YMA29-no-finite-density-strong,cor:YMA29-fixed-u-diverges,cor:YMA29-finite-density-small-target}.
Item~\textup{(V29.10)} is \Cref{thm:YMA29-calibration-alternative}.
\end{proof}

\begingroup\small
\begin{longtable}{>{\raggedright\arraybackslash}p{0.27\textwidth}
                  >{\raggedright\arraybackslash}p{0.18\textwidth}
                  >{\raggedright\arraybackslash}p{0.47\textwidth}}
\caption{Version V29 proof-status ledger.}\label{tab:YMA29-ledger}\\
\toprule
\textbf{Row}&\textbf{Status}&\textbf{Exact advance or limitation}\\
\midrule
\endfirsthead
\toprule
\textbf{Row}&\textbf{Status}&\textbf{Exact advance or limitation}\\
\midrule
\endhead
V1--V28 prefix
& \pass
& Complete preterminal source preserved byte for byte.\\[0.35em]
Rectangle passport
& \pass
& Lattice spacing, loop counts, physical side lengths, target scaling, and use
  semantics are separated.\\[0.35em]
Dyadic Creutz transport
& \pass
& One coarse response is exactly the sum of four fine response cells; density
  is their average.\\[0.35em]
Four-cell detail
& \pass
& Positive scalar variance with exact Pythagoras and a sharp
  \(\sqrt{3\mathcal V}\) child deviation bound.\\[0.35em]
Continuum rectangle interpretation
& Conditional theorem
& Requires a declared \(C^4\) renormalized continuum loop free energy.\\[0.35em]
String-tension interpretation
& \stopstatus
& Requires an additional large-physical-rectangle and uniform surface-pressure
  limit.\\[0.35em]
Fixed-loop strong-coupling use on growing loops
& \obstruction
& The explicit \(4A\zeta^4\) correction is nonuniform unless
  \(A\zeta^4\to0\).\\[0.35em]
Projected character-area density
& \pass
& Exactly invariant under one dyadic projected activity step.\\[0.35em]
Fixed or large response as finite continuum density
& \obstruction
& The response per physical cell area diverges like at least \(a^{-2}\).\\[0.35em]
Finite physical density target
& Reclassified
& Necessarily \(u_r=O(a_r^2)\to0\); the V28 large-response expansion does not
  apply.\\[0.35em]
Native response cofactors
& \stopstatus
& Must be computed only after the scale-complete loop and target passport is
  frozen.\\[0.35em]
Selected-shell Wilson mode
& \stopstatus
& No same-action selected-root Hessian or differentiated Ward tail has been
  populated.\\[0.35em]
Local-vacuum or continuum mass
& Not inferred
& Requires the complete signed head, conductance, ancestry, tail and cutoff
  transport, and reflection-positive continuum.\\
\bottomrule
\end{longtable}
\endgroup

% =============================================================================
\section{Revised V30 acquisition}
\label{sec:YMA29-next}
% =============================================================================

\begin{openproblem}[V30: populate one scale-complete response root]
\label{op:YMA29-V30}
Preserve the complete present source.  Do not append another bare-\(\beta\)
coefficient or comparison compiler before choosing one of the following
physical branches.
\begin{enumerate}[label=\textup{(V30.\arabic*)},leftmargin=5.7em]
\item \emph{Fixed lattice-cell scheme.}  Supply the numerical target \(u_*\),
      loop quartet, and finite root interval.  Report the result only as a
      finite lattice-scheme card.
\item \emph{Fixed physical rectangle.}  Supply \(a_r,m_r,n_r\) with fixed
      physical side limits, predeclare the target density \(\kappa_*\), and
      solve
      \begin{equation}
       \chi_{a_r}(m_r,n_r;\beta_r)
       =a_r^2\kappa_*+o(a_r^2).
       \label{eq:YMA29-V30-fixed-rectangle-root}
      \end{equation}
      Acquire the exact four-child response panel at one adjacent dyadic pair.
\item \emph{Large-loop area density.}  Prove or populate a uniform surface
      pressure and its mixed large-loop limit before calling the selected
      scalar a string tension.
\item On the chosen branch, compute the two native response cofactors for that
      precise quartet and target, load the complete same-action selected root,
      and evaluate one retained momentum--colour Hessian with its differentiated
      Ward tail.
\item Return \pass, an attained \obstruction, or \stopstatus{} at the first
      unresolved calibration, root, action, Hessian, or tail row.
\end{enumerate}
The V28 large-response root remains useful only for the fixed lattice-cell or
divergent-area-density strong-coupling branch.  It is not extrapolated into
\eqref{eq:YMA29-V30-fixed-rectangle-root}.
\end{openproblem}

\begin{assessment}[Version V29 termination]
\label{ass:YMA29-termination}
The upstream obstruction is no longer an unspecified choice between a fixed
scheme number and an area coordinate.  It is the missing physical rectangle
passport and, under cutoff transport, the missing four-child Creutz panel.
The projected B\"ottcher coordinate is the exact homogeneous scalar branch,
while the full physical response is an average with a positive cell-detail
residual.  A finite physical area-density continuum target necessarily leaves
the large-response strong-coupling regime.  The next advance must therefore
populate one scale-complete root rather than extend the strong-coupling
coefficient ladder.
\end{assessment}

% =============================================================================
\section{Exact replay and append-only continuation record}
\label{sec:YMA29-replay}
% =============================================================================

The accompanying exact replay verifies:
\begin{enumerate}[label=\textup{(R29.\arabic*)},leftmargin=5.5em]
\item the block-factor mixed-difference telescoping identity for symbolic loop
      free energies;
\item the dyadic four-cell response and density-average formulas;
\item the four-child Pythagoras identity and the sharp scalar deviation bound;
\item the positive-table counterexample with parent responses four and seven;
\item the exact cell-average formula on polynomial continuum free energies and
      the midpoint constant \(1/24\);
\item exact cancellation of constant and one-side terms;
\item character-area invariance under
      \((a,\mathcal U)\mapsto(2a,4\mathcal U)\);
\item fixed-target and finite-density scaling alternatives.
\end{enumerate}
The replay is an algebraic theorem check.  It does not evaluate a Wilson path
integral, a physical root, a surface pressure, a response cofactor, or a mode
margin.

The accumulated V29 source:
\begin{itemize}[leftmargin=2.2em]
\item preserves the first \texttt{1297114} bytes of V28 exactly;
\item contains \texttt{33800} newline-terminated source lines and
      \texttt{1337214} bytes;
\item contains \texttt{2458} labels, all unique;
\item has exactly one active terminal \verb|\end{document}|;
\item compiles through repeated \texttt{pdflatex} passes with no unresolved
      references or duplicate labels;
\item produces a \texttt{426}-page PDF, with V29 beginning on page
      \texttt{414};
\item passes PDF preflight and rendered-page inspection.
\end{itemize}

The complete Version~V28 source preceding its terminal document-closing
command is preserved byte for byte.  Its full-source SHA--256 digest is
\begin{center}\scriptsize\ttfamily
026eada02221c7c8259ea6b5a7fed74a4a6df7173749e38801a72afcfcf67bb3.
\end{center}
Its preterminal prefix has SHA--256
\begin{center}\scriptsize\ttfamily
c8c7de209c9a5869a2e218e0de0b76ea592073de1513de50f83e846e0220574b.
\end{center}
For Version~V30, preserve this accumulated source, append new material
immediately before the final active document-closing command, and increment
the filename to \texttt{\_V30.tex}.

% =============================================================================
% END APPENDED VERSION V29 MATERIAL
% =============================================================================


% =============================================================================
% START APPENDED VERSION V30 MATERIAL
% =============================================================================

\clearpage
\hypersetup{
 pdftitle={Renewal Yang--Mills Reduced-Fibre Wilson Shell Symbol and Local-Vacuum Programme V30},
 pdfsubject={Restoration of the fixed-physical macro-Creutz passport, exact source transport, and the first nonlinear alias read}
}

\pdfinfo{
 /Title (Renewal Yang--Mills Reduced-Fibre Wilson Shell Symbol and Local-Vacuum Programme V30)
 /Subject (Restoration of the fixed-physical macro-Creutz passport, exact source transport, and the first nonlinear alias read)
}

\section{Version V30: restoration of the fixed-physical macro-Creutz passport}
\label{sec:YMA30-direction}

Version~V29 is retained verbatim above.  Its unit-cell mixed-difference
identities are correct, but their application to the inherited Branch--R
selector used the wrong response cell.  The already frozen Yang--Mills
trajectory does not use a one-lattice-step Creutz cell.  It fixes a physical
box of side length \(L\), the four physical rectangles
\[
 R=T=\frac L4,
 \qquad
 \Delta=\frac L8,
\]
and, at lattice side \(M\in16\mathbb N\), represents them by
\[
 m_M=n_M=\frac M4,
 \qquad
 d_M=\frac M8,
 \qquad
 a_M=\frac LM.
\]
The Creutz response cell in the two side-length variables therefore has
physical side \(d_Ma_M=\Delta\), not \(a_M\).  Under the adjacent dyadic
refinement \(M\mapsto2M\), all three lattice counts double and the physical
rectangle is unchanged.

The correct upstream order is consequently
\begin{equation}
 \boxed{
 \begin{gathered}
 \text{the inherited fixed-physical macrocell and one fixed response }u_*\\
 \Downarrow\\
 \text{exact fine-to-trace transport of the four loop writers and their source}\\
 \Downarrow\\
 \text{natural-coarse versus trace-coarse law-alignment residual}\\
 \Downarrow\\
 \text{the first nonlinear selected-surface alias Read}\\
 \Downarrow\\
 \text{the complete same-action Wilson Hessian and differentiated Ward tail}.
 \end{gathered}}
 \label{eq:YMA30-corrected-order}
\end{equation}
In particular, the physical rectangle passport requested at the end of
Version~V29 was already present in the inherited programme.  The missing
quantity is not another rectangle convention or another four-child root.  It
is the Wilson-specific law-aligned response of the complete generated action.

\begin{firewall}[Scope of the Version V30 correction]
\label{fire:YMA30-scope}
The present continuation corrects the use of Version~V29 on the inherited
fixed-physical branch.  It does not retract the unit-cell telescope, the
four-child Pythagoras identity, the continuum cell-average formula, or the
large-loop warning in their stated domains.  It proves exact macrocell and
source-incidence identities and gives a finite nonlinear alias acquisition.
It does not populate a numerical response value, a cofinal root family, the
complete nonlinear Wilson run, a negative momentum--colour mode, or a mass
gap.
\end{firewall}

% =============================================================================
\section{Arbitrary-step Creutz macrocells}
\label{sec:YMA30-macrocell}
% =============================================================================

Retain the positive rectangular-loop expectations and logarithmic loop free
energy of Version~V29,
\(
 F_a(m,n)=-\log W_a(m,n)
\).
For positive integers \(d_1,d_2\), define the forward macrocell mixed
difference
\begin{equation}
 \boxed{
 \begin{aligned}
 \mathfrak C_{d_1,d_2}F(m,n)
 :={}&F(m+d_1,n+d_2)-F(m+d_1,n)\\
 &-F(m,n+d_2)+F(m,n).
 \end{aligned}}
 \label{eq:YMA30-macrocell-def}
\end{equation}
Its physical side increments are
\begin{equation}
 \Delta_1=d_1a,
 \qquad
 \Delta_2=d_2a,
 \label{eq:YMA30-macrocell-increments}
\end{equation}
and its physical mixed-response density is
\begin{equation}
 \boxed{
 \Sigma_{a;d_1,d_2}(m,n)
 :=\frac{\mathfrak C_{d_1,d_2}F_a(m,n)}
          {(d_1a)(d_2a)}.}
 \label{eq:YMA30-macrocell-density}
\end{equation}
The Version~V29 quantity \(\chi_a(m,n)\) is the translated special case
\(d_1=d_2=1\).

\begin{theorem}[Exact arbitrary-step mixed-difference telescope]
\label{thm:YMA30-macro-telescope}
With the unit backward difference of
\eqref{eq:YMA29-unit-mixed}, one has
\begin{equation}
 \boxed{
 \mathfrak C_{d_1,d_2}F(m,n)
 =\sum_{i=1}^{d_1}\sum_{j=1}^{d_2}
   (\Delta_1\Delta_2F)(m+i,n+j).}
 \label{eq:YMA30-macro-telescope}
\end{equation}
Consequently, the macrocell density is the arithmetic average of the
\(d_1d_2\) unit-cell densities:
\begin{equation}
 \boxed{
 \Sigma_{a;d_1,d_2}(m,n)
 =\frac1{d_1d_2}
  \sum_{i=1}^{d_1}\sum_{j=1}^{d_2}
  \Sigma_a(m+i,n+j).}
 \label{eq:YMA30-macro-density-average}
\end{equation}
\end{theorem}

\begin{proof}
Telescope first in the first variable:
\[
 F(m+d_1,t)-F(m,t)
 =\sum_{i=1}^{d_1}\bigl(F(m+i,t)-F(m+i-1,t)\bigr).
\]
Take the difference between \(t=n+d_2\) and \(t=n\), then telescope each
summand in the second variable.  The resulting summand is
\((\Delta_1\Delta_2F)(m+i,n+j)\).  Division by
\(d_1d_2a^2\) gives \eqref{eq:YMA30-macro-density-average}.
\end{proof}

\begin{theorem}[Exact continuum meaning of a fixed physical macrocell]
\label{thm:YMA30-continuum-macrocell}
Assume that on the relevant rectangle
\begin{equation}
 F_a(m,n)=\mathscr F(ma,na)
 \label{eq:YMA30-continuum-sampling}
\end{equation}
for a \(C^2\) function \(\mathscr F\).  Put
\(
 x=ma
\),
\(
 y=na
\), and retain \(\Delta_i=d_ia\).  Then
\begin{equation}
 \boxed{
 \mathfrak C_{d_1,d_2}F_a(m,n)
 =\int_x^{x+\Delta_1}\int_y^{y+\Delta_2}
   \partial_1\partial_2\mathscr F(s,t)\,\dd t\,\dd s,}
 \label{eq:YMA30-continuum-macro-integral}
\end{equation}
and hence
\begin{equation}
 \boxed{
 \Sigma_{a;d_1,d_2}(m,n)
 =\frac1{\Delta_1\Delta_2}
  \int_x^{x+\Delta_1}\int_y^{y+\Delta_2}
  \partial_1\partial_2\mathscr F(s,t)\,\dd t\,\dd s.}
 \label{eq:YMA30-continuum-macro-average}
\end{equation}
If \(\Delta_1,\Delta_2\) remain fixed as \(a\to0\), the limit is a fixed
physical cell average, not a point derivative obtained by dividing by
\(a^2\).
\end{theorem}

\begin{proof}
Apply the fundamental theorem of calculus in the second variable to the
difference
\(
 \mathscr F(x+\Delta_1,y+\Delta_2)
 -\mathscr F(x+\Delta_1,y)
\), do the same at \(x\), subtract, and apply the fundamental theorem in the
first variable.  The density formula is division by the physical cell area.
The final statement follows because the integration domain does not shrink.
\end{proof}

\begin{corollary}[Area-law calibration of a macrocell]
\label{cor:YMA30-area-law-macrocell}
If on the selected rectangle
\begin{equation}
 \mathscr F(x,y)=\sigma xy+p_1(x)+p_2(y)+c+\mathscr R(x,y),
 \label{eq:YMA30-macro-area-law}
\end{equation}
then
\begin{equation}
 \boxed{
 \mathfrak C_{d_1,d_2}F_a
 =\sigma\Delta_1\Delta_2
  +\int_x^{x+\Delta_1}\int_y^{y+\Delta_2}
    \partial_1\partial_2\mathscr R(s,t)\,\dd t\,\dd s.}
 \label{eq:YMA30-macro-area-response}
\end{equation}
The constant and one-side terms cancel exactly.  Thus a fixed nonzero
macro-response is compatible with a finite physical area density whenever
the physical response-cell area is fixed.
\end{corollary}

\begin{proof}
Insert \eqref{eq:YMA30-macro-area-law} into
\eqref{eq:YMA30-continuum-macro-integral}.  The mixed derivative of the first
term is \(\sigma\), and those of \(p_1,p_2,c\) vanish.
\end{proof}

% =============================================================================
\section{The inherited Branch--R macrocell and its exact dyadic transport}
\label{sec:YMA30-inherited-branch}
% =============================================================================

For the inherited card put
\begin{equation}
 a_M=\frac LM,
 \qquad
 m_M=n_M=\frac M4,
 \qquad
 d_M=\frac M8.
 \label{eq:YMA30-inherited-counts}
\end{equation}
The selected response is
\begin{equation}
 \boxed{
 \chi_M
 =\mathfrak C_{d_M,d_M}F_M(m_M,n_M).}
 \label{eq:YMA30-inherited-response}
\end{equation}

\begin{theorem}[The inherited response cell has fixed physical area]
\label{thm:YMA30-fixed-physical-cell}
For every \(M\in16\mathbb N\),
\begin{equation}
 \boxed{
 m_Ma_M=n_Ma_M=\frac L4,
 \qquad
 d_Ma_M=\frac L8=\Delta.}
 \label{eq:YMA30-fixed-physical-identities}
\end{equation}
Consequently, the physical mixed-response density of the inherited card is
\begin{equation}
 \boxed{
 \Sigma_M^{\rm phys}
 =\frac{\chi_M}{\Delta^2}
 =\frac{64}{L^2}\chi_M.}
 \label{eq:YMA30-inherited-density}
\end{equation}
On a response-matched family \(\chi_M=u_*\),
\begin{equation}
 \boxed{
 \Sigma_M^{\rm phys}=\frac{64u_*}{L^2}}
 \label{eq:YMA30-fixed-u-density}
\end{equation}
for every cutoff on which the root exists.  In particular, fixed \(u_*\)
does not create an \(a_M^{-2}\) divergence on this branch.
\end{theorem}

\begin{proof}
Substitute \eqref{eq:YMA30-inherited-counts} and \(a_M=L/M\).  The response
cell in side-length space has area
\((d_Ma_M)^2=L^2/64\).  Dividing \(\chi_M\) by that area gives
\eqref{eq:YMA30-inherited-density}; the matched formula is immediate.
\end{proof}

\begin{theorem}[Exact macrocell transport under a straight block]
\label{thm:YMA30-exact-macro-block}
Let a block of factor \(b\) satisfy
\begin{equation}
 F_{ba}^{\rm tr}(r,t)=F_a^{\rm f}(br,bt)
 \label{eq:YMA30-free-energy-block}
\end{equation}
for every rectangular loop in the selected family.  Then
\begin{equation}
 \boxed{
 \mathfrak C_{d_1,d_2}F_{ba}^{\rm tr}(m,n)
 =\mathfrak C_{bd_1,bd_2}F_a^{\rm f}(bm,bn).}
 \label{eq:YMA30-macro-block-identity}
\end{equation}
The corresponding physical densities are equal:
\begin{equation}
 \boxed{
 \Sigma_{ba;d_1,d_2}^{\rm tr}(m,n)
 =\Sigma_{a;bd_1,bd_2}^{\rm f}(bm,bn).}
 \label{eq:YMA30-macro-density-block}
\end{equation}
For the inherited dyadic pair,
\begin{equation}
 \boxed{
 (m_{2M},n_{2M},d_{2M})
 =2(m_M,n_M,d_M),
 \qquad
 \chi_M^{\rm tr}=\chi_{2M}^{\rm f}.}
 \label{eq:YMA30-inherited-dyadic-response}
\end{equation}
\end{theorem}

\begin{proof}
Substitute \eqref{eq:YMA30-free-energy-block} into the four terms of
\eqref{eq:YMA30-macrocell-def}.  The coarse physical increments are
\(d_i(ba)\), while the fine increments are \((bd_i)a\); these are equal, so
division gives \eqref{eq:YMA30-macro-density-block}.  The last statement is
the specialization \(b=2\) together with
\eqref{eq:YMA30-inherited-counts}.
\end{proof}

\begin{theorem}[Internal quadrant decomposition of the fine macrocell]
\label{thm:YMA30-quadrant-decomposition}
For every integer \(b\ge1\),
\begin{equation}
 \boxed{
 \begin{aligned}
 &\mathfrak C_{bd_1,bd_2}F(bm,bn)\\
 &\quad=
 \sum_{r=0}^{b-1}\sum_{s=0}^{b-1}
 \mathfrak C_{d_1,d_2}F(bm+rd_1,bn+sd_2).
 \end{aligned}}
 \label{eq:YMA30-quadrant-decomposition}
\end{equation}
For \(b=2\), the four terms on the right have physical side increments
\(\Delta_1/2,\Delta_2/2\) when the left side represents the same physical
macrocell as the coarse card.  Their density average is the parent density.
They are therefore an optional internal-resolution panel; they are not four
alternative roots for the inherited adjacent response match.
\end{theorem}

\begin{proof}
Partition the double sum in \eqref{eq:YMA30-macro-telescope} into the
\(b^2\) disjoint rectangles
\[
 rd_1<i\le(r+1)d_1,
 \qquad
 sd_2<j\le(s+1)d_2.
\]
Telescope each rectangle back into a macrocell.  Under a factor-\(b\) block,
the fine lattice spacing is \(a\) and the parent physical increment is
\(bd_i a\), while each child increment is \(d_i a\).  The density statement
follows after division by the respective areas.  The root statement follows
because the selected parent response is the sum, not any one child.
\end{proof}

\begin{assessment}[What the four-child panel now means]
\label{ass:YMA30-four-child-scope}
The Version~V29 four-child variance remains a valid positive measure of
inhomogeneity inside the fixed physical response cell.  It is required when
one wants a homogeneous scalarization, a resolved response density, or a
bound on variation inside that cell.  It is not a missing premise for the
already declared parent-macrocell equation
\(
 \chi_M(\beta_M^*)=u_*=\chi_{2M}(\beta_{2M}^*)
\).
\end{assessment}

% =============================================================================
\section{Exact transport of the macro-Creutz source}
\label{sec:YMA30-source-transport}
% =============================================================================

Index the four loop writers by
\(
 \alpha=(\epsilon,\eta)\in\{0,1\}^2
\)
and use the Creutz signs
\begin{equation}
 \sigma_{11}=\sigma_{00}=-1,
 \qquad
 \sigma_{10}=\sigma_{01}=1.
 \label{eq:YMA30-Creutz-signs}
\end{equation}
For a law \(\lambda\) under which all four loop means
\(
 W_\alpha^\lambda=\mathbb E_\lambda\mathcal O_\alpha
\)
are positive, define the centered macro-Creutz response source
\begin{equation}
 \boxed{
 \mathfrak s_{\rm Cr}^{\lambda}
 :=\sum_{\alpha}\sigma_\alpha
    \frac{\mathcal O_\alpha}{W_\alpha^\lambda}.}
 \label{eq:YMA30-macro-source}
\end{equation}
It is centered because \(\sum_\alpha\sigma_\alpha=0\).

\begin{theorem}[Pointwise fine-to-trace source incidence]
\label{thm:YMA30-source-incidence}
Let \(B:X_{2M}\to X_M\) be the exact dyadic block, let
\(
 \mu_{2M}^{\rm f}
\)
be the fine law, and let
\(
 \nu_M^{\rm tr}=B_*\mu_{2M}^{\rm f}
\)
be its induced coarse trace law.  For the same four physical rectangles,
\begin{equation}
 \boxed{
 W_{\alpha}^{M,{\rm tr}}
 =W_{\alpha}^{2M,{\rm f}},
 \qquad
 \mathfrak s_{{\rm Cr},M}^{\rm tr}\circ B
 =\mathfrak s_{{\rm Cr},2M}^{\rm f}.}
 \label{eq:YMA30-source-pointwise}
\end{equation}
Consequently,
\begin{equation}
 \boxed{
 \operatorname{Var}_{\nu_M^{\rm tr}}
   (\mathfrak s_{{\rm Cr},M}^{\rm tr})
 =\operatorname{Var}_{\mu_{2M}^{\rm f}}
   (\mathfrak s_{{\rm Cr},2M}^{\rm f}),}
 \label{eq:YMA30-source-stock-transport}
\end{equation}
and, for every bounded coarse writer \(X\),
\begin{equation}
 \boxed{
 \operatorname{Cov}_{\nu_M^{\rm tr}}
 (\mathfrak s_{{\rm Cr},M}^{\rm tr},X)
 =\operatorname{Cov}_{\mu_{2M}^{\rm f}}
 (\mathfrak s_{{\rm Cr},2M}^{\rm f},X\circ B).}
 \label{eq:YMA30-source-cov-transport}
\end{equation}
No stochastic fine/coarse coupling beyond the occurring deterministic block
is needed for this source incidence.
\end{theorem}

\begin{proof}
The block map sends each doubled fine physical loop to the corresponding
coarse loop pointwise.  Pushforward therefore gives equality of the four
means.  Substitute those two equalities into
\eqref{eq:YMA30-macro-source}; every summand agrees pointwise after pullback.
The variance and covariance identities are then the elementary pushforward
formula for expectations.
\end{proof}

\begin{corollary}[The selected scalar is already common to the natural and trace laws]
\label{cor:YMA30-common-selected-response}
Let \(\mu_M^{\rm nat}\) be the natural coarse Wilson law at the selected root
\(\beta_M^*\), and let \(\nu_M^{\rm tr}\) be the trace of the fine selected
law at \(\beta_{2M}^*\).  The inherited adjacent response match gives
\begin{equation}
 \boxed{
 \chi[\mu_M^{\rm nat}]=u_*=\chi[\nu_M^{\rm tr}].}
 \label{eq:YMA30-natural-trace-same-response}
\end{equation}
Thus the first missing coarse/fine row is not the selected Creutz scalar.  It
is the difference between the two laws, their complete source descendants,
and the Wilson Hessian they induce away from the one fitted response.
\end{corollary}

\begin{proof}
The first equality is the coarse selected-root equation.  By
\Cref{thm:YMA30-source-incidence}, the response of the trace law is the fine
response, which equals \(u_*\) by the fine selected-root equation.
\end{proof}

% =============================================================================
\section{What one matched Creutz scalar leaves unresolved}
\label{sec:YMA30-affine-loop-residual}
% =============================================================================

Let
\(
 W_\alpha^{\rm nat}
\)
and
\(
 W_\alpha^{\rm tr}
\)
be the four positive loop means under the two coarse laws of
\Cref{cor:YMA30-common-selected-response}, and define
\begin{equation}
 D_{\epsilon\eta}
 :=\log W_{\epsilon\eta}^{\rm nat}
   -\log W_{\epsilon\eta}^{\rm tr}.
 \label{eq:YMA30-log-loop-defect}
\end{equation}

\begin{theorem}[Exact three-coordinate loop-shape residual]
\label{thm:YMA30-affine-shape-residual}
Equality of the two Creutz responses is equivalent to
\begin{equation}
 \boxed{
 D_{11}=D_{10}+D_{01}-D_{00}.}
 \label{eq:YMA30-zero-mixed-log-defect}
\end{equation}
Equivalently, there are unique real numbers \(A,B,C\) such that
\begin{equation}
 \boxed{
 D_{\epsilon\eta}=A+B\epsilon+C\eta
 \qquad(\epsilon,\eta\in\{0,1\}),}
 \label{eq:YMA30-affine-shape-form}
\end{equation}
namely
\begin{equation}
 A=D_{00},
 \qquad
 B=D_{10}-D_{00},
 \qquad
 C=D_{01}-D_{00}.
 \label{eq:YMA30-affine-shape-coordinates}
\end{equation}
Thus one matched Creutz scalar removes exactly the mixed corner of the
four-log-loop table and leaves three affine shape coordinates.
\end{theorem}

\begin{proof}
The response is
\[
 \chi=-\log W_{11}-\log W_{00}
      +\log W_{10}+\log W_{01}.
\]
Subtract the natural and trace values.  Equality is precisely
\eqref{eq:YMA30-zero-mixed-log-defect}.  A function on the Boolean square has
zero mixed difference exactly when it is affine in the two Boolean
coordinates.  The displayed values of \(A,B,C\) reconstruct the four entries,
and uniqueness follows from the entries at \((0,0),(1,0),(0,1)\).
\end{proof}

\begin{theorem}[Exact macro-source mismatch from the three shape coordinates]
\label{thm:YMA30-source-shape-mismatch}
With \(D_\alpha\) as above,
\begin{equation}
 \boxed{
 \mathfrak s_{\rm Cr}^{\rm nat}
 -\mathfrak s_{\rm Cr}^{\rm tr}
 =\sum_\alpha\sigma_\alpha
   \frac{\mathcal O_\alpha}{W_\alpha^{\rm tr}}
   \bigl(e^{-D_\alpha}-1\bigr).}
 \label{eq:YMA30-exact-source-shape-mismatch}
\end{equation}
If
\begin{equation}
 w_*:=\min_\alpha W_\alpha^{\rm tr}>0,
 \qquad
 d_*:=\max_\alpha|D_\alpha|,
 \label{eq:YMA30-shape-bounds}
\end{equation}
then, using \(|\mathcal O_\alpha|\le1\),
\begin{equation}
 \boxed{
 \left\|
 \mathfrak s_{\rm Cr}^{\rm nat}
 -\mathfrak s_{\rm Cr}^{\rm tr}
 \right\|_\infty
 \le\frac{4(e^{d_*}-1)}{w_*}.}
 \label{eq:YMA30-source-shape-bound}
\end{equation}
The right side is a sufficient source-alignment bound.  Its failure is
unresolved unless a direct source-distance lower witness is acquired.
\end{theorem}

\begin{proof}
By definition,
\(
 W_\alpha^{\rm nat}=e^{D_\alpha}W_\alpha^{\rm tr}
\), so
\(
 (W_\alpha^{\rm nat})^{-1}
 =e^{-D_\alpha}(W_\alpha^{\rm tr})^{-1}
\).
Insert this into \eqref{eq:YMA30-macro-source} and subtract.  Since
\(
 |e^{-D}-1|\le e^{|D|}-1
\), each summand is at most
\((e^{d_*}-1)/w_*\) in absolute value.  Sum the four bounds.
\end{proof}

\begin{countertheorem}[A matched Creutz value does not align its source]
\label{cth:YMA30-response-no-source-alignment}
There are two full-support positive laws on the same four bounded loop writers
with identical Creutz response and different macro-Creutz source functions.
\end{countertheorem}

\begin{proof}
Let \(\Omega=\{0,1\}^4\), let the four coordinate writers be
\(\mathcal O_\alpha(\omega)=\omega_\alpha\), and use product Bernoulli laws.
In the order \((11,10,01,00)\), take the mean vectors
\[
 W^{\rm tr}=\left(\frac12,\frac12,\frac12,\frac12\right),
 \qquad
 W^{\rm nat}=\left(\frac13,\frac12,\frac12,\frac34\right).
\]
Both have
\[
 \frac{W_{11}W_{00}}{W_{10}W_{01}}=1,
\]
so both Creutz responses vanish.  All Bernoulli parameters lie strictly
between zero and one, so both laws have full support.  The coefficients of
\(\mathcal O_{11}\) in \eqref{eq:YMA30-macro-source} are respectively
\(-2\) and \(-3\); hence the source functions differ.
\end{proof}

\begin{assessment}[The first law-alignment packet]
\label{ass:YMA30-law-alignment-packet}
For the inherited selected root, the four-loop scalar has already been
matched.  The first finite loop-level law-alignment packet is therefore the
three-vector \((A,B,C)\), or a target-native source-distance Read derived from
\eqref{eq:YMA30-exact-source-shape-mismatch}.  Equality of the selected scalar
cannot replace that packet, and the packet itself still does not determine
the absolute Wilson Hessian.
\end{assessment}

% =============================================================================
\section{Reconciliation with the projected character-tension coordinate}
\label{sec:YMA30-character-reconciliation}
% =============================================================================

Version~V28 constructed, on the projected homogeneous fundamental branch, an
elementary character-tension coordinate satisfying
\begin{equation}
 \mathcal U_M=4\mathcal U_{2M}
 \label{eq:YMA30-U-coarse-fine}
\end{equation}
when \(M\) is the coarse lattice and \(2M\) the fine lattice.

\begin{theorem}[Elementary tension scaling and invariant macro-response]
\label{thm:YMA30-character-macro-invariance}
Let the coarse macrocell have lattice increment \(d_M\), the fine macrocell
increment \(d_{2M}=2d_M\), and suppose the homogeneous elementary response
obeys \eqref{eq:YMA30-U-coarse-fine}.  Then
\begin{equation}
 \boxed{
 d_M^2\mathcal U_M
 =(2d_M)^2\mathcal U_{2M}.}
 \label{eq:YMA30-U-macro-invariance}
\end{equation}
Moreover,
\begin{equation}
 \boxed{
 \frac{\mathcal U_M}{a_M^2}
 =\frac{\mathcal U_{2M}}{a_{2M}^2}.}
 \label{eq:YMA30-U-density-invariance}
\end{equation}
Thus the projected rule \(\mathcal U_M=4\mathcal U_{2M}\) is exactly
compatible with one fixed macro-response at fixed physical geometry.  The
factor four belongs to the elementary response per lattice cell, not to the
total inherited macrocell target.
\end{theorem}

\begin{proof}
Substitute \(d_{2M}=2d_M\) into the right side and use
\(\mathcal U_M=4\mathcal U_{2M}\).  Since
\(a_{2M}=a_M/2\), the density identity is the same algebra.
\end{proof}

\begin{corollary}[Fixed response is a vanishing elementary lattice tension]
\label{cor:YMA30-vanishing-elementary-tension}
On the inherited matched branch define
\begin{equation}
 \tau_M^*:=\frac{u_*}{d_M^2}.
 \label{eq:YMA30-elementary-tension}
\end{equation}
Then
\begin{equation}
 \boxed{
 \tau_M^*=\frac{64u_*}{M^2},
 \qquad
 \tau_{2M}^*=\frac14\tau_M^*,
 \qquad
 \frac{\tau_M^*}{a_M^2}=\frac{64u_*}{L^2}.}
 \label{eq:YMA30-tension-scaling}
\end{equation}
In particular,
\(
 \tau_M^*\to0
\)
along every cofinal fixed-response family.  Hence such a family must leave
any regime which carries a cutoff-uniform positive lower bound on the
elementary lattice tension.  This is a critical-window or trajectory
statement, not an area-calibration failure.
\end{corollary}

\begin{proof}
Use \(d_M=M/8\), \(a_M=L/M\), and the definition.  The final alternative is
immediate: a sequence tending to zero cannot remain above a fixed positive
floor.
\end{proof}

\begin{firewall}[Fixed-card strong coupling does not locate the cofinal root]
\label{fire:YMA30-strong-coupling-scope}
The fixed-\(M\) expansion
\(
 \chi_M(\beta)=d_M^2\log(1/\beta)+b_M+O_M(\beta)
\)
remains exact at each finite card.  It does not give a cofinal inverse chart
without estimates uniform in \(M\), and the finite-tower theorem explicitly
does not assert that one fixed \(u_*\) lies in every strong-response chart.
The correct cofinal obstruction is possible loss of the fixed-response root
or escape from the strong-coupling chart, not divergence of
\(u_*/\Delta^2\).
\end{firewall}

% =============================================================================
\section{Explicit qualification of Version V29}
\label{sec:YMA30-V29-qualification}
% =============================================================================

\begin{theorem}[Domain-corrected status of the Version V29 statements]
\label{thm:YMA30-V29-qualification}
The following qualifications are exact.
\begin{enumerate}[label=\textup{(Q30.\arabic*)},leftmargin=5.6em]
\item The unit-cell telescope
      \eqref{eq:YMA29-b-telescope}, density average
      \eqref{eq:YMA29-density-average}, four-child Pythagoras, continuum
      unit-cell formula, and perimeter cancellation remain valid as written.
\item The target rule
      \(
       u_{\rm out}=4u_{\rm in}
      \)
      in \eqref{eq:YMA29-homogeneous-target} applies when the two \(u\)'s are
      elementary one-step Creutz cells.  It does not compare the inherited
      coarse macrocell \(d_M\times d_M\) with the fine macrocell
      \(2d_M\times2d_M\).
\item The statements
      \eqref{eq:YMA29-u-density-diverges} and
      \eqref{eq:YMA29-small-u-target} remain true for the quantity
      \(u_r/a_r^2\).  That quantity is not the physical density of the
      inherited Branch--R cell, whose denominator is
      \((d_Ma_M)^2=\Delta^2\).
\item The applications of
      \Cref{cor:YMA29-fixed-u-diverges,cor:YMA29-finite-density-small-target}
      to the inherited fixed-physical response family are withdrawn and
      replaced by \eqref{eq:YMA30-fixed-u-density} and
      \eqref{eq:YMA30-tension-scaling}.
\item The physical-rectangle branch of
      \Cref{op:YMA29-V30} is already selected by the inherited data.  Its root
      equation is
      \begin{equation}
       \boxed{\chi_M(\beta_M^*)=u_*,}
       \label{eq:YMA30-correct-root-equation}
      \end{equation}
      not \(\chi_M=a_M^2\kappa_*+o(a_M^2)\).
\item The four-child panel is an optional internal-resolution or homogeneous-
      scalarization test.  It is not a prerequisite for the parent response
      match \eqref{eq:YMA30-natural-trace-same-response}.
\item The fixed-loop strong-coupling coefficients of Versions~V25--V28 remain
      valid local blocked-action calibrations.  They cannot be substituted
      directly into the cofinally growing fixed-physical rectangles without a
      uniform surface resummation.
\end{enumerate}
\end{theorem}

\begin{proof}
Item~\textup{(Q30.1)} is the stated domain of Version~V29.
Items~\textup{(Q30.2)--(Q30.3)} follow from
\Cref{thm:YMA30-exact-macro-block,thm:YMA30-fixed-physical-cell}.
Item~\textup{(Q30.4)} is their direct specialization to the inherited counts.
Item~\textup{(Q30.5)} is the already frozen response definition.
Item~\textup{(Q30.6)} is \Cref{thm:YMA30-quadrant-decomposition}.
Item~\textup{(Q30.7)} retains the nonuniform growing-loop warning already
proved in Version~V29.
\end{proof}

\begin{assessment}[Corrected upstream stopping boundary]
\label{ass:YMA30-corrected-stop}
The rectangle geometry and dimensionless response calibration are no longer
open.  The genuinely open upstream rows are:
\[
 \boxed{
 \begin{gathered}
 \text{cofinal existence of the fixed-}u_*\text{ roots}\\
 \big|\quad\text{natural/trace law alignment}\\
 \big|\quad\text{the Wilson-specific nonlinear mode response}
 \end{gathered}}
\]
No further rectangle convention, area-density redefinition, or unit-cell
coefficient ladder advances those rows.
\end{assessment}

% =============================================================================
\section{The first nonlinear selected-surface alias Read}
\label{sec:YMA30-alias-read}
% =============================================================================

The related fixed-scale colour packet already evaluates the complete tangent
alias trace.  For a nonzero axial momentum \(K\), put
\begin{equation}
 s(K)=\sin^2(K/2),
 \qquad
 \alpha(K)=\frac{8+s(K)}{8(2+s(K))}.
 \label{eq:YMA30-alpha-general}
\end{equation}
Let \(c_c,c_f\) be the coarse and fine tangent curvature coefficients and let
\(r^{\rm obs}(K)\) be the scalar law-aligned observed run.  Define the
nonlinear residual by
\begin{equation}
 \boxed{
 \varepsilon_{\rm nl}(K)
 :=r^{\rm obs}(K)-\bigl(c_c-\alpha(K)c_f\bigr).}
 \label{eq:YMA30-nonlinear-residual}
\end{equation}
This includes the nonconstant law incidence, perimeter-resummed non-Abelian
loop and BCH terms, same-cutoff Feshbach return, and non-Gaussian fine
response already identified in the fixed-scale packet.

\begin{theorem}[Two-mode contrast eliminates the coarse tangent coefficient]
\label{thm:YMA30-two-mode-contrast}
For any two retained nonzero axial momenta \(K_1,K_2\),
\begin{equation}
 \boxed{
 \begin{aligned}
 \varepsilon_{\rm nl}(K_2)-\varepsilon_{\rm nl}(K_1)
 ={}&r^{\rm obs}(K_2)-r^{\rm obs}(K_1)\\
 &-\bigl(\alpha(K_1)-\alpha(K_2)\bigr)c_f.
 \end{aligned}}
 \label{eq:YMA30-general-nonlinear-contrast}
\end{equation}
The coarse coefficient \(c_c\) cancels exactly.

On the literal low and held-out alias cards,
\begin{equation}
 \alpha_{\rm low}=\frac{169+12\sqrt2}{392},
 \qquad
 \alpha_{\rm hold}=\frac{17}{40},
 \label{eq:YMA30-literal-alphas}
\end{equation}
so
\begin{equation}
 \boxed{
 \delta_\alpha
 :=\alpha_{\rm low}-\alpha_{\rm hold}
 =\frac{3+15\sqrt2}{490}
 =0.049414700888972\ldots>0,}
 \label{eq:YMA30-delta-alpha}
\end{equation}
and
\begin{equation}
 \boxed{
 \Delta_{\rm nl}^{\rm hold/low}
 :=\varepsilon_{\rm nl}(K_{\rm hold})
   -\varepsilon_{\rm nl}(K_{\rm low})
 =r_{\rm hold}^{\rm obs}-r_{\rm low}^{\rm obs}
  -\delta_\alpha c_f.}
 \label{eq:YMA30-literal-nonlinear-contrast}
\end{equation}
A nonzero contrast forces
\begin{equation}
 \boxed{
 \max\left\{
 |\varepsilon_{\rm nl}(K_{\rm low})|,
 |\varepsilon_{\rm nl}(K_{\rm hold})|
 \right\}
 \ge\frac12|\Delta_{\rm nl}^{\rm hold/low}|.}
 \label{eq:YMA30-one-residual-lower}
\end{equation}
\end{theorem}

\begin{proof}
Subtract \eqref{eq:YMA30-nonlinear-residual} at \(K_1\) from the same identity
at \(K_2\).  The two copies of \(c_c\) cancel.  Substitution of the two exact
alias factors gives
\[
 \frac{169+12\sqrt2}{392}-\frac{17}{40}
 =\frac{3+15\sqrt2}{490}.
\]
Finally,
\(
 |x-y|\le|x|+|y|
\)
implies \(\max\{|x|,|y|\}\ge|x-y|/2\).
\end{proof}

\begin{corollary}[Outward interval for the nonlinear alias contrast]
\label{cor:YMA30-contrast-interval}
Suppose validated scalar intervals are supplied as
\begin{equation}
 |r_{\rm low}^{\rm obs}-\widehat r_{\rm low}|
 \le\eta_{\rm low},
 \qquad
 |r_{\rm hold}^{\rm obs}-\widehat r_{\rm hold}|
 \le\eta_{\rm hold},
 \qquad
 |c_f-\widehat c_f|\le\eta_f.
 \label{eq:YMA30-input-radii}
\end{equation}
Put
\begin{equation}
 \widehat\Delta_{\rm nl}
 =\widehat r_{\rm hold}-\widehat r_{\rm low}
  -\delta_\alpha\widehat c_f,
 \qquad
 \boxed{
 \eta_{\rm nl}
 =\eta_{\rm hold}+\eta_{\rm low}+\delta_\alpha\eta_f.}
 \label{eq:YMA30-contrast-radius}
\end{equation}
Then
\begin{equation}
 \boxed{
 |\Delta_{\rm nl}^{\rm hold/low}-\widehat\Delta_{\rm nl}|
 \le\eta_{\rm nl}.}
 \label{eq:YMA30-contrast-interval}
\end{equation}
Hence:
\begin{enumerate}[label=\textup{(A30.\arabic*)},leftmargin=5.2em]
\item if \(\widehat\Delta_{\rm nl}-\eta_{\rm nl}>0\), the held-out residual
      exceeds the low residual by a strictly positive certified amount;
\item if \(\widehat\Delta_{\rm nl}+\eta_{\rm nl}<0\), the reverse strict
      inequality holds;
\item otherwise the nonlinear alias contrast is \(\stopstatus\).
\end{enumerate}
A nonzero contrast is a Wilson-specific nonlinear acquisition.  It is not by
itself a negative Hessian mode.
\end{corollary}

\begin{proof}
Subtract the centers in \eqref{eq:YMA30-literal-nonlinear-contrast} and apply
the triangle inequality.  Positivity of \(\delta_\alpha\) gives the displayed
radius.  The three interval branches are immediate.
\end{proof}

\begin{assessment}[Why this is the next target-native Read]
\label{ass:YMA30-next-read-meaning}
The tangent alias trace and both literal factors are already exact.  The
three-scalar packet
\[
 \boxed{
 \bigl(r_{\rm low}^{\rm obs},r_{\rm hold}^{\rm obs},c_f\bigr)}
\]
therefore isolates the first momentum-dependent nonlinear law-alignment
quantity without reconstructing a complete boundary process, phase simplex,
or Witten-current atlas.  If the contrast vanishes, an absolute residual or
the complete Hessian is still required; vanishing of one contrast does not
prove that both residuals vanish.
\end{assessment}

% =============================================================================
\section{Integrated Version V30 theorem and revised acquisition}
\label{sec:YMA30-integrated}
% =============================================================================

\begin{theorem}[Fixed-physical macrocell and first nonlinear-read theorem]
\label{thm:YMA30-integrated}
Preserving Versions~V1--V29, Version~V30 proves the following statements at
their displayed scopes.
\begin{enumerate}[label=\textup{(V30.\arabic*)},leftmargin=5.7em]
\item The inherited response is a \(d_M\times d_M\) macrocell with
      \(d_M=M/8\), not a one-lattice-step Creutz cell.
\item Its physical side increment is exactly \(L/8\), and a fixed response
      \(u_*\) has the cutoff-independent physical mixed density
      \(64u_*/L^2\).
\item Arbitrary-step Creutz macrocells telescope into unit cells and have an
      exact fixed-physical continuum cell-average interpretation.
\item Exact dyadic blocking transports the parent macrocell directly to the
      doubled fine macrocell; the four half-size children are an internal
      subdivision, not four competing selected roots.
\item The complete macro-Creutz source transports pointwise from the fine law
      to its induced coarse trace law under the deterministic block map.
\item The natural coarse and trace coarse selected responses are already both
      \(u_*\).  Their four log-loop tables may nevertheless differ by exactly
      three affine shape coordinates, which can change the response source.
\item The projected character rule \(\mathcal U_M=4\mathcal U_{2M}\) is
      consistent with invariant macro-response because the fine macrocell
      contains four times as many elementary response cells.
\item A cofinal fixed-response family has elementary lattice tension
      \(64u_*/M^2\to0\); the open issue is root/critical-window transport, not
      divergence of the fixed physical response density.
\item The inherited-branch applications of the Version~V29 unit-cell density
      alternatives are explicitly corrected by
      \Cref{thm:YMA30-V29-qualification}.
\item Two observed alias runs and the fine tangent coefficient determine the
      exact nonlinear residual contrast
      \eqref{eq:YMA30-literal-nonlinear-contrast}, with a terminating outward
      interval.
\end{enumerate}
No item supplies the numerical fixed-response roots, the complete absolute
nonlinear residual, the signed Wilson Hessian, or its differentiated tail.
\end{theorem}

\begin{proof}
Items~\textup{(V30.1)--(V30.3)} are
\Cref{thm:YMA30-macro-telescope,thm:YMA30-continuum-macrocell,thm:YMA30-fixed-physical-cell}.
Item~\textup{(V30.4)} is
\Cref{thm:YMA30-exact-macro-block,thm:YMA30-quadrant-decomposition}.
Items~\textup{(V30.5)--(V30.6)} are
\Cref{thm:YMA30-source-incidence,cor:YMA30-common-selected-response,thm:YMA30-affine-shape-residual,thm:YMA30-source-shape-mismatch}.
Items~\textup{(V30.7)--(V30.8)} are
\Cref{thm:YMA30-character-macro-invariance,cor:YMA30-vanishing-elementary-tension}.
Item~\textup{(V30.9)} is \Cref{thm:YMA30-V29-qualification}.
Item~\textup{(V30.10)} is
\Cref{thm:YMA30-two-mode-contrast,cor:YMA30-contrast-interval}.
\end{proof}

\begingroup\small
\begin{longtable}{>{\raggedright\arraybackslash}p{0.27\textwidth}
                  >{\raggedright\arraybackslash}p{0.18\textwidth}
                  >{\raggedright\arraybackslash}p{0.47\textwidth}}
\caption{Version V30 proof-status ledger.}\label{tab:YMA30-ledger}\\
\toprule
\textbf{Row}&\textbf{Status}&\textbf{Exact advance or limitation}\\
\midrule
\endfirsthead
\toprule
\textbf{Row}&\textbf{Status}&\textbf{Exact advance or limitation}\\
\midrule
\endhead
V1--V29 prefix
& \pass
& Complete preterminal source preserved byte for byte.\\[0.35em]
Inherited physical rectangle passport
& \pass
& Already fixed by \(R=T=L/4\), \(\Delta=L/8\),
  \(m=n=M/4\), and \(d=M/8\).\\[0.35em]
General macrocell calculus
& \pass
& Exact arbitrary-step telescope, physical density, and continuum cell average.\\[0.35em]
Dyadic parent transport
& \pass
& Coarse macrocell equals the doubled fine macrocell exactly under the trace law.\\[0.35em]
Four-child panel
& Retyped
& Optional internal-resolution/homogeneity card, not a premise of the selected
  parent root.\\[0.35em]
Fixed response physical density
& \pass
& Equal to \(64u_*/L^2\); no \(a^{-2}\) divergence on Branch--R.\\[0.35em]
Elementary lattice tension
& \pass
& Equal to \(64u_*/M^2\) and therefore vanishes cofinally.\\[0.35em]
Fine-to-trace Creutz source
& \pass
& Exact pointwise deterministic-block incidence and exact stock/covariance
  transport.\\[0.35em]
Natural/trace loop-shape residual
& Exact finite reduction
& One matched response leaves three affine log-loop coordinates.  Their values
  remain unpopulated.\\[0.35em]
V29 inherited-branch density conclusion
& Corrected
& Unit-cell statement retained; application to the fixed-physical macrocell
  withdrawn.\\[0.35em]
Strong-coupling cofinal root
& \stopstatus
& Fixed-card inverse charts do not provide a cutoff-uniform fixed-\(u_*\) root
  family.\\[0.35em]
Nonlinear alias contrast
& Exact acquisition formula
& Requires two observed law-aligned runs and the fine tangent coefficient.\\[0.35em]
Complete selected-shell Wilson Hessian
& \stopstatus
& No populated same-action nonlinear run or differentiated Ward tail.\\[0.35em]
Local-vacuum or continuum mass
& Not inferred
& Requires the complete signed head, source ancestry, conductance, tail,
  adjacent-cutoff transport, and reflection-positive continuum.\\
\bottomrule
\end{longtable}
\endgroup

\begin{openproblem}[Version V31: execute the inherited fixed-response card]
\label{op:YMA30-V31}
Preserve the complete present source.  Do not introduce another rectangle
passport, unit-cell density convention, bare strong-coupling coefficient, or
general source compiler before executing the following finite card.
\begin{enumerate}[label=\textup{(V31.\arabic*)},leftmargin=5.7em]
\item Retain the inherited physical quartet
      \((R,T,\Delta)=(L/4,L/4,L/8)\) and supply one numerical response value
      \(u_*\), or a certified interval, chosen before spectral inspection.
\item Isolate the natural coarse and fine root intervals at one adjacent pair
      using the existing logarithm-free Wilson evaluator.  If the fixed
      response is outside either chart, return the root-domain \(\stopstatus\)
      rather than changing the target.
\item Push the fine selected law to the coarse trace law and acquire the three
      affine log-loop defects \((A,B,C)\), or one direct positive/source-distance
      replacement controlling \eqref{eq:YMA30-exact-source-shape-mismatch}.
\item On the same fixed action and law-aligned quotient, populate
      \(r_{\rm low}^{\rm obs}\), \(r_{\rm hold}^{\rm obs}\), and \(c_f\), and
      apply the nonlinear contrast interval
      \eqref{eq:YMA30-contrast-interval}.
\item Acquire one absolute nonlinear residual or directly acquire the complete
      same-action selected momentum--colour Hessian.  Retain the harmonic and
      differentiated Ward-relative tail exactly once.
\item Return \pass, an attained \obstruction, or \stopstatus{} at the first
      unresolved root, law-alignment, nonlinear-run, Hessian, or tail row.
\end{enumerate}
The internal four-child response panel is reopened only if a homogeneous
scalarization, a resolved response-density claim, or a bound on within-cell
variation is explicitly intended.
\end{openproblem}

\begin{assessment}[Version V30 termination]
\label{ass:YMA30-termination}
The upstream circle caused by repeated recalibration is closed.  The inherited
programme had already fixed a legitimate physical rectangle and a
cutoff-independent dimensionless response.  Exact blocking carries that
macrocell and its response source to the trace law.  The first unpopulated
Wilson-specific values are now the natural/trace law-alignment coordinates
and the nonlinear low/held-out alias runs.  No further geometric
redefinition is an advance before those values are read.
\end{assessment}

% =============================================================================
\section{Exact replay and append-only continuation record}
\label{sec:YMA30-replay}
% =============================================================================

The accompanying exact replay verifies:
\begin{enumerate}[label=\textup{(R30.\arabic*)},leftmargin=5.5em]
\item the arbitrary-step mixed-difference telescope on symbolic integer grids;
\item exact dyadic macrocell transport and the internal quadrant decomposition;
\item the inherited identities
      \(m=M/4\), \(d=M/8\), \(da=L/8\), and
      \(\chi/\Delta^2=64\chi/L^2\);
\item the scaling
      \(\tau_{2M}=\tau_M/4\) and invariance of
      \(\tau_M/a_M^2\);
\item compatibility of
      \(\mathcal U_M=4\mathcal U_{2M}\) with invariant macro-response;
\item the exact three-coordinate affine representation of a zero mixed
      log-loop defect;
\item a full-support finite counterexample with equal Creutz response and
      unequal response sources;
\item the exact literal alias difference
      \(\delta_\alpha=(3+15\sqrt2)/490\) and the nonlinear contrast identity.
\end{enumerate}
The replay is a theorem check.  It does not evaluate the Wilson path integral,
a selected root, the natural/trace shape coordinates, a nonlinear run, or a
physical symbol margin.

The complete Version~V29 source preceding its terminal document-closing
command is preserved byte for byte.  Its full-source SHA--256 digest is
\begin{center}\scriptsize\ttfamily
 e2968e076b01a030bd21da93f98bdfc96e9eefe5c49bdae0663a6de5132a3c11.
\end{center}
The accumulated Version~V30 source:
\begin{itemize}[leftmargin=2.2em]
\item preserves the first \texttt{1337199} bytes of Version~V29 exactly;
\item contains \texttt{34960} newline-terminated source lines;
\item contains \texttt{2544} labels, all unique;
\item has exactly one active terminal \verb|\end{document}|;
\item compiles through four \texttt{pdflatex} passes with no warnings,
      unresolved references, duplicate labels, overfull boxes, or underfull
      boxes;
\item produces a \texttt{439}-page PDF, with Version~V30 beginning on page
      \texttt{427};
\item passes PDF parsing, text-layer, metadata, and rendered-page inspection.
\end{itemize}
For Version~V31, preserve the accumulated source, remove only the terminal
active \verb|\end{document}|, append new material immediately before it,
restore one terminal command, and increment the filename to end in
\texttt{\_V31.tex}.

% =============================================================================
% END APPENDED VERSION V30 MATERIAL
% =============================================================================

\end{document}
